\input{preamble}

% OK, start here.
%
\begin{document}

\title{Cohomologie \'etale}


\maketitle

\phantomsection
\label{section-phantom}

\tableofcontents


\section{Introduction}
\label{section-introduction}

\noindent
Ce chapitre est le premier d'une s\'erie de chapitres sur la cohomologie
\'etale des sch\'emas. Nous y pr\'esentons les rudiments de la topologie
\'etale et de la cohomologie des faisceaux ab\'eliens pour cette topologie. Une
grande partie des sujets abord\'es peut sans inconv\'enient \^etre omise lors
d'une premi\`ere lecture; on consultera la section suivante pour d\'ecider des
passages \`a laisser de c\^ot\'e.

\medskip\noindent
La version initiale de ce chapitre a \`a l'origine \'et\'e constitu\'ee par les notes
de la premi\`ere partie d'un cours de cohomologie \'etale donn\'e par Johan de Jong
\`a l'Universit\'e Columbia \`a l'automne 2009. Les premiers auteurs des notes
\'etaient Thibaut Pugin, Zachary Maddock et Min Lee. La seconde partie du
cours se trouve dans le chapitre sur la formule des traces; voir
La formule des traces, section \ref{trace-section-introduction}.




\section{Quelles sections omettre lors d'une premi\`ere lecture?}
\label{section-skip}

\noindent
Nous voulons utiliser le contenu de ce chapitre pour d\'evelopper la th\'eorie
des espaces alg\'ebriques, des champs de Deligne--Mumford, des champs
alg\'ebriques, etc. Nous avons donc ajout\'e \`a l'expos\'e initial de la
cohomologie \'etale des sch\'emas certains d\'eveloppements assez techniques.
Le lecteur les reconna\^itra \`a la fr\'equence du mot ``topos'', aux discussions
relevant de la th\'eorie des ensembles ou encore aux d\'emonstrations portant
sur des propri\'et\'es tr\`es g\'en\'erales des morphismes de sch\'emas. Certaines de
ces discussions peuvent \^etre omises lors d'une premi\`ere lecture.

\medskip\noindent
En particulier, nous sugg\'erons au lecteur d'omettre les sections suivantes:
\begin{enumerate}
\item Comparaison des grands et petits topos,
section \ref{section-compare}.
\item Reconstruction des morphismes,
section \ref{section-morphisms}.
\item Image directe et image inverse,
section \ref{section-monomorphisms}.
\item Propri\'et\'e (A),
section \ref{section-A}.
\item Propri\'et\'e (B),
section \ref{section-B}.
\item Propri\'et\'e (C),
section \ref{section-C}.
\item Invariance topologique du petit site \'etale,
section \ref{section-topological-invariance}.
\item Morphismes entiers universellement injectifs,
section \ref{section-integral-universally-injective}.
\item Grands sites et image directe,
section \ref{section-big}.
\item Exactitude du grand foncteur image directe \`a support propre,
section \ref{section-exactness-lower-shriek}.
\end{enumerate}
Outre ces sections, quelques r\'esultats isol\'es peuvent \^etre omis; le lecteur
les reconna\^itra aux mots-cl\'es indiqu\'es ci-dessus.



%9.08.09
\section{Prologue}
\label{section-prologue}

\noindent
Ces le\c{c}ons portent sur une autre th\'eorie cohomologique. Il faut tout
d'abord observer que la topologie de Zariski n'est pas enti\`erement
satisfaisante. L'une des principales raisons pour lesquelles elle ne donne pas
les r\'esultats souhait\'es est que, si $X$ est une vari\'et\'e complexe et
$\mathcal{F}$ un faisceau constant, alors
$$
H^i(X, \mathcal{F}) = 0, \quad \text{ pour tout } i > 0.
$$
En voici la raison. Dans un sch\'ema irr\'eductible (en particulier dans une
vari\'et\'e), deux ouverts non vides quelconques se rencontrent, de sorte que les
applications de restriction d'un faisceau constant sont surjectives. On dit
que le faisceau est {\it flasque}. Dans ce cas, tous les groupes de cohomologie
de {\v C}ech sup\'erieurs sont nuls, de m\^eme que tous les groupes de
cohomologie de Zariski sup\'erieurs. Autrement dit, la topologie de Zariski ne
poss\`ede ``pas assez'' d'ouverts pour d\'etecter cette cohomologie sup\'erieure.

\medskip\noindent
En revanche, si $X$ est une vari\'et\'e complexe projective lisse, alors
$$
H_{Betti}^{2 \dim X}(X (\mathbf{C}), \Lambda) = \Lambda \quad \text{ pour }
\Lambda = \mathbf{Z}, \ \mathbf{Z}/n\mathbf{Z},
$$
o\`u $X(\mathbf{C})$ d\'esigne l'ensemble des points complexes de $X$. Il serait
souhaitable de retrouver cette propri\'et\'e en g\'eom\'etrie alg\'ebrique, tout
particuli\`erement en caract\'eristique positive.




\section{La topologie \'etale}
\label{section-etale-topology}

\noindent
Il est tr\`es difficile de simplement ``ajouter'' des ouverts suppl\'ementaires
pour raffiner la topologie de Zariski. Une mani\`ere efficace de d\'efinir une
topologie consiste \`a consid\'erer non seulement des ouverts, mais aussi certains
sch\'emas au-dessus de ceux-ci. Pour d\'efinir la topologie \'etale, on consid\`ere
tous les morphismes \'etales $\varphi : U \to X$. Si $X$ est une vari\'et\'e
projective lisse sur $\mathbf{C}$, cela signifie que
\begin{enumerate}
\item $U$ est une r\'eunion disjointe de vari\'et\'es lisses, et
\item $\varphi$ est localement un isomorphisme (analytiquement).
\end{enumerate}
Le mot ``analytiquement'' renvoie \`a la topologie usuelle (transcendante) sur
$\mathbf{C}$. La seconde condition signifie donc que la diff\'erentielle de
$\varphi$ est partout de rang maximal (et, en particulier, que toutes les
composantes de $U$ ont la m\^eme dimension que $X$).

\medskip\noindent
Un rev\^etement double --- d\'efini de mani\`ere informelle comme une application
finie de degr\'e $2$ entre vari\'et\'es --- par exemple
$$
\Spec(\mathbf{C}[t])
\longrightarrow
\Spec(\mathbf{C}[t]),
\quad t \longmapsto t^2
$$
n'est pas un morphisme \'etale s'il poss\`ede une fibre r\'eduite \`a un seul point.
Dans cet exemple, cela se produit pour $t = 0$. Pour qu'une application finie
entre vari\'et\'es sur $\mathbf{C}$ soit \'etale, toutes ses fibres doivent avoir le
m\^eme nombre de points. En retirant le point $t = 0$ de la source de
l'application dans l'exemple, on rend le morphisme \'etale. Mais on peut aussi
retirer d'autres points de la source et le morphisme demeure \'etale. Pour
consid\'erer la topologie \'etale, il faut examiner tous ces morphismes. \`A la
diff\'erence de la topologie de Zariski, ils ne sont pas n\'ecessairement de
simples ouverts de $X$, bien que leurs images le soient toujours.

\begin{definition}
\label{definition-etale-covering-initial}
Une famille de morphismes $\{ \varphi_i : U_i \to X\}_{i \in I}$ est appel\'ee
un {\it recouvrement \'etale} si chaque $\varphi_i$ est un morphisme \'etale et
si leurs images recouvrent $X$, c'est-\`a-dire si
$X = \bigcup_{i \in I} \varphi_i(U_i)$.
\end{definition}

\noindent
Ceci ``d\'efinit'' la topologie \'etale. Autrement dit, nous pouvons maintenant
pr\'eciser ce que sont les faisceaux. Un {\it faisceau \'etale} $\mathcal{F}$
d'ensembles (resp.\ de groupes ab\'eliens, d'espaces vectoriels, etc.) sur $X$
est la donn\'ee de:
\begin{enumerate}
\item pour chaque morphisme \'etale $\varphi : U \to X$, un ensemble
(resp.\ un groupe ab\'elien, un espace vectoriel, etc.) $\mathcal{F}(U)$;
\item pour chaque paire $U, \ U'$ de sch\'emas \'etales sur $X$, et chaque
morphisme $U \to U'$ au-dessus de $X$ (qui est automatiquement \'etale), une
application de restriction
$\rho^{U'}_U : \mathcal{F}(U') \to \mathcal{F}(U)$
\end{enumerate}
Ces donn\'ees doivent satisfaire les conditions $\rho^U_U = \text{id}$ pour le
morphisme identique $U \to U$ et
$\rho^{U'}_U \circ \rho^{U''}_{U'} = \rho^{U''}_U$ pour des morphismes
$U \to U' \to U''$ de sch\'emas \'etales sur $X$, ainsi que l'{\it axiome de
faisceau} suivant:
\begin{itemize}
\item[(*)] pour tout recouvrement \'etale $\{ \varphi_i : U_i \to U\}_{i \in
I}$, le diagramme
$$
\xymatrix{
\mathcal{F} (U) \ar[r] &
\Pi_{i \in I} \mathcal{F} (U_i) \ar@<1ex>[r] \ar@<-1ex>[r] &
\Pi_{i, j \in I} \mathcal{F} (U_i \times_U U_j)
}
$$
est un diagramme d'\'egalisateur dans la cat\'egorie des ensembles
(resp.\ des groupes ab\'eliens, des espaces vectoriels, etc.).
\end{itemize}

\begin{remark}
\label{remark-i-is-j}
Dans le dernier \'enonc\'e, il est essentiel de ne pas oublier le cas $i = j$,
qui constitue en g\'en\'eral une condition loin d'\^etre triviale (contrairement
\`a ce qui se passe pour la topologie de Zariski). De fait, des recouvrements
importants ne comportent souvent qu'un seul \`el\'ement.
\end{remark}

\noindent
Puisque l'identit\'e est un morphisme \'etale, nous pouvons calculer les sections
globales d'un faisceau \'etale, et la cohomologie sera simplement donn\'ee par les
foncteurs d\'eriv\'es \`a droite correspondants. Autrement dit, une fois la th\'eorie
davantage d\'evelopp\'ee et les \'enonc\'es rendus pr\'ecis, rien ne s'opposera \`a la
d\'efinition de la cohomologie.




\section{Prouesses de la topologie \'etale}
\label{section-feats}

\noindent
Pour tout entier naturel $n \in \mathbf{N} = \{1, 2, 3, 4, \ldots\}$, on a
$$
H_\etale^2 (\mathbf{P}^1_\mathbf{C}, \mathbf{Z}/n\mathbf{Z}) =
\mathbf{Z}/n\mathbf{Z}.
$$
Plus g\'en\'eralement, si $X$ est une vari\'et\'e complexe, ses nombres de Betti
\'etales \`a coefficients dans un corps fini co\"incident avec les nombres de
Betti usuels de $X(\mathbf{C})$, c'est-\`a-dire que
$$
\dim_{\mathbf{F}_q} H_\etale^{2i} (X, \mathbf{F}_q) =
\dim_{\mathbf{F}_q} H_{Betti}^{2i} (X(\mathbf{C}), \mathbf{F}_q).
$$
C'est tout \`a fait satisfaisant. Toutefois, ces \'egalit\'es ne valent que pour des
coefficients de torsion, et non en g\'en\'eral. \`A coefficients entiers, on a
$$
H_\etale^2 (\mathbf{P}^1_\mathbf{C}, \mathbf{Z}) = 0.
$$
En revanche,
$H_{Betti}^2(\mathbf{P}^1(\mathbf{C}), \mathbf{Z}) = \mathbf{Z}$, car l'espace
topologique $\mathbf{P}^1(\mathbf{C})$ est hom\'eomorphe \`a une sph\`ere de
dimension $2$. Il existe des moyens de retrouver des coefficients non de
torsion \`a partir de coefficients de torsion par un proc\'ed\'e de passage \`a la
limite, que nous aborderons bient\^ot.




\section{Un calcul}
\label{section-computation}

\noindent
Comment calculer la cohomologie de $\mathbf{P}^1_\mathbf{C}$ \`a coefficients
dans $\Lambda = \mathbf{Z}/n\mathbf{Z}$?
Nous utilisons la cohomologie de {\v C}ech. Un recouvrement de
$\mathbf{P}^1_\mathbf{C}$ est donn\'e par les deux ouverts standards $U_0, U_1$,
tous deux isomorphes \`a $\mathbf{A}^1_\mathbf{C}$, et dont l'intersection est
isomorphe \`a
$\mathbf{A}^1_\mathbf{C} \setminus \{0\} = \mathbf{G}_{m, \mathbf{C}}$.
Il se trouve que la suite de Mayer--Vietoris vaut en cohomologie \'etale. On en
d\'eduit une suite exacte
$$
H_\etale^{i-1}(U_0\cap U_1, \Lambda) \to
H_\etale^i(\mathbf{P}^1_\mathbf{C}, \Lambda) \to
H_\etale^i(U_0, \Lambda) \oplus
H_\etale^i(U_1, \Lambda) \to H_\etale^i(U_0\cap U_1,
\Lambda).
$$
Pour obtenir la r\'eponse attendue, il faudrait montrer que la somme directe
figurant dans le troisi\`eme terme est nulle. En fait, comme pour la topologie
usuelle, on a bien
$$
H_\etale^q (\mathbf{A}^1_\mathbf{C}, \Lambda) = 0
\quad \text{ pour } q \geq 1,
$$
and
$$
H_\etale^q (\mathbf{A}^1_\mathbf{C} \setminus \{0\}, \Lambda) = \left\{
\begin{matrix}
\Lambda & \text{ si }q = 1\text{, et} \\
0 & \text{ pour }q \geq 2.
\end{matrix}
\right.
$$
Ces r\'esultats sont d\'ej\`a assez difficiles (quelle en serait une preuve
\'el\'ementaire?). Expliquons comment effectuer ce calcul une fois que nous
disposerons du formalisme de la cohomologie \'etale.

\medskip\noindent
{\bf Cohomologie sup\'erieure.} Elle est trait\'ee par le fait g\'en\'eral suivant:
si $X$ est une courbe affine sur $\mathbf{C}$, alors
$$
H_\etale^q (X, \mathbf{Z}/n\mathbf{Z}) = 0 \quad \text{ pour } q \geq 2.
$$
On le d\'emontre en consid\'erant le point g\'en\'erique de la courbe et en faisant
un peu de cohomologie galoisienne. Il ne reste donc \`a examiner que la
cohomologie en degr\'e 1.

\medskip\noindent
{\bf Cohomologie en degr\'e 1.} Nous utilisons les identifications suivantes:
\begin{eqnarray*}
H_\etale^1 (X, \mathbf{Z}/n\mathbf{Z}) = \left\{
\begin{matrix}
\text{faisceaux d'ensembles }\mathcal{F}\text{ sur le site \'etale }X_\etale
\text{ munis d'une} \\
\text{action }\mathbf{Z}/n\mathbf{Z} \times \mathcal{F} \to \mathcal{F}
\text{ telle que }\mathcal{F}\text{ soit un }\mathbf{Z}/n\mathbf{Z}\text{-torseur.}
\end{matrix}
\right\}
\Big/ \cong
\\
 = \left\{
\begin{matrix}
\text{morphismes }Y \to X\text{ finis \'etales, munis} \\
\text{ d'une action libre de }\mathbf{Z}/n\mathbf{Z}\text{ tels que }
X = Y/(\mathbf{Z}/n\mathbf{Z}).
\end{matrix}
\right\}
\Big/ \cong.
\end{eqnarray*}
La premi\`ere identification est tr\`es g\'en\'erale (elle vaut pour toute th\'eorie
cohomologique sur un site) et n'a rien \`a voir avec la topologie \'etale. La
seconde identification est une cons\'equence de la th\'eorie de la descente. Le
dernier ensemble d\'ecrit une collection d'objets g\'eom\'etriques que nous pouvons
manipuler concr\`etement.

\medskip\noindent
La courbe $\mathbf{A}^1_\mathbf{C}$ n'a aucun rev\^etement \'etale fini non
trivial; par cons\'equent,
$H_\etale^1 (\mathbf{A}^1_\mathbf{C}, \mathbf{Z}/n\mathbf{Z}) = 0$.
On peut le voir soit topologiquement, soit au moyen de l'argument du paragraphe
suivant.

\medskip\noindent
D\'ecrivons les rev\^etements \'etales finis
$\varphi : Y \to \mathbf{A}^1_\mathbf{C} \setminus \{0\}$.
Il suffit de consid\'erer le cas o\`u $Y$ est connexe, ce que nous supposons. Nous
allons d\'eterminer les possibilit\'es pour $Y$ en appliquant la formule de
Riemann--Hurwitz (c'est bien s\^ur un peu superflu, et le lecteur peut passer la
section suivante s'il le souhaite). Supposons que ce morphisme soit de degr\'e
$n$ et consid\'erons une compactification projective
$$
\xymatrix{
{Y\ } \ar@{^{(}->}[r] \ar[d]^\varphi &
{\bar Y} \ar[d]^{\bar\varphi} \\
{\mathbf{A}^1_\mathbf{C} \setminus \{0\}} \ar@{^{(}->}[r] &
{\mathbf{P}^1_\mathbf{C}}
}
$$
Bien que $\varphi$ soit \'etale et non ramifi\'e, $\bar{\varphi}$ peut \^etre
ramifi\'e en 0 et en $\infty$. Supposons que les ant\'ec\'edents de 0 soient les
points $y_1, \ldots, y_r$, d'indices de ramification $e_1, \ldots e_r$, et que
les ant\'ec\'edents de $\infty$ soient les points $y_1', \ldots, y_s'$, d'indices
de ramification $d_1, \ldots d_s$. En particulier,
$\sum e_i = n = \sum d_j$. La formule de Riemann--Hurwitz donne
$$
2 g_Y - 2 = -2n + \sum (e_i - 1) + \sum (d_j - 1)
$$
et donc $g_Y = 0$, $r = s = 1$ et $e_1 = d_1 = n$.
Ainsi $Y \cong {\mathbf{A}^1_\mathbf{C} \setminus \{0\}}$, et l'on voit
facilement que $\varphi(z) = \lambda z^n$ pour un certain
$\lambda \in \mathbf{C}^*$. Apr\`es avoir reparam\'etr\'e $Y$, nous pouvons
supposer $\lambda = 1$. Notre rev\^etement s'obtient donc en extrayant une
racine $n$-i\`eme de la coordonn\'ee sur
$\mathbf{A}^1_{\mathbf{C}} \setminus \{0\}$.

\medskip\noindent
Rappelons qu'il nous faut classifier les rev\^etements de
${\mathbf{A}^1_\mathbf{C} \setminus \{0\}}$ munis d'une action libre de
$\mathbf{Z}/n\mathbf{Z}$. Dans notre cas, toute action de ce type correspond
\`a un automorphisme de $Y$ envoyant $z$ sur $\zeta_n z$, o\`u $\zeta_n$ est une
racine primitive $n$-i\`eme de l'unit\'e. Il existe $\phi(n)$ telles actions (ici,
$\phi(n)$ d\'esigne l'indicatrice d'Euler). Il existe donc exactement $\phi(n)$
rev\^etements \'etales finis connexes munis d'une action libre donn\'ee de
$\mathbf{Z}/n\mathbf{Z}$, chacun correspondant \`a une racine primitive
$n$-i\`eme de l'unit\'e. Nous laissons au lecteur le soin de v\'erifier que les
rev\^etements \'etales finis non connexes de degr\'e $n$ de
$\mathbf{A}^1_{\mathbf{C}} \setminus \{0\}$, munis d'une action libre donn\'ee
de $\mathbf{Z}/n\mathbf{Z}$, correspondent bijectivement aux racines
$n$-i\`emes non primitives de $1$. Autrement dit, ce calcul montre que
$$
H_\etale^1 (\mathbf{A}^1_\mathbf{C} \setminus \{0\},
\mathbf{Z}/n\mathbf{Z}) =
\Hom(\mu_n(\mathbf{C}), \mathbf{Z}/n\mathbf{Z}) \cong \mathbf{Z}/n\mathbf{Z}.
$$
La premi\`ere identification est canonique, la seconde ne l'est pas; voir la
remarque \ref{remark-normalize-H1-Gm}.
Puisque la preuve de la formule de Riemann--Hurwitz ne fait pas appel au calcul
de la cohomologie, ce qui pr\'ec\`ede constitue effectivement une preuve (\`a
condition de compl\'eter les d\'etails concernant l'annulation, etc.).




\section{Coefficients non de torsion}
\label{section-nontorsion}

\noindent
Pour \'etudier les coefficients non de torsion, on pose la d\'efinition suivante:
$$
H_\etale^i (X, \mathbf{Q}_\ell) :=
\left( \lim_n H_\etale^i(X, \mathbf{Z}/\ell^n\mathbf{Z}) \right)
\otimes_{\mathbf{Z}_\ell} \mathbf{Q}_\ell.
$$
Le symbole $\lim_n$ d\'esigne la {\it limite} du syst\`eme de groupes de
cohomologie $H_\etale^i(X, \mathbf{Z}/\ell^n\mathbf{Z})$ index\'e par $n$; voir
Cat\'egories, section \ref{categories-section-posets-limits}.
Nous devrons donc \'etudier des syst\`emes de faisceaux satisfaisant certaines
conditions de compatibilit\'e.




\section{Th\'eorie des faisceaux}
\label{section-sheaf-theory}
%9.10.09

\noindent
Nous abordons maintenant v\'eritablement les sites et les faisceaux.
Une quantit\'e impressionnante de r\'esultats abstraits utiles pourrait trouver place
dans les quelques sections qui suivent. Une partie de ces r\'esultats est d\'evelopp\'ee dans des
chapitres ant\'erieurs, tels que ceux sur les sites, les modules sur les sites et la cohomologie
sur les sites. Nous nous effor\c{c}ons de ne pas ajouter trop de mati\`ere ici, mais seulement
ce qu'il faut pour que la suite de ce chapitre soit compr\'ehensible.




\section{Pr\'efaisceaux}
\label{section-presheaves}

\noindent
Une r\'ef\'erence pour cette section est
Sites, section \ref{sites-section-presheaves}.

\begin{definition}
\label{definition-presheaf}
Soit $\mathcal{C}$ une cat\'egorie. Un {\it pr\'efaisceau d'ensembles} (respectivement, un
{\it pr\'efaisceau ab\'elien}) sur $\mathcal{C}$ est un foncteur $\mathcal{C}^{opp} \to
\textit{Sets}$ (resp.\ $\textit{Ab}$).
\end{definition}

\noindent
{\bf Terminologie.} Si $U \in \Ob(\mathcal{C})$, les \'el\'ements de
$\mathcal{F}(U)$ sont appel\'es les {\it sections} de $\mathcal{F}$ au-dessus de
$U$. Pour $\varphi : V \to U$ dans $\mathcal{C}$, l'application
$\mathcal{F}(\varphi) : \mathcal{F}(U) \to \mathcal{F}(V)$
est appel\'ee l'{\it application de restriction} et est souvent not\'ee $s \mapsto s|_V$
ou parfois $s \mapsto \varphi^*s$. La notation $s|_V$ est ambigu\"e
car l'application de restriction d\'epend de $\varphi$, mais il s'agit d'un abus
de notation usuel. Nous employons aussi la notation
$\Gamma(U, \mathcal{F}) = \mathcal{F}(U)$.

\medskip\noindent
Dire que $\mathcal{F}$ est un foncteur signifie que, si
$W \to V \to U$ sont des morphismes de $\mathcal{C}$ et si
$s \in \Gamma(U, \mathcal{F})$, alors
$(s|_V)|_W = s |_W$, avec l'abus de
notation que nous venons de voir. De plus, les applications de restriction correspondant aux
morphismes identit\'e $\text{id}_U : U \to U$ sont l'identit\'e.

\medskip\noindent
La cat\'egorie des pr\'efaisceaux d'ensembles (respectivement, des pr\'efaisceaux ab\'eliens) sur
$\mathcal{C}$ est not\'ee $\textit{PSh} (\mathcal{C})$ (resp. $\textit{PAb}
(\mathcal{C})$). C'est la cat\'egorie des foncteurs de $\mathcal{C}^{opp}$ vers
$\textit{Sets}$ (resp. $\textit{Ab}$), c'est-\`a-dire que les morphismes de
pr\'efaisceaux sont les transformations naturelles de foncteurs. Nous ne consid\'erons les
cat\'egories $\textit{PSh}(\mathcal{C})$ et $\textit{PAb}(\mathcal{C})$
que lorsque la cat\'egorie $\mathcal{C}$ est petite. (Notre convention est qu'une cat\'egorie
est petite sauf mention contraire et, si elle ne l'est pas, elle doit \^etre
r\'epertori\'ee dans Cat\'egories, remarque \ref{categories-remark-big-categories}.)

\begin{example}
\label{example-representable-presheaf}
\'Etant donn\'e un objet $X \in \Ob(\mathcal{C})$, consid\'erons la r\`egle
$$
\begin{matrix}
U & \longmapsto & h_X(U) = \Mor_\mathcal{C}(U, X) \\
V \xrightarrow{\varphi} U & \longmapsto &
(\psi \mapsto \psi \circ \varphi) : h_X(U) \to h_X(V).
\end{matrix}
$$
Elle d\'efinit un foncteur $h_X : \mathcal{C}^{opp} \to \textit{Sets}$
et donc un pr\'efaisceau. On l'appelle le
{\it pr\'efaisceau repr\'esentable associ\'e \`a $X$.}
Il n'est pas vrai que les pr\'efaisceaux repr\'esentables soient des faisceaux pour toute topologie sur
tout site.
\end{example}

\begin{lemma}[Yoneda]
\label{lemma-yoneda}
\begin{slogan}
Les morphismes entre objets sont en bijection avec les transformations naturelles
entre les foncteurs qu'ils repr\'esentent.
\end{slogan}
Soit $\mathcal{C}$ une cat\'egorie, et soient $X, Y \in
\Ob(\mathcal{C})$. Il existe une bijection naturelle
$$
\begin{matrix}
\Mor_\mathcal{C}(X, Y) &
\longrightarrow &
\Mor_{\textit{PSh}(\mathcal{C})} (h_X, h_Y) \\
\psi &
\longmapsto &
h_\psi = \psi \circ - : h_X \to h_Y.
\end{matrix}
$$
\end{lemma}

\begin{proof}
Voir
Cat\'egories, lemme \ref{categories-lemma-yoneda}.
\end{proof}




\section{Sites}
\label{section-sites}


\begin{definition}
\label{definition-family-morphisms-fixed-target}
Soit $\mathcal{C}$ une cat\'egorie. Une {\it famille de morphismes \`a cible fix\'ee}
$\mathcal{U} = \{\varphi_i : U_i \to U\}_{i\in I}$ est la donn\'ee de
\begin{enumerate}
\item un objet $U \in \mathcal{C}$,
\item un ensemble $I$ (\'eventuellement vide), et
\item pour tout $i\in I$, un morphisme $\varphi_i : U_i \to U$ de $\mathcal{C}$
de cible $U$.
\end{enumerate}
\end{definition}

\noindent
Il existe une notion de {\it morphisme de familles de morphismes \`a cible
fix\'ee}. Un cas particulier est la notion de {\it raffinement}.
On pourra consulter \`a ce sujet
Sites, section \ref{sites-section-refinements}.

\begin{definition}
\label{definition-site}
Un {\it site}\footnote{Ce que nous appelons un site est appel\'e une cat\'egorie munie d'une
pr\'etopologie dans \cite[Expos\'e II, D\'efinition 1.3]{SGA4}.
Dans \cite{ArtinTopologies}, on l'appelle une cat\'egorie munie d'une topologie de
Grothendieck.} est la donn\'ee d'une cat\'egorie $\mathcal{C}$ et d'un ensemble
$\text{Cov}(\mathcal{C})$ form\'e de familles de morphismes \`a cible fix\'ee
appel\'ees {\it recouvrements}, qui satisfont aux axiomes suivants :
\begin{enumerate}
\item (isomorphisme) si $\varphi : V \to U$ est un isomorphisme dans $\mathcal{C}$,
alors $\{\varphi : V \to U\}$ est un recouvrement,
\item (localit\'e) si $\{\varphi_i : U_i \to U\}_{i\in I}$ est un recouvrement et si
pour tout $i \in I$ on se donne un recouvrement
$\{\psi_{ij} : U_{ij} \to U_i \}_{j\in I_i}$, alors
$$
\{
\varphi_i \circ \psi_{ij} : U_{ij} \to U
\}_{(i, j)\in \prod_{i\in I} \{i\} \times I_i}
$$
est aussi un recouvrement, et
\item (changement de base) si $\{U_i \to U\}_{i\in I}$
est un recouvrement et si $V \to U$ est un morphisme dans $\mathcal{C}$, alors
\begin{enumerate}
\item pour tout $i \in I$, le produit fibr\'e
$U_i \times_U V$ existe dans $\mathcal{C}$, et
\item $\{U_i \times_U V \to V\}_{i\in I}$ est un recouvrement.
\end{enumerate}
\end{enumerate}
\end{definition}

\noindent
Pour nous, la cat\'egorie sous-jacente \`a un site est toujours ``petite'', c.-\`a-d. que sa
collection d'objets forme un ensemble et que la collection des recouvrements d'un
site est aussi un ensemble (comme dans la d\'efinition ci-dessus). Dans ce chapitre,
nous omettrons le plus souvent les arguments qui ram\`enent \`a un ensemble la collection
des objets et des recouvrements. Pour une discussion plus approfondie, voir
Sites, remarque \ref{sites-remark-no-big-sites}.

\begin{example}
\label{example-site-topological-space}
Si $X$ est un espace topologique, il poss\`ede un site associ\'e $X_{Zar}$
d\'efini comme suit : les objets de $X_{Zar}$ sont les ouverts de $X$,
les morphismes entre eux sont les applications d'inclusion, et les recouvrements sont
les recouvrements topologiques usuels (surjectifs). Observons que si
$U, V \subset W \subset X$ sont des ouverts, alors $U \times_W V = U \cap V$
existe : cette cat\'egorie admet les produits fibr\'es. Toutes les v\'erifications sont imm\'ediates et
tout fonctionne comme pr\'evu.
\end{example}




\section{Faisceaux}
\label{section-sheaves}

\begin{definition}
\label{definition-sheaf}
Un pr\'efaisceau $\mathcal{F}$ d'ensembles (resp. un pr\'efaisceau ab\'elien) sur un site
$\mathcal{C}$ est dit {\it s\'epar\'e} si, pour tout recouvrement
$\{\varphi_i : U_i \to U\}_{i\in I} \in \text{Cov} (\mathcal{C})$,
l'application
$$
\mathcal{F}(U) \longrightarrow \prod\nolimits_{i\in I} \mathcal{F}(U_i)
$$
est injective. Ici, l'application est $s \mapsto (s|_{U_i})_{i\in I}$.
Le pr\'efaisceau $\mathcal{F}$ est un {\it faisceau} si, pour tout recouvrement
$\{\varphi_i : U_i \to U\}_{i\in I} \in \text{Cov} (\mathcal{C})$, le
diagramme
\begin{equation}
\label{equation-sheaf-axiom}
\xymatrix{
\mathcal{F}(U) \ar[r] &
\prod_{i\in I} \mathcal{F}(U_i) \ar@<1ex>[r] \ar@<-1ex>[r] &
\prod_{i, j \in I} \mathcal{F}(U_i \times_U U_j),
}
\end{equation}
o\`u la premi\`ere application est $s \mapsto (s|_{U_i})_{i\in I}$ et les deux
applications de droite sont
$(s_i)_{i\in I} \mapsto (s_i |_{U_i \times_U U_j})$ et
$(s_i)_{i\in I} \mapsto (s_j |_{U_i \times_U U_j})$,
est un diagramme \'egalisateur dans la cat\'egorie des ensembles (resp.\ des groupes ab\'eliens).
\end{definition}

\begin{remark}
\label{remark-empty-covering}
Pour le recouvrement vide (o\`u $I = \emptyset$), ceci implique que
$\mathcal{F}(\emptyset)$ est un produit vide, qui est un objet final de la
cat\'egorie correspondante (un singleton, tant pour $\textit{Sets}$ que pour
$\textit{Ab}$).
\end{remark}

\begin{example}
\label{example-sheaf-site-space}
En explicitant cette d\'efinition pour le site $X_{Zar}$ associ\'e \`a un espace
topologique, voir l'exemple \ref{example-site-topological-space}, on retrouve la notion usuelle
de faisceau.
\end{example}

\begin{definition}
\label{definition-category-sheaves}
Nous notons $\Sh(\mathcal{C})$ (resp.\ $\textit{Ab}(\mathcal{C})$)
la sous-cat\'egorie pleine de $\textit{PSh}(\mathcal{C})$
(resp.\ $\textit{PAb}(\mathcal{C})$) dont les objets sont les faisceaux. C'est la
{\it cat\'egorie des faisceaux d'ensembles} (resp.\ celle des {\it faisceaux ab\'eliens}) sur
$\mathcal{C}$.
\end{definition}




\section{L'exemple des G-ensembles}
\label{section-G-sets}

\noindent
Soit $G$ un groupe et d\'efinissons un site $\mathcal{T}_G$ comme suit : la cat\'egorie
sous-jacente est la cat\'egorie des $G$-ensembles, c'est-\`a-dire que ses objets sont des ensembles munis
d'une action \`a gauche de $G$ et que les morphismes sont les applications \'equivariantes ; les
recouvrements de $\mathcal{T}_G$ sont les familles
$\{\varphi_i : U_i \to U\}_{i\in I}$ satisfaisant
$U = \bigcup_{i\in I} \varphi_i(U_i)$.

\medskip\noindent
Il existe un objet particulier du site $\mathcal{T}_G$, \`a savoir le $G$-ensemble $G$
muni de son action naturelle par translations \`a gauche. Nous le notons ${}_G G$.
Observons qu'il existe un isomorphisme naturel de groupes
$$
\begin{matrix}
\rho : & G^{opp} & \longrightarrow & \text{Aut}_{G\textit{-Sets}}({}_G G) \\
& g & \longmapsto & (h \mapsto hg).
\end{matrix}
$$
En particulier, pour tout pr\'efaisceau $\mathcal{F}$, l'ensemble $\mathcal{F}({}_G G)$
h\'erite d'une action de $G$ au moyen de $\rho$. (Notons que, par contravariance de
$\mathcal{F}$, l'ensemble $\mathcal{F}({}_G G)$ est encore un $G$-ensemble \`a gauche.) En fait,
le foncteur
$$
\begin{matrix}
\Sh(\mathcal{T}_G) & \longrightarrow & G\textit{-Sets} \\
\mathcal{F} & \longmapsto & \mathcal{F}({}_G G)
\end{matrix}
$$
est une \'equivalence de cat\'egories. Son quasi-inverse est le foncteur $X \mapsto
h_X$. Sans donner la d\'emonstration compl\`ete (que l'on trouvera dans
Sites, section \ref{sites-section-example-sheaf-G-sets}),
essayons d'expliquer pourquoi il en est ainsi.
\begin{enumerate}
\item
Si $S$ est un $G$-ensemble, nous pouvons le d\'ecomposer en orbites $S = \coprod_{i\in I}
O_i$. L'axiome des faisceaux appliqu\'e au recouvrement $\{O_i \to S\}_{i\in I}$ affirme que
$$
\xymatrix{
\mathcal{F}(S) \ar[r] &
\prod_{i\in I} \mathcal{F}(O_i) \ar@<1ex>[r] \ar@<-1ex>[r] &
\prod_{i, j \in I} \mathcal{F}(O_i \times_S O_j)
}
$$
est un diagramme \'egalisateur. Comme les produits fibr\'es dans $G\textit{-Sets}$
sont induits par les produits fibr\'es dans $\textit{Sets}$, et comme
$\mathcal{F}(\emptyset)$ est un $G$-ensemble r\'eduit \`a un point, nous obtenons
$$
\prod_{i, j \in I} \mathcal{F}(O_i \times_S O_j) = \prod_{i \in I}
\mathcal{F}(O_i)
$$
et les deux applications ci-dessus sont en fait identiques. Ainsi, l'axiome des faisceaux affirme simplement
que $\mathcal{F}(S) = \prod_{i\in I} \mathcal{F}(O_i)$.
\item
Si $S$ est le $G$-ensemble $S= G/H$ et si $\mathcal{F}$ est un faisceau sur $\mathcal{T}_G$,
nous affirmons que
$$
\mathcal{F}(G/H) = \mathcal{F}({}_G G)^H
$$
et en particulier que $\mathcal{F}(\{*\}) = \mathcal{F}({}_G G)^G$. Pour le voir,
appliquons l'axiome des faisceaux au recouvrement $\{ {}_G G \to G/H \}$ de $S$. Nous
avons
\begin{eqnarray*}
{}_G G \times_{G/H} {}_G G & \cong & G \times H \\
(g_1, g_2) & \longmapsto & (g_1, g_1^{-1} g_2)
\end{eqnarray*}
Le membre de gauche est une r\'eunion disjointe de copies de ${}_G G$ (comme $G$-ensemble). L'axiome des faisceaux
s'\'ecrit donc
$$
\xymatrix{
\mathcal{F} (G/H) \ar[r] &
\mathcal{F}({}_G G) \ar@<1ex>[r] \ar@<-1ex>[r] &
\prod_{h\in H} \mathcal{F}({}_G G)
}
$$
o\`u les deux applications de droite sont $s \mapsto (s)_{h \in H}$ et $s \mapsto
(hs)_{h \in H}$. Par cons\'equent, $\mathcal{F}(G/H) = \mathcal{F}({}_G G)^H$, comme
annonc\'e.
\end{enumerate}
Ceci ne d\'emontre pas tout \`a fait l'\'equivalence de cat\'egories annonc\'ee, mais montre au
moins qu'un faisceau $\mathcal{F}$ est enti\`erement d\'etermin\'e par ses sections au-dessus de
${}_G G$. On trouvera les d\'etails (ainsi que des remarques ensemblistes) dans
Sites, section \ref{sites-section-example-sheaf-G-sets}.




\section{Passage au faisceau associ\'e}
\label{section-sheafification}

\begin{definition}
\label{definition-0-cech}
Soit $\mathcal{F}$ un pr\'efaisceau sur le site $\mathcal{C}$ et soit
$\mathcal{U} = \{U_i \to U\} \in \text{Cov} (\mathcal{C})$.
Nous d\'efinissons le {\it groupe de cohomologie de {\v C}ech de degr\'e z\'ero} de
$\mathcal{F}$ relativement \`a $\mathcal{U}$ par
$$
\check H^0 (\mathcal{U}, \mathcal{F}) =
\left\{
(s_i)_{i\in I} \in \prod\nolimits_{i\in I }\mathcal{F}(U_i)
\text{ tels que }
s_i|_{U_i \times_U U_j} = s_j |_{U_i \times_U U_j}
\right\}.
$$
\end{definition}

\noindent
Il existe une application canonique
$\mathcal{F}(U) \to \check H^0 (\mathcal{U}, \mathcal{F})$,
$s \mapsto (s |_{U_i})_{i\in I}$.
Nous appelons {\it morphisme de recouvrements} d'un recouvrement
$\mathcal{V} = \{V_j \to V\}_{j \in J}$ vers $\mathcal{U}$ un triplet
$(\chi, \alpha, \chi_j)$, o\`u
$\chi : V \to U$ est un morphisme,
$\alpha : J \to I$ est une application d'ensembles, et, pour tout
$j \in J$, le morphisme $\chi_j$ s'ins\`ere dans un diagramme commutatif
$$
\xymatrix{
V_j \ar[rr]_{\chi_j} \ar[d] & & U_{\alpha(j)} \ar[d] \\
V \ar[rr]^\chi & & U.
}
$$
\`A partir des donn\'ees $\chi, \alpha, \{\chi_j\}_{j \in J}$, nous d\'efinissons
\begin{eqnarray*}
\check H^0(\mathcal{U}, \mathcal{F}) & \longrightarrow &
\check H^0(\mathcal{V}, \mathcal{F}) \\
(s_i)_{i\in I} & \longmapsto &
\left(\chi_j^*\left(s_{\alpha(j)}\right)\right)_{j\in J}.
\end{eqnarray*}
Nous affirmons alors que
\begin{enumerate}
\item l'application est bien d\'efinie, et
\item elle ne d\'epend que de $\chi$ et est ind\'ependante du choix de
$\alpha, \{\chi_j\}_{j \in J}$.
\end{enumerate}
Nous omettons la d\'emonstration du premier point.
Pour voir (2), consid\'erons un autre triplet $(\psi, \beta, \psi_j)$ tel que
$\chi = \psi$. On a alors le diagramme commutatif
$$
\xymatrix{
V_j \ar[rrr]_{(\chi_j, \psi_j)} \ar[dd] & & &
U_{\alpha(j)} \times_U U_{\beta(j)} \ar[dl] \ar[dr] \\
& & U_{\alpha(j)} \ar[dr] & &
U_{\beta(j)} \ar[dl] \\
V \ar[rrr]^{\chi = \psi} & & & U.
}
$$
\'Etant donn\'ee une section $s \in \mathcal{F}(\mathcal{U})$, son image dans
$\mathcal{F}(V_j)$ par l'application d\'efinie par
$(\chi, \alpha, \{\chi_j\}_{j \in J})$
est $\chi_j^*s_{\alpha(j)}$, et
son image par l'application d\'efinie par $(\psi, \beta, \{\psi_j\}_{j \in J})$
est $\psi_j^*s_{\beta(j)}$. Ces deux \'el\'ements
sont \'egaux car, par hypoth\`ese, $s \in \check H^0(\mathcal{U}, \mathcal{F})$ ;
ils sont donc tous deux \'egaux au tir\'e en arri\`ere de la valeur commune
$$
s_{\alpha(j)}|_{U_{\alpha(j)} \times_U U_{\beta(j)}} =
s_{\beta(j)}|_{U_{\alpha(j)} \times_U U_{\beta(j)}}
$$
par l'application $(\chi_j, \psi_j)$ du diagramme.

\begin{theorem}
\label{theorem-sheafification}
Soient $\mathcal{C}$ un site et $\mathcal{F}$ un pr\'efaisceau sur $\mathcal{C}$.
\begin{enumerate}
\item La r\`egle
$$
U \mapsto \mathcal{F}^+(U) :=
\colim_{\mathcal{U} \text{ recouvrement de }U}
\check H^0(\mathcal{U}, \mathcal{F})
$$
est un pr\'efaisceau. De plus, la limite inductive est filtrante.
\item Il existe un morphisme canonique de pr\'efaisceaux $\mathcal{F} \to \mathcal{F}^+$.
\item Si $\mathcal{F}$ est un pr\'efaisceau s\'epar\'e, alors $\mathcal{F}^+$ est un faisceau
et le morphisme de (2) est injectif.
\item $\mathcal{F}^+$ est un pr\'efaisceau s\'epar\'e.
\item $\mathcal{F}^\# = (\mathcal{F}^+)^+$ est un faisceau, et le morphisme canonique
induit un isomorphisme fonctoriel
$$
\Hom_{\textit{PSh}(\mathcal{C})}(\mathcal{F}, \mathcal{G}) =
\Hom_{\Sh(\mathcal{C})}(\mathcal{F}^\#, \mathcal{G})
$$
pour tout $\mathcal{G} \in \Sh(\mathcal{C})$.
\end{enumerate}
\end{theorem}

\begin{proof}
Voir Sites, th\'eor\`eme \ref{sites-theorem-plus}.
\end{proof}

\noindent
Autrement dit, le passage au faisceau associ\'e, donn\'e par le morphisme naturel
$\mathcal{F} \to \mathcal{F}^\#$, est adjoint \`a gauche au foncteur d'oubli
$\Sh(\mathcal{C}) \to \textit{PSh}(\mathcal{C})$.




\section{Cohomologie}
\label{section-cohomology}

\noindent
Le r\'esultat fondamental suivant permet de d\'efinir la cohomologie
des faisceaux ab\'eliens sur les sites.

\begin{theorem}
\label{theorem-enough-injectives}
La cat\'egorie des faisceaux ab\'eliens sur un site est une cat\'egorie ab\'elienne
qui poss\`ede assez d'injectifs.
\end{theorem}

\begin{proof}
Voir
Modules sur les sites, lemme \ref{sites-modules-lemma-abelian-abelian} et
Injectifs, th\'eor\`eme \ref{injectives-theorem-sheaves-injectives}.
\end{proof}

\noindent
On peut donc d\'efinir la cohomologie au moyen des foncteurs d\'eriv\'es \`a droite du foncteur
des sections : si $U \in \Ob(\mathcal{C})$ et
$\mathcal{F} \in \textit{Ab}(\mathcal{C})$,
$$
H^p(U, \mathcal{F}) :=
R^p\Gamma(U, \mathcal{F}) =
H^p(\Gamma(U, \mathcal{I}^\bullet))
$$
o\`u $\mathcal{F} \to \mathcal{I}^\bullet$ est une r\'esolution injective. Pour ce faire,
il faut v\'erifier que le foncteur $\Gamma(U, -)$ est exact \`a gauche. C'est
le cas, et c'est en partie pourquoi la cat\'egorie $\textit{Ab}(\mathcal{C})$ est ab\'elienne ;
voir
Modules sur les sites, lemme \ref{sites-modules-lemma-abelian-abelian}.
Pour une \'etude plus g\'en\'erale de la cohomologie sur les sites (y compris le
foncteur des sections globales et ses foncteurs d\'eriv\'es \`a droite), voir
Cohomologie sur les sites, section \ref{sites-cohomology-section-cohomology-sheaves}.



\section{La topologie fpqc}
\label{section-fpqc}
%9.15.09

\noindent
Avant d'aborder la cohomologie \'etale, nous \'etudions quelque peu la topologie fpqc, car
elle se comporte bien avec les faisceaux quasi-coh\'erents.

\begin{definition}
\label{definition-fpqc-covering}
Soit $T$ un sch\'ema. Un {\it recouvrement fpqc} de $T$ est une famille
$\{ \varphi_i : T_i \to T\}_{i \in I}$ telle que
\begin{enumerate}
\item chaque $\varphi_i$ est un morphisme plat et
$\bigcup_{i\in I} \varphi_i(T_i) = T$, et
\item pour tout ouvert affine $U \subset T$, il existe un ensemble fini
$K$, une application $\mathbf{i} : K \to I$ et des ouverts affines
$U_{\mathbf{i}(k)} \subset T_{\mathbf{i}(k)}$ tels que
$U = \bigcup_{k \in K} \varphi_{\mathbf{i}(k)}(U_{\mathbf{i}(k)})$.
\end{enumerate}
\end{definition}

\begin{remark}
\label{remark-fpqc}
La premi\`ere condition correspond \`a fp, sigle fran\c{c}ais de
{\it fid\`element plat}, qui signifie ``faithfully flat'' en anglais, et
la seconde \`a qc, {\it quasi-compact}.
La seconde partie de la premi\`ere condition est superflue lorsque la seconde est satisfaite.
\end{remark}

\begin{example}
\label{example-fpqc-coverings}
Voici des exemples de recouvrements fpqc.
\begin{enumerate}
\item Tout recouvrement ouvert de Zariski de $T$ est un recouvrement fpqc.
\item La famille $\{\Spec(B) \to \Spec(A)\}$ est un recouvrement fpqc
si et seulement si $A \to B$ est un morphisme d'anneaux fid\`element plat.
\item Si $f: X \to Y$ est plat, surjectif et quasi-compact, alors $\{ f: X\to
Y\}$ est un recouvrement fpqc.
\item Le morphisme
$\varphi :
\coprod_{x \in \mathbf{A}^1_k} \Spec(\mathcal{O}_{\mathbf{A}^1_k, x})
\to \mathbf{A}^1_k$,
o\`u $k$ est un corps, est plat et surjectif. Il n'est pas quasi-compact, et
en fait la famille $\{\varphi\}$ n'est pas un recouvrement fpqc.
\item \'Ecrivons
$\mathbf{A}^2_k = \Spec(k[x, y])$. Notons $i_x : D(x) \to \mathbf{A}^2_k$
et $i_y : D(y) \to \mathbf{A}^2_k$ les immersions ouvertes standard.
Alors les familles
$\{i_x, i_y, \Spec(k[[x, y]]) \to \mathbf{A}^2_k\}$
et
$\{i_x, i_y, \Spec(\mathcal{O}_{\mathbf{A}^2_k, 0}) \to \mathbf{A}^2_k\}$
sont des recouvrements fpqc.
\end{enumerate}
\end{example}

\begin{lemma}
\label{lemma-site-fpqc}
La collection des recouvrements fpqc de la cat\'egorie des sch\'emas
satisfait les axiomes d'un site.
\end{lemma}

\begin{proof}
Voir Topologies, lemme \ref{topologies-lemma-fpqc}.
\end{proof}

\noindent
Il semble que ce lemme permette de d\'efinir le site fpqc de la cat\'egorie
des sch\'emas. Cependant, un probl\`eme ensembliste surgit lorsqu'on
consid\`ere la topologie fpqc ; voir
Topologies, section \ref{topologies-section-fpqc}.
Il provient de ce que nous exigeons que les sites soient ``petits'', alors qu'aucune
petite cat\'egorie de sch\'emas ne peut contenir un syst\`eme cofinal de recouvrements fpqc d'un
sch\'ema non vide donn\'e. Bien que cela ne nous emp\^eche pas, \`a proprement parler,
de d\'efinir des sites fpqc ``partiels'',
il ne semble pas prudent de le faire. La solution consiste \`a autoriser
la notion de faisceau pour la topologie fpqc (voir ci-dessous), mais \`a interdire de
consid\'erer la cat\'egorie de tous les faisceaux fpqc.

\begin{definition}
\label{definition-sheaf-property-fpqc}
Soit $S$ un sch\'ema. La cat\'egorie des sch\'emas sur $S$ est not\'ee
$\Sch/S$. Consid\'erons un foncteur
$\mathcal{F} : (\Sch/S)^{opp} \to \textit{Sets}$, autrement dit
un pr\'efaisceau d'ensembles. Nous dirons que $\mathcal{F}$
{\it v\'erifie la propri\'et\'e de faisceau pour la topologie fpqc}
si, pour tout recouvrement fpqc $\{U_i \to U\}_{i \in I}$ de sch\'emas sur $S$,
le diagramme (\ref{equation-sheaf-axiom}) est un diagramme \'egalisateur.
\end{definition}

\noindent
Nous dirons de m\^eme que $\mathcal{F}$
{\it v\'erifie la propri\'et\'e de faisceau pour la topologie de Zariski} si, pour
tout recouvrement ouvert $U = \bigcup_{i \in I} U_i$, le diagramme
(\ref{equation-sheaf-axiom}) est un diagramme \'egalisateur. Voir
Sch\'emas, d\'efinition \ref{schemes-definition-representable-by-open-immersions}.
Il est clair que cela revient \`a dire que, pour tout sch\'ema $T$ sur $S$, la
restriction de $\mathcal{F}$ aux ouverts de $T$ est un faisceau (au sens usuel).

\begin{lemma}
\label{lemma-fpqc-sheaves}
Soit $\mathcal{F}$ un pr\'efaisceau sur $\Sch/S$. Alors
$\mathcal{F}$ v\'erifie la propri\'et\'e de faisceau pour la topologie fpqc
si et seulement si
\begin{enumerate}
\item $\mathcal{F}$ v\'erifie la propri\'et\'e de faisceau relativement \`a la
topologie de Zariski, et
\item pour tout morphisme fid\`element plat $\Spec(B) \to \Spec(A)$
de sch\'emas affines sur $S$, l'axiome de faisceau est v\'erifi\'e pour le recouvrement
$\{\Spec(B) \to \Spec(A)\}$. Plus pr\'ecis\'ement, cela signifie que
$$
\xymatrix{
\mathcal{F}(\Spec(A)) \ar[r] &
\mathcal{F}(\Spec(B)) \ar@<1ex>[r] \ar@<-1ex>[r] &
\mathcal{F}(\Spec(B \otimes_A B))
}
$$
est un diagramme \'egalisateur.
\end{enumerate}
\end{lemma}

\begin{proof}
Voir Topologies, lemme \ref{topologies-lemma-sheaf-property-fpqc}.
\end{proof}

\noindent
Une autre mani\`ere de consid\'erer un pr\'efaisceau $\mathcal{F}$ sur
$\Sch/S$ satisfaisant la condition de faisceau pour la
topologie fpqc consiste \`a se donner les donn\'ees suivantes :
\begin{enumerate}
\item pour tout $T/S$, un faisceau usuel (c'est-\`a-dire pour la topologie de Zariski) $\mathcal{F}_T$ sur
$T_{Zar}$,
\item pour tout morphisme $f : T' \to T$ sur $S$, un morphisme de restriction
$f^{-1}\mathcal{F}_T \to \mathcal{F}_{T'}$
\end{enumerate}
satisfaisant
\begin{enumerate}
\item[(a)] les morphismes de restriction sont fonctoriels,
\item[(b)] si $f : T' \to T$ est une immersion ouverte, alors le morphisme de restriction
$f^{-1}\mathcal{F}_T \to \mathcal{F}_{T'}$ est un isomorphisme, et
\item[(c)] pour tout morphisme fid\`element plat
$\Spec(B) \to \Spec(A)$ sur $S$, le diagramme
$$
\xymatrix{
\mathcal{F}_{\Spec(A)}(\Spec(A)) \ar[r] &
\mathcal{F}_{\Spec(B)}(\Spec(B)) \ar@<1ex>[r] \ar@<-1ex>[r] &
\mathcal{F}_{\Spec(B \otimes_A B)}(\Spec(B \otimes_A B))
}
$$
est un diagramme \'egalisateur.
\end{enumerate}
Les donn\'ees (1) et (2), avec les conditions (a), (b), d\'efinissent un pr\'efaisceau
sur $\Sch/S$ satisfaisant la condition de faisceau pour la topologie de Zariski.
D'apr\`es le lemme \ref{lemma-fpqc-sheaves}, la condition (c) suffit alors pour obtenir la
condition de faisceau pour la topologie fpqc.

\begin{example}
\label{example-quasi-coherent}
Consid\'erons le pr\'efaisceau
$$
\begin{matrix}
\mathcal{F} : & (\Sch/S)^{opp} & \longrightarrow & \textit{Ab} \\
& T/S & \longmapsto & \Gamma(T, \Omega_{T/S}).
\end{matrix}
$$
La compatibilit\'e des diff\'erentielles avec la localisation entra\^ine que
$\mathcal{F}$ est un faisceau sur le site de Zariski.
Il ne v\'erifie toutefois pas la condition de faisceau pour la topologie fpqc.
Consid\'erons en effet le cas
$S = \Spec(\mathbf{F}_p)$ et le morphisme
$$
\varphi :
V = \Spec(\mathbf{F}_p[v])
\to
U = \Spec(\mathbf{F}_p[u])
$$
d\'efini par l'envoi de $u$ sur $v^p$. La famille $\{\varphi\}$ est un recouvrement fpqc,
mais le morphisme de restriction
$\mathcal{F}(U) \to \mathcal{F}(V)$
envoie le g\'en\'erateur $\text{d}u$ sur $\text{d}(v^p) = 0$, donc
c'est le morphisme nul, et le diagramme
$$
\xymatrix{
\mathcal{F}(U) \ar[r]^{0} &
\mathcal{F}(V) \ar@<1ex>[r] \ar@<-1ex>[r] &
\mathcal{F}(V \times_U V)
}
$$
n'est pas un diagramme \'egalisateur. Nous verrons plus loin que $\mathcal{F}$ donne en fait
lieu \`a un faisceau sur les sites \'etale et lisse.
\end{example}

\begin{lemma}
\label{lemma-representable-sheaf-fpqc}
Tout pr\'efaisceau repr\'esentable sur $\Sch/S$ satisfait \`a la
condition de faisceau pour la topologie fpqc.
\end{lemma}

\begin{proof}
Voir
Descente, lemme \ref{descent-lemma-fpqc-universal-effective-epimorphisms}.
\end{proof}

\noindent
Nous y reviendrons plus loin, car la d\'emonstration de ce fait utilise
la descente des faisceaux quasi-coh\'erents, que nous \'etudierons dans la
section suivante. Une mani\`ere \'el\'egante d'exprimer le lemme consiste \`a dire que
{\it la topologie fpqc est moins fine que la topologie canonique}, ou
que la topologie fpqc est {\it sous-canonique}. Dans le cadre des sites,
ceci est \'etudi\'e dans
Sites, section \ref{sites-section-representable-sheaves}.

\begin{remark}
\label{remark-fpqc-finest}
La topologie fpqc est plus fine que les topologies de Zariski, \'etale,
lisse, syntomique et fppf. Par cons\'equent, tout pr\'efaisceau
satisfaisant \`a la condition de faisceau pour la topologie fpqc est un
faisceau sur les sites de Zariski, \'etale, lisse, syntomique et fppf.
En particulier, les pr\'efaisceaux repr\'esentables sont des faisceaux sur
le site \'etale d'un sch\'ema,
par exemple.
\end{remark}

\begin{example}
\label{example-additive-group-sheaf}
Soit $S$ un sch\'ema.
Consid\'erons le sch\'ema en groupes additif $\mathbf{G}_{a, S} = \mathbf{A}^1_S$
sur $S$, voir
Groupo\"ides, exemple \ref{groupoids-example-additive-group}.
Le pr\'efaisceau repr\'esentable associ\'e est donn\'e par
$$
h_{\mathbf{G}_{a, S}}(T) =
\Mor_S(T, \mathbf{G}_{a, S}) =
\Gamma(T, \mathcal{O}_T).
$$
D'apr\`es ce qui pr\'ec\`ede, nous savons qu'il s'agit d'un pr\'efaisceau
d'ensembles satisfaisant \`a la condition de faisceau pour la topologie fpqc. Il s'agit manifestement
aussi d'un pr\'efaisceau d'anneaux. Nous pouvons donc le consid\'erer comme le foncteur
$$
\begin{matrix}
\mathcal{O} : &
(\Sch/S)^{opp} &
\longrightarrow &
\textit{Rings} \\
&
T/S &
\longmapsto &
\Gamma(T, \mathcal{O}_T)
\end{matrix}
$$
qui satisfait \`a la condition de faisceau pour la topologie fpqc.
On dispose de m\^eme d'une notion de $\mathcal{O}$-module, et ainsi de
suite.
\end{example}




\section{Descente fid\`element plate}
\label{section-fpqc-descent}

\noindent
Dans cette section, nous \'etudions la descente fid\`element plate des modules
quasi-coh\'erents. Plus pr\'ecis\'ement, nous d\'emontrerons que les modules quasi-coh\'erents
satisfont \`a la descente effective relativement aux recouvrements fpqc.

\begin{definition}
\label{definition-descent-datum}
Soit $\mathcal{U} = \{ t_i : T_i \to T\}_{i \in I}$ une famille de
morphismes de sch\'emas \`a cible fix\'ee. Une {\it donn\'ee de descente} pour les
faisceaux quasi-coh\'erents relativement \`a $\mathcal{U}$ est une famille
$((\mathcal{F}_i)_{i \in I}, (\varphi_{ij})_{i, j \in I})$ telle que
\begin{enumerate}
\item $\mathcal{F}_i$ est un faisceau quasi-coh\'erent sur $T_i$, et
\item $\varphi_{ij} : \text{pr}_0^* \mathcal{F}_i \to
\text{pr}_1^* \mathcal{F}_j$ est un isomorphisme de modules
sur $T_i \times_T T_j$,
\end{enumerate}
et que la {\it condition de cocycle} soit satisfaite : les diagrammes
$$
\xymatrix{
\text{pr}_0^*\mathcal{F}_i \ar[dr]_{\text{pr}_{02}^*\varphi_{ik}}
\ar[rr]^{\text{pr}_{01}^*\varphi_{ij}} & &
\text{pr}_1^*\mathcal{F}_j \ar[dl]^{\text{pr}_{12}^*\varphi_{jk}} \\
& \text{pr}_2^*\mathcal{F}_k
}
$$
commutent sur $T_i \times_T T_j \times_T T_k$.
Cette donn\'ee de descente est dite {\it effective} s'il existe un faisceau
quasi-coh\'erent $\mathcal{F}$ sur $T$ et des isomorphismes de $\mathcal{O}_{T_i}$-modules
$\varphi_i : t_i^* \mathcal{F} \cong \mathcal{F}_i$ compatibles avec
les applications $\varphi_{ij}$, \`a savoir
$$
\varphi_{ij} = \text{pr}_1^* (\varphi_j) \circ \text{pr}_0^* (\varphi_i)^{-1}.
$$
\end{definition}

\noindent
Dans cette section et la suivante, nous \'etudions quelques ingr\'edients de la
d\'emonstration du th\'eor\`eme ci-dessous, ainsi que des r\'esultats connexes.

\begin{theorem}
\label{theorem-descent-quasi-coherent}
Si $\mathcal{V} = \{T_i \to T\}_{i\in I}$ est un recouvrement fpqc, alors toutes
les donn\'ees de descente pour les faisceaux quasi-coh\'erents relativement \`a $\mathcal{V}$
sont effectives.
\end{theorem}

\begin{proof}
Voir
Descente, proposition \ref{descent-proposition-fpqc-descent-quasi-coherent}.
\end{proof}

\noindent
Autrement dit, la cat\'egorie fibr\'ee des faisceaux quasi-coh\'erents est un champ
sur le site fpqc.
La d\'emonstration du th\'eor\`eme comporte deux \'etapes. La premi\`ere consiste \`a observer
que, pour les recouvrements de Zariski, le r\'esultat est facile (ou bien connu), par
le recollement usuel des faisceaux (voir
Faisceaux, section \ref{sheaves-section-glueing-sheaves})
et le caract\`ere local de la quasi-coh\'erence. La seconde \'etape traite le cas d'un
recouvrement fpqc de la forme $\{\Spec(B) \to \Spec(A)\}$,
o\`u $A \to B$ est un homomorphisme d'anneaux fid\`element plat.
Il s'agit d'un lemme d'alg\`ebre, que nous \'enon\c{c}ons maintenant.

\medskip\noindent
{\bf Descente des modules.}
Si $A \to B$ est un homomorphisme d'anneaux, consid\'erons le complexe
$$
(B/A)_\bullet : B \to B \otimes_A B \to B \otimes_A B \otimes_A B \to \ldots
$$
o\`u $B$ est en degr\'e 0, $B \otimes_A B$ en degr\'e 1, etc., les applications \'etant
donn\'ees par
\begin{eqnarray*}
b & \mapsto & 1 \otimes b - b \otimes 1, \\
b_0 \otimes b_1 & \mapsto & 1 \otimes b_0 \otimes b_1 - b_0 \otimes 1 \otimes
b_1 + b_0 \otimes b_1 \otimes 1, \\
& \text{etc.}
\end{eqnarray*}

\begin{lemma}
\label{lemma-algebra-descent}
Si $A \to B$ est fid\`element plat, alors le complexe $(B/A)_\bullet$ est exact en
degr\'es strictement positifs, et $H^0((B/A)_\bullet) = A$.
\end{lemma}

\begin{proof}
Voir Descente, lemme \ref{descent-lemma-ff-exact}.
\end{proof}

\noindent
Grothendieck d\'emontre ceci en trois \'etapes. Premi\`erement, il suppose que l'application $A
\to B$ admet une section, et construit une homotopie explicite avec le complexe dont
$A$ est l'unique terme non nul, en degr\'e 0. Deuxi\`emement, il observe que, pour d\'emontrer
le r\'esultat, il suffit de le faire apr\`es un changement de base fid\`element plat $A \to
A'$, en rempla\c{c}ant $B$ par $B' = B \otimes_A A'$. Troisi\`emement, il effectue le
changement de base fid\`element plat $A \to A' = B$ et remarque que l'application
$A' = B \to B' = B \otimes_A B$ poss\`ede une section naturelle.

\medskip\noindent
La m\^eme strat\'egie d\'emontre le lemme suivant.

\begin{lemma}
\label{lemma-descent-modules}
Si $A \to B$ est fid\`element plat et si $M$ est un $A$-module, alors le
complexe $(B/A)_\bullet \otimes_A M$ est exact en degr\'es strictement positifs, et
$H^0((B/A)_\bullet \otimes_A M) = M$.
\end{lemma}

\begin{proof}
Voir Descente, lemme \ref{descent-lemma-ff-exact}.
\end{proof}

\begin{definition}
\label{definition-descent-datum-modules}
Soient $A \to B$ un homomorphisme d'anneaux et $N$ un $B$-module. Une {\it donn\'ee de descente} pour
$N$ relativement \`a $A \to B$ est un isomorphisme
$\varphi : N \otimes_A B \cong B \otimes_A N$ de $B \otimes_A B$-modules tel
que le diagramme de $B \otimes_A B \otimes_A B$-modules
$$
\xymatrix{
{N \otimes_A B \otimes_A B} \ar[dr]_{\varphi_{02}} \ar[rr]^{\varphi_{01}} & &
{B \otimes_A N \otimes_A B} \ar[dl]^{\varphi_{12}} \\
& {B \otimes_A B \otimes_A N}
}
$$
commute, o\`u $\varphi_{01} = \varphi \otimes \text{id}_B$, et de m\^eme
pour $\varphi_{12}$ et $\varphi_{02}$.
\end{definition}

\noindent
Si $N' = B \otimes_A M$ pour un certain $A$-module M, alors il est muni d'une donn\'ee de descente
canonique donn\'ee par l'application
$$
\begin{matrix}
\varphi_\text{can}: & N' \otimes_A B & \to & B \otimes_A N' \\
& b_0 \otimes m \otimes b_1 & \mapsto & b_0 \otimes b_1 \otimes m.
\end{matrix}
$$

\begin{definition}
\label{definition-effective-modules}
Une donn\'ee de descente $(N, \varphi)$ est dite {\it effective} s'il existe un
$A$-module $M$ tel que $(N, \varphi) \cong (B \otimes_A M,
\varphi_\text{can})$, avec la notion \'evidente d'isomorphisme de donn\'ees de descente.
\end{definition}

\noindent
Le th\'eor\`eme \ref{theorem-descent-quasi-coherent} est une cons\'equence du
r\'esultat suivant.

\begin{theorem}
\label{theorem-descent-modules}
Si $A \to B$ est fid\`element plat, alors les donn\'ees de descente relativement \`a $A\to B$
sont effectives.
\end{theorem}

\begin{proof}
Voir
Descente, proposition \ref{descent-proposition-descent-module}.
Voir aussi
Descente, remarque \ref{descent-remark-homotopy-equivalent-cosimplicial-algebras},
pour un autre point de vue sur la d\'emonstration.
\end{proof}

\begin{remarks}
\label{remarks-theorem-modules-exactness}
Les r\'esultats sur la descente des modules ont plusieurs applications :
\begin{enumerate}
\item L'exactitude du complexe de {\v C}ech en degr\'es strictement positifs pour
le recouvrement $\{\Spec(B) \to \Spec(A)\}$, o\`u $A \to B$ est
fid\`element plat. Ceci donnera des r\'esultats d'annulation en cohomologie.
\item Si $(N, \varphi)$ est une donn\'ee de descente relativement \`a un homomorphisme
fid\`element plat $A \to B$, alors l'$A$-module correspondant est donn\'e par
$$
M = \Ker \left(
\begin{matrix}
N & \longrightarrow & B \otimes_A N \\
n & \longmapsto & 1 \otimes n - \varphi(n \otimes 1)
\end{matrix}
\right).
$$
Voir
Descente, proposition \ref{descent-proposition-descent-module}.
\end{enumerate}
\end{remarks}




%9.17.09
\section{Faisceaux quasi-coh\'erents}
\label{section-quasi-coherent}

\noindent
Nous pouvons appliquer la descente des modules \`a l'\'etude des faisceaux quasi-coh\'erents.

\begin{proposition}
\label{proposition-quasi-coherent-sheaf-fpqc}
Pour tout faisceau quasi-coh\'erent $\mathcal{F}$ sur $S$, le pr\'efaisceau
$$
\begin{matrix}
\mathcal{F}^a : & \Sch/S & \to & \textit{Ab}\\
& (f: T \to S) & \mapsto & \Gamma(T, f^*\mathcal{F})
\end{matrix}
$$
est un $\mathcal{O}$-module qui satisfait \`a la condition de faisceau pour la
topologie fpqc.
\end{proposition}

\begin{proof}
Ceci est d\'emontr\'e dans
Descente, lemme \ref{descent-lemma-sheaf-condition-holds}.
Indiquons ici la d\'emonstration. Comme \'etabli dans le
lemme \ref{lemma-fpqc-sheaves},
il suffit de v\'erifier la propri\'et\'e de faisceau
sur les recouvrements de Zariski et les morphismes fid\`element plats de sch\'emas affines. La
propri\'et\'e de faisceau pour les recouvrements de Zariski rel\`eve de la th\'eorie usuelle des sch\'emas, puisque
$\Gamma(U, i^\ast \mathcal{F}) = \mathcal{F}(U)$ lorsque
$i : U \hookrightarrow S$ est une immersion ouverte.

\medskip\noindent
Pour $\left\{\Spec(B)\to \Spec(A)\right\}$, o\`u $A\to B$ est fid\`element
plat et o\`u
$\mathcal{F}|_{\Spec(A)} = \widetilde{M}$,
cela correspond au fait que
$M = H^0\left((B/A)_\bullet \otimes_A M \right)$, c'est-\`a-dire que
\begin{align*}
0 \to M \to B \otimes_A M \to B \otimes_A B \otimes_A M
\end{align*}
est exacte d'apr\`es le
lemme \ref{lemma-descent-modules}.
\end{proof}

\noindent
Il existe une notion abstraite de faisceau quasi-coh\'erent sur un site annel\'e.
Nous l'introduisons bri\`evement ici. Pour plus d'informations, consulter
Modules sur les sites, section \ref{sites-modules-section-local}.
Soient $\mathcal{C}$ une cat\'egorie et $U$ un objet de $\mathcal{C}$.
Alors $\mathcal{C}/U$ d\'esigne la cat\'egorie des objets au-dessus de $U$, voir
Cat\'egories, exemple \ref{categories-example-category-over-X}.
Si $\mathcal{C}$ est un site, alors $\mathcal{C}/U$ est \'egalement un site :
les recouvrements de $V/U$ sont les familles $\{V_i/U \to V/U\}$ de morphismes
de $\mathcal{C}/U$ \`a cible fix\'ee telles que
$\{V_i \to V\}$ soit un recouvrement de $\mathcal{C}$. De plus, pour tout
faisceau $\mathcal{F}$ sur $\mathcal{C}$, la {\it restriction}
$\mathcal{F}|_{\mathcal{C}/U}$ (d\'efinie de mani\`ere \'evidente)
est encore un faisceau. Voir
Sites, section \ref{sites-section-localize},
pour plus de d\'etails.

\begin{definition}
\label{definition-ringed-site}
Soit $\mathcal{C}$ un {\it site annel\'e}, c'est-\`a-dire un site muni d'un
faisceau d'anneaux $\mathcal{O}$. Un faisceau de $\mathcal{O}$-modules $\mathcal{F}$ sur
$\mathcal{C}$ est dit {\it quasi-coh\'erent} si, pour tout
$U \in \Ob(\mathcal{C})$, il existe un recouvrement
$\{U_i \to U\}_{i\in I}$ de $\mathcal{C}$ tel que la restriction
$\mathcal{F}|_{\mathcal{C}/U_i}$ soit isomorphe au conoyau d'une
application $\mathcal{O}$-lin\'eaire entre $\mathcal{O}$-modules libres
$$
\bigoplus\nolimits_{k \in K} \mathcal{O}|_{\mathcal{C}/U_i}
\longrightarrow
\bigoplus\nolimits_{l \in L} \mathcal{O}|_{\mathcal{C}/U_i}.
$$
La somme directe index\'ee par $K$ est le faisceau associ\'e au pr\'efaisceau
$V \mapsto \bigoplus_{k \in K} \mathcal{O}(V)$, et de m\^eme pour l'autre.
\end{definition}

\noindent
Bien qu'il soit utile de pouvoir donner une d\'efinition g\'en\'erale comme ci-dessus,
cette notion ne se comporte pas bien en g\'en\'eral.

\begin{remark}
\label{remark-final-object}
Dans le cas o\`u $\mathcal{C}$ poss\`ede un objet final, par exemple $S$, il
suffit de v\'erifier la condition de la d\'efinition pour
$U = S$ dans l'\'enonc\'e ci-dessus. Voir
Modules on Sites, Lemma \ref{sites-modules-lemma-local-final-object}.
\end{remark}

\begin{theorem}[M\'eta-th\'eor\`eme sur les faisceaux quasi-coh\'erents]
\label{theorem-quasi-coherent}
Soit $S$ un sch\'ema.
Soit $\mathcal{C}$ un site. Supposons que
\begin{enumerate}
\item la cat\'egorie sous-jacente $\mathcal{C}$ soit une
sous-cat\'egorie pleine de $\Sch/S$,
\item tout recouvrement de Zariski de $T \in \Ob(\mathcal{C})$
puisse \^etre raffin\'e par un recouvrement de $\mathcal{C}$,
\item $S/S$ soit un objet de $\mathcal{C}$,
\item tout recouvrement de $\mathcal{C}$ soit un recouvrement fpqc de sch\'emas.
\end{enumerate}
Alors le pr\'efaisceau $\mathcal{O}$ est un faisceau sur $\mathcal{C}$ et
tout $\mathcal{O}$-module quasi-coh\'erent sur $(\mathcal{C}, \mathcal{O})$
est de la forme $\mathcal{F}^a$ pour un certain faisceau quasi-coh\'erent
$\mathcal{F}$ sur $S$.
\end{theorem}

\begin{proof}
Apr\`es quelques arguments formels, c'est exactement le th\'eor\`eme
\ref{theorem-descent-quasi-coherent}. Nous omettons les d\'etails. Dans
Descent, Proposition \ref{descent-proposition-equivalence-quasi-coherent},
nous d\'emontrons une version plus pr\'ecise du th\'eor\`eme pour les
grands sites de Zariski, fppf, \'etale, lisse et syntomique de $S$,
ainsi que pour les petits sites de Zariski et \'etale de $S$.
\end{proof}

\noindent
Autrement dit, il n'y a aucune diff\'erence entre les modules quasi-coh\'erents
sur le sch\'ema $S$ et les $\mathcal{O}$-modules quasi-coh\'erents
sur les sites $\mathcal{C}$ comme dans le th\'eor\`eme. On trouvera des \'enonc\'es plus pr\'ecis
pour les grands et petits sites $(\Sch/S)_{fppf}$, $S_\etale$, etc.,
dans Descent, Sections
\ref{descent-section-quasi-coherent-sheaves},
\ref{descent-section-quasi-coherent-cohomology}, et
\ref{descent-section-quasi-coherent-sheaves-bis}.
Dans ce chapitre, nous parlerons parfois d'un
``site comme dans le th\'eor\`eme \ref{theorem-quasi-coherent}''
afin d'\'enoncer commod\'ement des r\'esultats valables dans chacune de ces
situations.






\section{Cohomologie de {\v C}ech}
\label{section-cech-cohomology}

\noindent
Notre prochain objectif est d'utiliser la th\'eorie de la descente pour montrer que
$H^i(\mathcal{C}, \mathcal{F}^a) = H_{Zar}^i(S, \mathcal{F})$
pour tout faisceau quasi-coh\'erent $\mathcal{F}$ sur $S$, et
tout site $\mathcal{C}$ comme dans le th\'eor\`eme \ref{theorem-quasi-coherent}.
\`A cette fin, nous introduisons la cohomologie de {\v C}ech sur les sites.
Voir \cite{ArtinTopologies} et
Cohomology on Sites, Sections \ref{sites-cohomology-section-cech},
\ref{sites-cohomology-section-cech-functor}
et \ref{sites-cohomology-section-cech-cohomology-cohomology}
pour plus de d\'etails.

\begin{definition}
\label{definition-cech-complex}
Soit $\mathcal{C}$ une cat\'egorie. Soit
$\mathcal{U} = \{U_i \to U\}_{i \in I}$ une famille de morphismes de
$\mathcal{C}$ \`a cible fix\'ee. Supposons que tous les produits fibr\'es
$U_{i_0} \times_U \ldots \times_U U_{i_p}$ existent dans $\mathcal{C}$.
Soit $\mathcal{F} \in \textit{PAb}(\mathcal{C})$ un
pr\'efaisceau ab\'elien. Nous d\'efinissons le {\it complexe de {\v C}ech}
$\check{\mathcal{C}}^\bullet(\mathcal{U}, \mathcal{F})$ par
$$
\prod_{i_0 \in I} \mathcal{F}(U_{i_0}) \to
\prod_{i_0, i_1 \in I} \mathcal{F}(U_{i_0} \times_U U_{i_1}) \to
\prod_{i_0, i_1, i_2 \in I}
\mathcal{F}(U_{i_0} \times_U U_{i_1} \times_U U_{i_2}) \to \ldots
$$
o\`u le premier terme est en degr\'e 0 et les applications sont usuelles.
Les {\it groupes de cohomologie de {\v C}ech} sont d\'efinis par
$$
\check{H}^p(\mathcal{U}, \mathcal{F}) =
H^p(\check{\mathcal{C}}^\bullet(\mathcal{U}, \mathcal{F})).
$$
\end{definition}

\noindent
Dans la d\'efinition ci-dessus, il est essentiel d'autoriser les indices
$(i_0, \ldots, i_p)$ \`a se r\'ep\'eter.

\begin{lemma}
\label{lemma-cech-presheaves}
Notations et hypoth\`eses comme dans la d\'efinition \ref{definition-cech-complex}.
Le foncteur $\check{\mathcal{C}}^\bullet(\mathcal{U}, -)$
est exact sur la cat\'egorie $\textit{PAb}(\mathcal{C})$.
\end{lemma}

\noindent
Autrement dit, si $0\to \mathcal{F}_1\to \mathcal{F}_2\to \mathcal{F}_3\to 0$
est une suite exacte courte de pr\'efaisceaux de groupes ab\'eliens, alors
$$
0 \to \check{\mathcal{C}}^\bullet\left(\mathcal{U}, \mathcal{F}_1\right)
\to\check{\mathcal{C}}^\bullet(\mathcal{U}, \mathcal{F}_2) \to
\check{\mathcal{C}}^\bullet(\mathcal{U}, \mathcal{F}_3)\to 0
$$
est une suite exacte courte de complexes.

\begin{proof}
Cela r\'esulte imm\'ediatement de la d\'efinition d'une suite exacte courte de
pr\'efaisceaux. En effet, puisque la cat\'egorie des pr\'efaisceaux ab\'eliens est la cat\'egorie des
foncteurs d'une certaine cat\'egorie \`a valeurs dans $\textit{Ab}$, c'est automatiquement une
cat\'egorie ab\'elienne : une suite $\mathcal{F}_1\to \mathcal{F}_2\to \mathcal{F}_3$
est exacte dans $\textit{PAb}$ si et seulement si, pour tout
$U \in \Ob(\mathcal{C})$, la suite
$\mathcal{F}_1(U) \to \mathcal{F}_2(U) \to \mathcal{F}_3(U)$ est exacte dans
$\textit{Ab}$. Ainsi, le complexe ci-dessus est simplement un produit de suites exactes courtes
en chaque degr\'e. Voir aussi
Cohomology on Sites, Lemma \ref{sites-cohomology-lemma-cech-exact-presheaves}.
\end{proof}

\noindent
Cela montre que $\check{H}^\bullet(\mathcal{U}, -)$ est un $\delta$-foncteur.
Nous allons maintenant montrer qu'il s'agit d'un $\delta$-foncteur universel. Il nous faut donc
montrer que c'est un foncteur {\it effa\c{c}able}. Commen\c{c}ons par rappeler le
lemme de Yoneda.

\begin{lemma}[Lemme de Yoneda]
\label{lemma-yoneda-presheaf}
Pour tout pr\'efaisceau $\mathcal{F}$ sur une cat\'egorie $\mathcal{C}$
et tout $U \in \Ob(\mathcal{C})$, il existe un isomorphisme fonctoriel
$$
\Hom_{\textit{PSh}(\mathcal{C})}(h_U, \mathcal{F}) =
\mathcal{F}(U).
$$
\end{lemma}

\begin{proof}
Voir Categories, Lemma \ref{categories-lemma-yoneda}.
\end{proof}

\noindent
Pour un ensemble $E$, nous notons (dans cette section)
$\mathbf{Z}[E]$ le groupe ab\'elien libre engendr\'e par $E$. Sous forme de formule,
$\mathbf{Z}[E] = \bigoplus_{e \in E} \mathbf{Z}$, c'est-\`a-dire que $\mathbf{Z}[E]$ est
un $\mathbf{Z}$-module libre dont une base est form\'ee des \'el\'ements de $E$.
Avec cette notation, nous introduisons le pr\'efaisceau ab\'elien libre engendr\'e par un
pr\'efaisceau d'ensembles.

\begin{definition}
\label{definition-free-abelian-presheaf}
Soit $\mathcal{C}$ une cat\'egorie.
\'Etant donn\'e un pr\'efaisceau d'ensembles $\mathcal{G}$, nous d\'efinissons le
{\it pr\'efaisceau ab\'elien libre engendr\'e par $\mathcal{G}$},
not\'e $\mathbf{Z}_\mathcal{G}$, par la r\`egle
$$
\mathbf{Z}_\mathcal{G}(U)
=
\mathbf{Z}[\mathcal{G}(U)]
$$
pour $U \in \Ob(\mathcal{C})$,
les applications de restriction \'etant induites par celles de $\mathcal{G}$.
Dans le cas particulier o\`u $\mathcal{G} = h_U$, nous \'ecrivons simplement
$\mathbf{Z}_U = \mathbf{Z}_{h_U}$.
\end{definition}

\noindent
Le foncteur $\mathcal{G} \mapsto \mathbf{Z}_\mathcal{G}$ est adjoint \`a gauche du
foncteur d'oubli $\textit{PAb}(\mathcal{C}) \to \textit{PSh}(\mathcal{C})$.
Ainsi, pour tout pr\'efaisceau $\mathcal{F}$, il existe un isomorphisme canonique
$$
\Hom_{\textit{PAb}(\mathcal{C})}(\mathbf{Z}_U, \mathcal{F})
=
\Hom_{\textit{PSh}(\mathcal{C})}(h_U, \mathcal{F})
=
\mathcal{F}(U)
$$
o\`u la derni\`ere \'egalit\'e r\'esulte du lemme de Yoneda. En particulier, nous obtenons le
r\'esultat suivant.

\begin{lemma}
\label{lemma-cech-complex-describe}
Notations et hypoth\`eses comme dans la d\'efinition \ref{definition-cech-complex}.
Le complexe de {\v C}ech $\check{\mathcal{C}}^\bullet(\mathcal{U}, \mathcal{F})$
peut se d\'ecrire explicitement comme suit :
\begin{eqnarray*}
\check{\mathcal{C}}^\bullet(\mathcal{U}, \mathcal{F})
& = &
\left(
\prod_{i_0 \in I}
\Hom_{\textit{PAb}(\mathcal{C})}(\mathbf{Z}_{U_{i_0}}, \mathcal{F}) \to
\prod_{i_0, i_1 \in I}
\Hom_{\textit{PAb}(\mathcal{C})}(
\mathbf{Z}_{U_{i_0} \times_U U_{i_1}}, \mathcal{F}) \to \ldots
\right) \\
& = &
\Hom_{\textit{PAb}(\mathcal{C})}\left(
\left(
\bigoplus_{i_0 \in I} \mathbf{Z}_{U_{i_0}} \leftarrow
\bigoplus_{i_0, i_1 \in I} \mathbf{Z}_{U_{i_0} \times_U U_{i_1}} \leftarrow
\ldots
\right), \mathcal{F}\right)
\end{eqnarray*}
\end{lemma}

\begin{proof}
Cela r\'esulte de la formule ci-dessus. Voir
Cohomology on Sites, Lemma \ref{sites-cohomology-lemma-cech-map-into}.
\end{proof}

\noindent
Nous sommes ainsi ramen\'es \`a l'\'etude du seul complexe situ\'e dans le premier argument du
dernier $\Hom$.

\begin{lemma}
\label{lemma-exact}
Notations et hypoth\`eses comme dans la d\'efinition \ref{definition-cech-complex}.
Le complexe de pr\'efaisceaux ab\'eliens
\begin{align*}
\mathbf{Z}_\mathcal{U}^\bullet \quad : \quad
\bigoplus_{i_0 \in I} \mathbf{Z}_{U_{i_0}} \leftarrow
\bigoplus_{i_0, i_1 \in I} \mathbf{Z}_{U_{i_0} \times_U U_{i_1}} \leftarrow
\bigoplus_{i_0, i_1, i_2 \in I}
\mathbf{Z}_{U_{i_0} \times_U U_{i_1} \times_U U_{i_2}} \leftarrow
\ldots
\end{align*}
est exact en tout degr\'e sauf en degr\'e $0$ dans $\textit{PAb}(\mathcal{C})$.
\end{lemma}

\begin{proof}
Pour tout $V\in \Ob(\mathcal{C})$, le complexe de groupes ab\'eliens
$\mathbf{Z}_\mathcal{U}^\bullet(V)$ est
$$
\begin{matrix}
\mathbf{Z}\left[
\coprod_{i_0\in I} \Mor_\mathcal{C}(V, U_{i_0})\right]
\leftarrow
\mathbf{Z}\left[
\coprod_{i_0, i_1 \in I}
\Mor_\mathcal{C}(V, U_{i_0} \times_U U_{i_1})\right]
\leftarrow \ldots = \\
\bigoplus_{\varphi : V \to U}
\left(
\mathbf{Z}\left[\coprod_{i_0 \in I} \Mor_\varphi(V, U_{i_0})\right]
\leftarrow
\mathbf{Z}\left[\coprod_{i_0, i_1\in I} \Mor_\varphi(V, U_{i_0}) \times
\Mor_\varphi(V, U_{i_1})\right]
\leftarrow
\ldots
\right)
\end{matrix}
$$
o\`u
$$
\Mor_{\varphi}(V, U_i)
=
\{ V \to U_i \text{ tels que } V \to U_i \to U \text{ est \'egal \`a } \varphi \}.
$$
Posons $S_\varphi = \coprod_{i\in I} \Mor_\varphi(V, U_i)$, de sorte que
$$
\mathbf{Z}_\mathcal{U}^\bullet(V)
=
\bigoplus_{\varphi : V \to U}
\left(
\mathbf{Z}[S_\varphi] \leftarrow
\mathbf{Z}[S_\varphi \times S_\varphi] \leftarrow
\mathbf{Z}[S_\varphi \times S_\varphi \times S_\varphi] \leftarrow
\ldots \right).
$$
Il suffit donc de montrer que, pour chaque $S = S_\varphi$, le complexe
\begin{align*}
\mathbf{Z}[S] \leftarrow
\mathbf{Z}[S \times S] \leftarrow
\mathbf{Z}[S \times S \times S] \leftarrow \ldots
\end{align*}
est exact en degr\'es n\'egatifs. Pour le voir, nous pouvons donner une homotopie explicite.
Fixons $s\in S$ et d\'efinissons $K: n_{(s_0, \ldots, s_p)} \mapsto n_{(s, s_0,
\ldots, s_p)}.$ On v\'erifie ais\'ement que $K$ est une homotopie \`a z\'ero pour l'op\'erateur
$$
\delta :
\eta_{(s_0, \ldots, s_p)}
\mapsto
\sum\nolimits_{i = 0}^p (-1)^p \eta_{(s_0, \ldots, \hat s_i, \ldots, s_p)}.
$$
Voir
Cohomology on Sites, Lemma \ref{sites-cohomology-lemma-homology-complex},
pour plus de d\'etails.
\end{proof}

\begin{lemma}
\label{lemma-hom-injective}
Notations et hypoth\`eses comme dans la d\'efinition \ref{definition-cech-complex}.
Si $\mathcal{I}$ est un objet injectif de $\textit{PAb}(\mathcal{C})$,
alors $\check H^p(\mathcal{U}, \mathcal{I}) = 0$ pour tout $p > 0$.
\end{lemma}

\begin{proof}
Le complexe de {\v C}ech s'obtient en appliquant le foncteur
$\Hom_{\textit{PAb}(\mathcal{C})}(-, \mathcal{I}) $ au complexe $
\mathbf{Z}^\bullet_\mathcal{U} $, c'est-\`a-dire
$$
\check H^p(\mathcal{U}, \mathcal{I}) = H^p
(\Hom_{\textit{PAb}(\mathcal{C})} (\mathbf{Z}^\bullet_\mathcal{U},
\mathcal{I})).
$$
Mais nous venons de voir que $\mathbf{Z}^\bullet_\mathcal{U}$ est exact en
degr\'es n\'egatifs, et le foncteur $\Hom_{\textit{PAb}(\mathcal{C})}(-,
\mathcal{I})$ est exact ; par cons\'equent, $\Hom_{\textit{PAb}(\mathcal{C})}
(\mathbf{Z}^\bullet_\mathcal{U}, \mathcal{I})$ est exact en degr\'es positifs.
\end{proof}

\begin{theorem}
\label{theorem-cech-derived}
Notations et hypoth\`eses comme dans la d\'efinition \ref{definition-cech-complex}.
Sur $\textit{PAb}(\mathcal{C})$, les foncteurs $\check{H}^p(\mathcal{U}, -)$ sont
les foncteurs d\'eriv\'es \`a droite de $\check{H}^0(\mathcal{U}, -)$.
\end{theorem}

\begin{proof}
D'apr\`es le lemme \ref{lemma-hom-injective}, les foncteurs
$\check H^p(\mathcal{U}, -)$ sont des
$\delta$-foncteurs universels, puisqu'ils sont effa\c{c}ables.
Il en va de m\^eme des foncteurs d\'eriv\'es \`a droite de $\check H^0(\mathcal{U}, -)$. Puisqu'ils
co\"incident en degr\'e $0$, ils co\"incident par la propri\'et\'e universelle des
$\delta$-foncteurs universels. Pour plus de d\'etails, voir
Cohomology on Sites,
Lemma \ref{sites-cohomology-lemma-cech-cohomology-derived-presheaves}.
\end{proof}

\begin{remark}
\label{remark-presheaves-no-topology}
Notons que tous les \'enonc\'es pr\'ec\'edents portent sur les pr\'efaisceaux ; nous n'avons donc
pas encore utilis\'e la topologie.
\end{remark}




\section{La suite spectrale de {\v C}ech \`a la cohomologie}
\label{section-cech-ss}

\noindent
Cette suite spectrale est essentielle pour \`etablir les r\'esultats fondamentaux sur la
cohomologie des faisceaux.

\begin{lemma}
\label{lemma-forget-injectives}
Le foncteur d'oubli $\textit{Ab}(\mathcal{C})\to \textit{PAb}(\mathcal{C})$
transforme les objets injectifs en objets injectifs.
\end{lemma}

\begin{proof}
C'est formel, puisque le foncteur d'oubli admet un adjoint \`a gauche,
\`a savoir le passage au faisceau associ\'e, qui est un foncteur exact. Pour plus de d\'etails, voir
Cohomology on Sites,
Lemma \ref{sites-cohomology-lemma-injective-abelian-sheaf-injective-presheaf}.
\end{proof}

\begin{theorem}
\label{theorem-cech-ss}
Soit $\mathcal{C}$ un site. Pour tout recouvrement
$\mathcal{U} = \{U_i \to U\}_{i \in I}$ de $U \in \Ob(\mathcal{C})$
et tout faisceau ab\'elien $\mathcal{F}$ sur $\mathcal{C}$,
il existe une suite spectrale
$$
E_2^{p, q}
=
\check H^p(\mathcal{U}, \underline{H}^q(\mathcal{F}))
\Rightarrow
H^{p+q}(U, \mathcal{F}),
$$
o\`u $\underline{H}^q(\mathcal{F})$ est le pr\'efaisceau ab\'elien
$V \mapsto H^q(V, \mathcal{F})$.
\end{theorem}

\begin{proof}
Choisissons une r\'esolution injective $\mathcal{F}\to \mathcal{I}^\bullet$ dans
$\textit{Ab}(\mathcal{C})$ et consid\'erons le bicomplexe
$\check{\mathcal{C}}^\bullet(\mathcal{U}, \mathcal{I}^\bullet)$
et les applications
$$
\xymatrix{
\Gamma(U, I^\bullet) \ar[r] &
\check{\mathcal{C}}^\bullet(\mathcal{U}, \mathcal{I}^\bullet) \\
& \check{\mathcal{C}}^\bullet(\mathcal{U}, \mathcal{F}) \ar[u]
}
$$
Ici, l'application horizontale est l'application naturelle
$\Gamma(U, I^\bullet) \to
\check{\mathcal{C}}^0(\mathcal{U}, \mathcal{I}^\bullet)$
vers la colonne de gauche, et l'application verticale est induite par
$\mathcal{F}\to \mathcal{I}^0$ et aboutit \`a la ligne inf\'erieure.
Par hypoth\`ese, $\mathcal{I}^\bullet$ est un complexe d'injectifs dans
$\textit{Ab}(\mathcal{C})$ ; par le
lemme \ref{lemma-forget-injectives}, c'est donc un complexe d'injectifs dans
$\textit{PAb}(\mathcal{C})$. Ainsi, les lignes du bicomplexe sont
exactes en degr\'es positifs (lemme \ref{lemma-hom-injective}), et
le noyau de $\check{\mathcal{C}}^0(\mathcal{U}, \mathcal{I}^\bullet)
\to \check{\mathcal{C}}^1(\mathcal{U}, \mathcal{I}^\bullet)$
est \'egal \`a
$\Gamma(U, \mathcal{I}^\bullet)$, puisque $\mathcal{I}^\bullet$
est un complexe de faisceaux. En particulier, la cohomologie du complexe total
est la cohomologie usuelle
du foncteur des sections globales $H^0(U, \mathcal{F})$.

\medskip\noindent
Dans la direction verticale, le $q$-i\`eme groupe de cohomologie de la $p$-i\`eme colonne est
$$
\prod_{i_0, \ldots, i_p}
H^q(U_{i_0} \times_U \ldots \times_U U_{i_p}, \mathcal{F})
=
\prod_{i_0, \ldots, i_p}
\underline{H}^q(\mathcal{F})(U_{i_0} \times_U \ldots \times_U U_{i_p})
$$
dans le terme $E_1^{p, q}$. Il s'agit donc d'une suite spectrale standard de bicomplexe,
et sa page $E_2$ est celle indiqu\'ee. Pour plus de d\'etails, voir
Cohomology on Sites,
Lemma \ref{sites-cohomology-lemma-cech-spectral-sequence}.
\end{proof}

\begin{remark}
\label{remark-grothendieck-ss}
C'est une suite spectrale de Grothendieck pour la compos\'ee des foncteurs
$$
\textit{Ab}(\mathcal{C}) \longrightarrow
\textit{PAb}(\mathcal{C}) \xrightarrow{\check H^0} \textit{Ab}.
$$
\end{remark}








\section{Grands et petits sites des sch\'emas}
\label{section-big-small}

\noindent
Soit $S$ un sch\'ema.
Soit $\tau$ l'une des topologies dont nous allons parler.
Ainsi, $\tau \in \{fppf, syntomic, smooth, \etale, Zariski\}$.
Bien entendu, si seule la topologie \'etale vous int\'eresse, vous pouvez
simplement supposer $\tau = \etale$ partout. En outre, nous traiterons
plus en d\'etail des morphismes \'etales, des recouvrements \'etales et des sites \'etales
\`a partir de la section \ref{section-etale-site}.
Pour poursuivre l'\'etude de la cohomologie des
faisceaux quasi-coh\'erents, il est commode d'introduire le
grand $\tau$-site et, lorsque $\tau \in \{\etale, Zariski\}$, le
petit $\tau$-site de $S$. Pour ce faire, nous introduisons d'abord
la notion de $\tau$-recouvrement.

\begin{definition}
\label{definition-tau-covering}
(Voir
Topologies, Definitions
\ref{topologies-definition-fppf-covering},
\ref{topologies-definition-syntomic-covering},
\ref{topologies-definition-smooth-covering},
\ref{topologies-definition-etale-covering}, et
\ref{topologies-definition-zariski-covering}.)
Soit $\tau \in \{fppf, syntomic, smooth, \etale, Zariski\}$.
Une famille de morphismes de sch\'emas $\{f_i : T_i \to T\}_{i \in I}$ \`a cible
fix\'ee est appel\'ee un {\it $\tau$-recouvrement} si et seulement si
chaque $f_i$ est plat et de pr\'esentation finie, syntomique, lisse, \'etale,
respectivement une immersion ouverte, et si $\bigcup f_i(T_i) = T$.
\end{definition}

\noindent
La classe de tous les $\tau$-recouvrements satisfait aux axiomes
(1), (2) et (3) de la
d\'efinition \ref{definition-site} (notre d\'efinition d'un site), voir
Topologies, Lemmas
\ref{topologies-lemma-fppf},
\ref{topologies-lemma-syntomic},
\ref{topologies-lemma-smooth},
\ref{topologies-lemma-etale}, et
\ref{topologies-lemma-zariski}.

\medskip\noindent
Introduisons les sites avec lesquels nous travaillerons.
Contrairement \`a ce qui se passe dans
\cite{SGA4}, nous ne voulons pas choisir un univers. Nous choisissons plut\^ot un ``univers
partiel'' (qui est un ensemble suffisamment grand comme dans
Sets, Section \ref{sets-section-sets-hierarchy}) et consid\'erons tous les sch\'emas
contenus dans cet ensemble. Bien entendu, nous veillons \`a ce que notre sch\'ema de base privil\'egi\'e
$S$ appartienne \`a l'univers partiel. Apr\`es avoir choisi la cat\'egorie sous-jacente,
nous choisissons un ensemble suffisamment grand de $\tau$-recouvrements qui en fait un site.
Les d\'etails se trouvent dans le chapitre sur les topologies des sch\'emas ; les choix
laissent une grande libert\'e, mais, en fin de compte, les choix effectivement faits n'auront
aucune incidence sur la cohomologie \'etale (ou autre) de $S$ (tout comme, dans \cite{SGA4},
le choix effectif de l'univers n'a finalement aucune importance). En outre, la mani\`ere dont le
texte est r\'edig\'e permet au lecteur qui accepte les cardinaux fortement inaccessibles
(c'est-\`a-dire les univers) de les utiliser en remplacement.

\begin{definition}
\label{definition-tau-site}
Soit $S$ un sch\'ema.
Soit $\tau \in \{fppf, syntomic, smooth, \etale, \linebreak[0] Zariski\}$.
\begin{enumerate}
\item Un {\it grand $\tau$-site de $S$} est l'un quelconque des sites
$(\Sch/S)_\tau$ construits comme expliqu\'e ci-dessus et, plus en d\'etail, dans
Topologies, Definitions
\ref{topologies-definition-big-small-fppf},
\ref{topologies-definition-big-small-syntomic},
\ref{topologies-definition-big-small-smooth},
\ref{topologies-definition-big-small-etale}, et
\ref{topologies-definition-big-small-Zariski}.
\item Si $\tau \in \{\etale, Zariski\}$, alors le
{\it petit $\tau$-site de $S$}
est la sous-cat\'egorie pleine $S_\tau$ de $(\Sch/S)_\tau$ dont les objets
sont les sch\'emas $T$ sur $S$ dont le morphisme structural $T \to S$ est \'etale,
respectivement une immersion ouverte. Un recouvrement dans $S_\tau$ est un recouvrement
$\{U_i \to U\}$ dans $(\Sch/S)_\tau$
tel que $U$ soit un objet de $S_\tau$.
\end{enumerate}
\end{definition}

\noindent
La cat\'egorie sous-jacente au site $(\Sch/S)_\tau$ poss\`ede des propri\'et\'es
raisonnables de ``stabilit\'e'' : \'etant donn\'e un sch\'ema $T$ qui lui appartient, tout sous-sch\'ema
localement ferm\'e de $T$ est isomorphe \`a un objet de $(\Sch/S)_\tau$.
Parmi les autres propri\'et\'es de ce type : stabilit\'e par produits fibr\'es de sch\'emas,
par sommes disjointes d\'enombrables,
par passage aux sch\'emas de type fini sur un sch\'ema donn\'e ; \'etant donn\'e un sch\'ema affine
$\Spec(R)$, on peut compl\'eter ou localiser $R$, ou en prendre le quotient
par un id\'eal, tout en restant dans la cat\'egorie, etc.
En revanche, par exemple, des sommes disjointes arbitraires
de sch\'emas de $(\Sch/S)_\tau$ feront sortir de cette cat\'egorie.
Notons aussi que, pour un objet $T$ de $(\Sch/S)_\tau$, il existe des
$\tau$-recouvrements $\{T_i \to T\}_{i \in I}$ (comme dans la
d\'efinition \ref{definition-tau-covering})
qui ne sont pas des recouvrements dans $(\Sch/S)_\tau$, par exemple
parce que les sch\'emas $T_i$ ne sont pas des objets de la cat\'egorie
$(\Sch/S)_\tau$. Mais notre choix des sites $(\Sch/S)_\tau$
est tel qu'il existe toujours
un recouvrement $\{U_j \to T\}_{j \in J}$ de $(\Sch/S)_\tau$ qui raffine
le recouvrement $\{T_i \to T\}_{i \in I}$, voir
Topologies, Lemmas
\ref{topologies-lemma-fppf-induced},
\ref{topologies-lemma-syntomic-induced},
\ref{topologies-lemma-smooth-induced},
\ref{topologies-lemma-etale-induced}, et
\ref{topologies-lemma-zariski-induced}.
Dans ce chapitre, nous ignorerons pour l'essentiel ces questions.

\medskip\noindent
Si $\mathcal{F}$ est un faisceau sur $(\Sch/S)_\tau$ ou $S_\tau$, alors
nous notons
$$
H^p_\tau(U, \mathcal{F}), \text{ en particulier }
H^p_\tau(S, \mathcal{F})
$$
les groupes de cohomologie de $\mathcal{F}$ sur l'objet $U$ du site, voir
la section \ref{section-cohomology}. Nous avons ainsi
$H^p_{fppf}(S, \mathcal{F})$,
$H^p_{syntomic}(S, \mathcal{F})$,
$H^p_{smooth}(S, \mathcal{F})$,
$H^p_\etale(S, \mathcal{F})$, et
$H^p_{Zar}(S, \mathcal{F})$. Les deux derniers sont potentiellement ambigus, puisqu'ils
peuvent renvoyer soit au grand, soit au petit site \'etale ou de Zariski. Cependant,
cette ambigu\"it\'e est sans cons\'equence gr\^ace au lemme suivant.

\begin{lemma}
\label{lemma-compare-cohomology-big-small}
Soit $\tau \in \{\etale, Zariski\}$.
Si $\mathcal{F}$ est un faisceau ab\'elien d\'efini sur
$(\Sch/S)_\tau$, alors
les groupes de cohomologie de $\mathcal{F}$ sur $S$ co\"incident avec les groupes de cohomologie
de $\mathcal{F}|_{S_\tau}$ sur $S$.
\end{lemma}

\begin{proof}
D'apr\`es
Topologies, Lemmas \ref{topologies-lemma-at-the-bottom} et
\ref{topologies-lemma-at-the-bottom-etale},
les foncteurs $S_\tau \to (\Sch/S)_\tau$
satisfont aux hypoth\`eses de
Sites, Lemma \ref{sites-lemma-bigger-site}.
Notre lemme r\'esulte donc de
Cohomology on Sites, Lemma \ref{sites-cohomology-lemma-cohomology-bigger-site}.
\end{proof}

\noindent
La cat\'egorie des faisceaux sur le grand ou le petit site \'etale de $S$ ne d\'epend que
de la sous-cat\'egorie pleine de $(\Sch/S)_\etale$ ou de $S_\etale$ form\'ee des
sch\'emas affines, et il suffit de consid\'erer entre eux les recouvrements \'etales standards
(d\'efinis ci-dessous). Cela donne les sites
$(\textit{Aff}/S)_\etale$ et $S_{affine, \etale}$, voir
Topologies, Definition \ref{topologies-definition-big-small-etale}.
Les r\'esultats de comparaison sont d\'emontr\'es dans Topologies,
Lemmas \ref{topologies-lemma-affine-big-site-etale} et
\ref{topologies-lemma-alternative}. Voici notre d\'efinition des
recouvrements standards pour certaines des topologies consid\'er\'ees dans ce chapitre.

\begin{definition}
\label{definition-standard-tau}
(Voir
Topologies, Definitions
\ref{topologies-definition-standard-fppf},
\ref{topologies-definition-standard-syntomic},
\ref{topologies-definition-standard-smooth},
\ref{topologies-definition-standard-etale}, et
\ref{topologies-definition-standard-Zariski}.)
Soit $\tau \in \{fppf, syntomic, smooth, \etale, Zariski\}$.
Soit $T$ un sch\'ema affine.
Un {\it $\tau$-recouvrement standard} de $T$ est une famille
$\{f_j : U_j \to T\}_{j = 1, \ldots, m}$ o\`u chaque $U_j$ est affine,
et chaque $f_j$ est plat et de pr\'esentation finie,
syntomique standard, lisse standard, \'etale, respectivement l'immersion d'un
ouvert principal standard dans $T$, et $T = \bigcup f_j(U_j)$.
\end{definition}

\begin{lemma}
\label{lemma-tau-affine}
Soit $\tau \in \{fppf, syntomic, smooth, \etale, Zariski\}$.
Tout $\tau$-recouvrement d'un sch\'ema affine peut \^etre raffin\'e par un
$\tau$-recouvrement standard.
\end{lemma}

\begin{proof}
Voir
Topologies, Lemmas
\ref{topologies-lemma-fppf-affine},
\ref{topologies-lemma-syntomic-affine},
\ref{topologies-lemma-smooth-affine},
\ref{topologies-lemma-etale-affine}, et
\ref{topologies-lemma-zariski-affine}.
\end{proof}

\noindent
Par souci d'exhaustivit\'e, nous \'enon\c{c}ons et d\'emontrons l'invariance, par rapport au choix de l'univers
partiel, des groupes de cohomologie que nous consid\'erons. Nous d\'emontrerons l'invariance
du petit topos \'etale dans le
lemme \ref{lemma-etale-topos-independent-partial-universe} ci-dessous.
Pour les notations et la terminologie employ\'ees dans ce lemme, nous renvoyons \`a
Topologies, Section \ref{topologies-section-change-alpha}.

\begin{lemma}
\label{lemma-cohomology-enlarge-partial-universe}
Soit $\tau \in \{fppf, syntomic, smooth, \etale, Zariski\}$.
Soit $S$ un sch\'ema.
Soient $(\Sch/S)_\tau$ et $(\Sch'/S)_\tau$ deux
grands $\tau$-sites de $S$, et supposons que le premier soit contenu dans le second.
Dans ce cas,
\begin{enumerate}
\item pour tout faisceau ab\'elien $\mathcal{F}'$ d\'efini sur $(\Sch'/S)_\tau$ et
tout objet $U$ de $(\Sch/S)_\tau$, nous avons
$$
H^p_\tau(U, \mathcal{F}'|_{(\Sch/S)_\tau}) =
H^p_\tau(U, \mathcal{F}')
$$
Autrement dit : la cohomologie de $\mathcal{F}'$ sur $U$, calcul\'ee dans le site le plus grand,
co\"incide avec celle de la restriction de $\mathcal{F}'$ au site le plus petit,
sur $U$.
\item pour tout faisceau ab\'elien $\mathcal{F}$ sur $(\Sch/S)_\tau$, il existe un
faisceau ab\'elien $\mathcal{F}'$ sur $(\Sch/S)_\tau'$ dont la restriction \`a
$(\Sch/S)_\tau$ est isomorphe \`a $\mathcal{F}$.
\end{enumerate}
\end{lemma}

\begin{proof}
D'apr\`es Topologies, Lemma \ref{topologies-lemma-change-alpha}, le foncteur d'inclusion
$(\Sch/S)_\tau \to (\Sch'/S)_\tau$ satisfait aux hypoth\`eses de
Sites, Lemma \ref{sites-lemma-bigger-site}. Cela implique (2), et (1)
r\'esulte de
Cohomology on Sites, Lemma \ref{sites-cohomology-lemma-cohomology-bigger-site}.
\end{proof}




\section{Le topos \'etale}
\label{section-etale-topos}

\noindent
Un {\it topos} est la cat\'egorie des faisceaux d'ensembles sur un site, voir
Sites, Definition \ref{sites-definition-topos}. Il est donc d'usage
d'employer l'expression ``topos \'etale d'un sch\'ema'' pour d\'esigner
la cat\'egorie des faisceaux sur le petit site \'etale d'un sch\'ema.
En voici la d\'efinition formelle.

\begin{definition}
\label{definition-etale-topos}
Soit $S$ un sch\'ema.
\begin{enumerate}
\item Le {\it topos \'etale}, ou le {\it petit topos \'etale}
de $S$, est la cat\'egorie $\Sh(S_\etale)$ des faisceaux d'ensembles sur
le petit site \'etale de $S$.
\item Le {\it topos de Zariski}, ou le {\it petit topos de Zariski}
de $S$, est la cat\'egorie $\Sh(S_{Zar})$ des faisceaux d'ensembles sur le
petit site de Zariski de $S$.
\item Pour $\tau \in \{fppf, syntomic, smooth, \etale, Zariski\}$, un
{\it grand $\tau$-topos} est la cat\'egorie des faisceaux d'ensembles sur un
grand $\tau$-topos de $S$.
\end{enumerate}
\end{definition}

\noindent
Notons que le petit topos de Zariski de $S$ est simplement la cat\'egorie des faisceaux
d'ensembles sur l'espace topologique sous-jacent \`a $S$, voir
Topologies, Lemma \ref{topologies-lemma-Zariski-usual}.
Alors que le petit topos \'etale ne d\'epend pas des choix faits lors de la
construction du petit site \'etale, les grands topos d\'ependent en g\'en\'eral
de ces choix.

\medskip\noindent
Il se trouve que le grand ou le petit topos \'etale ne d\'epend que de la
sous-cat\'egorie pleine de $(\Sch/S)_\etale$ ou de $S_\etale$ form\'ee des
sch\'emas affines, voir
Topologies, Lemmas \ref{topologies-lemma-affine-big-site-etale} et
\ref{topologies-lemma-alternative}.
Nous l'utiliserons par exemple dans la d\'emonstration du lemme suivant.

\begin{lemma}
\label{lemma-etale-topos-independent-partial-universe}
Soit $S$ un sch\'ema. Le topos \'etale de $S$ est ind\'ependant
(\`a \'equivalence canonique pr\`es) de la construction du petit
site \'etale de la d\'efinition \ref{definition-tau-site}.
\end{lemma}

\begin{proof}
Nous devons montrer qu'\'etant donn\'es deux grands sites \'etales
$\Sch_\etale$ et $\Sch_\etale'$ contenant
$S$, on a $\Sh(S_\etale) \cong \Sh(S_\etale')$
avec les notations \'evidentes. D'apr\`es Topologies, Lemma \ref{topologies-lemma-contained-in},
nous pouvons supposer $\Sch_\etale \subset \Sch_\etale'$.
D'apr\`es Sets, Lemma \ref{sets-lemma-what-is-in-it},
tout sch\'ema affine \'etale sur $S$ est isomorphe \`a un objet
\`a la fois de $\Sch_\etale$ et de $\Sch_\etale'$.
Ainsi, le foncteur induit
$S_{affine, \etale} \to S_{affine, \etale}'$
est une \'equivalence. En outre, il est clair que ce foncteur
et un quasi-inverse transforment les recouvrements \'etales standards en
recouvrements \'etales standards.
Le r\'esultat d\'ecoule donc de
Topologies, Lemma \ref{topologies-lemma-alternative}.
\end{proof}





\section{Cohomologie des faisceaux quasi-coh\'erents}
\label{section-cohomology-quasi-coherent}
%9.22.09

\noindent
Commen\c{c}ons par un lemme simple (qui reste vrai dans une g\'en\'eralit\'e plus grande que
celle de l'\'enonc\'e). Il affirme que le complexe de {\v C}ech d'un recouvrement standard est
\'egal au complexe de {\v C}ech d'un recouvrement fpqc de la forme
$\{\Spec(B) \to \Spec(A)\}$ o\`u $A \to B$ est fid\`element plat.

\begin{lemma}
\label{lemma-cech-complex}
Soit $\tau \in \{fppf, syntomic, smooth, \etale, Zariski\}$.
Soit $S$ un sch\'ema.
Soit $\mathcal{F}$ un faisceau ab\'elien sur $(\Sch/S)_\tau$, ou sur
$S_\tau$ dans le cas o\`u $\tau = \etale$, et soit
$\mathcal{U} = \{U_i \to U\}_{i \in I}$
un $\tau$-recouvrement standard de ce site.
Soit $V = \coprod_{i \in I} U_i$. Alors
\begin{enumerate}
\item $V$ est un sch\'ema affine,
\item $\mathcal{V} = \{V \to U\}$ est un recouvrement fpqc
et aussi un $\tau$-recouvrement, sauf si $\tau = Zariski$,
\item les complexes de {\v C}ech
$\check{\mathcal{C}}^\bullet (\mathcal{U}, \mathcal{F})$ et
$\check{\mathcal{C}}^\bullet (\mathcal{V}, \mathcal{F})$ co\"incident.
\end{enumerate}
\end{lemma}

\begin{proof}
La d\'efinition d'un $\tau$-recouvrement standard est donn\'ee dans
Topologies, Definition
\ref{topologies-definition-standard-Zariski},
\ref{topologies-definition-standard-etale},
\ref{topologies-definition-standard-smooth},
\ref{topologies-definition-standard-syntomic}, et
\ref{topologies-definition-standard-fppf}.
Par d\'efinition, chacun des sch\'emas
$U_i$ est affine et $I$ est un ensemble fini. Ainsi, $V$ est un sch\'ema affine.
Il est clair que $V \to U$ est plat et surjectif ; par cons\'equent,
$\mathcal{V}$ est un recouvrement fpqc, voir
l'exemple \ref{example-fpqc-coverings}.
Sauf dans le cas de Zariski, le recouvrement $\mathcal{V}$
est aussi un $\tau$-recouvrement, voir
Topologies, Definition
\ref{topologies-definition-etale-covering},
\ref{topologies-definition-smooth-covering},
\ref{topologies-definition-syntomic-covering}, et
\ref{topologies-definition-fppf-covering}.

\medskip\noindent
Notons que $\mathcal{U}$ est un raffinement de $\mathcal{V}$ ;
il existe donc un morphisme de complexes de {\v C}ech
$\check{\mathcal{C}}^\bullet (\mathcal{V}, \mathcal{F}) \to
\check{\mathcal{C}}^\bullet (\mathcal{U}, \mathcal{F})$, voir
Cohomology on Sites,
Equation (\ref{sites-cohomology-equation-map-cech-complexes}).
Ensuite, observons que si $T = \coprod_{j \in J} T_j$ est une
somme disjointe de sch\'emas du site sur lequel $\mathcal{F}$ est d\'efini,
alors la famille de morphismes \`a cible fix\'ee
$\{T_j \to T\}_{j \in J}$ est un recouvrement de Zariski, et donc
\begin{equation}
\label{equation-sheaf-coprod}
\mathcal{F}(T) =
\mathcal{F}(\coprod\nolimits_{j \in J} T_j) =
\prod\nolimits_{j \in J} \mathcal{F}(T_j)
\end{equation}
d'apr\`es la condition de faisceau pour $\mathcal{F}$.
Cela implique que le morphisme de complexes de {\v C}ech ci-dessus est un isomorphisme
en chaque degr\'e, car
$$
V \times_U \ldots \times_U V
=
\coprod\nolimits_{i_0, \ldots i_p} U_{i_0} \times_U \ldots \times_U U_{i_p}
$$
en tant que sch\'emas.
\end{proof}

\noindent
Notons que l'\'egalit\'e (\ref{equation-sheaf-coprod})
est fausse pour un pr\'efaisceau g\'en\'eral. M\^eme pour les faisceaux, elle n'est pas vraie sur tout
site, car les coproduits peuvent ne pas donner des recouvrements et ne pas \^etre disjoints.
Mais elle est vraie pour tous les sites usuels (du moins pour tous ceux que nous \'etudierons).

\begin{remark}
\label{remark-refinement}
Dans l'\'enonc\'e du lemme \ref{lemma-cech-complex}, le recouvrement $\mathcal{U}$
est un raffinement de $\mathcal{V}$, mais non l'inverse. Les recouvrements
de la forme $\{V \to U\}$ ne forment pas une sous-cat\'egorie initiale de la
cat\'egorie de tous les recouvrements de $U$. Il reste pourtant vrai que
nous pouvons calculer la cohomologie de {\v C}ech $\check H^n(U, \mathcal{F})$ (qui
est d\'efinie comme la limite inductive index\'ee par la cat\'egorie oppos\'ee \`a celle des
recouvrements $\mathcal{U}$ de $U$ des groupes de cohomologie de {\v C}ech de
$\mathcal{F}$ relativement \`a $\mathcal{U}$) au moyen des recouvrements
$\{V \to U\}$. Nous formulerons un lemme pr\'ecis (il ne vaut que pour les faisceaux)
et l'ajouterons ici si nous en avons un jour besoin.
\end{remark}

\begin{lemma}[Localit\'e de la cohomologie]
\label{lemma-locality-cohomology}
Soient $\mathcal{C}$ un site, $\mathcal{F}$ un faisceau ab\'elien sur $\mathcal{C}$,
$U$ un objet de $\mathcal{C}$, $p > 0$ un entier et $\xi \in
H^p(U, \mathcal{F})$. Alors il existe un recouvrement
$\mathcal{U} = \{U_i \to U\}_{i \in I}$ de $U$ dans $\mathcal{C}$
tel que $\xi |_{U_i} = 0$ pour tout $i \in I$.
\end{lemma}

\begin{proof}
Choisissons une r\'esolution injective $\mathcal{F} \to \mathcal{I}^\bullet$. Alors
$\xi$ est repr\'esent\'ee par un cocycle $\tilde{\xi} \in \mathcal{I}^p(U)$
tel que $d^p(\tilde{\xi}) = 0$. Par hypoth\`ese, la suite
$\mathcal{I}^{p - 1} \to \mathcal{I}^p \to \mathcal{I}^{p + 1}$ est exacte dans
$\textit{Ab}(\mathcal{C})$, ce qui signifie qu'il existe un recouvrement
$\mathcal{U} = \{U_i \to U\}_{i \in I}$ tel que
$\tilde{\xi}|_{U_i} = d^{p - 1}(\xi_i)$ pour un certain
$\xi_i \in \mathcal{I}^{p-1}(U_i)$. Puisque
la classe de cohomologie $\xi|_{U_i}$ est repr\'esent\'ee par le cocycle
$\tilde{\xi}|_{U_i}$, qui est un cobord, elle s'annule.
Pour plus de d\'etails, voir
Cohomology on Sites,
Lemma \ref{sites-cohomology-lemma-kill-cohomology-class-on-covering}.
\end{proof}

\begin{theorem}
\label{theorem-zariski-fpqc-quasi-coherent}
Soit $S$ un sch\'ema et soit $\mathcal{F}$ un $\mathcal{O}_S$-module quasi-coh\'erent.
Soit $\mathcal{C}$ l'un des sites $(\Sch/S)_\tau$ pour
$\tau \in \{fppf, syntomic, smooth, \etale, Zariski\}$, ou
$S_\etale$. Alors
$$
H^p(S, \mathcal{F}) = H^p_\tau(S, \mathcal{F}^a)
$$
pour tout $p \geq 0$, o\`u
\begin{enumerate}
\item le membre de gauche d\'esigne la cohomologie usuelle du faisceau
$\mathcal{F}$ sur l'espace topologique sous-jacent au sch\'ema $S$, et
\item le membre de droite d\'esigne la cohomologie
du faisceau ab\'elien $\mathcal{F}^a$ (voir
la proposition \ref{proposition-quasi-coherent-sheaf-fpqc})
sur le site $\mathcal{C}$.
\end{enumerate}
\end{theorem}

\begin{proof}
Nous allons montrer que
$H^p(U, f^*\mathcal{F}) = H^p_\tau(U, \mathcal{F}^a)$
pour tout objet $f : U \to S$ du site $\mathcal{C}$.
Le r\'esultat est vrai pour $p = 0$ par la propri\'et\'e de faisceau.

\medskip\noindent
Supposons que $U$ soit affine. Nous voulons alors prouver que
$H^p_\tau(U, \mathcal{F}^a) = 0$ pour tout $p > 0$. Nous raisonnons par r\'ecurrence sur $p$.
\begin{enumerate}
\item[p = 1]
Prenons $\xi \in H^1_\tau(U, \mathcal{F}^a)$.
D'apr\`es le lemme \ref{lemma-locality-cohomology},
il existe un recouvrement fpqc $\mathcal{U} = \{U_i \to U\}_{i \in I}$
tel que $\xi|_{U_i} = 0$ pour tout $i \in I$. Quitte \`a raffiner
$\mathcal{U}$, nous pouvons supposer que $\mathcal{U}$ est un
$\tau$-recouvrement standard. En appliquant la suite spectrale du
th\'eor\`eme \ref{theorem-cech-ss},
nous voyons que $\xi$ provient d'une classe de cohomologie
$\check \xi \in \check H^1(\mathcal{U}, \mathcal{F}^a)$.
Consid\'erons le recouvrement $\mathcal{V} = \{\coprod_{i\in I} U_i \to U\}$. D'apr\`es le
lemme \ref{lemma-cech-complex},
$\check H^\bullet(\mathcal{U}, \mathcal{F}^a) =
\check H^\bullet(\mathcal{V}, \mathcal{F}^a)$.
D'autre part, puisque $\mathcal{V}$ est un recouvrement de la forme
$\{\Spec(B) \to \Spec(A)\}$ et $f^*\mathcal{F} = \widetilde{M}$
pour un certain $A$-module $M$, nous voyons que le complexe de {\v C}ech
$\check{\mathcal{C}}^\bullet(\mathcal{V}, \mathcal{F})$
n'est autre que le complexe $(B/A)_\bullet \otimes_A M$.
D'apr\`es le lemme \ref{lemma-descent-modules},
$H^p((B/A)_\bullet \otimes_A M) = 0$ pour $p > 0$ ; ainsi, $\check \xi = 0$
et donc $\xi = 0$.
\item[p > 1]
Prenons $\xi \in H^p_\tau(U, \mathcal{F}^a)$. D'apr\`es le
lemme \ref{lemma-locality-cohomology},
il existe un recouvrement fpqc $\mathcal{U} = \{U_i \to U\}_{i \in I}$
tel que $\xi|_{U_i} = 0$ pour tout $i \in I$. Quitte \`a raffiner
$\mathcal{U}$, nous pouvons supposer que $\mathcal{U}$ est un
$\tau$-recouvrement standard. Nous appliquons la suite spectrale du
th\'eor\`eme \ref{theorem-cech-ss}.
Observons que les intersections $U_{i_0} \times_U \ldots \times_U U_{i_p}$
sont affines ; par l'hypoth\`ese de r\'ecurrence, les groupes de cohomologie
$$
E_2^{p, q} = \check H^p(\mathcal{U}, \underline{H}^q(\mathcal{F}^a))
$$
s'annulent donc pour tout $0 < q < p$. Nous voyons que $\xi$ doit provenir d'un
$\check \xi \in \check H^p(\mathcal{U}, \mathcal{F}^a)$. En rempla\c{c}ant
$\mathcal{U}$ par le recouvrement $\mathcal{V}$ qui ne contient qu'un seul morphisme
et en utilisant de nouveau le lemme \ref{lemma-descent-modules},
nous voyons que la classe de cohomologie de {\v C}ech $\check \xi$ doit \^etre nulle,
donc $\xi = 0$.
\end{enumerate}
Supposons ensuite que $U$ soit s\'epar\'e. Choisissons des ouverts affines
$U = \bigcup_{i \in I} U_i$ formant un recouvrement de $U$. La famille
$\mathcal{U} = \{U_i \to U\}_{i \in I}$ est alors un recouvrement fpqc,
et toutes les intersections
$U_{i_0} \times_U \ldots \times_U U_{i_p}$ sont affines,
puisque $U$ est s\'epar\'e. Toutes les lignes de la suite spectrale du
th\'eor\`eme \ref{theorem-cech-ss}
sont donc nulles, sauf la ligne de degr\'e z\'ero. Par cons\'equent,
$$
H^p_\tau(U, \mathcal{F}^a) =
\check H^p(\mathcal{U}, \mathcal{F}^a) =
\check H^p(\mathcal{U}, \mathcal{F}) = H^p(U, \mathcal{F})
$$
o\`u la derni\`ere \'egalit\'e r\'esulte de la th\'eorie usuelle des sch\'emas, voir
Cohomology of Schemes, Lemma
\ref{coherent-lemma-cech-cohomology-quasi-coherent}.

\medskip\noindent
Le cas g\'en\'eral est technique et (pour prolonger la d\'emonstration donn\'ee ici)
n\'ecessite une discussion sur les morphismes de suites spectrales ; nous ne le traiterons donc pas.
Il r\'esulte de
Descent, Proposition \ref{descent-proposition-same-cohomology-quasi-coherent},
(dont la d\'emonstration adopte une approche l\'eg\`erement diff\'erente), combin\'ee avec
Cohomology on Sites, Lemma \ref{sites-cohomology-lemma-cohomology-of-open}.
\end{proof}

\begin{remark}
\label{remark-right-derived-global-sections}
Commentaire sur le th\'eor\`eme \ref{theorem-zariski-fpqc-quasi-coherent}.
Puisque $S$ est un objet final de la cat\'egorie $\mathcal{C}$, les groupes de cohomologie
du membre de droite sont simplement les foncteurs d\'eriv\'es \`a droite du
foncteur des sections globales. En fait, la d\'emonstration montre que
$H^p(U, f^*\mathcal{F}) = H^p_\tau(U, \mathcal{F}^a)$
pour tout objet $f : U \to S$ du site $\mathcal{C}$.
\end{remark}





\section{Exemples de faisceaux}
\label{section-examples-sheaves}

\noindent
Soient $S$ et $\tau$ comme dans la section \ref{section-big-small}.
Nous avons d\'ej\`a vu que tout pr\'efaisceau repr\'esentable est un faisceau sur
$(\Sch/S)_\tau$ ou $S_\tau$, voir
le lemme \ref{lemma-representable-sheaf-fpqc}
et la
remarque \ref{remark-fpqc-finest}.
Voici quelques cas particuliers.

\begin{definition}
\label{definition-additive-sheaf}
Sur l'un quelconque des sites $(\Sch/S)_\tau$ ou $S_\tau$ de la
section \ref{section-big-small} :
\begin{enumerate}
\item Le faisceau $T \mapsto \Gamma(T, \mathcal{O}_T)$ est not\'e
$\mathcal{O}_S$, ou $\mathbf{G}_a$, ou $\mathbf{G}_{a, S}$ si nous
voulons indiquer le sch\'ema de base.
\item De m\^eme, le faisceau
$T \mapsto \Gamma(T, \mathcal{O}^*_T)$ est not\'e $\mathcal{O}_S^*$, ou
$\mathbf{G}_m$, ou $\mathbf{G}_{m, S}$ si nous voulons
indiquer le sch\'ema de base.
\item Le {\it faisceau constant} $\underline{\mathbf{Z}/n\mathbf{Z}}$ sur un
site quelconque est le faisceau associ\'e au pr\'efaisceau constant
$U \mapsto \mathbf{Z}/n\mathbf{Z}$.
\end{enumerate}
\end{definition}

\noindent
Le premier est un faisceau, par exemple d'apr\`es le
th\'eor\`eme \ref{theorem-quasi-coherent}.
Le deuxi\`eme est un sous-pr\'efaisceau du premier, dont on v\'erifie facilement
qu'il est lui-m\^eme un faisceau. Le troisi\`eme est un faisceau par d\'efinition.
Notons que chacun de ces faisceaux est repr\'esentable.
Le premier et le deuxi\`eme le sont par les sch\'emas $\mathbf{G}_{a, S}$ et
$\mathbf{G}_{m, S}$, voir
Groupoids, Section \ref{groupoids-section-group-schemes}.
Le troisi\`eme l'est par le sch\'ema en groupes fini \'etale $\mathbf{Z}/n\mathbf{Z}_S$,
parfois not\'e $(\mathbf{Z}/n\mathbf{Z})_S$,
qui consiste simplement en $n$ copies de $S$ munies
de l'\'evidente structure de sch\'ema en groupes sur $S$, voir
Groupoids, Example \ref{groupoids-example-constant-group}
et la remarque suivante.

\begin{remark}
\label{remark-constant-locally-constant-maps}
Soit $G$ un groupe abstrait.
Sur l'un quelconque des sites $(\Sch/S)_\tau$ ou $S_\tau$ de la
section \ref{section-big-small},
le faisceau associ\'e $\underline{G}$
au pr\'efaisceau constant associ\'e \`a $G$ pour la
{\it topologie de Zariski} du site donne d\'ej\`a
$$
\Gamma(U, \underline{G}) =
\{\text{applications localement constantes pour la topologie de Zariski }U \to G\}
$$
Ce faisceau de Zariski est repr\'esentable par le sch\'ema en groupes $G_S$ d'apr\`es
Groupoids, Example \ref{groupoids-example-constant-group}.
D'apr\`es le
lemme \ref{lemma-representable-sheaf-fpqc},
tout pr\'efaisceau repr\'esentable satisfait aussi la condition de faisceau pour la
topologie $\tau$ ; nous concluons donc que le faisceau associ\'e de Zariski
$\underline{G}$ ci-dessus est aussi le faisceau associ\'e pour la topologie $\tau$.
\end{remark}

\begin{definition}
\label{definition-structure-sheaf}
Soit $S$ un sch\'ema. Le {\it faisceau structural} de $S$ est le faisceau d'anneaux
$\mathcal{O}_S$
sur l'un quelconque des sites $S_{Zar}$, $S_\etale$ ou $(\Sch/S)_\tau$
consid\'er\'es ci-dessus.
\end{definition}

\noindent
S'il peut y avoir une ambigu\"it\'e quant au site sur lequel nous travaillons,
nous l'indiquerons au moyen d'indices. Par exemple, nous pouvons utiliser
$\mathcal{O}_{S_\etale}$ pour souligner que nous travaillons sur le
petit site \'etale de $S$.

\begin{remark}
\label{remark-special-case-fpqc-cohomology-quasi-coherent}
Dans la terminologie introduite ci-dessus, un cas particulier du
th\'eor\`eme \ref{theorem-zariski-fpqc-quasi-coherent}
est
$$
H_{fppf}^p(X, \mathbf{G}_a) =
H_\etale^p(X, \mathbf{G}_a) =
H_{Zar}^p(X, \mathbf{G}_a) =
H^p(X, \mathcal{O}_X)
$$
pour tout $p \geq 0$. De plus, nous pourrions utiliser la notation
$H^p_{fppf}(X, \mathcal{O}_X)$ pour d\'esigner la cohomologie du
faisceau structural sur le grand site fppf de $X$.
\end{remark}




\section{Groupes de Picard}
\label{section-picard-groups}

\noindent
Le th\'eor\`eme suivant est parfois appel\'e ``th\'eor\`eme 90 de Hilbert''.

\begin{theorem}
\label{theorem-picard-group}
Pour tout sch\'ema $X$, nous avons des identifications canoniques
\begin{align*}
H_{fppf}^1(X, \mathbf{G}_m) & = H^1_{syntomic}(X, \mathbf{G}_m) \\
& = H^1_{smooth}(X, \mathbf{G}_m) \\
& = H_\etale^1(X, \mathbf{G}_m) \\
& = H^1_{Zar}(X, \mathbf{G}_m) \\
& = \Pic(X) \\
& = H^1(X, \mathcal{O}_X^*)
\end{align*}
\end{theorem}

\begin{proof}
Soit $\tau$ l'une des topologies consid\'er\'ees dans la
section \ref{section-big-small}.
D'apr\`es
Cohomology on Sites, Lemma
\ref{sites-cohomology-lemma-h1-invertible},
nous voyons que
$H^1_\tau(X, \mathbf{G}_m) =
H^1_\tau(X, \mathcal{O}_\tau^*) =
\Pic(\mathcal{O}_\tau)$,
o\`u $\mathcal{O}_\tau$ est le faisceau structural du site
$(\Sch/X)_\tau$. Or un $\mathcal{O}_\tau$-module inversible
est un $\mathcal{O}_\tau$-module quasi-coh\'erent.
D'apr\`es le th\'eor\`eme \ref{theorem-quasi-coherent}, ou plus pr\'ecis\'ement
Descent, Proposition \ref{descent-proposition-equivalence-quasi-coherent},
nous voyons que $\Pic(\mathcal{O}_\tau) = \Pic(X)$.
La derni\`ere \'egalit\'e se d\'emontre de la m\^eme mani\`ere.
\end{proof}






\section{Le site \'etale}
\label{section-etale-site}

\noindent
Nous commen\c{c}ons maintenant \`a examiner plus en d\'etail le site \'etale d'un
sch\'ema. Dans un premier temps, nous examinons bri\`evement la notion de
morphisme \'etale.





\section{Morphismes \'etales}
\label{section-etale-morphism}

\noindent
Pour plus de d\'etails, voir
Morphisms, Section \ref{morphisms-section-etale}
pour la d\'efinition formelle, et
\'Etale Morphisms, Sections
\ref{etale-section-etale-morphisms},
\ref{etale-section-structure-etale-map},
\ref{etale-section-etale-smooth},
\ref{etale-section-topological-etale},
\ref{etale-section-functorial-etale}, et
\ref{etale-section-properties-permanence}
pour un aper\c{c}u de propri\'et\'es int\'eressantes des morphismes \'etales.

\medskip\noindent
Rappelons qu'une alg\`ebre $A$ sur un corps alg\'ebriquement clos $k$ est
{\it lisse} si elle est de type fini et si le module des diff\'erentielles
$\Omega_{A/k}$ est localement libre de type fini, de rang \`egal \`a la dimension.
Un sch\'ema $X$ sur $k$ est {\it lisse} sur $k$ s'il est localement de type
fini et si chaque ouvert affine est le spectre d'une $k$-alg\`ebre lisse.
Si $k$ n'est pas alg\'ebriquement clos, alors une $k$-alg\`ebre $A$ est
une $k$-alg\`ebre lisse si $A \otimes_k \overline{k}$ est une
$\overline{k}$-alg\`ebre lisse. Un morphisme d'anneaux $A \to B$ est lisse s'il est
plat, de pr\'esentation finie et si, pour tout id\'eal premier $\mathfrak p \subset A$,
l'anneau de la fibre $\kappa(\mathfrak p) \otimes_A B$ est lisse sur le corps
r\'esiduel $\kappa(\mathfrak p)$. Plus g\'en\'eralement, un morphisme de sch\'emas est
{\it lisse} s'il est plat, localement de pr\'esentation finie et si ses
fibres g\'eom\'etriques sont lisses.

\medskip\noindent
Pour ces faits, voir
Morphisms, Section \ref{morphisms-section-smooth}.
Cela nous permet de d\'efinir un morphisme \'etale comme suit.

\begin{definition}
\label{definition-etale-morphism}
Un morphisme de sch\'emas est {\it \'etale} s'il est lisse de dimension relative 0.
\end{definition}

\noindent
En particulier, un morphisme de sch\'emas $X \to S$ est \'etale s'il est lisse
et si $\Omega_{X/S} = 0$.

\begin{proposition}
\label{proposition-etale-morphisms}
Propri\'et\'es des morphismes \'etales.
\begin{enumerate}
\item Soit $k$ un corps. Un morphisme de sch\'emas $U \to \Spec(k)$ est
\'etale si et seulement si $U \cong \coprod_{i \in I} \Spec(k_i)$,
o\`u, pour chaque $i \in I$,
le corps $k_i$ est une extension finie s\'eparable de $k$.
\item Soit $\varphi : U \to S$ un morphisme de sch\'emas. Les
conditions suivantes sont \`equivalentes :
\begin{enumerate}
\item $\varphi$ est \'etale,
\item $\varphi$ est localement de pr\'esentation finie, plat, et toutes ses fibres sont
\'etales,
\item $\varphi$ est plat, non ramifi\'e et localement de pr\'esentation finie.
\end{enumerate}
\item Un morphisme d'anneaux $A \to B$ est \'etale si et seulement si
$B \cong A[x_1, \ldots, x_n]/(f_1, \ldots, f_n)$
tel que $\Delta = \det \left( \frac{\partial f_i}{\partial x_j} \right)$
soit inversible dans $B$.
\item Le changement de base d'un morphisme \'etale est \'etale.
\item Les compos\'ees de morphismes \'etales sont \'etales.
\item Les produits fibr\'es et les produits de morphismes \'etales sont \'etales.
\item Un morphisme \'etale a pour dimension relative 0.
\item Soit $Y \to X$ un morphisme \'etale.
Si $X$ est r\'eduit (respectivement r\'egulier), alors $Y$ l'est aussi.
\item Les morphismes \'etales sont ouverts.
\item Si $X \to S$ et $Y \to S$ sont \'etales, alors tout
$S$-morphisme $X \to Y$ est lui aussi \'etale.
\end{enumerate}
\end{proposition}

\begin{proof}
Nous avons d\'emontr\'e ces faits (et davantage) dans les chapitres pr\'ec\'edents.
Voici une liste de r\'ef\'erences :
(1) Morphisms, Lemma \ref{morphisms-lemma-etale-over-field}.
(2) Morphisms, Lemmas \ref{morphisms-lemma-etale-flat-etale-fibres}
et \ref{morphisms-lemma-flat-unramified-etale}.
(3) Algebra, Lemma \ref{algebra-lemma-etale-standard-smooth}.
(4) Morphisms, Lemma \ref{morphisms-lemma-base-change-etale}.
(5) Morphisms, Lemma \ref{morphisms-lemma-composition-etale}.
(6) D\'ecoule formellement de (4) et (5).
(7) Morphisms, Lemmas \ref{morphisms-lemma-etale-locally-quasi-finite}
et \ref{morphisms-lemma-locally-quasi-finite-rel-dimension-0}.
(8) Voir Algebra, Lemmas \ref{algebra-lemma-reduced-goes-up} et
\ref{algebra-lemma-Rk-goes-up} ; voir aussi d'autres r\'esultats de ce type
dans \'Etale Morphisms, Section \ref{etale-section-properties-permanence}.
(9) Voir Morphisms, Lemma \ref{morphisms-lemma-fppf-open} et
\ref{morphisms-lemma-etale-flat}.
(10) Voir Morphisms, Lemma \ref{morphisms-lemma-etale-permanence}.
\end{proof}

\begin{definition}
\label{definition-standard-etale}
Un morphisme d'anneaux $A \to B$ est dit {\it \'etale standard} si
$B \cong \left(A[t]/(f)\right)_g$ avec $f, g \in A[t]$, o\`u $f$ est unitaire,
et si $\text{d}f/\text{d}t$ est inversible dans $B$.
\end{definition}

\noindent
Il est vrai qu'un morphisme d'anneaux \'etale standard est \'etale. En effet, supposons
que $B = \left(A[t]/(f)\right)_g$ avec $f, g \in A[t]$, o\`u $f$ est unitaire,
et que $\text{d}f/\text{d}t$ soit inversible dans $B$. Alors $A[t]/(f)$ est un
$A$-module libre de rang \'egal au degr\'e du polyn\^ome unitaire $f$.
Par cons\'equent, $B$, en tant que localisation de cette alg\`ebre qui est libre comme module, est de pr\'esentation finie
et plat sur $A$. Pour achever la preuve que $B$ est \'etale, il suffit
de montrer que les anneaux des fibres
$$
\kappa(\mathfrak p) \otimes_A B
\cong
\kappa(\mathfrak p) \otimes_A (A[t]/(f))_g
\cong
\kappa(\mathfrak p)[t, 1/\overline{g}]/(\overline{f})
$$
sont des produits finis d'extensions de corps finies s\'eparables.
Ici, $\overline{f}, \overline{g} \in \kappa(\mathfrak p)[t]$ sont
les images de $f$ et $g$. Soit
$$
\overline{f} = \overline{f}_1 \ldots \overline{f}_a
\overline{f}_{a + 1}^{e_1} \ldots \overline{f}_{a + b}^{e_b}
$$
la factorisation de $\overline{f}$ en puissances de facteurs unitaires irr\'eductibles
deux \`a deux distincts $\overline{f}_i$, avec $e_1, \ldots, e_b > 1$.
Par hypoth\`ese, $\text{d}\overline{f}/\text{d}t$ est inversible dans
$\kappa(\mathfrak p)[t, 1/\overline{g}]$. Nous voyons donc que
au moins les $\overline{f}_i$, $i > a$, sont inversibles. Nous concluons
que
$$
\kappa(\mathfrak p)[t, 1/\overline{g}]/(\overline{f})
\cong
\prod\nolimits_{i \in I} \kappa(\mathfrak p)[t]/(\overline{f}_i)
$$
o\`u $I \subset \{1, \ldots, a\}$ est le sous-ensemble des indices $i$ tels que
$\overline{f}_i$ ne divise pas $\overline{g}$. De plus, l'image de
$\text{d}\overline{f}/\text{d}t$ dans le facteur
$\kappa(\mathfrak p)[t]/(\overline{f}_i)$ est clairement \'egale au produit d'une
unit\'e par $\text{d}\overline{f}_i/\text{d}t$. Nous en concluons que
$\kappa_i = \kappa(\mathfrak p)[t]/(\overline{f}_i)$ est une extension de corps
finie de $\kappa(\mathfrak p)$, engendr\'ee par un \'el\'ement dont le
polyn\^ome minimal est s\'eparable, c'est-\`a-dire que l'extension de corps
$\kappa_i/\kappa(\mathfrak p)$ est finie s\'eparable, comme souhait\'e.

\medskip\noindent
Il se trouve que tout morphisme d'anneaux \'etale est localement \'etale standard.
Pour formuler cela, introduisons la notation suivante.
Un morphisme d'anneaux $A \to B$ est {\it \'etale en un id\'eal premier $\mathfrak q$} de $B$ s'il
existe $h \in B$, $h \not \in \mathfrak q$, tel que $A \to B_h$ soit \'etale.
Voici le r\'esultat.

\begin{theorem}
\label{theorem-standard-etale}
Un morphisme d'anneaux $A \to B$ est \'etale en un id\'eal premier $\mathfrak q$ si et seulement s'il
existe $g \in B$, $g \not \in \mathfrak q$, tel que $B_g$ soit \'etale
standard sur $A$.
\end{theorem}

\begin{proof}
Voir
Algebra, Proposition \ref{algebra-proposition-etale-locally-standard}.
\end{proof}





\section{Recouvrements \'etales}
\label{section-etale-covering}

\noindent
Rappelons la d\'efinition.

\begin{definition}
\label{definition-etale-covering}
Un {\it recouvrement \'etale} d'un sch\'ema $U$ est une famille de morphismes
de sch\'emas
$\{\varphi_i : U_i \to U\}_{i \in I}$ telle que
\begin{enumerate}
\item chaque $\varphi_i$ est un morphisme \'etale,
\item les $U_i$ recouvrent $U$, c'est-\`a-dire $U = \bigcup_{i\in I}\varphi_i(U_i)$.
\end{enumerate}
\end{definition}

\begin{lemma}
\label{lemma-etale-fpqc}
Tout recouvrement \'etale est un recouvrement fpqc.
\end{lemma}

\begin{proof}
(Voir aussi
Topologies,
Lemma \ref{topologies-lemma-zariski-etale-smooth-syntomic-fppf-fpqc}.)
Soit $\{\varphi_i : U_i \to U\}_{i \in I}$ un recouvrement \'etale.
Puisqu'un morphisme \'etale est plat et que les \'el\'ements du recouvrement doivent
recouvrir sa cible, la propri\'et\'e fp (fid\`element plat) est satisfaite.
Pour v\'erifier la propri\'et\'e qc (quasi-compacit\'e), soit $V \subset U$ un ouvert
affine, et \'ecrivons $\varphi_i^{-1}(V) = \bigcup_{j \in J_i} V_{ij}$
pour certains ouverts affines $V_{ij} \subset U_i$. Puisque $\varphi_i$ est ouvert
(car les morphismes \'etales sont ouverts), nous voyons que
$V = \bigcup_{i\in I} \bigcup_{j \in J_i} \varphi_i(V_{ij})$
est un recouvrement ouvert de $V$.
De plus, puisque $V$ est quasi-compact, ce recouvrement admet un
raffinement fini.
\end{proof}

\noindent
Ainsi, tout \'enonc\'e vrai pour les recouvrements fpqc
reste vrai {\it a fortiori} pour les recouvrements \'etales. Par
exemple, le site \'etale est sous-canonique.

\begin{definition}
\label{definition-big-etale-site}
(Pour plus de d\'etails, voir la section \ref{section-big-small}, ou
Topologies, Section \ref{topologies-section-etale}.)
Soit $S$ un sch\'ema.
Le {\it grand site \'etale de $S$} est le site
$(\Sch/S)_\etale$, voir
la d\'efinition \ref{definition-tau-site}.
Le {\it petit site \'etale de $S$} est le site $S_\etale$, voir
la d\'efinition \ref{definition-tau-site}.
Nous d\'efinissons de m\^eme le {\it grand} et le {\it petit site de Zariski} de $S$,
not\'es $(\Sch/S)_{Zar}$ et $S_{Zar}$.
\end{definition}

\noindent
Grosso modo, le grand site \'etale de $S$ est constitu\'e des sch\'emas sur $S$
et ses recouvrements sont les recouvrements \'etales. Le petit site \'etale de $S$ est constitu\'e
des sch\'emas \'etales sur $S$, avec pour recouvrements les recouvrements \'etales.
En fait, tout morphisme entre objets de $S_\etale$ est \'etale, en
vertu de la
proposition \ref{proposition-etale-morphisms} ;
par cons\'equent, pour v\'erifier que $\{U_i \to U\}_{i \in I}$ dans $S_\etale$
est un recouvrement, il suffit de v\'erifier que $\coprod U_i \to U$ est surjectif.

\medskip\noindent
Le petit site \'etale a moins d'objets que le grand site \'etale ; il
ne contient que les ``ouverts'' de la topologie \'etale sur $S$. C'est une sous-cat\'egorie
pleine du grand site \'etale, et sa topologie est induite par la
topologie du grand site. Ainsi, le foncteur de restriction
du grand site \'etale au petit est exact et transforme les objets injectifs en
objets injectifs. Cela a la cons\'equence suivante.

\begin{proposition}
\label{proposition-cohomology-restrict-small-site}
Soit $S$ un sch\'ema et soit $\mathcal{F}$ un faisceau ab\'elien sur
$(\Sch/S)_\etale$.
Alors $\mathcal{F}|_{S_\etale}$ est un faisceau sur $S_\etale$ et
$$
H^p_\etale(S, \mathcal{F}|_{S_\etale}) =
H^p_\etale(S, \mathcal{F})
$$
pour tout $p \geq 0$.
\end{proposition}

\begin{proof}
C'est un cas particulier du lemme \ref{lemma-compare-cohomology-big-small}.
\end{proof}

\noindent
Conform\'ement \`a la notation g\'en\'erale introduite dans la
section \ref{section-big-small},
nous notons $H_\etale^p(S, \mathcal{F})$ le groupe de cohomologie ci-dessus.





%9.24.09
\section{Th\'eorie de Kummer}
\label{section-kummer}

\noindent
Soit $n \in \mathbf{N}$ et consid\'erons le foncteur $\mu_n$ d\'efini par
$$
\begin{matrix}
\Sch^{opp} & \longrightarrow & \textit{Ab} \\
S & \longmapsto &
\mu_n(S)
=
\{t \in \Gamma(S, \mathcal{O}_S^*) \mid t^n = 1 \}.
\end{matrix}
$$
D'apr\`es
Groupoids, Example \ref{groupoids-example-roots-of-unity},
ce foncteur est repr\'esentable, et le sch\'ema qui le repr\'esente
est lui aussi not\'e $\mu_n$. D'apr\`es le
lemme \ref{lemma-representable-sheaf-fpqc},
ce foncteur satisfait \`a la condition de faisceau pour la topologie fpqc
(en particulier, il satisfait aussi \`a la condition de faisceau pour les
topologies \'etale, de Zariski, etc.).

\begin{lemma}
\label{lemma-kummer-sequence}
Si $n\in \mathcal{O}_S^*$, alors
$$
0 \to
\mu_{n, S} \to
\mathbf{G}_{m, S} \xrightarrow{(\cdot)^n}
\mathbf{G}_{m, S} \to 0
$$
est une suite exacte courte de faisceaux aussi bien sur le petit que sur le
grand site \'etale de $S$.
\end{lemma}

\begin{proof}
Par d\'efinition, le faisceau $\mu_{n, S}$ est le noyau de l'application
$(\cdot)^n$. Il suffit donc de montrer que la derni\`ere application est surjective.
Soit $U$ un sch\'ema sur $S$. Soit
$f \in \mathbf{G}_m(U) = \Gamma(U, \mathcal{O}_U^*)$.
Nous devons montrer que nous pouvons trouver un recouvrement \'etale de
$U$ tel que la restriction de $f$ \`a chacun de ses membres soit une puissance $n$-i\`eme.
Posons
$$
U' =
\underline{\Spec}_U(\mathcal{O}_U[T]/(T^n-f))
\xrightarrow{\pi}
U.
$$
(Voir
Constructions, Section \ref{constructions-section-spec-via-glueing} ou
\ref{constructions-section-spec}
pour une pr\'esentation du spectre relatif.)
Soit $\Spec(A) \subset U$ un ouvert affine, et supposons que $f|_{\Spec(A)}$ corresponde
\`a l'unit\'e $a \in A^*$. Alors $\pi^{-1}(\Spec(A)) = \Spec(B)$ avec
$B = A[T]/(T^n - a)$. Le morphisme d'anneaux $A \to B$ est fini libre de rang $n$,
donc fid\`element plat, et nous en concluons que
$\Spec(B) \to \Spec(A)$ est surjectif. Comme cela vaut pour tout
ouvert affine de $U$, nous concluons que $\pi$ est surjectif.
De plus, $n$ et $T^{n - 1}$ sont inversibles dans $B$, donc
$nT^{n-1} \in B^*$ et le morphisme d'anneaux $A \to B$ est \'etale standard,
en particulier \'etale. Comme cela vaut pour tout ouvert affine de $U$,
nous concluons que $\pi$ est \'etale. Par cons\'equent,
$\mathcal{U} = \{\pi : U' \to U\}$ est un recouvrement \'etale.
De plus, $f|_{U'} = (f')^n$, o\`u $f'$ est la classe de $T$
dans $\Gamma(U', \mathcal{O}_{U'}^*)$ ; ainsi, $\mathcal{U}$ poss\`ede la propri\'et\'e voulue.
\end{proof}

\begin{remark}
\label{remark-no-kummer-sequence-zariski}
Le lemme \ref{lemma-kummer-sequence} est faux lorsque ``\'etale'' est remplac\'e
par ``Zariski''.
Comme la topologie \'etale est moins fine que la topologie lisse, voir
Topologies, Lemma \ref{topologies-lemma-zariski-etale-smooth},
il s'ensuit que la suite est aussi exacte pour la topologie lisse.
\end{remark}

\noindent
D'apr\`es le
th\'eor\`eme \ref{theorem-picard-group},
le
lemme \ref{lemma-kummer-sequence}
et les propri\'et\'es g\'en\'erales de la cohomologie, nous obtenons
la suite exacte longue de cohomologie
$$
\xymatrix{
0 \ar[r] &
H_\etale^0(S, \mu_{n, S}) \ar[r] &
\Gamma(S, \mathcal{O}_S^*) \ar[r]^{(\cdot)^n} &
\Gamma(S, \mathcal{O}_S^*) \ar@(rd, ul)[rdllllr]
\\
& H_\etale^1(S, \mu_{n, S}) \ar[r] &
\Pic(S) \ar[r]^{(\cdot)^n} &
\Pic(S) \ar@(rd, ul)[rdllllr] \\
& H_\etale^2(S, \mu_{n, S}) \ar[r] &
\ldots
}
$$
du moins si $n$ est inversible sur $S$. Lorsque $n$ n'est pas inversible sur $S$,
nous pouvons appliquer le lemme suivant.

\begin{lemma}
\label{lemma-kummer-sequence-syntomic}
Pour tout $n \in \mathbf{N}$, la suite
$$
0 \to
\mu_{n, S} \to
\mathbf{G}_{m, S} \xrightarrow{(\cdot)^n}
\mathbf{G}_{m, S} \to 0
$$
est une suite exacte courte de faisceaux sur les sites
$(\Sch/S)_{fppf}$ et $(\Sch/S)_{syntomic}$.
\end{lemma}

\begin{proof}
Par d\'efinition, le faisceau $\mu_{n, S}$ est le noyau de l'application
$(\cdot)^n$. Il suffit donc de montrer que la derni\`ere application est surjective.
Comme la topologie syntomique est moins fine que la topologie fppf, voir
Topologies, Lemma \ref{topologies-lemma-zariski-etale-smooth-syntomic-fppf},
il suffit de le d\'emontrer pour la topologie syntomique.
Soit $U$ un sch\'ema sur $S$. Soit
$f \in \mathbf{G}_m(U) = \Gamma(U, \mathcal{O}_U^*)$.
Nous devons montrer que nous pouvons trouver un recouvrement syntomique de
$U$ tel que la restriction de $f$ \`a chacun de ses membres soit une puissance $n$-i\`eme.
Posons
$$
U' =
\underline{\Spec}_U(\mathcal{O}_U[T]/(T^n-f))
\xrightarrow{\pi}
U.
$$
(Voir
Constructions, Section \ref{constructions-section-spec-via-glueing} ou
\ref{constructions-section-spec}
pour une pr\'esentation du spectre relatif.)
Soit $\Spec(A) \subset U$ un ouvert affine, et supposons que $f|_{\Spec(A)}$ corresponde
\`a l'unit\'e $a \in A^*$. Alors $\pi^{-1}(\Spec(A)) = \Spec(B)$ avec
$B = A[T]/(T^n - a)$. Le morphisme d'anneaux $A \to B$ est fini libre de rang $n$,
donc fid\`element plat, et nous en concluons que
$\Spec(B) \to \Spec(A)$ est surjectif. Comme cela vaut pour tout
ouvert affine de $U$, nous concluons que $\pi$ est surjectif.
De plus, $B$ est une intersection compl\`ete globale relative sur $A$, donc
le morphisme d'anneaux $A \to B$ est syntomique standard,
en particulier syntomique. Comme cela vaut pour tout ouvert affine de $U$,
nous concluons que $\pi$ est syntomique. Par cons\'equent,
$\mathcal{U} = \{\pi : U' \to U\}$ est un recouvrement syntomique.
De plus, $f|_{U'} = (f')^n$, o\`u $f'$ est la classe de $T$
dans $\Gamma(U', \mathcal{O}_{U'}^*)$ ; ainsi, $\mathcal{U}$ poss\`ede la propri\'et\'e voulue.
\end{proof}

\begin{remark}
\label{remark-no-kummer-sequence-smooth-etale-zariski}
Le lemme \ref{lemma-kummer-sequence-syntomic}
est faux pour les topologies lisse, \'etale ou de Zariski.
\end{remark}

\noindent
D'apr\`es le
th\'eor\`eme \ref{theorem-picard-group},
le
lemme \ref{lemma-kummer-sequence-syntomic}
et les propri\'et\'es g\'en\'erales de la cohomologie, nous obtenons
la suite exacte longue de cohomologie
$$
\xymatrix{
0 \ar[r] &
H_{fppf}^0(S, \mu_{n, S}) \ar[r] &
\Gamma(S, \mathcal{O}_S^*) \ar[r]^{(\cdot)^n} &
\Gamma(S, \mathcal{O}_S^*) \ar@(rd, ul)[rdllllr]
\\
& H_{fppf}^1(S, \mu_{n, S}) \ar[r] &
\Pic(S) \ar[r]^{(\cdot)^n} &
\Pic(S) \ar@(rd, ul)[rdllllr] \\
& H_{fppf}^2(S, \mu_{n, S}) \ar[r] &
\ldots
}
$$
pour tout sch\'ema $S$ et tout entier $n$. Bien entendu, il existe une suite analogue
en cohomologie syntomique.

\medskip\noindent
Soit $n \in \mathbf{N}$ et soit $S$ un sch\'ema quelconque.
Il existe une autre fa\c{c}on, plus directe, de d\'ecrire le premier groupe de cohomologie \`a
valeurs dans $\mu_n$. Consid\'erons les paires
$(\mathcal{L}, \alpha)$ o\`u $\mathcal{L}$ est un faisceau inversible sur $S$
et $\alpha : \mathcal{L}^{\otimes n} \to \mathcal{O}_S$ est une trivialisation
de la $n$-i\`eme puissance tensorielle de $\mathcal{L}$.
Soit $(\mathcal{L}', \alpha')$ une seconde paire de ce type.
Un isomorphisme $\varphi : (\mathcal{L}, \alpha) \to (\mathcal{L}', \alpha')$
est un isomorphisme $\varphi : \mathcal{L} \to \mathcal{L}'$ de faisceaux
inversibles tel que le diagramme
$$
\xymatrix{
\mathcal{L}^{\otimes n} \ar[d]_{\varphi^{\otimes n}} \ar[r]_\alpha &
\mathcal{O}_S \ar[d]^1 \\
(\mathcal{L}')^{\otimes n} \ar[r]^{\alpha'} &
\mathcal{O}_S \\
}
$$
soit commutatif. Nous avons donc
\begin{equation}
\label{equation-isomorphisms-pairs}
\mathit{Isom}_S((\mathcal{L}, \alpha), (\mathcal{L}', \alpha'))
=
\left\{
\begin{matrix}
\emptyset & \text{si} & \text{elles ne sont pas isomorphes} \\
H^0(S, \mu_{n, S})\cdot \varphi & \text{si} &
\varphi \text{ est un isomorphisme de paires}
\end{matrix}
\right.
\end{equation}
De plus, \'etant donn\'ees deux paires $(\mathcal{L}, \alpha)$, $(\mathcal{L}', \alpha')$,
le produit tensoriel
$$
(\mathcal{L}, \alpha) \otimes (\mathcal{L}', \alpha')
=
(\mathcal{L} \otimes \mathcal{L}', \alpha \otimes \alpha')
$$
est encore une paire. La paire $(\mathcal{O}_S, 1)$ est un \'el\'ement neutre pour cette
op\'eration de produit tensoriel, et un inverse est donn\'e par
$$
(\mathcal{L}, \alpha)^{-1} = (\mathcal{L}^{\otimes -1}, \alpha^{\otimes -1}).
$$
Ainsi, l'ensemble des classes d'isomorphisme de paires forme un groupe ab\'elien.
Notons que
$$
(\mathcal{L}, \alpha)^{\otimes n}
=
(\mathcal{L}^{\otimes n}, \alpha^{\otimes n})
\xrightarrow{\alpha}
(\mathcal{O}_S, 1)
$$
est un isomorphisme ;
par cons\'equent, tout \'el\'ement de ce groupe est d'ordre divisant $n$. Nous avertissons le lecteur
que ce groupe n'est en g\'en\'eral {\bf pas} la $n$-torsion de $\Pic(S)$.

\begin{lemma}
\label{lemma-describe-h1-mun}
Soit $S$ un sch\'ema. Il existe une identification canonique
$$
H_\etale^1(S, \mu_n) =
\text{groupe des paires }(\mathcal{L}, \alpha)\text{ \`a isomorphisme pr\`es comme ci-dessus}
$$
si $n$ est inversible sur $S$. En g\'en\'eral, nous avons
$$
H_{fppf}^1(S, \mu_n) =
\text{groupe des paires }(\mathcal{L}, \alpha)\text{ \`a isomorphisme pr\`es comme ci-dessus}.
$$
Le m\^eme r\'esultat vaut en rempla\c{c}ant fppf par syntomique.
\end{lemma}

\begin{proof}
Nous d\'emontrons d'abord le second isomorphisme.
Soit $(\mathcal{L}, \alpha)$ une paire comme ci-dessus.
Choisissons un recouvrement de $S = \bigcup U_i$ par des ouverts affines tel que
$\mathcal{L}|_{U_i} \cong \mathcal{O}_{U_i}$. Choisissons un g\'en\'erateur $s_i \in \mathcal{L}(U_i)$.
Alors $\alpha(s_i^{\otimes n}) = f_i \in \mathcal{O}_S^*(U_i)$.
En \'ecrivant $U_i = \Spec(A_i)$, nous voyons qu'il existe une
intersection compl\`ete globale relative $A_i \to B_i = A_i[T]/(T^n - f_i)$
telle que $f_i$ devienne une puissance $n$-i\`eme dans $B_i$. Autrement dit, en posant
$V_i = \Spec(B_i)$, nous obtenons un recouvrement syntomique
$\mathcal{V} = \{V_i \to S\}_{i \in I}$ et des trivialisations
$\varphi_i : (\mathcal{L}, \alpha)|_{V_i} \to (\mathcal{O}_{V_i}, 1)$.

\medskip\noindent
Nous utiliserons ce r\'esultat (l'existence du recouvrement $\mathcal{V}$)
pour associer \`a cette paire une classe de cohomologie dans
$H^1_{syntomic}(S, \mu_{n, S})$. Nous donnons deux constructions (\'equivalentes).

\medskip\noindent
Premi\`ere construction : au moyen de la cohomologie de {\v C}ech.
Sur les intersections doubles $V_i \times_S V_j$, nous avons l'isomorphisme
$$
(\mathcal{O}_{V_i \times_S V_j}, 1)
\xrightarrow{\text{pr}_0^*\varphi_i^{-1}}
(\mathcal{L}|_{V_i \times_S V_j}, \alpha|_{V_i \times_S V_j})
\xrightarrow{\text{pr}_1^*\varphi_j}
(\mathcal{O}_{V_i \times_S V_j}, 1)
$$
de paires. D'apr\`es (\ref{equation-isomorphisms-pairs}), il est donn\'e par un
\'el\'ement $\zeta_{ij} \in \mu_n(V_i \times_S V_j)$. Nous omettons la v\'erification
que ces $\zeta_{ij}$ forment un $1$-cocycle, c'est-\`a-dire qu'ils donnent
un \'el\'ement $(\zeta_{i_0i_1}) \in \check C(\mathcal{V}, \mu_n)$
tel que $d(\zeta_{i_0i_1}) = 0$. Sa classe est donc un \'el\'ement de
$\check H^1(\mathcal{V}, \mu_n)$ et, d'apr\`es le
th\'eor\`eme \ref{theorem-cech-ss},
elle a pour image une classe de cohomologie dans $H^1_{syntomic}(S, \mu_{n, S})$.

\medskip\noindent
Seconde construction : au moyen des torseurs. Consid\'erons le pr\'efaisceau
$$
\mu_n(\mathcal{L}, \alpha) :
U
\longmapsto
\mathit{Isom}_U((\mathcal{O}_U, 1), (\mathcal{L}, \alpha)|_U)
$$
sur $(\Sch/S)_{syntomic}$.
Nous pouvons le voir comme un sous-pr\'efaisceau de
$\SheafHom_\mathcal{O}(\mathcal{O}, \mathcal{L})$ (faisceau Hom interne,
voir
Modules on Sites, Section \ref{sites-modules-section-internal-hom}).
Comme les conditions qui d\'efinissent ce sous-pr\'efaisceau sont locales, nous voyons que c'est
un faisceau.
D'apr\`es (\ref{equation-isomorphisms-pairs}), ce faisceau est muni d'une action libre du
faisceau $\mu_{n, S}$. Il reste donc seulement \`a v\'erifier qu'il
admet localement des sections. C'est vrai gr\^ace \`a l'existence du
recouvrement trivialisant $\mathcal{V}$. Ainsi, $\mu_n(\mathcal{L}, \alpha)$
est un $\mu_{n, S}$-torseur et, d'apr\`es
Cohomology on Sites, Lemma \ref{sites-cohomology-lemma-torsors-h1},
nous obtenons un \'el\'ement correspondant de $H^1_{syntomic}(S, \mu_{n, S})$.

\medskip\noindent
Il nous reste maintenant \`a montrer les points suivants :
\begin{enumerate}
\item Les deux constructions donnent la m\^eme classe de cohomologie.
\item Des paires isomorphes donnent la m\^eme classe de cohomologie.
\item La classe de cohomologie de
$(\mathcal{L}, \alpha) \otimes (\mathcal{L}', \alpha')$
est la somme des classes de cohomologie de
$(\mathcal{L}, \alpha)$ et $(\mathcal{L}', \alpha')$.
\item Si la classe de cohomologie est triviale, alors la paire est triviale.
\item Tout \'el\'ement de $H^1_{syntomic}(S, \mu_{n, S})$ est la
classe de cohomologie d'une paire.
\end{enumerate}
Nous omettons la preuve de (1). Le point (2) ressort clairement de la seconde construction,
puisque des torseurs isomorphes donnent la m\^eme classe de cohomologie.
Le point (3) ressort clairement de la premi\`ere construction, puisque les classes de
{\v C}ech obtenues s'additionnent. Le point (4) ressort clairement de la seconde construction,
puisqu'un torseur est trivial si et seulement s'il poss\`ede une section globale, voir
Cohomology on Sites, Lemma \ref{sites-cohomology-lemma-trivial-torsor}.

\medskip\noindent
Le point (5) se voit comme suit (bien qu'une preuve directe soit
pr\'ef\'erable). Supposons que $\xi \in H^1_{syntomic}(S, \mu_{n, S})$.
Alors $\xi$ s'envoie sur un \'el\'ement
$\overline{\xi} \in H^1_{syntomic}(S, \mathbf{G}_{m, S})$
tel que $n \overline{\xi} = 0$. D'apr\`es le
th\'eor\`eme \ref{theorem-picard-group},
nous voyons que $\overline{\xi}$ correspond \`a un faisceau inversible $\mathcal{L}$
dont la $n$-i\`eme puissance tensorielle est isomorphe \`a $\mathcal{O}_S$.
Il existe donc une paire $(\mathcal{L}, \alpha')$ dont la classe de cohomologie
$\xi'$ a la m\^eme image $\overline{\xi'}$ dans
$H^1_{syntomic}(S, \mathbf{G}_{m, S})$. Il suffit donc de montrer
que $\xi - \xi'$ est la classe d'une paire. Par construction et d'apr\`es la
suite exacte longue de cohomologie ci-dessus, nous voyons que
$\xi - \xi' = \partial(f)$ pour un certain $f \in H^0(S, \mathcal{O}_S^*)$.
Consid\'erons la paire $(\mathcal{O}_S, f)$. Nous omettons la v\'erification
que la classe de cohomologie de cette paire est $\partial(f)$, ce qui
ach\`eve la preuve de la premi\`ere identification (avec fppf remplac\'ee
par syntomique).

\medskip\noindent
Pour voir la premi\`ere identification, notons que si $n$ est inversible sur $S$, alors le
recouvrement $\mathcal{V}$ construit dans la premi\`ere partie de la preuve
est en fait un recouvrement \'etale (comparer avec la preuve du
lemme \ref{lemma-kummer-sequence}). Le reste de la preuve est ind\'ependant
de la topologie, \`a l'exception du tout dernier argument, qui utilise l'exactitude de
la suite de Kummer, c'est-\`a-dire le lemme \ref{lemma-kummer-sequence}.
\end{proof}






\section{Voisinages, fibres et points}
\label{section-stalks}

\noindent
Nous pouvons associer \`a tout point g\'eom\'etrique de $S$ un foncteur fibre qui est
exact. Un morphisme de faisceaux sur $S_\etale$ est un isomorphisme si et seulement
s'il
induit un isomorphisme sur toutes ces fibres. Un complexe de faisceaux ab\'eliens est
exact si et seulement si le complexe des fibres est exact en tout point g\'eom\'etrique.
Ces faits, pris ensemble, signifient que le petit site \'etale d'un sch\'ema $S$
poss\`ede assez de points. Il se trouve aussi que tout point du petit topos \'etale
de $S$ (notion abstraite) est donn\'e par un point g\'eom\'etrique.
Ainsi, en un certain sens, le petit topos \'etale de $S$ peut se comprendre en
termes de points g\'eom\'etriques et de voisinages.

\begin{definition}
\label{definition-geometric-point}
Soit $S$ un sch\'ema.
\begin{enumerate}
\item Un {\it point g\'eom\'etrique} de $S$ est un morphisme
$\Spec(k) \to S$ o\`u $k$ est alg\'ebriquement clos.
Un tel point se note habituellement $\overline{s}$, autrement dit au moyen d'une
lettre minuscule surmont\'ee d'une barre. Nous employons souvent $\overline{s}$ pour d\'esigner aussi bien le sch\'ema
$\Spec(k)$ que le morphisme, et nous employons $\kappa(\overline{s})$
pour d\'esigner $k$.
\item Nous disons que $\overline{s}$ {\it est au-dessus de} $s$
pour indiquer que $s \in S$ est l'image de $\overline{s}$.
\item Un {\it voisinage \'etale} d'un point g\'eom\'etrique $\overline{s}$
de $S$ est un diagramme commutatif
$$
\xymatrix{
& U \ar[d]^\varphi \\
{\overline{s}} \ar[r]^{\overline{s}} \ar[ur]^{\bar u} & S
}
$$
o\`u $\varphi$ est un morphisme \'etale de sch\'emas.
Nous \'ecrivons $(U, \overline{u}) \to (S, \overline{s})$.
\item Un {\it morphisme de voisinages \'etales}
$(U, \overline{u}) \to (U', \overline{u}')$
est un $S$-morphisme $h: U \to U'$
tel que $\overline{u}' = h \circ \overline{u}$.
\end{enumerate}
\end{definition}

\begin{remark}
\label{remark-etale-between-etale}
Puisque $U$ et $U'$ sont \'etales sur $S$, tout $S$-morphisme
entre eux est lui aussi \'etale, voir la
proposition \ref{proposition-etale-morphisms}.
En particulier, tous les morphismes de voisinages \'etales sont \'etales.
\end{remark}

\begin{remark}
\label{remark-etale-neighbourhoods}
Soient $S$ un sch\'ema et $s \in S$ un point. Dans
More on Morphisms,
Definition \ref{more-morphisms-definition-etale-neighbourhood},
nous avons d\'efini la notion de voisinage \'etale $(U, u) \to (S, s)$
de $(S, s)$. Si $\overline{s}$ est un point g\'eom\'etrique de $S$ au-dessus de
$s$, alors tout voisinage \'etale $(U, \overline{u}) \to (S, \overline{s})$
donne lieu \`a un voisinage \'etale $(U, u)$ de $(S, s)$ en prenant pour
$u \in U$ l'unique point de $U$ tel que $\overline{u}$
soit au-dessus de $u$. R\'eciproquement, \'etant donn\'e un voisinage \'etale $(U, u)$
de $(S, s)$, l'extension de corps r\'esiduels $\kappa(u)/\kappa(s)$
est finie s\'eparable (voir la
proposition \ref{proposition-etale-morphisms}) ;
nous pouvons donc trouver un plongement $\kappa(u) \subset \kappa(\overline{s})$
qui induit l'identit\'e sur $\kappa(s)$. Autrement dit, nous pouvons trouver un point g\'eom\'etrique
$\overline{u}$ de $U$ au-dessus de $u$ tel que $(U, \overline{u})$
soit un voisinage \'etale de $(S, \overline{s})$.
Nous utiliserons ces observations pour passer d'un type de
voisinage \'etale \`a l'autre.
\end{remark}

\begin{lemma}
\label{lemma-cofinal-etale}
Soit $S$ un sch\'ema, et soit $\overline{s}$ un point g\'eom\'etrique de $S$.
La cat\'egorie des voisinages \'etales est cofiltrante. Plus pr\'ecis\'ement :
\begin{enumerate}
\item Soient $(U_i, \overline{u}_i)_{i = 1, 2}$ deux voisinages \'etales de
$\overline{s}$ dans $S$. Alors il existe un troisi\`eme voisinage \'etale
$(U, \overline{u})$ et des morphismes
$(U, \overline{u}) \to (U_i, \overline{u}_i)$, $i = 1, 2$.
\item Soient $h_1, h_2: (U, \overline{u}) \to (U', \overline{u}')$ deux
morphismes entre voisinages \'etales de $\overline{s}$. Alors il existe un
voisinage \'etale $(U'', \overline{u}'')$ et un morphisme
$h : (U'', \overline{u}'') \to (U, \overline{u})$
qui \'egalise $h_1$ et $h_2$, c'est-\`a-dire tel que
$h_1 \circ h = h_2 \circ h$.
\end{enumerate}
\end{lemma}

\begin{proof}
Pour le point (1), consid\'erons le produit fibr\'e $U = U_1 \times_S U_2$.
Il est \'etale sur $U_1$ et sur $U_2$, puisque les morphismes \'etales sont
pr\'eserv\'es par changement de base, voir la
proposition \ref{proposition-etale-morphisms}.
Le morphisme $\overline{s} \to U$ d\'efini par $(\overline{u}_1, \overline{u}_2)$
lui conf\`ere la structure d'un voisinage \'etale muni de morphismes vers
$U_1$ et $U_2$. Pour le point (2), d\'efinissons $U''$ comme le produit fibr\'e
$$
\xymatrix{
U'' \ar[r] \ar[d] & U \ar[d]^{(h_1, h_2)} \\
U' \ar[r]^-\Delta & U' \times_S U'.
}
$$
Puisque les compos\'ees de $\overline{u}$ et $\overline{u}'$ vers $S$ sont toutes deux \'egales \`a $\overline{s}$,
nous voyons que $\overline{u}'' = (\overline{u}, \overline{u}')$ est un point g\'eom\'etrique
de $U''$. En particulier, $U'' \not = \emptyset$.
De plus, puisque $U'$ est \'etale sur $S$, le produit fibr\'e
$U'\times_S U'$ l'est aussi (voir la
proposition \ref{proposition-etale-morphisms}).
Ainsi, la fl\`eche verticale $(h_1, h_2)$ est \'etale d'apr\`es la
remarque \ref{remark-etale-between-etale} ci-dessus.
Par cons\'equent, $U''$ est \'etale sur $U'$ par changement de base, et donc aussi
\'etale sur $S$ (car les compos\'ees de morphismes \'etales sont \'etales).
Ainsi, $(U'', \overline{u}'')$ est une solution au probl\`eme.
\end{proof}

\begin{lemma}
\label{lemma-geometric-lift-to-cover}
Soit $S$ un sch\'ema.
Soit $\overline{s}$ un point g\'eom\'etrique de $S$.
Soit $(U, \overline{u})$ un voisinage \'etale de $\overline{s}$.
Soit $\mathcal{U} = \{\varphi_i : U_i \to U \}_{i\in I}$ un recouvrement \'etale.
Alors il existe $i \in I$ et $\overline{u}_i : \overline{s} \to U_i$
tels que $\varphi_i : (U_i, \overline{u}_i) \to (U, \overline{u})$
soit un morphisme de voisinages \'etales.
\end{lemma}

\begin{proof}
Comme $U = \bigcup_{i\in I} \varphi_i(U_i)$, le produit fibr\'e
$\overline{s} \times_{\overline{u}, U, \varphi_i} U_i$ n'est pas vide
pour un certain $i$. Consid\'erons alors le diagramme cart\'esien
$$
\xymatrix{
\overline{s} \times_{\overline{u}, U, \varphi_i} U_i
\ar[d]^{\text{pr}_1} \ar[r]_-{\text{pr}_2} & U_i
\ar[d]^{\varphi_i} \\
\Spec(k) = \overline{s} \ar@/^1pc/[u]^\sigma
\ar[r]^-{\overline{u}} & U
}
$$
La projection $\text{pr}_1$ est le changement de base d'un morphisme \'etale ; elle
est donc \'etale, voir la
proposition \ref{proposition-etale-morphisms}.
Par cons\'equent, $\overline{s} \times_{\overline{u}, U, \varphi_i} U_i$
est une somme disjointe d'extensions finies s\'eparables de $k$, d'apr\`es la
proposition \ref{proposition-etale-morphisms}. Ici,
$\overline{s} = \Spec(k)$. Mais $k$ est alg\'ebriquement clos, donc toutes
ces extensions sont triviales, et il existe une section $\sigma$ de
$\text{pr}_1$. La compos\'ee
$\text{pr}_2 \circ \sigma$ donne un morphisme compatible avec $\overline{u}$.
\end{proof}

\begin{definition}
\label{definition-stalk}
Soit $S$ un sch\'ema.
Soit $\mathcal{F}$ un pr\'efaisceau sur $S_\etale$.
Soit $\overline{s}$ un point g\'eom\'etrique de $S$.
La {\it fibre} de $\mathcal{F}$ en $\overline{s}$ est
$$
\mathcal{F}_{\overline{s}}
=
\colim_{(U, \overline{u})} \mathcal{F}(U)
$$
o\`u $(U, \overline{u})$ parcourt tous les voisinages
\'etales de $\overline{s}$ dans $S$.
\end{definition}

\noindent
D'apr\`es le lemme \ref{lemma-cofinal-etale}, cette limite inductive est index\'ee par une cat\'egorie
filtrante, \`a savoir l'oppos\'ee de la cat\'egorie des voisinages \'etales.
Autrement dit, un \'el\'ement de $\mathcal{F}_{\overline{s}}$ peut \'etre
consid\'er\'e comme un triplet $(U, \overline{u}, \sigma)$ o\`u
$\sigma \in \mathcal{F}(U)$. Deux triplets
$(U, \overline{u}, \sigma)$, $(U', \overline{u}', \sigma')$
d\'efinissent le m\^eme \'el\'ement de la fibre s'il existe un troisi\`eme
voisinage \'etale $(U'', \overline{u}'')$ et des morphismes de voisinages
\'etales $h : (U'', \overline{u}'') \to (U, \overline{u})$,
$h' : (U'', \overline{u}'') \to (U', \overline{u}')$ tels que
$h^*\sigma = (h')^*\sigma'$ dans $\mathcal{F}(U'')$. Voir
Categories, Section \ref{categories-section-directed-colimits}.

\begin{lemma}
\label{lemma-stalk-gives-point}
Soient $S$ un sch\'ema et $\overline{s}$ un point g\'eom\'etrique de $S$.
Consid\'erons le foncteur
\begin{align*}
u : S_\etale & \longrightarrow \textit{Sets}, \\
U & \longmapsto
|U_{\overline{s}}|
=
\{\overline{u} \text{ tel que }(U, \overline{u})
\text{ soit un voisinage \'etale de }\overline{s}\}.
\end{align*}
Ici, $|U_{\overline{s}}|$ d\'esigne l'ensemble sous-jacent \`a la fibre g\'eom\'etrique.
Alors $u$ d\'efinit un point $p$ du site $S_\etale$
(Sites, Definition \ref{sites-definition-point})
et le foncteur fibre associ\'e $\mathcal{F} \mapsto \mathcal{F}_p$
(Sites, Equation \ref{sites-equation-stalk})
est le foncteur $\mathcal{F} \mapsto \mathcal{F}_{\overline{s}}$
d\'efini ci-dessus.
\end{lemma}

\begin{proof}
Dans la preuve du
lemme \ref{lemma-geometric-lift-to-cover},
nous avons vu que le sch\'ema $U_{\overline{s}}$ est une somme disjointe de
sch\'emas isomorphes \`a $\overline{s}$. Nous pouvons donc aussi consid\'erer
$|U_{\overline{s}}|$ comme l'ensemble des points g\'eom\'etriques de $U$ situ\'es au-dessus de
$\overline{s}$, c'est-\`a-dire comme l'ensemble des morphismes
$\overline{u} : \overline{s} \to U$ qui s'inscrivent dans le diagramme de la
d\'efinition \ref{definition-geometric-point}.
Il s'ensuit que $u(S)$ est un singleton, et que
$u(U \times_V W) = u(U) \times_{u(V)} u(W)$
chaque fois que $U \to V$ et $W \to V$ sont des morphismes dans $S_\etale$.
Et, \'etant donn\'e un recouvrement $\{U_i \to U\}_{i \in I}$ dans $S_\etale$,
nous voyons que $\coprod u(U_i) \to u(U)$ est surjectif par le
lemme \ref{lemma-geometric-lift-to-cover}.
Par cons\'equent,
Sites, Proposition \ref{sites-proposition-point-limits}
s'applique, de sorte que $p$ est un point du site $S_\etale$.
Enfin, notre foncteur $\mathcal{F} \mapsto \mathcal{F}_{\overline{s}}$
est donn\'e par exactement la m\^eme limite inductive que le foncteur
$\mathcal{F} \mapsto \mathcal{F}_p$ associ\'e \`a $p$ dans
Sites, Equation \ref{sites-equation-stalk},
ce qui d\'emontre la derni\`ere assertion.
\end{proof}

\begin{remark}
\label{remark-map-stalks}
Soit $S$ un sch\'ema et soient $\overline{s} : \Spec(k) \to S$
et $\overline{s}' : \Spec(k') \to S$ deux points g\'eom\'etriques de
$S$. Un {\it morphisme $a : \overline{s} \to \overline{s}'$ de points g\'eom\'etriques}
est simplement un morphisme $a : \Spec(k) \to \Spec(k')$ tel que
$\overline{s}' \circ a = \overline{s}$. \'Etant donn\'e un tel morphisme, nous obtenons
un foncteur de la cat\'egorie des voisinages \'etales de $\overline{s}'$
vers la cat\'egorie des voisinages \'etales de $\overline{s}$ par la r\`egle
$(U, \overline{u}') \mapsto (U, \overline{u}' \circ a)$. Nous obtenons donc
une application canonique
$$
\mathcal{F}_{\overline{s}'}
=
\colim_{(U, \overline{u}')} \mathcal{F}(U)
\longrightarrow
\colim_{(U, \overline{u})} \mathcal{F}(U)
=
\mathcal{F}_{\overline{s}}
$$
donn\'ee par Categories, Lemma \ref{categories-lemma-functorial-colimit}. En utilisant la
description des \'el\'ements des fibres par des triplets, elle envoie l'\'el\'ement de
$\mathcal{F}_{\overline{s}'}$ repr\'esent\'e par le triplet
$(U, \overline{u}', \sigma)$ sur l'\'el\'ement de $\mathcal{F}_{\overline{s}}$
repr\'esent\'e par le triplet $(U, \overline{u}' \circ a, \sigma)$.
Comme le foncteur ci-dessus est manifestement une \'equivalence, nous concluons que cette
application canonique est un isomorphisme de foncteurs fibres.

\medskip\noindent
Assurons-nous que l'application de fibres correspondant \`a $a$ est orient\'ee
dans le bon sens. Notons que ce qui pr\'ec\`ede signifie, d'apr\`es
Sites, Definition \ref{sites-definition-morphism-points},
que $a$ d\'efinit un morphisme $a : p \to p'$ entre les points $p, p'$ du
site $S_\etale$ associ\'es \`a $\overline{s}, \overline{s}'$ par le
lemme \ref{lemma-stalk-gives-point}. Il existe des morphismes plus g\'en\'eraux de
points (correspondant \`a des sp\'ecialisations de points de $S$), que nous
d\'ecrirons plus loin et qui ne seront pas des isomorphismes, voir la section
\ref{section-specialization}.
\end{remark}

\begin{lemma}
\label{lemma-stalk-exact}
Soit $S$ un sch\'ema. Soit $\overline{s}$ un point g\'eom\'etrique de $S$.
\begin{enumerate}
\item Le foncteur fibre
$\textit{PAb}(S_\etale) \to \textit{Ab}$,
$\mathcal{F}  \mapsto  \mathcal{F}_{\overline{s}}$
est exact.
\item On a $(\mathcal{F}^\#)_{\overline{s}} = \mathcal{F}_{\overline{s}}$
pour tout pr\'efaisceau d'ensembles $\mathcal{F}$ sur $S_\etale$.
\item Le foncteur
$\textit{Ab}(S_\etale) \to \textit{Ab}$,
$\mathcal{F} \mapsto \mathcal{F}_{\overline{s}}$ est exact.
\item De m\^eme, les foncteurs
$\textit{PSh}(S_\etale) \to \textit{Sets}$ et
$\Sh(S_\etale) \to \textit{Sets}$ donn\'es par le foncteur fibre
$\mathcal{F} \mapsto \mathcal{F}_{\overline{s}}$ sont exacts (voir
Categories, Definition \ref{categories-definition-exact})
et commutent aux limites inductives arbitraires.
\end{enumerate}
\end{lemma}

\begin{proof}
Avant d'indiquer comment le d\'emontrer par des arguments directs,
notons que le r\'esultat d\'ecoule de la th\'eorie g\'en\'erale de
Modules on Sites, Section \ref{sites-modules-section-stalks}.
Cela est vrai parce que $\mathcal{F} \mapsto \mathcal{F}_{\overline{s}}$
provient d'un point du petit site \'etale de $S$, voir le
lemme \ref{lemma-stalk-gives-point}.
Nous ne donnerons une preuve directe que de (1), (2) et (3), et omettrons
une preuve directe de (4).

\medskip\noindent
L'exactitude comme foncteur sur $\textit{PAb}(S_\etale)$ d\'ecoule formellement du
fait que les limites inductives filtrantes commutent \`a toutes les limites inductives et aux
limites finies. L'identification des fibres dans (2) se fait au moyen de l'application
$$
\kappa :
\mathcal{F}_{\overline{s}}
\longrightarrow
(\mathcal{F}^\#)_{\overline{s}}
$$
induite par le morphisme naturel $\mathcal{F}\to \mathcal{F}^\#$, voir le
th\'eor\`eme \ref{theorem-sheafification}.
Nous affirmons que cette application est un isomorphisme de groupes ab\'eliens. Nous montrerons
l'injectivit\'e et omettrons la preuve de la surjectivit\'e.

\medskip\noindent
Soit $\sigma\in \mathcal{F}_{\overline{s}}$.
Il existe un voisinage \'etale
$(U, \overline{u})\to (S, \overline{s})$ tel que $\sigma$ soit l'image d'une
section $s \in \mathcal{F}(U)$. Si $\kappa(\sigma) = 0$ dans
$(\mathcal{F}^\#)_{\overline{s}}$, alors il existe un morphisme de voisinages \'etales
$(U', \overline{u}')\to (U, \overline{u})$ tel que
$s|_{U'}$ soit nul dans $\mathcal{F}^\#(U')$. Il s'ensuit qu'il
existe un recouvrement \'etale
$\{U_i'\to U'\}_{i\in I}$ tel que $s|_{U_i'}=0$ dans
$\mathcal{F}(U_i')$ pour tout $i$. Par le lemme \ref{lemma-geometric-lift-to-cover},
il existe $i \in I$ et un morphisme
$\overline{u}_i': \overline{s} \to U_i'$ tels que
$(U_i', \overline{u}_i') \to (U', \overline{u}')\to (U, \overline{u})$
soient des morphismes de voisinages \'etales. Ainsi $\sigma = 0$
puisque $(U_i', \overline{u}_i') \to (U, \overline{u})$
est un morphisme de voisinages \'etales tel que
l'on ait $s|_{U'_i}=0$. Cela d\'emontre que $\kappa$ est injectif.

\medskip\noindent
Pour montrer que le foncteur $\textit{Ab}(S_\etale) \to \textit{Ab}$ est
exact, consid\'erons une suite exacte courte quelconque dans $\textit{Ab}(S_\etale)$ :
$
0\to \mathcal{F}\to \mathcal{G}\to \mathcal H \to 0.
$
Elle nous donne la suite exacte de pr\'efaisceaux
$$
0 \to \mathcal{F} \to \mathcal{G} \to \mathcal H \to
\mathcal H/^p\mathcal{G} \to 0,
$$
o\`u $/^p$ d\'esigne le quotient dans $\textit{PAb}(S_\etale)$.
En prenant les fibres en
$\overline{s}$, nous voyons que $(\mathcal H /^p\mathcal{G})_{\bar{s}} =
(\mathcal H /\mathcal{G})_{\bar{s}} = 0$, puisque le faisceau associ\'e \`a
$\mathcal H/^p\mathcal{G}$ est $0$.
Par cons\'equent,
$$
0\to \mathcal{F}_{\overline{s}} \to \mathcal{G}_{\overline{s}} \to
\mathcal{H}_{\overline{s}} \to 0 = (\mathcal H/^p\mathcal{G})_{\overline{s}}
$$
est exacte, puisque prendre les fibres est exact comme foncteur sur les pr\'efaisceaux.
\end{proof}

\begin{theorem}
\label{theorem-exactness-stalks}
Soit $S$ un sch\'ema.
Un morphisme $a : \mathcal{F} \to \mathcal{G}$ de faisceaux d'ensembles est injectif
(resp.\ surjectif) si et seulement si le morphisme sur les fibres
$a_{\overline{s}} : \mathcal{F}_{\overline{s}} \to \mathcal{G}_{\overline{s}}$
est injectif (resp.\ surjectif) pour tout point g\'eom\'etrique de $S$.
Une suite de faisceaux ab\'eliens sur $S_\etale$ est exacte
si et seulement si elle est exacte sur toutes les fibres aux points g\'eom\'etriques de $S$.
\end{theorem}

\begin{proof}
La n\'ecessit\'e de l'exactitude sur les fibres d\'ecoule du
lemme \ref{lemma-stalk-exact}.
R\'eciproquement, il suffit de montrer qu'un morphisme de faisceaux est surjectif
(respectivement injectif) si et seulement s'il est surjectif (respectivement
injectif) sur toutes les fibres. Nous le d\'emontrons dans le cas de la surjectivit\'e et omettons
la preuve dans le cas de l'injectivit\'e.

\medskip\noindent
Soit $\alpha : \mathcal{F} \to \mathcal{G}$ un morphisme de faisceaux tel
que $\mathcal{F}_{\overline{s}} \to \mathcal{G}_{\overline{s}}$
soit surjectif pour tout point g\'eom\'etrique. Fixons
$U \in \Ob(S_\etale)$
et $s \in \mathcal{G}(U)$. Pour tout $u \in U$, choisissons un
$\overline{u} \to U$ situ\'e au-dessus de $u$ et un voisinage \'etale
$(V_u , \overline{v}_u) \to (U, \overline{u})$ tel que
$s|_{V_u} = \alpha(s_{V_u})$ pour un certain
$s_{V_u} \in \mathcal{F}(V_u)$.
C'est possible puisque $\alpha$ est surjectif sur les
fibres. Alors $\{V_u \to U\}_{u \in U}$
est un recouvrement \'etale sur lequel les restrictions de $s$
appartiennent \`a l'image du morphisme $\alpha$.
Ainsi, $\alpha$ est surjectif, voir
Sites, Section \ref{sites-section-sheaves-injective}.
\end{proof}

\begin{remarks}
\label{remarks-enough-points}
Quelques remarques sur les points des sites g\'eom\'etriques.
\begin{enumerate}
\item Le th\'eor\`eme \ref{theorem-exactness-stalks} dit que la famille de points
de $S_\etale$ donn\'ee par les points g\'eom\'etriques de $S$
(lemme \ref{lemma-stalk-gives-point}) est conservative, voir
Sites, Definition \ref{sites-definition-enough-points}.
En particulier, $S_\etale$ a assez de points.
\item Supposons que $\mathcal{F}$ soit un faisceau sur le grand site \'etale
\label{item-stalks-big}
de $S$. Soit $T \to S$ un objet du grand site \'etale de $S$,
et soit $\overline{t}$ un point g\'eom\'etrique de $T$. Nous d\'efinissons alors
$\mathcal{F}_{\overline{t}}$ comme la fibre
de la restriction $\mathcal{F}|_{T_\etale}$ de $\mathcal{F}$
au petit site \'etale de $T$. Autrement dit, nous pouvons d\'efinir
la fibre de $\mathcal{F}$ en tout point g\'eom\'etrique de tout
sch\'ema $T/S \in \Ob((\Sch/S)_\etale)$.
\item Le grand site \'etale de $S$ a lui aussi assez de points, en
consid\'erant tous les points g\'eom\'etriques de tous les objets de ce site, voir
(\ref{item-stalks-big}).
\end{enumerate}
\end{remarks}

\noindent
Le lemme suivant devrait \^etre omis en premi\`ere lecture.

\begin{lemma}
\label{lemma-points-small-etale-site}
Soit $S$ un sch\'ema.
\begin{enumerate}
\item Soit $p$ un point du petit site \'etale
$S_\etale$ de $S$ donn\'e par un foncteur
$u : S_\etale \to \textit{Sets}$.
Alors il existe un point g\'eom\'etrique $\overline{s}$ de $S$ tel que
$p$ soit isomorphe au point de $S_\etale$ associ\'e \`a
$\overline{s}$ dans le
lemme \ref{lemma-stalk-gives-point}.
\item Soit $p : \Sh(pt) \to \Sh(S_\etale)$ un point
du petit topos \'etale de $S$. Alors $p$ provient d'un point g\'eom\'etrique
de $S$, c'est-\`a-dire que le foncteur fibre $\mathcal{F} \mapsto \mathcal{F}_p$
est isomorphe \`a un foncteur fibre tel que d\'efini dans la
d\'efinition \ref{definition-stalk}.
\end{enumerate}
\end{lemma}

\begin{proof}
D'apr\`es
Sites, Lemma \ref{sites-lemma-point-site-topos},
il existe une correspondance bijective entre les points du site et les points
du topos associ\'e ; il suffit donc de d\'emontrer (1).
D'apr\`es
Sites, Proposition \ref{sites-proposition-point-limits},
le foncteur $u$ a les propri\'et\'es suivantes :
(a) $u(S) = \{*\}$, (b) $u(U \times_V W) = u(U) \times_{u(V)} u(W)$, et
(c) si $\{U_i \to U\}$ est un recouvrement \'etale, alors
$\coprod u(U_i) \to u(U)$ est surjectif.
En particulier, si $U' \subset U$ est un sous-sch\'ema ouvert,
alors $u(U') \subset u(U)$. De plus, d'apr\`es
Sites, Lemma \ref{sites-lemma-point-site-topos},
nous pouvons \'ecrire $u(U) = p^{-1}(h_U^\#)$ ; autrement dit, $u(U)$ est la
fibre du faisceau repr\'esentable $h_U$. Si
$U = V \amalg W$, alors nous voyons que $h_U = (h_V \amalg h_W)^\#$ et obtenons
$u(U) = u(V) \amalg u(W)$ puisque $p^{-1}$ est exact.

\medskip\noindent
Consid\'erons la restriction de $u$ \`a $S_{Zar}$. D'apr\`es
Sites, Examples \ref{sites-example-point-topological} et
\ref{sites-example-point-topology},
il existe un unique point $s \in S$ tel que, pour $S' \subset S$ ouvert, nous
ayons $u(S') = \{*\}$ si $s \in S'$ et $u(S') = \emptyset$ si $s \not \in S'$.
Notons que si $\varphi : U \to S$ est un objet de $S_\etale$, alors
$\varphi(U) \subset S$ est ouvert (voir la
proposition \ref{proposition-etale-morphisms})
et $\{U \to \varphi(U)\}$ est un recouvrement \'etale. Nous concluons donc que
$u(U) = \emptyset \Leftrightarrow s \not \in \varphi(U)$.

\medskip\noindent
Choisissons un point g\'eom\'etrique $\overline{s} : \overline{s} \to S$ au-dessus de $s$, voir la
d\'efinition \ref{definition-geometric-point}
pour l'abus de notation usuel. Supposons que $\varphi : U \to S$ soit un objet
de $S_\etale$ avec $U$ affine. Notons que $\varphi$ est s\'epar\'e et
que la fibre $U_s$ de $\varphi$ au-dessus de $s$ est un sch\'ema affine sur
$\Spec(\kappa(s))$ et qu'elle est le spectre d'un produit fini
d'extensions finies s\'eparables $k_i$ de $\kappa(s)$. Nous pouvons donc appliquer
\'Etale Morphisms, Lemma \ref{etale-lemma-etale-etale-local-technical}
pour obtenir un voisinage \'etale $(V, \overline{v})$ de $(S, \overline{s})$
tel que
$$
U \times_S V = U_1 \amalg \ldots \amalg U_n \amalg W
$$
o\`u $U_i \to V$ est un isomorphisme et $W$ n'a aucun point au-dessus de
$\overline{v}$. Nous concluons donc que
$$
u(U) \times u(V) =
u(U \times_S V) =
u(U_1) \amalg \ldots \amalg u(U_n) \amalg u(W)
$$
et, bien s\^ur, aussi $u(U_i) = u(V)$. Apr\`es avoir un peu r\'etr\'eci $V$, nous pouvons
supposer que $V$ a exactement un point au-dessus de $s$ et que, par cons\'equent, $W$ n'a aucun
point au-dessus de $s$. D'apr\`es ce qui pr\'ec\`ede, cela donne alors $u(W) = \emptyset$.
Nous obtenons donc
$$
u(U) \times u(V) =
u(U_1) \amalg \ldots \amalg u(U_n) =
\coprod\nolimits_{i = 1, \ldots, n} u(V)
$$
o\`u les signes d'\'egalit\'e sont compatibles avec les applications vers l'ensemble $u(V)$.
Notons que $u(V) \not = \emptyset$ puisque $s$ est dans l'image de $V \to S$.
En particulier, nous voyons que, dans cette situation, $u(U)$ est un
ensemble fini \`a $n$ \'el\'ements : en effet, le membre de gauche a pour fibre $u(U)$ au-dessus de
tout \'el\'ement de $u(V)$ et le membre de droite a une fibre de cardinal $n$.

\medskip\noindent
Consid\'erons la limite
$$
\lim_{(V, \overline{v})} u(V)
$$
sur la cat\'egorie des voisinages \'etales $(V, \overline{v})$ de
$\overline{s}$. Il est clair que nous obtenons la m\^eme valeur en prenant
la limite sur la sous-cat\'egorie des $(V, \overline{v})$ tels que $V$ soit affine.
D'apr\`es le paragraphe pr\'ec\'edent (appliqu\'e en \'echangeant les r\^oles de $V$ et $U$),
nous voyons que, dans ce cas, $u(V)$ est toujours un ensemble fini non vide.
De plus, la cat\'egorie d'indexation est cofiltrante, voir le
lemme \ref{lemma-cofinal-etale}.
Par cons\'equent, d'apr\`es
Categories, Section \ref{categories-section-codirected-limits},
la limite est non vide. Choisissons un \'el\'ement $x$ de cette limite.
Cela signifie que nous obtenons un $x_{V, \overline{v}} \in u(V)$ pour
chaque voisinage \'etale $(V, \overline{v})$ de $(S, \overline{s})$
tel que, pour tout morphisme de voisinages \'etales
$\varphi : (V', \overline{v}') \to (V, \overline{v})$, on ait
$u(\varphi)(x_{V', \overline{v}'}) = x_{V, \overline{v}}$.

\medskip\noindent
Nous utiliserons le choix de $x$ pour construire une application bijective fonctorielle
$$
c : |U_{\overline{s}}| \longrightarrow u(U)
$$
pour $U \in \Ob(S_\etale)$, ce qui conclura la preuve. Voir le
lemme \ref{lemma-stalk-gives-point}
et sa preuve pour une description de $|U_{\overline{s}}|$.
Nous affirmons d'abord qu'il suffit de construire l'application pour $U$ affine.
Nous omettons la preuve de cette affirmation.
Supposons que $U \to S$ soit un objet de $S_\etale$, avec $U$ affine, et soit
$\overline{u} : \overline{s} \to U$ un \'el\'ement de $|U_{\overline{s}}|$.
Choisissons $(V, \overline{v})$ tel que $U \times_S V$ se d\'ecompose
comme dans le troisi\`eme paragraphe de la preuve.
Alors le couple $(\overline{u}, \overline{v})$ donne un point g\'eom\'etrique de
$U \times_S V$ au-dessus de $\overline{v}$ et d\'etermine l'une des
composantes $U_i$ de $U \times_S V$. Plus pr\'ecis\'ement, il existe
une section $\sigma : V \to U \times_S V$ de la projection $\text{pr}_U$
telle que $(\overline{u}, \overline{v}) = \sigma \circ \overline{v}$. Posons
$c(\overline{u}) = u(\text{pr}_U)(u(\sigma)(x_{V, \overline{v}})) \in u(U)$.
Nous devons v\'erifier que cela est ind\'ependant du choix de $(V, \overline{v})$. D'apr\`es le
lemme \ref{lemma-cofinal-etale},
la cat\'egorie des voisinages \'etales est cofiltrante.
Il suffit donc de montrer
que, \'etant donn\'es un morphisme de voisinages \'etales
$\varphi : (V', \overline{v}') \to (V, \overline{v})$ et le choix d'une
section $\sigma' : V' \to U \times_S V'$ de la projection telle que
$(\overline{u}, \overline{v'}) = \sigma' \circ \overline{v}'$,
on a $u(\sigma')(x_{V', \overline{v}'}) = u(\sigma)(x_{V, \overline{v}})$.
Consid\'erons le diagramme
$$
\xymatrix{
V' \ar[d]^{\sigma'} \ar[r]_\varphi & V \ar[d]^\sigma \\
U \times_S V' \ar[r]^{1 \times \varphi} &
U \times_S V
}
$$
Il se peut que ce diagramme ne soit pas commutatif. La raison en est
que les sch\'emas $V'$ et $V$ ne sont peut-\^etre pas connexes, et donc que
les d\'ecompositions utilis\'ees ci-dessus pour construire $\sigma'$ et $\sigma$ ne sont
peut-\^etre pas uniques. Mais nous savons que
$\sigma \circ \varphi \circ \overline{v}' =
(1 \times \varphi) \circ \sigma' \circ \overline{v}'$
par construction. Par cons\'equent, puisque $U \times_S V$ est \'etale sur $S$,
il existe un voisinage ouvert
$V'' \subset V'$ de $\overline{v'}$ tel que le diagramme
soit commutatif apr\`es restriction \`a $V''$, voir
Morphisms, Lemma \ref{morphisms-lemma-value-at-one-point}.
Cela signifie que nous pouvons compl\'eter le diagramme ci-dessus en
$$
\xymatrix{
V'' \ar[r] \ar[d]^{\sigma'|_{V''}} &
V' \ar[d]^{\sigma'} \ar[r]_\varphi &
V \ar[d]^\sigma \\
U \times_S V'' \ar[r] &
U \times_S V' \ar[r]^{1 \times \varphi} &
U \times_S V
}
$$
de sorte que le carr\'e de gauche et le rectangle ext\'erieur soient commutatifs.
Comme $u$ est un foncteur, cela implique que
$x_{V'', \overline{v}'}$ s'envoie sur le m\^eme \'el\'ement de
$u(U \times_S V)$ quel que soit le chemin suivi dans le
diagramme. D'autre part, il s'envoie sur les \'el\'ements
$x_{V', \overline{v}'}$ et $x_{V, \overline{v}}$ dans
$u(V')$ et $u(V)$. Cela implique l'\'egalit\'e voulue
$u(\sigma')(x_{V', \overline{v}'}) = u(\sigma)(x_{V, \overline{v}})$.

\medskip\noindent
De mani\`ere analogue, on d\'emontre que la construction
$c : |U_{\overline{s}}| \to u(U)$ est fonctorielle en $U$ ;
nous omettons les d\'etails. Enfin, d'apr\`es les r\'esultats du
troisi\`eme paragraphe, il est clair que l'application $c$ est bijective,
ce qui termine la preuve du lemme.
\end{proof}







\section{Points dans d'autres topologies}
\label{section-points-topologies}

\noindent
Dans cette section, nous discutons bri\`evement de l'existence de points pour certains
sites autres que le site \'etale d'un sch\'ema. Nous renvoyons \`a
Sites, Section \ref{sites-section-sites-enough-points}
et \`a
Topologies, Section \ref{topologies-section-procedure} et suivantes
pour la terminologie employ\'ee dans cette section.
Tous les sites g\'eom\'etriques ont assez de points.

\begin{lemma}
\label{lemma-points-fppf}
Soit $S$ un sch\'ema. Tous les sites suivants ont assez de points :
$S_{affine, Zar}$, $S_{Zar}$, $S_{affine, \etale}$, $S_\etale$,
$(\Sch/S)_{Zar}$, $(\textit{Aff}/S)_{Zar}$,
$(\Sch/S)_\etale$, $(\textit{Aff}/S)_\etale$,
$(\Sch/S)_{smooth}$, $(\textit{Aff}/S)_{smooth}$,
$(\Sch/S)_{syntomic}$, $(\textit{Aff}/S)_{syntomic}$,
$(\Sch/S)_{fppf}$ et $(\textit{Aff}/S)_{fppf}$.
\end{lemma}

\begin{proof}
Pour chacun des grands sites, le topos associ\'e est \'equivalent au
topos d\'efini par le site $(\textit{Aff}/S)_\tau$, voir
Topologies, Lemmas \ref{topologies-lemma-affine-big-site-Zariski},
\ref{topologies-lemma-affine-big-site-etale},
\ref{topologies-lemma-affine-big-site-smooth},
\ref{topologies-lemma-affine-big-site-syntomic} et
\ref{topologies-lemma-affine-big-site-fppf}.
Le r\'esultat pour les sites $(\textit{Aff}/S)_\tau$ d\'ecoule imm\'ediatement
du r\'esultat de Deligne,
Sites, Lemma \ref{sites-lemma-criterion-points}.

\medskip\noindent
Le r\'esultat pour $S_{Zar}$ est clair. Celui pour $S_{affine, Zar}$
d\'ecoule du r\'esultat de Deligne. Le r\'esultat pour $S_\etale$
d\'ecoule soit (de la preuve) du
th\'eor\`eme \ref{theorem-exactness-stalks},
soit de Topologies, Lemma \ref{topologies-lemma-alternative}
et du r\'esultat de Deligne appliqu\'e \`a $S_{affine, \etale}$.
\end{proof}

\noindent
Le lemme ci-dessus garantit l'existence de points, mais ne nous
dit pas \`a quoi ils ressemblent. Nous pouvons construire explicitement
{\it certains} points comme suit.
Supposons que $\overline{s} : \Spec(k) \to S$ soit un point
g\'eom\'etrique, avec $k$ alg\'ebriquement clos. Consid\'erons le foncteur
$$
u : (\Sch/S)_{fppf} \longrightarrow \textit{Sets},
\quad
u(U) = U(k) = \Mor_S(\Spec(k), U).
$$
Notons que $U \mapsto U(k)$ commute aux limites finies puisque
$S(k) = \{\overline{s}\}$ et
$(U_1 \times_U U_2)(k) = U_1(k) \times_{U(k)} U_2(k)$.
De plus, si $\{U_i \to U\}$ est un recouvrement fppf, alors
$\coprod U_i(k) \to U(k)$ est surjectif.
D'apr\`es
Sites, Proposition \ref{sites-proposition-point-limits},
nous voyons que $u$ d\'efinit un point $p$ de $(\Sch/S)_{fppf}$ dont les
fibres sont
$$
\mathcal{F}_p = \colim_{(U, x)} \mathcal{F}(U)
$$
o\`u la limite inductive porte, comme d'habitude, sur les couples $U \to S$, $x \in U(k)$.
Mais... cette cat\'egorie a un objet initial, \`a savoir
$(\Spec(k), \text{id})$ ; nous voyons donc que
$$
\mathcal{F}_p = \mathcal{F}(\Spec(k))
$$
ce qui n'est gu\`ere int\'eressant ! En fait, en g\'en\'eral, ces points ne
forment pas une famille conservative de points. Un type de point plus int\'eressant
est d\'ecrit dans la remarque suivante.

\begin{remark}
\label{remark-points-fppf-site}
\begin{reference}
Ceci est discut\'e dans \cite{Schroeer}.
\end{reference}
Soit $S = \Spec(A)$ un sch\'ema affine. Soit $(p, u)$ un point du
site $(\textit{Aff}/S)_{fppf}$, voir
Sites, Sections \ref{sites-section-points} et
\ref{sites-section-construct-points}. Soit $B = \mathcal{O}_p$ la fibre
du faisceau structural au point $p$. Rappelons que
$$
B = \colim_{(U, x)} \mathcal{O}(U) =
\colim_{(\Spec(C), x_C)} C
$$
o\`u $x_C \in u(\Spec(C))$. Il peut arriver que
$\Spec(B)$ soit un objet de $(\textit{Aff}/S)_{fppf}$
et qu'il existe un \'el\'ement $x_B \in u(\Spec(B))$ qui s'envoie sur
le syst\`eme compatible $x_C$. Dans ce cas, le syst\`eme de voisinages
a un objet initial et il s'ensuit que
$\mathcal{F}_p = \mathcal{F}(\Spec(B))$ pour tout faisceau $\mathcal{F}$
sur $(\textit{Aff}/S)_{fppf}$. Il est imm\'ediat
de voir que, si $\mathcal{F} \mapsto \mathcal{F}(\Spec(B))$ d\'efinit un point
de $\Sh((\textit{Aff}/S)_{fppf})$, alors
$B$ doit \^etre une $A$-alg\`ebre locale telle que, pour tout morphisme d'anneaux
fid\`element plat et de pr\'esentation finie $B \to B'$, il existe une section $B' \to B$.
R\'eciproquement, pour toute $A$-alg\`ebre $B$ de ce type, le foncteur
$\mathcal{F} \mapsto \mathcal{F}(\Spec(B))$ est le foncteur fibre
d'un point. Nous omettons les d\'etails. On ne sait pas clairement \`a quoi ressemble un point quelconque du
site $(\textit{Aff}/S)_{fppf}$.
\end{remark}











\section{Supports des faisceaux ab\'eliens}
\label{section-support}

\noindent
Nous commen\c{c}ons par les supports des sections locales.

\begin{lemma}
\label{lemma-support-subsheaf-final}
Soit $S$ un sch\'ema. Soit $\mathcal{F}$ un sous-faisceau de l'objet
final du topos \'etale de $S$ (voir
Sites, Example \ref{sites-example-singleton-sheaf}).
Alors il existe un unique ouvert
$W \subset S$ tel que $\mathcal{F} = h_W$.
\end{lemma}

\begin{proof}
La condition signifie que $\mathcal{F}(U)$ est un singleton ou est
vide pour tout $\varphi : U \to S$ dans $\Ob(S_\etale)$.
En particulier, les sections locales se recollent toujours. Si
$\mathcal{F}(U) \not = \emptyset$, alors
$\mathcal{F}(\varphi(U)) \not = \emptyset$ parce que
$\{\varphi : U \to \varphi(U)\}$ est un recouvrement.
Nous pouvons donc prendre
$W = \bigcup_{\varphi : U \to S, \mathcal{F}(U) \not = \emptyset} \varphi(U)$.
\end{proof}

\begin{lemma}
\label{lemma-zero-over-image}
Soit $S$ un sch\'ema.
Soit $\mathcal{F}$ un faisceau ab\'elien sur $S_\etale$.
Soit $\sigma \in \mathcal{F}(U)$ une section locale.
Il existe un ouvert $W \subset U$ tel que
\begin{enumerate}
\item $W \subset U$ soit le plus grand ouvert de Zariski de $U$ tel
que $\sigma|_W = 0$,
\item pour tout $\varphi : V \to U$ dans $S_\etale$, on ait
$$
\sigma|_V = 0 \Leftrightarrow \varphi(V) \subset W,
$$
\item pour tout point g\'eom\'etrique $\overline{u}$ de $U$, on ait
$$
(U, \overline{u}, \sigma) = 0\text{ dans }\mathcal{F}_{\overline{s}}
\Leftrightarrow
\overline{u} \in W
$$
o\`u $\overline{s} = (U \to S) \circ \overline{u}$.
\end{enumerate}
\end{lemma}

\begin{proof}
Comme $\mathcal{F}$ est un faisceau pour la topologie \'etale, la restriction de
$\mathcal{F}$ \`a $U_{Zar}$ est un faisceau sur $U$ pour la topologie de Zariski.
Il existe donc un ouvert de Zariski $W$ ayant la propri\'et\'e (1), voir
Modules, Lemma \ref{modules-lemma-support-section-closed}. Soit
$\varphi : V \to U$ un morphisme de $S_\etale$. Notons que
$\varphi(V) \subset U$ est un ouvert et que
$\{V \to \varphi(V)\}$ est un recouvrement \'etale. Ainsi, si
$\sigma|_V = 0$, alors la condition de faisceau pour $\mathcal{F}$ montre
que $\sigma|_{\varphi(V)} = 0$. Cela d\'emontre (2).
Pour d\'emontrer (3), nous devons montrer que si $(U, \overline{u}, \sigma)$
d\'efinit l'\'el\'ement nul de $\mathcal{F}_{\overline{s}}$, alors
$\overline{u} \in W$. Cela est vrai parce que l'hypoth\`ese signifie
qu'il existe un morphisme de voisinages \'etales
$(V, \overline{v}) \to (U, \overline{u})$ tel que
$\sigma|_V = 0$. D'apr\`es (2), nous voyons donc que $V \to U$ est \`a valeurs dans $W$, et
par cons\'equent que $\overline{u} \in W$.
\end{proof}

\noindent
Soit $S$ un sch\'ema. Soit $s \in S$.
Soit $\mathcal{F}$ un faisceau sur $S_\etale$. D'apr\`es la
remarque \ref{remark-map-stalks},
la classe d'isomorphisme de la fibre du faisceau $\mathcal{F}$
en un point g\'eom\'etrique au-dessus de $s$ est bien d\'efinie.

\begin{definition}
\label{definition-support}
Soit $S$ un sch\'ema.
Soit $\mathcal{F}$ un faisceau ab\'elien sur $S_\etale$.
\begin{enumerate}
\item Le {\it support de $\mathcal{F}$} est l'ensemble des
points $s \in S$ tels que $\mathcal{F}_{\overline{s}} \not = 0$
pour tout (ou un) point g\'eom\'etrique $\overline{s}$ au-dessus de $s$.
\item Soit $\sigma \in \mathcal{F}(U)$ une section.
Le {\it support de $\sigma$} est le ferm\'e $U \setminus W$, o\`u
$W \subset U$ est le plus grand ouvert de $U$ sur lequel $\sigma$
est nulle (voir le
lemme \ref{lemma-zero-over-image}).
\end{enumerate}
\end{definition}

\noindent
En g\'en\'eral, le support d'un faisceau ab\'elien n'est pas ferm\'e.
Par exemple, supposons que $S = \mathbf{A}^1_{\mathbf{C}}$.
Soit $i_t : \Spec(\mathbf{C}) \to S$ l'inclusion du
point $t \in \mathbf{C}$.
Nous verrons plus loin que $\mathcal{F}_t = i_{t, *}(\mathbf{Z}/2\mathbf{Z})$
est un faisceau ab\'elien dont le support est exactement $\{t\}$, voir la
section \ref{section-closed-immersions}.
Alors
$$
\bigoplus\nolimits_{n \in \mathbf{N}} \mathcal{F}_n
$$
est un faisceau ab\'elien de support $\{1, 2, 3, \ldots\} \subset S$.
Cela est vrai parce que prendre les fibres commute aux limites inductives, voir le
lemme \ref{lemma-stalk-exact}.
Nous avons ainsi un exemple de faisceau ab\'elien dont le support n'est pas ferm\'e.
Voici quelques faits fondamentaux sur les supports des faisceaux et des sections.

\begin{lemma}
\label{lemma-support-section-closed}
Soit $S$ un sch\'ema.
Soit $\mathcal{F}$ un faisceau ab\'elien sur $S_\etale$.
Soient $U \in \Ob(S_\etale)$ et $\sigma \in \mathcal{F}(U)$.
\begin{enumerate}
\item Le support de $\sigma$ est ferm\'e dans $U$.
\item Le support de $\sigma + \sigma'$ est contenu dans la r\'eunion des
supports de $\sigma, \sigma' \in \mathcal{F}(U)$.
\item Si $\varphi : \mathcal{F} \to \mathcal{G}$ est un morphisme de
faisceaux ab\'eliens sur $S_\etale$, alors le support de $\varphi(\sigma)$
est contenu dans le support de $\sigma \in \mathcal{F}(U)$.
\item Le support de $\mathcal{F}$ est la r\'eunion des images des
supports de toutes les sections locales de $\mathcal{F}$.
\item Si $\mathcal{F} \to \mathcal{G}$ est surjectif, alors le support
de $\mathcal{G}$ est un sous-ensemble du support de $\mathcal{F}$.
\item Si $\mathcal{F} \to \mathcal{G}$ est injectif, alors le support
de $\mathcal{F}$ est un sous-ensemble du support de $\mathcal{G}$.
\end{enumerate}
\end{lemma}

\begin{proof}
L'assertion (1) est vraie par d\'efinition.
Les assertions (2) et (3) sont vraies parce qu'elles le sont pour les restrictions de
$\mathcal{F}$ et $\mathcal{G}$ \`a $U_{Zar}$, voir
Modules, Lemma \ref{modules-lemma-support-section-closed}.
L'assertion (4) est une cons\'equence directe de la partie (3) du
lemme \ref{lemma-zero-over-image}.
Les assertions (5) et (6) d\'ecoulent des autres.
\end{proof}

\begin{lemma}
\label{lemma-support-sheaf-rings-closed}
Le support d'un faisceau d'anneaux sur $S_\etale$ est ferm\'e.
\end{lemma}

\begin{proof}
Cela est vrai parce que (selon nos conventions)
un anneau est $0$ si et seulement si
$1 = 0$ ; par cons\'equent, le support d'un faisceau d'anneaux
est le support de la section unit\'e.
\end{proof}




\section{Anneaux hens\'eliens}
\label{section-henselian-ring}

\noindent
Nous commen\c{c}ons par \'enoncer un th\'eor\`eme qui a d\'ej\`a \'et\'e utilis\'e de nombreuses fois
dans le projet Stacks. Il existe de nombreuses versions de ce r\'esultat ; ici, nous
n'\'enon\c{c}ons que la version alg\'ebrique.

\begin{theorem}
\label{theorem-quasi-finite-etale-locally}
Soient $A\to B$ un morphisme d'anneaux de type fini et $\mathfrak p \subset A$ un
id\'eal premier. Alors il existe un morphisme d'anneaux \'etale $A \to A'$ et un id\'eal premier
$\mathfrak p' \subset A'$ au-dessus de $\mathfrak p$ tels que
\begin{enumerate}
\item
$\kappa(\mathfrak p) = \kappa(\mathfrak p')$,
\item
$ B \otimes_A A' = B_1\times \ldots \times B_r \times C$,
\item
$ A'\to B_i$ est fini et il existe un unique id\'eal premier $q_i\subset B_i$ au-dessus
de $\mathfrak p'$, et
\item toutes les composantes irr\'eductibles de la fibre
$\Spec(C \otimes_{A'} \kappa(\mathfrak p'))$ de $C$ au-dessus de $\mathfrak p'$
sont de dimension au moins 1.
\end{enumerate}
\end{theorem}

\begin{proof}
Voir Algebra, Lemma \ref{algebra-lemma-etale-makes-quasi-finite-finite}, ou
voir \cite[Th\'eor\`eme 18.12.1]{EGA4}. Pour de nombreuses versions en termes de
morphismes de sch\'emas, voir
More on Morphisms, Section \ref{more-morphisms-section-etale-localization}.
\end{proof}

\noindent
Rappelons le lemme de Hensel.
Il existe de nombreuses versions de ce lemme. En voici deux :
\begin{enumerate}
\item[(f)] si $f\in \mathbf{Z}_p[T]$ est unitaire et
$f \bmod p = g_0 h_0$ avec $gcd(g_0, h_0) = 1$, alors $f$ se factorise
sous la forme $f = gh$ avec $\bar g = g_0$ et $\bar h = h_0$,
\item[(r)] si $f \in \mathbf{Z}_p[T]$ est unitaire, $a_0 \in \mathbf{F}_p$,
$\bar f(a_0) =0$ mais $\bar f'(a_0) \neq 0$,
alors il existe $a \in \mathbf{Z}_p$ tel que
$f(a) = 0$ et $\bar a = a_0$.
\end{enumerate}
Les deux versions sont vraies (nous le verrons plus loin). La premi\`ere version
demande des rel\`evements de factorisations en facteurs premiers entre eux,
et la seconde demande des rel\`evements de racines simples
modulo l'id\'eal maximal. Il se trouve que, pour un anneau local g\'en\'eral,
ces deux conditions sont \'equivalentes, et qu'elles sont
\'equivalentes \`a beaucoup d'autres conditions. Nous utilisons la propri\'et\'e de rel\`evement
des racines comme d\'efinition d'un anneau local hens\'elien, car c'est
souvent la plus facile \`a v\'erifier.

%10.01.09
\begin{definition}
\label{definition-henselian}
(Voir Algebra, Definition \ref{algebra-definition-henselian}.)
Un anneau local $(R, \mathfrak m, \kappa)$ est dit
{\it hens\'elien} si, pour tout
$f \in R[T]$ unitaire et tout $a_0 \in \kappa$ tel que
$\bar f(a_0) = 0$ et $\bar f'(a_0) \neq 0$, il existe
un $a \in R$ tel que $f(a) = 0$ et $a \bmod \mathfrak m = a_0$.
\end{definition}

\noindent
Les anneaux locaux complets constituent un bon exemple
d'anneaux locaux hens\'eliens \`a garder \`a l'esprit.
Rappelons
(Algebra, Definition \ref{algebra-definition-complete-local-ring})
qu'un anneau local complet est un anneau local $(R, \mathfrak m)$ tel que
$R \cong \lim_n R/\mathfrak m^n$, c'est-\`a-dire complet et s\'epar\'e
pour la topologie $\mathfrak m$-adique.

\begin{theorem}
\label{theorem-hensel}
Les anneaux locaux complets sont hens\'eliens.
\end{theorem}

\begin{proof}
M\'ethode de Newton. Voir
Algebra, Lemma \ref{algebra-lemma-complete-henselian}.
\end{proof}

\begin{theorem}
\label{theorem-henselian}
Soit $(R, \mathfrak m, \kappa)$ un anneau local. Les propri\'et\'es suivantes sont \'equivalentes :
\begin{enumerate}
\item $R$ est hens\'elien,
\item pour tout $f\in R[T]$ et toute factorisation $\bar f = g_0 h_0$ dans
$\kappa[T]$ avec $\gcd(g_0, h_0)=1$, il existe une factorisation $f = gh$ dans
$R[T]$ avec $\bar g = g_0$ et $\bar h = h_0$,
\item toute $R$-alg\`ebre finie $S$ est isomorphe \`a un produit fini
d'anneaux locaux finis sur $R$,
\item toute $R$-alg\`ebre de type fini $A$ est isomorphe \`a un produit
$A \cong A' \times C$ o\`u $A' \cong A_1 \times \ldots \times A_r$
est un produit de $R$-alg\`ebres locales finies et toutes les composantes
irr\'eductibles de $C \otimes_R \kappa$ sont de dimension au moins 1,
\item si $A$ est une $R$-alg\`ebre \'etale et $\mathfrak n$ un id\'eal maximal de
$A$ au-dessus de $\mathfrak m$ tel que $\kappa \cong A/\mathfrak n$, alors il
existe un isomorphisme $\varphi : A \cong R \times A'$ tel que
$\varphi(\mathfrak n) = \mathfrak m \times A' \subset R \times A'$.
\end{enumerate}
\end{theorem}

\begin{proof}
Il s'agit simplement d'une partie des r\'esultats de
Algebra, Lemma \ref{algebra-lemma-characterize-henselian}.
Notons que la partie (5) ci-dessus correspond \`a la partie (8) de
Algebra, Lemma \ref{algebra-lemma-characterize-henselian},
mais qu'elle est formul\'ee l\'eg\`erement diff\'eremment.
\end{proof}

\begin{lemma}
\label{lemma-finite-over-henselian}
Si $R$ est hens\'elien et $A$ est une $R$-alg\`ebre finie, alors $A$ est un produit
fini d'anneaux locaux hens\'eliens.
\end{lemma}

\begin{proof}
Voir
Algebra, Lemma \ref{algebra-lemma-finite-over-henselian}.
\end{proof}

\begin{definition}
\label{definition-strictly-henselian}
Un anneau local $R$ est dit {\it strictement hens\'elien} s'il est hens\'elien et si son
corps r\'esiduel est s\'eparablement clos.
\end{definition}

\begin{example}
\label{example-powerseries}
Dans le cas $R = \mathbf{C}[[t]]$, les $R$-alg\`ebres \'etales sont des produits finis
de l'extension triviale $R \to R$ et des extensions
$R \to R[X, X^{-1}]/(X^n-t)$.
Ces derni\`eres se factorisent par l'ouvert $D(t) \subset \Spec(R)$, de sorte que tout
recouvrement \'etale peut \^etre raffin\'e par le recouvrement
$\{\text{id} : \Spec(R) \to \Spec(R)\}$. Nous verrons ci-dessous que
c'est un fait assez g\'en\'eral sur les recouvrements \'etales des spectres d'anneaux
hens\'eliens. Cela montrera que la cohomologie \'etale sup\'erieure du spectre d'un anneau
strictement hens\'elien est nulle.
\end{example}

\begin{theorem}
\label{theorem-henselization}
Soient $(R, \mathfrak m, \kappa)$ un anneau local et
$\kappa\subset\kappa^{sep}$ une cl\^oture alg\'ebrique s\'eparable.
Il existe des morphismes locaux plats canoniques d'anneaux $R \to R^h \to R^{sh}$ tels que
\begin{enumerate}
\item $R^h$, $R^{sh}$ soient des limites inductives filtrantes de $R$-alg\`ebres \'etales,
\item $R^h$ soit hens\'elien, $R^{sh}$ soit strictement hens\'elien,
\item $\mathfrak m R^h$ (resp.\ $\mathfrak m R^{sh}$) soit
l'id\'eal maximal de $R^h$ (resp.\ de $R^{sh}$), et
\item $\kappa = R^h/\mathfrak m R^h$, et
$\kappa^{sep} = R^{sh}/\mathfrak m R^{sh}$ comme extensions de $\kappa$.
\end{enumerate}
\end{theorem}

\begin{proof}
La structure de $R^h$ et de $R^{sh}$ est d\'ecrite dans
Algebra, Lemmas \ref{algebra-lemma-henselization} et
\ref{algebra-lemma-strict-henselization}.
\end{proof}

\noindent
Les anneaux construits dans le th\'eor\`eme \ref{theorem-henselization}
sont appel\'es respectivement {\it l'hens\'elisation} et
{\it l'hens\'elisation stricte} de l'anneau local $R$, voir
Algebra, Definition \ref{algebra-definition-henselization}.
De nombreuses propri\'et\'es de $R$ se refl\`etent dans son hens\'elisation (stricte),
voir More on Algebra,
Section \ref{more-algebra-section-permanence-henselization}.




\section{Fibres du faisceau structural}
\label{section-stalks-structure-sheaf}

\noindent
Dans cette section, nous identifions la fibre du faisceau structural en un point
g\'eom\'etrique \`a l'hens\'elisation stricte de l'anneau local au point
``usuel'' correspondant.

\begin{lemma}
\label{lemma-describe-etale-local-ring}
\begin{slogan}
La fibre du faisceau structural d'un sch\'ema
pour la topologie \'etale est l'hens\'elisation stricte.
\end{slogan}
Soit $S$ un sch\'ema.
Soit $\overline{s}$ un point g\'eom\'etrique de $S$ au-dessus de $s \in S$.
Posons $\kappa = \kappa(s)$ et notons
$\kappa \subset \kappa^{sep} \subset \kappa(\overline{s})$
la cl\^oture alg\'ebrique s\'eparable de $\kappa$ dans $\kappa(\overline{s})$.
Alors il existe une identification canonique
$$
(\mathcal{O}_{S, s})^{sh}
\cong
(\mathcal{O}_S)_{\overline{s}}
$$
o\`u le membre de gauche est l'hens\'elisation stricte de l'anneau local
$\mathcal{O}_{S, s}$ telle qu'elle est d\'ecrite dans le
th\'eor\`eme \ref{theorem-henselization},
et le membre de droite est la fibre du faisceau structural
$\mathcal{O}_S$ sur $S_\etale$ au
point g\'eom\'etrique $\overline{s}$.
\end{lemma}

\begin{proof}
Soit $\Spec(A) \subset S$ un voisinage affine de $s$.
Soit $\mathfrak p \subset A$ l'id\'eal premier correspondant \`a $s$.
Avec ces choix, nous avons des isomorphismes canoniques
$\mathcal{O}_{S, s} = A_{\mathfrak p}$ et $\kappa(s) = \kappa(\mathfrak p)$.
Nous avons donc
$\kappa(\mathfrak p) \subset \kappa^{sep} \subset \kappa(\overline{s})$.
Rappelons que
$$
(\mathcal{O}_S)_{\overline{s}} =
\colim_{(U, \overline{u})} \mathcal{O}(U)
$$
o\`u la limite inductive porte sur les voisinages \'etales de $(S, \overline{s})$.
Un syst\`eme cofinal est donn\'e par les voisinages \'etales $(U, \overline{u})$
tels que $U$ soit affine et que $U \to S$ se factorise par $\Spec(A)$.
Autrement dit, nous voyons que
$$
(\mathcal{O}_S)_{\overline{s}} = \colim_{(B, \mathfrak q, \phi)} B
$$
o\`u la limite inductive porte sur les $A$-alg\`ebres \'etales $B$ munies d'un id\'eal premier
$\mathfrak q$ au-dessus de $\mathfrak p$ et d'un
morphisme de $\kappa(\mathfrak p)$-alg\`ebres
$\phi : \kappa(\mathfrak q) \to \kappa(\overline{s})$.
Notons que, puisque $\kappa(\mathfrak q)$ est une extension finie s\'eparable de
$\kappa(\mathfrak p)$, l'image de $\phi$ est contenue dans $\kappa^{sep}$.
Moyennant ces identifications, le r\'esultat du lemme est \'equivalent
au r\'esultat de
Algebra, Lemma \ref{algebra-lemma-strict-henselization-different}.
\end{proof}

\begin{definition}
\label{definition-etale-local-rings}
Soit $S$ un sch\'ema. Soit $\overline{s}$ un point g\'eom\'etrique de $S$
au-dessus du point $s \in S$.
\begin{enumerate}
\item L'{\it anneau local \'etale de $S$ en $\overline{s}$}
est la fibre du faisceau structural $\mathcal{O}_S$ sur $S_\etale$
en $\overline{s}$. Nous appelons parfois cet anneau
{\it l'hens\'elisation stricte de $\mathcal{O}_{S, s}$} relative
au point g\'eom\'etrique $\overline{s}$.
Notation utilis\'ee : $\mathcal{O}_{S, \overline{s}}^{sh}$.
\item L'{\it hens\'elisation de $\mathcal{O}_{S, s}$} d\'esigne
l'hens\'elisation de l'anneau local de $S$ en $s$. Voir
Algebra, Definition \ref{algebra-definition-henselization},
et le
th\'eor\`eme \ref{theorem-henselization}.
Notation : $\mathcal{O}_{S, s}^h$.
\item L'{\it hens\'elisation stricte de $S$ en $\overline{s}$}
est le sch\'ema $\Spec(\mathcal{O}_{S, \overline{s}}^{sh})$.
\item L'{\it hens\'elisation de $S$ en $s$} est le sch\'ema
$\Spec(\mathcal{O}_{S, s}^h)$.
\end{enumerate}
\end{definition}

\noindent
Soit $f : T \to S$ un morphisme de sch\'emas. Soit $\overline{t}$
un point g\'eom\'etrique de $T$ d'image $\overline{s}$ dans $S$.
Soient $t \in T$ et $s \in S$ leurs images.
Nous obtenons alors un diagramme commutatif canonique
$$
\xymatrix{
\Spec(\mathcal{O}^{sh}_{T, \overline{t}}) \ar[r] \ar[d] &
\Spec(\mathcal{O}^h_{T, t}) \ar[r] \ar[d] &
T \ar[d]^f \\
\Spec(\mathcal{O}^{sh}_{S, \overline{s}}) \ar[r] &
\Spec(\mathcal{O}^h_{S, s}) \ar[r] &
S
}
$$
des hens\'elisations et hens\'elisations strictes de $T$ et de $S$. On peut
le d\'emontrer en choisissant des voisinages affines de $t$ et de $s$ et en utilisant
la fonctorialit\'e des hens\'elisations (strictes) donn\'ee par
Algebra, Lemmas \ref{algebra-lemma-henselian-functorial-improve} et
\ref{algebra-lemma-strictly-henselian-functorial-improve}.

\begin{lemma}
\label{lemma-describe-henselization}
Soit $S$ un sch\'ema. Soit $s \in S$. Alors
$$
\mathcal{O}_{S, s}^h =
\colim_{(U, u)} \mathcal{O}(U)
$$
o\`u la limite inductive porte sur la cat\'egorie filtrante des
voisinages \'etales $(U, u)$ de $(S, s)$ tels que
$\kappa(s) = \kappa(u)$.
\end{lemma}

\begin{proof}
Ce lemme est une copie de
More on Morphisms, Lemma \ref{more-morphisms-lemma-describe-henselization}.
\end{proof}

\begin{remark}
\label{remark-henselization-Noetherian}
Soit $S$ un sch\'ema. Soit $s \in S$.
Si $S$ est localement noeth\'erien, alors $\mathcal{O}_{S, s}^h$
est lui aussi noeth\'erien et il a le m\^eme compl\'et\'e :
$$
\widehat{\mathcal{O}_{S, s}} \cong \widehat{\mathcal{O}_{S, s}^h}.
$$
En particulier,
$\mathcal{O}_{S, s} \subset
\mathcal{O}_{S, s}^h \subset
\widehat{\mathcal{O}_{S, s}}$.
L'hens\'elisation de $\mathcal{O}_{S, s}$ est en g\'en\'eral beaucoup
plus petite que son compl\'et\'e et h\'erite de nombre de ses propri\'et\'es.
Par exemple, si $\mathcal{O}_{S, s}$ est r\'eduit, alors
$\mathcal{O}_{S, s}^h$ l'est aussi, mais ce n'est en g\'en\'eral pas vrai du compl\'et\'e.
Ins\'erer ici des r\'ef\'erences futures.
\end{remark}

\begin{lemma}
\label{lemma-etale-site-locally-ringed}
Soit $S$ un sch\'ema. Le petit site \'etale $S_\etale$, muni de
son faisceau structural $\mathcal{O}_S$, est un site localement annel\'e, voir
Modules on Sites, Definition \ref{sites-modules-definition-locally-ringed}.
\end{lemma}

\begin{proof}
Cela r\'esulte du fait que les fibres
$(\mathcal{O}_S)_{\overline{s}} = \mathcal{O}^{sh}_{S, \overline{s}}$ sont
des anneaux locaux, et du fait que $S_\etale$ a assez de points, voir
le lemme \ref{lemma-describe-etale-local-ring},
le th\'eor\`eme \ref{theorem-exactness-stalks},
et les
remarques \ref{remarks-enough-points}.
Voir
Modules on Sites, Lemmas \ref{sites-modules-lemma-locally-ringed-stalk} et
\ref{sites-modules-lemma-ringed-stalk-not-zero}
pour le fait que cela implique que le petit site \'etale est localement annel\'e.
\end{proof}






\section{Fonctorialit\'e du petit topos \'etale}
\label{section-functoriality}

\noindent
Jusqu'ici, nous n'avons pas abord\'e la fonctorialit\'e du site \'etale,
autrement dit ce qui se passe lorsque l'on se donne un morphisme de sch\'emas. Une
discussion formelle pr\'ecise se trouve dans
Topologies, Section \ref{topologies-section-etale}.
Dans cette section et les suivantes, nous examinons bri\`evement ces questions,
plus pr\'ecis\'ement dans le cadre des petits sites \'etales.

\medskip\noindent
Soit $f : X \to Y$ un morphisme de sch\'emas. Nous obtenons un foncteur
\begin{equation}
\label{equation-functorial}
u : Y_\etale \longrightarrow X_\etale, \quad
V/Y \longmapsto X \times_Y V/X.
\end{equation}
Ce foncteur poss\`ede les propri\'et\'es importantes suivantes :
\begin{enumerate}
\item $u(\text{objet final}) = \text{objet final}$,
\item $u$ pr\'eserve les produits fibr\'es,
\item si $\{V_j \to V\}$ est un recouvrement dans $Y_\etale$, alors
$\{u(V_j) \to u(V)\}$ est un recouvrement dans $X_\etale$.
\end{enumerate}
Chacune d'elles se v\'erifie facilement (nous omettons la v\'erification). Par cons\'equent, nous obtenons ce
que l'on appelle un {\it morphisme de sites}
$$
f_{small} : X_\etale \longrightarrow Y_\etale,
$$
voir
Sites, Definition \ref{sites-definition-morphism-sites}
et
Sites, Proposition \ref{sites-proposition-get-morphism}.
Il n'est pas n\'ecessaire de conna\^itre en d\'etail la notion abstraite
pour travailler avec les faisceaux \'etales et la cohomologie \'etale.
Il suffit habituellement de savoir qu'il existe des foncteurs
$f_{small, *}$ (image directe) et $f_{small}^{-1}$ (image inverse)
sur les faisceaux \'etales, et de conna\^itre quelques-unes de leurs propri\'et\'es simples.
Nous examinerons ces propri\'et\'es dans les sections suivantes, mais nous nous
r\'ef\'ererons parfois au cadre plus abstrait pour les d\'emonstrations, car
c'est souvent le cadre naturel pour les \`etablir.


\section{Images directes}
\label{section-direct-image}

\noindent
D\'efinissons l'image directe d'un pr\'efaisceau.

\begin{definition}
\label{definition-direct-image-presheaf}
Soit $f: X\to Y$ un morphisme de sch\'emas.
Soit $\mathcal{F} $ un pr\'efaisceau d'ensembles sur $X_\etale$.
L'{\it image directe}, ou {\it pouss\'ee en avant}, de $\mathcal{F}$
(par $f$) est
$$
f_*\mathcal{F} : Y_\etale^{opp} \longrightarrow \textit{Sets}, \quad
(V/Y) \longmapsto \mathcal{F}(X \times_Y V/X).
$$
Nous \`ecrivons parfois $f_* = f_{small, *}$ pour distinguer ce foncteur des autres
foncteurs d'image directe (comme l'image directe usuelle pour la topologie de Zariski ou $f_{big, *}$).
\end{definition}

\noindent
C'est bien un pr\'efaisceau \'etale, puisque le changement de base d'un morphisme
\'etale est encore \'etale. Une fa\c{c}on plus cat\'egorique de l'exprimer consiste \`a dire que
$f_*\mathcal{F}$ est le foncteur compos\'e $\mathcal{F} \circ u$
o\`u $u$ est d\'efini par l'\'equation (\ref{equation-functorial}). Cela montre
clairement que la construction d\'epend fonctoriellement du pr\'efaisceau
$\mathcal{F}$ et nous obtenons donc un foncteur
$$
f_* = f_{small, *} :
\textit{PSh}(X_\etale)
\longrightarrow
\textit{PSh}(Y_\etale)
$$
Notons que si $\mathcal{F}$ est un pr\'efaisceau de groupes ab\'eliens, alors
$f_*\mathcal{F}$ est aussi un pr\'efaisceau de groupes ab\'eliens et nous obtenons
$$
f_* = f_{small, *} :
\textit{PAb}(X_\etale)
\longrightarrow
\textit{PAb}(Y_\etale)
$$
comme pr\'ec\'edemment (c'est-\`a-dire d\'efini par exactement la m\^eme r\`egle).

\begin{remark}
\label{remark-direct-image-sheaf}
Nous affirmons que l'image directe d'un faisceau est un faisceau.
En effet, si $\{V_j \to V\}$ est un recouvrement \'etale dans $Y_\etale$,
alors $\{X \times_Y V_j \to X \times_Y V\}$ est un recouvrement \'etale dans
$X_\etale$. Par cons\'equent, la condition de faisceau pour $\mathcal{F}$ relativement
\`a $\{X \times_Y V_i \to X \times_Y V\}$
est \`equivalente \`a la condition de faisceau pour $f_*\mathcal{F}$ relativement \`a
$\{V_i \to V\}$. Ainsi, si $\mathcal{F}$ est un faisceau,
$f_*\mathcal{F}$ l'est aussi.
\end{remark}

\begin{definition}
\label{definition-direct-image-sheaf}
Soit $f: X\to Y$ un morphisme de sch\'emas.
Soit $\mathcal{F} $ un faisceau d'ensembles sur $X_\etale$.
L'{\it image directe}, ou {\it pouss\'ee en avant}, de $\mathcal{F}$
(par $f$) est
$$
f_*\mathcal{F} : Y_\etale^{opp} \longrightarrow \textit{Sets}, \quad
(V/Y) \longmapsto \mathcal{F}(X \times_Y V/X)
$$
qui est un faisceau d'apr\`es la
remarque \ref{remark-direct-image-sheaf}.
Nous \`ecrivons parfois $f_* = f_{small, *}$ pour distinguer ce foncteur des autres
foncteurs d'image directe (comme l'image directe usuelle pour la topologie de Zariski ou $f_{big, *}$).
\end{definition}

\noindent
La discussion pr\'ec\'edente s'applique mot pour mot et nous obtenons des foncteurs
$$
f_* = f_{small, *} :
\Sh(X_\etale)
\longrightarrow
\Sh(Y_\etale)
$$
et
$$
f_* = f_{small, *} :
\textit{Ab}(X_\etale)
\longrightarrow
\textit{Ab}(Y_\etale)
$$
\`a nouveau appel\'es {\it foncteurs d'image directe}.

\medskip\noindent
Le foncteur $f_*$ sur les faisceaux ab\'eliens est exact \`a gauche. (Voir
Homology, Section \ref{homology-section-functors}
pour savoir ce que signifie qu'un foncteur entre cat\'egories ab\'eliennes est exact \`a gauche.)
En effet, si
$0 \to \mathcal{F}_1 \to \mathcal{F}_2 \to \mathcal{F}_3$
est exacte sur $X_\etale$, alors, pour tout
$U/X \in \Ob(X_\etale)$,
la suite de groupes ab\'eliens
$0 \to \mathcal{F}_1(U) \to \mathcal{F}_2(U) \to \mathcal{F}_3(U)$
est exacte. Par cons\'equent, pour tout $V/Y \in \Ob(Y_\etale)$,
la suite de groupes ab\'eliens
$0 \to f_*\mathcal{F}_1(V) \to f_*\mathcal{F}_2(V) \to f_*\mathcal{F}_3(V)$
est exacte, car c'est la suite pr\'ec\'edente avec $U = X \times_Y V$.

\begin{definition}
\label{definition-higher-direct-images}
Soit $f: X \to Y$ un morphisme de sch\'emas.
Les foncteurs d\'eriv\'es droits $\{R^pf_*\}_{p \geq 1}$ de
$f_* : \textit{Ab}(X_\etale) \to \textit{Ab}(Y_\etale)$
sont appel\'es {\it images directes sup\'erieures}.
\end{definition}

\noindent
Les images directes sup\'erieures et leurs variantes dans la cat\'egorie d\'eriv\'ee
sont examin\'ees plus en d\'etail dans (ins\'erer ici une r\'ef\'erence future).



\section{Image inverse}
\label{section-inverse-image}

\noindent
Dans cette section, nous examinons bri\`evement l'image inverse des faisceaux sur les petits
sites \'etales. Sa construction pr\'ecise se trouve dans
Topologies, Section \ref{topologies-section-etale}.

\begin{definition}
\label{definition-inverse-image}
Soit $f: X\to Y$ un morphisme de sch\'emas. Les foncteurs d'{\it image inverse}, ou
de {\it tir\'e en arri\`ere}\footnote{Nous utilisons la notation $f^{-1}$ pour les images inverses de
faisceaux d'ensembles ou de groupes ab\'eliens, et nous r\'eservons $f^*$ aux
images inverses de faisceaux de modules par un morphisme de sites ou de topos annel\'es.}
sont les foncteurs
$$
f^{-1} = f_{small}^{-1} :
\Sh(Y_\etale)
\longrightarrow
\Sh(X_\etale)
$$
et
$$
f^{-1} = f_{small}^{-1} :
\textit{Ab}(Y_\etale)
\longrightarrow
\textit{Ab}(X_\etale)
$$
qui sont adjoints \`a gauche de $f_* = f_{small, *}$. Ainsi,
$f^{-1}$ est caract\'eris\'e par le fait que
$$
\Hom_{{\Sh(X_\etale)}} (f^{-1}\mathcal{G}, \mathcal{F})
=
\Hom_{\Sh(Y_\etale)} (\mathcal{G}, f_*\mathcal{F})
$$
fonctoriellement, pour tout $\mathcal{F} \in \Sh(X_\etale)$ et tout
$\mathcal{G} \in \Sh(Y_\etale)$. De m\^eme, nous avons
$$
\Hom_{{\textit{Ab}(X_\etale)}} (f^{-1}\mathcal{G}, \mathcal{F})
=
\Hom_{\textit{Ab}(Y_\etale)} (\mathcal{G}, f_*\mathcal{F})
$$
pour $\mathcal{F} \in \textit{Ab}(X_\etale)$ et
$\mathcal{G} \in \textit{Ab}(Y_\etale)$.
\end{definition}

\noindent
L'existence d'un tel adjoint n'est pas triviale.
En revanche, il existe dans un cadre assez g\'en\'eral, voir la
remarque \ref{remark-functoriality-general}
ci-dessous. Le formalisme g\'en\'eral montre que $f^{-1}\mathcal{G}$
est le faisceau associ\'e au pr\'efaisceau
\begin{equation}
\label{equation-pullback}
U/X
\longmapsto
\colim_{U \to X \times_Y V} \mathcal{G}(V/Y)
\end{equation}
o\`u la limite inductive est prise sur la cat\'egorie des couples
$(V/Y, \varphi : U/X \to X \times_Y V/X)$.
Pour le voir, appliquons
Sites, Proposition \ref{sites-proposition-get-morphism}
au foncteur $u$ de l'\'equation (\ref{equation-functorial})
et utilisons la description de $u_s = (u_p\ )^\#$ dans
Sites, Sections \ref{sites-section-continuous-functors} et
\ref{sites-section-functoriality-PSh}.
Nous utiliserons occasionnellement cette formule pour l'image inverse
afin de d\'emontrer certaines de ses propri\'et\'es fondamentales.

\begin{lemma}
\label{lemma-stalk-pullback}
Soit $f : X \to Y$ un morphisme de sch\'emas.
\begin{enumerate}
\item Le foncteur
$f^{-1} : \textit{Ab}(Y_\etale) \to \textit{Ab}(X_\etale)$
est exact.
\item Le foncteur
$f^{-1} : \Sh(Y_\etale) \to \Sh(X_\etale)$
est exact, c'est-\`a-dire qu'il commute aux limites et colimites finies, voir
Categories, Definition \ref{categories-definition-exact}.
\item Soit $\overline{x} \to X$ un point g\'eom\'etrique.
Soit $\mathcal{G}$ un faisceau sur $Y_\etale$.
Il existe alors une identification canonique
$$
(f^{-1}\mathcal{G})_{\overline{x}} = \mathcal{G}_{\overline{y}}.
$$
o\`u $\overline{y} = f \circ \overline{x}$.
\item Pour tout $V \to Y$ \'etale, nous avons $f^{-1}h_V = h_{X \times_Y V}$.
\end{enumerate}
\end{lemma}

\begin{proof}
L'exactitude de $f^{-1}$ sur les faisceaux d'ensembles r\'esulte de
Sites, Proposition \ref{sites-proposition-get-morphism}
appliqu\'ee \`a notre foncteur $u$ de l'\'equation (\ref{equation-functorial}).
En fait, l'exactitude de l'image inverse fait partie de la d\'efinition d'un
morphisme de topos (ou de sites, si l'on pr\'ef\`ere). Ainsi, (2) est v\'erifi\'e.
Cela implique (1), car, pour tout faisceau ab\'elien $\mathcal{G}$ sur
$Y_\etale$,
le faisceau d'ensembles sous-jacent \`a $f^{-1}\mathcal{F}$ est le m\^eme
que l'image par $f^{-1}$ du faisceau d'ensembles sous-jacent \`a $\mathcal{F}$, voir
Sites, Section \ref{sites-section-sheaves-algebraic-structures}.
Voir aussi
Modules on Sites, Lemma \ref{sites-modules-lemma-flat-pullback-exact}.
Dans la litt\'erature, (1) et (2) sont parfois d\'eduits de (3) au moyen du
th\'eor\`eme \ref{theorem-exactness-stalks}.

\medskip\noindent
L'assertion (3) est un fait g\'en\'eral sur les fibres des images inverses, voir
Sites, Lemma \ref{sites-lemma-point-morphism-sites}.
Nous allons aussi d\'emontrer directement (3) comme suit. Notons, d'apr\`es le
lemme \ref{lemma-stalk-exact},
que la prise des fibres commute \`a la faisceautisation.
Rappelons maintenant que $f^{-1}\mathcal{G}$ est le faisceau
associ\'e au pr\'efaisceau
$$
U \longrightarrow \colim_{U \to X \times_Y V} \mathcal{G}(V),
$$
voir l'\'equation (\ref{equation-pullback}).
Ainsi, nous avons
\begin{align*}
(f^{-1}\mathcal{G})_{\overline{x}}
& = \colim_{(U, \overline{u})} f^{-1}\mathcal{G}(U) \\
& = \colim_{(U, \overline{u})}
\colim_{a : U \to X \times_Y V} \mathcal{G}(V) \\
& = \colim_{(V, \overline{v})} \mathcal{G}(V) \\
& = \mathcal{G}_{\overline{y}}
\end{align*}
Dans la troisi\`eme \`egalit\'e, le couple $(U, \overline{u})$ et le morphisme
$a : U \to X \times_Y V$ correspondent au couple $(V, a \circ \overline{u})$.

\medskip\noindent
L'assertion (4) se d\'emontre de mani\`ere analogue en identifiant les limites inductives
qui d\'efinissent $f^{-1}h_V$. On peut aussi utiliser
le lemme de Yoneda (Categories, Lemma \ref{categories-lemma-yoneda})
et les \`egalit\'es fonctorielles
$$
\Mor_{\Sh(X_\etale)}(f^{-1}h_V, \mathcal{F}) =
\Mor_{\Sh(Y_\etale)}(h_V, f_*\mathcal{F}) =
f_*\mathcal{F}(V) = \mathcal{F}(X \times_Y V)
$$
ainsi que le fait que les pr\'efaisceaux repr\'esentables sont des faisceaux. Voir aussi
Sites, Lemma \ref{sites-lemma-pullback-representable-sheaf}
pour un r\'esultat enti\`erement g\'en\'eral.
\end{proof}

\noindent
Le couple de foncteurs $(f_*, f^{-1})$ d\'efinit un morphisme de petits topos \'etales
$$
f_{small} :
\Sh(X_\etale)
\longrightarrow
\Sh(Y_\etale)
$$
De nombreux r\'esultats g\'en\'eraux sur la cohomologie des faisceaux valent pour les topos et les
morphismes de topos. Nous nous efforcerons de signaler ceux qui sont
g\'en\'eraux et ceux qui sont propres au topos \'etale.

\begin{remark}
\label{remark-functoriality-general}
Plus g\'en\'eralement, soient $\mathcal{C}_1, \mathcal{C}_2$ des sites, et
supposons qu'ils poss\`edent des objets finals et des produits fibr\'es. Soit
$u: \mathcal{C}_2 \to \mathcal{C}_1$ un foncteur satisfaisant aux conditions suivantes :
\begin{enumerate}
\item si $\{V_i \to V\}$ est un recouvrement dans $\mathcal{C}_2$, alors
$\{u(V_i) \to u(V)\}$ est un recouvrement dans $\mathcal{C}_1$ (nous
disons que $u$ est {\it continu}), et
\item $u$ commute aux limites finies (c'est-\`a-dire que $u$ est exact \`a gauche, ou encore que
$u$ pr\'eserve les produits fibr\'es et les objets finals).
\end{enumerate}
On peut alors d\'efinir
$f_*: \Sh(\mathcal{C}_1) \to \Sh(\mathcal{C}_2)$
par $ f_* \mathcal{F}(V) = \mathcal{F}(u(V))$.
De plus, il existe un foncteur exact $f^{-1}$ qui
est adjoint \`a gauche de $f_*$, voir
Sites, Definition \ref{sites-definition-morphism-sites} et
Proposition \ref{sites-proposition-get-morphism}.
Attention : demander simplement que $u$ soit continu
et commute aux produits fibr\'es ne suffit pas pour obtenir un morphisme de topos.
\end{remark}




\section{Fonctorialit\'e des grands topos}
\label{section-functoriality-big-topoi}

\noindent
Pour tout morphisme de sch\'emas $f : X \to Y$, on dispose de toute une famille de
morphismes de topos associ\'es \`a $f$, voir
Topologies, Section \ref{topologies-section-change-topologies}
pour une liste. Les plus utilis\'es sont peut-\^etre les morphismes de topos
$$
f_{big} = f_{big, \tau} :
\Sh((\Sch/X)_\tau)
\longrightarrow
\Sh((\Sch/Y)_\tau)
$$
o\`u $\tau \in \{Zariski, \etale, smooth, syntomic, fppf\}$.
Chacun d'eux correspond \`a un foncteur continu
$$
(\Sch/Y)_\tau \longrightarrow (\Sch/X)_\tau, \quad
V/Y \longmapsto X \times_Y V/X
$$
qui pr\'eserve les objets finals, les produits fibr\'es et les recouvrements, et d\'efinit donc
un morphisme de sites
$$
f_{big} : (\Sch/X)_\tau \longrightarrow (\Sch/Y)_\tau.
$$
Voir
Topologies, Sections \ref{topologies-section-zariski},
\ref{topologies-section-etale},
\ref{topologies-section-smooth},
\ref{topologies-section-syntomic}, et
\ref{topologies-section-fppf}.
En particulier, l'image directe par $f_{big}$ est donn\'ee par la r\`egle
$$
(f_{big, *}\mathcal{F})(V/Y) = \mathcal{F}(X \times_Y V/X)
$$
Il se trouve que ces morphismes de topos poss\`edent un foncteur image
inverse $f_{big}^{-1}$ tr\`es facile \`a d\'ecrire.
Plus pr\'ecis\'ement, nous avons
$$
(f_{big}^{-1}\mathcal{G})(U/X) = \mathcal{G}(U/Y)
$$
o\`u le morphisme structural de $U/Y$ est la compos\'ee du
morphisme structural $U \to X$ et de $f$, voir
Topologies, Lemmas \ref{topologies-lemma-morphism-big},
\ref{topologies-lemma-morphism-big-etale},
\ref{topologies-lemma-morphism-big-smooth},
\ref{topologies-lemma-morphism-big-syntomic}, et
\ref{topologies-lemma-morphism-big-fppf}.







\section{Fonctorialit\'e et faisceaux de modules}
\label{section-morphisms-modules}

\noindent
Dans cette section, nous allons reformuler une partie des r\'esultats expliqu\'es dans
Descent, Sections \ref{descent-section-quasi-coherent-sheaves},
\ref{descent-section-quasi-coherent-cohomology}, et
\ref{descent-section-quasi-coherent-sheaves-bis}
dans le cadre des topologies \'etales. Soit $f : X \to Y$ un morphisme de
sch\'emas. Nous avons vu plus haut, voir les
sections \ref{section-functoriality}, \ref{section-direct-image} et
\ref{section-inverse-image},
qu'il induit un morphisme $f_{small}$ de petits sites \'etales. Dans
Descent, Remark \ref{descent-remark-change-topologies-ringed},
nous avons vu que $f$ induit aussi un morphisme naturel
$$
f_{small}^\sharp :
\mathcal{O}_{Y_\etale}
\longrightarrow
f_{small, *}\mathcal{O}_{X_\etale}
$$
de faisceaux d'anneaux sur
$Y_\etale$ tel que $(f_{small}, f_{small}^\sharp)$
soit un morphisme de sites annel\'es. Voir
Modules on Sites, Definition \ref{sites-modules-definition-ringed-site}
pour la d\'efinition d'un morphisme de sites annel\'es.
Rappelons simplement ici que $f_{small}^\sharp$ est d\'efini par le
syst\`eme compatible de morphismes
$$
\text{pr}_V^\sharp : \mathcal{O}(V) \longrightarrow \mathcal{O}(X \times_Y V)
$$
lorsque $V$ parcourt les objets de $Y_\etale$.

\medskip\noindent
Il est clair que cette construction est compatible avec les compositions de
morphismes de sch\'emas. Plus pr\'ecis\'ement, si $f : X \to Y$ et $g : Y \to Z$
sont des morphismes de sch\'emas, alors nous avons
$$
(g_{small}, g_{small}^\sharp)
\circ
(f_{small}, f_{small}^\sharp)
=
((g \circ f)_{small}, (g \circ f)_{small}^\sharp)
$$
comme morphismes de topos annel\'es. De plus, d'apr\`es
Modules on Sites, Definition \ref{sites-modules-definition-pushforward},
nous voyons que tout morphisme de sch\'emas $f : X \to Y$
donne des foncteurs bien d\'efinis d'image inverse et d'image directe
\begin{align*}
f_{small}^* :
\textit{Mod}(\mathcal{O}_{Y_\etale})
\longrightarrow
\textit{Mod}(\mathcal{O}_{X_\etale}), \\
f_{small, *} :
\textit{Mod}(\mathcal{O}_{X_\etale})
\longrightarrow
\textit{Mod}(\mathcal{O}_{Y_\etale})
\end{align*}
qui sont adjoints de la mani\`ere habituelle. Si $g : Y \to Z$ est un autre morphisme
de sch\'emas, alors nous avons
$(g \circ f)_{small}^* = f_{small}^* \circ g_{small}^*$
et $(g \circ f)_{small, *} = g_{small, *} \circ f_{small, *}$
en vertu de ce que nous avons dit sur les compositions.

\medskip\noindent
Il y a une diff\'erence assez importante entre la cat\'egorie
de tous les $\mathcal{O}_X$-modules sur $X$ et la cat\'egorie de tous les
$\mathcal{O}_{X_\etale}$-modules sur $X_\etale$. Mais les
r\'esultats de
Descent, Sections \ref{descent-section-quasi-coherent-sheaves},
\ref{descent-section-quasi-coherent-cohomology}, et
\ref{descent-section-quasi-coherent-sheaves-bis}
nous disent qu'il y a peu de diff\'erence entre consid\'erer les modules quasi-coh\'erents
sur $S$ et les modules quasi-coh\'erents sur $S_\etale$.
(Nous l'avons d\'ej\`a vu dans le
th\'eor\`eme \ref{theorem-quasi-coherent},
par exemple.)
En particulier, si $f : X \to Y$ est un morphisme quelconque de sch\'emas, alors
les foncteurs d'image inverse $f_{small}^*$ et $f^*$ s'identifient pour les
faisceaux quasi-coh\'erents, voir
Descent,
Proposition \ref{descent-proposition-equivalence-quasi-coherent-functorial}.
De plus, il en va de m\^eme pour l'image directe pourvu que $f$ soit
quasi-compact et quasi-s\'epar\'e, voir
Descent, Lemma \ref{descent-lemma-higher-direct-images-small-etale}.

\medskip\noindent
Disons quelques mots de la fonctorialit\'e du faisceau structural sur les grands sites.
Soit $f : X \to Y$ un morphisme de sch\'emas. Choisissons l'une des
topologies $\tau \in \{Zariski,\linebreak[0]
\etale,\linebreak[0] smooth,\linebreak[0] syntomic,\linebreak[0]
fppf\}$. Alors le morphisme
$f_{big} : (\Sch/X)_\tau \to (\Sch/Y)_\tau$
devient un morphisme de sites annel\'es au moyen d'un morphisme
$$
f_{big}^\sharp : \mathcal{O}_Y \longrightarrow f_{big, *}\mathcal{O}_X
$$
voir Descent, Remark \ref{descent-remark-change-topologies-ringed}.
En fait, il est donn\'e par la m\^eme construction que celle expos\'ee pour les petits
sites plus haut.












\section{Comparaison de topologies}
\label{section-compare-topologies}

\noindent
Dans cette section, nous commen\c{c}ons \`a \`etudier ce qui se produit lorsque l'on compare
les faisceaux relativement \`a diff\'erentes topologies.

\begin{lemma}
\label{lemma-where-sections-are-equal}
Soit $S$ un sch\'ema. Soit $\mathcal{F}$ un faisceau d'ensembles sur $S_\etale$.
Soient $s, t \in \mathcal{F}(S)$. Il existe alors un ouvert $W \subset S$
caract\'eris\'e par la propri\'et\'e suivante : un morphisme $f : T \to S$
se factorise par $W$ si et seulement si $s|_T = t|_T$ (la restriction est
l'image inverse par $f_{small}$).
\end{lemma}

\begin{proof}
Consid\'erons le pr\'efaisceau qui associe \`a $U \in \Ob(S_\etale)$ l'ensemble vide
si $s|_U \not = t|_U$, et un singleton sinon. Il est clair qu'il s'agit
d'un sous-faisceau de l'objet final de $\Sh(S_\etale)$. D'apr\`es le
lemme \ref{lemma-support-subsheaf-final},
nous trouvons un ouvert $W \subset S$ qui repr\'esente ce pr\'efaisceau.
Pour un point g\'eom\'etrique $\overline{x}$ de $S$, nous voyons que $\overline{x} \in W$
si et seulement si les germes de $s$ et de $t$ en $\overline{x}$ sont identiques.
D'apr\`es la description des fibres des images inverses donn\'ee dans le
lemme \ref{lemma-stalk-pullback},
nous voyons que $W$ poss\`ede la propri\'et\'e voulue.
\end{proof}

\begin{lemma}
\label{lemma-describe-pullback}
Soit $S$ un sch\'ema. Soit $\tau \in \{Zariski, \etale\}$. Consid\'erons le morphisme
$$
\pi_S : (\Sch/S)_\tau \longrightarrow S_\tau
$$
de Topologies, Lemma \ref{topologies-lemma-at-the-bottom} ou
\ref{topologies-lemma-at-the-bottom-etale}. Soit $\mathcal{F}$ un faisceau sur
$S_\tau$. Alors $\pi_S^{-1}\mathcal{F}$ est donn\'e par la r\`egle
$$
(\pi_S^{-1}\mathcal{F})(T) = \Gamma(T_\tau, f_{small}^{-1}\mathcal{F})
$$
o\`u $f : T \to S$. De plus, $\pi_S^{-1}\mathcal{F}$ satisfait \`a la
condition de faisceau relativement aux recouvrements fpqc.
\end{lemma}

\begin{proof}
Observons que nous avons un morphisme $i_f : \Sh(T_\tau) \to \Sh(\Sch/S)_\tau)$
tel que $\pi_S \circ i_f = f_{small}$ comme morphismes
$T_\tau \to S_\tau$, voir
Topologies, Lemmas \ref{topologies-lemma-put-in-T},
\ref{topologies-lemma-morphism-big-small},
\ref{topologies-lemma-put-in-T-etale}, et
\ref{topologies-lemma-morphism-big-small-etale}.
Comme l'image inverse est transitive, nous voyons que
$i_f^{-1} \pi_S^{-1}\mathcal{F} = f_{small}^{-1}\mathcal{F}$ comme voulu.

\medskip\noindent
Soit $\{g_i : T_i \to T\}_{i \in I}$ un recouvrement fpqc.
La derni\`ere assertion signifie ce qui suit : \`etant donn\'e un faisceau $\mathcal{G}$
sur $T_\tau$ et des sections
$s_i \in \Gamma(T_i, g_{i, small}^{-1}\mathcal{G})$ dont les images inverses
sur $T_i \times_T T_j$ concordent, il existe une unique section $s$ de $\mathcal{G}$
sur $T$ dont l'image inverse sur $T_i$ concorde avec $s_i$.

\medskip\noindent
Soit $V \to T$ un objet de $T_\tau$ et soit $t \in \mathcal{G}(V)$.
Pour tout $i$, il existe un plus grand ouvert $W_i \subset T_i \times_T V$
tel que les images inverses de $s_i$ et de $t$ concordent comme sections de l'image
inverse de $\mathcal{G}$ sur $W_i \subset T_i \times_T V$, voir le
lemme \ref{lemma-where-sections-are-equal}.
Comme $s_i$ et $s_j$ concordent sur $T_i \times_T T_j$, nous constatons
que les images inverses de $W_i$ et de $W_j$ sont le m\^eme ouvert au-dessus de
$T_i \times_T T_j \times_T V$. D'apr\`es
Descent, Lemma \ref{descent-lemma-open-fpqc-covering},
nous trouvons un ouvert $W \subset V$ dont l'image inverse sur $T_i \times_T V$
redonne $W_i$.

\medskip\noindent
Par construction de $g_{i, small}^{-1}\mathcal{G}$, il existe
un recouvrement pour $\tau$ $\{T_{ij} \to T_i\}_{j \in J_i}$, pour chaque $j$ une
immersion ouverte ou un morphisme \'etale $V_{ij} \to T$, une section
$t_{ij} \in \mathcal{G}(V_{ij})$, et des diagrammes commutatifs
$$
\xymatrix{
T_{ij} \ar[r] \ar[d] & V_{ij} \ar[d] \\
T_i \ar[r] & T
}
$$
tels que $s_i|_{T_{ij}}$ soit l'image inverse de $t_{ij}$. Autrement dit,
apr\`es avoir remplac\'e le recouvrement $\{T_i \to T\}$ par $\{T_{ij} \to T\}$,
nous pouvons supposer qu'il existe des factorisations $T_i \to V_i \to T$ avec
$V_i \in \Ob(T_\tau)$ et des sections $t_i \in \mathcal{G}(V_i)$
dont les images inverses sont $s_i$ sur $T_i$.
D'apr\`es le r\'esultat du paragraphe pr\'ec\'edent, nous trouvons des ouverts $W_i \subset V_i$
tels que $t_i|_{W_i}$ ``concorde avec'' chaque $s_j$ sur $T_j \times_T W_i$.
Notons que $T_i \to V_i$ se factorise par $W_i$.
Ainsi, $\{W_i \to T\}$ est un recouvrement pour $\tau$ et le lemme est d\'emontr\'e.
\end{proof}

\begin{lemma}
\label{lemma-sections-upstairs}
Soit $S$ un sch\'ema. Soit $f : T \to S$ un morphisme tel que
\begin{enumerate}
\item $f$ est plat et quasi-compact, et
\item les fibres g\'eom\'etriques de $f$ sont connexes.
\end{enumerate}
Soit $\mathcal{F}$ un faisceau sur $S_\etale$.
Alors $\Gamma(S, \mathcal{F}) = \Gamma(T, f^{-1}_{small}\mathcal{F})$.
\end{lemma}

\begin{proof}
Il existe un morphisme canonique
$\Gamma(S, \mathcal{F}) \to \Gamma(T, f_{small}^{-1}\mathcal{F})$.
Comme $f$ est surjectif (puisque ses fibres sont connexes), nous voyons que
ce morphisme est injectif.

\medskip\noindent
Pour montrer que le morphisme est surjectif, soit
$\alpha \in \Gamma(T, f_{small}^{-1}\mathcal{F})$.
Comme $\{T \to S\}$ est un recouvrement fpqc, nous pouvons utiliser le
lemme \ref{lemma-describe-pullback} pour voir qu'il suffit de montrer que
les deux images inverses de $\alpha$ sur $T \times_S T$ sont la m\^eme section
par les deux projections. Soit $\overline{s} \to S$ un point g\'eom\'etrique.
Il suffit de montrer l'\'egalit\'e sur $(T \times_S T)_{\overline{s}}$,
car tout point g\'eom\'etrique de $T \times_S T$ est contenu dans l'une de
ces fibres g\'eom\'etriques. Autrement dit, nous voulons montrer que
$\alpha|_{T_{\overline{s}}}$ donne la m\^eme section par les deux projections sur
$$
(T \times_S T)_{\overline{s}} =
T_{\overline{s}} \times_{\overline{s}} T_{\overline{s}}
$$
vers $T_{\overline{s}}$.
Cependant, comme $\mathcal{F}|_{T_{\overline{s}}}$ est l'image
inverse de $\mathcal{F}|_{\overline{s}}$, c'est un faisceau constant
de valeur $\mathcal{F}_{\overline{s}}$. Comme $T_{\overline{s}}$
est connexe par hypoth\`ese, toute section d'un faisceau constant est constante.
Ainsi, $\alpha|_{T_{\overline{s}}}$ correspond \`a un \`el\'ement
de $\mathcal{F}_{\overline{s}}$. Les deux images inverses sur
$(T \times_S T)_{\overline{s}}$ correspondent donc toutes deux
\`a ce m\^eme \`el\'ement, d'o\`u la conclusion.
\end{proof}

\noindent
Voici une variante du lemme \ref{lemma-sections-upstairs}
o\`u nous ne supposons pas que le morphisme est plat.

\begin{lemma}
\label{lemma-sections-upstairs-submersive}
Soit $S$ un sch\'ema. Soit $f : X \to S$ un morphisme tel que
\begin{enumerate}
\item $f$ est submersif, et
\item les fibres g\'eom\'etriques de $f$ sont connexes.
\end{enumerate}
Soit $\mathcal{F}$ un faisceau sur $S_\etale$.
Alors $\Gamma(S, \mathcal{F}) = \Gamma(X, f^{-1}_{small}\mathcal{F})$.
\end{lemma}

\begin{proof}
Il existe un morphisme canonique
$\Gamma(S, \mathcal{F}) \to \Gamma(X, f_{small}^{-1}\mathcal{F})$.
Comme $f$ est surjectif (puisque ses fibres sont connexes), nous voyons que
ce morphisme est injectif.

\medskip\noindent
Pour montrer que le morphisme est surjectif, soit
$\tau \in \Gamma(X, f_{small}^{-1}\mathcal{F})$.
Il suffit de trouver
un recouvrement \'etale $\{U_i \to S\}$ et des
sections $\sigma_i \in \mathcal{F}(U_i)$
telles que l'image inverse de $\sigma_i$ soit $\tau|_{X \times_S U_i}$.
En effet, l'injectivit\'e montr\'ee ci-dessus garantit
que $\sigma_i$ et $\sigma_j$ induisent la m\^eme
section de $\mathcal{F}$ sur $U_i \times_S U_j$.
Nous obtenons donc une unique section $\sigma \in \mathcal{F}(S)$
qui se restreint \`a $\sigma_i$ sur $U_i$.
L'image inverse de $\sigma$ sur $X$ est alors $\tau$,
car c'est vrai localement.

\medskip\noindent
Soit $\overline{x}$ un point g\'eom\'etrique de $X$ d'image $\overline{s}$
dans $S$. Consid\'erons l'image de $\tau$ dans la fibre
$$
(f_{small}^{-1}\mathcal{F})_{\overline{x}} = \mathcal{F}_{\overline{s}}
$$
Voir le lemme \ref{lemma-stalk-pullback}.
Nous pouvons trouver un voisinage \'etale $U \to S$ de $\overline{s}$
et une section $\sigma \in \mathcal{F}(U)$ qui s'envoie sur cette image
dans la fibre. Ainsi, apr\`es avoir remplac\'e $S$ par $U$ et $X$ par $X \times_S U$,
nous pouvons supposer qu'il existe une section $\sigma$ de $\mathcal{F}$ sur $S$
dont l'image dans $(f_{small}^{-1}\mathcal{F})_{\overline{x}}$ est la
m\^eme que celle de $\tau$.

\medskip\noindent
D'apr\`es le lemme \ref{lemma-where-sections-are-equal},
il existe un ouvert maximal $W \subset X$ sur lequel
$f_{small}^{-1}\sigma$ et $\tau$ concordent sur $W$,
et la formation de $W$ commute aux changements de base ult\'erieurs.
Observons que l'image inverse de $\mathcal{F}$ sur la
fibre g\'eom\'etrique $X_{\overline{s}}$ est l'image inverse
de $\mathcal{F}_{\overline{s}}$, consid\'er\'e comme faisceau sur
$\overline{s}$, par $X_{\overline{s}} \to \overline{s}$.
Nous voyons donc que $\tau$ et $\sigma$ donnent des sections
du faisceau constant de valeur $\mathcal{F}_{\overline{s}}$
sur $X_{\overline{s}}$ qui concordent en un point. Comme
$X_{\overline{s}}$ est connexe par hypoth\`ese, nous concluons
que $W$ contient $X_s$. Le m\^eme argument appliqu\'e aux autres
fibres g\'eom\'etriques montre que $W$ contient toute fibre qu'il rencontre.
Comme $f$ est submersif, nous concluons que $W$ est l'image
inverse d'un voisinage ouvert de $s$ dans $S$.
Cela ach\`eve la d\'emonstration.
\end{proof}

\begin{lemma}
\label{lemma-sections-base-field-extension}
Soit $K/k$ une extension de corps o\`u $k$ est
s\'eparablement clos. Soit $S$ un sch\'ema sur $k$. Notons
$p : S_K = S \times_{\Spec(k)} \Spec(K) \to S$ la projection.
Soit $\mathcal{F}$ un faisceau sur $S_\etale$.
Alors $\Gamma(S, \mathcal{F}) = \Gamma(S_K, p^{-1}_{small}\mathcal{F})$.
\end{lemma}

\begin{proof}
Cela r\'esulte du lemme \ref{lemma-sections-upstairs}. En effet, il est clair
que $p$ est plat et quasi-compact comme changement de base de
$\Spec(K) \to \Spec(k)$. D'autre part, si $\overline{s} : \Spec(L) \to S$
est un point g\'eom\'etrique, alors la fibre de $p$ en $\overline{s}$
est le spectre de $K \otimes_k L$, qui est irr\'eductible, donc connexe, d'apr\`es
Algebra, Lemma \ref{algebra-lemma-separably-closed-irreducible}.
\end{proof}









\section{Reconstruction des morphismes}
\label{section-morphisms}

\noindent
Dans cette section, nous d\'emontrons que la r\`egle qui associe \`a un sch\'ema
son petit topos \'etale localement annel\'e est pleinement fid\`ele en un sens
convenable, voir le
th\'eor\`eme \ref{theorem-fully-faithful}.

\begin{lemma}
\label{lemma-morphism-locally-ringed}
Soit $f : X \to Y$ un morphisme de sch\'emas.
Le morphisme de sites annel\'es $(f_{small}, f_{small}^\sharp)$
associ\'e \`a $f$ est un morphisme de sites localement annel\'es, voir
Modules on Sites,
Definition \ref{sites-modules-definition-morphism-locally-ringed-topoi}.
\end{lemma}

\begin{proof}
Notons que l'assertion a un sens, puisque nous avons vu que
$(X_\etale, \mathcal{O}_{X_\etale})$ et
$(Y_\etale, \mathcal{O}_{Y_\etale})$
sont des sites localement annel\'es, voir le
lemme \ref{lemma-etale-site-locally-ringed}.
De plus, nous savons que $X_\etale$ a assez de points, voir le
th\'eor\`eme \ref{theorem-exactness-stalks}
et les
remarques \ref{remarks-enough-points}.
Il suffit donc de d\'emontrer que $(f_{small}, f_{small}^\sharp)$
satisfait \`a la condition (3) de
Modules on Sites,
Lemma \ref{sites-modules-lemma-locally-ringed-morphism}.
Pour le voir, prenons un point $p$ de $X_\etale$. D'apr\`es le
lemme \ref{lemma-points-small-etale-site},
$p$ correspond \`a un point g\'eom\'etrique $\overline{x}$ de $X$.
D'apr\`es le
lemme \ref{lemma-stalk-pullback},
le point $q = f_{small} \circ p$ correspond au
point g\'eom\'etrique $\overline{y} = f \circ \overline{x}$ de $Y$.
L'assertion \`a d\'emontrer est donc que le morphisme induit
sur les fibres
$$
(\mathcal{O}_Y)_{\overline{y}} \longrightarrow (\mathcal{O}_X)_{\overline{x}}
$$
est un morphisme local d'anneaux. Supposons que $a \in (\mathcal{O}_Y)_{\overline{y}}$
soit un \`el\'ement du membre de gauche qui s'envoie sur un \`el\'ement de l'id\'eal
maximal du membre de droite. Supposons que $a$ soit la classe d'\'equivalence
d'un triplet $(V, \overline{v}, a)$ o\`u $V \to Y$ est \'etale,
$\overline{v} : \overline{x} \to V$ au-dessus de $Y$, et $a \in \mathcal{O}(V)$.
Il s'envoie sur la classe d'\'equivalence de
$(X \times_Y V, \overline{x} \times \overline{v}, \text{pr}_V^\sharp(a))$
dans l'anneau local $(\mathcal{O}_X)_{\overline{x}}$. Mais il est clair que
le fait d'appartenir \`a l'id\'eal maximal signifie que l'image de $\text{pr}_V^\sharp(a)$
dans $\kappa(\overline{x})$ est nulle. Par cons\'equent, l'image de
$a$ dans $\kappa(\overline{x})$ est elle aussi nulle. Cela signifie que $a$ appartient \`a
l'id\'eal maximal de $(\mathcal{O}_Y)_{\overline{y}}$.
\end{proof}

\begin{lemma}
\label{lemma-2-morphism}
Soient $X$, $Y$ des sch\'emas. Soit $f : X \to Y$ un morphisme de sch\'emas.
Soit $t$ un $2$-morphisme de $(f_{small}, f_{small}^\sharp)$ vers lui-m\^eme, voir
Modules on Sites,
Definition \ref{sites-modules-definition-2-morphism-ringed-topoi}.
Alors $t = \text{id}$.
\end{lemma}

\begin{proof}
Cela signifie que $t : f^{-1}_{small} \to f^{-1}_{small}$
est une transformation naturelle telle que le diagramme
$$
\xymatrix{
f_{small}^{-1}\mathcal{O}_Y
\ar[rd]_{f_{small}^\sharp}  & &
f_{small}^{-1}\mathcal{O}_Y \ar[ll]^t \ar[ld]^{f_{small}^\sharp} \\
& \mathcal{O}_X
}
$$
soit commutatif. Supposons que $V \to Y$ soit \'etale avec $V$ affine. D'apr\`es
Morphisms, Lemma \ref{morphisms-lemma-quasi-affine-finite-type-over-S},
nous pouvons choisir une immersion $i : V \to \mathbf{A}^n_Y$ au-dessus de $Y$.
En termes de faisceaux, cela signifie que $i$ induit une injection
$h_i : h_V \to \prod_{j = 1, \ldots, n} \mathcal{O}_Y$ de faisceaux.
Le changement de base $i'$ de $i$ \`a $X$ est une immersion
(Schemes, Lemma \ref{schemes-lemma-base-change-immersion}).
Ainsi, $i' : X \times_Y V \to \mathbf{A}^n_X$ est une immersion, ce qui
signifie \`a son tour que
$h_{i'} : h_{X \times_Y V} \to \prod_{j = 1, \ldots, n} \mathcal{O}_X$
est une injection de faisceaux.
Via l'identification $f_{small}^{-1}h_V = h_{X \times_Y V}$ du
lemme \ref{lemma-stalk-pullback},
le morphisme $h_{i'}$ est \`egal \`a
$$
\xymatrix{
f_{small}^{-1}h_V \ar[r]^-{f^{-1}h_i} &
\prod_{j = 1, \ldots, n} f_{small}^{-1}\mathcal{O}_Y
\ar[r]^{\prod f^\sharp} &
\prod_{j = 1, \ldots, n} \mathcal{O}_X
}
$$
(nous omettons la v\'erification). Cela signifie que le morphisme
$t : f_{small}^{-1}h_V \to f_{small}^{-1}h_V$
s'ins\`ere dans le diagramme commutatif
$$
\xymatrix{
f_{small}^{-1}h_V \ar[r]^-{f^{-1}h_i} \ar[d]^t &
\prod_{j = 1, \ldots, n} f_{small}^{-1}\mathcal{O}_Y
\ar[r]^-{\prod f^\sharp} \ar[d]^{\prod t} &
\prod_{j = 1, \ldots, n} \mathcal{O}_X \ar[d]^{\text{id}}\\
f_{small}^{-1}h_V \ar[r]^-{f^{-1}h_i} &
\prod_{j = 1, \ldots, n} f_{small}^{-1}\mathcal{O}_Y
\ar[r]^-{\prod f^\sharp} &
\prod_{j = 1, \ldots, n} \mathcal{O}_X
}
$$
Le carr\'e de droite est commutatif par notre hypoth\`ese sur $t$
expliqu\'ee ci-dessus.
Comme la compos\'ee des fl\`eches horizontales est injective
d'apr\`es la discussion pr\'ec\'edente, nous concluons que la fl\`eche verticale de gauche
est elle aussi l'identit\'e. Tout faisceau d'ensembles sur
$Y_\etale$ est quotient d'un coproduit (\'enorme) de faisceaux
de la forme $h_V$ avec $V$ affine (combiner
Topologies, Lemma \ref{topologies-lemma-alternative}
avec
Sites, Lemma \ref{sites-lemma-sheaf-coequalizer-representable}).
Nous concluons donc que $t : f_{small}^{-1} \to f_{small}^{-1}$
est la transformation identit\'e, comme voulu.
\end{proof}

\begin{lemma}
\label{lemma-faithful}
Soient $X$, $Y$ des sch\'emas.
Deux morphismes de sch\'emas $a, b : X \to Y$
pour lesquels il existe un $2$-isomorphisme
$(a_{small}, a_{small}^\sharp) \cong (b_{small}, b_{small}^\sharp)$
dans la $2$-cat\'egorie des topos annel\'es sont \`egaux.
\end{lemma}

\begin{proof}
Donnons l'argument avec soin, car il est quelque peu d\'eroutant.
Soit $t : a_{small}^{-1} \to b_{small}^{-1}$ le $2$-isomorphisme.
Consid\'erons un ouvert quelconque $V \subset Y$. Notons que $h_V$ est un sous-faisceau du
faisceau final $*$. Ainsi, $a_{small}^{-1}h_V = h_{a^{-1}(V)}$
et $b_{small}^{-1}h_V = h_{b^{-1}(V)}$ sont tous deux des sous-faisceaux du faisceau final.
L'isomorphisme
$$
t : a_{small}^{-1}h_V = h_{a^{-1}(V)} \to b_{small}^{-1}h_V = h_{b^{-1}(V)}
$$
est donc n\'ecessairement l'identit\'e, et $a^{-1}(V) = b^{-1}(V)$.
Il s'ensuit que $a$ et $b$ sont \`egaux sur les espaces topologiques sous-jacents.
Prenons ensuite une section $f \in \mathcal{O}_Y(V)$. Elle d\'etermine, et est
d\'etermin\'ee par, un morphisme de faisceaux d'ensembles
$f : h_V \to \mathcal{O}_Y$.
Prenons son image inverse et appliquons $t$ pour obtenir un diagramme commutatif
$$
\xymatrix{
h_{b^{-1}(V)} \ar@{=}[r] &
b_{small}^{-1}h_V \ar[d]^{b_{small}^{-1}(f)} & &
a_{small}^{-1}h_V \ar[d]^{a_{small}^{-1}(f)} \ar[ll]^t &
h_{a^{-1}(V)} \ar@{=}[l]
\\
& b_{small}^{-1}\mathcal{O}_Y
\ar[rd]_{b^\sharp}  & &
a_{small}^{-1}\mathcal{O}_Y \ar[ll]^t \ar[ld]^{a^\sharp} \\
& & \mathcal{O}_X
}
$$
o\`u le triangle est commutatif par d\'efinition d'un $2$-isomorphisme dans
Modules on Sites, Section \ref{sites-modules-section-2-category}.
Nous avons vu plus haut que la compos\'ee des fl\`eches horizontales
sup\'erieures provient de l'identit\'e $a^{-1}(V) = b^{-1}(V)$.
La commutativit\'e du diagramme nous dit donc que
$a_{small}^\sharp(f) = b_{small}^\sharp(f)$ dans
$\mathcal{O}_X(a^{-1}(V)) = \mathcal{O}_X(b^{-1}(V))$.
Comme cela vaut pour tout ouvert $V$ et tout $f \in \mathcal{O}_Y(V)$,
nous concluons que $a = b$ comme morphismes de sch\'emas.
\end{proof}

\begin{lemma}
\label{lemma-morphism-ringed-etale-topoi-affines}
Soient $X$, $Y$ des sch\'emas affines.
Soit
$$
(g, g^\#) :
(\Sh(X_\etale), \mathcal{O}_X)
\longrightarrow
(\Sh(Y_\etale), \mathcal{O}_Y)
$$
un morphisme de topos localement annel\'es. Il existe alors un
unique morphisme de sch\'emas $f : X \to Y$ tel que
$(g, g^\#)$ soit $2$-isomorphe \`a $(f_{small}, f_{small}^\sharp)$,
voir
Modules on Sites,
Definition \ref{sites-modules-definition-2-morphism-ringed-topoi}.
\end{lemma}

\begin{proof}
Dans cette d\'emonstration, nous notons $\mathcal{O}_X$ le faisceau structural
du petit site \'etale $X_\etale$, et de m\^eme pour
$\mathcal{O}_Y$. Posons $Y = \Spec(B)$ et $X = \Spec(A)$. Comme
$B = \Gamma(Y_\etale, \mathcal{O}_Y)$,
$A = \Gamma(X_\etale, \mathcal{O}_X)$,
nous voyons que $g^\sharp$ induit un morphisme d'anneaux $\varphi : B \to A$.
Soit $f = \Spec(\varphi) : X \to Y$ le morphisme correspondant
de sch\'emas affines. Nous allons montrer que ce $f$ convient.

\medskip\noindent
Soit $V \to Y$ un sch\'ema affine \'etale sur $Y$. Nous pouvons donc \`ecrire
$V = \Spec(C)$, o\`u $C$ est une $B$-alg\`ebre \'etale. Nous pouvons \`ecrire
$$
C = B[x_1, \ldots, x_n]/(P_1, \ldots, P_n)
$$
avec des polyn\^omes $P_i$ tels que $\Delta = \det(\partial P_i/ \partial x_j)$
soit inversible dans $C$, voir par exemple
Algebra, Lemma \ref{algebra-lemma-etale-standard-smooth}.
Si $T$ est un sch\'ema sur $Y$, alors un $T$-point de $V$ est donn\'e par
$n$ sections de $\Gamma(T, \mathcal{O}_T)$ qui satisfont aux \`equations polynomiales
$P_1 = 0, \ldots, P_n = 0$. Autrement dit, le faisceau $h_V$
sur $Y_\etale$ est l'\'egalisateur des deux morphismes
$$
\xymatrix{
\prod\nolimits_{i = 1, \ldots, n} \mathcal{O}_Y
\ar@<1ex>[r]^a \ar@<-1ex>[r]_b &
\prod\nolimits_{j = 1, \ldots, n} \mathcal{O}_Y
}
$$
o\`u $b(h_1, \ldots, h_n) = 0$ et
$a(h_1, \ldots, h_n) =
(P_1(h_1, \ldots, h_n), \ldots, P_n(h_1, \ldots, h_n))$.
Comme $g^{-1}$ est exact, nous concluons que la ligne sup\'erieure de la
partie du diagramme commutatif suivant form\'ee par les fl\`eches pleines est elle aussi un diagramme d'\'egalisateur :
$$
\xymatrix{
g^{-1}h_V \ar[r] \ar@{..>}[d] &
\prod\nolimits_{i = 1, \ldots, n} g^{-1}\mathcal{O}_Y
\ar@<1ex>[r]^{g^{-1}a} \ar@<-1ex>[r]_{g^{-1}b} \ar[d]^{\prod g^\sharp} &
\prod\nolimits_{j = 1, \ldots, n} g^{-1}\mathcal{O}_Y \ar[d]^{\prod g^\sharp}\\
h_{X \times_Y V} \ar[r] &
\prod\nolimits_{i = 1, \ldots, n} \mathcal{O}_X
\ar@<1ex>[r]^{a'} \ar@<-1ex>[r]_{b'} &
\prod\nolimits_{j = 1, \ldots, n} \mathcal{O}_X  \\
}
$$
Ici, $b'$ est l'application nulle et $a'$ est l'application d\'efinie par les
images $P'_i = \varphi(P_i) \in A[x_1, \ldots, x_n]$ suivant la m\^eme
r\`egle
$a'(h_1, \ldots, h_n) =
(P'_1(h_1, \ldots, h_n), \ldots, P'_n(h_1, \ldots, h_n))$.
que celle qui d\'efinit $a$. La commutativit\'e du diagramme r\'esulte de
l'\'egalit\'e $\varphi = g^\sharp$ sur les sections globales. La ligne
inf\'erieure est \'egalement un diagramme \'egalisateur, exactement par les m\^emes arguments que
pr\'ec\'edemment, puisque $X \times_Y V$ est le sch\'ema affine
$\Spec(A \otimes_B C)$ et
$A \otimes_B C = A[x_1, \ldots, x_n]/(P'_1, \ldots, P'_n)$.
Nous obtenons donc une unique fl\`eche pointill\'ee
$g^{-1}h_V \to h_{X \times_Y V}$ qui s'ins\`ere dans le diagramme.

\medskip\noindent
Nous affirmons que le morphisme de faisceaux $g^{-1}h_V \to h_{X \times_Y V}$
est un isomorphisme. Comme le petit site \'etale de $X$ poss\`ede suffisamment de points
(Th\'eor\`eme \ref{theorem-exactness-stalks}),
il suffit de le d\'emontrer sur les germes. Soit donc $\overline{x}$ un
point g\'eom\'etrique de $X$, et notons $p$ le point associ\'e du
petit topos \'etale de $X$. Posons $q = g \circ p$. C'est un point du
petit topos \'etale de $Y$. D'apr\`es le
Lemme \ref{lemma-points-small-etale-site},
nous voyons que $q$ correspond \`a un point g\'eom\'etrique $\overline{y}$ de
$Y$. Consid\'erons le morphisme de germes
$$
(g^\sharp)_p :
(\mathcal{O}_Y)_{\overline{y}} =
\mathcal{O}_{Y, q} =
(g^{-1}\mathcal{O}_Y)_p
\longrightarrow
\mathcal{O}_{X, p} =
(\mathcal{O}_X)_{\overline{x}}
$$
Comme $(g, g^\sharp)$ est un morphisme de topos {\it localement} annel\'es,
$(g^\sharp)_p$ est un homomorphisme local entre anneaux locaux strictement
hens\'eliens. En localisant le grand diagramme commutatif ci-dessus et en appliquant
Alg\`ebre, Lemme \ref{algebra-lemma-strictly-henselian-solutions},
nous concluons que $(g^{-1}h_V)_p \to (h_{X \times_Y V})_p$ est un isomorphisme,
comme voulu.

\medskip\noindent
Nous affirmons que les isomorphismes $g^{-1}h_V \to h_{X \times_Y V}$ sont
fonctoriels. Plus pr\'ecis\'ement, supposons que $V_1 \to V_2$ soit un morphisme de sch\'emas affines
\'etales sur $Y$. \'Ecrivons
$V_i = \Spec(C_i)$ avec
$$
C_i = B[x_{i, 1}, \ldots, x_{i, n_i}]/(P_{i, 1}, \ldots, P_{i, n_i})
$$
Le morphisme $V_1 \to V_2$ est donn\'e par un homomorphisme de $B$-alg\`ebres $C_2 \to C_1$
qui, \`a son tour, est donn\'e par des polyn\^omes
$Q_j \in B[x_{1, 1}, \ldots, x_{1, n_1}]$, pour $j = 1, \ldots, n_2$.
Il est alors ais\'e de montrer que le diagramme de faisceaux
$$
\xymatrix{
h_{V_1} \ar[d] \ar[r] & \prod_{i = 1, \ldots, n_1} \mathcal{O}_Y
\ar[d]^{Q_1, \ldots, Q_{n_2}}\\
h_{V_2} \ar[r] & \prod_{i = 1, \ldots, n_2} \mathcal{O}_Y
}
$$
est commutatif et que, par image inverse sur $X_\etale$, nous obtenons le
diagramme commutatif \`a fl\`eches pleines
$$
\xymatrix{
g^{-1}h_{V_1} \ar@{..>}[dd] \ar[rrd] \ar[r] &
\prod_{i = 1, \ldots, n_1} g^{-1}\mathcal{O}_Y
\ar[dd]^{g^\sharp}
\ar[rrd]^{Q_1, \ldots, Q_{n_2}} \\
& & g^{-1}h_{V_2} \ar@{..>}[dd] \ar[r] &
\prod_{i = 1, \ldots, n_2} g^{-1}\mathcal{O}_Y
\ar[dd]^{g^\sharp} \\
h_{X \times_Y V_1} \ar[r] \ar[rrd] &
\prod\nolimits_{i = 1, \ldots, n_1} \mathcal{O}_X
\ar[rrd]^{Q'_1, \ldots, Q'_{n_2}} \\
& & h_{X \times_Y V_2} \ar[r] &
\prod\nolimits_{i = 1, \ldots, n_2} \mathcal{O}_X
}
$$
o\`u $Q'_j \in A[x_{1, 1}, \ldots, x_{1, n_1}]$ est l'image de
$Q_j$ par $\varphi$. Puisque les fl\`eches pointill\'ees existent, rendent les
deux carr\'es commutatifs, et que les fl\`eches horizontales sont injectives,
nous voyons que le diagramme entier est commutatif. Cela prouve la fonctorialit\'e
(ainsi que l'ind\'ependance de la construction de $g^{-1}h_V \to h_{X \times_Y V}$
\`a l'\'egard du choix de la pr\'esentation, bien qu'\`a
strictement parler nous n'ayons pas besoin de le montrer).

\medskip\noindent
\`A ce stade, nous pouvons montrer que $f_{small, *} \cong g_*$.
En effet, soit $\mathcal{F}$ un faisceau sur $X_\etale$. Pour tout objet affine
$V \in \Ob(X_\etale)$, nous avons
\begin{align*}
(g_*\mathcal{F})(V)
& =
\Mor_{\Sh(Y_\etale)}(h_V, g_*\mathcal{F}) \\
& =
\Mor_{\Sh(X_\etale)}(g^{-1}h_V, \mathcal{F}) \\
& =
\Mor_{\Sh(X_\etale)}(h_{X \times_Y V}, \mathcal{F}) \\
& =
\mathcal{F}(X \times_Y V) \\
& =
f_{small, *}\mathcal{F}(V)
\end{align*}
o\`u, dans la troisi\`eme \'egalit\'e, nous utilisons l'isomorphisme
$g^{-1}h_V \cong h_{X \times_Y V}$ construit ci-dessus. Ces isomorphismes
sont manifestement fonctoriels en $\mathcal{F}$ et en $V$,
car les isomorphismes $g^{-1}h_V \cong h_{X \times_Y V}$ sont fonctoriels.
Or tout faisceau sur $Y_\etale$ est d\'etermin\'e par sa restriction
\`a la sous-cat\'egorie des sch\'emas affines
(Topologies, Lemme \ref{topologies-lemma-alternative}) ;
nous obtenons donc un isomorphisme de foncteurs $f_{small, *} \cong g_*$,
comme voulu.

\medskip\noindent
Enfin, nous devons v\'erifier que, par l'isomorphisme
$f_{small, *} \cong g_*$ ci-dessus, les morphismes $f_{small}^\sharp$ et
$g^\sharp$ co\"incident. Par construction, c'est d\'ej\`a le cas sur les
sections globales de $\mathcal{O}_Y$, c'est-\`a-dire sur les \'el\'ements de $B$.
Il suffit de v\'erifier le r\'esultat sur les
sections au-dessus d'un ouvert affine $V$ \'etale sur $Y$ (encore d'apr\`es
Topologies, Lemme \ref{topologies-lemma-alternative}).
En \'ecrivant
$V = \Spec(C)$, $C = B[x_i]/(P_j)$ comme pr\'ec\'edemment, il suffit
de v\'erifier que les fonctions coordonn\'ees $x_i$ sont envoy\'ees sur
les m\^emes sections de $\mathcal{O}_X$ au-dessus de $X \times_Y V$.
Et c'est exactement ce que signifie la commutativit\'e du diagramme
$$
\xymatrix{
g^{-1}h_V \ar[r] \ar@{..>}[d] &
\prod\nolimits_{i = 1, \ldots, n} g^{-1}\mathcal{O}_Y
\ar[d]^{\prod g^\sharp} \\
h_{X \times_Y V} \ar[r] &
\prod\nolimits_{i = 1, \ldots, n} \mathcal{O}_X
}
$$
Le lemme est donc d\'emontr\'e.
\end{proof}

\noindent
Voici une version pour des sch\'emas quelconques.

\begin{theorem}
\label{theorem-fully-faithful}
Soient $X$, $Y$ des sch\'emas. Soit
$$
(g, g^\#) :
(\Sh(X_\etale), \mathcal{O}_X)
\longrightarrow
(\Sh(Y_\etale), \mathcal{O}_Y)
$$
un morphisme de topos localement annel\'es. Alors il existe un unique
morphisme de sch\'emas $f : X \to Y$ tel que
$(g, g^\#)$ soit isomorphe \`a $(f_{small}, f_{small}^\sharp)$.
Autrement dit, la construction
$$
\Sch \longrightarrow \textit{Topos localement annel\'es},
\quad
X \longrightarrow (X_\etale, \mathcal{O}_X)
$$
est pleinement fid\`ele (les morphismes du membre de droite sont consid\'er\'es \`a $2$-isomorphisme pr\`es).
\end{theorem}

\begin{proof}
On peut d\'emontrer ce th\'eor\`eme en adaptant soigneusement les arguments de
la d\'emonstration du
Lemme \ref{lemma-morphism-ringed-etale-topoi-affines}
au cadre global. Nous voulons toutefois indiquer comment
recoller le r\'esultat de ce lemme pour obtenir un morphisme global,
gr\^ace \`a la rigidit\'e fournie par le r\'esultat du
Lemme \ref{lemma-2-morphism}.
Malheureusement, les d\'etails sont quelque peu laborieux.

\medskip\noindent
Montrons l'existence lorsque $Y$ est affine. Choisissons alors un
recouvrement ouvert affine $X = \bigcup U_i$. Pour chaque $i$, le morphisme
d'inclusion $j_i : U_i \to X$ induit un morphisme de topos localement annel\'es
$(j_{i, small}, j_{i, small}^\sharp) :
(\Sh(U_{i, \etale}), \mathcal{O}_{U_i})
\to
(\Sh(X_\etale), \mathcal{O}_X)$
d'apr\`es le
Lemme \ref{lemma-morphism-locally-ringed}.
Nous pouvons le composer avec $(g, g^\sharp)$ pour obtenir un morphisme
de topos localement annel\'es
$$
(g, g^\sharp) \circ (j_{i, small}, j_{i, small}^\sharp) :
(\Sh(U_{i, \etale}), \mathcal{O}_{U_i})
\to
(\Sh(Y_\etale), \mathcal{O}_Y)
$$
voir
Modules on Sites,
Lemme \ref{sites-modules-lemma-composition-morphisms-locally-ringed-topoi}.
D'apr\`es le
Lemme \ref{lemma-morphism-ringed-etale-topoi-affines},
il existe un unique morphisme de sch\'emas $f_i : U_i \to Y$
et un $2$-isomorphisme
$$
t_i :
(f_{i, small}, f_{i, small}^\sharp)
\longrightarrow
(g, g^\sharp) \circ (j_{i, small}, j_{i, small}^\sharp).
$$
Posons $U_{i, i'} = U_i \cap U_{i'}$, et notons $j_{i, i'} : U_{i, i'} \to U_i$
le morphisme d'inclusion. Puisque
$j_i \circ j_{i, i'} = j_{i'} \circ j_{i', i}$
nous avons
\begin{align*}
(g, g^\sharp) \circ
(j_{i, small}, j_{i, small}^\sharp) \circ
(j_{i, i', small}, j_{i, i', small}^\sharp)
= \\
(g, g^\sharp) \circ
(j_{i', small}, j_{i', small}^\sharp) \circ
(j_{i', i, small}, j_{i', i, small}^\sharp)
\end{align*}
Par unicit\'e (voir le
Lemme \ref{lemma-faithful}),
nous en concluons que
$f_i \circ j_{i, i'} = f_{i'} \circ j_{i', i}$ ; autrement dit, les
morphismes de sch\'emas $f_i = f \circ j_i$ sont les restrictions d'un
morphisme de sch\'emas global $f : X \to Y$. Consid\'erons le diagramme
de $2$-isomorphismes (dans lequel nous omettons les composantes ${}^\sharp$ pour all\'eger la
notation)
$$
\xymatrix{
g \circ j_{i, small} \circ j_{i, i', small}
\ar[rr]^{t_i \star \text{id}_{j_{i, i', small}}}
\ar@{=}[d] & &
f_{small} \circ j_{i, small} \circ j_{i, i', small} \ar@{=}[d] \\
g \circ j_{i', small} \circ j_{i', i, small}
\ar[rr]^{t_{i'} \star \text{id}_{j_{i', i, small}}} & &
f_{small} \circ j_{i', small} \circ j_{i', i, small}
}
$$
La notation $\star$ d\'esigne la composition horizontale ; voir
Categories, D\'efinition \ref{categories-definition-2-category}
dans le cas g\'en\'eral et
Sites, Section \ref{sites-section-2-category}
dans notre cas particulier. D'apr\`es le r\'esultat du
Lemme \ref{lemma-2-morphism},
ce diagramme est commutatif. Ainsi, pour tout faisceau $\mathcal{G}$
sur $Y_\etale$, les isomorphismes
$t_i : f_{small}^{-1}\mathcal{G}|_{U_i} \to g^{-1}\mathcal{G}|_{U_i}$
co\"incident au-dessus de $U_{i, i'}$ et nous obtenons un isomorphisme global
$t : f_{small}^{-1}\mathcal{G} \to g^{-1}\mathcal{G}$.
Il est clair que cet isomorphisme est fonctoriel en $\mathcal{G}$
et compatible avec les morphismes $f_{small}^\sharp$ et $g^\sharp$
(puisqu'il est localement compatible avec ces morphismes).
Cela d\'emontre le th\'eor\`eme lorsque $Y$ est affine.

\medskip\noindent
Dans le cas g\'en\'eral, soit $V \subset Y$ un ouvert affine.
Alors $h_V$ est un sous-faisceau du faisceau final $*$ sur $Y_\etale$.
Comme $g$ est exact, nous voyons que $g^{-1}h_V$ est un sous-faisceau du faisceau
final sur $X_\etale$. Ainsi, d'apr\`es le
Lemme \ref{lemma-support-subsheaf-final},
il existe un sous-sch\'ema ouvert $W \subset X$ tel que $g^{-1}h_V = h_W$. D'apr\`es
Modules on Sites,
Lemme \ref{sites-modules-lemma-localize-morphism-locally-ringed-topoi},
il existe un diagramme commutatif de morphismes de topos localement
annel\'es
$$
\xymatrix{
(\Sh(W_\etale), \mathcal{O}_W) \ar[r] \ar[d]_{g'} &
(\Sh(X_\etale), \mathcal{O}_X) \ar[d]^g \\
(\Sh(V_\etale), \mathcal{O}_V) \ar[r] &
(\Sh(Y_\etale), \mathcal{O}_Y)
}
$$
o\`u les fl\`eches horizontales sont les morphismes de localisation
(induits par les morphismes d'inclusion $V \to Y$ et $W \to X$)
et o\`u $g'$ est induit par $g$. Le paragraphe pr\'ec\'edent
nous fournit un morphisme de sch\'emas $f' : W \to V$ et
un $2$-isomorphisme
$t : (f'_{small}, (f'_{small})^\sharp) \to (g', (g')^\sharp)$.
Exactement comme pr\'ec\'edemment, ces morphismes $f'$ (pour les divers ouverts affines $V \subset Y$)
co\"incident sur les intersections par unicit\'e, et donnent donc un morphisme $f : X \to Y$.
De plus, les $2$-isomorphismes $t$ sont compatibles sur les intersections, encore d'apr\`es le
Lemme \ref{lemma-2-morphism},
et nous obtenons un $2$-isomorphisme global
$(f_{small}, (f_{small})^\sharp) \to (g, (g)^\sharp)$.
comme voulu. Nous omettons certains d\'etails.
\end{proof}










\section{Image directe et image inverse}
\label{section-monomorphisms}

\noindent
Soit $f : X \to Y$ un morphisme de sch\'emas.
Voici une liste de conditions que nous consid\'ererons dans la suite :
\begin{enumerate}
\item[(A)] Pour tout morphisme \'etale $U \to X$ et tout $u \in U$, il existe
un morphisme \'etale $V \to Y$, une d\'ecomposition en union disjointe
$X \times_Y V = W \amalg W'$ et un morphisme $h : W \to U$ au-dessus de $X$
tel que $u$ appartienne \`a l'image de $h$.
\item[(B)] Pour tout morphisme \'etale $V \to Y$ et tout recouvrement \'etale
$\{U_i \to X \times_Y V\}$, il existe un recouvrement \'etale
$\{V_j \to V\}$ tel que, pour chaque $j$, on ait
$X \times_Y V_j = \coprod W_{ij}$, o\`u $W_{ij} \to X \times_Y V$
se factorise par $U_i \to X \times_Y V$ pour un certain $i$.
\item[(C)] Pour tout morphisme \'etale $U \to X$, il existe un morphisme \'etale $V \to Y$
et un morphisme surjectif $X \times_Y V \to U$ au-dessus de $X$.
\end{enumerate}
Il s'av\`ere que chacune de ces propri\'et\'es s'interpr\`ete en termes du
comportement du foncteur $f_{small, *}$. Nous expliciterons cela
dans les quelques sections suivantes.



\section{Propri\'et\'e (A)}
\label{section-A}

\noindent
Voir la Section \ref{section-monomorphisms} pour la d\'efinition de la propri\'et\'e
(A).

\begin{lemma}
\label{lemma-property-A-implies}
Soit $f : X \to Y$ un morphisme de sch\'emas.
Supposons (A).
\begin{enumerate}
\item
$f_{small, *} :
\textit{Ab}(X_\etale)
\to
\textit{Ab}(Y_\etale)$
refl\`ete les injections et les surjections,
\item $f_{small}^{-1}f_{small, *}\mathcal{F} \to \mathcal{F}$
est surjectif pour tout faisceau ab\'elien $\mathcal{F}$ sur $X_\etale$,
\item
$f_{small, *} :
\textit{Ab}(X_\etale)
\to
\textit{Ab}(Y_\etale)$
est fid\`ele.
\end{enumerate}
\end{lemma}

\begin{proof}
Soit $\mathcal{F}$ un faisceau ab\'elien sur $X_\etale$.
Soit $U$ un objet de $X_\etale$. Par hypoth\`ese, nous pouvons trouver un
recouvrement $\{W_i \to U\}$ dans $X_\etale$ tel que chaque $W_i$ soit
un sous-sch\'ema ouvert et ferm\'e de $X \times_Y V_i$ pour un certain objet
$V_i$ de $Y_\etale$. La condition de faisceau montre que
$$
\mathcal{F}(U) \subset \prod \mathcal{F}(W_i)
$$
et que $\mathcal{F}(W_i)$ est un facteur direct de
$\mathcal{F}(X \times_Y V_i) = f_{small, *}\mathcal{F}(V_i)$.
Il est donc clair que $f_{small, *}$ refl\`ete les injections.

\medskip\noindent
Supposons ensuite que $a : \mathcal{G} \to \mathcal{F}$ soit un morphisme de
faisceaux ab\'eliens tel que $f_{small, *}a$ soit surjectif. Soit
$s \in \mathcal{F}(U)$, avec $U$ comme ci-dessus. Avec les $W_i$, $V_i$ ci-dessus,
nous voyons qu'il suffit de montrer que $s|_{W_i}$ est, localement pour la topologie \'etale,
l'image d'une section de $\mathcal{G}$ par $a$. Comme $\mathcal{F}(W_i)$
est un facteur direct de $\mathcal{F}(X \times_Y V_i)$,
il suffit de montrer que, pour tout $V \in \Ob(Y_\etale)$,
tout \'el\'ement $s \in \mathcal{F}(X \times_Y V)$
est, localement pour la topologie \'etale sur $X \times_Y V$, l'image d'une section de
$\mathcal{G}$ par $a$. Puisque
$\mathcal{F}(X \times_Y V) = f_{small, *}\mathcal{F}(V)$
nous voyons par hypoth\`ese qu'il existe un recouvrement $\{V_j \to V\}$ tel que
$s$ soit l'image de
$s_j \in f_{small, *}\mathcal{G}(V_j) = \mathcal{G}(X \times_Y V_j)$.
Cela prouve que $f_{small, *}$ refl\`ete les surjections.

\medskip\noindent
Les assertions (2) et (3) d\'ecoulent formellement de l'assertion (1) ; voir
Modules on Sites, Lemme \ref{sites-modules-lemma-reflect-surjections}.
\end{proof}

\begin{lemma}
\label{lemma-locally-quasi-finite-A}
Soit $f : X \to Y$ un morphisme de sch\'emas s\'epar\'e et localement quasi-fini.
Alors la propri\'et\'e (A) ci-dessus est satisfaite.
\end{lemma}

\begin{proof}
Soient $U \to X$ un morphisme \'etale et $u \in U$.
L'\'enonc\'e g\'eom\'etrique (A) se r\'eduit directement au cas o\`u $U$ et $Y$
sont des sch\'emas affines. Notons $x \in X$ et $y \in Y$ les
images de $u$. Comme $X \to Y$ est localement quasi-fini et que $U \to X$ est
localement quasi-fini (voir
Morphisms, Lemme \ref{morphisms-lemma-etale-locally-quasi-finite}),
nous voyons que $U \to Y$ est localement quasi-fini (voir
Morphisms, Lemme \ref{morphisms-lemma-composition-quasi-finite}).
De plus, $X \to Y$ et $U \to Y$ sont tous deux s\'epar\'es. Ainsi, le
More on Morphisms, Lemme
\ref{more-morphisms-lemma-etale-splits-off-quasi-finite-part-technical-variant}
s'applique aux deux morphismes. Nous pouvons donc choisir un voisinage \'etale
$(V, v) \to (Y, y)$ tel que
$$
X \times_Y V = W \amalg R, \quad
U \times_Y V = W' \amalg R'
$$
et des points $w \in W$, $w' \in W'$ tels que
\begin{enumerate}
\item $W$, $R$ sont ouverts et ferm\'es dans $X \times_Y V$,
\item $W'$, $R'$ sont ouverts et ferm\'es dans $U \times_Y V$,
\item $W \to V$ et $W' \to V$ sont finis,
\item $w$, $w'$ s'envoient sur $v$,
\item $\kappa(v) \subset \kappa(w)$ et $\kappa(v) \subset \kappa(w')$
sont purement ins\'eparables, et
\item aucun autre point de $W$ ou de $W'$ ne s'envoie sur $v$.
\end{enumerate}
Voici un diagramme commutatif
$$
\xymatrix{
U \ar[d] & U \times_Y V \ar[l] \ar[d] & W' \amalg R' \ar[d] \ar[l] \\
X \ar[d] & X \times_Y V \ar[l] \ar[d] & W \amalg R \ar[l] \\
Y & V \ar[l]
}
$$
Apr\`es avoir r\'etr\'eci $V$, nous pouvons supposer que $W'$ s'envoie dans $W$ :
il suffit d'enlever l'image de l'image inverse de $R$ dans $W'$ ; c'est
un ferm\'e (puisque $W' \to V$ est fini) qui ne contient pas $v$.
Alors $W' \to W$ est fini, puisque $W \to V$ et $W' \to V$ sont finis.
Ainsi $W' \to W$ est fini \'etale, et la
fibre au-dessus de $w$ contient exactement un point, avec $\kappa(w) = \kappa(w')$. Par cons\'equent, $W' \to W$ est un
isomorphisme sur un voisinage ouvert $W^\circ$ de $w$ ; voir
\'Etale Morphisms, Lemme \ref{etale-lemma-finite-etale-one-point}.
Comme $W \to V$ est fini, l'image de $W \setminus W^\circ$ est un sous-ensemble
ferm\'e $T$ de $V$ qui ne contient pas $v$. Apr\`es avoir remplac\'e $V$ par
$V \setminus T$, nous pouvons donc supposer que $W' \to W$ est un isomorphisme.
La d\'ecomposition $X \times_Y V = W \amalg R$ et le morphisme
$W \to U$ sont alors ceux recherch\'es, ce qui ach\`eve la d\'emonstration.
\end{proof}

\begin{lemma}
\label{lemma-integral-A}
Soit $f : X \to Y$ un morphisme entier de sch\'emas.
Alors la propri\'et\'e (A) est satisfaite.
\end{lemma}

\begin{proof}
Soient $U \to X$ \'etale et $u \in U$ un point.
Nous devons trouver un morphisme \'etale $V \to Y$, une d\'ecomposition en union disjointe
$X \times_Y V = W \amalg W'$ et un $X$-morphisme $W \to U$
tel que $u$ appartienne \`a son image. Nous pouvons r\'etr\'ecir $U$ et $Y$ et supposer
que $U$ et $Y$ sont affines. Alors $X$ est \'egalement affine, puisqu'un
morphisme entier est affine par d\'efinition. \'Ecrivons $Y = \Spec(A)$,
$X = \Spec(B)$ et $U = \Spec(C)$. Alors $A \to B$ est un
homomorphisme d'anneaux entier, et $B \to C$ un homomorphisme d'anneaux \'etale. D'apr\`es
Algebra, Lemme \ref{algebra-lemma-etale},
nous pouvons trouver une sous-$A$-alg\`ebre finie $B' \subset B$ et un homomorphisme d'anneaux \'etale
$B' \to C'$ tels que $C = B \otimes_{B'} C'$. La question
se ram\`ene donc au morphisme \'etale
$U' = \Spec(C') \to X' = \Spec(B')$
au-dessus du morphisme fini $X' \to Y$. Dans ce cas, le r\'esultat d\'ecoule du
Lemme \ref{lemma-locally-quasi-finite-A}.
\end{proof}

\begin{lemma}
\label{lemma-when-push-pull-surjective}
Soit $f : X \to Y$ un morphisme de sch\'emas. Notons
$f_{small} :
\Sh(X_\etale)
\to
\Sh(Y_\etale)$
le morphisme associ\'e de petits topos \'etales. Supposons qu'au moins l'une
des hypoth\`eses suivantes soit satisfaite :
\begin{enumerate}
\item $f$ est entier, ou
\item $f$ est s\'epar\'e et localement quasi-fini.
\end{enumerate}
Alors le foncteur
$f_{small, *} :
\textit{Ab}(X_\etale)
\to
\textit{Ab}(Y_\etale)$
poss\`ede les propri\'et\'es suivantes :
\begin{enumerate}
\item le morphisme
$f_{small}^{-1}f_{small, *}\mathcal{F} \to \mathcal{F}$
est toujours surjectif,
\item $f_{small, *}$ est fid\`ele, et
\item $f_{small, *}$ refl\`ete les injections et les surjections.
\end{enumerate}
\end{lemma}

\begin{proof}
Il suffit de combiner les
Lemmes \ref{lemma-locally-quasi-finite-A},
\ref{lemma-integral-A}, et
\ref{lemma-property-A-implies}.
\end{proof}



\section{Propri\'et\'e (B)}
\label{section-B}

\noindent
Voir la Section \ref{section-monomorphisms} pour la d\'efinition de la propri\'et\'e
(B).

\begin{lemma}
\label{lemma-property-B-implies}
Soit $f : X \to Y$ un morphisme de sch\'emas. Supposons la propri\'et\'e (B) satisfaite.
Alors le foncteur
$f_{small, *} :
\Sh(X_\etale)
\to
\Sh(Y_\etale)$
envoie les morphismes surjectifs sur des morphismes surjectifs.
\end{lemma}

\begin{proof}
Cela d\'ecoule de
Sites, Lemme \ref{sites-lemma-weaker}.
\end{proof}

\begin{lemma}
\label{lemma-simplify-B}
Soit $f : X \to Y$ un morphisme de sch\'emas. Supposons que
\begin{enumerate}
\item $V \to Y$ soit un morphisme \'etale de sch\'emas,
\item $\{U_i \to X \times_Y V\}$ soit un recouvrement \'etale, et
\item $v \in V$ soit un point.
\end{enumerate}
Supposons que, pour toute donn\'ee de ce type, il existe un voisinage \'etale
$(V', v') \to (V, v)$, une d\'ecomposition en union disjointe
$X \times_Y V' = \coprod W'_i$ et des morphismes $W'_i \to U_i$
au-dessus de $X \times_Y V$. Alors la propri\'et\'e (B) est satisfaite.
\end{lemma}

\begin{proof}
D\'emonstration omise.
\end{proof}

\begin{lemma}
\label{lemma-finite-B}
Soit $f : X \to Y$ un morphisme fini de sch\'emas.
Alors la propri\'et\'e (B) est satisfaite.
\end{lemma}

\begin{proof}
Consid\'erons un morphisme \'etale $V \to Y$, un recouvrement \'etale $\{U_i \to X \times_Y V\}$ et
un point $v \in V$. Nous devons trouver $V' \to V$, une d\'ecomposition et des morphismes comme dans le
Lemme \ref{lemma-simplify-B}.
Nous pouvons r\'etr\'ecir $V$ et $Y$, donc supposer que $V$ et $Y$ sont affines.
Comme $X$ est fini sur $Y$, il en r\'esulte aussi que $X$ est affine.
Au cours de la d\'emonstration, nous pouvons remplacer (un nombre fini de fois) $(V, v)$ par un
voisinage \'etale $(V', v')$ et, de mani\`ere correspondante, le recouvrement
$\{U_i \to X \times_Y V\}$ par $\{V' \times_V U_i \to X \times_Y V'\}$.

\medskip\noindent
Comme $X \times_Y V \to V$ est fini, il existe un nombre fini de
points deux \`a deux distincts $x_1, \ldots, x_n \in X \times_Y V$
qui s'envoient sur $v$. Nous pouvons appliquer
More on Morphisms, Lemme
\ref{more-morphisms-lemma-etale-splits-off-quasi-finite-part-technical-variant}
\`a $X \times_Y V \to V$ et aux points $x_1, \ldots, x_n$ situ\'es au-dessus de
$v$, et trouver un voisinage \'etale $(V', v') \to (V, v)$
tel que
$$
X \times_Y V' = R \amalg \coprod T_a
$$
o\`u $T_a \to V'$ est fini, avec exactement un point $p_a$ au-dessus de $v'$,
o\`u de plus $\kappa(v') \subset \kappa(p_a)$ est purement ins\'eparable, et
tel que la fibre de $R \to V'$ au-dessus de $v'$ soit vide.
Puisque $X \to Y$ est fini, $R \to V'$ l'est aussi. Apr\`es avoir
r\'etr\'eci $V'$, nous pouvons donc supposer que $R = \emptyset$. Ainsi, nous pouvons
supposer que $X \times_Y V = X_1 \amalg \ldots \amalg X_n$, avec
exactement un point $x_l \in X_l$ au-dessus de $v$, et tel que, de plus,
$\kappa(v) \subset \kappa(x_l)$ soit purement ins\'eparable. Remarquons que cette
propri\'et\'e est pr\'eserv\'ee par tout raffinement du voisinage \'etale
$(V, v)$.

\medskip\noindent
Pour chaque $l$, choisissons un $i_l$ et un point $u_l \in U_{i_l}$ qui s'envoie sur $x_l$.
Appliquons maintenant la propri\'et\'e (A) au morphisme fini
$X \times_Y V \to V$, aux morphismes \'etales
$U_{i_l} \to X \times_Y V$ et aux points $u_l$.
Cela est permis par le
Lemme \ref{lemma-integral-A}.
Nous obtenons ainsi un voisinage \'etale $(V', v') \to (V, v)$
et des d\'ecompositions
$$
X \times_Y V' = W_l \amalg R_l
$$
et des $X$-morphismes $a_l : W_l \to U_{i_l}$ dont l'image contient $u_{i_l}$.
Voici une figure :
$$
\xymatrix{
& & & U_{i_l} \ar[d] & \\
W_l \ar[rrru] \ar[r] & W_l \amalg R_l \ar@{=}[r] &
X \times_Y V' \ar[r] \ar[d] &
X \times_Y V \ar[r] \ar[d] & X \ar[d] \\
& & V' \ar[r] & V \ar[r] & Y
}
$$
Apr\`es avoir remplac\'e $(V, v)$ par $(V', v')$, nous concluons que chaque
$x_l$ appartient \`a un voisinage ouvert et ferm\'e $W_l$ tel que
le morphisme d'inclusion $W_l \to X \times_Y V$ se factorise par
$U_i \to X \times_Y V$ pour un certain $i$. En rempla\c{c}ant $W_l$ par $W_l \cap X_l$,
nous voyons que ces ouverts et ferm\'es sont disjoints et, de plus,
que $\{x_1, \ldots, x_n\} \subset W_1 \cup \ldots \cup W_n$.
Comme $X \times_Y V \to V$ est fini, nous pouvons r\'etr\'ecir $V$ et supposer que
$X \times_Y V = W_1 \amalg \ldots \amalg W_n$, comme voulu.
\end{proof}

\begin{lemma}
\label{lemma-integral-B}
Soit $f : X \to Y$ un morphisme entier de sch\'emas.
Alors la propri\'et\'e (B) est satisfaite.
\end{lemma}

\begin{proof}
Consid\'erons un morphisme \'etale $V \to Y$, un recouvrement \'etale $\{U_i \to X \times_Y V\}$ et
un point $v \in V$. Nous devons trouver $V' \to V$, une d\'ecomposition et des morphismes comme dans le
Lemme \ref{lemma-simplify-B}.
Nous pouvons r\'etr\'ecir $V$ et $Y$, donc supposer que $V$ et $Y$ sont affines.
Comme $X$ est entier sur $Y$, il en r\'esulte aussi que $X$ et
$X \times_Y V$ sont affines. Nous pouvons raffiner le recouvrement
$\{U_i \to X \times_Y V\}$ et donc supposer que
$\{U_i \to X \times_Y V\}_{i = 1, \ldots, n}$ est un recouvrement \'etale standard.
\'Ecrivons $Y = \Spec(A)$, $X = \Spec(B)$,
$V = \Spec(C)$ et $U_i = \Spec(B_i)$.
Alors $A \to B$ est un homomorphisme d'anneaux entier, et les $B \otimes_A C \to B_i$ sont des
homomorphismes d'anneaux \'etales. D'apr\`es
Algebra, Lemme \ref{algebra-lemma-etale},
nous pouvons trouver une sous-$A$-alg\`ebre finie $B' \subset B$ et un homomorphisme d'anneaux \'etale
$B' \otimes_A C \to B'_i$, pour $i = 1, \ldots, n$,
tels que $B_i = B \otimes_{B'} B'_i$. La question
se ram\`ene donc au recouvrement \'etale
$\{\Spec(B'_i) \to X' \times_Y V\}_{i = 1, \ldots, n}$
o\`u $X' = \Spec(B')$ est fini sur $Y$.
Dans ce cas, le r\'esultat d\'ecoule du
Lemme \ref{lemma-finite-B}.
\end{proof}

\begin{lemma}
\label{lemma-what-integral}
Soit $f : X \to Y$ un morphisme de sch\'emas.
Supposons $f$ entier (par exemple fini).
Alors
\begin{enumerate}
\item $f_{small, *}$ envoie les morphismes surjectifs sur des morphismes surjectifs (pour les faisceaux
d'ensembles et pour les faisceaux ab\'eliens),
\item $f_{small}^{-1}f_{small, *}\mathcal{F} \to \mathcal{F}$
est surjectif pour tout faisceau ab\'elien $\mathcal{F}$ sur $X_\etale$,
\item
$f_{small, *} :
\textit{Ab}(X_\etale)
\to
\textit{Ab}(Y_\etale)$
est fid\`ele et refl\`ete les injections et les surjections, et
\item
$f_{small, *} :
\textit{Ab}(X_\etale)
\to
\textit{Ab}(Y_\etale)$
est exact.
\end{enumerate}
\end{lemma}

\begin{proof}
Nous avons vu les assertions (2) et (3) dans le
Lemme \ref{lemma-when-push-pull-surjective}.
L'assertion (1) d\'ecoule des
Lemmes \ref{lemma-integral-B} et \ref{lemma-property-B-implies}.
L'assertion (4) est une cons\'equence de l'assertion (1) ; voir
Modules on Sites, Lemme \ref{sites-modules-lemma-exactness}.
\end{proof}








\section{Propri\'et\'e (C)}
\label{section-C}

\noindent
Voir la Section \ref{section-monomorphisms} pour la d\'efinition de la propri\'et\'e
(C).

\begin{lemma}
\label{lemma-property-C-implies}
Soit $f : X \to Y$ un morphisme de sch\'emas. Supposons (C) satisfaite. Alors le foncteur
$f_{small, *} :
\Sh(X_\etale)
\to
\Sh(Y_\etale)$
refl\`ete les injections et les surjections.
\end{lemma}

\begin{proof}
Cela d\'ecoule de
Sites, Lemme \ref{sites-lemma-cover-from-below}.
Nous omettons la v\'erification du fait que la propri\'et\'e (C) implique que le foncteur
$Y_\etale \to X_\etale$, $V \mapsto X \times_Y V$
satisfait l'hypoth\`ese de
Sites, Lemme \ref{sites-lemma-cover-from-below}.
\end{proof}

\begin{remark}
\label{remark-property-C-strong}
La propri\'et\'e (C) est satisfaite si $f : X \to Y$ est une immersion ouverte. En effet, si
$U \in \Ob(X_\etale)$, nous pouvons aussi consid\'erer $U$ comme un objet
de $Y_\etale$ et $U \times_Y X = U$. Ainsi, la propri\'et\'e (C)
n'implique pas que $f_{small, *}$ soit exact, car ce n'est pas
le cas pour les immersions ouvertes (en g\'en\'eral).
\end{remark}

\begin{lemma}
\label{lemma-property-C-closed-implies}
Soit $f : X \to Y$ un morphisme de sch\'emas. Supposons que,
pour tout morphisme \'etale $V \to Y$, on ait
\begin{enumerate}
\item $X \times_Y V \to V$ poss\`ede la propri\'et\'e (C), et
\item $X \times_Y V \to V$ est ferm\'e.
\end{enumerate}
Alors le foncteur
$Y_\etale \to X_\etale$, $V \mapsto X \times_Y V$
est presque cocontinu ; voir
Sites, D\'efinition \ref{sites-definition-almost-cocontinuous}.
\end{lemma}

\begin{proof}
Soit $V \to Y$ un objet de $Y_\etale$ et soit
$\{U_i \to X \times_Y V\}_{i \in I}$ un recouvrement de $X_\etale$.
Par l'hypoth\`ese (1), pour chaque $i$, nous pouvons trouver un morphisme \'etale
$h_i : V_i \to V$ et un morphisme surjectif $X \times_Y V_i \to U_i$
au-dessus de $X \times_Y V$. Remarquons que $\bigcup h_i(V_i) \subset V$ est un
ouvert qui contient le ferm\'e $Z = \Im(X \times_Y V \to V)$.
Soit $h_0 : V_0 = V \setminus Z \to V$ l'immersion ouverte.
Il est clair que $\{V_i \to V\}_{i \in I \cup \{0\}}$ est un
recouvrement \'etale tel que, pour tout $i \in I \cup \{0\}$, on ait
soit $V_i \times_Y X = \emptyset$ (si $i = 0$), soit
que $V_i \times_Y X \to V \times_Y X$ se factorise par $U_i \to X \times_Y V$
(si $i \not = 0$). Le foncteur $Y_\etale \to X_\etale$ est donc
presque cocontinu.
\end{proof}

\begin{lemma}
\label{lemma-integral-homeo-onto-image-C}
Soit $f : X \to Y$ un morphisme entier de sch\'emas qui d\'efinit
un hom\'eomorphisme de $X$ sur un sous-ensemble ferm\'e de $Y$.
Alors la propri\'et\'e (C) est satisfaite.
\end{lemma}

\begin{proof}
Soit $g : U \to X$ un morphisme \'etale. Nous devons trouver un objet
$V \to Y$ de $Y_\etale$ et un morphisme surjectif $X \times_Y V \to U$
au-dessus de $X$. Supposons que, pour tout $u \in U$, nous puissions trouver un objet
$V_u \to Y$ de $Y_\etale$ et un morphisme $h_u : X \times_Y V_u \to U$
au-dessus de $X$ avec $u \in \Im(h_u)$. Nous pouvons alors prendre $V = \coprod V_u$
et $h = \coprod h_u$, ce qui donne le r\'esultat. Ainsi, \'etant donn\'e un point
$u \in U$, construisons une paire $(V_u, h_u)$ comme ci-dessus. Pour ce faire, nous pouvons
r\'etr\'ecir $U$ et supposer que $U$ est affine. Dans ce cas,
$g : U \to X$ est localement quasi-fini. Posons
$g^{-1}(g(\{u\})) = \{u, u_2, \ldots, u_n\}$. Puisqu'il n'existe aucune
sp\'ecialisations $u_i \leadsto u$, nous pouvons remplacer $U$ par un voisinage affine
de sorte que $g^{-1}(g(\{u\})) = \{u\}$.

\medskip\noindent
L'image $g(U) \subset X$ est ouverte ;
par cons\'equent, $f(g(U))$ est localement ferm\'e dans $Y$. Choisissons un ouvert $V \subset Y$ tel
que $f(g(U)) = f(X) \cap V$. Il en r\'esulte que $g$ se factorise par
$X \times_Y V$ et que le morphisme $\{U \to X \times_Y V\}$ ainsi obtenu est un recouvrement
\'etale. Comme $f$ poss\`ede la propri\'et\'e (B), voir le
Lemme \ref{lemma-integral-B},
nous voyons qu'il existe un recouvrement \'etale $\{V_j \to V\}$ tel que
$X \times_Y V_j \to X \times_Y V$ se factorise par $U$.
Il en r\'esulte que $V' = \coprod V_j$ est \'etale sur $Y$ et qu'il existe un
morphisme $h : X \times_Y V' \to U$ dont l'image
se projette surjectivement sur $g(U)$. Comme $u$ est le seul point de sa fibre, il doit
appartenir \`a l'image de $h$, ce qui donne le r\'esultat.
\end{proof}

\noindent
Nous invitons le lecteur \`a consid\'erer le lemme suivant comme une
\'etape\footnote{Une \'etape est un lieu o\`u l'on s'arr\^ete pour manger
et se reposer au cours d'un long voyage.} sur le chemin qui m\`ene au
r\'esultat ultime concernant $f_{small, *}$ pour les morphismes entiers et universellement
injectifs.

\begin{lemma}
\label{lemma-integral-universally-injective}
Soit $f : X \to Y$ un morphisme de sch\'emas. Supposons que $f$ soit
universellement injectif et entier (par exemple une immersion ferm\'ee).
Alors
\begin{enumerate}
\item
$f_{small, *} :
\Sh(X_\etale)
\to
\Sh(Y_\etale)$
refl\`ete les injections et les surjections,
\item
$f_{small, *} :
\Sh(X_\etale)
\to
\Sh(Y_\etale)$
commute aux sommes amalgam\'ees et aux co\'egalisateurs (et, plus g\'en\'eralement,
aux colimites connexes finies),
\item $f_{small, *}$ envoie les morphismes surjectifs sur des morphismes surjectifs (pour les faisceaux
d'ensembles et pour les faisceaux ab\'eliens),
\item le morphisme
$f_{small}^{-1}f_{small, *}\mathcal{F} \to \mathcal{F}$
est surjectif pour tout faisceau (d'ensembles ou de groupes ab\'eliens)
$\mathcal{F}$ sur $X_\etale$,
\item le foncteur $f_{small, *}$ est fid\`ele (sur les faisceaux d'ensembles et
sur les faisceaux ab\'eliens),
\item
$f_{small, *} :
\textit{Ab}(X_\etale)
\to
\textit{Ab}(Y_\etale)$
est exact, et
\item le foncteur
$Y_\etale \to X_\etale$, $V \mapsto X \times_Y V$ est
presque cocontinu.
\end{enumerate}
\end{lemma}

\begin{proof}
D'apr\`es les
Lemmes \ref{lemma-integral-A},
\ref{lemma-integral-B} et
\ref{lemma-integral-homeo-onto-image-C},
nous savons que le morphisme $f$ poss\`ede les propri\'et\'es (A), (B) et (C).
De plus, d'apr\`es le
Lemme \ref{lemma-property-C-closed-implies},
nous savons que le foncteur $Y_\etale \to X_\etale$ est
presque cocontinu. Nous avons maintenant :
\begin{enumerate}
\item la propri\'et\'e (C) implique (1) d'apr\`es le
Lemme \ref{lemma-property-C-implies},
\item la presque cocontinuit\'e implique (2) d'apr\`es
Sites, Lemme \ref{sites-lemma-morphism-of-sites-almost-cocontinuous},
\item la propri\'et\'e (B) implique (3) d'apr\`es le
Lemme \ref{lemma-property-B-implies}.
\end{enumerate}
Les propri\'et\'es (4), (5) et (6) d\'ecoulent formellement des trois premi\`eres ; voir
Sites, Lemme \ref{sites-lemma-exactness-properties}
et
Modules on Sites, Lemme \ref{sites-modules-lemma-exactness}.
Nous avons vu la propri\'et\'e (7) ci-dessus.
\end{proof}




\section{Invariance topologique du petit site \'etale}
\label{section-topological-invariance}

\noindent
Dans le th\'eor\`eme suivant, nous montrons que le petit site \'etale est un invariant
topologique au sens suivant : si $f : X \to Y$ est un morphisme de sch\'emas
qui est un hom\'eomorphisme universel, alors $X_\etale \cong Y_\etale$
en tant que sites. Cela am\'eliore le r\'esultat de
\'Etale Morphisms, Th\'eor\`eme \ref{etale-theorem-remarkable-equivalence}.
Nous d\'emontrons d'abord le r\'esultat pour les morphismes, puis nous \'enon\c{c}ons celui
pour les cat\'egories.

\begin{theorem}
\label{theorem-etale-topological}
Soient $X$ et $Y$ deux sch\'emas sur un sch\'ema de base $S$. Soit
$S' \to S$ un hom\'eomorphisme universel.
Notons $X'$ (resp.\ $Y'$) le changement de base \`a $S'$.
Si $X$ est \'etale sur $S$, alors l'application
$$
\Mor_S(Y, X) \longrightarrow \Mor_{S'}(Y', X')
$$
est bijective.
\end{theorem}

\begin{proof}
Apr\`es changement de base par $Y \to S$, nous pouvons supposer que $Y = S$.
Ainsi, nous pouvons et allons supposer que $X$ et $Y$ sont tous deux \'etales sur $S$.
Autrement dit, le th\'eor\`eme affirme que le foncteur de changement de base
est un foncteur pleinement fid\`ele de la cat\'egorie des sch\'emas \'etales
sur $S$ vers la cat\'egorie des sch\'emas \'etales sur $S'$.

\medskip\noindent
Consid\'erons le foncteur d'oubli
\begin{equation}
\label{equation-descent-etale-forget}
\begin{matrix}
\text{donn\'ees de descente }(X', \varphi')\text{ relativement \`a }S'/S \\
\text{ avec }X'\text{ \'etale sur }S'
\end{matrix}
\longrightarrow
\text{sch\'emas }X'\text{ \'etales sur }S'
\end{equation}
Nous affirmons que ce foncteur est une \'equivalence. D'autre part, le
foncteur
\begin{equation}
\label{equation-descent-etale}
\text{sch\'emas }X\text{ \'etales sur }S \longrightarrow
\begin{matrix}
\text{donn\'ees de descente }(X', \varphi')\text{ relativement \`a }S'/S \\
\text{ avec }X'\text{ \'etale sur }S'
\end{matrix}
\end{equation}
est pleinement fid\`ele d'apr\`es \'Etale Morphisms, Lemme
\ref{etale-lemma-fully-faithful-cases}.
L'affirmation entra\^ine donc le th\'eor\`eme.

\medskip\noindent
D\'emonstration de l'affirmation.
Rappelons qu'un hom\'eomorphisme universel est la m\^eme chose qu'un morphisme
entier, universellement injectif et surjectif ; voir
Morphisms, Lemme \ref{morphisms-lemma-universal-homeomorphism}.
En particulier, la diagonale $\Delta : S' \to S' \times_S S'$ est un \'epaississement
d'apr\`es Morphisms, Lemme \ref{morphisms-lemma-universally-injective}.
Ainsi, d'apr\`es \'Etale Morphisms, Th\'eor\`eme
\ref{etale-theorem-etale-topological},
nous voyons que, pour tout $X' \to S'$ \'etale, il existe un unique isomorphisme
$$
\varphi' : X' \times_S S' \to S' \times_S X'
$$
de sch\'emas \'etales sur $S' \times_S S'$ dont l'image inverse par
$\Delta$ est $\text{id} : X' \to X'$ sur $S'$.
Comme $S' \to S' \times_S S' \times_S S'$
est aussi un \'epaississement (c'est une immersion ferm\'ee bijective),
nous concluons que $(X', \varphi')$ est une donn\'ee de descente relativement \`a $S'/S$.
Le caract\`ere canonique de la construction de $\varphi'$ montre
qu'elle est compatible aux morphismes entre sch\'emas \'etales sur $S'$.
Autrement dit, nous obtenons un quasi-inverse
$X' \mapsto (X', \varphi')$ du foncteur
(\ref{equation-descent-etale-forget}). Cela d\'emontre l'affirmation et
ach\`eve la d\'emonstration du th\'eor\`eme.
\end{proof}

\begin{theorem}
\label{theorem-topological-invariance}
\begin{reference}
\cite[IV Theorem 18.1.2]{EGA}
\end{reference}
Soit $f : X \to Y$ un morphisme de sch\'emas.
Supposons $f$ entier, universellement injectif et surjectif
(c'est-\`a-dire que $f$ est un hom\'eomorphisme universel ; voir
Morphisms, Lemme \ref{morphisms-lemma-universal-homeomorphism}).
Le foncteur
$$
V \longmapsto V_X = X \times_Y V
$$
d\'efinit une \'equivalence de cat\'egories
$$
\{
\text{sch\'emas }V\text{ \'etales sur }Y
\}
\leftrightarrow
\{
\text{sch\'emas }U\text{ \'etales sur }X
\}
$$
\end{theorem}

\noindent
Nous donnons deux d\'emonstrations. La premi\`ere utilise l'effectivit\'e de la descente
des morphismes quasi-compacts, s\'epar\'es et \'etales relativement
aux morphismes entiers surjectifs. La seconde utilise les r\'esultats
sur les propri\'et\'es (A), (B) et (C) discut\'es plus haut dans le chapitre.

\begin{proof}[Premi\`ere d\'emonstration]
D'apr\`es le Th\'eor\`eme \ref{theorem-etale-topological},
nous voyons que le foncteur est pleinement fid\`ele.
Il reste \`a montrer que le foncteur est essentiellement surjectif.
Soit $U \to X$ un morphisme \'etale de sch\'emas.

\medskip\noindent
Supposons le r\'esultat vrai lorsque $U$ et $Y$ sont affines.
Dans ce cas, choisissons un recouvrement ouvert affine
$U = \bigcup U_i$ tel que chaque $U_i$ s'envoie
dans un ouvert affine de $Y$. Par hypoth\`ese (cas affine), nous pouvons
trouver des morphismes \'etales $V_i \to Y$ tels que $X \times_Y V_i \cong U_i$
comme sch\'emas sur $X$. Soit $V_{i, i'} \subset V_i$
le sous-sch\'ema ouvert dont l'espace topologique sous-jacent
correspond \`a $U_i \cap U_{i'}$. Puisque nous avons des isomorphismes
$$
X \times_Y V_{i, i'} \cong U_i \cap U_{i'} \cong X \times_Y V_{i', i}
$$
comme sch\'emas sur $X$, la pleine fid\'elit\'e montre que nous
obtenons des isomorphismes
$\theta_{i, i'} : V_{i, i'} \to V_{i', i}$ de sch\'emas sur $Y$.
Nous omettons de v\'erifier que ces isomorphismes satisfont \`a la
condition de cocycle de Schemes, Section \ref{schemes-section-glueing-schemes}.
En appliquant Schemes, Lemme \ref{schemes-lemma-glue-schemes}
nous obtenons un sch\'ema $V \to Y$ par
recollement des sch\'emas $V_i$ suivant les identifications $\theta_{i, i'}$.
Il est clair que $V \to Y$ est \'etale et que $X \times_Y V \cong U$
par construction.

\medskip\noindent
Ainsi, il suffit de d\'emontrer le lemme dans le cas o\`u $U$ et $Y$ sont affines.
Rappelons que, dans la d\'emonstration du Th\'eor\`eme \ref{theorem-etale-topological},
nous avons montr\'e que $U$ est muni d'une unique donn\'ee de descente
$(U, \varphi)$ relativement \`a $X/Y$. D'apr\`es
\'Etale Morphisms, Proposition \ref{etale-proposition-effective}
(qui s'applique car $U \to X$ est quasi-compact et s\'epar\'e
ainsi qu'\'etale, d'apr\`es notre r\'eduction au cas affine),
il existe un morphisme \'etale $V \to Y$ tel que
$X \times_Y V \cong U$, ce qui ach\`eve la d\'emonstration.
\end{proof}

\begin{proof}[Seconde d\'emonstration]
D'apr\`es le Th\'eor\`eme \ref{theorem-etale-topological},
nous voyons que le foncteur est pleinement fid\`ele.
Il reste \`a montrer que le foncteur est essentiellement surjectif.
Soit $U \to X$ un morphisme \'etale de sch\'emas.

\medskip\noindent
Supposons le r\'esultat vrai lorsque $U$ et $Y$ sont affines.
Dans ce cas, choisissons un recouvrement ouvert affine
$U = \bigcup U_i$ tel que chaque $U_i$ s'envoie
dans un ouvert affine de $Y$. Par hypoth\`ese (cas affine), nous pouvons
trouver des morphismes \'etales $V_i \to Y$ tels que $X \times_Y V_i \cong U_i$
comme sch\'emas sur $X$. Soit $V_{i, i'} \subset V_i$
le sous-sch\'ema ouvert dont l'espace topologique sous-jacent
correspond \`a $U_i \cap U_{i'}$. Puisque nous avons des isomorphismes
$$
X \times_Y V_{i, i'} \cong U_i \cap U_{i'} \cong X \times_Y V_{i', i}
$$
comme sch\'emas sur $X$, la pleine fid\'elit\'e montre que nous
obtenons des isomorphismes
$\theta_{i, i'} : V_{i, i'} \to V_{i', i}$ de sch\'emas sur $Y$.
Nous omettons de v\'erifier que ces isomorphismes satisfont \`a la
condition de cocycle de Schemes, Section \ref{schemes-section-glueing-schemes}.
En appliquant Schemes, Lemme \ref{schemes-lemma-glue-schemes}
nous obtenons un sch\'ema $V \to Y$ par
recollement des sch\'emas $V_i$ suivant les identifications $\theta_{i, i'}$.
Il est clair que $V \to Y$ est \'etale et que $X \times_Y V \cong U$
par construction.

\medskip\noindent
Ainsi, il suffit de d\'emontrer que le foncteur
\begin{equation}
\label{equation-affine-etale}
\{
\text{sch\'emas affines }V\text{ \'etales sur }Y
\}
\leftrightarrow
\{
\text{sch\'emas affines }U\text{ \'etales sur }X
\}
\end{equation}
est essentiellement surjectif lorsque $X$ et $Y$ sont affines.

\medskip\noindent
Soit $U \to X$ un sch\'ema affine \'etale sur $X$.
Nous devons trouver $V \to Y$ \'etale (et affine) tel que $X \times_Y V$
soit isomorphe \`a $U$ sur $X$. Remarquons qu'un morphisme \'etale de sch\'emas affines
a des fibres universellement born\'ees, voir
Morphisms,
Lemmes \ref{morphisms-lemma-etale-locally-quasi-finite} et
\ref{morphisms-lemma-locally-quasi-finite-qc-source-universally-bounded}.
Nous pouvons donc raisonner par r\'ecurrence sur l'entier $n$ qui borne le degr\'e des fibres
de $U \to X$. Voir
Morphisms, Lemme \ref{morphisms-lemma-etale-universally-bounded}
pour une description de cet entier dans le cas d'un morphisme \'etale.
Si $n = 1$, alors $U \to X$ est une immersion ouverte (voir
\'Etale Morphisms, Th\'eor\`eme \ref{etale-theorem-etale-radicial-open}),
et le r\'esultat est clair. Supposons $n > 1$.

\medskip\noindent
D'apr\`es
le Lemme \ref{lemma-integral-homeo-onto-image-C},
il existe un morphisme \'etale de sch\'emas $W \to Y$ et un
morphisme surjectif $W_X \to U$ sur $X$.
Comme $U$ est quasi-compact, nous pouvons remplacer $W$ par une somme disjointe d'un
nombre fini d'ouverts affines de $W$ ; nous pouvons donc supposer que $W$
est lui aussi affine. Voici un diagramme
$$
\xymatrix{
U \ar[d] & U \times_Y W \ar[l] \ar[d] & W_X \amalg R \ar@{=}[l]\\
X \ar[d] & W_X \ar[l] \ar[d] \\
Y & W \ar[l]
}
$$
La d\'ecomposition en somme disjointe provient du fait que, par construction, le
morphisme \'etale de sch\'emas affines $U \times_Y W \to W_X$ admet une section.
Nous voyons alors que le morphisme $R \to X \times_Y W$ est un morphisme \'etale
de sch\'emas affines dont les fibres ont un degr\'e universellement born\'e
par $n - 1$. Par l'hypoth\`ese de r\'ecurrence, il existe donc un sch\'ema
$V' \to W$ \'etale tel que $R \cong W_X \times_W V'$.
En posant $V'' = W \amalg V'$, nous obtenons un sch\'ema $V''$ \'etale sur $W$ dont le
changement de base \`a $W_X$ est isomorphe \`a $U \times_Y W$
sur $X \times_Y W$.

\medskip\noindent
\`A ce stade, nous pouvons utiliser la descente pour trouver $V$ sur $Y$ dont le
changement de base \`a $X$ est isomorphe \`a $U$ sur $X$. En effet, par la pleine
fid\'elit\'e du foncteur (\ref{equation-affine-etale})
correspondant \`a l'hom\'eomorphisme universel
$X \times_Y (W \times_Y W) \to (W \times_Y W)$
il existe un unique isomorphisme $\varphi : V'' \times_Y W \to W \times_Y V''$
dont le changement de base \`a $X \times_Y (W \times_Y W)$ est la donn\'ee de
descente canonique de $U \times_Y W$ sur $X \times_Y W$. En particulier,
$\varphi$ satisfait \`a la condition de cocycle. Ainsi, d'apr\`es
Descent, Lemme \ref{descent-lemma-affine},
nous voyons que $\varphi$ est effective (rappelons que tous les sch\'emas ci-dessus sont affines).
Nous obtenons donc $V \to Y$ et un isomorphisme $V'' \cong W \times_Y V$
tel que la donn\'ee de descente canonique sur $W \times_Y V/W/Y$ co\"incide
avec $\varphi$. Remarquons que $V \to Y$ est \'etale, d'apr\`es
Descent, Lemme \ref{descent-lemma-descending-property-etale}.
De plus, il existe un isomorphisme $V_X \cong U$ obtenu par
descente de l'isomorphisme
$$
V_X \times_X W_X =
X \times_Y V \times_Y W =
(X \times_Y W) \times_W (W \times_Y V) \cong
W_X  \times_W V'' \cong U \times_Y W
$$
dont nous disposons par construction. Nous omettons certains d\'etails.
\end{proof}

\begin{remark}
\label{remark-affine-inside-equivalence}
Dans la situation du
Th\'eor\`eme \ref{theorem-topological-invariance},
il est \`egalement vrai que $V \mapsto V_X$ induit une \'equivalence
entre les morphismes \'etales $V \to Y$ pour lesquels $V$ est affine et
les morphismes \'etales $U \to X$ pour lesquels $U$ est affine.
Cela r\'esulte par exemple de
Limits, Proposition \ref{limits-proposition-affine}.
\end{remark}

\begin{proposition}[Invariance topologique de la cohomologie \'etale]
\label{proposition-topological-invariance}
Soit $X_0 \to X$ un hom\'eomorphisme universel de sch\'emas
(par exemple l'immersion ferm\'ee d\'efinie par un faisceau d'id\'eaux nilpotent).
Alors
\begin{enumerate}
\item les sites \'etales $X_\etale$ et $(X_0)_\etale$ sont isomorphes,
\item les topos \'etales $\Sh(X_\etale)$ et $\Sh((X_0)_\etale)$
sont \'equivalents, et
\item $H^q_\etale(X, \mathcal{F}) = H^q_\etale(X_0, \mathcal{F}|_{X_0})$
pour tout $q$ et
pour tout faisceau ab\'elien $\mathcal{F}$ sur $X_\etale$.
\end{enumerate}
\end{proposition}

\begin{proof}
L'\'equivalence de cat\'egories $X_\etale \to (X_0)_\etale$ est
donn\'ee par le Th\'eor\`eme \ref{theorem-topological-invariance}. Nous omettons
de d\'emontrer que, sous cette \'equivalence, les recouvrements \'etales se correspondent.
Ainsi, (1) est d\'emontr\'e. Les assertions (2) et (3) d\'ecoulent formellement de (1).
\end{proof}






\section{Immersions ferm\'ees et image directe}
\label{section-closed-immersions}

\noindent
Avant d'\'enoncer et de d\'emontrer
Proposition \ref{proposition-closed-immersion-pushforward}
dans toute sa g\'en\'eralit\'e, nous l'\'enon\c{c}ons et la d\'emontrons bri\`evement pour les
immersions ferm\'ees. En effet, certains des arguments pr\'ec\'edents
sont beaucoup plus faciles \`a suivre dans le cas d'une immersion ferm\'ee et
nous les r\'ep\'etons donc ici sous leur forme simplifi\'ee.

\medskip\noindent
Dans la suite de cette section, $i : Z \to X$ est une immersion ferm\'ee.
Le foncteur
$$
\Sch/X \longrightarrow \Sch/Z, \quad
U \longmapsto U_Z = Z \times_X U
$$
sera not\'e $U \mapsto U_Z$, comme indiqu\'e. Comme la propri\'et\'e d'\^etre une immersion ferm\'ee
se conserve par tout changement de base, le sch\'ema $U_Z$ est un sous-sch\'ema ferm\'e
de $U$.

\begin{lemma}
\label{lemma-closed-immersion-almost-full}
Soit $i : Z \to X$ une immersion ferm\'ee de sch\'emas.
Soient $U, U'$ des sch\'emas \'etales sur $X$. Soit $h : U_Z \to U'_Z$
un morphisme sur $Z$. Alors il existe un diagramme
$$
\xymatrix{
U & W \ar[l]_a \ar[r]^b & U'
}
$$
dans $X_\etale$ tel que $a_Z : W_Z \to U_Z$
soit un isomorphisme et $h = b_Z \circ (a_Z)^{-1}$.
\end{lemma}

\begin{proof}
Consid\'erons le sch\'ema $M = U \times_X U'$. Le graphe $\Gamma_h \subset M_Z$
de $h$ est ouvert. Cela r\'esulte, par exemple, de ce que $\Gamma_h$ est l'image d'une
section du morphisme \'etale $\text{pr}_{1, Z} : M_Z \to U_Z$, voir
\'Etale Morphisms, Proposition \ref{etale-proposition-properties-sections}.
Il existe donc un sous-sch\'ema ouvert $W \subset M$ dont l'intersection avec
la partie ferm\'ee $M_Z$ est $\Gamma_h$. Posons $a = \text{pr}_1|_W$
et $b = \text{pr}_2|_W$.
\end{proof}

\begin{lemma}
\label{lemma-closed-immersion-almost-essentially-surjective}
Soit $i : Z \to X$ une immersion ferm\'ee de sch\'emas.
Soit $V \to Z$ un morphisme \'etale de sch\'emas.
Il existe des morphismes \'etales $U_i \to X$ et des morphismes
$U_{i, Z} \to V$ tels que $\{U_{i, Z} \to V\}$
soit un recouvrement de Zariski de $V$.
\end{lemma}

\begin{proof}
Puisqu'il nous suffit de trouver un recouvrement de Zariski de $V$ form\'e de sch\'emas
de la forme $U_Z$, avec $U$ \'etale sur $X$, nous pouvons raisonner localement pour la topologie de Zariski sur $X$
et sur $V$. Nous pouvons donc supposer $X$ et $V$ affines. Dans le cas affine, cela r\'esulte de
Algebra, Lemme \ref{algebra-lemma-lift-etale}.
\end{proof}

\noindent
Si $\overline{x} : \Spec(k) \to X$ est un point g\'eom\'etrique de $X$, alors
ou bien $\overline{x}$ se factorise (de fa\c{c}on unique) par le sous-sch\'ema ferm\'e $Z$, ou bien
$Z_{\overline{x}} = \emptyset$. Si $\overline{x}$ se factorise par $Z$,
nous disons que $\overline{x}$ est un point g\'eom\'etrique de $Z$ (ce qu'il est) et
nous employons la notation ``$\overline{x} \in Z$'' pour l'indiquer.

\begin{lemma}
\label{lemma-stalk-pushforward-closed-immersion}
Soit $i : Z \to X$ une immersion ferm\'ee de sch\'emas.
Soit $\mathcal{G}$ un faisceau d'ensembles sur $Z_\etale$.
Soit $\overline{x}$ un point g\'eom\'etrique de $X$.
Alors
$$
(i_{small, *}\mathcal{G})_{\overline{x}} =
\left\{
\begin{matrix}
* & \text{si} & \overline{x} \not \in Z \\
\mathcal{G}_{\overline{x}} & \text{si} & \overline{x} \in Z
\end{matrix}
\right.
$$
o\`u $*$ d\'esigne un ensemble r\'eduit \`a un \'el\'ement.
\end{lemma}

\begin{proof}
Posons $U = X \setminus Z$.
Notons que $i_{small, *}\mathcal{G}|_{U_\etale} = *$ est l'objet final
de la cat\'egorie des faisceaux \'etales sur $U$, c'est-\`a-dire le faisceau
qui associe un ensemble r\'eduit \`a un \'el\'ement \`a chaque sch\'ema \'etale sur $U$.
Cela explique la valeur de $(i_{small, *}\mathcal{G})_{\overline{x}}$
lorsque $\overline{x} \not \in Z$.

\medskip\noindent
Supposons maintenant que $\overline{x} \in Z$. Notons que
$$
(i_{small, *}\mathcal{G})_{\overline{x}}
=
\colim_{(U, \overline{u})} \mathcal{G}(U_Z)
$$
et, d'autre part,
$$
\mathcal{G}_{\overline{x}}
=
\colim_{(V, \overline{v})} \mathcal{G}(V).
$$
Soit $\mathcal{C}_1 = \{(U, \overline{u})\}^{opp}$ la cat\'egorie oppos\'ee \`a la
cat\'egorie des voisinages \'etales de $\overline{x}$ dans $X$, et soit
$\mathcal{C}_2 = \{(V, \overline{v})\}^{opp}$ la cat\'egorie oppos\'ee \`a la
cat\'egorie des voisinages \'etales de $\overline{x}$ dans $Z$. L'application canonique
$$
\mathcal{G}_{\overline{x}}
\longrightarrow
(i_{small, *}\mathcal{G})_{\overline{x}}
$$
correspond au foncteur $F : \mathcal{C}_1 \to \mathcal{C}_2$,
$F(U, \overline{u}) = (U_Z, \overline{x})$. Or les
lemmes \ref{lemma-closed-immersion-almost-essentially-surjective} et
\ref{lemma-closed-immersion-almost-full}
impliquent que $\mathcal{C}_1$ est cofinale dans $\mathcal{C}_2$, voir
Categories, Definition \ref{categories-definition-cofinal}.
Il en r\'esulte que la fl\`eche ci-dessus est un isomorphisme, voir
Categories, Lemma \ref{categories-lemma-cofinal}.
\end{proof}

\begin{proposition}
\label{proposition-closed-immersion-pushforward}
Soit $i : Z \to X$ une immersion ferm\'ee de sch\'emas.
\begin{enumerate}
\item Le foncteur
$$
i_{small, *} :
\Sh(Z_\etale)
\longrightarrow
\Sh(X_\etale)
$$
est pleinement fid\`ele et son image essentielle est form\'ee des faisceaux d'ensembles
$\mathcal{F}$ sur $X_\etale$ dont la restriction \`a $X \setminus Z$ est
isomorphe \`a $*$, et
\item le foncteur
$$
i_{small, *} :
\textit{Ab}(Z_\etale)
\longrightarrow
\textit{Ab}(X_\etale)
$$
est pleinement fid\`ele et son image essentielle est form\'ee des faisceaux ab\'eliens sur
$X_\etale$ dont le support est contenu dans $Z$.
\end{enumerate}
Dans les deux cas, $i_{small}^{-1}$ est un inverse \`a gauche du foncteur
$i_{small, *}$.
\end{proposition}

\begin{proof}
Traitons le cas des faisceaux d'ensembles.
Pour tout faisceau $\mathcal{G}$ sur $Z$, le morphisme
$i_{small}^{-1}i_{small, *}\mathcal{G} \to \mathcal{G}$
est un isomorphisme d'apr\`es
le lemme \ref{lemma-stalk-pushforward-closed-immersion}
(et
le th\'eor\`eme \ref{theorem-exactness-stalks}).
Il en r\'esulte formellement que $i_{small, *}$ est pleinement fid\`ele, voir
Sites, Lemma \ref{sites-lemma-exactness-properties}.
Il est clair que $i_{small, *}\mathcal{G}|_{U_\etale} \cong *$
o\`u $U = X \setminus Z$. R\'eciproquement, supposons que $\mathcal{F}$
soit un faisceau d'ensembles sur $X$ tel que $\mathcal{F}|_{U_\etale} \cong *$.
Consid\'erons le morphisme d'adjonction
$$
\mathcal{F} \longrightarrow i_{small, *}i_{small}^{-1}\mathcal{F}
$$
En combinant
les lemmes \ref{lemma-stalk-pushforward-closed-immersion} et
\ref{lemma-stalk-pullback},
nous voyons que c'est un isomorphisme. Cela ach\`eve la d\'emonstration de (1).
La d\'emonstration de (2) est identique.
\end{proof}





\section{Morphismes entiers universellement injectifs}
\label{section-integral-universally-injective}

\noindent
Voici la version g\'en\'erale de la
proposition \ref{proposition-closed-immersion-pushforward}.

\begin{proposition}
\label{proposition-integral-universally-injective-pushforward}
Soit $f : X \to Y$ un morphisme entier de sch\'emas
et universellement injectif.
\begin{enumerate}
\item Le foncteur
$$
f_{small, *} :
\Sh(X_\etale)
\longrightarrow
\Sh(Y_\etale)
$$
est pleinement fid\`ele et son image essentielle est form\'ee des faisceaux d'ensembles
$\mathcal{F}$ sur $Y_\etale$ dont la restriction \`a $Y \setminus f(X)$ est
isomorphe \`a $*$, et
\item le foncteur
$$
f_{small, *} :
\textit{Ab}(X_\etale)
\longrightarrow
\textit{Ab}(Y_\etale)
$$
est pleinement fid\`ele et son image essentielle est form\'ee des faisceaux ab\'eliens sur
$Y_\etale$ dont le support est contenu dans $f(X)$.
\end{enumerate}
Dans les deux cas, $f_{small}^{-1}$ est un inverse \`a gauche du foncteur
$f_{small, *}$.
\end{proposition}

\begin{proof}
Nous pouvons factoriser $f$ sous la forme
$$
\xymatrix{
X \ar[r]^h & Z \ar[r]^i & Y
}
$$
o\`u $h$ est entier, universellement injectif et surjectif
et $i : Z \to Y$ est une immersion ferm\'ee.
Appliquons la
proposition \ref{proposition-closed-immersion-pushforward}
\`a $i$ et le
th\'eor\`eme \ref{theorem-topological-invariance}
\`a $h$.
\end{proof}









\section{Grands sites et image directe}
\label{section-big}

\noindent
Dans cette section, nous d\'emontrons quelques r\'esultats techniques sur $f_{big, *}$ pour
certains types de morphismes de sch\'emas.

\begin{lemma}
\label{lemma-monomorphism-big-push-pull}
Soit $\tau \in \{Zariski, \etale, smooth, syntomic, fppf\}$.
Soit $f : X \to Y$ un monomorphisme de sch\'emas.
Alors le morphisme canonique
$f_{big}^{-1}f_{big, *}\mathcal{F} \to \mathcal{F}$
est un isomorphisme pour tout faisceau $\mathcal{F}$ sur
$(\Sch/X)_\tau$.
\end{lemma}

\begin{proof}
Dans ce cas, le foncteur $(\Sch/X)_\tau \to (\Sch/Y)_\tau$
est continu, cocontinu et pleinement fid\`ele. Le r\'esultat d\'ecoule donc de
Sites, Lemma \ref{sites-lemma-back-and-forth}.
\end{proof}

\begin{remark}
\label{remark-push-pull-shriek}
Dans la situation du
lemme \ref{lemma-monomorphism-big-push-pull},
le morphisme canonique
$\mathcal{F} \to f_{big}^{-1}f_{big!}\mathcal{F}$
est un isomorphisme pour tout faisceau d'ensembles $\mathcal{F}$ sur
$(\Sch/X)_\tau$. La d\'emonstration est la m\^eme. Cela
vaut aussi pour les faisceaux ab\'eliens. Cependant, notons
que le foncteur $f_{big!}$ pour les faisceaux ab\'eliens est d\'efini dans
Modules on Sites, Section \ref{sites-modules-section-exactness-lower-shriek}
et est en g\'en\'eral diff\'erent de $f_{big!}$ sur les faisceaux d'ensembles.
Le r\'esultat pour les faisceaux ab\'eliens d\'ecoule de
Modules on Sites, Lemma \ref{sites-modules-lemma-back-and-forth}.
\end{remark}

\begin{lemma}
\label{lemma-closed-immersion-cover-from-below}
Soit $f : X \to Y$ une immersion ferm\'ee de sch\'emas.
Soit $U \to X$ un morphisme syntomique (resp.\ lisse, resp.\ \'etale).
Il existe alors des morphismes syntomiques (resp.\ lisses, resp.\ \'etales)
$V_i \to Y$ et des morphismes $V_i \times_Y X \to U$ tels que
$\{V_i \times_Y X \to U\}$ forme un recouvrement de Zariski de $U$.
\end{lemma}

\begin{proof}
D\'emontrons le lemme dans le cas o\`u $\tau = syntomic$.
La question est locale sur $U$. Nous pouvons donc supposer que $U$ est
un sch\'ema affine dont l'image est contenue dans un ouvert affine de $Y$.
Nous sommes donc ramen\'es \`a d\'emontrer le cas suivant :
$Y = \Spec(A)$, $X = \Spec(A/I)$, et
$U = \Spec(\overline{B})$, o\`u
$A/I \to \overline{B}$ est un morphisme d'anneaux syntomique.
D'apr\`es Algebra, Lemma \ref{algebra-lemma-lift-syntomic}
nous pouvons trouver des \'el\'ements $\overline{g}_i \in \overline{B}$
tels que
$\overline{B}_{\overline{g}_i} = A_i/IA_i$ pour certains morphismes d'anneaux syntomiques
$A \to A_i$.
Cela d\'emontre le lemme dans le cas syntomique.
La d\'emonstration du cas lisse est identique, \`a ceci pr\`es qu'elle utilise
Algebra, Lemma \ref{algebra-lemma-lift-smooth}.
Dans le cas \'etale, on utilise
Algebra, Lemma \ref{algebra-lemma-lift-etale}.
\end{proof}

\begin{lemma}
\label{lemma-prepare-closed-immersion-almost-cocontinuous}
Soit $f : X \to Y$ une immersion ferm\'ee de sch\'emas.
Soit $\{U_i \to X\}$ un recouvrement syntomique (resp.\ lisse, resp.\ \'etale).
Il existe un recouvrement syntomique (resp.\ lisse, resp.\ \'etale) $\{V_j \to Y\}$
tel que, pour chaque $j$, soit $V_j \times_Y X = \emptyset$, soit le
morphisme $V_j \times_Y X \to X$ se factorise par $U_i$ pour un certain $i$.
\end{lemma}

\begin{proof}
Pour chaque $i$, nous pouvons choisir des morphismes syntomiques (resp.\ lisses, resp.\ \'etales)
$g_{ij} : V_{ij} \to Y$ et des morphismes $V_{ij} \times_Y X \to U_i$ au-dessus de $X$,
de sorte que $\{V_{ij} \times_Y X \to U_i\}$ forme un recouvrement de Zariski, voir
le lemme \ref{lemma-closed-immersion-cover-from-below}.
Cela implique en particulier que
$\bigcup_{ij} g_{ij}(V_{ij})$ contient le sous-ensemble ferm\'e $f(X)$.
Par cons\'equent, la famille de morphismes syntomiques (resp.\ lisses, resp.\ \'etales) $g_{ij}$
ainsi que l'immersion ouverte $Y \setminus f(X) \to Y$ forment le recouvrement
syntomique (resp.\ lisse, resp.\ \'etale) voulu de $Y$.
\end{proof}

\begin{lemma}
\label{lemma-closed-immersion-almost-cocontinuous}
Soit $f : X \to Y$ une immersion ferm\'ee de sch\'emas.
Soit $\tau \in \{syntomic, smooth, \etale\}$.
Le foncteur $V \mapsto X \times_Y V$ d\'efinit un foncteur presque
cocontinu (voir
Sites, D\'efinition \ref{sites-definition-almost-cocontinuous})
$(\Sch/Y)_\tau \to (\Sch/X)_\tau$ entre les
grands $\tau$-sites.
\end{lemma}

\begin{proof}
Il faut montrer ce qui suit : \'etant donn\'e un morphisme $V \to Y$
et un recouvrement syntomique quelconque (resp.\ lisse, resp.\ \'etale)
$\{U_i \to X \times_Y V\}$, il existe un
recouvrement lisse (resp.\ lisse, resp.\ \'etale) $\{V_j \to V\}$
tel que, pour tout $j$, soit $X \times_Y V_j$ est vide, soit
$X \times_Y V_j \to Z \times_Y V$ se factorise par l'un des
$U_i$. Cela r\'esulte de l'application du
Lemme \ref{lemma-prepare-closed-immersion-almost-cocontinuous}
ci-dessus \`a l'immersion ferm\'ee $X \times_Y V \to V$.
\end{proof}

\begin{lemma}
\label{lemma-closed-immersion-pushforward-exact}
Soit $f : X \to Y$ une immersion ferm\'ee de sch\'emas.
Soit $\tau \in \{syntomic, smooth, \etale\}$.
\begin{enumerate}
\item Le foncteur d'image directe
$f_{big, *} :
\Sh((\Sch/X)_\tau)
\to
\Sh((\Sch/Y)_\tau)$
commute aux co\'egalisateurs et aux sommes amalgam\'ees.
\item Le foncteur d'image directe
$f_{big, *} :
\textit{Ab}((\Sch/X)_\tau)
\to
\textit{Ab}((\Sch/Y)_\tau)$
est exact.
\end{enumerate}
\end{lemma}

\begin{proof}
Cela r\'esulte de
Sites, Lemme \ref{sites-lemma-morphism-of-sites-almost-cocontinuous},
Modules sur les sites,
Lemme \ref{sites-modules-lemma-morphism-ringed-sites-almost-cocontinuous},
et du
Lemme \ref{lemma-closed-immersion-almost-cocontinuous}
ci-dessus.
\end{proof}

\begin{remark}
\label{remark-fppf-closed-immersion-not-closed}
Dans le Lemme \ref{lemma-closed-immersion-pushforward-exact}, le cas $\tau = fppf$
n'est pas trait\'e. La raison en est qu'\'etant donn\'es un anneau $A$, un id\'eal $I$ et un
homomorphisme d'anneaux $A/I \to \overline{B}$ fid\`element plat et de pr\'esentation finie, il
n'y a aucune raison de penser que l'on puisse trouver {\it un quelconque} homomorphisme d'anneaux
$A \to B$, plat et de pr\'esentation finie, avec $B/IB \not = 0$, tel que $A/I \to B/IB$ se factorise par
$\overline{B}$. Ainsi, la preuve du
Lemme \ref{lemma-closed-immersion-almost-cocontinuous}
ne fonctionne pas pour la topologie fppf.
En fait, il est probablement faux que
$f_{big, *} : \textit{Ab}((\Sch/X)_{fppf})
\to \textit{Ab}((\Sch/Y)_{fppf})$
soit exact lorsque $f$ est une immersion ferm\'ee.
Si vous connaissez un exemple, veuillez envoyer un courriel \`a
\href{mailto:stacks.project@gmail.com}{stacks.project@gmail.com}.
\end{remark}












\section{Exactitude du grand foncteur image directe \`a support propre}
\label{section-exactness-lower-shriek}

\noindent
Il s'agit simplement du r\'esultat technique suivant. Notons que le foncteur $f_{big!}$
n'a absolument rien \`a voir avec la cohomologie \`a support compact en
g\'en\'eral.

\begin{lemma}
\label{lemma-exactness-lower-shriek}
Soit $\tau \in \{Zariski, \etale, smooth, syntomic, fppf\}$.
Soit $f : X \to Y$ un morphisme de sch\'emas. Soit
$$
f_{big} :
\Sh((\Sch/X)_\tau)
\longrightarrow
\Sh((\Sch/Y)_\tau)
$$
le morphisme de topos correspondant comme dans
Topologies, Lemme
\ref{topologies-lemma-morphism-big},
\ref{topologies-lemma-morphism-big-etale},
\ref{topologies-lemma-morphism-big-smooth},
\ref{topologies-lemma-morphism-big-syntomic}, ou
\ref{topologies-lemma-morphism-big-fppf}.
\begin{enumerate}
\item Le foncteur
$f_{big}^{-1} : \textit{Ab}((\Sch/Y)_\tau) \to \textit{Ab}((\Sch/X)_\tau)$
admet un adjoint \`a gauche
$$
f_{big!} : \textit{Ab}((\Sch/X)_\tau) \to \textit{Ab}((\Sch/Y)_\tau)
$$
qui est exact.
\item Le foncteur
$f_{big}^* :
\textit{Mod}((\Sch/Y)_\tau, \mathcal{O})
\to
\textit{Mod}((\Sch/X)_\tau, \mathcal{O})$
admet un adjoint \`a gauche
$$
f_{big!} :
\textit{Mod}((\Sch/X)_\tau, \mathcal{O})
\to
\textit{Mod}((\Sch/Y)_\tau, \mathcal{O})
$$
qui est exact.
\end{enumerate}
De plus, les deux foncteurs $f_{big!}$ co\"incident sur les faisceaux sous-jacents
de groupes ab\'eliens.
\end{lemma}

\begin{proof}
Rappelons que $f_{big}$ est le morphisme de topos associ\'e au
foncteur continu et cocontinu
$u : (\Sch/X)_\tau \to (\Sch/Y)_\tau$, $U/X \mapsto U/Y$.
De plus, on a $f_{big}^{-1}\mathcal{O} = \mathcal{O}$.
Ainsi, l'existence de $f_{big!}$ r\'esulte de
Modules sur les sites, Lemme \ref{sites-modules-lemma-g-shriek-adjoint},
et, respectivement, de
Modules sur les sites, Lemme \ref{sites-modules-lemma-lower-shriek-modules}.
Notons que, si $U$ est un objet de $(\Sch/X)_\tau$, alors le foncteur
$u$ induit une \'equivalence de cat\'egories
$$
u' :
(\Sch/X)_\tau/U
\longrightarrow
(\Sch/Y)_\tau/U
$$
car les deux membres de la fl\`eche sont \'egaux \`a $(\Sch/U)_\tau$.
Ainsi, la co\"incidence de $f_{big!}$ sur les faisceaux ab\'eliens sous-jacents
r\'esulte de la discussion dans
Modules sur les sites, Remarque \ref{sites-modules-remark-when-shriek-equal}.
L'exactitude de $f_{big!}$ r\'esulte de
Modules sur les sites, Lemme \ref{sites-modules-lemma-exactness-lower-shriek}
puisque le foncteur $u$ ci-dessus commute aux produits fibr\'es et aux \'egalisateurs.
\end{proof}

\noindent
D\'emontrons ensuite un lemme technique qui sera utile plus loin pour comparer
les faisceaux de modules sur les diff\'erents sites associ\'es aux champs alg\'ebriques.

\begin{lemma}
\label{lemma-compare-structure-sheaves}
Soit $X$ un sch\'ema. Soit
$\tau \in \{Zariski, \etale, smooth, syntomic, fppf\}$.
Soient $\mathcal{C}_1 \subset \mathcal{C}_2 \subset (\Sch/X)_\tau$ des sous-cat\'egories pleines
ayant les propri\'et\'es suivantes :
\begin{enumerate}
\item Pour tout objet $U/X$ de $\mathcal{C}_t$,
\begin{enumerate}
\item si $\{U_i \to U\}$ est un recouvrement de $(\Sch/X)_\tau$, alors
$U_i/X$ est un objet de $\mathcal{C}_t$,
\item $U \times \mathbf{A}^1/X$ est un objet de $\mathcal{C}_t$.
\end{enumerate}
\item $X/X$ est un objet de $\mathcal{C}_t$.
\end{enumerate}
Munissons $\mathcal{C}_t$ de la structure de site dont les recouvrements sont
exactement les recouvrements $\{U_i \to U\}$ de $(\Sch/X)_\tau$ tels que
$U \in \Ob(\mathcal{C}_t)$. Alors
\begin{enumerate}
\item[(a)] Le foncteur $\mathcal{C}_1 \to \mathcal{C}_2$
est pleinement fid\`ele, continu et cocontinu.
\end{enumerate}
Notons $g : \Sh(\mathcal{C}_1) \to \Sh(\mathcal{C}_2)$ le
morphisme de topos correspondant. Notons $\mathcal{O}_t$ la restriction de $\mathcal{O}$
\`a $\mathcal{C}_t$. Notons $g_!$ le foncteur de
Modules sur les sites, D\'efinition \ref{sites-modules-definition-g-shriek}.
\begin{enumerate}
\item[(b)] Le morphisme canonique $g_!\mathcal{O}_1 \to \mathcal{O}_2$
est un isomorphisme.
\end{enumerate}
\end{lemma}

\begin{proof}
L'assertion (a) est imm\'ediate d'apr\`es les d\'efinitions.
Dans cette preuve, tous les sch\'emas sont des sch\'emas sur $X$ et tous les morphismes de
sch\'emas sont des morphismes de sch\'emas sur $X$. Notons que $g^{-1}$ est
donn\'e par restriction, de sorte que, pour tout objet $U$ de $\mathcal{C}_1$
on a $\mathcal{O}_1(U) = \mathcal{O}_2(U) = \mathcal{O}(U)$.
Rappelons que $g_!\mathcal{O}_1$ est le faisceau associ\'e au pr\'efaisceau
$g_{p!}\mathcal{O}_1$ qui associe \`a $V$ dans $\mathcal{C}_2$ le groupe
$$
\colim_{V \to U} \mathcal{O}(U)
$$
o\`u $U$ parcourt les objets de $\mathcal{C}_1$ et o\`u la colimite est
prise dans la cat\'egorie des groupes ab\'eliens. Nous utiliserons souvent ci-dessous
le fait que, si
$$
V \to U \to U'
$$
sont des morphismes avec $U, U' \in \Ob(\mathcal{C}_1)$
et si $f' \in \mathcal{O}(U')$ se restreint en $f \in \mathcal{O}(U)$,
alors $(V \to U, f)$ et $(V \to U', f')$ d\'efinissent le m\^eme \'el\'ement de la
colimite. De plus, $g_!\mathcal{O}_1 \to \mathcal{O}_2$ envoie l'\'el\'ement
$(V \to U, f)$ simplement sur la restriction de $f$ \`a $V$.

\medskip\noindent
Surjectivit\'e. Soit $V$ un sch\'ema et soit $h \in \mathcal{O}(V)$.
On obtient alors un morphisme $V \to X \times \mathbf{A}^1$ induit par $h$
et le morphisme structural $V \to X$. En \'ecrivant
$\mathbf{A}^1 = \Spec(\mathbf{Z}[x])$, on voit que l'\'el\'ement
$x \in \mathcal{O}(X \times \mathbf{A}^1)$ a pour image inverse
$h$. Puisque $X \times \mathbf{A}^1$ est un objet de $\mathcal{C}_1$
d'apr\`es les hypoth\`eses (1)(b) et (2), on obtient la surjectivit\'e voulue.

\medskip\noindent
Injectivit\'e. Soit $V$ un sch\'ema. Soit
$s = \sum_{i = 1, \ldots, n} (V \to U_i, f_i)$ un \'el\'ement de la colimite
affich\'ee ci-dessus. Pour tout $i$, on peut utiliser le morphisme
$f_i : U_i \to X \times \mathbf{A}^1$
pour voir que $(V \to U_i, f_i)$ d\'efinit le m\^eme \'el\'ement de la colimite que
$(f_i : V \to X \times \mathbf{A}^1, x)$. Nous pouvons alors consid\'erer
$$
f_1 \times \ldots \times f_n : V \to X \times \mathbf{A}^n
$$
et l'on voit que $s$ est \'equivalent dans la colimite \`a
$$
\sum\nolimits_{i = 1, \ldots, n}
(f_1 \times \ldots \times f_n : V \to X \times \mathbf{A}^n, x_i) =
(f_1 \times \ldots \times f_n : V \to X \times \mathbf{A}^n,
x_1 + \ldots + x_n)
$$
Maintenant, si $x_1 + \ldots + x_n$ se restreint \`a z\'ero sur $V$, alors on voit
que $f_1 \times \ldots \times f_n$ se factorise par
$X \times \mathbf{A}^{n - 1} = V(x_1 + \ldots + x_n)$. Ainsi, nous voyons
que $s$ est \'equivalent \`a z\'ero dans la limite inductive.
\end{proof}











%9.29.09
\section{Cohomologie \'etale}
\label{section-etale-cohomology}

\noindent
Dans les sections suivantes, nous d\'emontrons quelques r\'esultats \`el\'ementaires sur la cohomologie \'etale.
Voici une propri\'et\'e connue de la cohomologie des espaces
topologiques et qui vaut \`egalement pour la cohomologie \'etale.

\begin{lemma}[Mayer-Vietoris en cohomologie \'etale]
\label{lemma-mayer-vietoris}
Soit $X$ un sch\'ema. Supposons que $X = U \cup V$ soit la
r\'eunion de deux ouverts. Pour tout faisceau ab\'elien $\mathcal{F}$ sur $X_\etale$
il existe une suite exacte longue de cohomologie
$$
\begin{matrix}
0 \to
H^0_\etale(X, \mathcal{F}) \to
H^0_\etale(U, \mathcal{F}) \oplus H^0_\etale(V, \mathcal{F}) \to
H^0_\etale(U \cap V, \mathcal{F}) \phantom{\to \ldots} \\
\phantom{0} \to H^1_\etale(X, \mathcal{F}) \to
H^1_\etale(U, \mathcal{F}) \oplus H^1_\etale(V, \mathcal{F}) \to
H^1_\etale(U \cap V, \mathcal{F}) \to \ldots
\end{matrix}
$$
Cette suite exacte longue est fonctorielle en $\mathcal{F}$.
\end{lemma}

\begin{proof}
Observons que si $\mathcal{I}$ est un faisceau ab\'elien injectif, alors
$$
0 \to \mathcal{I}(X) \to \mathcal{I}(U) \oplus \mathcal{I}(V) \to
\mathcal{I}(U \cap V) \to 0
$$
est exacte. L'exactitude aux premier et second termes vient de ce que $\mathcal{I}$
est un faisceau. L'exactitude \`a droite vient de ce que
$\mathcal{I}(U) \to \mathcal{I}(U \cap V)$ est surjective d'apr\`es
Cohomologie sur les sites, Lemme
\ref{sites-cohomology-lemma-restriction-along-monomorphism-surjective}.
Une autre mani\`ere de le d\'emontrer serait de montrer que le conoyau
de l'application $\mathcal{I}(U) \oplus \mathcal{I}(V) \to
\mathcal{I}(U \cap V)$ est le premier groupe de cohomologie de {\v C}ech
de $\mathcal{I}$ relativement au recouvrement
$X = U \cup V$, lequel s'annule d'apr\`es les
lemmes \ref{lemma-hom-injective} et \ref{lemma-forget-injectives}.
Ainsi, si $\mathcal{F} \to \mathcal{I}^\bullet$ est une
r\'esolution injective, alors
$$
0 \to \mathcal{I}^\bullet(X) \to
\mathcal{I}^\bullet(U) \oplus \mathcal{I}^\bullet(V) \to
\mathcal{I}^\bullet(U \cap V) \to 0
$$
est une suite exacte courte de complexes et la suite exacte longue
de cohomologie associ\'ee est celle de l'\'enonc\'e du lemme.
\end{proof}

\begin{lemma}[Mayer-Vietoris relatif]
\label{lemma-relative-mayer-vietoris}
Soit $f : X \to Y$ un morphisme de sch\'emas. Supposons que $X = U \cup V$
soit une r\'eunion de deux sous-sch\'emas ouverts. Notons
$a = f|_U : U \to Y$, $b = f|_V : V \to Y$, et
$c = f|_{U \cap V} : U \cap V \to Y$.
Pour tout faisceau ab\'elien $\mathcal{F}$ sur $X_\etale$
il existe une suite exacte longue
$$
0 \to
f_*\mathcal{F} \to
a_*(\mathcal{F}|_U) \oplus b_*(\mathcal{F}|_V) \to
c_*(\mathcal{F}|_{U \cap V}) \to
R^1f_*\mathcal{F} \to \ldots
$$
sur $Y_\etale$.
Cette suite exacte longue est fonctorielle en $\mathcal{F}$.
\end{lemma}

\begin{proof}
Soit $\mathcal{F} \to \mathcal{I}^\bullet$ une r\'esolution injective
de $\mathcal{F}$ sur $X_\etale$. Nous affirmons qu'il existe
une suite exacte courte de complexes
$$
0 \to
f_*\mathcal{I}^\bullet \to
a_*\mathcal{I}^\bullet|_U \oplus b_*\mathcal{I}^\bullet|_V \to
c_*\mathcal{I}^\bullet|_{U \cap V} \to
0.
$$
En effet, pour tout $W$ de $Y_\etale$ et tout $n \geq 0$, la
suite correspondante des groupes de sections sur $W$
$$
0 \to
\mathcal{I}^n(W \times_Y X) \to
\mathcal{I}^n(W \times_Y U)
\oplus \mathcal{I}^n(W \times_Y V) \to
\mathcal{I}^n(W \times_Y (U \cap V)) \to
0
$$
est une suite exacte courte, comme on l'a montr\'e dans la d\'emonstration du lemme \ref{lemma-mayer-vietoris}.
Le lemme r\'esulte alors du passage aux faisceaux de cohomologie et du fait que
$\mathcal{I}^\bullet|_U$ est une r\'esolution injective de $\mathcal{F}|_U$
et de m\^eme pour $\mathcal{I}^\bullet|_V$, $\mathcal{I}^\bullet|_{U \cap V}$.
\end{proof}







\section{Limites inductives}
\label{section-colimit}

\noindent
Rappelons que si $(\mathcal{F}_i, \varphi_{ii'})$ est un diagramme de faisceaux sur
un site $\mathcal{C}$, sa limite inductive (dans la cat\'egorie des faisceaux) est le
faisceau associ\'e au pr\'efaisceau $U \mapsto \colim_i \mathcal{F}_i(U)$. Voir
Sites, Lemme \ref{sites-lemma-colimit-sheaves}.
Si le syst\`eme est filtrant, si $U$ est un objet quasi-compact de
$\mathcal{C}$ qui poss\`ede un syst\`eme cofinal de recouvrements par des objets
quasi-compacts, alors $\mathcal{F}(U) = \colim \mathcal{F}_i(U)$, voir
Sites, Lemme \ref{sites-lemma-directed-colimits-sections}.
Voir Cohomologie sur les sites, Lemme
\ref{sites-cohomology-lemma-colim-works-over-collection}
pour un r\'esultat concernant les groupes de cohomologie sup\'erieure des limites inductives
de faisceaux ab\'eliens.

\medskip\noindent
Dans Cohomologie sur les sites, Lemme \ref{sites-cohomology-lemma-colimit},
nous g\'en\'eralisons ce r\'esultat \`a un syst\`eme de faisceaux
sur un syst\`eme projectif de sites. Voici la notion correspondante
dans le cas d'un syst\`eme de faisceaux \'etales d\'efinis
sur un syst\`eme projectif de sch\'emas.

\begin{definition}
\label{definition-inverse-system-sheaves}
Soit $I$ un ensemble pr\'eordonn\'e. Soit $(X_i, f_{i'i})$ un syst\`eme
projectif de sch\'emas index\'e par $I$.
Un {\it syst\`eme $(\mathcal{F}_i, \varphi_{i'i})$ de faisceaux
sur $(X_i, f_{i'i})$} est la donn\'ee de
\begin{enumerate}
\item un faisceau $\mathcal{F}_i$ sur $(X_i)_\etale$ pour tout $i \in I$,
\item pour $i' \geq i$, un morphisme
$\varphi_{i'i} : f_{i'i}^{-1}\mathcal{F}_i \to \mathcal{F}_{i'}$
de faisceaux sur $(X_{i'})_\etale$
\end{enumerate}
de telle sorte que $\varphi_{i''i} = \varphi_{i''i'} \circ f_{i'' i'}^{-1}\varphi_{i'i}$
lorsque $i'' \geq i' \geq i$.
\end{definition}

\noindent
Dans la situation de la d\'efinition \ref{definition-inverse-system-sheaves},
supposons que $I$ soit filtrant et que les morphismes de transition $f_{i'i}$ soient affines.
Soit $X = \lim X_i$ la limite dans la cat\'egorie des sch\'emas, voir
Limites, Section \ref{limits-section-limits}.
Notons $f_i : X \to X_i$ les morphismes de projection et consid\'erons les morphismes
$$
f_i^{-1}\mathcal{F}_i = f_{i'}^{-1}f_{i'i}^{-1}\mathcal{F}_i
\xrightarrow{f_{i'}^{-1}\varphi_{i'i}}
f_{i'}^{-1}\mathcal{F}_{i'}
$$
Cela fait des $f_i^{-1}\mathcal{F}_i$ un syst\`eme de faisceaux sur $X_\etale$
index\'e par $I$ (c'est un bon exercice de le v\'erifier).
Nous voulons souvent savoir si l'on a un isomorphisme
$$
H^q_\etale(X, \colim f_i^{-1}\mathcal{F}_i) =
\colim H^q_\etale(X_i, \mathcal{F}_i)
$$
Il s'av\'erera que cela est vrai si $X_i$ est quasi-compact et quasi-s\'epar\'e
pour tout $i$, voir le th\'eor\`eme \ref{theorem-colimit}.

\begin{lemma}
\label{lemma-colimit-affine-sites}
Soit $I$ un ensemble pr\'eordonn\'e filtrant. Soit $(X_i, f_{i'i})$ un syst\`eme
projectif de sch\'emas index\'e par $I$ dont les morphismes de transition sont affines.
Soit $X = \lim_{i \in I} X_i$. Avec les
notations de Topologies, Lemme \ref{topologies-lemma-alternative}, on a
$$
X_{affine, \etale} = \colim (X_i)_{affine, \etale}
$$
comme sites au sens de
Sites, Lemme \ref{sites-lemma-colimit-sites}.
\end{lemma}

\begin{proof}
D\'emontrons d'abord ceci lorsque $X$ et les $X_i$ sont quasi-compacts et
quasi-s\'epar\'es pour tout $i$ (c'est le cas dans toutes les situations qui
nous int\'eressent). Dans ce cas, tout objet de
$X_{affine, \etale}$, resp.\ de $(X_i)_{affine, \etale}$, est
de pr\'esentation finie sur $X$. De plus, la
cat\'egorie des sch\'emas de pr\'esentation finie sur $X$ est la
limite inductive des cat\'egories de sch\'emas de pr\'esentation finie
sur les $X_i$, voir
Limites, Lemme \ref{limits-lemma-descend-finite-presentation}.
Il en va de m\^eme des sous-cat\'egories d'objets affines et \'etales
sur $X$, d'apr\`es les lemmes de Limites
\ref{limits-lemma-limit-affine} et \ref{limits-lemma-descend-etale}.
Enfin, si $\{U^j \to U\}$ est un recouvrement de $X_{affine, \etale}$
et si $U_i^j \to U_i$ est un morphisme de sch\'emas affines \'etales sur
$X_i$ dont le changement de base \`a $X$ est $U^j \to U$, alors nous voyons que
le changement de base de $\{U^j_i \to U_i\}$ \`a un certain $X_{i'}$ est
un recouvrement pour $i'$ assez grand, voir
Limites, Lemme \ref{limits-lemma-descend-surjective}.

\medskip\noindent
Dans le cas g\'en\'eral, soit $U$ un objet de $X_{affine, \etale}$.
Alors $U \to X$ est \'etale et s\'epar\'e (puisque $U$ est s\'epar\'e),
mais n'est en g\'en\'eral pas quasi-compact. Toutefois, $U \to X$ est localement
de pr\'esentation finie et donc, d'apr\`es le
lemme \ref{limits-lemma-descend-finite-presentation-variant} de Limites,
il existe un $i$, un sch\'ema $U_i$ quasi-compact et quasi-s\'epar\'e, et
un morphisme $U_i \to X_i$ qui est localement de pr\'esentation finie
dont le changement de base \`a $X$ est $U \to X$. Alors $U = \lim_{i' \geq i} U_{i'}$
o\`u $U_{i'} = U_i \times_{X_i} X_{i'}$.
Quitte \`a augmenter $i$, nous pouvons supposer $U_i$ affine, voir
Limites, Lemme \ref{limits-lemma-limit-affine}.
Pour v\'erifier que $U_i \to X_i$ est \'etale pour $i$ assez grand,
choisissons un recouvrement ouvert affine fini $U_i = U_{i, 1} \cup \ldots \cup U_{i, m}$
tel que $U_{i, j} \to U_i \to X_i$ se factorise par un ouvert affine
$W_{i, j} \subset X_i$. On peut alors appliquer
Limites, Lemme \ref{limits-lemma-descend-etale},
pour voir que $U_{i, j} \to W_{i, j}$ est \'etale
apr\`es avoir, si n\'ecessaire, augment\'e $i$.
Nous voyons ainsi que le foncteur
$\colim (X_i)_{affine, \etale} \to X_{affine, \etale}$
est essentiellement surjectif. Sa pleine fid\'elit\'e r\'esulte
directement du r\'esultat d\'ej\`a utilis\'e :
Limites, Lemme \ref{limits-lemma-descend-finite-presentation-variant}.
L'assertion concernant les recouvrements se d\'emontre exactement de la m\^eme
mani\`ere que dans le premier paragraphe de la d\'emonstration.
\end{proof}

\noindent
Ce qui pr\'ec\`ede donne le r\'esultat g\'en\'eral suivant sur les limites inductives et la cohomologie.

\begin{theorem}
\label{theorem-colimit}
Soit $X = \lim_{i \in I} X_i$ la limite d'un syst\`eme projectif de sch\'emas
dont les morphismes de transition $f_{i'i} : X_{i'} \to X_i$ sont affines. Nous supposons
que $X_i$ est quasi-compact et quasi-s\'epar\'e pour tout $i \in I$.
Soit $(\mathcal{F}_i, \varphi_{i'i})$ un syst\`eme de faisceaux ab\'eliens
sur $(X_i, f_{i'i})$. Notons $f_i : X \to X_i$ la projection et posons
$\mathcal{F} = \colim f_i^{-1}\mathcal{F}_i$. Alors
$$
\colim_{i\in I} H_\etale^p(X_i, \mathcal{F}_i) = H_\etale^p(X, \mathcal{F}).
$$
pour tout $p \geq 0$.
\end{theorem}

\begin{proof}
D'apr\`es Topologies, lemme \ref{topologies-lemma-alternative},
nous pouvons calculer la cohomologie
de $\mathcal{F}$ sur $X_{affine, \etale}$.
Le r\'esultat d\'ecoule donc de la combinaison du
lemme \ref{lemma-colimit-affine-sites}
et du
lemme \ref{sites-cohomology-lemma-colimit} de Cohomologie sur les sites.
\end{proof}

\noindent
Les deux r\'esultats suivants sont des cas particuliers du th\'eor\`eme ci-dessus.

\begin{lemma}
\label{lemma-colimit}
Soit $X$ un sch\'ema quasi-compact et quasi-s\'epar\'e. Soit $I$
un ensemble ordonn\'e filtrant. Soit $(\mathcal{F}_i, \varphi_{ij})$ un syst\`eme
de faisceaux ab\'eliens sur $X_\etale$ index\'e par $I$. Alors
$$
\colim_{i\in I} H_\etale^p(X, \mathcal{F}_i) = H_\etale^p(X,
\colim_{i\in I} \mathcal{F}_i).
$$
\end{lemma}

\begin{proof}
C'est un cas particulier du th\'eor\`eme \ref{theorem-colimit}.
Nous esquissons \'egalement une d\'emonstration directe.
Nous le d\'emontrons simultan\'ement pour tout $X$, par r\'ecurrence sur $p$.
\begin{enumerate}
\item
Pour tout sch\'ema quasi-compact et quasi-s\'epar\'e $X$ et tout recouvrement \'etale
$\mathcal{U}$ de $X$, montrer qu'il existe un raffinement
$\mathcal{V} = \{V_j \to X\}_{j\in J}$ avec $J$ fini et chaque $V_j$
quasi-compact et quasi-s\'epar\'e, de telle sorte que tous les
$V_{j_0} \times_X \ldots \times_X V_{j_p}$ soient eux aussi
quasi-compacts et quasi-s\'epar\'es.
\item
En utilisant l'\'etape pr\'ec\'edente et la d\'efinition des limites inductives dans la cat\'egorie des
faisceaux, montrer que le th\'eor\`eme vaut pour $p = 0$ et tout $X$.
\item
En utilisant la localit\'e de la cohomologie
(lemme \ref{lemma-locality-cohomology}),
la suite spectrale de {\v C}ech vers la cohomologie
(th\'eor\`eme \ref{theorem-cech-ss}) et le fait que l'hypoth\`ese de r\'ecurrence
s'applique \`a tous les
$V_{j_0} \times_X \ldots \times_X V_{j_p}$
dans la situation ci-dessus, d\'emontrer l'\'etape de r\'ecurrence $p \to p + 1$.
\end{enumerate}
\end{proof}

\begin{lemma}
\label{lemma-directed-colimit-cohomology}
Soient $A$ un anneau, $(I, \leq)$ un ensemble ordonn\'e filtrant et $(B_i, \varphi_{ij})$ un
syst\`eme de $A$-alg\`ebres. Posons $B = \colim_{i\in I} B_i$. Soit $X \to \Spec(A)$
un morphisme de sch\'emas quasi-compact et quasi-s\'epar\'e. Soit
$\mathcal{F}$ un faisceau ab\'elien sur $X_\etale$.
Notons $Y_i = X \times_{\Spec(A)} \Spec(B_i)$,
$Y = X \times_{\Spec(A)} \Spec(B)$,
$\mathcal{G}_i = (Y_i \to X)^{-1}\mathcal{F}$ et
$\mathcal{G} = (Y \to X)^{-1}\mathcal{F}$. Alors
$$
H_\etale^p(Y, \mathcal{G}) =
\colim_{i\in I} H_\etale^p (Y_i, \mathcal{G}_i).
$$
\end{lemma}

\begin{proof}
C'est un cas particulier du th\'eor\`eme \ref{theorem-colimit}.
Nous indiquons \'egalement une d\'emonstration directe comme suit.
\begin{enumerate}
\item \'Etant donn\'e $V \to Y$ \'etale, avec $V$ quasi-compact et
quasi-s\'epar\'e, il existe $i\in I$ et $V_i \to Y_i$ tels que
$V = V_i \times_{Y_i} Y$.
Si tous les sch\'emas consid\'er\'es \'etaient affines, cela correspondrait \`a
l'\'enonc\'e alg\'ebrique suivant : si $B = \colim B_i$ et si $B \to C$ est \'etale,
alors il existe $i \in I$ et $B_i \to C_i$ \'etale tels que
$C \cong B \otimes_{B_i} C_i$.
Ceci est d\'emontr\'e dans Alg\`ebre, lemme \ref{algebra-lemma-etale}.
\item Dans la situation de (1), montrer que
$\mathcal{G}(V) = \colim_{i' \geq i} \mathcal{G}_{i'}(V_{i'})$
o\`u $V_{i'}$ est le changement de base de $V_i$ sur $Y_{i'}$.
\item D'apr\`es (1), on voit que pour tout recouvrement \'etale
$\mathcal{V} = \{V_j \to Y\}_{j\in J}$ avec $J$ fini et les
$V_j$ quasi-compacts et quasi-s\'epar\'es, il existe $i \in I$ et
un recouvrement \'etale $\mathcal{V}_i = \{V_{ij} \to Y_i\}_{j \in J}$
tel que $\mathcal{V} \cong \mathcal{V}_i \times_{Y_i} Y$.
\item Montrer que (2) et (3) impliquent
$$
\check H^*(\mathcal{V}, \mathcal{G})=
\colim_{i\in I} \check H^*(\mathcal{V}_i, \mathcal{G}_i).
$$
\item Utiliser judicieusement la suite spectrale de {\v C}ech vers la cohomologie
(th\'eor\`eme \ref{theorem-cech-ss}).
\end{enumerate}
\end{proof}

\begin{lemma}
\label{lemma-higher-direct-images}
Soit $f: X\to Y$ un morphisme de sch\'emas et soit $\mathcal{F}\in
\textit{Ab}(X_\etale)$. Alors $R^pf_*\mathcal{F}$ est le faisceau
associ\'e au pr\'efaisceau
$$
(V \to Y) \longmapsto H_\etale^p(X \times_Y V, \mathcal{F}|_{X \times_Y V}).
$$
Plus g\'en\'eralement, pour $K \in D(X_\etale)$, $R^pf_*K$ est le
faisceau associ\'e au pr\'efaisceau
$$
(V \to Y) \longmapsto H_\etale^p(X \times_Y V, K|_{X \times_Y V}).
$$
\end{lemma}

\begin{proof}
Ce lemme est valable pour les espaces topologiques, et la d\'emonstration dans ce cas est la
m\^eme. Voir Cohomologie sur les sites, lemme
\ref{sites-cohomology-lemma-higher-direct-images}
pour le cas d'un faisceau, et voir
Cohomologie sur les sites, lemme
\ref{sites-cohomology-lemma-sheafification-cohomology}
pour le cas d'un complexe de faisceaux ab\'eliens.
\end{proof}

\begin{lemma}
\label{lemma-relative-colimit}
Soit $S$ un sch\'ema. Soit $X = \lim_{i \in I} X_i$ une limite d'un
syst\`eme projectif de sch\'emas sur $S$ avec morphismes de transition affines
$f_{i'i} : X_{i'} \to X_i$. Nous supposons que les morphismes structuraux
$g_i : X_i \to S$ et $g : X \to S$ sont quasi-compacts et quasi-s\'epar\'es.
Soit $(\mathcal{F}_i, \varphi_{i'i})$ un syst\`eme de faisceaux ab\'eliens
sur $(X_i, f_{i'i})$. Notons $f_i : X \to X_i$ la projection et posons
$\mathcal{F} = \colim f_i^{-1}\mathcal{F}_i$. Alors
$$
\colim_{i\in I} R^p g_{i, *} \mathcal{F}_i = R^p g_* \mathcal{F}
$$
pour tout $p \geq 0$.
\end{lemma}

\begin{proof}
Rappelons (lemme \ref{lemma-higher-direct-images})
que $R^p g_{i, *} \mathcal{F}_i$ est le faisceau associ\'e au
pr\'efaisceau $U \mapsto H^p_\etale(U \times_S X_i, \mathcal{F}_i)$
et de m\^eme pour $R^pg_*\mathcal{F}$. De plus, la limite inductive d'un
syst\`eme de faisceaux est le faisceau associ\'e \`a la limite inductive calcul\'ee dans la cat\'egorie
des pr\'efaisceaux. Remarquons que tout objet de $S_\etale$ admet un recouvrement
par des objets quasi-compacts et quasi-s\'epar\'es (par exemple, des sch\'emas affines).
De plus, si $U$ est un objet quasi-compact et quasi-s\'epar\'e,
alors on a
$$
\colim H^p_\etale(U \times_S X_i, \mathcal{F}_i) =
H^p_\etale(U \times_S X, \mathcal{F})
$$
d'apr\`es le th\'eor\`eme \ref{theorem-colimit}. Le lemme en r\'esulte.
\end{proof}

\begin{lemma}
\label{lemma-relative-colimit-general}
Soit $I$ un ensemble ordonn\'e filtrant. Soit $g_i : X_i \to S_i$ un syst\`eme projectif de
morphismes de sch\'emas index\'e par $I$. Supposons que $g_i$ soit quasi-compact et
quasi-s\'epar\'e et que, pour $i' \geq i$, les morphismes de transition
$f_{i'i} : X_{i'} \to X_i$ et $h_{i'i} : S_{i'} \to S_i$ soient affines.
Soit $g : X \to S$ la limite des morphismes $g_i$, voir
Limites, section \ref{limits-section-limits}.
Notons $f_i : X \to X_i$ et $h_i : S \to S_i$ les projections.
Soit $(\mathcal{F}_i, \varphi_{i'i})$ un syst\`eme de faisceaux
sur $(X_i, f_{i'i})$. Posons $\mathcal{F} = \colim f_i^{-1}\mathcal{F}_i$. Alors
$$
R^p g_* \mathcal{F} =
\colim_{i \in I} h_i^{-1}R^p g_{i, *} \mathcal{F}_i
$$
pour tout $p \geq 0$.
\end{lemma}

\begin{proof}
Comment construit-on l'application du lemme ?
Pour $i' \geq i$, nous avons un diagramme commutatif
$$
\xymatrix{
X \ar[r]_{f_{i'}} \ar[d]_g &
X_{i'} \ar[r]_{f_{i'i}} \ar[d]_{g_{i'}} &
X_i \ar[d]^{g_i} \\
S \ar[r]^{h_{i'}} &
S_{i'} \ar[r]^{h_{i'i}} &
S_i
}
$$
Si nous composons l'application de changement de base
$h_{i'i}^{-1}Rg_{i, *}\mathcal{F}_i \to Rg_{i', *}f_{i'i}^{-1}\mathcal{F}_i$
(Cohomologie sur les sites, lemme
\ref{sites-cohomology-lemma-base-change-map-flat-case} ou
Remarque \ref{sites-cohomology-remark-base-change})
avec l'application $Rg_{i', *}\varphi_{i'i}$, nous obtenons alors
$\psi_{i'i} : h_{i' i}^{-1} R^p g_{i, *} \mathcal{F}_i \to
R^pg_{i', *} \mathcal{F}_{i'}$. De m\^eme, en utilisant le carr\'e de gauche
du diagramme, nous obtenons les applications
$\psi_i : h_i^{-1}R^pg_{i, *}\mathcal{F}_i \to R^pg_*\mathcal{F}$.
Les applications $h_{i'}^{-1}\psi_{i'i}$ et $\psi_i$ sont celles utilis\'ees dans
l'\'enonc\'e du lemme. Pour que cela ait un sens, nous devons v\'erifier que
$\psi_{i''i} = \psi_{i''i'} \circ h_{i''i'}^{-1}\psi_{i'i}$ et
$\psi_{i'} \circ h_{i'}^{-1}\psi_{i'i} = \psi_i$ ; cela r\'esulte
de Cohomologie sur les sites, remarque
\ref{sites-cohomology-remark-compose-base-change-horizontal}.

\medskip\noindent
D\'emonstration de l'\'egalit\'e. Premi\`ere d\'emonstration par
d\'ecalage de dimension\footnote{On peut aussi utiliser cette m\'ethode
pour construire les applications du lemme.}. Pour tout
$U$ affine et \'etale sur $X$, d'apr\`es le Th\'eor\`eme \ref{theorem-colimit},
nous avons
$$
g_*\mathcal{F}(U) =
H^0(U \times_S X, \mathcal{F}) =
\colim H^0(U_i \times_{S_i} X_i, \mathcal{F}_i) =
\colim g_{i, *}\mathcal{F}_i(U_i)
$$
o\`u la limite inductive porte sur les $i$ assez grands tels qu'
il existe un $i$ et un $U_i$ affine et \'etale sur $S_i$
dont le changement de base est $U$ sur $S$ (voir
Lemme \ref{lemma-colimit-affine-sites}).
Le membre de droite est \'egal \`a
$(\colim h_i^{-1}g_{i, *}\mathcal{F}_i)(U)$ d'apr\`es
Sites, Lemme \ref{sites-lemma-colimit}.
Cela d\'emontre le lemme pour $p = 0$.
Si $(\mathcal{G}_i, \varphi_{i'i})$ est un syst\`eme tel que
$\mathcal{G} = \colim f_i^{-1}\mathcal{G}_i$
et tel que $\mathcal{G}_i$ soit un faisceau ab\'elien injectif sur $X_i$
pour tout $i$, alors, pour tout $U$ affine et \'etale sur $X$, d'apr\`es
le Th\'eor\`eme \ref{theorem-colimit}, nous avons
$$
H^p(U \times_S X, \mathcal{G}) =
\colim H^p(U_i \times_{S_i} X_i, \mathcal{G}_i) = 0
$$
pour $p > 0$ (la m\^eme limite inductive que pr\'ec\'edemment). Ainsi $R^pg_*\mathcal{G} = 0$
et nous obtenons le r\'esultat pour $p > 0$ pour un tel syst\`eme.
En g\'en\'eral, nous pouvons choisir une suite exacte courte de syst\`emes
$$
0 \to (\mathcal{F}_i, \varphi_{i'i}) \to
(\mathcal{G}_i, \varphi_{i'i}) \to
(\mathcal{Q}_i, \varphi_{i'i}) \to 0
$$
o\`u $(\mathcal{G}_i, \varphi_{i'i})$ est comme ci-dessus, voir
Cohomologie sur les sites, Lemme
\ref{sites-cohomology-lemma-colim-sites-injective}.
Par r\'ecurrence, le lemme vaut pour $p - 1$ et, d'apr\`es ce qui pr\'ec\`ede, on a
l'annulation en degr\'e $p$ pour $(\mathcal{G}_i, \varphi_{i'i})$.
On en d\'eduit le r\'esultat pour $p$
et $(\mathcal{F}_i, \varphi_{i'i})$ par la suite exacte longue
de cohomologie.

\medskip\noindent
Seconde d\'emonstration.
Rappelons que $S_{affine, \etale} = \colim (S_i)_{affine, \etale}$, voir
le lemme \ref{lemma-colimit-affine-sites}. Ainsi, si $U$ est un objet de
$S_{affine, \etale}$, alors on peut \'ecrire $U = U_i \times_{S_i} S$
pour un certain $i$ et un certain $U_i$ de $(S_i)_{affine, \etale}$, et
$$
(\colim_{i \in I} h_i^{-1}R^p g_{i, *} \mathcal{F}_i)(U) =
\colim_{i' \geq i} (R^p g_{i', *}\mathcal{F}_{i'})(U_i \times_{S_i} S_{i'})
$$
d'apr\`es Sites, Lemme \ref{sites-lemma-colimit}, et la construction des
morphismes de transition du syst\`eme d\'ecrit ci-dessus. Puisque
$R^pg_{i', *}\mathcal{F}_{i'}$ est le faisceau associ\'e au pr\'efaisceau
$U_{i'} \mapsto H^p(U_{i'} \times_{S_{i'}} X_{i'}, \mathcal{F}_{i'})$
et puisque $R^pg_*\mathcal{F}$ est le faisceau associ\'e au pr\'efaisceau
$U \mapsto H^p(U \times_S X, \mathcal{F})$
(Lemme \ref{lemma-higher-direct-images})
on obtient un diagramme commutatif canonique
$$
\xymatrix{
\colim_{i' \geq i}
H^p(U_i \times_{S_i} X_{i'}, \mathcal{F}_{i'}) \ar[r] \ar[d] &
\colim_{i' \geq i}
(R^p g_{i', *}\mathcal{F}_{i'})(U_i \times_{S_i} S_{i'}) \ar[d] \\
H^p(U \times_S X, \mathcal{F}) \ar[r] &
R^pg_*\mathcal{F}(U)
}
$$
Remarquons que la fl\`eche verticale de gauche est un isomorphisme
d'apr\`es le th\'eor\`eme \ref{theorem-colimit}. Nous cherchons \`a montrer que
la fl\`eche verticale de droite est un isomorphisme. Cependant, nous
savons d\'ej\`a que la source et le but de cette fl\`eche
sont des faisceaux sur $S_{affine, \etale}$. Il suffit donc de
montrer : (1) tout \'el\'ement du but provient localement d'un
\'el\'ement de la source, et (2) tout \'el\'ement de la source
qui s'envoie sur z\'ero dans le but s'annule localement.
L'assertion (1) d\'ecoule imm\'ediatement de ce qui pr\'ec\`ede et du fait que
la fl\`eche horizontale inf\'erieure provient d'un morphisme de pr\'efaisceaux
qui devient un isomorphisme apr\`es faisceautisation.
Pour (2), soit $\xi \in  \colim_{i' \geq i}
(R^p g_{i', *}\mathcal{F}_{i'})(U_i \times_{S_i} S_{i'})$
dans le noyau. Choisissons $i' \geq i$ et
$\xi_{i'} \in (R^p g_{i', *}\mathcal{F}_{i'})(U_i \times_{S_i} S_{i'})$
repr\'esentant $\xi$.
Choisissons un recouvrement \'etale standard
$\{U_{i', k} \to U_i \times_{S_i} S_{i'}\}_{k = 1, \ldots, m}$
tel que $\xi_{i'}|_{U_{i', k}}$ provienne de
$\xi_{i', k} \in H^p(U_{i', k} \times_{S_{i'}} X_{i'}, \mathcal{F}_{i'})$.
Puisqu'il suffit de prouver que $\xi$ s'annule localement, nous
pouvons remplacer $U$ par les membres du recouvrement \'etale
$\{U_{i', k} \times_{S_{i'}} S \to U = U_i \times_{S_i} S\}$.
Apr\`es ce remplacement, on voit que $\xi$ est l'image
d'un \'el\'ement $\xi'$ du groupe
$\colim_{i' \geq i} H^p(U_i \times_{S_i} X_{i'}, \mathcal{F}_{i'})$
dans le diagramme ci-dessus. Puisque $\xi'$ s'envoie sur z\'ero dans $R^pg_*\mathcal{F}(U)$
nous pouvons effectuer un autre remplacement et supposer que $\xi'$ s'envoie
sur z\'ero dans $H^p(U \times_S X, \mathcal{F})$.
Toutefois, puisque la fl\`eche verticale de gauche est un isomorphisme,
on conclut alors que $\xi' = 0$, d'o\`u $\xi = 0$, comme voulu.
\end{proof}

\begin{lemma}
\label{lemma-linus-hamann}
Soit $X = \lim_{i \in I} X_i$ la limite d'un syst\`eme projectif de sch\'emas
dont les morphismes de transition $f_{i'i}$ sont affines et dont les morphismes de projection
sont $f_i : X \to X_i$. Soit $\mathcal{F}$ un faisceau sur $X_\etale$. Alors
\begin{enumerate}
\item il existe des morphismes canoniques
$\varphi_{i'i} : f_{i'i}^{-1}f_{i, *}\mathcal{F} \to f_{i', *}\mathcal{F}$
tels que $(f_{i, *}\mathcal{F}, \varphi_{i'i})$ soit un syst\`eme de
faisceaux sur $(X_i, f_{i'i})$ au sens de la
d\'efinition \ref{definition-inverse-system-sheaves}, et
\item $\mathcal{F} = \colim f_i^{-1}f_{i, *}\mathcal{F}$.
\end{enumerate}
\end{lemma}

\begin{proof}
Par Topologies, Lemme \ref{topologies-lemma-alternative}, et le
lemme \ref{lemma-colimit-affine-sites}, ceci est un cas particulier de
Sites, Lemme \ref{sites-lemma-colimit-push-pull}.
\end{proof}

\begin{lemma}
\label{lemma-compute-strangely}
Soit $I$ un ensemble ordonn\'e filtrant. Soit $g_i : X_i \to S_i$ un syst\`eme projectif de
morphismes de sch\'emas index\'e par $I$. Supposons que $g_i$ soit quasi-compact et
quasi-s\'epar\'e et que, pour $i' \geq i$, les morphismes de transition
$X_{i'} \to X_i$ et $S_{i'} \to S_i$ soient affines.
Soit $g : X \to S$ la limite des morphismes $g_i$, voir
Limites, Section \ref{limits-section-limits}.
Notons $f_i : X \to X_i$ et $h_i : S \to S_i$ les projections.
Soit $\mathcal{F}$ un faisceau ab\'elien sur $X$. Alors on a
$$
R^pg_*\mathcal{F} = \colim_{i \in I} h_i^{-1}R^pg_{i, *}(f_{i, *}\mathcal{F})
$$
\end{lemma}

\begin{proof}
Combinaison formelle des lemmes \ref{lemma-relative-colimit-general}
et \ref{lemma-linus-hamann}.
\end{proof}











\section{Limites inductives et complexes}
\label{section-colimit-complexes}

\noindent
Dans cette section, nous \'etudions la cohomologie des syst\`emes de complexes
dans divers contextes, en poursuivant l'\'etude des faisceaux commenc\'ee
\`a la section \ref{section-colimit}.
Nous conseillons vivement au lecteur de ne pas lire cette section, sauf si cela est absolument
n\'ecessaire.

\begin{lemma}
\label{lemma-colimit-variant-complexes}
Soit $X = \lim_{i \in I} X_i$ la limite d'un syst\`eme projectif de sch\'emas
dont les morphismes de transition $f_{i'i} : X_{i'} \to X_i$ sont affines. Supposons
que $X_i$ soit quasi-compact et quasi-s\'epar\'e pour tout $i \in I$.
Soit $\mathcal{F}_i^\bullet$ un complexe de faisceaux ab\'eliens sur
$X_{i, \etale}$. Soit $\varphi_{i'i} : f_{i'i}^{-1}\mathcal{F}_i^\bullet \to
\mathcal{F}_{i'}^\bullet$ un morphisme de complexes sur $X_{i, \etale}$
tel que $\varphi_{i''i} = \varphi_{i''i'} \circ f_{i'' i'}^{-1}\varphi_{i'i}$
lorsque $i'' \geq i' \geq i$. Supposons qu'il existe un entier $a$ tel que
$\mathcal{F}_i^n = 0$ pour $n < a$ et tout $i \in I$.
Alors on a
$$
H^p_\etale(X, \colim f_i^{-1}\mathcal{F}_i^\bullet) =
\colim H^p_\etale(X_i, \mathcal{F}^\bullet_i)
$$
o\`u $f_i : X \to X_i$ est la projection.
\end{lemma}

\begin{proof}
Cela r\'esulte du th\'eor\`eme \ref{theorem-colimit}. Posons
$\mathcal{F}^\bullet = \colim f_i^{-1}\mathcal{F}_i^\bullet$.
Le th\'eor\`eme affirme que
$$
\colim_{i \in I} H_\etale^p(X_i, \mathcal{F}_i^n) =
H_\etale^p(X, \mathcal{F}^n)
$$
pour tous $n, p \in \mathbf{Z}$. Utilisons les suites spectrales
$$
E_{1, i}^{s, t} = H_\etale^t(X_i, \mathcal{F}_i^s) \Rightarrow
H_\etale^{s + t}(X_i, \mathcal{F}_i^\bullet)
$$
et
$$
E_1^{s, t} = H_\etale^t(X, \mathcal{F}^s) \Rightarrow
H_\etale^{s + t}(X, \mathcal{F}^\bullet)
$$
de Cat\'egories d\'eriv\'ees, Lemme \ref{derived-lemma-two-ss-complex-functor}.
Puisque $\mathcal{F}_i^n = 0$ pour $n < a$ (avec $a$ ind\'ependant de $i$),
on voit que seul un nombre fini de termes $E_{1, i}^{s, t}$
(ind\'ependant de $i$) et $E_1^{s, t}$ contribuent \`a
$H^q_\etale(X_i, \mathcal{F}_i^\bullet)$ et
$H^q_\etale(X, \mathcal{F}^\bullet)$ et $E_1^{s, t} = \colim E_{i, i}^{s, t}$.
Cela d\'emontre l'assertion. Certains d\'etails sont omis.
(Il existe un autre argument utilisant les troncatures ``b\^etes''
des complexes, qui \'evite d'utiliser les suites spectrales.)
\end{proof}

\begin{lemma}
\label{lemma-direct-sum-bounded-below-cohomology}
Soit $X$ un sch\'ema quasi-compact et quasi-s\'epar\'e. Soit
$K_i \in D(X_\etale)$, $i \in I$, une famille d'objets.
Supposons donn\'e $a \in \mathbf{Z}$ tel que $H^n(K_i) = 0$ pour $n < a$
et $i \in I$. Alors $R\Gamma(X, \bigoplus_i K_i) =
\bigoplus_i R\Gamma(X, K_i)$.
\end{lemma}

\begin{proof}
Il faut montrer que $H^p(X, \bigoplus_i K_i) =
\bigoplus_i H^p(X, K_i)$ pour tout $p \in \mathbf{Z}$.
Choisissons des complexes $\mathcal{F}_i^\bullet$ repr\'esentant $K_i$
tels que $\mathcal{F}_i^n = 0$ pour $n < a$.
La somme directe des complexes $\mathcal{F}_i^\bullet$
repr\'esente l'objet $\bigoplus K_i$ d'apr\`es
Injectifs, Lemme \ref{injectives-lemma-derived-products}.
Puisque $\bigoplus \mathcal{F}^\bullet$ est la limite inductive filtrante
des sommes directes finies, le r\'esultat d\'ecoule
du lemme \ref{lemma-colimit-variant-complexes}.
\end{proof}

\begin{lemma}
\label{lemma-colimit-complexes}
Soit $S$ un sch\'ema. Soit $X = \lim_{i \in I} X_i$ la limite d'un
syst\`eme projectif de sch\'emas sur $S$ avec pour morphismes de transition affines
$f_{i'i} : X_{i'} \to X_i$. Supposons que $X_i$ soit quasi-compact et
quasi-s\'epar\'e pour tout $i \in I$. Soit $K \in D^+(S_\etale)$. Alors
$$
\colim_{i \in I} H_\etale^p(X_i, K|_{X_i}) = H_\etale^p(X, K|_X).
$$
pour tout $p \in \mathbf{Z}$, o\`u $K|_{X_i}$ et $K|_X$
sont les images inverses de $K$ sur $X_i$ et $X$.
\end{lemma}

\begin{proof}
On peut repr\'esenter $K$ par un complexe born\'e inf\'erieurement $\mathcal{G}^\bullet$
de faisceaux ab\'eliens sur $S_\etale$. Disons que $\mathcal{G}^n = 0$ pour $n < a$.
Notons $\mathcal{F}^\bullet_i$ et $\mathcal{F}^\bullet$ les images inverses
de ce complexe sur $X_i$ et $X$. Ces complexes repr\'esentent les
objets $K|_{X_i}$ et $K|_X$, et l'on a
$\mathcal{F}^\bullet = \colim f_i^{-1}\mathcal{F}_i^\bullet$
terme \`a terme. Le lemme d\'ecoule donc du
lemme \ref{lemma-colimit-variant-complexes}.
\end{proof}

\begin{lemma}
\label{lemma-relative-colimit-general-complexes}
Soient $I$, $g_i : X_i \to S_i$, $g : X \to S$, $f_i$, $g_i$, $h_i$ comme dans le
lemme \ref{lemma-relative-colimit-general}.
Soient $0 \in I$ et $K_0 \in D^+(X_{0, \etale})$.
Pour $i \geq 0$, notons $K_i$ l'image inverse de $K_0$ sur $X_i$.
Notons $K$ l'image inverse de $K_0$ sur $X$. Alors
$$
R^pg_*K = \colim_{i \geq 0} h_i^{-1}R^pg_{i, *}K_i
$$
pour tout $p \in \mathbf{Z}$.
\end{lemma}

\begin{proof}
Fixons un entier $p_0 \in \mathbf{Z}$.
Soit $a$ un entier tel que $H^j(K_0) = 0$ pour $j < a$.
Nous allons d\'emontrer que la formule vaut pour tout $p \leq p_0$
par r\'ecurrence descendante sur $a$. Si
$a > p_0$, alors on voit que les membres de gauche et de droite
de la formule sont nuls pour $p \leq p_0$ par annulation triviale, voir
Cat\'egories d\'eriv\'ees, Lemme \ref{derived-lemma-negative-vanishing}.
Supposons $a \leq p_0$. Consid\'erons le triangle distingu\'e
$$
H^a(K_0)[-a] \to K_0 \to \tau_{\geq a + 1}K_0
$$
Le tir\'e en arri\`ere de ce triangle distingu\'e sur $X_i$ et $X$
donne des triangles distingu\'es compatibles pour $K_i$ et $K$.
Pour $p \leq p_0$, consid\'erons le diagramme commutatif
$$
\xymatrix{
\colim_{i \geq 0} h_i^{-1}R^{p - 1}g_{i, *}(\tau_{\geq a + 1}K_i)
\ar[r]_-\alpha \ar[d] &
R^{p - 1}g_*(\tau_{\geq a + 1}K) \ar[d] \\
\colim_{i \geq 0} h_i^{-1}R^pg_{i, *}(H^a(K_i)[-a])
\ar[r]_-\beta \ar[d] &
R^pg_*(H^a(K)[-a]) \ar[d] \\
\colim_{i \geq 0} h_i^{-1}R^pg_{i, *}K_i
\ar[r]_-\gamma \ar[d] &
R^pg_*K \ar[d] \\
\colim_{i \geq 0} R^pg_{i, *}\tau_{\geq a + 1}K_i
\ar[r]_-\delta \ar[d] &
R^pg_*\tau_{\geq a + 1}K \ar[d] \\
\colim_{i \geq 0} R^{p + 1}g_{i, *}(H^a(K_i)[-a])
\ar[r]^-\epsilon &
R^{p + 1}g_*(H^a(K)[-a])
}
$$
\`a colonnes exactes. Les fl\`eches $\beta$ et $\epsilon$ sont des isomorphismes d'apr\`es le
lemme \ref{lemma-relative-colimit-general}.
Les fl\`eches $\alpha$ et $\delta$ sont des isomorphismes
par l'hypoth\`ese de r\'ecurrence.
Ainsi, $\gamma$ est un isomorphisme, comme voulu.
\end{proof}

\begin{lemma}
\label{lemma-relative-colimit-general-really-complexes}
Soient $I$, $g_i : X_i \to S_i$, $g : X \to S$, $f_{ii'}$, $f_i$, $g_i$, $h_i$
comme dans le lemme \ref{lemma-relative-colimit-general}.
Soit $\mathcal{F}_i^\bullet$ un complexe de faisceaux ab\'eliens sur
$X_{i, \etale}$. Soit $\varphi_{i'i} : f_{i'i}^{-1}\mathcal{F}_i^\bullet \to
\mathcal{F}_{i'}^\bullet$ un morphisme de complexes sur $X_{i, \etale}$
tel que $\varphi_{i''i} = \varphi_{i''i'} \circ f_{i'' i'}^{-1}\varphi_{i'i}$
lorsque $i'' \geq i' \geq i$. Supposons qu'il existe un entier $a$ tel que
$\mathcal{F}_i^n = 0$ pour $n < a$ et tout $i \in I$. Alors
$$
R^pg_*(\colim f_i^{-1}\mathcal{F}_i^\bullet) =
\colim_{i \geq 0} h_i^{-1}R^pg_{i, *}\mathcal{F}_i^\bullet
$$
pour tout $p \in \mathbf{Z}$.
\end{lemma}

\begin{proof}
Ceci r\'esulte du lemme \ref{lemma-relative-colimit-general}. Posons
$\mathcal{F}^\bullet = \colim f_i^{-1}\mathcal{F}_i^\bullet$.
Le lemme affirme que
$$
\colim_{i \in I} h_i^{-1}R^pg_{i, *}\mathcal{F}_i^n = R^pg_*\mathcal{F}^n
$$
pour tous $n, p \in \mathbf{Z}$. Utilisons les suites spectrales
$$
E_{1, i}^{s, t} = R^tg_{i, *}\mathcal{F}_i^s \Rightarrow
R^{s + t}g_{i, *}\mathcal{F}_i^\bullet
$$
et
$$
E_1^{s, t} = R^tg_*\mathcal{F}^s \Rightarrow
R^{s + t}g_*\mathcal{F}^\bullet
$$
de Cat\'egories d\'eriv\'ees, Lemme \ref{derived-lemma-two-ss-complex-functor}.
Puisque $\mathcal{F}_i^n = 0$ pour $n < a$ (avec $a$ ind\'ependant de $i$),
on voit que seul un nombre fini fix\'e de termes $E_{1, i}^{s, t}$
(ind\'ependant de $i$) et $E_1^{s, t}$ contribuent et que
$E_1^{s, t} = \colim E_{i, i}^{s, t}$.
Cela d\'emontre l'assertion. Certains d\'etails sont omis.
(Il existe un autre argument utilisant les troncatures ``b\^etes''
des complexes, qui \'evite d'utiliser les suites spectrales.)
\end{proof}

\begin{lemma}
\label{lemma-direct-sum-bounded-below-pushforward}
Soit $f : X \to Y$ un morphisme de sch\'emas quasi-compact et quasi-s\'epar\'e.
Soit $K_i \in D(X_\etale)$, $i \in I$, une famille d'objets.
Supposons donn\'e $a \in \mathbf{Z}$ tel que $H^n(K_i) = 0$ pour $n < a$
et $i \in I$. Alors $Rf_*(\bigoplus_i K_i) = \bigoplus_i Rf_*K_i$.
\end{lemma}

\begin{proof}
Il faut montrer que $R^pf_*(\bigoplus_i K_i) =
\bigoplus_i R^pf_*K_i$ pour tout $p \in \mathbf{Z}$.
Choisissons des complexes $\mathcal{F}_i^\bullet$ repr\'esentant $K_i$
tels que $\mathcal{F}_i^n = 0$ pour $n < a$.
La somme directe des complexes $\mathcal{F}_i^\bullet$
repr\'esente l'objet $\bigoplus K_i$ d'apr\`es
Injectifs, Lemme \ref{injectives-lemma-derived-products}.
Puisque $\bigoplus \mathcal{F}^\bullet$ est la limite inductive filtrante
des sommes directes finies, le r\'esultat d\'ecoule
du lemme \ref{lemma-relative-colimit-general-really-complexes}.
\end{proof}














\section{Fibres des images directes sup\'erieures}
\label{section-stalks-direct-image}

\noindent
Les fibres des images directes sup\'erieures se calculent souvent comme suit.

\begin{theorem}
\label{theorem-higher-direct-images}
Soit $f: X \to S$ un morphisme de sch\'emas quasi-compact et quasi-s\'epar\'e,
$\mathcal{F}$ un faisceau ab\'elien sur $X_\etale$, et $\overline{s}$ un
point g\'eom\'etrique de $S$ au-dessus de $s \in S$. Alors
$$
\left(R^nf_* \mathcal{F}\right)_{\overline{s}} =
H_\etale^n( X \times_S \Spec(\mathcal{O}_{S, \overline{s}}^{sh}),
p^{-1}\mathcal{F})
$$
o\`u $p : X \times_S \Spec(\mathcal{O}_{S, \overline{s}}^{sh}) \to X$
est la projection. Pour $K \in D^+(X_\etale)$ et $n \in \mathbf{Z}$,
on a
$$
\left(R^nf_*K\right)_{\overline{s}} =
H_\etale^n(X \times_S \Spec(\mathcal{O}_{S, \overline{s}}^{sh}), p^{-1}K)
$$
En fait, on a
$$
\left(Rf_*K\right)_{\overline{s}}
=
R\Gamma_\etale(X \times_S \Spec(\mathcal{O}_{S, \overline{s}}^{sh}), p^{-1}K)
$$
dans $D^+(\textit{Ab})$.
\end{theorem}

\begin{proof}
Soit $\mathcal{I}$ la cat\'egorie des voisinages \'etales de $\overline{s}$
sur $S$. D'apr\`es le lemme \ref{lemma-higher-direct-images},
on a
$$
(R^nf_*\mathcal{F})_{\overline{s}} =
\colim_{(V, \overline{v}) \in \mathcal{I}^{opp}}
H_\etale^n(X \times_S V, \mathcal{F}|_{X \times_S V}).
$$
Nous pouvons remplacer $\mathcal{I}$ par la sous-cat\'egorie initiale form\'ee
des voisinages \'etales affines de $\overline{s}$. Remarquons que
$$
\Spec(\mathcal{O}_{S, \overline{s}}^{sh}) =
\lim_{(V, \overline{v}) \in \mathcal{I}} V
$$
d'apr\`es le lemme \ref{lemma-describe-etale-local-ring} et
Limites, Lemme
\ref{limits-lemma-directed-inverse-system-affine-schemes-has-limit}.
Puisque les produits fibr\'es commutent aux limites, on obtient aussi
$$
X \times_S \Spec(\mathcal{O}_{S, \overline{s}}^{sh}) =
\lim_{(V, \overline{v}) \in \mathcal{I}} X \times_S V
$$
On conclut par le lemme \ref{lemma-directed-colimit-cohomology}.
Pour la seconde variante, on emploie le m\^eme argument avec
le lemme \ref{lemma-colimit-complexes} au lieu du
lemme \ref{lemma-directed-colimit-cohomology}.

\medskip\noindent
Pour voir que la derni\`ere assertion est vraie, il suffit de construire un morphisme
$\left(Rf_*K\right)_{\overline{s}} \to
R\Gamma_\etale(X \times_S \Spec(\mathcal{O}_{S, \overline{s}}^{sh}), p^{-1}K)$
dans $D^+(\textit{Ab})$ qui r\'ealise les isomorphismes ci-dessus sur les groupes de cohomologie
en degr\'e $n$ pour tout $n$. Pour cela, choisissons un
complexe born\'e inf\'erieurement $\mathcal{J}^\bullet$ de faisceaux ab\'eliens injectifs
sur $X_\etale$ repr\'esentant $K$. Le complexe
$f_*\mathcal{J}^\bullet$ repr\'esente $Rf_*K$. Ainsi, le complexe
$$
(f_*\mathcal{J}^\bullet)_{\overline{s}} =
\colim_{(V, \overline{v}) \in \mathcal{I}^{opp}}
(f_*\mathcal{J}^\bullet)(V)
$$
repr\'esente $(Rf_*K)_{\overline{s}}$. Pour chaque $V$, on a des morphismes
$$
(f_*\mathcal{J}^\bullet)(V) =
\Gamma(X \times_S V, \mathcal{J}^\bullet)
\longrightarrow
\Gamma(X \times_S \Spec(\mathcal{O}_{S, \overline{s}}^{sh}),
p^{-1}\mathcal{J}^\bullet)
$$
et le complexe cible repr\'esente
$R\Gamma_\etale(X \times_S \Spec(\mathcal{O}_{S, \overline{s}}^{sh}), p^{-1}K)$
dans $D^+(\textit{Ab})$.
En prenant la limite inductive de ces morphismes, on obtient le r\'esultat.
\end{proof}

\begin{remark}
\label{remark-stalk-fibre}
Soit $f : X \to S$ un morphisme de sch\'emas.
Soit $K \in D(X_\etale)$.
Soit $\overline{s}$ un point g\'eom\'etrique de $S$.
Il existe toujours des morphismes canoniques
$$
(Rf_*K)_{\overline{s}}
\longrightarrow
R\Gamma(X \times_S \Spec(\mathcal{O}_{S, \overline{s}}^{sh}), p^{-1}K)
\longrightarrow
R\Gamma(X_{\overline{s}}, K|_{X_{\overline{s}}})
$$
o\`u $p : X \times_S \Spec(\mathcal{O}_{S, \overline{s}}^{sh}) \to X$
est la projection. Consid\'erons en effet le diagramme commutatif
$$
\xymatrix{
X_{\overline{s}} \ar[r] \ar[d]^{f_{\overline{s}}} &
X \times_S \Spec(\mathcal{O}_{S, \overline{s}}^{sh}) \ar[r]_-p \ar[d]^{f'} &
X \ar[d]^f \\
\overline{s} \ar[r]^-i &
\Spec(\mathcal{O}_{S, \overline{s}}^{sh}) \ar[r]^-j &
S
}
$$
On dispose des morphismes de changement de base
$$
i^{-1}Rf'_*(p^{-1}K) \to Rf_{\overline{s}, *}(K|_{X_{\overline{s}}})
\quad\text{and}\quad
j^{-1}Rf_*K \to Rf'_*(p^{-1}K)
$$
(Cohomologie sur les sites, Remarque \ref{sites-cohomology-remark-base-change})
associ\'es aux deux carr\'es de ce diagramme.
En prenant les sections globales, on obtient les morphismes voulus.
D'apr\`es Cohomologie sur les sites, Remarque
\ref{sites-cohomology-remark-compose-base-change-horizontal}
la compos\'ee de ces deux morphismes est le
morphisme usuel de changement de base
$(Rf_*K)_{\overline{s}} \to R\Gamma(X_{\overline{s}}, K|_{X_{\overline{s}}})$.
\end{remark}






\section{La suite spectrale de Leray}
\label{section-leray}

\begin{lemma}
\label{lemma-prepare-leray}
Soient $f: X \to Y$ un morphisme et $\mathcal{I}$ un objet injectif de
$\textit{Ab}(X_\etale)$. Soit $V \in \Ob(Y_\etale)$. Alors
\begin{enumerate}
\item pour tout recouvrement $\mathcal{V} = \{V_j\to V\}_{j \in J}$, on a
$\check H^p(\mathcal{V}, f_*\mathcal{I}) = 0$ pour tout $p > 0$ ;
\item $f_*\mathcal{I}$ est acyclique pour le foncteur $\Gamma(V, -)$ ; et
\item si $g : Y \to Z$, alors $f_*\mathcal{I}$ est acyclique pour $g_*$. 
\end{enumerate}
\end{lemma}

\begin{proof}
Remarquons que $\check{\mathcal{C}}^\bullet(\mathcal{V}, f_*\mathcal{I}) =
\check{\mathcal{C}}^\bullet(\mathcal{V} \times_Y X, \mathcal{I})$
dont les groupes de cohomologie sup\'erieure sont nuls d'apr\`es le lemme \ref{lemma-hom-injective}.
Cela d\'emontre (1). La deuxi\`eme assertion d\'ecoule du fait qu'un faisceau dont
la cohomologie sup\'erieure de {\v C}ech s'annule pour tout recouvrement a des groupes de
cohomologie sup\'erieure nuls. C'est un excellent exercice d'utilisation de la
suite spectrale de {\v C}ech vers la cohomologie, mais voir
Cohomologie sur les sites, Lemme \ref{sites-cohomology-lemma-cech-vanish-collection}
pour les d\'etails et pour un \'enonc\'e plus pr\'ecis et plus g\'en\'eral.
L'assertion (3) r\'esulte de (2) et de la description de
$R^pg_*$ dans le lemme \ref{lemma-higher-direct-images}.
\end{proof}

\noindent
Le formalisme des suites spectrales de Grothendieck donne alors le
r\'esultat suivant.

\begin{proposition}[Suite spectrale de Leray]
\label{proposition-leray}
Soit $f: X \to Y$ un morphisme de sch\'emas et soit $\mathcal{F}$ un faisceau \'etale sur
$X$. Il existe alors une suite spectrale
$$
E_2^{p, q} = H_\etale^p(Y, R^qf_*\mathcal{F}) \Rightarrow
H_\etale^{p+q}(X, \mathcal{F}).
$$
\end{proposition}

\begin{proof}
Voir le lemme \ref{lemma-prepare-leray} et
Cat\'egories d\'eriv\'ees, Section
\ref{derived-section-composition-right-derived-functors}.
\end{proof}









\section{Annulation des images directes sup\'erieures finies}
\label{section-vanishing-finite-morphism}

\noindent
Le but suivant est de d\'emontrer que les images directes sup\'erieures d'un morphisme fini de
sch\'emas s'annulent.

\begin{lemma}
\label{lemma-vanishing-etale-cohomology-strictly-henselian}
Soit $R$ un anneau local strictement hens\'elien. Posons $S = \Spec(R)$ et soit
$\overline{s}$ son point ferm\'e. Alors le foncteur des
sections globales
$\Gamma(S, -) : \textit{Ab}(S_\etale) \to \textit{Ab}$
est exact. En fait, $\Gamma(S, \mathcal{F}) = \mathcal{F}_{\overline{s}}$
pour tout faisceau d'ensembles $\mathcal{F}$. En particulier,
$$
\forall p\geq 1, \quad H_\etale^p(S, \mathcal{F})=0
$$
pour tout $\mathcal{F}\in \textit{Ab}(S_\etale)$.
\end{lemma}

\begin{proof}
Si l'on montre que $\Gamma(S, \mathcal{F}) = \mathcal{F}_{\overline{s}}$,
alors $\Gamma(S, -)$ est exact, puisque le foncteur fibre est exact.
Soit $(U, \overline{u})$ un voisinage \'etale de $\overline{s}$.
Choisissons un voisinage ouvert affine $\Spec(A)$ de $\overline{u}$ dans $U$.
Alors $R \to A$ est \'etale et $\kappa(\overline{s}) = \kappa(\overline{u})$.
D'apr\`es le th\'eor\`eme \ref{theorem-henselian}, on voit que $A \cong R \times A'$
comme $R$-alg\`ebre, de mani\`ere compatible avec les morphismes vers
$\kappa(\overline{s}) = \kappa(\overline{u})$.
On obtient donc une section
$$
\xymatrix{
\Spec(A) \ar[r] & U \ar[d]\\
& S \ar[ul]
}
$$
Il s'ensuit que, dans le syst\`eme des voisinages \'etales de $\overline{s}$,
le morphisme identit\'e $(S, \overline{s}) \to (S, \overline{s})$ est cofinal.
Ainsi $\Gamma(S, \mathcal{F}) = \mathcal{F}_{\overline{s}}$.
La derni\`ere assertion du lemme r\'esulte du fait que les foncteurs d\'eriv\'es sup\'erieurs
d'un foncteur exact sont nuls, voir
Cat\'egories d\'eriv\'ees, Lemme \ref{derived-lemma-right-derived-exact-functor}.
\end{proof}

\begin{proposition}
\label{proposition-finite-higher-direct-image-zero}
Soit $f : X \to Y$ un morphisme fini de sch\'emas.
\begin{enumerate}
\item Pour tout point g\'eom\'etrique $\overline{y} : \Spec(k) \to Y$, on a
$$
(f_*\mathcal{F})_{\overline{y}} =
\prod\nolimits_{\overline{x} : \Spec(k) \to X,\ f(\overline{x}) =
\overline{y}} \mathcal{F}_{\overline{x}}.
$$
pour $\mathcal{F}$ dans $\Sh(X_\etale)$ et
$$
(f_*\mathcal{F})_{\overline{y}} =
\bigoplus\nolimits_{\overline{x} : \Spec(k) \to X,\ f(\overline{x}) =
\overline{y}} \mathcal{F}_{\overline{x}}.
$$
pour $\mathcal{F}$ dans $\textit{Ab}(X_\etale)$.
\item Pour tout $q \geq 1$, on a $R^q f_*\mathcal{F} = 0$
pour $\mathcal{F}$ dans $\textit{Ab}(X_\etale)$.
\end{enumerate}
\end{proposition}

\begin{proof}
Notons $X_{\overline{y}}^{sh}$ le produit fibr\'e
$X \times_Y \Spec(\mathcal{O}_{Y, \overline{y}}^{sh})$.
D'apr\`es le th\'eor\`eme \ref{theorem-higher-direct-images},
la fibre de $R^qf_*\mathcal{F}$ en $\overline{y}$ se calcule par
$H_\etale^q(X_{\overline{y}}^{sh}, \mathcal{F})$.
Puisque $f$ est fini, $X_{\bar y}^{sh}$ est fini sur
$\Spec(\mathcal{O}_{Y, \overline{y}}^{sh})$ ; ainsi,
$X_{\bar y}^{sh} = \Spec(A)$ pour un anneau $A$
fini sur $\mathcal{O}_{Y, \bar y}^{sh}$.
Comme ce dernier est strictement hens\'elien,
le lemme \ref{lemma-finite-over-henselian}
implique que $A$ est un produit fini d'anneaux locaux hens\'eliens
$A = A_1 \times \ldots \times A_r$. Puisque le corps r\'esiduel de
$\mathcal{O}_{Y, \overline{y}}^{sh}$ est s\'eparablement clos, il en va de
m\^eme pour chaque $A_i$. Ainsi $A_i$ est strictement hens\'elien.
Il s'ensuit que $X_{\overline{y}}^{sh} = \coprod_{i = 1}^r \Spec(A_i)$.
L'annulation du
lemme \ref{lemma-vanishing-etale-cohomology-strictly-henselian}
implique que $(R^qf_*\mathcal{F})_{\overline{y}} = 0$ pour $q > 0$ ;
cela implique (2) d'apr\`es le th\'eor\`eme \ref{theorem-exactness-stalks}.
L'assertion (1) r\'esulte de l'assertion correspondante du
lemme \ref{lemma-vanishing-etale-cohomology-strictly-henselian}.
\end{proof}

\begin{lemma}
\label{lemma-finite-pushforward-commutes-with-base-change}
Consid\'erons un carr\'e cart\'esien
$$
\xymatrix{
X' \ar[r]_{g'} \ar[d]_{f'} & X \ar[d]^f \\
Y' \ar[r]^g & Y
}
$$
de sch\'emas, avec $f$ un morphisme fini.
Pour tout faisceau d'ensembles $\mathcal{F}$ sur $X_\etale$, on a
$f'_*(g')^{-1}\mathcal{F} = g^{-1}f_*\mathcal{F}$.
\end{lemma}

\begin{proof}
Dans une grande g\'en\'eralit\'e, il existe un morphisme de changement de base
$g^{-1}f_*\mathcal{F} \to f'_*(g')^{-1}\mathcal{F}$, voir
Sites, section \ref{sites-section-pullback}.
Il suffit de v\'erifier l'assertion sur les fibres (th\'eor\`eme \ref{theorem-exactness-stalks}).
Soit $\overline{y}' : \Spec(k) \to Y'$ un point g\'eom\'etrique.
On a
\begin{align*}
(f'_*(g')^{-1}\mathcal{F})_{\overline{y}'}
& =
\prod\nolimits_{\overline{x}' : \Spec(k) \to X',\ f' \circ \overline{x}' =
\overline{y}'}
((g')^{-1}\mathcal{F})_{\overline{x}'} \\
& =
\prod\nolimits_{\overline{x}' : \Spec(k) \to X',\ f' \circ \overline{x}' =
\overline{y}'} \mathcal{F}_{g' \circ \overline{x}'} \\
& =
\prod\nolimits_{\overline{x} : \Spec(k) \to X,\ f \circ \overline{x} =
g \circ \overline{y}'} \mathcal{F}_{\overline{x}} \\
& =
(f_*\mathcal{F})_{g \circ \overline{y}'} \\
& =
(g^{-1}f_*\mathcal{F})_{\overline{y}'}
\end{align*}
La premi\`ere \'egalit\'e r\'esulte de la
proposition \ref{proposition-finite-higher-direct-image-zero}.
La deuxi\`eme \'egalit\'e r\'esulte du
lemme \ref{lemma-stalk-pullback}.
La troisi\`eme \'egalit\'e vaut parce que le diagramme est un carr\'e cart\'esien
et que, par cons\'equent, l'application
$$
\{\overline{x}' : \Spec(k) \to X',\ f' \circ \overline{x}' =
\overline{y}'\}
\longrightarrow
\{\overline{x} : \Spec(k) \to X,\ f \circ \overline{x} =
g \circ \overline{y}'\}
$$
qui envoie $\overline{x}'$ sur $g' \circ \overline{x}'$ est bijective.
La quatri\`eme \'egalit\'e r\'esulte de la
proposition \ref{proposition-finite-higher-direct-image-zero}.
La cinqui\`eme \'egalit\'e r\'esulte du
lemme \ref{lemma-stalk-pullback}.
\end{proof}

\begin{lemma}
\label{lemma-integral-pushforward-commutes-with-base-change}
Consid\'erons un carr\'e cart\'esien
$$
\xymatrix{
X' \ar[r]_{g'} \ar[d]_{f'} & X \ar[d]^f \\
Y' \ar[r]^g & Y
}
$$
de sch\'emas, avec $f$ un morphisme entier.
Pour tout faisceau d'ensembles $\mathcal{F}$ sur $X_\etale$, on a
$f'_*(g')^{-1}\mathcal{F} = g^{-1}f_*\mathcal{F}$.
\end{lemma}

\begin{proof}
La question est locale sur $Y$ ; nous pouvons donc supposer $Y$ affine.
On peut alors \'ecrire $X = \lim X_i$, o\`u $f_i : X_i \to Y$ est fini
(c'est imm\'ediat dans le cas affine, mais voir
Limites, Lemme \ref{limits-lemma-integral-limit-finite-and-finite-presentation}
pour une r\'ef\'erence). Notons $p_{i'i} : X_{i'} \to X_i$
les morphismes de transition et $p_i : X \to X_i$ les projections.
En posant $\mathcal{F}_i = p_{i, *}\mathcal{F}$, on obtient
du lemme \ref{lemma-linus-hamann}
un syst\`eme $(\mathcal{F}_i, \varphi_{i'i})$
tel que $\mathcal{F} = \colim p_i^{-1}\mathcal{F}_i$.
On obtient $f_*\mathcal{F} = \colim f_{i, *}\mathcal{F}_i$
par le lemme \ref{lemma-relative-colimit}.
Posons $X'_i = Y' \times_Y X_i$, de projections $f'_i$ et $g'_i$.
Alors $X' = \lim X'_i$, puisque les limites commutent aux limites.
Notons $p'_i : X' \to X'_i$ les projections. On a
\begin{align*}
g^{-1}f_*\mathcal{F}
& =
g^{-1} \colim f_{i, *}\mathcal{F}_i \\
& =
\colim g^{-1}f_{i, *}\mathcal{F}_i \\
& =
\colim f'_{i, *}(g'_i)^{-1}\mathcal{F}_i \\
& =
f'_*(\colim (p'_i)^{-1}(g'_i)^{-1}\mathcal{F}_i) \\
& =
f'_*(\colim (g')^{-1}p_i^{-1}\mathcal{F}_i) \\
& =
f'_*(g')^{-1} \colim p_i^{-1}\mathcal{F}_i \\
& =
f'_*(g')^{-1}\mathcal{F}
\end{align*}
comme voulu. Pour la premi\`ere \'egalit\'e, voir ci-dessus.
Pour la deuxi\`eme, on utilise le fait que l'image inverse commute aux limites inductives.
Pour la troisi\`eme, on utilise le cas fini, voir
le lemme \ref{lemma-finite-pushforward-commutes-with-base-change}.
Pour la quatri\`eme, on utilise le lemme \ref{lemma-relative-colimit}.
Pour la cinqui\`eme, on utilise $g'_i \circ p'_i = p_i \circ g'$.
Pour la sixi\`eme, on utilise le fait que l'image inverse commute aux limites inductives.
Pour la septi\`eme, on utilise $\mathcal{F} = \colim p_i^{-1}\mathcal{F}_i$.
\end{proof}

\noindent
Le lemme suivant est un cas de descente cohomologique portant sur
les faisceaux \'etales et les morphismes finis surjectifs. Nous g\'en\'eraliserons
consid\'erablement ce r\'esultat une fois d\'emontr\'e le th\'eor\`eme de changement de base propre.

\begin{lemma}
\label{lemma-cohomological-descent-finite}
Soit $f : X \to Y$ un morphisme fini surjectif de sch\'emas.
Posons $f_n : X_n \to Y$ \'egal au produit fibr\'e it\'er\'e $(n + 1)$ fois
de $X$ au-dessus de $Y$. Pour $\mathcal{F} \in \textit{Ab}(Y_\etale)$, posons
$\mathcal{F}_n = f_{n, *}f_n^{-1}\mathcal{F}$. Il existe une suite
exacte
$$
0 \to \mathcal{F} \to \mathcal{F}_0 \to \mathcal{F}_1 \to
\mathcal{F}_2 \to \ldots
$$
sur $Y_\etale$. De plus, il existe une suite spectrale
$$
E_1^{p, q} = H^q_\etale(X_p, f_p^{-1}\mathcal{F})
$$
convergeant vers $H^{p + q}(Y_\etale, \mathcal{F})$.
Cette suite spectrale est fonctorielle en $\mathcal{F}$.
\end{lemma}

\begin{proof}
Si nous d\'emontrons la premi\`ere assertion du lemme, nous obtenons une suite
spectrale de terme $E_1^{p, q} = H^q_\etale(Y, \mathcal{F}_p)$ convergeant
vers $H^{p + q}(Y_\etale, \mathcal{F})$, voir
Cat\'egories d\'eriv\'ees, lemme \ref{derived-lemma-two-ss-complex-functor}.
D'autre part, puisque
$R^if_{p, *}f_p^{-1}\mathcal{F} = 0$ pour $i > 0$
(proposition \ref{proposition-finite-higher-direct-image-zero})
on obtient
$$
H^q_\etale(X_p, f_p^{-1}\mathcal{F}) =
H^q_\etale(Y, f_{p, *}f_p^{-1} \mathcal{F}) =
H^q_\etale(Y, \mathcal{F}_p)
$$
d'apr\`es la proposition \ref{proposition-leray},
ce qui donne la suite spectrale du lemme.

\medskip\noindent
Pour d\'emontrer la premi\`ere assertion du lemme, observons que
les $X_n$ forment un sch\'ema simplicial au-dessus de $Y$, voir
Simplicial, exemple \ref{simplicial-example-fibre-products-simplicial-object}.
Observons de plus que, pour chacune des projections
$d_j : X_{n + 1} \to X_n$, il existe un morphisme
$d_j^{-1} f_n^{-1}\mathcal{F} \to f_{n + 1}^{-1}\mathcal{F}$.
Ces morphismes induisent des morphismes
$$
\delta_j : \mathcal{F}_n \to \mathcal{F}_{n + 1}
$$
pour $j = 0, \ldots, n + 1$. Nous utilisons la somme altern\'ee de ces morphismes
pour d\'efinir les diff\'erentielles $\mathcal{F}_n \to \mathcal{F}_{n + 1}$.
De m\^eme, il existe une augmentation canonique $\mathcal{F} \to \mathcal{F}_0$,
qui est simplement le morphisme canonique $\mathcal{F} \to f_*f^{-1}\mathcal{F}$.
Pour v\'erifier que cette suite de faisceaux est un complexe exact, il suffit
de le v\'erifier sur les fibres aux points g\'eom\'etriques (th\'eor\`eme \ref{theorem-exactness-stalks}).
Soit donc $\overline{y} : \Spec(k) \to Y$ un point g\'eom\'etrique. Posons
$E = \{\overline{x} : \Spec(k) \to X \mid f(\overline{x}) = \overline{y}\}$.
Alors $E$ est un ensemble fini non vide et l'on voit que
$$
(\mathcal{F}_n)_{\overline{y}} =
\bigoplus\nolimits_{e \in E^{n + 1}} \mathcal{F}_{\overline{y}}
$$
d'apr\`es la proposition \ref{proposition-finite-higher-direct-image-zero}
et le lemme \ref{lemma-stalk-pullback}.
Il reste donc \`a voir que, pour tout groupe ab\'elien $M$, la suite
$$
0 \to M \to \bigoplus\nolimits_{e \in E} M \to
\bigoplus\nolimits_{e \in E^2} M \to \ldots
$$
est exacte. Ici, le premier morphisme est le morphisme diagonal et le morphisme
$\bigoplus_{e \in E^{n + 1}} M  \to \bigoplus_{e \in E^{n + 2}} M$
est la somme altern\'ee des morphismes induits par les $(n + 2)$
projections $E^{n + 2} \to E^{n + 1}$. Cela se d\'emontre directement
ou se d\'eduit en appliquant Simplicial, lemme
\ref{simplicial-lemma-fibre-products-simplicial-object-w-section}
au morphisme $E \to \{*\}$.
\end{proof}

\begin{remark}
\label{remark-cohomological-descent-finite}
Dans la situation du lemme \ref{lemma-cohomological-descent-finite},
si $\mathcal{G}$ est un faisceau d'ensembles sur $Y_\etale$, alors on a
$$
\Gamma(Y, \mathcal{G}) =
\text{\'Egaliseur}(
\xymatrix{
\Gamma(X_0, f_0^{-1}\mathcal{G})
\ar@<1ex>[r] \ar@<-1ex>[r] &
\Gamma(X_1, f_1^{-1}\mathcal{G})
}
)
$$
Cela se d\'emontre exactement de la m\^eme mani\`ere, en montrant que
le faisceau $\mathcal{G}$ est l'\'egalisateur des deux morphismes
$f_{0, *}f_0^{-1}\mathcal{G} \to f_{1, *}f_1^{-1}\mathcal{G}$.
\end{remark}









%10.06.09
\section{Action galoisienne sur les fibres}
\label{section-galois-action-stalks}

\noindent
Dans cette section, nous d\'efinissons une action du groupe de Galois absolu du corps
r\'esiduel d'un point $s$ de $S$ sur le foncteur fibre en tout point g\'eom\'etrique situ\'e
au-dessus de $s$.

\medskip\noindent
Action galoisienne sur les fibres.
Soit $S$ un sch\'ema.
Soit $\overline{s}$ un point g\'eom\'etrique de $S$.
Soit $\sigma \in \text{Aut}(\kappa(\overline{s})/\kappa(s))$.
D\'efinissons comme suit une action de $\sigma$ sur la fibre $\mathcal{F}_{\overline{s}}$
d'un faisceau $\mathcal{F}$
\begin{equation}
\label{equation-galois-action}
\begin{matrix}
\mathcal{F}_{\overline{s}} &
\longrightarrow &
\mathcal{F}_{\overline{s}} \\
(U, \overline{u}, t) &
\longmapsto &
(U, \overline{u} \circ \Spec(\sigma), t).
\end{matrix}
\end{equation}
o\`u nous utilisons la description des \'el\'ements de la fibre au moyen de triplets,
comme dans la discussion qui suit la
d\'efinition \ref{definition-stalk}.
Il s'agit d'une action \`a gauche, car si
$\sigma_i \in \text{Aut}(\kappa(\overline{s})/\kappa(s))$
alors
\begin{align*}
\sigma_1 \cdot (\sigma_2 \cdot (U, \overline{u}, t))
& =
\sigma_1 \cdot (U, \overline{u} \circ \Spec(\sigma_2), t) \\
& =
(U, \overline{u} \circ \Spec(\sigma_2) \circ \Spec(\sigma_1), t) \\
& =
(U, \overline{u} \circ \Spec(\sigma_1 \circ \sigma_2), t) \\
& =
(\sigma_1 \circ \sigma_2) \cdot (U, \overline{u}, t)
\end{align*}
Il est clair que cette action est fonctorielle par rapport au faisceau $\mathcal{F}$.
Notons que nous aurions pu d\'efinir cette action en renvoyant directement \`a la
remarque \ref{remark-map-stalks}.

\begin{definition}
\label{definition-algebraic-geometric-point}
Soit $S$ un sch\'ema.
Soit $\overline{s}$ un point g\'eom\'etrique situ\'e au-dessus du point $s$ de $S$.
Notons $\kappa(s) \subset \kappa(s)^{sep} \subset \kappa(\overline{s})$
la cl\^oture s\'eparable de $\kappa(s)$ dans le corps alg\'ebriquement
clos $\kappa(\overline{s})$.
\begin{enumerate}
\item Dans cette situation, le {\it groupe de Galois absolu} de $\kappa(s)$
est $\text{Gal}(\kappa(s)^{sep}/\kappa(s))$. On le note parfois
$\text{Gal}_{\kappa(s)}$.
\item Le point g\'eom\'etrique $\overline{s}$ est dit
{\it alg\'ebrique} si $\kappa(s) \subset \kappa(\overline{s})$ est
une cl\^oture alg\'ebrique de $\kappa(s)$.
\end{enumerate}
\end{definition}

\begin{example}
\label{example-stupid}
Le point g\'eom\'etrique
$\Spec(\mathbf{C}) \to \Spec(\mathbf{Q})$
n'est pas alg\'ebrique.
\end{example}

\noindent
Soit $\kappa(s) \subset \kappa(s)^{sep} \subset \kappa(\overline{s})$
comme dans la d\'efinition. Puisque $\kappa(\overline{s})$ est alg\'ebriquement
clos, le morphisme
$$
\text{Aut}(\kappa(\overline{s})/\kappa(s))
\longrightarrow
\text{Gal}(\kappa(s)^{sep}/\kappa(s)) = \text{Gal}_{\kappa(s)}
$$
est surjectif. Supposons que $(U, \overline{u})$ soit un
voisinage \'etale de $\overline{s}$ et que $\overline{u}$ soit situ\'e au-dessus
du point $u$ de $U$. Puisque $U \to S$ est \'etale, l'extension de corps r\'esiduels
$\kappa(u)/\kappa(s)$ est finie s\'eparable.
Il s'ensuit ce qui suit
\begin{enumerate}
\item Si $\sigma \in \text{Aut}(\kappa(\overline{s})/\kappa(s)^{sep})$,
alors $\sigma$ agit trivialement sur $\mathcal{F}_{\overline{s}}$.
\item Plus pr\'ecis\'ement, l'action de
$\text{Aut}(\kappa(\overline{s})/\kappa(s))$
d\'etermine une action du groupe de Galois absolu
$\text{Gal}_{\kappa(s)}$ sur $\mathcal{F}_{\overline{s}}$, et est d\'etermin\'ee par elle.
\item Si $(U, \overline{u}, t)$ repr\'esente un \'el\'ement $\xi$ de
$\mathcal{F}_{\overline{s}}$, tout \'el\'ement de
$\text{Gal}(\kappa(s)^{sep}/K)$ agit trivialement, o\`u
$\kappa(s) \subset K \subset \kappa(s)^{sep}$ est l'image de
$\overline{u}^\sharp : \kappa(u) \to \kappa(\overline{s})$.
\end{enumerate}
Au total, nous voyons que $\mathcal{F}_{\overline{s}}$ devient un
$\text{Gal}_{\kappa(s)}$-ensemble (voir
Groupes fondamentaux, d\'efinition \ref{pione-definition-G-set-continuous}).
Nous pouvons donc consid\'erer le foncteur fibre comme un foncteur
$$
\Sh(S_\etale) \longrightarrow
\text{Gal}_{\kappa(s)}\textit{-Sets},
\quad
\mathcal{F} \longmapsto \mathcal{F}_{\overline{s}}
$$
et, dor\'enavant, c'est g\'en\'eralement ainsi que nous consid\'ererons le foncteur fibre.

\begin{theorem}
\label{theorem-equivalence-sheaves-point}
Soit $S = \Spec(K)$, o\`u $K$ est un corps.
Soit $\overline{s}$ un point g\'eom\'etrique de $S$.
Notons $G = \text{Gal}_{\kappa(s)}$ le groupe de Galois absolu.
Le passage aux fibres induit une \'equivalence de cat\'egories
$$
\Sh(S_\etale) \longrightarrow G\textit{-Sets},
\quad
\mathcal{F} \longmapsto \mathcal{F}_{\overline{s}}.
$$
\end{theorem}

\begin{proof}
Construisons un quasi-inverse de ce foncteur. Dans
Groupes fondamentaux, lemme \ref{pione-lemma-sheaves-point},
nous avons vu que, pour tout $G$-ensemble $M$, il existe un morphisme \'etale
$X \to \Spec(K)$
tel que $\Mor_K(\Spec(K^{sep}), X)$ soit
isomorphe \`a $M$ comme $G$-ensemble. Consid\'erons le faisceau
$\mathcal{F}$ sur $\Spec(K)_\etale$ d\'efini par
la r\`egle $U \mapsto \Mor_K(U, X)$. C'est un faisceau, car la topologie
\'etale est sous-canonique. On voit alors que
$\mathcal{F}_{\overline{s}} = \Mor_K(\Spec(K^{sep}), X) = M$
comme $G$-ensembles (d\'etails omis). Cela fournit un quasi-inverse du foncteur et
ach\`eve la d\'emonstration.
\end{proof}

\begin{remark}
\label{remark-every-sheaf-representable}
Une autre mani\`ere d'\'enoncer la conclusion du
th\'eor\`eme \ref{theorem-equivalence-sheaves-point} et de
Groupes fondamentaux, lemme \ref{pione-lemma-sheaves-point},
consiste \`a dire que tout faisceau sur $\Spec(K)_\etale$ est repr\'esentable
par un sch\'ema $X$ \'etale sur $\Spec(K)$.
Cela ne signifie pas que tout faisceau soit repr\'esentable au sens de
Sites, d\'efinition \ref{sites-definition-representable-sheaf}.
La raison est que, dans notre construction de $\Spec(K)_\etale$,
nous avons choisi un ensemble suffisamment grand de sch\'emas \'etales sur $\Spec(K)$,
tandis que les faisceaux sur $\Spec(K)_\etale$ forment une classe propre.
\end{remark}

\begin{lemma}
\label{lemma-global-sections-point}
Hypoth\`eses et notations comme dans le
th\'eor\`eme \ref{theorem-equivalence-sheaves-point}.
Il existe une bijection fonctorielle
$$
\Gamma(S, \mathcal{F}) = (\mathcal{F}_{\overline{s}})^G
$$
\end{lemma}

\begin{proof}
Nous pouvons le d\'emontrer par des arguments formels et \`a l'aide du
th\'eor\`eme \ref{theorem-equivalence-sheaves-point},
comme suit. Pour un faisceau $\mathcal{F}$ correspondant au
$G$-ensemble $M = \mathcal{F}_{\overline{s}}$, on a
\begin{eqnarray*}
\Gamma(S, \mathcal{F}) & = &
\Mor_{\Sh(S_\etale)}(h_{\Spec(K)},  \mathcal{F})
\\
& = & \Mor_{G\textit{-Sets}}(\{*\}, M) \\
& = & M^G
\end{eqnarray*}
Ici, la premi\`ere identification est expliqu\'ee dans
Sites, sections \ref{sites-section-presheaves} et
\ref{sites-section-representable-sheaves},
la deuxi\`eme r\'esulte du
th\'eor\`eme \ref{theorem-equivalence-sheaves-point}
et la troisi\`eme est claire. Nous donnerons aussi une preuve directe\footnote{Pour
les saint Thomas parmi nous.}.

\medskip\noindent
Supposons que $t \in \Gamma(S, \mathcal{F})$ soit une section globale.
Alors le triplet $(S, \overline{s}, t)$ d\'efinit un \'el\'ement de
$\mathcal{F}_{\overline{s}}$ qui est clairement invariant sous
l'action de $G$. R\'eciproquement, supposons que $(U, \overline{u}, t)$
d\'efinisse un \'el\'ement invariant de $\mathcal{F}_{\overline{s}}$.
Nous pouvons alors r\'etr\'ecir $U$ et supposer que $U = \Spec(L)$ pour une
extension de corps finie s\'eparable $L$ de $K$, voir la
proposition \ref{proposition-etale-morphisms}.
Dans ce cas, le morphisme $\mathcal{F}(U) \to \mathcal{F}_{\overline{s}}$
est injectif, car pour tout morphisme de voisinages \'etales
$(U', \overline{u}') \to (U, \overline{u})$, le morphisme de restriction
$\mathcal{F}(U) \to \mathcal{F}(U')$ est injectif puisque $U' \to U$
est un recouvrement de $S_\etale$.
Apr\`es avoir un peu agrandi $L$, nous pouvons supposer que $K \subset L$ est une extension
galoisienne finie. Nous utilisons alors l'\'egalit\'e
$$
\Spec(L) \times_{\Spec(K)} \Spec(L)
=
\coprod\nolimits_{\sigma \in \text{Gal}(L/K)} \Spec(L)
$$
o\`u les morphismes $\Spec(L) \to \Spec(L \otimes_K L)$
proviennent des morphismes d'anneaux $a \otimes b \mapsto a\sigma(b)$. Nous
voyons donc que l'invariance de $(U, \overline{u}, t)$
sous tout $G$ implique que $t \in \mathcal{F}(\Spec(L))$
a la m\^eme image dans
$\mathcal{F}(\Spec(L) \times_{\Spec(K)} \Spec(L))$
par restriction suivant chacune des deux projections (on utilise ici l'injectivit\'e mentionn\'ee
ci-dessus; d\'etails omis). La condition de faisceau pour $\mathcal{F}$
et le recouvrement \'etale $\{\Spec(L) \to \Spec(K)\}$ s'applique
donc, et nous concluons que $t$ provient d'une unique section de $\mathcal{F}$
sur $\Spec(K)$.
\end{proof}

\begin{remark}
\label{remark-stalk-pullback}
Soit $S$ un sch\'ema et soit $\overline{s} : \Spec(k) \to S$
un point g\'eom\'etrique de $S$. Par d\'efinition, cela signifie que $k$
est alg\'ebriquement clos. En particulier, le groupe de Galois absolu de $k$
est trivial. Par cons\'equent, d'apr\`es le
th\'eor\`eme \ref{theorem-equivalence-sheaves-point},
la cat\'egorie des faisceaux sur $\Spec(k)_\etale$ est \'equivalente
\`a la cat\'egorie des ensembles. L'\'equivalence est donn\'ee par la prise des
sections sur $\Spec(k)$. Cela nous fournit enfin une
autre d\'efinition du foncteur fibre. Plus pr\'ecis\'ement, le foncteur
$$
\Sh(S_\etale) \longrightarrow \textit{Sets}, \quad
\mathcal{F} \longmapsto \mathcal{F}_{\overline{s}}
$$
est isomorphe au foncteur
$$
\Sh(S_\etale)
\longrightarrow
\Sh(\Spec(k)_\etale) = \textit{Sets},
\quad
\mathcal{F} \longmapsto \overline{s}^*\mathcal{F}
$$
Pour le d\'emontrer rigoureusement, on peut utiliser
la partie (3) du lemme \ref{lemma-stalk-pullback}
avec $f = \overline{s}$. De plus, ceci \'etant dit, le cas g\'en\'eral de la
partie (3) du lemme \ref{lemma-stalk-pullback}
r\'esulte de la fonctorialit\'e des images inverses.
\end{remark}






\section{Cohomologie des groupes}
\label{section-group-cohomology}

\noindent
Dans la suite, lorsque nous \'ecrivons $H^i(G, M)$, nous entendons que $G$
est un groupe topologique, que $M$ est un $G$-module discret muni d'une
action continue de $G$, et que $H^i(G, -)$ est le $i$-i\`eme foncteur d\'eriv\'e
droit sur la cat\'egorie $\text{Mod}_G$ de ces $G$-modules, voir les
d\'efinitions \ref{definition-G-module-continuous} et
\ref{definition-galois-cohomology}. Cela comprend
le cas d'un groupe abstrait $G$, ce qui signifie simplement que l'on consid\`ere $G$
comme un groupe topologique muni de la topologie discr\`ete.

\medskip\noindent
Lorsque le module est muni d'une topologie non discr\`ete, nous emploierons la notation
$H^i_{cont}(G, M)$ pour d\'esigner les groupes de cohomologie continue
introduits dans \cite{Tate}, voir la
section \ref{section-continuous-group-cohomology}.

\begin{definition}
\label{definition-G-module-continuous}
Soit $G$ un groupe topologique.
\begin{enumerate}
\item Un {\it $G$-module}, parfois appel\'e {\it $G$-module discret},
est un groupe ab\'elien $M$ muni d'une action \`a gauche $a : G \times M \to M$
par homomorphismes de groupes, telle que $a$ soit continue lorsque $M$ est muni de la
topologie discr\`ete.
\item Un {\it morphisme de $G$-modules} $f : M \to N$ est un
homomorphisme $G$-\'equivariant de $M$ dans $N$.
\item La cat\'egorie des $G$-modules est not\'ee $\text{Mod}_G$.
\end{enumerate}
Soit $R$ un anneau.
\begin{enumerate}
\item Un {\it $R\text{-}G$-module} est un $R$-module $M$ muni
d'une action \`a gauche $a : G \times M \to M$ par applications $R$-lin\'eaires, telle que $a$
soit continue lorsque $M$ est muni de la topologie discr\`ete.
\item Un {\it morphisme de $R\text{-}G$-modules} $f : M \to N$ est une
application $R$-lin\'eaire $G$-\'equivariante de $M$ dans $N$.
\item La cat\'egorie des $R\text{-}G$-modules est not\'ee $\text{Mod}_{R, G}$.
\end{enumerate}
\end{definition}

\noindent
La continuit\'e de $a : G \times M \to M$ \'equivaut
\`a demander que le stabilisateur de tout $x \in M$ soit ouvert dans $G$.
Si $G$ est un groupe abstrait, cela correspond \`a la notion d'un
groupe ab\'elien muni d'une action de $G$, pourvu que l'on munisse $G$ de la
topologie discr\`ete. Observons que $\text{Mod}_{\mathbf{Z}, G} = \text{Mod}_G$.

\medskip\noindent
La cat\'egorie $\text{Mod}_G$ a assez d'injectifs, voir
Injectifs, lemme \ref{injectives-lemma-G-modules}.
Consid\'erons le foncteur exact \`a gauche
$$
\text{Mod}_G \longrightarrow \textit{Ab},
\quad
M \longmapsto M^G =
\{x \in M \mid g \cdot x = x\ \forall g \in G\}
$$
Nous notons parfois $M^G = H^0(G, M)$ et parfois
$M^G = \Gamma_G(M)$. Ce foncteur admet un foncteur d\'eriv\'e droit total
$R\Gamma_G(M)$ et, pour tout $i \geq 0$, un $i$-i\`eme foncteur d\'eriv\'e droit
$R^i\Gamma_G(M) = H^i(G, M)$.

\medskip\noindent
La m\^eme construction s'applique \`a
$H^0(G, -) : \text{Mod}_{R, G} \to \text{Mod}_R$. Nous verrons au
lemme \ref{lemma-modules-abelian} qu'elle co\"\i ncide avec la cohomologie
du $G$-module sous-jacent.

\begin{definition}
\label{definition-galois-cohomology}
Soit $G$ un groupe topologique. Soit $M$ un $G$-module discret
muni d'une action continue de $G$. Autrement dit, $M$ est un objet
de la cat\'egorie $\text{Mod}_G$ introduite dans la
d\'efinition \ref{definition-G-module-continuous}.
\begin{enumerate}
\item Les foncteurs d\'eriv\'es droits $H^i(G, M)$ de $H^0(G, M)$ sur la
cat\'egorie $\text{Mod}_G$ sont appel\'es les
{\it groupes de cohomologie continue de $G$ \`a coefficients dans $M$}.
\item Si $G$ est un groupe abstrait muni de la topologie discr\`ete,
les $H^i(G, M)$ sont appel\'es les {\it groupes de cohomologie de $G$ \`a coefficients dans $M$}.
\item Si $G$ est un groupe de Galois, les groupes $H^i(G, M)$ sont appel\'es
les {\it groupes de cohomologie galoisienne \`a coefficients dans $M$}.
\item Si $G$ est le groupe de Galois absolu d'un corps $K$, les groupes
$H^i(G, M)$ sont parfois appel\'es les {\it groupes de cohomologie galoisienne de $K$
\`a coefficients dans $M$}. Dans ce cas, nous \'ecrivons parfois
$H^i(K, M)$ au lieu de $H^i(G, M)$.
\end{enumerate}
\end{definition}

\begin{lemma}
\label{lemma-modules-abelian}
Soit $G$ un groupe topologique. Soit $R$ un anneau.
Pour tout $i \geq 0$, le diagramme
$$
\xymatrix{
\text{Mod}_{R, G} \ar[rr]_{H^i(G, -)} \ar[d] & &
\text{Mod}_R \ar[d] \\
\text{Mod}_G \ar[rr]^{H^i(G, -)} & &
\textit{Ab}
}
$$
dont les fl\`eches verticales sont les foncteurs d'oubli est commutatif.
\end{lemma}

\begin{proof}
Notons $F : \text{Mod}_{R, G} \to \text{Mod}_G$ le foncteur d'oubli.
Alors $F$ admet un adjoint \`a gauche $H : \text{Mod}_G \to \text{Mod}_{R, G}$
donn\'e par $H(M) = M \otimes_\mathbf{Z} R$. Observons que tout objet de
$\text{Mod}_G$ est un quotient d'une somme directe de modules de la forme
$\mathbf{Z}[G/U]$, o\`u $U \subset G$ est un sous-groupe ouvert.
Ici, $\mathbf{Z}[G/U]$ d\'esigne le $G$-module form\'e des
combinaisons $\mathbf{Z}$-lin\'eaires finies
des classes \`a droite modulo $U$ dans $G$, muni de l'action \`a gauche de $G$.
Ainsi, tout complexe born\'e sup\'erieurement dans $\text{Mod}_G$ est quasi-isomorphe
\`a un complexe born\'e sup\'erieurement dans $\text{Mod}_G$ dont les
termes sous-jacents sont des $\mathbf{Z}$-modules plats
(Cat\'egories d\'eriv\'ees, lemme \ref{derived-lemma-subcategory-left-resolution}).
Il est donc clair que $LH$ existe sur $D^-(\text{Mod}_G)$ et se calcule en
appliquant $H$ \`a tout complexe dont les termes sont des $\mathbf{Z}$-modules plats;
cela r\'esulte de
Cat\'egories d\'eriv\'ees, lemme \ref{derived-lemma-subcategory-left-acyclics} et
proposition \ref{derived-proposition-enough-acyclics}.
Nous concluons de Cat\'egories d\'eriv\'ees, lemme
\ref{derived-lemma-pre-derived-adjoint-functors}
que
$$
\text{Ext}^i(\mathbf{Z}, F(M)) = \text{Ext}^i(R, M)
$$
pour $M$ dans $\text{Mod}_{R, G}$.
Observons que $H^0(G, -) = \Hom(\mathbf{Z}, -)$ sur
$\text{Mod}_G$, o\`u $\mathbf{Z}$ d\'esigne le $G$-module
\`a action triviale. Par cons\'equent,
$H^i(G, -) = \text{Ext}^i(\mathbf{Z}, -)$ sur $\text{Mod}_G$.
De m\^eme, on a $H^i(G, -) = \text{Ext}^i(R, -)$ sur
$\text{Mod}_{R, G}$. En combinant ces faits, on obtient le lemme.
\end{proof}

\begin{lemma}
\label{lemma-ext-modules-hom}
Soit $G$ un groupe topologique. Soit $R$ un anneau.
Soient $M$, $N$ des $R\text{-}G$-modules. Si $M$ est projectif de type fini
comme $R$-module, alors
$\text{Ext}^i(M, N) = H^i(G, M^\vee \otimes_R N)$ (pour les notations,
voir la preuve).
\end{lemma}

\begin{proof}
Munissons le module $M^\vee = \Hom_R(M, R)$ de l'action contragr\'ediente
de $G$. Ainsi, $(g \cdot \lambda)(m) = \lambda(g^{-1} \cdot m)$
pour $g \in G$, $\lambda \in M^\vee$, $m \in M$. L'action de $G$ sur
$M^\vee \otimes_R N$ est l'action diagonale, c'est-\`a-dire qu'elle est donn\'ee par
$g \cdot (\lambda \otimes n) = g \cdot \lambda \otimes g \cdot n$.
Notons que, pour un troisi\`eme $R\text{-}G$-module $E$, on a
$\Hom(E, M^\vee \otimes_R N) = \Hom(M \otimes_R E, N)$.
En effet, cela est vrai au niveau des $R$-modules d'apr\`es
Alg\`ebre, lemmes \ref{algebra-lemma-hom-from-tensor-product} et
\ref{algebra-lemma-evaluation-map-iso-finite-projective}
et les d\'efinitions des actions de $G$ sont choisies de sorte que cela
reste vrai pour les $R\text{-}G$-modules. Il en r\'esulte que
$M^\vee \otimes_R N$ est un $R\text{-}G$-module injectif
si $N$ est un $R\text{-}G$-module injectif. Ainsi, si
$N \to N^\bullet$ est une r\'esolution injective, alors
$M^\vee \otimes_R N \to M^\vee \otimes_R N^\bullet$
est une r\'esolution injective. On a alors
$$
\Hom(M, N^\bullet) = \Hom(R, M^\vee \otimes_R N^\bullet) =
(M^\vee \otimes_R N^\bullet)^G
$$
Puisque le membre de gauche calcule $\text{Ext}^i(M, N)$ et celui de droite
calcule $H^i(G, M^\vee \otimes_R N)$, la preuve est termin\'ee.
\end{proof}

\begin{lemma}
\label{lemma-finite-dim-group-cohomology}
Soit $G$ un groupe topologique. Soit $k$ un corps.
Soit $V$ un $k\text{-}G$-module.
Si $G$ est topologiquement engendr\'e par un nombre fini d'\'el\'ements et si
$\dim_k(V) < \infty$, alors $\dim_k H^1(G, V) < \infty$.
\end{lemma}

\begin{proof}
Soient $g_1, \ldots, g_r \in G$ des \'el\'ements qui engendrent topologiquement $G$,
c'est-\`a-dire que le sous-groupe engendr\'e par $g_1, \ldots, g_r$ est dense.
D'apr\`es le lemme \ref{lemma-ext-modules-hom},
on voit que $H^1(G, V)$ est l'espace vectoriel sur $k$ des extensions
$$
0 \to V \to E \to k \to 0
$$
de $k\text{-}G$-modules. Choisissons $e \in E$ d'image $1 \in k$.
\'Ecrivons
$$
g_i \cdot e = v_i + e
$$
pour un $v_i \in V$. Cela est possible puisque $g_i \cdot 1 = 1$.
Nous affirmons que la liste des \'el\'ements $v_1, \ldots, v_r \in V$
d\'etermine la classe d'isomorphisme de l'extension $E$.
Une fois ceci d\'emontr\'e, le lemme en r\'esulte, car notre espace vectoriel
Ext est alors isomorphe \`a un sous-quotient de l'espace vectoriel sur $k$
$V^{\oplus r}$; quelques d\'etails sont omis.
Puisque $E$ est un objet de la cat\'egorie d\'efinie dans la
d\'efinition \ref{definition-G-module-continuous},
il existe un sous-groupe ouvert $U$ tel que
$u \cdot e = e$ pour tout $u \in U$.
Prenons maintenant un $g \in G$. Alors $gU$ contient un mot $w$ en
les \'el\'ements $g_1, \ldots, g_r$.
\'Ecrivons $gu = w$. Puisque l'\'el\'ement $w \cdot e$ est d\'etermin\'e par
$v_1, \ldots, v_r$, on voit que $g \cdot e = (gu) \cdot e = w \cdot e$
l'est lui aussi.
\end{proof}

\begin{lemma}
\label{lemma-profinite-group-cohomology-torsion}
Soit $G$ un groupe topologique profini.
Alors
\begin{enumerate}
\item $H^i(G, M)$ est de torsion pour $i > 0$ et tout $G$-module $M$, et
\item $H^i(G, M) = 0$ si $M$ est un $\mathbf{Q}$-espace vectoriel.
\end{enumerate}
\end{lemma}

\begin{proof}
Preuve de (1). Par d\'ecalage de dimension, il suffit
de montrer que $H^1(G, M)$ est de torsion pour tout $G$-module $M$.
Choisissons une suite exacte $0 \to M \to I \to N \to 0$, o\`u $I$
est un objet injectif de la cat\'egorie des $G$-modules.
Tout \'el\'ement de $H^1(G, M)$ est alors l'image d'un \'el\'ement
$y \in N^G$. Choisissons $x \in I$ d'image $y$.
Le stabilisateur $U \subset G$ de $x$ est ouvert, donc
d'indice fini $r$. Soit $g_1, \ldots, g_r \in G$ un syst\`eme
de repr\'esentants de $G/U$. Alors $\sum g_i(x)$ est un \'el\'ement
invariant de $I$ dont l'image est $ry$. Ainsi, $r$ annule l'\'el\'ement
de $H^1(G, M)$ dont nous \'etions partis. La partie (2) en r\'esulte, puisque
$H^i(G, M)$ est alors \`a la fois un $\mathbf{Q}$-espace vectoriel et de torsion.
\end{proof}





\section{Cohomologie continue de Tate}
\label{section-continuous-group-cohomology}

\noindent
La cohomologie continue de Tate (\cite{Tate}) est d\'efinie par le complexe des
cocha\^ines inhomog\`enes continues. On peut la d\'efinir lorsque $M$ est un
groupe ab\'elien topologique quelconque muni d'une action continue de $G$.
Plus pr\'ecis\'ement, consid\'erons le complexe
$$
C^\bullet_{cont}(G, M) :
M \to \text{Maps}_{cont}(G, M) \to
\text{Maps}_{cont}(G \times G, M) \to \ldots
$$
o\`u l'application de bord est d\'efinie, pour $n \geq 1$, par la formule
\begin{align*}
\text{d}(f)(g_1, \ldots, g_{n + 1})
& = g_1(f(g_2, \ldots, g_{n + 1})) \\
&
+ \sum\nolimits_{j = 1, \ldots, n}
(-1)^jf(g_1, \ldots, g_jg_{j + 1}, \ldots, g_{n + 1}) \\
&
+ (-1)^{n + 1}f(g_1, \ldots, g_n)
\end{align*}
et, pour $n = 0$, envoie $m \in M$ sur l'application $g \mapsto g(m) - m$. Posons
$$
H^i_{cont}(G, M) = H^i(C^\bullet_{cont}(G, M))
$$
Comme les termes du complexe font intervenir les applications continues de $G$ et
de ses puissances cart\'esiennes dans le module topologique $M$, il n'est pas clair
qu'une suite exacte courte de modules topologiques soit ainsi transform\'ee en
une suite exacte longue de cohomologie. Une autre difficult\'e tient \`a ce que la cat\'egorie
des groupes ab\'eliens topologiques n'est pas une cat\'egorie ab\'elienne !

\medskip\noindent
Cependant, une suite exacte courte de $G$-modules discrets fournit bien
une suite exacte courte de complexes de cocha\^ines continues
et donc une suite exacte longue de groupes de cohomologie
continue $H^i_{cont}(G, -)$.
Par cons\'equent, dans la cat\'egorie $\text{Mod}_G$ de la
d\'efinition \ref{definition-G-module-continuous}, les foncteurs
$H^i_{cont}(G, M)$ forment un $\delta$-foncteur cohomologique au sens de
Homologie, section \ref{homology-section-cohomological-delta-functor}.
Puisque la cohomologie $H^i(G, M)$
de la d\'efinition \ref{definition-galois-cohomology}
est un $\delta$-foncteur universel
(Cat\'egories d\'eriv\'ees, lemme \ref{derived-lemma-right-derived-delta-functor}),
on obtient des morphismes canoniques
$$
H^i(G, M) \longrightarrow H^i_{cont}(G, M)
$$
pour $M \in \text{Mod}_G$. On sait que ces morphismes sont des
isomorphismes lorsque $G$ est un groupe abstrait (c'est-\`a-dire muni de
la topologie discr\`ete) ou lorsque $G$ est un groupe profini
(ins\'erer ici une r\'ef\'erence ult\'erieure).
Si vous connaissez un exemple o\`u ce morphisme n'est pas un isomorphisme
pour un groupe topologique $G$ et $M \in \Ob(\text{Mod}_G)$,
veuillez l'envoyer par courriel \`a
\href{mailto:stacks.project@gmail.com}{stacks.project@gmail.com}.




\section{Cohomologie d'un point}
\label{section-cohomology-point}

\noindent
Comme cons\'equence de la discussion des sections pr\'ec\'edentes,
nous obtenons l'\'equivalence entre la cohomologie \'etale du spectre d'un
corps et la cohomologie galoisienne.

\begin{lemma}
\label{lemma-equivalence-abelian-sheaves-point}
Soit $S = \Spec(K)$, o\`u $K$ est un corps.
Soit $\overline{s}$ un point g\'eom\'etrique de $S$.
Notons $G = \text{Gal}_{\kappa(s)}$ le groupe de Galois absolu.
Le foncteur fibre induit une \'equivalence de cat\'egories
$$
\textit{Ab}(S_\etale) \longrightarrow \text{Mod}_G,
\quad
\mathcal{F} \longmapsto \mathcal{F}_{\overline{s}}.
$$
\end{lemma}

\begin{proof}
Dans le
th\'eor\`eme \ref{theorem-equivalence-sheaves-point},
nous avons vu l'\'equivalence entre les faisceaux d'ensembles et les $G$-ensembles.
Le pr\'esent lemme s'en d\'eduit formellement, car un faisceau ab\'elien n'est autre
qu'un faisceau d'ensembles muni d'une loi de groupe commutative, et un $G$-module
n'est autre qu'un $G$-ensemble muni d'une loi de groupe commutative.
\end{proof}

\begin{lemma}
\label{lemma-compare-cohomology-point}
Les notations et hypoth\`eses sont celles du
lemme \ref{lemma-equivalence-abelian-sheaves-point}.
Soit $\mathcal{F}$ un faisceau ab\'elien sur $\Spec(K)_\etale$
qui correspond au $G$-module $M$.
Alors
\begin{enumerate}
\item dans $D(\textit{Ab})$, on a un isomorphisme canonique
$R\Gamma(S, \mathcal{F}) = R\Gamma_G(M)$,
\item $H_\etale^0(S, \mathcal{F}) = M^G$, et
\item $H_\etale^q(S, \mathcal{F}) = H^q(G, M)$.
\end{enumerate}
\end{lemma}

\begin{proof}
Combiner le
lemme \ref{lemma-equivalence-abelian-sheaves-point}
avec le
lemme \ref{lemma-global-sections-point}.
\end{proof}

\begin{example}
\label{example-sheaves-point}
Faisceaux sur $\Spec(K)_\etale$.
Soit $G = \text{Gal}(K^{sep}/K)$ le groupe de Galois absolu de $K$.
\begin{enumerate}
\item Le faisceau constant $\underline{\mathbf{Z}/n\mathbf{Z}}$ correspond au
module $\mathbf{Z}/n\mathbf{Z}$ muni de l'action triviale de $G$,
\item le faisceau $\mathbf{G}_m|_{\Spec(K)_\etale}$ correspond \`a
$(K^{sep})^*$ muni de son action de $G$,
\item le faisceau $\mathbf{G}_a|_{\Spec(K^{sep})}$ correspond \`a
$(K^{sep}, +)$ muni de son action de $G$, et
\item le faisceau $\mu_n|_{\Spec(K^{sep})}$ correspond \`a
$\mu_n(K^{sep})$ muni de son action de $G$.
\end{enumerate}
D'apr\`es la
remarque \ref{remark-special-case-fpqc-cohomology-quasi-coherent}
et le
th\'eor\`eme \ref{theorem-picard-group},
on a les identifications suivantes pour les groupes de cohomologie :
\begin{align*}
H_\etale^0(S_\etale, \mathbf{G}_m) & =
\Gamma(S, \mathcal{O}_S^*) \\
H_\etale^1(S_\etale, \mathbf{G}_m) & =
H_{Zar}^1(S, \mathcal{O}_S^*) = \Pic(S) \\
H_\etale^i(S_\etale, \mathbf{G}_a) & =
H_{Zar}^i(S, \mathcal{O}_S)
\end{align*}
De plus, pour tout faisceau quasi-coh\'erent $\mathcal{F}$ sur $S_\etale$, on a
$$
H^i(S_\etale, \mathcal{F}) = H_{Zar}^i(S, \mathcal{F}),
$$
voir le
th\'eor\`eme \ref{theorem-zariski-fpqc-quasi-coherent}.
En particulier, on obtient la suite d'\'egalit\'es suivante
$$
0 =
\Pic(\Spec(K)) =
H_\etale^1(\Spec(K)_\etale, \mathbf{G}_m) =
H^1(G, (K^{sep})^*)
$$
qui n'est autre que le th\'eor\`eme 90 de Hilbert. De m\^eme, pour $i \geq 1$,
$$
0 = H^i(\Spec(K), \mathcal{O})
= H_\etale^i(\Spec(K)_\etale, \mathbf{G}_a)
= H^i(G, K^{sep})
$$
o\`u $K^{sep}$ est ici consid\'er\'e comme module galoisien, l'addition \'etant
sa loi de groupe. De la sorte, nous pouvons consid\'erer le travail accompli jusqu'ici comme
une mani\`ere compliqu\'ee de calculer des groupes de cohomologie galoisienne.
\end{example}

\noindent
Le r\'esultat suivant est une curiosit\'e et devrait \^etre omis lors d'une
premi\`ere lecture.

\begin{lemma}
\label{lemma-all-modules-quasi-coherent}
Soit $R$ un anneau local de dimension $0$. Soit $S = \Spec(R)$.
Alors tout $\mathcal{O}_S$-module sur $S_\etale$ est quasi-coh\'erent.
\end{lemma}

\begin{proof}
Soit $\mathcal{F}$ un $\mathcal{O}_S$-module sur $S_\etale$.
Nous devons montrer que $\mathcal{F}$ est d\'etermin\'e par le
$R$-module $M = \Gamma(S, \mathcal{F})$.
Plus pr\'ecis\'ement, si $\pi : X \to S$ est \'etale, nous devons
montrer que $\Gamma(X, \mathcal{F}) = \Gamma(X, \pi^*\widetilde{M})$.

\medskip\noindent
Soit $\mathfrak m \subset R$ l'id\'eal maximal et soit
$\kappa$ le corps r\'esiduel. D'apr\`es
Alg\`ebre, lemme \ref{algebra-lemma-local-dimension-zero-henselian},
l'anneau local $R$ est hens\'elien. Si $X \to S$ est \'etale,
alors l'espace topologique sous-jacent \`a $X$ est discret
d'apr\`es Morphismes, lemme \ref{morphisms-lemma-etale-over-field},
et $X$ est donc une r\'eunion disjointe de sch\'emas affines
ayant chacun un seul point. De plus, si $X = \Spec(A)$ est affine et
n'a qu'un point, alors $A$ est une $R$-alg\`ebre finie \'etale d'apr\`es
Alg\`ebre, lemme \ref{algebra-lemma-mop-up}.
Nous devons montrer que $\Gamma(X, \mathcal{F}) = M \otimes_R A$
dans ce cas.

\medskip\noindent
Le foncteur $A \mapsto A/\mathfrak m A$ d\'efinit une \'equivalence entre
la cat\'egorie des $R$-alg\`ebres finies \'etales
et celle des $\kappa$-alg\`ebres finies s\'eparables, d'apr\`es
Alg\`ebre, lemme \ref{algebra-lemma-henselian-cat-finite-etale}.
Consid\'erons d'abord le cas o\`u $A/\mathfrak m A$
est une extension galoisienne de $\kappa$, de groupe de Galois $G$.
Pour chaque $\sigma \in G$, notons $\sigma : A \to A$
l'automorphisme correspondant de $A$ au-dessus de $R$.
Posons $N = \Gamma(X, \mathcal{F})$.
Alors $\Spec(\sigma) : X \to X$ est un automorphisme au-dessus de $S$
et l'image inverse par celui-ci d\'efinit donc une application $\sigma : N \to N$
qui est $\sigma$-lin\'eaire : $\sigma(an) = \sigma(a) \sigma(n)$
pour $a \in A$ et $n \in N$.
Nous allons appliquer la descente galoisienne au module quasi-coh\'erent
$\widetilde{N}$ sur $X$, muni des isomorphismes
provenant de l'action de $\sigma$ sur $N$. Voir Descente, lemme
\ref{descent-lemma-galois-descent-more-general}.
Ce lemme fournit un isomorphisme $N = N^G \otimes_R A$.
D'autre part, il est clair que $N^G = M$ par la propri\'et\'e de faisceau
de $\mathcal{F}$. L'isomorphisme voulu est donc \'etabli.

\medskip\noindent
Le cas g\'en\'eral (o\`u $A$ est une $R$-alg\`ebre locale finie \'etale)
se d\'eduit du cas galoisien comme suit. Choisissons $A \to B$
finie \'etale, avec $B$ local et son corps r\'esiduel
galoisien sur $\kappa$. Posons $G = \text{Aut}(B/R) = \text{Gal}(\kappa_B/\kappa)$.
Soit $H \subset G$ le groupe de Galois correspondant \`a
l'extension galoisienne $\kappa_B/\kappa_A$. Alors, comme ci-dessus, on
montre que $\Gamma(X, \mathcal{F}) = \Gamma(\Spec(B), \mathcal{F})^H$.
En appliquant deux fois le r\'esultat relatif aux extensions galoisiennes, on obtient
$$
\Gamma(X, \mathcal{F}) = (M \otimes_R B)^H = M \otimes_R A
$$
comme voulu.
\end{proof}








\section{Cohomologie des courbes}
\label{section-cohomology-curves}

\noindent
La t\^ache suivante consiste \`a calculer la cohomologie \'etale d'une courbe lisse
sur un corps alg\'ebriquement clos, \`a coefficients de torsion, et en
particulier \`a montrer qu'elle s'annule en degr\'e au moins $3$. Pour cela, nous
calculerons la cohomologie au point g\'en\'erique, ce qui
revient \`a calculer de la cohomologie galoisienne.





\section{Groupes de Brauer}
\label{section-brauer-groups}

\noindent
Les groupes de Brauer des corps sont d\'efinis au moyen d'alg\`ebres centrales simples de dimension finie.
Dans cette section, nous rappelons les faits pertinents sur les groupes de Brauer, dont la plupart
sont expos\'es dans le chapitre
Groupes de Brauer, section \ref{brauer-section-introduction}.
Pour d'autres r\'ef\'erences, voir \cite{SerreCorpsLocaux},
\cite{SerreGaloisCohomology} ou \cite{Weil}.

\begin{theorem}
\label{theorem-central-simple-algebra}
Soit $K$ un corps. Pour une $K$-alg\`ebre $A$ unitaire et associative (non n\'ecessairement commutative),
les conditions suivantes sont \'equivalentes
\begin{enumerate}
\item $A$ est une $K$-alg\`ebre centrale simple de dimension finie,
\item $A$ est un $K$-espace vectoriel de dimension finie, $K$ est le centre de $A$,
et $A$ n'a pas d'id\'eal bilat\`ere non trivial,
\item il existe $d \geq 1$ tel que
$A \otimes_K \bar K \cong \text{Mat}(d \times d, \bar K)$,
\item il existe $d \geq 1$ tel que
$A \otimes_K K^{sep} \cong \text{Mat}(d \times d, K^{sep})$,
\item il existe $d \geq 1$ et une extension galoisienne finie $K'/K$
tels que
$A \otimes_K K' \cong \text{Mat}(d \times d, K')$,
\item il existe $n \geq 1$ et un corps gauche central $D$ de dimension finie
sur $K$ tel que $A \cong \text{Mat}(n \times n, D)$.
\end{enumerate}
L'entier $d$ est appel\'e le {\it degr\'e} de $A$.
\end{theorem}

\begin{proof}
Il s'agit d'une copie de
Groupes de Brauer, lemme \ref{brauer-lemma-finite-central-simple-algebra}.
\end{proof}

\begin{lemma}
\label{lemma-brauer-inverse}
Soit $A$ une alg\`ebre centrale simple de dimension finie sur $K$. Alors
$$
\begin{matrix}
A \otimes_K A^{opp} & \longrightarrow & \text{End}_K(A) \\
\ a \otimes a' & \longmapsto & (x \mapsto a x a')
\end{matrix}
$$
est un isomorphisme de $K$-alg\`ebres.
\end{lemma}

\begin{proof}
Voir
Groupes de Brauer, lemme \ref{brauer-lemma-inverse}.
\end{proof}

\begin{definition}
\label{definition-brauer-equivalent}
Deux alg\`ebres centrales simples de dimension finie $A_1$ et $A_2$ sur $K$ sont dites
{\it semblables}, ou {\it \'equivalentes}, s'il existe $m, n \geq 1$
tels que $\text{Mat}(n \times n, A_1)
\cong \text{Mat}(m \times m, A_2)$. On \'ecrit $A_1 \sim A_2$.
\end{definition}

\noindent
D'apr\`es Groupes de Brauer, lemme \ref{brauer-lemma-similar}, c'est une
relation d'\'equivalence.

\begin{definition}
\label{definition-brauer-group}
Soit $K$ un corps. Le {\it groupe de Brauer} de $K$ est l'ensemble $\text{Br} (K)$
des classes de similitude d'alg\`ebres centrales simples de dimension finie sur $K$, muni de
la loi de groupe induite par le produit tensoriel (sur $K$). La classe de $A$ dans
$\text{Br}(K)$ est not\'ee $[A]$. L'\'el\'ement neutre est
$[K] = [\text{Mat}(d \times d, K)]$ pour tout $d \geq 1$.
\end{definition}

\noindent
Le lemme pr\'ec\'edent entra\^ine l'existence des inverses et l'\'egalit\'e $-[A] = [A^{opp}]$.
Le groupe de Brauer d'un corps est toujours de torsion.
En fait, nous verrons que $[A]$ est d'ordre divisant $\deg(A)$
pour toute alg\`ebre centrale simple de dimension finie $A$ (voir le
lemme \ref{lemma-annihilated-by-degree}).
En g\'en\'eral, le groupe de Brauer n'est pas de type fini; par exemple,
le groupe de Brauer d'un corps local non archim\'edien est $\mathbf{Q}/\mathbf{Z}$.
Le groupe de Brauer de $\mathbf{C}(x, y)$ n'est pas d\'enombrable.

\begin{lemma}
\label{lemma-central-simple-algebra-pgln}
\begin{slogan}
Les alg\`ebres centrales simples sont class\'ees par la cohomologie galoisienne de PGL.
\end{slogan}
Soient $K$ un corps et $K^{sep}$ une cl\^oture s\'eparable.
Alors l'ensemble des classes d'isomorphisme d'alg\`ebres centrales simples de degr\'e
$d$ sur $K$ est en bijection avec la cohomologie non ab\'elienne
$H^1(\text{Gal}(K^{sep}/K), \text{PGL}_d(K^{sep}))$.
\end{lemma}

\begin{proof}[Esquisse de preuve.]
Le th\'eor\`eme de Skolem-Noether (voir
Groupes de Brauer, th\'eor\`eme \ref{brauer-theorem-skolem-noether})
entra\^ine que, pour tout corps $L$, le groupe
$\text{Aut}_{L\text{-Algebras}}(\text{Mat}_d(L))$
est \`egal \`a $\text{PGL}_d(L)$. D'apr\`es le
th\'eor\`eme \ref{theorem-central-simple-algebra}, on voit que les
alg\`ebres centrales simples de degr\'e $d$ correspondent
aux formes de la $K$-alg\`ebre $\text{Mat}_d(K)$.
En combinant ces faits, on voit que les classes d'isomorphisme d'alg\`ebres
centrales simples de degr\'e $d$ correspondent aux \'el\'ements de
$H^1(\text{Gal}(K^{sep}/K), \text{PGL}_d(K^{sep}))$.
Pour plus de d\'etails sur les formes tordues, voir par exemple
\cite{SilvermanEllipticCurves}.
\end{proof}

\noindent
Si $A$ est une alg\`ebre centrale simple de dimension finie et de degr\'e $d$ sur un corps $K$,
nous notons $\xi_A$ la classe de cohomologie correspondante dans
$H^1(\text{Gal}(K^{sep}/K), \text{PGL}_d(K^{sep}))$.
Consid\'erons la suite exacte courte
$$
1 \to (K^{sep})^* \to \text{GL}_d(K^{sep}) \to \text{PGL}_d(K^{sep}) \to 1,
$$
qui donne naissance \`a une suite exacte longue de cohomologie (jusqu'au degr\'e $2$), avec
pour homomorphisme bord
$$
\delta_d :
H ^1(\text{Gal}(K^{sep}/K), \text{PGL}_d(K^{sep}))
\longrightarrow
H^2(\text{Gal}(K^{sep}/K), (K^{sep})^*).
$$
Plus pr\'ecis\'ement, il est d\'efini comme suit : si $\xi$ est une classe de cohomologie
repr\'esent\'ee par le $1$-cocycle $(g_\sigma)$, alors $\delta_d(\xi)$ est la
classe du $2$-cocycle
\begin{equation}
\label{equation-two-cocycle}
(\sigma, \tau)
\longmapsto
\tilde g_\sigma^{-1} \tilde g_{\sigma \tau} \sigma(\tilde g_\tau^{-1})
\in (K^{sep})^*
\end{equation}
o\`u $\tilde g_\sigma \in \text{GL}_d(K^{sep})$ est un rel\`evement de $g_\sigma$.
Cela permet d'expliciter l'application
$$
\delta : \text{Br}(K) \longrightarrow H^2(\text{Gal}(K^{sep}/K), (K^{sep})^*),
\quad
[A] \longmapsto \delta_{\deg A} (\xi_A)
$$
comme suit. Supposons que $A$ soit de degr\'e $d$ sur $K$. Choisissons un isomorphisme
$\varphi : \text{Mat}_d(K^{sep}) \to A \otimes_K K^{sep}$. Pour
$\sigma \in \text{Gal}(K^{sep}/K)$, choisissons un \'el\'ement
$\tilde g_\sigma \in \text{GL}_d(K^{sep})$ tel que
$\varphi^{-1} \circ \sigma(\varphi)$ soit \'egal \`a l'application
$x \mapsto \tilde g_\sigma x \tilde g_\sigma^{-1}$. La classe dans $H^2$
est d\'efinie par le $2$-cocycle (\ref{equation-two-cocycle}).

\begin{theorem}
\label{theorem-brauer-delta}
Soit $K$ un corps de cl\^oture s\'eparable $K^{sep}$. L'application
$\delta : \text{Br}(K) \to H^2(\text{Gal}(K^{sep}/K), (K^{sep})^*)$
d\'efinie ci-dessus est un isomorphisme de groupes.
\end{theorem}

\begin{proof}[Esquisse de preuve]
Pour montrer que $\delta$ d\'efinit un homomorphisme de groupes, c'est-\`a-dire que
$\delta(A \otimes_K B) = \delta(A) + \delta(B)$, on calcule
directement sur les cocycles.

\medskip\noindent
Injectivit\'e de $\delta$. Dans le cas ab\'elien ($d = 1$), on a
l'identification
$$
H^1(\text{Gal}(K^{sep}/K), \text{GL}_d(K^{sep})) =
H_\etale^1(\Spec(K), \text{GL}_d(\mathcal{O}))
$$
dont le second membre est trivial par descente fpqc. Si cela \'etait vrai dans le
cas non ab\'elien, l'injectivit\'e de $\delta$ en r\'esulterait aussit\^ot. (Voir
\cite{SGA4.5}.) En fait, pour la d\'emontrer, on peut r\'einterpr\'eter $\delta([A])$ comme
l'obstruction \`a l'existence d'un $K$-espace vectoriel $V$ muni d'une structure de $A$-module
\`a gauche et tel que $\dim_K V = \deg A$. Lorsque $V$ existe, on
a $A \cong \text{End}_K(V)$.

\medskip\noindent
Pour la surjectivit\'e, choisissons une
classe de cohomologie $\xi \in H^2(\text{Gal}(K^{sep}/K), (K^{sep})^*)$ ;
il existe alors une extension galoisienne finie $K^{sep}/K'/K$
telle que $\xi$ soit l'image d'une classe
$\xi' \in H^2(\text{Gal}(K'|K), (K')^*)$. On construit alors
explicitement une alg\`ebre centrale simple sur $K$ \`a partir des donn\'ees $K', \xi'$.
\end{proof}










\section{Le groupe de Brauer d'un sch\'ema}
\label{section-brauer-scheme}

\noindent
Soit $S$ un sch\'ema. Une $\mathcal{O}_S$-alg\`ebre
$\mathcal{A}$ est dite {\it d'Azumaya} si elle est, localement pour la topologie \'etale, une alg\`ebre
de matrices, c'est-\`a-dire s'il existe un recouvrement \'etale
$\mathcal{U} = \{ \varphi_i : U_i \to S\}_{i \in I}$ tel que
$\varphi_i^*\mathcal{A} \cong \text{Mat}_{d_i}(\mathcal{O}_{U_i})$
pour un $d_i \geq 1$. Deux telles alg\`ebres
$\mathcal{A}$ et $\mathcal{B}$ sont dites {\it \'equivalentes} s'il existe
des $\mathcal{O}_S$-modules localement libres de type fini $\mathcal{F}$ et $\mathcal{G}$
de rang strictement positif en tout $s \in S$, tels que
$$
\mathcal{A} \otimes_{\mathcal{O}_S}
\SheafHom_{\mathcal{O}_S}(\mathcal{F}, \mathcal{F})
\cong
\mathcal{B} \otimes_{\mathcal{O}_S}
\SheafHom_{\mathcal{O}_S}(\mathcal{G}, \mathcal{G})
$$
comme $\mathcal{O}_S$-alg\`ebres. Le {\it groupe de Brauer} de
$S$ est l'ensemble $\text{Br}(S)$ des classes d'\'equivalence d'alg\`ebres
d'Azumaya sur $\mathcal{O}_S$, muni de l'op\'eration induite par le produit tensoriel (sur
$\mathcal{O}_S$).

\begin{lemma}
\label{lemma-end-unique-up-to-invertible}
Soit $S$ un sch\'ema. Soient $\mathcal{F}$ et $\mathcal{G}$ des faisceaux de
$\mathcal{O}_S$-modules localement libres de type fini et de rang strictement positif. S'il
existe un isomorphisme
$\SheafHom_{\mathcal{O}_S}(\mathcal{F}, \mathcal{F}) \cong
\SheafHom_{\mathcal{O}_S}(\mathcal{G}, \mathcal{G})$
de $\mathcal{O}_S$-alg\`ebres, alors il existe un faisceau inversible
$\mathcal{L}$ sur $S$ tel que
$\mathcal{F} \otimes_{\mathcal{O}_S} \mathcal{L} \cong \mathcal{G}$
et tel que cet isomorphisme induise l'isomorphisme donn\'e des
alg\`ebres d'endomorphismes.
\end{lemma}

\begin{proof}
Fixons un isomorphisme
$\SheafHom_{\mathcal{O}_S}(\mathcal{F}, \mathcal{F}) \to
\SheafHom_{\mathcal{O}_S}(\mathcal{G}, \mathcal{G})$.
Consid\'erons le faisceau $\mathcal{L} \subset \SheafHom(\mathcal{F}, \mathcal{G})$
engendr\'e, comme $\mathcal{O}_S$-module, par les isomorphismes locaux
$\varphi : \mathcal{F} \to \mathcal{G}$ tels que la conjugaison par
$\varphi$ soit l'isomorphisme donn\'e des alg\`ebres d'endomorphismes.
Un calcul local (qui se ram\`ene au cas o\`u $\mathcal{F}$ et $\mathcal{G}$
sont libres de type fini et $S$ est affine) montre que $\mathcal{L}$ est inversible.
Un autre calcul local montre que le morphisme d'\'evaluation
$$
\mathcal{F} \otimes_{\mathcal{O}_S} \mathcal{L} \longrightarrow \mathcal{G}
$$
est un isomorphisme.
\end{proof}

\noindent
L'argument donn\'e dans la preuve du lemme suivant se trouve dans
\cite{Saltman-torsion}.

\begin{lemma}
\label{lemma-annihilated-by-degree}
\begin{reference}
Argument tir\'e de \cite{Saltman-torsion}.
\end{reference}
Soit $S$ un sch\'ema. Soit $\mathcal{A}$ une alg\`ebre d'Azumaya qui est
localement libre de rang $d^2$ sur $S$. Alors la classe
de $\mathcal{A}$ dans le groupe de Brauer de $S$ est annul\'ee par $d$.
\end{lemma}

\begin{proof}
Choisissons un recouvrement \'etale $\{U_i \to S\}$ et des isomorphismes
$\mathcal{A}|_{U_i} \to \SheafHom(\mathcal{F}_i, \mathcal{F}_i)$
pour des $\mathcal{O}_{U_i}$-modules localement libres $\mathcal{F}_i$
de rang $d$. (On peut supposer $\mathcal{F}_i$ libre.) Consid\'erons la
compos\'ee
$$
p_i : \mathcal{F}_i^{\otimes d} \to
\wedge^d(\mathcal{F}_i) \to \mathcal{F}_i^{\otimes d}
$$
La premi\`ere fl\`eche est la projection usuelle et la seconde est
l'isomorphisme de la puissance ext\'erieure maximale de $\mathcal{F}_i$ sur
le sous-module des sections de $\mathcal{F}_i^{\otimes d}$ qui se transforment
suivant le caract\`ere signature sous l'action du groupe sym\'etrique
sur $d$ lettres. Alors $p_i^2 = d! p_i$ et le rang de $p_i$ vaut $1$.
En utilisant l'isomorphisme donn\'e
$\mathcal{A}|_{U_i} \to \SheafHom(\mathcal{F}_i, \mathcal{F}_i)$
et l'isomorphisme canonique
$$
\SheafHom(\mathcal{F}_i, \mathcal{F}_i)^{\otimes d} =
\SheafHom(\mathcal{F}_i^{\otimes d}, \mathcal{F}_i^{\otimes d})
$$
nous pouvons consid\'erer $p_i$ comme une section de $\mathcal{A}^{\otimes d}$
sur $U_i$. Nous affirmons que $p_i|_{U_i \times_S U_j} = p_j|_{U_i \times_S U_j}$
comme sections de $\mathcal{A}^{\otimes d}$. En effet, en appliquant le
lemme \ref{lemma-end-unique-up-to-invertible},
nous obtenons un faisceau inversible $\mathcal{L}_{ij}$ et un isomorphisme canonique
$$
\mathcal{F}_i|_{U_i \times_S U_j} \otimes \mathcal{L}_{ij}
\longrightarrow
\mathcal{F}_j|_{U_i \times_S U_j}.
$$
Cet isomorphisme montre que $p_i$ s'envoie sur $p_j$.
Comme $\mathcal{A}^{\otimes d}$ est un faisceau sur $S_\etale$
(proposition \ref{proposition-quasi-coherent-sheaf-fpqc}), on obtient une
section globale canonique $p \in \Gamma(S, \mathcal{A}^{\otimes d})$. Un calcul local
montre que
$$
\mathcal{H} =
\Im(\mathcal{A}^{\otimes d} \to \mathcal{A}^{\otimes d}, f \mapsto fp)
$$
est un module localement libre de rang $d^d$ et que la multiplication (\`a gauche)
par $\mathcal{A}^{\otimes d}$ induit un isomorphisme
$\mathcal{A}^{\otimes d} \to \SheafHom(\mathcal{H}, \mathcal{H})$.
Autrement dit, $\mathcal{A}^{\otimes d}$ est l'\'el\'ement neutre
du groupe de Brauer de $S$, comme voulu.
\end{proof}

\noindent
Dans ce cadre, l'analogue de l'isomorphisme $\delta$ du
th\'eor\`eme \ref{theorem-brauer-delta}
est une application
$$
\delta_S: \text{Br}(S) \to H_\etale^2(S, \mathbf{G}_m).
$$
L'application $\delta_S$ est injective. Si $S$ est quasi-compact ou
connexe, alors $\text{Br}(S)$ est un groupe de torsion; dans ce cas, l'image
de $\delta_S$ est contenue dans le {\it groupe de Brauer cohomologique} de $S$
$$
\text{Br}'(S) := H_\etale^2(S, \mathbf{G}_m)_\text{torsion}.
$$
Ainsi, si $S$ est quasi-compact ou connexe, on a une inclusion $\text{Br}(S)
\subset \text{Br}'(S)$. Ce n'est pas toujours une \'egalit\'e : il existe une
surface singuli\`ere non s\'epar\'ee $S$ pour laquelle $\text{Br}(S) \subset
\text{Br}'(S)$ est une inclusion stricte. Si $S$ est quasi-projectif, alors
$\text{Br}(S) = \text{Br}'(S)$. Cependant, on ignore si cela reste vrai pour
une vari\'et\'e lisse propre sur $\mathbf{C}$, par exemple.




\section{La suite d'Artin-Schreier}
\label{section-artin-schreier}

\noindent
Soit $p$ un nombre premier. Soit $S$ un sch\'ema de caract\'eristique $p$.
La suite {\it d'Artin-Schreier} est la suite exacte courte
$$
0 \longrightarrow \underline{\mathbf{Z}/p\mathbf{Z}}_S \longrightarrow
\mathbf{G}_{a, S} \xrightarrow{F-1} \mathbf{G}_{a, S} \longrightarrow 0
$$
o\`u $F - 1$ est l'application $x \mapsto x^p - x$.

\begin{lemma}
\label{lemma-vanishing-affine-char-p-p}
Soit $p$ un nombre premier. Soit $S$ un sch\'ema de caract\'eristique $p$.
\begin{enumerate}
\item Si $S$ est affine, alors
$H_\etale^q(S, \underline{\mathbf{Z}/p\mathbf{Z}}) = 0$ pour tout
$q \geq 2$.
\item Si $S$ est un sch\'ema quasi-compact et quasi-s\'epar\'e de
dimension $d$, alors $H_\etale^q(S, \underline{\mathbf{Z}/p\mathbf{Z}}) = 0$
pour tout $q \geq 2 + d$.
\end{enumerate}
\end{lemma}

\begin{proof}
Rappelons que la cohomologie \'etale du faisceau structural est \'egale
\`a sa cohomologie sur l'espace topologique sous-jacent
(Theorem \ref{theorem-zariski-fpqc-quasi-coherent}).
La premi\`ere assertion r\'esulte de la suite exacte d'Artin-Schreier
et de l'annulation de la cohomologie du faisceau structural sur un sch\'ema
affine (Cohomology of Schemes, Lemma
\ref{coherent-lemma-quasi-coherent-affine-cohomology-zero}).
La seconde assertion r\'esulte par le m\^eme argument de l'annulation
donn\'ee par Cohomology, Proposition
\ref{cohomology-proposition-cohomological-dimension-spectral}
et du fait que $S$ est un espace spectral
(Properties, Lemma
\ref{properties-lemma-quasi-compact-quasi-separated-spectral}).
\end{proof}

\begin{lemma}
\label{lemma-F-1}
Soit $k$ un corps alg\'ebriquement clos de caract\'eristique $p > 0$.
Soit $V$ un $k$-espace vectoriel de dimension finie. Soit $F : V \to V$
une application semi-lin\'eaire relativement au Frobenius, c'est-\`a-dire une application additive telle que
$F(\lambda v) = \lambda^p F(v)$ pour tous $\lambda \in k$ et $v \in V$.
Alors $F - 1 : V \to V$ est surjective, et son noyau est un
$\mathbf{F}_p$-espace vectoriel de dimension finie, celle-ci \'etant $\leq \dim_k(V)$.
\end{lemma}

\begin{proof}
Si $F = 0$, l'\'enonc\'e est vrai. Si $V$ est muni d'une filtration par des
sous-espaces vectoriels stables par $F$ telle que l'\'enonc\'e soit vrai pour chaque
quotient gradu\'e, alors il est vrai pour $(V, F)$. En combinant ces deux remarques,
nous pouvons supposer que le noyau de $F$ est nul.

\medskip\noindent
Choisissons une base $v_1, \ldots, v_n$ de $V$ et \'ecrivons
$F(v_i) = \sum a_{ij} v_j$. Remarquons que $v = \sum \lambda_i v_i$
appartient au noyau si et seulement si $\sum \lambda_i^p a_{ij} v_j = 0$.
Comme $k$ est alg\'ebriquement clos, cela implique que la matrice $(a_{ij})$
est inversible. Notons $(b_{ij})$ son inverse. Pour voir que $F - 1$
est surjective, choisissons $w = \sum \mu_i v_i \in V$ et cherchons \`a r\'esoudre
$$
(F - 1)(\sum \lambda_iv_i) =
\sum \lambda_i^p a_{ij} v_j - \sum \lambda_j v_j = \sum \mu_j v_j
$$
Cela \'equivaut \`a
$$
\sum \lambda_j^p v_j - \sum b_{ij} \lambda_i v_j = \sum b_{ij} \mu_i v_j
$$
autrement dit
$$
\lambda_j^p - \sum b_{ij} \lambda_i = \sum b_{ij} \mu_i,
\quad j = 1, \ldots, \dim(V).
$$
L'alg\`ebre
$$
A = k[x_1, \ldots, x_n]/
(x_j^p - \sum b_{ij} x_i - \sum b_{ij} \mu_i)
$$
est lisse standard sur $k$
(Algebra, Definition \ref{algebra-definition-standard-smooth})
car la matrice $(b_{ij})$ est inversible et les d\'eriv\'ees partielles
des $x_j^p$ sont nulles. Une base de $A$ sur $k$ est l'ensemble des mon\^omes
$x_1^{e_1} \ldots x_n^{e_n}$ avec $e_i < p$, donc $\dim_k(A) = p^n$.
Comme $k$ est alg\'ebriquement clos, on voit que $\Spec(A)$ poss\`ede exactement
$p^n$ points. Il s'ensuit que $F - 1$ est surjective et que chaque fibre
poss\`ede $p^n$ points, c'est-\`a-dire que le noyau de $F - 1$ est un groupe d'ordre $p^n$.
\end{proof}

\begin{lemma}
\label{lemma-F-2}
Dans la situation du lemme \ref{lemma-F-1}, soit $k'/k$ une extension
de corps alg\'ebriquement clos. Posons $V' = V \otimes_k k'$ et soit
$F' : V' \to V'$ l'application semi-lin\'eaire relativement au Frobenius induite par $F$.
Alors $\Ker(F - 1)$ s'envoie isomorphiquement sur $\Ker(F' - 1)$.
\end{lemma}

\begin{proof}
Soit $v_1, \ldots, v_n$ une base de $V$ et \'ecrivons $F(v_i) = \sum a_{ij}v_j$.
Alors $v = \sum \lambda_i v_i$ appartient \`a $\Ker(F - 1)$ si et seulement si
$\sum \lambda_i^pa_{ij} = \lambda_j$ pour $j = 1, \ldots, n$.
D'apr\`es le lemme \ref{lemma-F-1}, ces \'equations
d\'efinissent un sous-sch\'ema ferm\'e $Z$ de
$\mathbf{A}^n_k = \Spec(k[\lambda_1, \ldots, \lambda_n])$
fini sur $\Spec(k)$. Alors $Z(k) = Z(k')$, ce qui permet de conclure.
\end{proof}

\begin{lemma}
\label{lemma-top-cohomology-coherent}
Soit $X$ un sch\'ema s\'epar\'e de type fini sur un corps $k$.
Soit $\mathcal{F}$ un faisceau coh\'erent de $\mathcal{O}_X$-modules.
Alors $\dim_k H^d(X, \mathcal{F}) < \infty$, o\`u $d = \dim(X)$.
\end{lemma}

\begin{proof}
Nous allons le d\'emontrer par r\'ecurrence sur $d$. Le cas $d = 0$ est vrai car
dans ce cas $X$ est le spectre d'une $k$-alg\`ebre $A$ de dimension finie
(Varieties, Lemma \ref{varieties-lemma-algebraic-scheme-dim-0})
et tout faisceau coh\'erent $\mathcal{F}$ correspond \`a un $A$-module de type fini
$M = H^0(X, \mathcal{F})$ qui v\'erifie $\dim_k M < \infty$.

\medskip\noindent
Supposons $d > 0$ et le r\'esultat d\'emontr\'e pour les sch\'emas s\'epar\'es
de type fini de dimension $< d$. Le sch\'ema $X$ est noeth\'erien. Consid\'erons
la propri\'et\'e $\mathcal{P}$ des faisceaux coh\'erents sur $X$ d\'efinie par la r\`egle
$$
\mathcal{P}(\mathcal{F}) \Leftrightarrow
\dim_k H^d(X, \mathcal{F}) < \infty
$$
Nous allons utiliser le r\'esultat de
Cohomology of Schemes, Lemma \ref{coherent-lemma-property-initial}
pour montrer que $\mathcal{P}$ est satisfaite par tout faisceau coh\'erent sur $X$.

\medskip\noindent
Soit
$$
0 \to \mathcal{F}_1 \to \mathcal{F} \to \mathcal{F}_2 \to 0
$$
une suite exacte courte de faisceaux coh\'erents sur $X$.
Consid\'erons la suite exacte longue de cohomologie
$$
H^d(X, \mathcal{F}_1) \to
H^d(X, \mathcal{F}) \to
H^d(X, \mathcal{F}_2)
$$
Ainsi, si $\mathcal{P}$ est satisfaite par $\mathcal{F}_1$ et $\mathcal{F}_2$,
elle l'est par $\mathcal{F}$.

\medskip\noindent
Soit $Z \subset X$ un sous-sch\'ema ferm\'e int\`egre. Soit $\mathcal{I}$
un faisceau coh\'erent d'id\'eaux sur $Z$. Pour achever la d\'emonstration, il faut montrer
que $H^d(X, i_*\mathcal{I}) = H^d(Z, \mathcal{I})$ est de dimension finie.
Si $\dim(Z) < d$, le r\'esultat est vrai car ce groupe de cohomologie
est nul (Cohomology, Proposition
\ref{cohomology-proposition-vanishing-Noetherian}).
Nous sommes ainsi ramen\'es \`a la situation examin\'ee au paragraphe suivant.

\medskip\noindent
Supposons que $X$ soit une vari\'et\'e de dimension $d$ et que
$\mathcal{F} = \mathcal{I}$ soit un faisceau coh\'erent d'id\'eaux. Dans ce
cas, nous avons une suite exacte courte
$$
0 \to \mathcal{I} \to \mathcal{O}_X \to i_*\mathcal{O}_Z \to 0
$$
o\`u $i : Z \to X$ est l'immersion ferm\'ee d\'efinie par $\mathcal{I}$.
L'hypoth\`ese de r\'ecurrence montre que
$H^{d - 1}(Z, \mathcal{O}_Z) = H^{d - 1}(X, i_*\mathcal{O}_Z)$ est
de dimension finie. Il suffit donc de d\'emontrer le r\'esultat
pour le faisceau structural.

\medskip\noindent
Nous pouvons appliquer le lemme de Chow
(Cohomology of Schemes, Lemma \ref{coherent-lemma-chow-Noetherian})
au morphisme $X \to \Spec(k)$. Nous obtenons ainsi un diagramme
$$
\xymatrix{
X \ar[rd]_g & X' \ar[d]^-{g'} \ar[l]^\pi \ar[r]_i & \mathbf{P}^n_k \ar[dl] \\
& \Spec(k) &
}
$$
comme dans l'\'enonc\'e du lemme de Chow. Soit en outre $U \subset X$
le sous-sch\'ema ouvert dense tel que $\pi^{-1}(U) \to U$ soit un isomorphisme.
Nous pouvons aussi supposer que $X'$ est une vari\'et\'e, voir
Cohomology of Schemes, Remark \ref{coherent-remark-chow-Noetherian}.
Le morphisme $i' = (i, \pi) : X' \to \mathbf{P}^n_X$ est
une immersion ferm\'ee (loc. cit.). Par cons\'equent
$$
\mathcal{L} = i^*\mathcal{O}_{\mathbf{P}^n_k}(1) \cong
(i')^*\mathcal{O}_{\mathbf{P}^n_X}(1)
$$
est ample relativement \`a $\pi$ (par exemple d'apr\`es
Morphisms, Lemma \ref{morphisms-lemma-characterize-ample-on-finite-type}).
Ainsi, d'apr\`es Cohomology of Schemes, Lemma \ref{coherent-lemma-kill-by-twisting},
il existe un $n \geq 0$ tel que
$R^p\pi_*\mathcal{L}^{\otimes n} = 0$ pour tout $p > 0$.
Posons $\mathcal{G} = \pi_*\mathcal{L}^{\otimes n}$.
Choisissons une section globale non nulle $s$ de $\mathcal{L}^{\otimes n}$.
Puisque $\mathcal{G} = \pi_*\mathcal{L}^{\otimes n}$, la section $s$
correspond \`a une section de $\mathcal{G}$, c'est-\`a-dire \`a un morphisme
$\mathcal{O}_X \to \mathcal{G}$.
Comme $s|_U \not = 0$, puisque $X'$ est une vari\'et\'e et $\mathcal{L}$ est
inversible, on voit que $\mathcal{O}_X|_U \to \mathcal{G}|_U$
est non nul. Comme $\mathcal{G}|_U = \mathcal{L}^{\otimes n}|_{\pi^{-1}(U)}$
est inversible, nous obtenons une suite exacte courte
$$
0 \to \mathcal{O}_X \to \mathcal{G} \to \mathcal{Q} \to 0
$$
o\`u $\mathcal{Q}$ est coh\'erent et \`a support dans un sous-sch\'ema
ferm\'e strict de $X$. En raisonnant comme ci-dessus avec notre hypoth\`ese de
r\'ecurrence, on voit qu'il
suffit de d\'emontrer $\dim H^d(X, \mathcal{G}) < \infty$.

\medskip\noindent
La suite spectrale de Leray
(Cohomology, Lemma \ref{cohomology-lemma-apply-Leray})
montre que $H^d(X, \mathcal{G}) = H^d(X', \mathcal{L}^{\otimes n})$.
Soit $\overline{X}' \subset \mathbf{P}^n_k$ l'adh\'erence
de $X'$. Alors $\overline{X}'$ est une vari\'et\'e projective de dimension $d$
sur $k$ et $X' \subset \overline{X}'$ est un ouvert dense.
Le faisceau inversible $\mathcal{L}^{\otimes n}$ est la restriction de
$\mathcal{O}_{\overline{X}'}(n)$ \`a $X'$. D'apr\`es
Cohomology, Proposition
\ref{cohomology-proposition-cohomological-dimension-spectral}
l'application
$$
H^d(\overline{X}', \mathcal{O}_{\overline{X}'}(n))
\longrightarrow
H^d(X', \mathcal{L}^{\otimes n})
$$
est surjective. Comme le groupe de cohomologie de gauche est
de dimension finie d'apr\`es
Cohomology of Schemes, Lemma \ref{coherent-lemma-coherent-projective}
la d\'emonstration est achev\'ee.
\end{proof}

\begin{lemma}
\label{lemma-vanishing-variety-char-p-p}
Soit $X$ un sch\'ema s\'epar\'e de type fini sur un corps alg\'ebriquement clos
$k$ de caract\'eristique $p > 0$. Alors
$H_\etale^q(X, \underline{\mathbf{Z}/p\mathbf{Z}}) = 0$ pour
$q \geq dim(X) + 1$.
\end{lemma}

\begin{proof}
Posons $d = \dim(X)$. D'apr\`es l'annulation \'etablie dans le
lemme \ref{lemma-vanishing-affine-char-p-p},
il suffit de montrer que
$H_\etale^{d + 1}(X, \underline{\mathbf{Z}/p\mathbf{Z}}) = 0$.
D'apr\`es le lemme \ref{lemma-top-cohomology-coherent}, on voit que
$H^d(X, \mathcal{O}_X)$ est un espace vectoriel de dimension finie sur $k$.
Par cons\'equent, la suite exacte longue de cohomologie associ\'ee \`a la
suite d'Artin-Schreier se termine par
$$
H^d(X, \mathcal{O}_X) \xrightarrow{F - 1}
H^d(X, \mathcal{O}_X) \to H^{d + 1}_\etale(X, \mathbf{Z}/p\mathbf{Z}) \to 0
$$
D'apr\`es le lemme \ref{lemma-F-1}, l'application $F - 1$ de cette suite est surjective.
Cela d\'emontre le lemme.
\end{proof}

\begin{lemma}
\label{lemma-finiteness-proper-variety-char-p-p}
Soit $X$ un sch\'ema propre sur un corps alg\'ebriquement clos
$k$ de caract\'eristique $p > 0$. Alors
\begin{enumerate}
\item $H_\etale^q(X, \underline{\mathbf{Z}/p\mathbf{Z}})$
est un $\mathbf{Z}/p\mathbf{Z}$-module de type fini pour tout $q$, et
\item $H^q_\etale(X, \underline{\mathbf{Z}/p\mathbf{Z}}) \to
H^q_\etale(X_{k'}, \underline{\mathbf{Z}/p\mathbf{Z}})$
est un isomorphisme si $k'/k$ est une extension de corps alg\'ebriquement
clos.
\end{enumerate}
\end{lemma}

\begin{proof}
D'apr\`es Cohomology of Schemes, Lemma
\ref{coherent-lemma-proper-over-affine-cohomology-finite}
et la comparaison de la cohomologie donn\'ee par le
th\'eor\`eme \ref{theorem-zariski-fpqc-quasi-coherent}, les groupes de cohomologie
$H^q_\etale(X, \mathbf{G}_a) = H^q(X, \mathcal{O}_X)$ sont
des espaces vectoriels de dimension finie sur $k$. Par cons\'equent, d'apr\`es le
lemme \ref{lemma-F-1}, la suite exacte longue de cohomologie
associ\'ee \`a la suite d'Artin-Schreier se d\'ecompose en
suites exactes courtes
$$
0 \to H_\etale^q(X, \underline{\mathbf{Z}/p\mathbf{Z}}) \to
H^q(X, \mathcal{O}_X) \xrightarrow{F - 1} H^q(X, \mathcal{O}_X) \to 0
$$
et, de plus, la dimension sur $\mathbf{F}_p$ des groupes de cohomologie
$H_\etale^q(X, \underline{\mathbf{Z}/p\mathbf{Z}})$ est inf\'erieure ou \'egale \`a la
dimension sur $k$ de l'espace vectoriel $H^q(X, \mathcal{O}_X)$.
Cela d\'emontre la premi\`ere assertion. La seconde r\'esulte
du lemme \ref{lemma-F-2}
car $H^q(X, \mathcal{O}_X) \otimes_k k' \to H^q(X_{k'}, \mathcal{O}_{X_{k'}})$
est un isomorphisme par changement de base plat
(Cohomology of Schemes,
Lemma \ref{coherent-lemma-flat-base-change-cohomology}).
\end{proof}








\section{Faisceaux localement constants}
\label{section-locally-constant}

\noindent
Cette section est l'analogue de
Modules on Sites, Section \ref{sites-modules-section-locally-constant}
pour le site \'etale.

\begin{definition}
\label{definition-finite-locally-constant}
Soit $X$ un sch\'ema.
Soit $\mathcal{F}$ un faisceau d'ensembles sur $X_\etale$.
\begin{enumerate}
\item Soit $E$ un ensemble. Nous disons que $\mathcal{F}$ est le
{\it faisceau constant de valeur $E$} si $\mathcal{F}$ est le
faisceau associ\'e au pr\'efaisceau $U \mapsto E$.
Notation : $\underline{E}_X$ ou $\underline{E}$.
\item Nous disons que $\mathcal{F}$ est un {\it faisceau constant} s'il est
isomorphe \`a un faisceau comme en (1).
\item Nous disons que $\mathcal{F}$ est {\it localement constant} s'il existe un
recouvrement $\{U_i \to X\}$ tel que $\mathcal{F}|_{U_i}$ soit un faisceau constant.
\item Nous disons que $\mathcal{F}$ est {\it localement constant fini} s'il
est localement constant et si ses valeurs locales sont des ensembles finis.
\end{enumerate}
Soit $\mathcal{F}$ un faisceau de groupes ab\'eliens sur $X_\etale$.
\begin{enumerate}
\item Soit $A$ un groupe ab\'elien.
Nous disons que $\mathcal{F}$ est le {\it faisceau constant de valeur $A$} si
$\mathcal{F}$ est le faisceau associ\'e au pr\'efaisceau $U \mapsto A$.
Notation : $\underline{A}_X$ ou $\underline{A}$.
\item Nous disons que $\mathcal{F}$ est un {\it faisceau constant} s'il est isomorphe
comme faisceau ab\'elien \`a un faisceau comme en (1).
\item Nous disons que $\mathcal{F}$ est {\it localement constant} s'il existe un
recouvrement $\{U_i \to X\}$ tel que $\mathcal{F}|_{U_i}$ soit un faisceau constant.
\item Nous disons que $\mathcal{F}$ est {\it localement constant fini} s'il
est localement constant et si ses valeurs locales sont des groupes ab\'eliens finis.
\end{enumerate}
Soit $\Lambda$ un anneau. Soit $\mathcal{F}$ un faisceau de $\Lambda$-modules
sur $X_\etale$.
\begin{enumerate}
\item Soit $M$ un $\Lambda$-module.
Nous disons que $\mathcal{F}$ est le {\it faisceau constant de valeur $M$} si
$\mathcal{F}$ est le faisceau associ\'e au pr\'efaisceau $U \mapsto M$.
Notation : $\underline{M}_X$ ou $\underline{M}$.
\item Nous disons que $\mathcal{F}$ est un {\it faisceau constant} s'il est isomorphe
comme faisceau de $\Lambda$-modules \`a un faisceau comme en (1).
\item Nous disons que $\mathcal{F}$ est {\it localement constant} s'il existe un
recouvrement $\{U_i \to X\}$ tel que $\mathcal{F}|_{U_i}$ soit un faisceau constant.
\end{enumerate}
\end{definition}

\begin{lemma}
\label{lemma-pullback-locally-constant}
Soit $f : X \to Y$ un morphisme de sch\'emas. Si $\mathcal{G}$ est un
faisceau localement constant d'ensembles, de groupes ab\'eliens ou de $\Lambda$-modules
sur $Y_\etale$, alors $f^{-1}\mathcal{G}$ l'est aussi
sur $X_\etale$.
\end{lemma}

\begin{proof}
Cela vaut pour tout morphisme de topos, voir
Modules on Sites, Lemma \ref{sites-modules-lemma-pullback-locally-constant}.
\end{proof}

\begin{lemma}
\label{lemma-pushforward-locally-constant}
Soit $f : X \to Y$ un morphisme fini \'etale de sch\'emas.
Si $\mathcal{F}$ est un faisceau d'ensembles localement constant (fini),
un faisceau de groupes ab\'eliens localement constant (fini), ou
un faisceau de $\Lambda$-modules localement constant (de type fini)
sur $X_\etale$, alors $f_*\mathcal{F}$ l'est aussi
sur $Y_\etale$.
\end{lemma}

\begin{proof}
La construction de $f_*$ commute \`a la localisation \'etale.
Un morphisme fini \'etale est localement isomorphe \`a une somme disjointe
d'isomorphismes, voir
\'Etale Morphisms, Lemma \ref{etale-lemma-finite-etale-etale-local}.
Ainsi le lemme affirme que, si $\mathcal{F}_i$, $i = 1, \ldots, n$
sont des faisceaux d'ensembles localement constants (finis), alors
$\prod_{i = 1, \ldots, n} \mathcal{F}_i$ l'est aussi.
C'est clair. Il en va de m\^eme pour les faisceaux de groupes ab\'eliens et de modules.
\end{proof}

\begin{lemma}
\label{lemma-characterize-finite-locally-constant}
Soient $X$ un sch\'ema et $\mathcal{F}$ un faisceau d'ensembles sur $X_\etale$.
Alors les assertions suivantes sont \`equivalentes
\begin{enumerate}
\item $\mathcal{F}$ est localement constant fini, et
\item $\mathcal{F} = h_U$ pour un morphisme fini \'etale $U \to X$.
\end{enumerate}
\end{lemma}

\begin{proof}
Un morphisme fini \'etale est localement isomorphe \`a une somme disjointe
d'isomorphismes, voir
\'Etale Morphisms, Lemma \ref{etale-lemma-finite-etale-etale-local}.
Ainsi (2) implique (1). R\'eciproquement, si $\mathcal{F}$ est localement
constant fini, il existe un recouvrement \'etale $\{X_i \to X\}$ tel que
$\mathcal{F}|_{X_i}$ soit repr\'esentable par un morphisme fini \'etale $U_i \to X_i$.
En raisonnant exactement comme dans la d\'emonstration du
Descent, Lemma \ref{descent-lemma-descent-data-sheaves}
nous obtenons une donn\'ee de descente de sch\'emas $(U_i, \varphi_{ij})$ relativement \`a
$\{X_i \to X\}$ (d\'etails omis). Cette donn\'ee de descente est effective, par
exemple d'apr\`es Descent, Lemma \ref{descent-lemma-affine},
et le morphisme de sch\'emas obtenu $U \to X$ est fini \'etale
d'apr\`es Descent, Lemmas \ref{descent-lemma-descending-property-finite} et
\ref{descent-lemma-descending-property-etale}.
\end{proof}

\begin{lemma}
\label{lemma-morphism-locally-constant}
Soit $X$ un sch\'ema.
\begin{enumerate}
\item Soit $\varphi : \mathcal{F} \to \mathcal{G}$ un morphisme
de faisceaux d'ensembles localement constants sur $X_\etale$.
Si $\mathcal{F}$ est localement constant fini, il existe un
recouvrement \'etale $\{U_i \to X\}$ tel que
$\varphi|_{U_i}$ soit le morphisme de faisceaux constants associ\'e \`a
une application d'ensembles.
\item Soit $\varphi : \mathcal{F} \to \mathcal{G}$ un morphisme
de faisceaux de groupes ab\'eliens localement constants sur $X_\etale$.
Si $\mathcal{F}$ est localement constant fini, il existe un recouvrement \'etale
$\{U_i \to X\}$ tel que $\varphi|_{U_i}$ soit le morphisme de
faisceaux ab\'eliens constants associ\'e \`a un homomorphisme de groupes ab\'eliens.
\item Soit $\Lambda$ un anneau.
Soit $\varphi : \mathcal{F} \to \mathcal{G}$ un morphisme
de faisceaux de $\Lambda$-modules localement constants sur $X_\etale$.
Si $\mathcal{F}$ est de type fini, alors il existe un recouvrement \'etale
$\{U_i \to X\}$ tel que $\varphi|_{U_i}$ soit le morphisme de
faisceaux constants de $\Lambda$-modules associ\'e \`a un homomorphisme de $\Lambda$-modules.
\end{enumerate}
\end{lemma}

\begin{proof}
Cela vaut sur tout site, voir
Modules on Sites, Lemma \ref{sites-modules-lemma-morphism-locally-constant}.
\end{proof}

\begin{lemma}
\label{lemma-kernel-finite-locally-constant}
Soit $X$ un sch\'ema.
\begin{enumerate}
\item La cat\'egorie des faisceaux d'ensembles localement constants finis
est stable par limites projectives et inductives finies dans $\Sh(X_\etale)$.
\item La cat\'egorie des faisceaux ab\'eliens localement constants finis est une
sous-cat\'egorie de Serre faible de $\textit{Ab}(X_\etale)$.
\item Soit $\Lambda$ un anneau noeth\'erien. La cat\'egorie des
faisceaux de $\Lambda$-modules localement constants de type fini sur
$X_\etale$ est une sous-cat\'egorie de Serre faible de
$\textit{Mod}(X_\etale, \Lambda)$.
\end{enumerate}
\end{lemma}

\begin{proof}
Cela vaut sur tout site, voir
Modules on Sites, Lemma
\ref{sites-modules-lemma-kernel-finite-locally-constant}.
\end{proof}

\begin{lemma}
\label{lemma-tensor-product-locally-constant}
Soit $X$ un sch\'ema. Soit $\Lambda$ un anneau.
Le produit tensoriel de deux faisceaux de $\Lambda$-modules localement constants
sur $X_\etale$ est un faisceau de $\Lambda$-modules localement constant.
\end{lemma}

\begin{proof}
Cela vaut sur tout site, voir
Modules on Sites, Lemma
\ref{sites-modules-lemma-tensor-product-locally-constant}.
\end{proof}

\begin{lemma}
\label{lemma-connected-locally-constant}
Soit $X$ un sch\'ema connexe. Soit $\Lambda$ un anneau et soit
$\mathcal{F}$ un faisceau de $\Lambda$-modules localement constant.
Alors il existe un $\Lambda$-module $M$ et un recouvrement \'etale
$\{U_i \to X\}$ tels que $\mathcal{F}|_{U_i} \cong \underline{M}|_{U_i}$.
\end{lemma}

\begin{proof}
Choisissons un recouvrement \'etale
$\{U_i \to X\}$ tel que $\mathcal{F}|_{U_i}$ soit constant, disons
$\mathcal{F}|_{U_i} \cong \underline{M_i}_{U_i}$.
Remarquons que $U_i \times_X U_j$ est vide si $M_i$ n'est pas isomorphe
\`a $M_j$.
Pour chaque classe d'isomorphisme de $\Lambda$-modules repr\'esent\'ee par $M$, posons $I_M = \{i \in I \mid M_i \cong M\}$.
Comme les morphismes \'etales sont ouverts, nous voyons que
$U_M = \bigcup_{i \in I_M} \Im(U_i \to X)$
est un ouvert de $X$. Alors $X = \coprod U_M$ est un recouvrement
ouvert disjoint de $X$. Comme $X$ est connexe, un seul $U_M$ est non vide
et le lemme en r\'esulte.
\end{proof}






\section{Faisceaux localement constants et groupe fondamental}
\label{section-pione}

\noindent
Dans certains cas, nous pouvons relier les faisceaux localement constants
au groupe fondamental d'un sch\'ema.

\begin{lemma}
\label{lemma-locally-constant-on-connected}
Soit $X$ un sch\'ema connexe. Soit $\overline{x}$ un point g\'eom\'etrique de $X$.
\begin{enumerate}
\item Il existe une \`equivalence de cat\'egories
$$
\left\{
\begin{matrix}
\text{faisceaux d'ensembles localement constants finis}\\
\text{sur }X_\etale
\end{matrix}
\right\}
\longleftrightarrow
\left\{
\begin{matrix}
\text{ensembles finis munis d'une action de }\pi_1(X, \overline{x})
\end{matrix}
\right\}
$$
\item Il existe une \`equivalence de cat\'egories
$$
\left\{
\begin{matrix}
\text{faisceaux de groupes ab\'eliens localement constants finis}\\
\text{sur }X_\etale
\end{matrix}
\right\}
\longleftrightarrow
\left\{
\begin{matrix}
\pi_1(X, \overline{x})\text{-modules finis}
\end{matrix}
\right\}
$$
\item Soit $\Lambda$ un anneau fini. Il existe une \`equivalence de cat\'egories
$$
\left\{
\begin{matrix}
\text{faisceaux localement constants de type fini}\\
\text{de }\Lambda\text{-modules sur }X_\etale
\end{matrix}
\right\}
\longleftrightarrow
\left\{
\begin{matrix}
\pi_1(X, \overline{x})\text{-modules finis munis}\\
\text{d'une structure commutante de }\Lambda\text{-module}
\end{matrix}
\right\}
$$
\end{enumerate}
\end{lemma}

\begin{proof}
Remarquons que $\pi_1(X, \overline{x})$ est un groupe topologique
profini, voir Fundamental Groups, Definition
\ref{pione-definition-fundamental-group}.
Les cat\'egories des membres de gauche sont d\'efinies dans la
section \ref{section-locally-constant}.
La notation employ\'ee dans les membres de droite vient de
Fundamental Groups, Definition \ref{pione-definition-G-set-continuous}
pour les ensembles et de la
d\'efinition \ref{definition-G-module-continuous} pour les groupes ab\'eliens.
Cela explique la notation.

\medskip\noindent
L'assertion (1) r\'esulte du
lemme \ref{lemma-characterize-finite-locally-constant}
et de Fundamental Groups, Theorem \ref{pione-theorem-fundamental-group}.
Les parties (2) et (3) s'en d\'eduisent imm\'ediatement en munissant les
(faisceaux d')ensembles sous-jacents d'une structure suppl\'ementaire. Par exemple, un
faisceau de groupes ab\'eliens localement constant fini sur $X_\etale$ revient
\`a un faisceau d'ensembles localement constant fini $\mathcal{F}$
muni d'un morphisme $+ : \mathcal{F} \times \mathcal{F} \to \mathcal{F}$
satisfaisant aux axiomes usuels. L'\'equivalence de (1) pr\'eserve les produits
et transforme donc $+$ en une addition sur l'ensemble fini correspondant
muni d'une action de $\pi_1(X, \overline{x})$. Puisque les $\pi_1(X, \overline{x})$-modules
reviennent aux ensembles munis d'une action de $\pi_1(X, \overline{x})$ et d'une
structure compatible de groupe ab\'elien, nous obtenons (2). La partie (3) se d\'emontre
exactement de la m\^eme mani\`ere.
\end{proof}

\begin{lemma}
\label{lemma-locally-constant-on-connected-geom-unibranch}
Soit $X$ un sch\'ema irr\'eductible et g\'eom\'etriquement unibranche.
Soit $\overline{x}$ un point g\'eom\'etrique de $X$.
Soit $\Lambda$ un anneau. Il existe une \`equivalence de cat\'egories
$$
\left\{
\begin{matrix}
\text{faisceaux localement constants de type fini}\\
\text{de }\Lambda\text{-modules sur }X_\etale
\end{matrix}
\right\}
\longleftrightarrow
\left\{
\begin{matrix}
\Lambda\text{-modules de type fini }M\text{ munis}\\
\text{d'une action continue de }\pi_1(X, \overline{x})
\end{matrix}
\right\}
$$
\end{lemma}

\begin{proof}
La d\'emonstration du lemme \ref{lemma-locally-constant-on-connected}
ne s'applique pas, car un $\Lambda$-module de type fini $M$ peut ne pas avoir
un ensemble sous-jacent fini.

\medskip\noindent
Soit $\nu : X^\nu \to X$ le morphisme de normalisation. D'apr\`es Morphisms, Lemma
\ref{morphisms-lemma-normalization-geom-unibranch-univ-homeo}
il s'agit d'un hom\'eomorphisme universel. D'apr\`es
Fundamental Groups, Proposition \ref{pione-proposition-universal-homeomorphism}
il induit un isomorphisme
$\pi_1(X^\nu, \overline{x}) \to \pi_1(X, \overline{x})$
et, d'apr\`es le th\'eor\`eme \ref{theorem-topological-invariance}, nous obtenons
une \`equivalence de cat\'egories entre les faisceaux localement constants de type fini
de $\Lambda$-modules sur $X_\etale$ et sur $X^\nu_\etale$.
Nous sommes ainsi ramen\'es au cas o\`u $X$ est un sch\'ema int\`egre normal.

\medskip\noindent
Supposons que $X$ soit un sch\'ema int\`egre normal. Soit $\eta \in X$ le point g\'en\'erique.
Soit $\overline{\eta}$ un point g\'eom\'etrique au-dessus de $\eta$.
D'apr\`es Fundamental Groups, Proposition \ref{pione-proposition-normal},
nous avons une surjection continue
$$
\text{Gal}(\kappa(\eta)^{sep}/\kappa(\eta)) = \pi_1(\eta, \overline{\eta})
\longrightarrow
\pi_1(X, \overline{\eta})
$$
dont le noyau est d\'ecrit dans Fundamental Groups, Lemma
\ref{pione-lemma-normal-pione-quotient-inertia}.
Soit $\mathcal{F}$ un faisceau localement constant de type fini de $\Lambda$-modules
sur $X_\etale$. Soit $M = \mathcal{F}_{\overline{\eta}}$
la fibre de $\mathcal{F}$ en $\overline{\eta}$.
La section \ref{section-galois-action-stalks} nous donne une action continue de
$\text{Gal}(\kappa(\eta)^{sep}/\kappa(\eta))$ sur $M$.
Notre but est de montrer que cette action se factorise par la
surjection affich\'ee.
Puisque $\mathcal{F}$ est de type fini, $M$ est un $\Lambda$-module de type fini.
Puisque $\mathcal{F}$ est localement constant, pour tout $x \in X$,
la restriction de $\mathcal{F}$ \`a $\Spec(\mathcal{O}_{X, x}^{sh})$
est constante. L'action de $\text{Gal}(K^{sep}/K_x^{sh})$
(avec les notations de Fundamental Groups, Lemma
\ref{pione-lemma-normal-pione-quotient-inertia})
sur $M$ est donc triviale. Nous obtenons la factorisation voulue.

\medskip\noindent
R\'eciproquement, supposons donn\'e un $\Lambda$-module de type fini $M$
muni d'une action continue de $\pi_1(X, \overline{\eta})$.
Nous allons construire un $\mathcal{F}$ tel que
$M \cong \mathcal{F}_{\overline{\eta}}$ comme
$\Lambda[\pi_1(X, \overline{\eta})]$-modules.
Choisissons des g\'en\'erateurs $m_1, \ldots, m_r \in M$.
Comme l'action de $\pi_1(X, \overline{\eta})$ sur $M$ est
continue, il existe pour tout $i$ un sous-groupe ouvert
$H_i$ du groupe profini $\pi_1(X, \overline{\eta})$
tel que tout $\gamma \in H_i$ fixe $m_i$. Nous en d\'eduisons que
tout \`el\'ement du sous-groupe ouvert $H = \bigcap_{i = 1, \ldots, r} H_i$
fixe tout \`el\'ement de $M$. En rempla\c cant $H$ par un sous-groupe plus petit, nous pouvons supposer que $H$
est un sous-groupe ouvert normal de $\pi_1(X, \overline{\eta})$.
Posons $G = \pi_1(X, \overline{\eta})/H$. Soit $f : Y \to X$ le
rev\^etement galoisien fini \'etale de groupe $G$ correspondant.
Nous pouvons voir $f_*\underline{\mathbf{Z}}$ comme un faisceau de
$\mathbf{Z}[G]$-modules sur $X_\etale$. Il suffit alors de poser
$$
\mathcal{F} =
f_*\underline{\mathbf{Z}} \otimes_{\underline{\mathbf{Z}[G]}} \underline{M}
$$
Nous laissons au lecteur le soin de calculer
$\mathcal{F}_{\overline{\eta}}$.
Nous omettons aussi de v\'erifier que cette construction
est inverse de celle du paragraphe pr\'ec\'edent.
\end{proof}

\begin{remark}
\label{remark-functorial-locally-constant-on-connected}
Les \`equivalences des lemmes \ref{lemma-locally-constant-on-connected} et
\ref{lemma-locally-constant-on-connected-geom-unibranch}
sont compatibles aux images inverses. Par exemple, soit $f : Y \to X$
un morphisme de sch\'emas connexes. Soit $\overline{y}$ un point g\'eom\'etrique
de $Y$ et posons $\overline{x} = f(\overline{y})$.
Alors le diagramme
$$
\xymatrix{
\text{faisceaux d'ensembles localement constants finis sur }Y_\etale
\ar[r] &
\text{ensembles finis munis d'une action de }\pi_1(Y, \overline{y}) \\
\text{faisceaux d'ensembles localement constants finis sur }X_\etale
\ar[r] \ar[u]_{f^{-1}} &
\text{ensembles finis munis d'une action de }\pi_1(X, \overline{x}) \ar[u]
}
$$
est commutatif, o\`u la fl\`eche verticale de droite provient de
l'homomorphisme continu
$\pi_1(Y, \overline{y}) \to \pi_1(X, \overline{x})$
induit par $f$. Cela r\'esulte imm\'ediatement du
diagramme commutatif de
Fundamental Groups, Theorem \ref{pione-theorem-fundamental-group}.
Un r\'esultat analogue vaut dans les autres cas.
\end{remark}










\section{M\'ethode de la trace}
\label{section-trace-method}

\noindent
Une r\'ef\'erence pour cette section est \cite[Expos\'e IX, \S 5]{SGA4}.
Les r\'esultats qui suivent seront utilis\'es dans la d\'emonstration du
lemme \ref{lemma-vanishing-easier} ci-dessous.

\medskip\noindent
Soit $f : Y \to X$ un morphisme \'etale de sch\'emas. Il existe
une suite
$$
f_!, f^{-1}, f_*
$$
de foncteurs adjoints entre
$\textit{Ab}(X_\etale)$ et $\textit{Ab}(Y_\etale)$.
Le foncteur $f_!$ est \'etudi\'e dans la section \ref{section-extension-by-zero}.
Le morphisme d'adjonction $\text{id} \to f_* f^{-1}$ est appel\'e {\it restriction}.
Le morphisme d'adjonction $f_! f^{-1} \to \text{id}$ est souvent
appel\'e le {\it morphisme trace}. Si $f$ est fini \'etale, alors $f_* = f_!$
(lemme \ref{lemma-shriek-equals-star-finite-etale}) et
nous pouvons le voir comme un morphisme $f_*f^{-1} \to \text{id}$.

\begin{definition}
\label{definition-trace-map}
Soit $f : Y \to X$ un morphisme fini \'etale de sch\'emas.
Le morphisme $f_* f^{-1} \to \text{id}$ d\'ecrit ci-dessus et explicit\'e ci-dessous
est appel\'e la {\it trace}.
\end{definition}

\noindent
Soit $f : Y \to X$ un morphisme fini \'etale de sch\'emas. Le morphisme trace est
caract\'eris\'e par les deux propri\'et\'es suivantes :
\begin{enumerate}
\item il commute \`a la localisation \'etale sur $X$, et
\item si $Y = \coprod_{i = 1}^d X$, alors le morphisme trace est
l'application somme $f_*f^{-1} \mathcal{F} = \mathcal{F}^{\oplus d} \to \mathcal{F}$.
\end{enumerate}
D'apr\`es \'Etale Morphisms, Lemma \ref{etale-lemma-finite-etale-etale-local},
tout morphisme fini \'etale $f : Y \to X$ est, localement pour la topologie \'etale sur $X$,
de la forme donn\'ee en (2) pour un entier $d \geq 0$. Nous pouvons donc
d\'efinir le morphisme trace par cette caract\'erisation; en particulier,
nous n'avons besoin ni de l'existence de $f_!$, ni de l'\'egalit\'e
de $f_!$ et $f_*$ pour construire le morphisme trace.
Cette description montre que, si $f$ est de degr\'e constant $d$, alors
la compos\'ee
$$
\mathcal{F} \xrightarrow{res}
f_* f^{-1} \mathcal{F} \xrightarrow{trace}
\mathcal{F}
$$
est la multiplication par $d$. La « m\'ethode de la trace »
est l'observation suivante : si $\mathcal{F}$
est un faisceau ab\'elien sur $X_\etale$ tel que la multiplication par $d$
sur $\mathcal{F}$ soit un isomorphisme, alors l'application
$$
H^n_\etale(X, \mathcal{F}) \longrightarrow H^n_\etale(Y, f^{-1}\mathcal{F})
$$
est injective. En effet, nous avons
$$
H^n_\etale(Y, f^{-1}\mathcal{F}) = H^n_\etale(X, f_*f^{-1}\mathcal{F})
$$
par l'annulation des images directes sup\'erieures
(proposition \ref{proposition-finite-higher-direct-image-zero})
et la suite spectrale de Leray
(proposition \ref{proposition-leray}).
Nous pouvons donc consid\'erer les applications
$$
H^n_\etale(X, \mathcal{F}) \to
H^n_\etale(Y, f^{-1}\mathcal{F})= H^n_\etale(X, f_*f^{-1}\mathcal{F})
\xrightarrow{trace}
H^n_\etale(X, \mathcal{F})
$$
et leur compos\'ee est un isomorphisme (sous notre hypoth\`ese sur $\mathcal{F}$
et $f$). En particulier, si
$H_\etale^q(Y, f^{-1}\mathcal{F}) = 0$, alors
$H_\etale^q(X, \mathcal{F}) = 0$ aussi.
En effet, la multiplication par $d$ induit un
isomorphisme de $H_\etale^q(X, \mathcal{F})$ qui se factorise par
$H_\etale^q(Y, f^{-1}\mathcal{F})= 0$.

\medskip\noindent
On combine souvent ceci avec le lemme suivant.

\begin{lemma}
\label{lemma-pullback-filtered}
Soit $S$ un sch\'ema connexe. Soit $\ell$ un nombre premier. Soit
$\mathcal{F}$ un faisceau localement constant de type fini de
$\mathbf{F}_\ell$-espaces vectoriels sur $S_\etale$.
Alors il existe un morphisme fini \'etale
$f : T \to S$ de degr\'e premier \`a $\ell$ tel que $f^{-1}\mathcal{F}$
admette une filtration finie dont les quotients successifs sont
$\underline{\mathbf{Z}/\ell\mathbf{Z}}_T$.
\end{lemma}

\begin{proof}
Choisissons un point g\'eom\'etrique $\overline{s}$ de $S$.
Par l'\'equivalence du lemme \ref{lemma-locally-constant-on-connected},
le faisceau $\mathcal{F}$ correspond \`a un espace vectoriel
$V$ de dimension finie sur $\mathbf{F}_\ell$, muni d'une action continue de
$\pi_1(S, \overline{s})$.
Soit $G \subset \text{Aut}(V)$ l'image de l'homomorphisme
$\rho : \pi_1(S, \overline{s}) \to \text{Aut}(V)$ qui donne l'action.
Remarquons que $G$ est fini.
L'homomorphisme continu surjectif
$\overline{\rho} : \pi_1(S, \overline{s}) \to G$
correspond \`a un objet galoisien $Y \to S$ de
$\textit{F\'Et}_S$, de groupe d'automorphismes $G = \text{Aut}(Y/S)$, voir
Fundamental Groups, Section \ref{pione-section-finite-etale-under-galois}.
Soit $H \subset G$ un $\ell$-sous-groupe de Sylow.
Nous affirmons que $T = Y/H \to S$ convient. En effet, soit $\overline{t} \in T$
un point g\'eom\'etrique au-dessus de $\overline{s}$. L'image de
$\pi_1(T, \overline{t}) \to \pi_1(S, \overline{s})$
est $(\overline{\rho})^{-1}(H)$ par fonctorialit\'e
des groupes fondamentaux. Par cons\'equent, l'action de $\pi_1(T, \overline{t})$
sur $V$ correspondant \`a $f^{-1}\mathcal{F}$ se fait par
l'application $\pi_1(T, \overline{t}) \to H$, voir la
remarque \ref{remark-functorial-locally-constant-on-connected}. Comme
$H$ est un $\ell$-groupe fini, les constituants irr\'eductibles de la
repr\'esentation $\rho|_{\pi_1(T, \overline{t})}$
sont tous triviaux de rang $1$ (c'est un lemme \`el\'ementaire de
th\'eorie des repr\'esentations des groupes finis).
Par l'\'equivalence du
lemme \ref{lemma-locally-constant-on-connected},
cela signifie que $f^{-1}\mathcal{F}$ est une extension successive de
faisceaux constants $\underline{\mathbf{Z}/\ell\mathbf{Z}}_T$.
De plus, le degr\'e de $T = Y/H \to S$ est premier \`a $\ell$,
car il est \'egal \`a l'indice de $H$ dans $G$.
\end{proof}

\begin{lemma}
\label{lemma-action-l-group}
Soit $\Lambda$ un anneau noeth\'erien.
Soit $\ell$ un nombre premier et $n \geq 1$.
Soit $H$ un $\ell$-groupe fini.
Soit $M$ un $\Lambda[H]$-module de type fini annul\'e par $\ell^n$.
Alors il existe une filtration finie
$0 = M_0 \subset M_1 \subset \ldots \subset M_t = M$
par des sous-$\Lambda[H]$-modules
telle que $H$ agisse trivialement sur $M_{i + 1}/M_i$ pour tout
$i = 0, \ldots, t - 1$.
\end{lemma}

\begin{proof}
Omis. Indication : montrer que l'id\'eal d'augmentation $\mathfrak m$
de l'anneau non commutatif $\mathbf{Z}/\ell^n\mathbf{Z}[H]$
est nilpotent.
\end{proof}

\begin{lemma}
\label{lemma-pullback-filtered-modules}
Soit $S$ un sch\'ema irr\'eductible et g\'eom\'etriquement unibranche.
Soit $\ell$ un nombre premier et $n \geq 1$. Soit $\Lambda$
un anneau noeth\'erien. Soit $\mathcal{F}$ un faisceau localement
constant de type fini de $\Lambda$-modules sur $S_\etale$
qui est annul\'e par $\ell^n$. Alors il existe un
morphisme fini \'etale $f : T \to S$ de degr\'e premier \`a $\ell$
tel que $f^{-1}\mathcal{F}$ admette une filtration finie dont les
quotients successifs sont de la forme $\underline{M}_T$
pour des $\Lambda$-modules de type fini $M$.
\end{lemma}

\begin{proof}
Choisissons un point g\'eom\'etrique $\overline{s}$ de $S$.
Par l'\'equivalence du
lemme \ref{lemma-locally-constant-on-connected-geom-unibranch},
le faisceau $\mathcal{F}$ correspond \`a un $\Lambda$-module de type fini
$M$ muni d'une action continue de $\pi_1(S, \overline{s})$.
Soit $G \subset \text{Aut}(M)$ l'image de l'homomorphisme
$\rho : \pi_1(S, \overline{s}) \to \text{Aut}(M)$ qui donne l'action.
Remarquons que $G$ est fini par continuit\'e et parce que $M$ est engendr\'e
par un nombre fini d'\'el\'ements (voir la preuve du lemme \ref{lemma-locally-constant-on-connected-geom-unibranch}).
L'homomorphisme continu surjectif
$\overline{\rho} : \pi_1(S, \overline{s}) \to G$
correspond \`a un objet galoisien $Y \to S$ de
$\textit{F\'Et}_S$, de groupe d'automorphismes $G = \text{Aut}(Y/S)$, voir
Fundamental Groups, Section \ref{pione-section-finite-etale-under-galois}.
Soit $H \subset G$ un $\ell$-sous-groupe de Sylow.
Nous affirmons que $T = Y/H \to S$ convient. En effet, soit $\overline{t} \in T$
un point g\'eom\'etrique au-dessus de $\overline{s}$. L'image de
$\pi_1(T, \overline{t}) \to \pi_1(S, \overline{s})$
est $(\overline{\rho})^{-1}(H)$ par fonctorialit\'e
des groupes fondamentaux. Par cons\'equent, l'action de $\pi_1(T, \overline{t})$
sur $M$ correspondant \`a $f^{-1}\mathcal{F}$ se fait par
l'application $\pi_1(T, \overline{t}) \to H$, voir la
remarque \ref{remark-functorial-locally-constant-on-connected}.
Soit $0 = M_0 \subset M_1 \subset \ldots \subset M_t = M$
comme dans le lemme \ref{lemma-action-l-group}.
Cela induit une filtration
$0 = \mathcal{F}_0 \subset \mathcal{F}_1 \subset \ldots
\subset \mathcal{F}_t = f^{-1}\mathcal{F}$
telle que les quotients successifs soient constants de
valeur $M_{i + 1}/M_i$.
Enfin, le degr\'e de $T = Y/H \to S$ est premier \`a $\ell$,
car il est \'egal \`a l'indice de $H$ dans $G$.
\end{proof}






\section{Galois cohomology}
\label{section-galois-cohomology}

\noindent
In this section we prove a result on Galois cohomology
(Proposition \ref{proposition-serre-galois})
using \'etale cohomology and the trick from
Section \ref{section-trace-method}.
This will allow us to prove vanishing of higher \'etale cohomology groups
over the spectrum of a field.

\begin{lemma}
\label{lemma-nonvanishing-inherited}
Let $\ell$ be a prime number and $n$ an integer $> 0$.
Let $S$ be a quasi-compact and quasi-separated scheme.
Let $X = \lim_{i \in I} X_i$ be the limit of a
directed system of $S$-schemes each $X_i \to S$
being finite \'etale of constant degree relatively prime to $\ell$.
The following are equivalent:
\begin{enumerate}
\item there exists an $\ell$-power torsion sheaf
$\mathcal{G}$ on $S$ such that $H_\etale^n(S, \mathcal{G}) \neq 0$ and
\item there exists an $\ell$-power torsion sheaf $\mathcal{F}$ on $X$
such that $H_\etale^n(X, \mathcal{F}) \neq 0$.
\end{enumerate}
In fact, given
$\mathcal{G}$ we can take $\mathcal{F} = g^{-1}\mathcal{F}$
and given
$\mathcal{F}$ we can take $\mathcal{G} = g_*\mathcal{F}$.
\end{lemma}

\begin{proof}
Let $g : X \to S$ and $g_i : X_i \to S$ denote the structure morphisms.
Fix an $\ell$-power torsion sheaf $\mathcal{G}$ on $S$
with $H^n_\etale(S, \mathcal{G}) \not = 0$.
The system given by $\mathcal{G}_i = g_i^{-1}\mathcal{G}$
satisfy the conditions of Theorem \ref{theorem-colimit}
with colimit sheaf given by $g^{-1}\mathcal{G}$. This tells 
us that:
$$
\colim_{i\in I} H^n_\etale(X_i, g_i^{-1}\mathcal{G}) =
H^n_\etale(X, \mathcal{G})
$$
By virtue of the $g_i$ being finite \'etale morphism of degree prime
to $\ell$ we can apply ``la m\'ethode de la trace'' and we find
the maps
$$
H^n_\etale(S, \mathcal{G}) \to H^n_\etale(X_i, g_i^{-1}\mathcal{G})
$$
are all injective (and compatible with the transition maps).
See Section \ref{section-trace-method}. Thus, the colimit is non-zero, i.e.,
$H^n(X,g^{-1}\mathcal{G}) \neq 0$, giving us the desired result with 
$\mathcal{F} = g^{-1}\mathcal{G}$.

\medskip\noindent
Conversely, suppose given an $\ell$-power torsion sheaf $\mathcal{F}$ on $X$
with $H^n_\etale(X, \mathcal{F}) \not = 0$. We note that since the $g_i$
are finite morphisms the higher direct images vanish
(Proposition \ref{proposition-finite-higher-direct-image-zero}).
Then, by applying Lemma \ref{lemma-relative-colimit}
we may also conclude the  same for $g$.
The vanishing of the higher direct images tells us that
$H^n_\etale(X, \mathcal{F}) = H^n(S, g_*\mathcal{F}) \neq 0$
by Leray (Proposition \ref{proposition-leray})
giving us what we want with $\mathcal{G} = g_*\mathcal{F}$.
\end{proof}

\begin{lemma}
\label{lemma-reduce-to-l-group}
Let $\ell$ be a prime number and $n$ an integer $> 0$.
Let $K$ be a field with $G = Gal(K^{sep}/K)$ and let
$H \subset G$ be a maximal pro-$\ell$ subgroup with $L/K$
being the corresponding field extension. Then
$H^n_\etale(\Spec(K), \mathcal{F}) = 0$ for all
$\ell$-power torsion $\mathcal{F}$ if and only if
$H^n_\etale(\Spec(L), \underline{\mathbf{Z}/\ell\mathbf{Z}}) = 0$.
\end{lemma}

\begin{proof}
Write $L = \bigcup L_i$ as the union of its finite subextensions over $K$.
Our choice of $H$ implies that $[L_i : K]$ is prime to $\ell$.
Thus $\Spec(L) = \lim_{i \in I} \Spec(L_i)$ as in
Lemma \ref{lemma-nonvanishing-inherited}.
Thus we may replace $K$ by $L$ and assume that
the absolute Galois group $G$ of $K$ is a
profinite pro-$\ell$ group.

\medskip\noindent
Assume $H^n(\Spec(K), \underline{\mathbf{Z}/\ell\mathbf{Z}}) = 0$.
Let $\mathcal{F}$ be an $\ell$-power torsion sheaf on $\Spec(K)_\etale$.
We will show that $H^n_\etale(\Spec(K), \mathcal{F}) = 0$.
By the correspondence specified in
Lemma \ref{lemma-equivalence-abelian-sheaves-point}
our sheaf $\mathcal{F}$ corresponds to an $\ell$-power torsion
$G$-module $M$. Any finite set of elements $x_1, \ldots, x_m \in M$
must be fixed by an open subgroup $U$ by continuity.
Let $M'$ be the module spanned by the orbits of $x_1, \ldots, x_m$.
This is a finite abelian $\ell$-group
as each $x_i$ is killed by a power of $\ell$
and the orbits are finite. Since $M$ is the filtered colimit of
these submodules $M'$, we see that $\mathcal{F}$ is the filtered
colimit of the corresponding subsheaves $\mathcal{F}' \subset \mathcal{F}$.
Applying Theorem \ref{theorem-colimit} to this colimit, we reduce
to the case where $\mathcal{F}$ is a finite locally constant sheaf.

\medskip\noindent
Let $M$ be a finite abelian $\ell$-group with a continuous action
of the profinite pro-$\ell$ group $G$. Then there is a $G$-invariant
filtration
$$
0 = M_0 \subset M_1 \subset \ldots \subset M_r = M
$$
such that $M_{i + 1}/M_i \cong \mathbf{Z}/\ell \mathbf{Z}$ with
trivial $G$-action (this is a simple lemma on representation
theory of finite groups; insert future reference here).
Thus the corresponding sheaf $\mathcal{F}$ has a filtration
$$
0 = \mathcal{F}_0 \subset \mathcal{F}_1 \subset \ldots \subset
\mathcal{F}_r = \mathcal{F}
$$
with successive quotients isomorphic to
$\underline{\mathbf{Z}/\ell \mathbf{Z}}$.
Thus by induction and the long exact cohomology
sequence we conclude.
\end{proof}

\begin{lemma}
\label{lemma-reduce-to-l-group-higher}
Let $\ell$ be a prime number and $n$ an integer $> 0$.
Let $K$ be a field with $G = Gal(K^{sep}/K)$ and let
$H \subset G$ be a maximal pro-$\ell$ subgroup 
with $L/K$ being the corresponding field extension.
Then $H^q_\etale(\Spec(K),\mathcal{F}) = 0$ for $q \geq n$ and all
$\ell$-torsion sheaves $\mathcal{F}$ if  and only if
$H^n_\etale(\Spec(L), \underline{\mathbf{Z}/\ell\mathbf{Z}}) = 0$.
\end{lemma}

\begin{proof}
The forward direction is trivial, so we need only prove the reverse direction. 
We proceed by induction on $q$. The case of $q = n$ is
Lemma \ref{lemma-reduce-to-l-group}. Now let 
$\mathcal{F}$ be an $\ell$-power torsion sheaf on $\Spec(K)$.
Let $f : \Spec(K^{sep}) \rightarrow \Spec(K)$
be the inclusion of a geometric point.
Then consider the exact sequence:
$$
0 \rightarrow \mathcal{F} \xrightarrow{res}
f_* f^{-1} \mathcal{F} \rightarrow f_* f^{-1} \mathcal{F}/\mathcal{F} 
\rightarrow 0
$$
Note that $K^{sep}$ may be written as the filtered colimit of finite 
separable extensions. Thus $f$
is the limit of a directed system of finite \'etale  morphisms.
We may, as was seen in the proof of
Lemma \ref{lemma-nonvanishing-inherited}, conclude that $f$ has 
vanishing higher direct images. Thus, we may express the higher cohomology of 
$f_* f^{-1} \mathcal{F}$ as the higher cohomology on the geometric point which 
clearly vanishes. Hence, as everything here is still $\ell$-torsion, we may use 
the inductive hypothesis in conjunction with the long-exact cohomology sequence 
to conclude the result for $q + 1$.
\end{proof}

\begin{proposition}
\label{proposition-serre-galois}
\begin{reference}
\cite[Chapter II, Section 3, Proposition 5]{SerreGaloisCohomology}
\end{reference}
Let $K$ be a field with separable algebraic closure $K^{sep}$.
Assume that for any finite extension $K'$ of $K$ we have
$\text{Br}(K') = 0$. Then
\begin{enumerate}
\item $H^q(\text{Gal}(K^{sep}/K), (K^{sep})^*) = 0$
for all $q \geq 1$, and
\item $H^q(\text{Gal}(K^{sep}/K), M) = 0$
for any torsion $\text{Gal}(K^{sep}/K)$-module $M$ and any $q \geq 2$,
\end{enumerate}
\end{proposition}

\begin{proof}
Set $p = \text{char}(K)$.
By Lemma \ref{lemma-compare-cohomology-point},
Theorem \ref{theorem-brauer-delta},
and Example \ref{example-sheaves-point}
the proposition is equivalent to showing that if 
$H^2(\Spec(K'),\mathbf{G}_m|_{\Spec(K')_\etale}) = 0$
for all finite extensions $K'/K$ then:
\begin{itemize}
\item $H^q(\Spec(K),\mathbf{G}_m|_{\Spec(K)_\etale}) = 0$
for all $q \geq 1$, and
\item $H^q(\Spec(K),\mathcal{F}) = 0$
for any torsion sheaf $\mathcal{F}$ and any $q \geq 2$.
\end{itemize}
We prove the second part first.
Since $\mathcal{F}$ is a torsion sheaf, we may use the 
$\ell$-primary decomposition as well as the compatibility of
cohomology with colimits (i.e, direct sums, see Theorem \ref{theorem-colimit})
to reduce to showing  $H^q(\Spec(K),\mathcal{F}) = 0$, $q \geq 2$
for all $\ell$-power torsion sheaves for every prime $\ell$. This 
allows us to analyze each prime individually.

\medskip\noindent
Suppose that $\ell \neq p$. For any extension $K'/K$
consider the Kummer sequence (Lemma \ref{lemma-kummer-sequence})
$$
0 \to
\mu_{\ell, \Spec{K'}} \to
\mathbf{G}_{m, \Spec{K'}} \xrightarrow{(\cdot)^{\ell}}
\mathbf{G}_{m, \Spec{K'}} \to 0
$$
Since $H^q(\Spec{K'},\mathbf{G}_m|_{\Spec(K')_\etale}) = 0$
for $q = 2$ by assumption  and for $q = 1$ by
Theorem \ref{theorem-picard-group} combined with
$\Pic(K) = (0)$. Thus, by the long-exact cohomology 
sequence we may conclude that $H^2(\Spec{K'}, \mu_\ell) = 0$ for 
any separable $K'/K$. Now let $H$ be a maximal pro-$\ell$ subgroup
of the absolute Galois group of $K$ and let $L$ be the 
corresponding extension. We can write $L$ as the colimit of finite extensions, 
applying Theorem \ref{theorem-colimit} to this colimit we see that
$H^2(\Spec(L), \mu_\ell) = 0$.
Now $\mu_\ell$ must be the constant sheaf. If it weren't, that would imply
there exists a Galois extension of degree relatively prime to 
$\ell$ of $L$ which is not true by definition of $L$ (namely, the extension
one gets by adjoining the  $\ell$th roots of unity to $L$).
Hence, via Lemma \ref{lemma-reduce-to-l-group-higher},
we conclude the result for $\ell \neq p$. 

\medskip\noindent
Now suppose that $\ell = p$. We consider the
Artin-Schrier exact sequence (Section \ref{section-artin-schreier})
$$
0 \longrightarrow \underline{\mathbf{Z}/p\mathbf{Z}}_{\Spec{K}} \longrightarrow
\mathbf{G}_{a, \Spec{K}} \xrightarrow{F-1} \mathbf{G}_{a, \Spec{K}} 
\longrightarrow 0
$$
where $F - 1$ is the map $x \mapsto x^p - x$. Then note that the higher 
Cohomology of $\mathbf{G}_{a, \Spec{K}}$ vanishes, by
Remark \ref{remark-special-case-fpqc-cohomology-quasi-coherent} and the 
vanishing of the higher cohomology of the structure sheaf of an affine scheme
(Cohomology of Schemes, Lemma
\ref{coherent-lemma-quasi-coherent-affine-cohomology-zero}).
Note this can be applied to any field of 
characteristic $p$. In particular, we can apply it to the field extension $L$ 
defined by a maximal pro-$p$ subgroup $H$. This allows us to conclude 
$H^n(\Spec{L}, \underline{\mathbf{Z}/p\mathbf{Z}}_{\Spec{L}}) = 0$
for $n \geq 2$, from which the result follows for $\ell = p$, by
Lemma \ref{lemma-reduce-to-l-group-higher}.

\medskip\noindent
To finish the proof we still have to show that
$H^q(\text{Gal}(K^{sep}/K), (K^{sep})^*) = 0$ for all $q \geq 1$.
Set $G = \text{Gal}(K^{sep}/K)$ and set $M = (K^{sep})^*$
viewed as a $G$-module. We have already shown (above) that
$H^1(G, M) = 0$ and $H^2(G, M) = 0$. Consider the exact sequence
$$
0 \to A \to M \to M \otimes \mathbf{Q} \to B \to 0
$$
of $G$-modules. By the above we have $H^i(G, A) = 0$
and $H^i(G, B) = 0$ for $i > 1$ since $A$ and $B$ are
torsion $G$-modules. By
Lemma \ref{lemma-profinite-group-cohomology-torsion}
we have $H^i(G, M \otimes \mathbf{Q}) = 0$ for $i > 0$.
It is a pleasant exercise to see that this implies that
$H^i(G, M) = 0$ also for $i \geq 3$.
\end{proof}

%10.08.09

\begin{definition}
\label{definition-Cr}
A field $K$ is called {\it $C_r$}
if for every $0 < d^r < n$ and every $f \in K[T_1,
\ldots, T_n]$ homogeneous of degree $d$, there exist $\alpha = (\alpha_1,
\ldots, \alpha_n)$, $\alpha_i \in K$ not all zero, such that $f(\alpha) = 0$.
Such an $\alpha$ is called a {\it nontrivial solution} of $f$.
\end{definition}

\begin{example}
\label{example-algebraically-closed-field-Cr}
An algebraically closed field is $C_r$.
\end{example}

\noindent
In fact, we have the following simple lemma.

\begin{lemma}
\label{lemma-algebraically-closed-find-solutions}
Let $k$ be an algebraically closed field. Let
$f_1, \ldots, f_s \in k[T_1, \ldots, T_n]$
be homogeneous polynomials of degree $d_1, \ldots, d_s$ with $d_i
> 0$. If $s < n$, then $f_1 = \ldots = f_s = 0$ have a common nontrivial
solution.
\end{lemma}

\begin{proof}
This follows from dimension theory, for example in the form of
Varieties, Lemma \ref{varieties-lemma-intersection-in-affine-space}
applied $s - 1$ times.
\end{proof}

\noindent
The following result computes the Brauer group of $C_1$ fields.

\begin{theorem}
\label{theorem-C1-brauer-group-zero}
Let $K$ be a $C_1$ field. Then $\text{Br}(K) = 0$.
\end{theorem}

\begin{proof}
Let $D$ be a finite dimensional division algebra over $K$ with center $K$. We
have seen that
$$
D \otimes_K K^{sep} \cong \text{Mat}_d(K^{sep})
$$
uniquely up to inner isomorphism. Hence the determinant $\det :
\text{Mat}_d(K^{sep}) \to K^{sep}$ is Galois invariant and descends to a
homogeneous degree $d$ map
$$
\det = N_\text{red} : D \longrightarrow K
$$
called the {\it reduced norm}. Since $K$ is $C_1$, if $d > 1$, then there
exists a nonzero $x \in D$ with $N_\text{red}(x) = 0$. This clearly implies
that $x$ is not invertible, which is a contradiction. Hence $\text{Br}(K) = 0$.
\end{proof}

\begin{definition}
\label{definition-variety}
Let $k$ be a field. A {\it variety} is separated, integral scheme of
finite type over $k$. A {\it curve} is a variety of dimension $1$.
\end{definition}

\begin{theorem}[Tsen's theorem]
\label{theorem-tsen}
The function field of a variety of dimension $r$ over an algebraically closed
field $k$ is $C_r$.
\end{theorem}

\begin{proof}
For projective space one can show directly that the field
$k(x_1, \ldots, x_r)$ is $C_r$ (exercise).

\medskip\noindent
General case. Without loss of generality, we may assume $X$ to be projective.
Let $f \in k(X)[T_1, \ldots, T_n]_d$ with $0 < d^r < n$.
Say the coefficients of $f$ are in $\Gamma(X, \mathcal{O}_X(H))$
for some ample $H \subset X$. Let
$\mathbf{\alpha} = (\alpha_1, \ldots, \alpha_n)$ with $\alpha_i \in \Gamma(X,
\mathcal{O}_X(eH))$. Then $f(\mathbf{\alpha}) \in \Gamma(X,
\mathcal{O}_X((de + 1)H))$. Consider the system of equations $f(\mathbf{\alpha})
=0$. Then by asymptotic Riemann-Roch
(Varieties, Proposition \ref{varieties-proposition-asymptotic-riemann-roch})
there exists a $c > 0$ such that
\begin{itemize}
\item the number of variables is
$n\dim_k \Gamma(X, \mathcal{O}_X(eH)) \sim n e^r c$, and
\item the number of equations is
$\dim_k \Gamma(X, \mathcal{O}_X((de + 1)H)) \sim (de + 1)^r c$.
\end{itemize}
Since $n > d^r$, there are more variables than equations. The equations are
homogeneous hence there is a solution by
Lemma \ref{lemma-algebraically-closed-find-solutions}.
\end{proof}

\begin{lemma}
\label{lemma-curve-brauer-zero}
Let $C$ be a curve over an algebraically closed field $k$. Then
the Brauer group of the function field of $C$ is zero:
$\text{Br}(k(C)) = 0$.
\end{lemma}

\begin{proof}
This is clear from Tsen's theorem,
Theorem \ref{theorem-tsen} and
Theorem \ref{theorem-C1-brauer-group-zero}.
\end{proof}

\begin{lemma}
\label{lemma-cohomology-Gm-function-field-curve}
Let $k$ be an algebraically closed field and $K/k$ a field extension
of transcendence degree 1. Then for all $q \geq 1$,
$H_\etale^q(\Spec(K), \mathbf{G}_m) = 0$.
\end{lemma}

\begin{proof}
Recall that
$H_\etale^q(\Spec(K), \mathbf{G}_m) = H^q(\text{Gal}(K^{sep}/K), (K^{sep})^*)$
by Lemma \ref{lemma-compare-cohomology-point}.
Thus by Proposition \ref{proposition-serre-galois}
it suffices to show that if $K'/K$ is a finite field extension, then
$\text{Br}(K') = 0$. Now observe that $K' = \colim K''$, where $K''$ runs
over the finitely generated subextensions of $k$ contained in $K'$ of
transcendence degree $1$.
Note that $\text{Br}(K') = \colim \text{Br}(K'')$ which reduces us
to a finitely generated field extension $K''/k$ of transcendence
degree $1$. Such a field is the function field of a curve over $k$,
hence has trivial Brauer group by
Lemma \ref{lemma-curve-brauer-zero}.
\end{proof}






\section{Higher vanishing for the multiplicative group}
\label{section-higher-Gm}

\noindent
In this section, we fix an algebraically closed field $k$ and a smooth curve
$X$ over $k$. We denote $i_x : x \hookrightarrow X$ the inclusion of a closed
point of $X$ and $j : \eta \hookrightarrow X$ the inclusion of the generic
point. We also denote $X_0$ the set of closed points of $X$.

\begin{theorem}[The Fundamental Exact Sequence]
\label{theorem-fundamental-exact-sequence}
There is a short exact sequence of \'etale sheaves on $X$
$$
0 \longrightarrow
\mathbf{G}_{m, X} \longrightarrow
j_* \mathbf{G}_{m, \eta} \longrightarrow
\bigoplus\nolimits_{x \in X_0} {i_x}_* \underline{\mathbf{Z}}
\longrightarrow 0.
$$
\end{theorem}

\begin{proof}
Let $\varphi : U \to X$ be an \'etale morphism. Then by properties of
\'etale morphisms (Proposition \ref{proposition-etale-morphisms}),
$U = \coprod_i U_i$ where each $U_i$ is a smooth curve mapping to $X$.
The above sequence for $U$ is a product of the corresponding sequences
for each $U_i$, so it suffices to treat the case where $U$ is connected,
hence irreducible. In this case, there is a well known exact sequence
$$
1 \longrightarrow
\Gamma(U, \mathcal{O}_U^*) \longrightarrow
k(U)^* \longrightarrow \bigoplus\nolimits_{y \in U_0} \mathbf{Z}_y.
$$
This amounts to a sequence
$$
0 \longrightarrow \Gamma(U, \mathcal{O}_U^*) \longrightarrow
\Gamma(\eta \times_X U, \mathcal{O}_{\eta \times_X U}^*) \longrightarrow
\bigoplus\nolimits_{x \in X_0}
\Gamma(x \times_X U, \underline{\mathbf{Z}})
$$
which, unfolding definitions, is nothing but a sequence
$$
0 \longrightarrow \mathbf{G}_m(U) \longrightarrow
j_* \mathbf{G}_{m, \eta}(U) \longrightarrow
\left(\bigoplus\nolimits_{x \in X_0} {i_x}_* \underline{\mathbf{Z}}\right)(U).
$$
This defines the maps in the Fundamental Exact Sequence and shows it is exact
except possibly at the last step. To see surjectivity, let us recall that if
$U$ is a nonsingular curve and $D$ is a divisor on $U$,
then there exists a Zariski open covering $\{U_j \to U\}$ of $U$
such that $D|_{U_j} = \text{div}(f_j)$ for some $f_j \in k(U)^*$.
\end{proof}

\begin{lemma}
\label{lemma-higher-direct-jstar-Gm}
For any $q \geq 1$, $R^q j_*\mathbf{G}_{m, \eta} = 0$.
\end{lemma}

\begin{proof}
We need to show that $(R^q j_*\mathbf{G}_{m, \eta})_{\bar x} = 0$ for every
geometric point $\bar x$ of $X$.

\medskip\noindent
Assume that $\bar x$ lies over a closed point $x$ of $X$. Let $\Spec(A)$
be an affine open neighbourhood of $x$ in $X$, and $K$ the fraction field
of $A$. Then
$$
\Spec(\mathcal{O}^{sh}_{X, \bar x}) \times_X \eta =
\Spec(\mathcal{O}^{sh}_{X, \bar x} \otimes_A K).
$$
The ring $\mathcal{O}^{sh}_{X, \bar x} \otimes_A K$ is a localization of
the discrete valuation ring $\mathcal{O}^{sh}_{X, \bar x}$, so it is either
$\mathcal{O}^{sh}_{X, \bar x}$ again, or its fraction field
$K^{sh}_{\bar x}$. But since some local uniformizer gets inverted, it must
be the latter. Hence
$$
(R^q j_*\mathbf{G}_{m, \eta})_{(X, \bar x)} = H_\etale^q(\Spec
K^{sh}_{\bar x}, \mathbf{G}_m).
$$
Now recall that $\mathcal{O}^{sh}_{X, \bar x} =
\colim_{(U, \bar u) \to \bar x} \mathcal{O}(U) = \colim_{A \subset B} B$
where $A \to B$ is \'etale, hence $K^{sh}_{\bar x}$ is an algebraic
extension of $K = k(X)$, and we may apply
Lemma \ref{lemma-cohomology-Gm-function-field-curve}
to get the vanishing.

\medskip\noindent
Assume that $\bar x = \bar \eta$ lies over the generic point $\eta$ of $X$ (in
fact, this case is superfluous). Then
$\mathcal{O}^{sh}_{X, \bar \eta} = \kappa(\eta)^{sep}$ and thus
\begin{eqnarray*}
(R^q j_*\mathbf{G}_{m, \eta})_{\bar \eta}
& = &
H_\etale^q(\Spec(\kappa(\eta)^{sep}) \times_X \eta, \mathbf{G}_m) \\
& = & H_\etale^q (\Spec(\kappa(\eta)^{sep}), \mathbf{G}_m) \\
& = & 0 \ \ \text{ for } q \geq 1
\end{eqnarray*}
since the corresponding Galois group is trivial.
\end{proof}

\begin{lemma}
\label{lemma-cohomology-jstar-Gm}
For all $p \geq 1$, $H_\etale^p(X, j_*\mathbf{G}_{m, \eta}) = 0$.
\end{lemma}

\begin{proof}
The Leray spectral sequence reads
$$
E_2^{p, q} = H_\etale^p(X, R^qj_*\mathbf{G}_{m, \eta}) \Rightarrow
H_\etale^{p+q}(\eta, \mathbf{G}_{m, \eta}),
$$
The right hand side vanishes for $p + q \geq 1$ by
Lemma \ref{lemma-cohomology-Gm-function-field-curve}.
The left hand side vanishes for $q \geq 1$ by
Lemma \ref{lemma-higher-direct-jstar-Gm}. Thus we get the desired vanishing.
\end{proof}

\begin{lemma}
\label{lemma-cohomology-istar-Z}
For all $q \geq 1$, $H_\etale^q(X, \bigoplus_{x \in X_0} {i_x}_*
\underline{\mathbf{Z}}) = 0$.
\end{lemma}

\begin{proof}
For $X$ quasi-compact and quasi-separated, cohomology commutes with colimits,
so it suffices to show the vanishing of $H_\etale^q(X, {i_x}_*
\underline{\mathbf{Z}})$. But then the inclusion $i_x$ of a closed point is
finite so $R^p {i_x}_* \underline{\mathbf{Z}} = 0$ for all $p \geq 1$ by
Proposition \ref{proposition-finite-higher-direct-image-zero}.
Applying the Leray spectral sequence, we see that
$H_\etale^q(X, {i_x}_* \underline{\mathbf{Z}}) =
H_\etale^q(x, \underline{\mathbf{Z}})$.
Finally, since $x$ is the spectrum of an
algebraically closed field, all higher cohomology on $x$ vanishes.
\end{proof}

\noindent
Concluding this series of lemmata, we get the following result.

\begin{theorem}
\label{theorem-vanishing-cohomology-Gm-curve}
Let $X$ be a smooth curve over an algebraically closed field. Then
$$
H_\etale^q(X, \mathbf{G}_m) = 0 \ \ \text{ for all } q \geq 2.
$$
\end{theorem}

\begin{proof}
See discussion above.
\end{proof}

\noindent
We also get the cohomology long exact sequence
$$
0 \to
H_\etale^0(X, \mathbf{G}_m) \to
H_\etale^0(X, j_*\mathbf{G}_{m\eta}) \to
H_\etale^0(X, \bigoplus {i_x}_*\underline{\mathbf{Z}}) \to
H_\etale^1(X, \mathbf{G}_m) \to 0
$$
although this is the familiar
$$
0 \to H_{Zar}^0(X, \mathcal{O}_X^*) \to k(X)^* \to \text{Div}(X)
\to \Pic(X) \to 0.
$$





\section{Picard groups of curves}
\label{section-pic-curves}

\noindent
Our next step is to use the Kummer sequence to deduce some information
about the cohomology group of a curve with finite coefficients. In order
to get vanishing in the long exact sequence, we review some facts about
Picard groups.

\medskip\noindent
Let $X$ be a smooth projective curve over an algebraically closed field $k$.
Let $g = \dim_k H^1(X, \mathcal{O}_X)$ be the genus of $X$.
There exists a short exact sequence
$$
0 \to \Pic^0(X) \to \Pic(X) \xrightarrow{\deg} \mathbf{Z} \to 0.
$$
The abelian group $\Pic^0(X)$ can be identified with
$\Pic^0(X) = \underline{\Picardfunctor}^0_{X/k}(k)$, i.e.,
the $k$-valued points of an abelian variety
$\underline{\Picardfunctor}^0_{X/k}$ over $k$ of dimension $g$.
Consequently, if $n \in k^*$ then
$\Pic^0(X)[n] \cong (\mathbf{Z}/n\mathbf{Z})^{2g}$
as abelian groups. See
Picard Schemes of Curves, Section \ref{pic-section-picard-curve}
and
Groupoids, Section \ref{groupoids-section-abelian-varieties}.
This key fact, namely the description of the torsion in the
Picard group of a smooth projective curve over an algebraically
closed field does not appear to have an elementary proof.

\begin{lemma}
\label{lemma-cohomology-smooth-projective-curve}
Let $X$ be a smooth projective curve of genus $g$ over an
algebraically closed field $k$ and let $n \geq 1$ be invertible in $k$.
Then there are canonical identifications
$$
H_\etale^q(X, \mu_n) =
\left\{
\begin{matrix}
\mu_n(k) & \text{ if }q = 0, \\
\Pic^0(X)[n] & \text{ if }q = 1, \\
\mathbf{Z}/n\mathbf{Z} & \text{ if }q = 2, \\
0 & \text{ if }q \geq 3.
\end{matrix}
\right.
$$
Since $\mu_n \cong \underline{\mathbf{Z}/n\mathbf{Z}}$, this gives
(noncanonical) identifications
$$
H_\etale^q(X, \underline{\mathbf{Z}/n\mathbf{Z}}) \cong
\left\{
\begin{matrix}
\mathbf{Z}/n\mathbf{Z} & \text{ if }q = 0, \\
(\mathbf{Z}/n\mathbf{Z})^{2g} & \text{ if }q = 1, \\
\mathbf{Z}/n\mathbf{Z} & \text{ if }q = 2, \\
0 & \text{ if }q \geq 3.
\end{matrix}
\right.
$$
\end{lemma}

\begin{proof}
Theorems \ref{theorem-picard-group} and
\ref{theorem-vanishing-cohomology-Gm-curve}
determine the \'etale cohomology of $\mathbf{G}_m$ on $X$
in terms of the Picard group of $X$.
The Kummer sequence $0\to \mu_{n, X} \to \mathbf{G}_{m, X}
\to \mathbf{G}_{m, X}\to 0$ (Lemma \ref{lemma-kummer-sequence})
then gives us the long exact cohomology sequence
$$
\xymatrix{
0 \ar[r] & \mu_n(k) \ar[r] &
k^* \ar[r]^{(\cdot)^n} &
k^* \ar@(rd, ul)[rdllllr] \\
& H_\etale^1(X, \mu_n) \ar[r] &
\Pic(X) \ar[r]^{(\cdot)^n} &
\Pic(X) \ar@(rd, ul)[rdllllr] \\
& H_\etale^2(X, \mu_n) \ar[r] & 0 \ar[r] & 0 \ldots
}
$$
The $n$th power map $k^* \to k^*$ is surjective since $k$ is algebraically
closed. So we need to compute the kernel and cokernel of the map
$n : \Pic(X) \to \Pic(X)$. Consider the commutative diagram with
exact rows
$$
\xymatrix{
0 \ar[r] &
\Pic^0(X) \ar[r] \ar@{>>}[d]^{(\cdot)^n} &
\Pic(X) \ar[r]_-\deg \ar[d]^{(\cdot)^n} &
\mathbf{Z} \ar[r] \ar@{^{(}->}[d]^n & 0 \\
0 \ar[r] &
\Pic^0(X) \ar[r] &
\Pic(X) \ar[r]^-\deg &
\mathbf{Z} \ar[r] & 0
}
$$
The group $\Pic^0(X)$ is the $k$-points of
the group scheme $\underline{\Picardfunctor}^0_{X/k}$, see
Picard Schemes of Curves, Lemma \ref{pic-lemma-picard-pieces}.
The same lemma tells us that $\underline{\Picardfunctor}^0_{X/k}$
is a $g$-dimensional abelian variety over $k$ as defined in
Groupoids, Definition \ref{groupoids-definition-abelian-variety}.
Hence the left vertical map is surjective by
Groupoids, Proposition \ref{groupoids-proposition-review-abelian-varieties}.
Applying the snake lemma gives canonical identifications as stated
in the lemma.

\medskip\noindent
To get the noncanonical identifications of the lemma we need to
show the kernel of $n : \Pic^0(X) \to \Pic^0(X)$
is isomorphic to $(\mathbf{Z}/n\mathbf{Z})^{\oplus 2g}$.
This is also part of Groupoids, Proposition
\ref{groupoids-proposition-review-abelian-varieties}.
\end{proof}

\begin{lemma}
\label{lemma-pullback-on-h2-curve}
Let $\pi : X \to Y$ be a nonconstant morphism of smooth projective curves
over an algebraically closed field $k$ and let $n \geq 1$ be invertible in $k$.
The map
$$
\pi^* : H^2_\etale(Y, \mu_n) \longrightarrow H^2_\etale(X, \mu_n)
$$
is given by multiplication by the degree of $\pi$.
\end{lemma}

\begin{proof}
Observe that the statement makes sense as we have identified both
cohomology groups $ H^2_\etale(Y, \mu_n)$ and $H^2_\etale(X, \mu_n)$
with $\mathbf{Z}/n\mathbf{Z}$ in
Lemma \ref{lemma-cohomology-smooth-projective-curve}.
In fact, if $\mathcal{L}$ is a line bundle of degree $1$
on $Y$ with class $[\mathcal{L}] \in H^1_\etale(Y, \mathbf{G}_m)$,
then the coboundary of $[\mathcal{L}]$ is the generator of
$H^2_\etale(Y, \mu_n)$. Here the coboundary is the coboundary
of the long exact sequence of cohomology associated to the Kummer
sequence. Thus the result of the lemma follows from the fact that
the degree of the line bundle $\pi^*\mathcal{L}$ on $X$ is $\deg(\pi)$.
Some details omitted.
\end{proof}

\begin{lemma}
\label{lemma-vanishing-cohomology-mu-smooth-curve}
Let $X$ be an affine smooth curve over an algebraically closed field $k$ and
$n \in k^*$. Let $X \subset \overline{X}$ be a smooth projective
compactification
(Varieties, Remark \ref{varieties-remark-smooth-projective-compactification}).
Let $g$ be the genus of $\overline{X}$ and let $r$ be the number of
points of $\overline{X} \setminus X$. Then
\begin{enumerate}
\item $H_\etale^0(X, \mu_n) = \mu_n(k)$;
\item $H_\etale^1(X, \mu_n) \cong (\mathbf{Z}/n\mathbf{Z})^{2g+r-1}$, and
\item $H_\etale^q(X, \mu_n) = 0$ for all $q \geq 2$.
\end{enumerate}
\end{lemma}

\begin{proof}
Write $X = \overline{X} - \{ x_1, \ldots, x_r\}$. Then
$\Pic(X) = \Pic(\overline{X})/ R$, where $R$ is the subgroup generated by
$\mathcal{O}_{\overline{X}}(x_i)$, $1 \leq i \leq r$.
Since $r \geq 1$, we see that $\Pic^0(\overline{X}) \to \Pic(X)$
is surjective, hence $\Pic(X)$ is divisible (see discussion in proof
of Lemma \ref{lemma-cohomology-smooth-projective-curve}).
Applying the Kummer sequence, we get (1) and (3). For (2) we use
\begin{align*}
H_\etale^1(X, \mu_n)
& =
\{(\mathcal L, \alpha)\}/\cong \\
& =
(\{(\bar{\mathcal L}, D, \bar \alpha)\}/\cong)/\tilde{R}
\end{align*}
For the first equality sign and notation see Lemma \ref{lemma-describe-h1-mun}.
The triples $(\bar{\mathcal L}, D, \bar \alpha)$ in the second line
consist of an invertible module $\bar{\mathcal L}$ on $\overline{X}$
of degree $0$, a divisor $D$ on $\overline{X}$ of degree $0$ supported on
$\left\{x_1, \ldots, x_r\right\}$ and an isomorphism
$\bar{\alpha}: \bar{\mathcal L}^{\otimes n} \cong \mathcal{O}_{\bar{X}}(D)$.
Isomorphisms of triples are defined as in the discussion preceding
Lemma \ref{lemma-describe-h1-mun}.
Further $\tilde{R}$ is the subgroup generated by triples of the form
$(\mathcal{O}_{\overline{X}}(D'), n D', 1^{\otimes n})$
where $D'$ is supported on $\left\{x_1, \ldots, x_r\right\}$ and has degree 0.
The discussion above shows we have the second equality sign via
the map sending the triple $(\bar{\mathcal L},\ D,\ \bar \alpha)$
to the pair $(\bar{\mathcal L}|_X, \bar \alpha|_X)$. Thus,
we get an exact sequence
$$
0 \longrightarrow
H_\etale^1(\overline{X}, \mu_n) \longrightarrow
H_\etale^1(X, \mu_n) \longrightarrow
\bigoplus_{i = 1}^r \mathbf{Z}/n\mathbf{Z}
\xrightarrow{\ \sum\ }
\mathbf{Z}/n\mathbf{Z} \longrightarrow 0
$$
where the middle map sends the class of a triple
$(\bar{ \mathcal L}, D, \bar \alpha)$ with $D = \sum_{i = 1}^r a_i (x_i)$
to the $r$-tuple $(a_i)_{i = 1}^r$. It now suffices to use
Lemma \ref{lemma-cohomology-smooth-projective-curve}
to count ranks.
\end{proof}

\begin{remark}
\label{remark-natural-proof}
The ``natural'' way to prove the previous corollary is to excise $X$ from $\bar
X$. This is possible, we just haven't developed that theory.
\end{remark}

\begin{remark}
\label{remark-normalize-H1-Gm}
Let $k$ be an algebraically closed field. Let $n$ be an integer prime to
the characteristic of $k$. Recall that
$$
\mathbf{G}_{m, k} = \mathbf{A}^1_k \setminus \{0\} =
\mathbf{P}^1_k \setminus \{0, \infty\}
$$
We claim there is a canonical isomorphism
$$
H^1_\etale(\mathbf{G}_{m, k}, \mu_n) = \mathbf{Z}/n\mathbf{Z}
$$
What does this mean? This means there is an element $1_k$ in
$H^1_\etale(\mathbf{G}_{m, k}, \mu_n)$ such that for
every morphism $\Spec(k') \to \Spec(k)$ the pullback map on
\'etale cohomology for the map $\mathbf{G}_{m, k'} \to \mathbf{G}_{m, k}$
maps $1_k$ to $1_{k'}$. (In particular this element is
fixed under all automorphisms of $k$.) To see this, consider the
$\mu_{n, \mathbf{Z}}$-torsor
$\mathbf{G}_{m, \mathbf{Z}} \to \mathbf{G}_{m, \mathbf{Z}}$,
$x \mapsto x^n$. By the identification of torsors with
first cohomology, this pulls back to give our canonical elements $1_k$.
Twisting back we see that there are canonical identifications
$$
H^1_\etale(\mathbf{G}_{m, k}, \mathbf{Z}/n\mathbf{Z}) =
\Hom(\mu_n(k), \mathbf{Z}/n\mathbf{Z}),
$$
i.e., these isomorphisms are compatible with respect to maps of
algebraically closed fields, in particular with respect to
automorphisms of $k$.
\end{remark}










\section{Extension by zero}
\label{section-extension-by-zero}

\noindent
The general material in
Modules on Sites, Section \ref{sites-modules-section-localize}
allows us to make the following definition.

\begin{definition}
\label{definition-extension-zero}
Let $j : U \to X$ be an \'etale morphism of schemes.
\begin{enumerate}
\item The restriction functor
$j^{-1} : \Sh(X_\etale) \to \Sh(U_\etale)$
has a left adjoint
$j_!^{Sh} : \Sh(U_\etale) \to \Sh(X_\etale)$.
\item The restriction functor
$j^{-1} : \textit{Ab}(X_\etale) \to \textit{Ab}(U_\etale)$
has a left adjoint which is denoted
$j_! : \textit{Ab}(U_\etale) \to \textit{Ab}(X_\etale)$
and called {\it extension by zero}.
\item Let $\Lambda$ be a ring. The restriction functor
$j^{-1} : \textit{Mod}(X_\etale, \Lambda) \to
\textit{Mod}(U_\etale, \Lambda)$
has a left adjoint which is denoted
$j_! : \textit{Mod}(U_\etale, \Lambda) \to
\textit{Mod}(X_\etale, \Lambda)$
and called {\it extension by zero}.
\end{enumerate}
\end{definition}

\noindent
If $\mathcal{F}$ is an abelian sheaf on $X_\etale$, then
$j_!\mathcal{F} \not = j_!^{Sh}\mathcal{F}$ in general. On the other hand
$j_!$ for sheaves of $\Lambda$-modules agrees with $j_!$ on underlying
abelian sheaves
(Modules on Sites, Remark \ref{sites-modules-remark-localize-shriek-equal}).
The functor $j_!$ is characterized by the functorial isomorphism
$$
\Hom_X(j_!\mathcal{F}, \mathcal{G}) = \Hom_U(\mathcal{F}, j^{-1}\mathcal{G})
$$
for all $\mathcal{F} \in \textit{Ab}(U_\etale)$ and
$\mathcal{G} \in \textit{Ab}(X_\etale)$. Similarly for
sheaves of $\Lambda$-modules.

\medskip\noindent
To describe the functors in Definition \ref{definition-extension-zero}
more explicitly, recall that $j^{-1}$ is just the restriction
via the functor $U_\etale \to X_\etale$. In other words,
$j^{-1}\mathcal{G}(U') = \mathcal{G}(U')$ for $U'$ \'etale over $U$.
On the other hand, for $\mathcal{F} \in \textit{Ab}(U_\etale)$
we consider the presheaf
\begin{equation}
\label{equation-j-p-shriek}
j_{p!}\mathcal{F} : X_\etale \longrightarrow \textit{Ab},
\quad
V \longmapsto \bigoplus\nolimits_{V \to U} \mathcal{F}(V \to U)
\end{equation}
Then $j_!\mathcal{F}$ is the sheafification of $j_{p!}\mathcal{F}$.
This is proven in Modules on Sites, Lemma
\ref{sites-modules-lemma-extension-by-zero};
more generally see the discussion in
Modules on Sites, Sections \ref{sites-modules-section-localize} and
\ref{sites-modules-section-exactness-lower-shriek}.

\begin{exercise}
\label{exercise-jshriek-direct}
Prove directly that the functor $j_!$ defined as the sheafification of
the functor $j_{p!}$ given in
(\ref{equation-j-p-shriek}) is a left adjoint to $j^{-1}$.
\end{exercise}

\begin{proposition}
\label{proposition-describe-jshriek}
Let $j : U \to X$ be an \'etale morphism of schemes.
Let $\mathcal{F}$ in $\textit{Ab}(U_\etale)$.
If $\overline{x} : \Spec(k) \to X$ is a geometric point of $X$, then 
$$
(j_!\mathcal{F})_{\overline{x}} =
\bigoplus\nolimits_{\overline{u} : \Spec(k) \to U,\ j(\overline{u}) =
\overline{x}} \mathcal{F}_{\bar{u}}.
$$
In particular, $j_!$ is an exact functor.
\end{proposition}

\begin{proof}
Exactness of $j_!$ is very general, see Modules on Sites, 
Lemma \ref{sites-modules-lemma-extension-by-zero-exact}.
Of course it does also follow from the description of stalks.
The formula for the stalk follows from
Modules on Sites, Lemma \ref{sites-modules-lemma-stalk-j-shriek}
and the description of points of the small \'etale site
in terms of geometric points, see Lemma \ref{lemma-points-small-etale-site}.

\medskip\noindent
For later use we note that the isomorphism
\begin{align*}
(j_!\mathcal{F})_{\overline{x}}
& =
(j_{p!}\mathcal{F})_{\overline{x}} \\
& =
\colim_{(V, \overline{v})} j_{p!}\mathcal{F}(V) \\
& =
\colim_{(V, \overline{v})}
\bigoplus\nolimits_{\varphi : V \to U}
\mathcal{F}(V \xrightarrow{\varphi} U) \\
& \to
\bigoplus\nolimits_{\overline{u} : \Spec(k) \to U,\ j(\overline{u}) =
\overline{x}} \mathcal{F}_{\bar{u}}.
\end{align*}
constructed in Modules on Sites, Lemma \ref{sites-modules-lemma-stalk-j-shriek}
sends $(V, \overline{v}, \varphi, s)$ to the class of $s$ in the stalk
of $\mathcal{F}$ at $\overline{u} = \varphi(\overline{v})$.
\end{proof}

\begin{lemma}
\label{lemma-jshriek-open}
Let $j : U \to X$ be an open immersion of schemes. For any
abelian sheaf $\mathcal{F}$ on $U_\etale$, the adjunction mappings
$j^{-1}j_*\mathcal{F} \to \mathcal{F}$ and
$\mathcal{F} \to j^{-1}j_!\mathcal{F}$ are isomorphisms.
In fact, $j_!\mathcal{F}$ is the unique abelian sheaf on $X_\etale$
whose restriction to $U$ is $\mathcal{F}$ and whose stalks at
geometric points of $X \setminus U$ are zero.
\end{lemma}

\begin{proof}
We encourage the reader to prove the first statement by working through the
definitions, but here we just use that it is a special case of the very general
Modules on Sites, Lemma \ref{sites-modules-lemma-restrict-back}.
For the second statement, observe that if $\mathcal{G}$ is an abelian sheaf
on $X_\etale$ whose restriction to $U$ is $\mathcal{F}$, then we obtain
by adjointness a map $j_!\mathcal{F} \to \mathcal{G}$. This map
is then an isomorphism at stalks of geometric points of $U$ by
Proposition \ref{proposition-describe-jshriek}.
Thus if $\mathcal{G}$ has vanishing stalks at geometric points
of $X \setminus U$, then $j_!\mathcal{F} \to \mathcal{G}$ is an
isomorphism by Theorem \ref{theorem-exactness-stalks}.
\end{proof}

\begin{lemma}[Extension by zero commutes with base change]
\label{lemma-shriek-base-change}
Let $f: Y \to X$ be a morphism of schemes. Let $j: V \to X$ be an \'etale
morphism. Consider the fibre product
$$
\xymatrix{
V' = Y \times_X V \ar[d]_{f'} \ar[r]_-{j'} & Y \ar[d]^f \\
V \ar[r]^j & X
}
$$
Then we have $j'_! f'^{-1} = f^{-1} j_!$ on abelian sheaves and on
sheaves of modules.
\end{lemma}

\begin{proof}
This is true because $j'_! f'^{-1}$ is left adjoint to
$f'_* (j')^{-1}$ and $f^{-1} j_!$ is left adjoint to $j^{-1}f_*$.
Further $f'_* (j')^{-1} = j^{-1}f_*$ because $f_*$ commutes with
\'etale localization (by construction). In fact, the lemma holds very generally
in the setting of a morphism of sites, see
Modules on Sites, Lemma
\ref{sites-modules-lemma-localize-morphism-ringed-sites}.
\end{proof}

\begin{lemma}
\label{lemma-shriek-into-star-separated-etale}
Let $j : U \to X$ be separated and \'etale. Then there is a functorial
injective map $j_!\mathcal{F} \to j_*\mathcal{F}$
on abelian sheaves and sheaves of $\Lambda$-modules.
\end{lemma}

\begin{proof}
We prove this in the case of abelian sheaves. Let us construct a canonical map
$$
j_{p!}\mathcal{F} \to j_*\mathcal{F}
$$
of abelian presheaves on $X_\etale$ for any abelian sheaf
$\mathcal{F}$ on $U_\etale$ where $j_{p!}$ is as in
(\ref{equation-j-p-shriek}). Sheafification of this map will
be the desired map $j_!\mathcal{F} \to j_*\mathcal{F}$.
Evaluating both sides on $V \to X$ \'etale we obtain
$$
j_{p!}\mathcal{F}(V) =
\bigoplus\nolimits_{\varphi : V \to U} \mathcal{F}(V \xrightarrow{\varphi} U)
\quad\text{and}\quad
j_*\mathcal{F}(V) = \mathcal{F}(V \times_X U)
$$
For each $\varphi$ we have an open and closed immersion
$$
\Gamma_\varphi = (1, \varphi) : V \longrightarrow V \times_X U
$$
over $U$. It is open as it is a morphism between schemes \'etale
over $U$ and it is closed as it is a section of a scheme separated
over $V$ (Schemes, Lemma \ref{schemes-lemma-section-immersion}).
Thus for a section $s_\varphi \in \mathcal{F}(V \xrightarrow{\varphi} U)$
there exists a unique section $s'_\varphi$ in $\mathcal{F}(V \times_X U)$
which pulls back to $s_\varphi$ by $\Gamma_\varphi$
and which restricts to zero on the complement of the image of $\Gamma_\varphi$.

\medskip\noindent
To show that our map is injective suppose that
$\sum_{i = 1, \ldots, n} s_{\varphi_i}$ is an element of
$j_{p!}\mathcal{F}(V)$ in the formula above maps to zero
in $j_*\mathcal{F}(V)$. Our task is to show that
$\sum_{i = 1, \ldots, n} s_{\varphi_i}$ restricts to zero
on the members of an \'etale covering of $V$.
Looking at all pairwise equalizers
(which are open and closed in $V$) of the morphisms
$\varphi_i : V \to U$ and working locally on $V$, we
may assume the images of the morphisms
$\Gamma_{\varphi_1}, \ldots, \Gamma_{\varphi_n}$
are pairwise disjoint. Since our assumption is that
$\sum_{i = 1, \ldots, n} s'_{\varphi_i} = 0$
we then immediately conclude that $s'_{\varphi_i} = 0$
for each $i$ (by the disjointness of the supports of these
sections), whence $s_{\varphi_i} = 0$ for all $i$ as desired.
\end{proof}

\begin{lemma}
\label{lemma-shriek-equals-star-finite-etale}
Let $j : U \to X$ be finite and \'etale. Then the map
$j_! \to j_*$ of Lemma \ref{lemma-shriek-into-star-separated-etale}
is an isomorphism
on abelian sheaves and sheaves of $\Lambda$-modules.
\end{lemma}

\begin{proof}
It suffices to check $j_!\mathcal{F} \to j_*\mathcal{F}$
is an isomorphism \'etale locally on $X$.
Thus we may assume $U \to X$ is a finite disjoint union
of isomorphisms, see
\'Etale Morphisms, Lemma \ref{etale-lemma-finite-etale-etale-local}.
We omit the proof in this case.
\end{proof}

\begin{lemma}
\label{lemma-ses-associated-to-open}
Let $X$ be a scheme. Let $Z \subset X$ be a closed subscheme and let
$U \subset X$ be the complement. Denote $i : Z \to X$ and $j : U \to X$
the inclusion morphisms. For every abelian sheaf $\mathcal{F}$ on $X_\etale$
there is a canonical short exact sequence
$$
0 \to j_!j^{-1}\mathcal{F} \to \mathcal{F} \to i_*i^{-1}\mathcal{F} \to 0
$$
on $X_\etale$.
\end{lemma}

\begin{proof}
We obtain the maps by the adjointness properties of the functors
involved. For a geometric point $\overline{x}$ in $X$ we have either
$\overline{x} \in U$ in which case the map on the left hand side
is an isomorphism on stalks and the stalk of $i_*i^{-1}\mathcal{F}$
is zero or $\overline{x} \in Z$ in which case the map on the right hand side
is an isomorphism on stalks and the stalk of $j_!j^{-1}\mathcal{F}$
is zero. Here we have used the description of stalks of
Lemma \ref{lemma-stalk-pushforward-closed-immersion} and
Proposition \ref{proposition-describe-jshriek}.
\end{proof}

\begin{lemma}
\label{lemma-compatible-shriek-push-finite}
Consider a cartesian diagram of schemes
$$
\xymatrix{
U \ar[d]_g \ar[r]_{j'} & X \ar[d]^f \\
V \ar[r]^j & Y
}
$$
where $f$ is finite and $j$ is an open immersion.
Then $f_* \circ j'_! = j_! \circ g_*$ as functors
$\textit{Ab}(U_\etale) \to \textit{Ab}(Y_\etale)$.
\end{lemma}

\begin{proof}
Let $\mathcal{F}$ be an object of $\textit{Ab}(U_\etale)$.
Let $\overline{y}$ be a geometric point of $Y$ not contained
in the open $V$. Then
$$
(f_*j'_!\mathcal{F})_{\overline{y}} =
\bigoplus\nolimits_{\overline{x},\ f(\overline{x}) = \overline{y}}
(j'_!\mathcal{F})_{\overline{x}} = 0
$$
by Proposition \ref{proposition-finite-higher-direct-image-zero} and
because the stalk of $j'_!\mathcal{F}$ at $\overline{x} \not \in U$
are zero by Lemma \ref{lemma-jshriek-open}. On the other hand, we have
$$
j^{-1}f_*j'_!\mathcal{F} =
g_*(j')^{-1}j'_!\mathcal{F} =
g_*\mathcal{F}
$$
by Lemmas \ref{lemma-finite-pushforward-commutes-with-base-change} and
\ref{lemma-jshriek-open}.
Hence by the characterization of
$j_!$ in Lemma \ref{lemma-jshriek-open} we see that
$f_*j'_!\mathcal{F} = j_!g_*\mathcal{F}$.
We omit the verification that this identification
is functorial in $\mathcal{F}$.
\end{proof}







\section{Constructible sheaves}
\label{section-constructible}

\noindent
Let $X$ be a scheme. A {\it constructible locally closed subscheme} of $X$
is a locally closed subscheme $T \subset X$ such that the underlying
topological space of $T$ is a constructible subset of $X$.
If $T, T' \subset X$ are locally
closed subschemes with the same underlying topological space, then
$T_\etale \cong T'_\etale$ by the topological
invariance of the \'etale site (Theorem \ref{theorem-topological-invariance}).
Thus in the following definition we may assume our locally closed
subschemes are reduced.

\begin{definition}
\label{definition-constructible}
Let $X$ be a scheme.
\begin{enumerate}
\item A sheaf of sets on $X_\etale$ is {\it constructible}
if for every affine open $U \subset X$ there exists a finite decomposition
of $U$ into constructible locally closed subschemes $U = \coprod_i U_i$
such that $\mathcal{F}|_{U_i}$ is finite locally constant for all $i$.
\item A sheaf of abelian groups on $X_\etale$ is {\it constructible}
if for every affine open $U \subset X$ there exists a finite decomposition
of $U$ into constructible locally closed subschemes $U = \coprod_i U_i$
such that $\mathcal{F}|_{U_i}$ is finite locally constant for all $i$.
\item Let $\Lambda$ be a Noetherian ring. A sheaf of $\Lambda$-modules
on $X_\etale$ is {\it constructible} if for every affine open
$U \subset X$ there exists a finite decomposition
of $U$ into constructible locally closed subschemes
$U = \coprod_i U_i$ such that
$\mathcal{F}|_{U_i}$ is of finite type and locally constant for all $i$.
\end{enumerate}
\end{definition}

\noindent
It seems that this is the accepted definition. An alternative, which lends
itself more readily to generalizations beyond the \'etale site of a scheme,
would have been to define constructible sheaves by starting with
$h_U$, $j_{U!}\mathbf{Z}/n\mathbf{Z}$, and $j_{U!}\underline{\Lambda}$
where $U$ runs over all quasi-compact and quasi-separated objects
of $X_\etale$, and then take the smallest full subcategory of
$\Sh(X_\etale)$, $\textit{Ab}(X_\etale)$, and
$\textit{Mod}(X_\etale, \underline{\Lambda})$ containing these
and closed under finite limits and colimits. It follows from
Lemma \ref{lemma-constructible-abelian}
and
Lemmas \ref{lemma-category-constructible-sets},
\ref{lemma-category-constructible-abelian}, and
\ref{lemma-category-constructible-modules}
that this produces the same category if $X$ is quasi-compact and
quasi-separated. In general this does not produce the same
category however.

\medskip\noindent
A disjoint union decomposition $U = \coprod U_i$ of a scheme by
locally closed subschemes will be called a {\it partition} of $U$
(compare with Topology, Section \ref{topology-section-stratifications}).

\begin{lemma}
\label{lemma-constructible-quasi-compact-quasi-separated}
Let $X$ be a quasi-compact and quasi-separated scheme. Let $\mathcal{F}$
be a sheaf of sets on $X_\etale$. The following are equivalent
\begin{enumerate}
\item $\mathcal{F}$ is constructible,
\item there exists an open covering $X = \bigcup U_i$ such that
$\mathcal{F}|_{U_i}$ is constructible, and
\item there exists a partition $X = \bigcup X_i$ by constructible
locally closed subschemes such that $\mathcal{F}|_{X_i}$ is finite
locally constant.
\end{enumerate}
A similar statement holds for abelian sheaves and sheaves of
$\Lambda$-modules if $\Lambda$ is Noetherian.
\end{lemma}

\begin{proof}
It is clear that (1) implies (2).

\medskip\noindent
Assume (2). For every $x \in X$ we can find an $i$ and an affine open
neighbourhood $V_x \subset U_i$ of $x$. Hence we can find a finite
affine open covering $X = \bigcup V_j$ such that for each $j$ there
exists a finite decomposition $V_j = \coprod V_{j, k}$ by locally closed
constructible subsets such that $\mathcal{F}|_{V_{j, k}}$ is finite
locally constant. By Topology, Lemma
\ref{topology-lemma-quasi-compact-open-immersion-constructible-image}
each $V_{j, k}$ is constructible as a subset of $X$.
By Topology, Lemma
\ref{topology-lemma-constructible-partition-refined-by-stratification}
we can find a finite stratification $X = \coprod X_l$ with constructible
locally closed strata such that each
$V_{j, k}$ is a union of $X_l$. Thus (3) holds.

\medskip\noindent
Assume (3) holds. Let $U \subset X$ be an affine open.
Then $U \cap X_i$ is a constructible locally closed subset of $U$
(for example by Properties, Lemma \ref{properties-lemma-locally-constructible})
and $U = \coprod U \cap X_i$ is a partition of $U$ as in
Definition \ref{definition-constructible}. Thus (1) holds.
\end{proof}

\begin{lemma}
\label{lemma-constructible-constructible}
Let $X$ be a quasi-compact and quasi-separated scheme.
Let $\mathcal{F}$ be a sheaf of sets, abelian groups,
$\Lambda$-modules (with $\Lambda$ Noetherian) on $X_\etale$.
If there exist constructible locally closed subschemes $T_i \subset X$
such that (a) $X = \bigcup T_j$ and (b) $\mathcal{F}|_{T_j}$ is
constructible, then $\mathcal{F}$ is constructible.
\end{lemma}

\begin{proof}
First, we can assume the covering is finite as $X$ is quasi-compact
in the spectral topology
(Topology, Lemma \ref{topology-lemma-constructible-hausdorff-quasi-compact}
and
Properties, Lemma
\ref{properties-lemma-quasi-compact-quasi-separated-spectral}).
Observe that each $T_i$ is a quasi-compact and quasi-separated
scheme in its own right (because it is constructible in $X$; details omitted).
Thus we can find a finite partition $T_i = \coprod T_{i, j}$ into
locally closed constructible parts of $T_i$ such that
$\mathcal{F}|_{T_{i, j}}$ is finite locally constant
(Lemma \ref{lemma-constructible-quasi-compact-quasi-separated}).
By Topology, Lemma \ref{topology-lemma-constructible-in-constructible}
we see that $T_{i, j}$ is a constructible locally closed subscheme of $X$.
Then we can apply Topology, Lemma
\ref{topology-lemma-constructible-partition-refined-by-stratification}
to $X = \bigcup T_{i, j}$ to find the desired partition of $X$.
\end{proof}

\begin{lemma}
\label{lemma-constructible-local}
Let $X$ be a scheme. Checking constructibility of a sheaf
of sets, abelian groups, $\Lambda$-modules (with $\Lambda$ Noetherian)
can be done Zariski locally on $X$.
\end{lemma}

\begin{proof}
The statement means if $X = \bigcup U_i$ is an open covering
such that $\mathcal{F}|_{U_i}$ is constructible, then $\mathcal{F}$
is constructible. If $U \subset X$ is affine open, then
$U = \bigcup U \cap U_i$ and $\mathcal{F}|_{U \cap U_i}$ is constructible
(it is trivial that the restriction of a constructible sheaf to
an open is constructible). It follows from
Lemma \ref{lemma-constructible-quasi-compact-quasi-separated}
that $\mathcal{F}|_U$ is constructible, i.e., a suitable partition
of $U$ exists.
\end{proof}

\begin{lemma}
\label{lemma-pullback-constructible}
Let $f : X \to Y$ be a morphism of schemes. If $\mathcal{F}$ is a
constructible sheaf of sets, abelian groups, or $\Lambda$-modules
(with $\Lambda$ Noetherian) on $Y_\etale$, the same
is true for $f^{-1}\mathcal{F}$ on $X_\etale$.
\end{lemma}

\begin{proof}
By Lemma \ref{lemma-constructible-local} this reduces to the case
where $X$ and $Y$ are affine. By
Lemma \ref{lemma-constructible-quasi-compact-quasi-separated}
it suffices to find a finite partition of $X$ by constructible
locally closed subschemes such that $f^{-1}\mathcal{F}$ is finite locally
constant on each of them.
To find it we just pull back the partition of $Y$ adapted to
$\mathcal{F}$ and use
Lemma \ref{lemma-pullback-locally-constant}.
\end{proof}

\begin{lemma}
\label{lemma-constructible-abelian}
Let $X$ be a scheme.
\begin{enumerate}
\item The category of constructible sheaves of sets
is closed under finite limits and colimits inside $\Sh(X_\etale)$.
\item The category of constructible abelian sheaves is a
weak Serre subcategory of $\textit{Ab}(X_\etale)$.
\item Let $\Lambda$ be a Noetherian ring. The category of
constructible sheaves of $\Lambda$-modules on
$X_\etale$ is a weak Serre subcategory of
$\textit{Mod}(X_\etale, \Lambda)$.
\end{enumerate}
\end{lemma}

\begin{proof}
We prove (3). We will use the criterion of
Homology, Lemma \ref{homology-lemma-characterize-weak-serre-subcategory}.
Suppose that $\varphi : \mathcal{F} \to \mathcal{G}$
is a map of constructible sheaves of $\Lambda$-modules.
We have to show that $\mathcal{K} = \Ker(\varphi)$ and
$\mathcal{Q} = \Coker(\varphi)$ are constructible. 
Similarly, suppose that
$0 \to \mathcal{F} \to \mathcal{E} \to \mathcal{G} \to 0$
is a short exact sequence of sheaves of $\Lambda$-modules
with $\mathcal{F}$, $\mathcal{G}$ constructible. We have to show
that $\mathcal{E}$ is constructible.
In both cases we can replace $X$ with the members of an
affine open covering. Hence we may assume $X$ is affine.
Then we may further replace $X$ by the members of a finite
partition of $X$ by constructible locally closed subschemes
on which $\mathcal{F}$ and $\mathcal{G}$ are of finite type and
locally constant. Thus we may apply
Lemma \ref{lemma-kernel-finite-locally-constant} to conclude.

\medskip\noindent
The proofs of (1) and (2) are very similar and are omitted.
\end{proof}

\begin{lemma}
\label{lemma-support-constructible}
Let $X$ be a quasi-compact and quasi-separated scheme.
\begin{enumerate}
\item Let $\mathcal{F} \to \mathcal{G}$ be a map of constructible
sheaves of sets on $X_\etale$. Then the set of points $x \in X$
where $\mathcal{F}_{\overline{x}} \to \mathcal{G}_{\overline{x}}$
is surjective, resp.\ injective, resp.\ is isomorphic to a given map
of sets, is constructible in $X$.
\item Let $\mathcal{F}$ be a constructible abelian sheaf on $X_\etale$.
The support of $\mathcal{F}$ is constructible.
\item Let $\Lambda$ be a Noetherian ring.
Let $\mathcal{F}$ be a constructible sheaf of $\Lambda$-modules on $X_\etale$.
The support of $\mathcal{F}$ is constructible.
\end{enumerate}
\end{lemma}

\begin{proof}
Proof of (1).
Let $X = \coprod X_i$ be a partition of $X$ by locally closed constructible
subschemes such that both $\mathcal{F}$ and $\mathcal{G}$ are
finite locally constant over the parts (use
Lemma \ref{lemma-constructible-quasi-compact-quasi-separated}
for both $\mathcal{F}$ and $\mathcal{G}$ and choose a common
refinement). Then apply Lemma \ref{lemma-morphism-locally-constant}
to the restriction of the map to each part.

\medskip\noindent
The proof of (2) and (3) is omitted.
\end{proof}

\noindent
The following lemma will turn out to be very useful later on.
It roughly says that the category of constructible sheaves
has a kind of weak ``Noetherian'' property.

\begin{lemma}
\label{lemma-colimit-constructible}
Let $X$ be a quasi-compact and quasi-separated scheme. Let
$\mathcal{F} = \colim_{i \in I} \mathcal{F}_i$ be a filtered
colimit of sheaves of sets, abelian sheaves, or sheaves of modules.
\begin{enumerate}
\item If $\mathcal{F}$ and $\mathcal{F}_i$ are constructible sheaves of
sets, then the ind-object $\mathcal{F}_i$ is essentially constant with
value $\mathcal{F}$.
\item If $\mathcal{F}$ and $\mathcal{F}_i$ are constructible sheaves of
abelian groups, then the ind-object $\mathcal{F}_i$ is essentially constant
with value $\mathcal{F}$.
\item Let $\Lambda$ be a Noetherian ring.
If $\mathcal{F}$ and $\mathcal{F}_i$ are constructible sheaves of
$\Lambda$-modules, then the ind-object $\mathcal{F}_i$ is essentially constant
with value $\mathcal{F}$.
\end{enumerate}
\end{lemma}

\begin{proof}
Proof of (1). We will use without further mention that finite limits
and colimits of constructible sheaves are constructible
(Lemma \ref{lemma-kernel-finite-locally-constant}).
For each $i$ let $T_i \subset X$ be the set of points $x \in X$
where $\mathcal{F}_{i, \overline{x}} \to \mathcal{F}_{\overline{x}}$
is not surjective. Because $\mathcal{F}_i$ and $\mathcal{F}$ are
constructible $T_i$ is a constructible subset of $X$
(Lemma \ref{lemma-support-constructible}).
Since the stalks of $\mathcal{F}$ are finite
and since $\mathcal{F} = \colim_{i \in I} \mathcal{F}_i$ we see
that for all $x \in X$ we have $x \not \in T_i$ for $i$ large enough.
Since $X$ is a spectral space by Properties, Lemma
\ref{properties-lemma-quasi-compact-quasi-separated-spectral}
the constructible topology on $X$ is quasi-compact by
Topology, Lemma \ref{topology-lemma-constructible-hausdorff-quasi-compact}.
Thus $T_i = \emptyset$ for $i$ large enough. Thus
$\mathcal{F}_i \to \mathcal{F}$ is surjective for $i$ large enough.
Assume now that $\mathcal{F}_i \to \mathcal{F}$ is surjective for all $i$.
Choose $i \in I$. For $i' \geq i$ denote $S_{i'} \subset X$ the set of
points $x$ such that the number of elements in
$\Im(\mathcal{F}_{i, \overline{x}} \to \mathcal{F}_{\overline{x}})$
is not equal to the number of elements in
$\Im(\mathcal{F}_{i, \overline{x}} \to \mathcal{F}_{i', \overline{x}})$.
Because $\mathcal{F}_i$, $\mathcal{F}_{i'}$ and $\mathcal{F}$ are
constructible $S_{i'}$ is a constructible subset of $X$
(details omitted; hint: use Lemma \ref{lemma-support-constructible}).
Since the stalks of $\mathcal{F}_i$ and $\mathcal{F}$
are finite and since $\mathcal{F} = \colim_{i' \geq i} \mathcal{F}_{i'}$
we see that for all $x \in X$ we have $x \not \in S_{i'}$ for $i'$
large enough. By the same argument as above we can find a large $i'$ such
that $S_{i'} = \emptyset$. Thus $\mathcal{F}_i \to \mathcal{F}_{i'}$
factors through $\mathcal{F}$ as desired.

\medskip\noindent
Proof of (2). Observe that a constructible abelian sheaf is a constructible
sheaf of sets. Thus case (2) follows from (1).

\medskip\noindent
Proof of (3). We will use without further mention that the category of
constructible sheaves of $\Lambda$-modules is abelian
(Lemma \ref{lemma-kernel-finite-locally-constant}).
For each $i$ let $\mathcal{Q}_i$ be the cokernel of the map
$\mathcal{F}_i \to \mathcal{F}$. The support $T_i$ of $\mathcal{Q}_i$
is a constructible subset of $X$ as $\mathcal{Q}_i$ is constructible
(Lemma \ref{lemma-support-constructible}).
Since the stalks of $\mathcal{F}$ are finite $\Lambda$-modules
and since $\mathcal{F} = \colim_{i \in I} \mathcal{F}_i$ we see
that for all $x \in X$ we have $x \not \in T_i$ for $i$ large enough.
Since $X$ is a spectral space by Properties, Lemma
\ref{properties-lemma-quasi-compact-quasi-separated-spectral}
the constructible topology on $X$ is quasi-compact by
Topology, Lemma \ref{topology-lemma-constructible-hausdorff-quasi-compact}.
Thus $T_i = \emptyset$ for $i$ large enough. This proves the first
assertion. For the second, assume now that
$\mathcal{F}_i \to \mathcal{F}$ is surjective for all $i$.
Choose $i \in I$. For $i' \geq i$ denote $\mathcal{K}_{i'}$ the
image of $\Ker(\mathcal{F}_i \to \mathcal{F})$ in $\mathcal{F}_{i'}$.
The support $S_{i'}$ of $\mathcal{K}_{i'}$
is a constructible subset of $X$ as $\mathcal{K}_{i'}$ is constructible.
Since the stalks of $\Ker(\mathcal{F}_i \to \mathcal{F})$
are finite $\Lambda$-modules and since
$\mathcal{F} = \colim_{i' \geq i} \mathcal{F}_{i'}$ we see
that for all $x \in X$ we have $x \not \in S_{i'}$ for $i'$ large enough.
By the same argument as above we can find a large $i'$ such
that $S_{i'} = \emptyset$. Thus $\mathcal{F}_i \to \mathcal{F}_{i'}$
factors through $\mathcal{F}$ as desired.
\end{proof}

\begin{lemma}
\label{lemma-tensor-product-constructible}
Let $X$ be a scheme. Let $\Lambda$ be a Noetherian ring.
The tensor product of two constructible sheaves of $\Lambda$-modules
on $X_\etale$ is a constructible sheaf of $\Lambda$-modules.
\end{lemma}

\begin{proof}
The question immediately reduces to the case where $X$ is affine.
Since any two partitions of $X$ with constructible locally
closed strata have a common refinement of the same type and
since pullbacks commute with tensor product we reduce to
Lemma \ref{lemma-tensor-product-locally-constant}.
\end{proof}

\begin{lemma}
\label{lemma-tensor-constructible}
Let $\Lambda \to \Lambda'$ be a homomorphism of Noetherian rings.
Let $X$ be a scheme. Let $\mathcal{F}$ be a constructible
sheaf of $\Lambda$-modules on $X_\etale$. Then
$\mathcal{F} \otimes_{\underline{\Lambda}} \underline{\Lambda'}$
is a constructible sheaf of $\Lambda'$-modules.
\end{lemma}

\begin{proof}
Omitted. Hint: affine locally you can use the same stratification.
\end{proof}










\section{Auxiliary lemmas on morphisms}
\label{section-stratify-morphisms}

\noindent
Some lemmas that are useful for proving functoriality properties
of constructible sheaves.

\begin{lemma}
\label{lemma-etale-stratified-finite}
Let $U \to X$ be an \'etale morphism of quasi-compact and quasi-separated
schemes (for example an \'etale morphism of Noetherian schemes). Then there
exists a partition $X = \coprod_i X_i$ by constructible locally closed
subschemes such that $X_i \times_X U \to X_i$ is finite \'etale for all $i$.
\end{lemma}

\begin{proof}
If $U \to X$ is separated, then this is
More on Morphisms, Lemma \ref{more-morphisms-lemma-stratify-flat-fp-qf}.
In general, we may assume $X$ is affine. Choose a finite affine open
covering $U = \bigcup U_j$. Apply the previous case to all the morphisms
$U_j \to X$ and $U_j \cap U_{j'} \to X$ and choose a common
refinement $X = \coprod X_i$ of the resulting partitions.
After refining the partition further we may assume $X_i$ affine as well.
Fix $i$ and set $V = U \times_X X_i$. The morphisms
$V_j = U_j \times_X X_i \to X_i$ and
$V_{jj'} = (U_j \cap U_{j'}) \times_X X_i \to X_i$ are finite \'etale.
Hence $V_j$ and $V_{jj'}$ are affine schemes and $V_{jj'} \subset V_j$
is closed as well as open (since $V_{jj'} \to X_i$ is proper, so
Morphisms, Lemma \ref{morphisms-lemma-image-proper-scheme-closed}
applies). Then $V = \bigcup V_j$ is separated because
$\mathcal{O}(V_j) \to \mathcal{O}(V_{jj'})$ is surjective, see
Schemes, Lemma \ref{schemes-lemma-characterize-separated}.
Thus the previous case applies to $V \to X_i$ and we can further
refine the partition if needed (it actually isn't but we don't
need this).
\end{proof}

\noindent
In the Noetherian case one can prove the preceding lemma by
Noetherian induction and the following amusing lemma.

\begin{lemma}
\label{lemma-generically-finite}
Let $f: X \to Y$ be a morphism of schemes which is quasi-compact,
quasi-separated, and locally of finite type. If $\eta$ is a generic point
of an irreducible component of $Y$ such that $f^{-1}(\eta)$ is finite, then
there exists an open $V \subset Y$ containing $\eta$ such that
$f^{-1}(V) \to V$ is finite.
\end{lemma}

\begin{proof}
This is Morphisms, Lemma \ref{morphisms-lemma-generically-finite}.
\end{proof}

\noindent
The statement of the following lemma can be strengthened a bit.

\begin{lemma}
\label{lemma-decompose-quasi-finite-morphism}
Let $f : Y \to X$ be a quasi-finite and finitely presented
morphism of affine schemes.
\begin{enumerate}
\item There exists a surjective morphism of affine schemes $X' \to X$ and a
closed subscheme $Z' \subset Y' = X' \times_X Y$ such that
\begin{enumerate}
\item $Z' \subset Y'$ is a thickening, and
\item $Z' \to X'$ is a finite \'etale morphism.
\end{enumerate}
\item There exists a finite partition $X = \coprod X_i$ by
locally closed, constructible, affine strata, and surjective finite locally
free morphisms $X'_i \to X_i$ such that the reduction of
$Y'_i = X'_i \times_X Y \to X'_i$ is isomorphic to
$\coprod_{j = 1}^{n_i} (X'_i)_{red} \to (X'_i)_{red}$ for some $n_i$.
\end{enumerate}
\end{lemma}

\begin{proof}
Setting $X' = \coprod X'_i$ we see that (2) implies (1).
Write $X = \Spec(A)$ and $Y = \Spec(B)$. Write $A$ as a filtered colimit
of finite type $\mathbf{Z}$-algebras $A_i$. Since $B$ is an $A$-algebra of
finite presentation, we see that there exists $0 \in I$ and a
finite type ring map $A_0 \to B_0$ such that $B = \colim B_i$ with
$B_i = A_i \otimes_{A_0} B_0$, see
Algebra, Lemma \ref{algebra-lemma-colimit-category-fp-algebras}.
For $i$ sufficiently large we see that $A_i \to B_i$ is
quasi-finite, see Limits, Lemma \ref{limits-lemma-descend-quasi-finite}.
Thus we reduce to the case of finite type algebras over $\mathbf{Z}$,
in particular we reduce to the Noetherian case. (Details omitted.)

\medskip\noindent
Assume $X$ and $Y$ Noetherian. In this case any locally closed
subset of $X$ is constructible. By Lemma \ref{lemma-generically-finite}
and Noetherian induction we see that
there is a finite partition $X = \coprod X_i$ of $X$
by locally closed strata such that $Y \times_X X_i \to X_i$ is finite.
We can refine this partition to get affine strata.
Thus after replacing $X$ by $X' = \coprod X_i$ we may assume
$Y \to X$ is finite.

\medskip\noindent
Assume $X$ and $Y$ Noetherian and $Y \to X$ finite.
Suppose that we can prove (2) after base change by a surjective,
flat, quasi-finite morphism $U \to X$. Thus we have a partition
$U = \coprod U_i$ and finite locally free morphisms $U'_i \to U_i$
such that $U'_i \times_X Y \to U'_i$ is isomorphic to
$\coprod_{j = 1}^{n_i} (U'_i)_{red} \to (U'_i)_{red}$ for some $n_i$.
Then, by the argument in the previous paragraph, we can find a
partition $X = \coprod X_j$ with locally closed affine strata such that
$X_j \times_X U_i \to X_j$ is finite for all $i, j$. By
Morphisms, Lemma \ref{morphisms-lemma-finite-flat}
each $X_j \times_X U_i \to X_j$ is finite locally free.
Hence $X_j \times_X U'_i \to X_j$ is finite locally free
(Morphisms, Lemma \ref{morphisms-lemma-composition-finite-locally-free}).
It follows that $X = \coprod X_j$ and $X_j' = \coprod_i X_j \times_X U'_i$
is a solution for $Y \to X$. Thus it suffices to prove
the result (in the Noetherian case) after a surjective flat quasi-finite
base change.

\medskip\noindent
Applying Morphisms, Lemma \ref{morphisms-lemma-massage-finite}
we see we may assume that $Y$ is a closed subscheme of an
affine scheme $Z$ which is (set theoretically) a finite union
$Z = \bigcup_{i \in I} Z_i$ of closed subschemes mapping isomorphically
to $X$. In this case we will find a finite partition of $X = \coprod X_j$
with affine locally closed strata that works (in other words $X'_j = X_j$).
Set $T_i = Y \cap Z_i$. This is a closed subscheme of $X$.
As $X$ is Noetherian we can find a finite partition of $X = \coprod X_j$
by affine locally closed subschemes, such that each
$X_j \times_X T_i$ is (set theoretically) a union of strata $X_j \times_X Z_i$.
Replacing $X$ by $X_j$ we see that we may assume $I = I_1 \amalg I_2$
with $Z_i \subset Y$ for $i \in I_1$ and $Z_i \cap Y = \emptyset$ for
$i \in I_2$. Replacing $Z$ by $\bigcup_{i \in I_1} Z_i$ we see that we
may assume $Y = Z$.
Finally, we can replace $X$ again by the members of a partition
as above such that for every $i, i' \subset I$ the intersection
$Z_i \cap Z_{i'}$ is either empty or (set theoretically) equal
to $Z_i$ and $Z_{i'}$. This clearly means that $Y$ is (set theoretically)
equal to a disjoint union of the $Z_i$ which is what we wanted to show.
\end{proof}








\section{More on constructible sheaves}
\label{section-more-constructible}

\noindent
Let $\Lambda$ be a Noetherian ring. Let $X$ be a scheme.
We often consider $X_\etale$ as a ringed site with
sheaf of rings $\underline{\Lambda}$. In case of abelian sheaves
we often take $\Lambda = \mathbf{Z}/n\mathbf{Z}$ for a suitable
integer $n$.

\begin{lemma}
\label{lemma-jshriek-constructible}
Let $j : U \to X$ be an \'etale morphism of quasi-compact and
quasi-separated schemes.
\begin{enumerate}
\item The sheaf $h_U$ is a constructible sheaf of sets.
\item The sheaf $j_!\underline{M}$ is a constructible abelian sheaf
for a finite abelian group $M$.
\item If $\Lambda$ is a Noetherian ring and $M$ is a finite $\Lambda$-module,
then $j_!\underline{M}$ is a constructible sheaf of $\Lambda$-modules
on $X_\etale$.
\end{enumerate}
\end{lemma}

\begin{proof}
By Lemma \ref{lemma-etale-stratified-finite} there is a partition
$\coprod_i X_i$ such that $\pi_i : j^{-1}(X_i) \to X_i$ is finite \'etale.
The restriction of $h_U$ to $X_i$ is $h_{j^{-1}(X_i)}$ which is finite
locally constant by Lemma \ref{lemma-characterize-finite-locally-constant}.
For cases (2) and (3) we note that
$$
j_!(\underline{M})|_{X_i} =
\pi_{i!}(\underline{M}) =
\pi_{i*}(\underline{M})
$$
by Lemmas \ref{lemma-shriek-base-change} and
\ref{lemma-shriek-equals-star-finite-etale}.
Thus it suffices to show the lemma for $\pi : Y \to X$ finite \'etale.
This is Lemma \ref{lemma-pushforward-locally-constant}.
\end{proof}

\begin{lemma}
\label{lemma-torsion-colimit-constructible}
Let $X$ be a quasi-compact and quasi-separated scheme.
\begin{enumerate}
\item Let $\mathcal{F}$ be a sheaf of sets on $X_\etale$.
Then $\mathcal{F}$ is a filtered colimit of constructible
sheaves of sets.
\item Let $\mathcal{F}$ be a torsion abelian sheaf on $X_\etale$.
Then $\mathcal{F}$ is a filtered colimit of constructible abelian sheaves.
\item Let $\Lambda$ be a Noetherian ring and $\mathcal{F}$ a sheaf
of $\Lambda$-modules on $X_\etale$. Then
$\mathcal{F}$ is a filtered colimit of constructible sheaves of
$\Lambda$-modules.
\end{enumerate}
\end{lemma}

\begin{proof}
Let $\mathcal{B}$ be the collection of quasi-compact and quasi-separated
objects of $X_\etale$. By Modules on Sites,
Lemma \ref{sites-modules-lemma-module-filtered-colimit-constructibles}
any sheaf of sets is a filtered colimit of sheaves of the form
$$
\text{Coequalizer}\left(
\xymatrix{
\coprod\nolimits_{j = 1, \ldots, m} h_{V_j}
\ar@<1ex>[r] \ar@<-1ex>[r] &
\coprod\nolimits_{i = 1, \ldots, n} h_{U_i}
}
\right)
$$
with $V_j$ and $U_i$ quasi-compact and quasi-separated objects
of $X_\etale$. By
Lemmas \ref{lemma-jshriek-constructible} and \ref{lemma-constructible-abelian}
these coequalizers are constructible. This proves (1).

\medskip\noindent
Let $\Lambda$ be a Noetherian ring.
By Modules on Sites,
Lemma \ref{sites-modules-lemma-module-filtered-colimit-constructibles}
$\Lambda$-modules $\mathcal{F}$ is a filtered colimit
of modules of the form
$$
\Coker\left(
\bigoplus\nolimits_{j = 1, \ldots, m} j_{V_j!}\underline{\Lambda}_{V_j}
\longrightarrow
\bigoplus\nolimits_{i = 1, \ldots, n} j_{U_i!}\underline{\Lambda}_{U_i}
\right)
$$
with $V_j$ and $U_i$ quasi-compact and quasi-separated objects
of $X_\etale$. By
Lemmas \ref{lemma-jshriek-constructible} and \ref{lemma-constructible-abelian}
these cokernels are constructible. This proves (3).

\medskip\noindent
Proof of (2). First write $\mathcal{F} = \bigcup \mathcal{F}[n]$ where
$\mathcal{F}[n]$ is the $n$-torsion subsheaf. Then we can view
$\mathcal{F}[n]$ as a sheaf of $\mathbf{Z}/n\mathbf{Z}$-modules
and apply (3).
\end{proof}

\begin{lemma}
\label{lemma-check-constructible}
Let $f : X \to Y$ be a surjective morphism of quasi-compact and
quasi-separated schemes.
\begin{enumerate}
\item Let $\mathcal{F}$ be a sheaf of sets on $Y_\etale$. Then $\mathcal{F}$
is constructible if and only if $f^{-1}\mathcal{F}$ is constructible.
\item Let $\mathcal{F}$ be an abelian sheaf on $Y_\etale$. Then $\mathcal{F}$
is constructible if and only if $f^{-1}\mathcal{F}$ is constructible.
\item Let $\Lambda$ be a Noetherian ring.
Let $\mathcal{F}$ be sheaf of $\Lambda$-modules on $Y_\etale$.
Then $\mathcal{F}$ is constructible if and only if $f^{-1}\mathcal{F}$
is constructible.
\end{enumerate}
\end{lemma}

\begin{proof}
One implication follows from Lemma \ref{lemma-pullback-constructible}.
For the converse, assume $f^{-1}\mathcal{F}$ is constructible.
Write $\mathcal{F} = \colim \mathcal{F}_i$ as a
filtered colimit of constructible sheaves (of sets, abelian groups, or modules)
using Lemma \ref{lemma-torsion-colimit-constructible}.
Since $f^{-1}$ is a left adjoint it commutes with colimits
(Categories, Lemma \ref{categories-lemma-adjoint-exact}) and we see that
$f^{-1}\mathcal{F} = \colim f^{-1}\mathcal{F}_i$.
By Lemma \ref{lemma-colimit-constructible} we see that
$f^{-1}\mathcal{F}_i \to f^{-1}\mathcal{F}$
is surjective for all $i$ large enough.
Since $f$ is surjective we conclude (by looking at stalks using
Lemma \ref{lemma-stalk-pullback} and
Theorem \ref{theorem-exactness-stalks})
that $\mathcal{F}_i \to \mathcal{F}$ is surjective for all $i$ large enough.
Thus $\mathcal{F}$ is the quotient of a constructible sheaf $\mathcal{G}$.
Applying the argument once more to
$\mathcal{G} \times_\mathcal{F} \mathcal{G}$ or
the kernel of $\mathcal{G} \to \mathcal{F}$
we conclude using that $f^{-1}$ is exact and that the category of
constructible sheaves (of sets, abelian groups, or modules) is
preserved under finite (co)limits or (co)kernels inside
$\Sh(Y_\etale)$, $\Sh(X_\etale)$, $\textit{Ab}(Y_\etale)$,
$\textit{Ab}(X_\etale)$, $\textit{Mod}(Y_\etale, \Lambda)$, and
$\textit{Mod}(X_\etale, \Lambda)$, see
Lemma \ref{lemma-constructible-abelian}.
\end{proof}

\begin{lemma}
\label{lemma-pushforward-constructible}
Let $f : X \to Y$ be a finite \'etale morphism of schemes. Let $\Lambda$ be a
Noetherian ring. If $\mathcal{F}$ is a constructible sheaf of sets,
constructible sheaf of abelian groups, or constructible sheaf of
$\Lambda$-modules on $X_\etale$, the same is true for
$f_*\mathcal{F}$ on $Y_\etale$.
\end{lemma}

\begin{proof}
By Lemma \ref{lemma-constructible-local} it suffices to check this
Zariski locally on $Y$ and by Lemma \ref{lemma-check-constructible}
we may replace $Y$ by an \'etale cover (the construction of $f_*$
commutes with \'etale localization). A finite \'etale morphism is
\'etale locally isomorphic to a disjoint union of isomorphisms, see
\'Etale Morphisms, Lemma \ref{etale-lemma-finite-etale-etale-local}.
Thus, in the case of sheaves of sets, the lemma says that if
$\mathcal{F}_i$, $i = 1, \ldots, n$ are constructible sheaves of sets, then
$\prod_{i = 1, \ldots, n} \mathcal{F}_i$ is too.
This is clear. Similarly for sheaves of abelian groups and modules.
\end{proof}

\begin{lemma}
\label{lemma-category-constructible-sets}
Let $X$ be a quasi-compact and quasi-separated scheme. The category of
constructible sheaves of sets is the full subcategory of $\Sh(X_\etale)$
consisting of sheaves $\mathcal{F}$ which are coequalizers
$$
\xymatrix{
\mathcal{F}_1
\ar@<1ex>[r] \ar@<-1ex>[r]
&
\mathcal{F}_0 \ar[r]
&
\mathcal{F}}
$$
such that $\mathcal{F}_i$, $i = 0, 1$ is a finite coproduct of sheaves of
the form $h_U$ with $U$ a quasi-compact and quasi-separated
object of $X_\etale$.
\end{lemma}

\begin{proof}
In the proof of Lemma \ref{lemma-torsion-colimit-constructible}
we have seen that sheaves of this form are constructible.
For the converse, suppose that for every constructible sheaf of
sets $\mathcal{F}$ we can find a surjection $\mathcal{F}_0 \to \mathcal{F}$
with $\mathcal{F}_0$ as in the lemma. Then we find our surjection
$\mathcal{F}_1 \to \mathcal{F}_0 \times_\mathcal{F} \mathcal{F}_0$
because the latter is constructible by Lemma \ref{lemma-constructible-abelian}.

\medskip\noindent
By Topology, Lemma
\ref{topology-lemma-constructible-partition-refined-by-stratification}
we may choose a finite stratification
$X = \coprod_{i \in I} X_i$ such that $\mathcal{F}$ is finite locally
constant on each stratum. We will prove the result by induction on
the cardinality of $I$. Let $i \in I$ be a minimal element in the
partial ordering of $I$. Then $X_i \subset X$ is closed.
By induction, there exist finitely many quasi-compact and quasi-separated
objects $U_\alpha$ of $(X \setminus X_i)_\etale$ and a surjective
map $\coprod h_{U_\alpha} \to \mathcal{F}|_{X \setminus X_i}$.
These determine a map
$$
\coprod h_{U_\alpha} \to \mathcal{F}
$$
which is surjective after restricting to $X \setminus X_i$. By
Lemma \ref{lemma-characterize-finite-locally-constant}
we see that $\mathcal{F}|_{X_i} = h_V$ for some scheme $V$ finite \'etale
over $X_i$. Let $\overline{v}$ be a geometric point of $V$ lying
over $\overline{x} \in X_i$. We may think of $\overline{v}$ as an
element of the stalk $\mathcal{F}_{\overline{x}} = V_{\overline{x}}$.
Thus we can find an \'etale neighbourhood $(U, \overline{u})$
of $\overline{x}$ and a section $s \in \mathcal{F}(U)$ whose stalk at
$\overline{x}$ gives $\overline{v}$. Thinking of $s$ as a map
$s : h_U \to \mathcal{F}$, restricting to $X_i$ we obtain a morphism
$s|_{X_i} : U \times_X X_i \to V$ over $X_i$ which maps $\overline{u}$
to $\overline{v}$. Since $V$ is quasi-compact (finite over the closed
subscheme $X_i$ of the quasi-compact scheme $X$) a finite number
$s^{(1)}, \ldots, s^{(m)}$ of these sections of $\mathcal{F}$ over
$U^{(1)}, \ldots, U^{(m)}$ will determine a jointly
surjective map
$$
\coprod s^{(j)}|_{X_i} : \coprod U^{(j)} \times_X X_i \longrightarrow V
$$
Then we obtain the surjection
$$
\coprod h_{U_\alpha} \amalg \coprod h_{U^{(j)}} \to \mathcal{F}
$$
as desired.
\end{proof}

\begin{lemma}
\label{lemma-category-constructible-modules}
Let $X$ be a quasi-compact and quasi-separated scheme. Let $\Lambda$
be a Noetherian ring. The category of constructible sheaves of
$\Lambda$-modules is exactly the category of modules of the form
$$
\Coker\left(
\bigoplus\nolimits_{j = 1, \ldots, m} j_{V_j!}\underline{\Lambda}_{V_j}
\longrightarrow
\bigoplus\nolimits_{i = 1, \ldots, n} j_{U_i!}\underline{\Lambda}_{U_i}
\right)
$$
with $V_j$ and $U_i$ quasi-compact and quasi-separated objects of
$X_\etale$. In fact, we can even assume $U_i$ and $V_j$ affine.
\end{lemma}

\begin{proof}
In the proof of Lemma \ref{lemma-torsion-colimit-constructible}
we have seen modules of this form are constructible. Since the
category of constructible modules is abelian
(Lemma \ref{lemma-constructible-abelian})
it suffices to prove that given a constructible module $\mathcal{F}$
there is a surjection
$$
\bigoplus\nolimits_{i = 1, \ldots, n} j_{U_i!}\underline{\Lambda}_{U_i}
\longrightarrow \mathcal{F}
$$
for some affine objects $U_i$ in $X_\etale$. By
Modules on Sites, Lemma
\ref{sites-modules-lemma-module-filtered-colimit-constructibles}
there is a surjection
$$
\Psi :
\bigoplus\nolimits_{i \in I} j_{U_i!}\underline{\Lambda}_{U_i}
\longrightarrow
\mathcal{F}
$$
with $U_i$ affine and the direct sum over a possibly infinite
index set $I$. For every finite subset $I' \subset I$ set
$$
T_{I'} = \text{Supp}(\Coker(
\bigoplus\nolimits_{i \in I'} j_{U_i!}\underline{\Lambda}_{U_i}
\longrightarrow \mathcal{F}))
$$
By the very definition of constructible sheaves, the set $T_{I'}$
is a constructible subset of $X$. We want to show that $T_{I'} = \emptyset$
for some $I'$. Since every stalk $\mathcal{F}_{\overline{x}}$ is
a finite type $\Lambda$-module and since $\Psi$ is surjective, for
every $x \in X$ there is an $I'$ such that $x \not \in T_{I'}$.
In other words we have
$\emptyset = \bigcap_{I' \subset I\text{ finite}} T_{I'}$. Since
$X$ is a spectral space by Properties, Lemma
\ref{properties-lemma-quasi-compact-quasi-separated-spectral}
the constructible topology on $X$ is quasi-compact by
Topology, Lemma \ref{topology-lemma-constructible-hausdorff-quasi-compact}.
Thus $T_{I'} = \emptyset$ for some $I' \subset I$ finite
as desired.
\end{proof}

\begin{lemma}
\label{lemma-category-constructible-abelian}
Let $X$ be a quasi-compact and quasi-separated scheme. The category of
constructible abelian sheaves is exactly the category of abelian
sheaves of the form
$$
\Coker\left(
\bigoplus\nolimits_{j = 1, \ldots, m}
j_{V_j!}\underline{\mathbf{Z}/m_j\mathbf{Z}}_{V_j}
\longrightarrow
\bigoplus\nolimits_{i = 1, \ldots, n}
j_{U_i!}\underline{\mathbf{Z}/n_i\mathbf{Z}}_{U_i}
\right)
$$
with $V_j$ and $U_i$ quasi-compact and quasi-separated objects of
$X_\etale$ and $m_j$, $n_i$ positive integers.
In fact, we can even assume $U_i$ and $V_j$ affine.
\end{lemma}

\begin{proof}
This follows from Lemma \ref{lemma-category-constructible-modules}
applied with $\Lambda = \mathbf{Z}/n\mathbf{Z}$
and the fact that, since $X$ is quasi-compact, every constructible
abelian sheaf is annihilated by some positive integer $n$ (details omitted).
\end{proof}

\begin{lemma}
\label{lemma-constructible-is-compact}
Let $X$ be a quasi-compact and quasi-separated scheme. Let $\Lambda$ be a
Noetherian ring. Let $\mathcal{F}$ be a constructible sheaf of sets, abelian
groups, or $\Lambda$-modules on $X_\etale$. Let
$\mathcal{G} = \colim \mathcal{G}_i$ be a filtered colimit of sheaves of
sets, abelian groups, or $\Lambda$-modules. Then
$$
\Mor(\mathcal{F}, \mathcal{G}) = \colim \Mor(\mathcal{F}, \mathcal{G}_i)
$$
in the category of sheaves of sets, abelian groups, or $\Lambda$-modules on
$X_\etale$.
\end{lemma}

\begin{proof}
The case of sheaves of sets. By Lemma \ref{lemma-category-constructible-sets}
it suffices to prove the lemma for $h_U$ where $U$ is a quasi-compact
and quasi-separated object of $X_\etale$. Recall that
$\Mor(h_U, \mathcal{G}) = \mathcal{G}(U)$. Hence the result
follows from Sites, Lemma \ref{sites-lemma-directed-colimits-sections}.

\medskip\noindent
In the case of abelian sheaves or sheaves of modules, the result follows
in the same way using
Lemmas \ref{lemma-category-constructible-abelian} and
\ref{lemma-category-constructible-modules}.
For the case of abelian sheaves, we add that
$\Mor(j_{U!}\underline{\mathbf{Z}/n\mathbf{Z}}, \mathcal{G})$
is equal to the $n$-torsion elements of $\mathcal{G}(U)$.
\end{proof}

\begin{lemma}
\label{lemma-finite-pushforward-constructible}
Let $f : X \to Y$ be a finite and finitely presented morphism of schemes.
Let $\Lambda$ be a Noetherian ring. If $\mathcal{F}$ is a constructible
sheaf of sets, abelian groups, or $\Lambda$-modules on $X_\etale$,
then $f_*\mathcal{F}$ is too.
\end{lemma}

\begin{proof}
It suffices to prove this when $X$ and $Y$ are affine by
Lemma \ref{lemma-constructible-local}.
By
Lemmas \ref{lemma-finite-pushforward-commutes-with-base-change} and
\ref{lemma-check-constructible} we may base change to any
affine scheme surjective over $X$. By
Lemma \ref{lemma-decompose-quasi-finite-morphism}
this reduces us to the case of a finite \'etale morphism
(because a thickening leads to an equivalence of \'etale topoi
and even small \'etale sites, see
Theorem \ref{theorem-topological-invariance}).
The finite \'etale case is
Lemma \ref{lemma-pushforward-constructible}.
\end{proof}

\begin{lemma}
\label{lemma-category-is-colimit}
Let $X = \lim_{i \in I} X_i$ be a limit of a directed
system of schemes with affine transition morphisms.
We assume that $X_i$ is quasi-compact and quasi-separated
for all $i \in I$.
\begin{enumerate}
\item The category of constructible sheaves of sets on $X_\etale$
is the colimit of the categories of constructible sheaves of sets
on $(X_i)_\etale$.
\item The category of constructible abelian sheaves on $X_\etale$
is the colimit of the categories of constructible abelian sheaves
on $(X_i)_\etale$.
\item Let $\Lambda$ be a Noetherian ring. The category of constructible
sheaves of $\Lambda$-modules on $X_\etale$ is the colimit of the
categories of constructible sheaves of $\Lambda$-modules on $(X_i)_\etale$.
\end{enumerate}
\end{lemma}

\begin{proof}
Proof of (1). Denote $f_i : X \to X_i$ the projection maps.
There are 3 parts to the proof corresponding to ``faithful'',
``fully faithful'', and ``essentially surjective''.

\medskip\noindent
Faithful. Choose $0 \in I$ and let $\mathcal{F}_0$, $\mathcal{G}_0$ be
constructible sheaves on $X_0$. Suppose that
$a, b : \mathcal{F}_0 \to \mathcal{G}_0$ are maps such that
$f_0^{-1}a = f_0^{-1}b$. Let $E \subset X_0$ be the set
of points $x \in X_0$ such that $a_{\overline{x}} = b_{\overline{x}}$.
By Lemma \ref{lemma-support-constructible} the subset
$E \subset X_0$ is constructible. By assumption $X \to X_0$ maps into $E$.
By Limits, Lemma \ref{limits-lemma-limit-contained-in-constructible}
we find an $i \geq 0$ such that $X_i \to X_0$ maps into $E$.
Hence $f_{i0}^{-1}a = f_{i0}^{-1}b$.

\medskip\noindent
Fully faithful. Choose $0 \in I$ and let $\mathcal{F}_0$, $\mathcal{G}_0$ be
constructible sheaves on $X_0$. Suppose that
$a : f_0^{-1}\mathcal{F}_0 \to f_0^{-1}\mathcal{G}_0$ is a map.
We claim there is an $i$ and a map
$a_i : f_{i0}^{-1}\mathcal{F}_0 \to f_{i0}^{-1}\mathcal{G}_0$
which pulls back to $a$ on $X$.
By Lemma \ref{lemma-category-constructible-sets}
we can replace $\mathcal{F}_0$
by a finite coproduct of sheaves represented by quasi-compact
and quasi-separated objects of $(X_0)_\etale$.
Thus we have to show: If $U_0 \to X_0$ is such an object
of $(X_0)_\etale$, then
$$
f_0^{-1}\mathcal{G}(U) = \colim_{i \geq 0} f_{i0}^{-1}\mathcal{G}(U_i)
$$
where $U = X \times_{X_0} U_0$ and $U_i = X_i \times_{X_0} U_0$.
This is a special case of Theorem \ref{theorem-colimit}.

\medskip\noindent
Essentially surjective. We have to show every constructible $\mathcal{F}$
on $X$ is isomorphic to $f_i^{-1}\mathcal{F}$ for some constructible
$\mathcal{F}_i$ on $X_i$. Applying
Lemma \ref{lemma-category-constructible-sets}
and using the results of the previous two paragraphs, we see that
it suffices to prove this for $h_U$ for some quasi-compact
and quasi-separated object $U$ of $X_\etale$.
In this case we have to show that $U$ is the base change of
a quasi-compact and quasi-separated scheme \'etale over $X_i$ for some $i$.
This follows from
Limits, Lemmas \ref{limits-lemma-descend-finite-presentation} and
\ref{limits-lemma-descend-etale}.

\medskip\noindent
Proof of (3). The argument is very similar to the argument for
sheaves of sets, but using
Lemma \ref{lemma-category-constructible-modules}
instead of
Lemma \ref{lemma-category-constructible-sets}. Details omitted.
Part (2) follows from part (3) because every constructible abelian
sheaf over a quasi-compact scheme is a constructible sheaf of
$\mathbf{Z}/n\mathbf{Z}$-modules for some $n$.
\end{proof}

\begin{lemma}
\label{lemma-category-loc-cst-is-colimit}
Let $X = \lim_{i \in I} X_i$ be a limit of a directed
system of schemes with affine transition morphisms.
We assume that $X_i$ is quasi-compact and quasi-separated
for all $i \in I$.
\begin{enumerate}
\item The category of finite locally constant sheaves on $X_\etale$
is the colimit of the categories of finite locally constant sheaves
on $(X_i)_\etale$.
\item The category of finite locally constant abelian sheaves on $X_\etale$
is the colimit of the categories of finite locally constant abelian sheaves
on $(X_i)_\etale$.
\item Let $\Lambda$ be a Noetherian ring. The category of finite type,
locally constant sheaves of $\Lambda$-modules on $X_\etale$
is the colimit of the categories of finite type, locally constant
sheaves of $\Lambda$-modules on $(X_i)_\etale$.
\end{enumerate}
\end{lemma}

\begin{proof}
By Lemma \ref{lemma-category-is-colimit} the functor in each case
is fully faithful. By the same lemma, all we have to show to finish
the proof in case (1) is the following: given a constructible sheaf
$\mathcal{F}_i$ on $X_i$ whose pullback $\mathcal{F}$ to $X$ is
finite locally constant, there exists an $i' \geq i$ such that the
pullback $\mathcal{F}_{i'}$ of $\mathcal{F}_i$ to $X_{i'}$
is finite locally constant. By assumption there exists an \'etale covering
$\mathcal{U} = \{U_j \to X\}_{j \in J}$
such that $\mathcal{F}|_{U_j} \cong \underline{S_j}$ for some finite set $S_j$.
We may assume $U_j$ is affine for all $j \in J$.
Since $X$ is quasi-compact, we may assume $J$ finite.
By Lemma \ref{lemma-colimit-affine-sites} we can find an $i' \geq i$
and an \'etale covering $\mathcal{U}_{i'} = \{U_{i', j} \to X_{i'}\}_{j \in J}$
whose base change to $X$ is $\mathcal{U}$. Then
$\mathcal{F}_{i'}|_{U_{i', j}}$ and $\underline{S_j}$ are
constructible sheaves on $(U_{i', j})_\etale$ whose pullbacks to
$U_j$ are isomorphic. Hence after increasing $i'$ we get
that $\mathcal{F}_{i'}|_{U_{i', j}}$ and $\underline{S_j}$
are isomorphic. Thus $\mathcal{F}_{i'}$ is finite locally constant.
The proof in cases (2) and (3) is exactly the same.
\end{proof}

\begin{lemma}
\label{lemma-irreducible-subsheaf-constant-zero}
Let $X$ be an irreducible scheme with generic point $\eta$.
\begin{enumerate}
\item Let $S' \subset S$ be an inclusion of sets. If we have
$\underline{S'} \subset \mathcal{G} \subset \underline{S}$
in $\Sh(X_\etale)$ and $S' = \mathcal{G}_{\overline{\eta}}$, then
$\mathcal{G} = \underline{S'}$.
\item Let $A' \subset A$ be an inclusion of abelian groups. If we have
$\underline{A'} \subset \mathcal{G} \subset \underline{A}$
in $\textit{Ab}(X_\etale)$ and $A' = \mathcal{G}_{\overline{\eta}}$, then
$\mathcal{G} = \underline{A'}$.
\item Let $M' \subset M$ be an inclusion of modules over a ring $\Lambda$.
If we have $\underline{M'} \subset \mathcal{G} \subset \underline{M}$
in $\textit{Mod}(X_\etale, \underline{\Lambda})$
and $M' = \mathcal{G}_{\overline{\eta}}$, then
$\mathcal{G} = \underline{M'}$.
\end{enumerate}
\end{lemma}

\begin{proof}
This is true because for every \'etale morphism $U \to X$
with $U \not = \emptyset$ the point $\eta$ is in the image.
\end{proof}

\begin{lemma}
\label{lemma-push-constant-sheaf-from-open}
Let $X$ be an integral normal scheme with function field $K$.
Let $E$ be a set.
\begin{enumerate}
\item Let $g : \Spec(K) \to X$ be the inclusion of the generic point.
Then $g_*\underline{E} = \underline{E}$.
\item Let $j : U \to X$ be the inclusion of a nonempty open. Then
$j_*\underline{E} = \underline{E}$.
\end{enumerate}
\end{lemma}

\begin{proof}
Proof of (1). Let $x \in X$ be a point. Let
$\mathcal{O}^{sh}_{X, \overline{x}}$
be a strict henselization of $\mathcal{O}_{X, x}$. 
By More on Algebra, Lemma \ref{more-algebra-lemma-henselization-normal}
we see that $\mathcal{O}^{sh}_{X, \overline{x}}$ is a normal domain.
Hence $\Spec(K) \times_X \Spec(\mathcal{O}^{sh}_{X, \overline{x}})$
is irreducible. It follows that
the stalk $(g_*\underline{E})_{\overline{x}}$ is equal to $E$,
see Theorem \ref{theorem-higher-direct-images}.

\medskip\noindent
Proof of (2). Since $g$ factors through $j$ there is a map
$j_*\underline{E} \to g_*\underline{E}$. This map is injective because
for every scheme $V$ \'etale over $X$ the set $\Spec(K) \times_X V$
is dense in $U \times_X V$. On the other hand, we have a map
$\underline{E} \to j_*\underline{E}$ and we conclude.
\end{proof}

\begin{lemma}
\label{lemma-zero-in-generic-point}
Let $X$ be a quasi-compact and quasi-separated scheme. Let
$\eta \in X$ be a generic point of an irreducible component of $X$.
\begin{enumerate}
\item Let $\mathcal{F}$ be a torsion abelian sheaf on $X_\etale$
whose stalk $\mathcal{F}_{\overline{\eta}}$ is zero.
Then $\mathcal{F} = \colim \mathcal{F}_i$ is a filtered colimit of
constructible abelian sheaves $\mathcal{F}_i$ such that for each $i$
the support of $\mathcal{F}_i$ is contained
in a closed subscheme not containing $\eta$.
\item Let $\Lambda$ be a Noetherian ring and $\mathcal{F}$ a sheaf
of $\Lambda$-modules on $X_\etale$ whose stalk
$\mathcal{F}_{\overline{\eta}}$ is zero. Then
$\mathcal{F} = \colim \mathcal{F}_i$
is a filtered colimit of constructible sheaves of
$\Lambda$-modules $\mathcal{F}_i$ such that for each $i$
the support of $\mathcal{F}_i$ is contained in a closed subscheme
not containing $\eta$.
\end{enumerate}
\end{lemma}

\begin{proof}
Proof of (1). We can write $\mathcal{F} = \colim_{i \in I} \mathcal{F}_i$
with $\mathcal{F}_i$ constructible abelian by
Lemma \ref{lemma-torsion-colimit-constructible}.
Choose $i \in I$. Since $\mathcal{F}|_\eta$ is zero by assumption, we
see that there exists an $i'(i) \geq i$ such that
$\mathcal{F}_i|_\eta \to \mathcal{F}_{i'(i)}|_\eta$ is zero, see
Lemma \ref{lemma-colimit-constructible}.
Then $\mathcal{G}_i = \Im(\mathcal{F}_i \to \mathcal{F}_{i'(i)})$
is a constructible abelian sheaf (Lemma \ref{lemma-constructible-abelian})
whose stalk at $\eta$ is zero.
Hence the support $E_i$ of $\mathcal{G}_i$ is a constructible
subset of $X$ not containing $\eta$. Since
$\eta$ is a generic point of an irreducible component of
$X$, we see that $\eta \not \in Z_i = \overline{E_i}$ by
Topology, Lemma \ref{topology-lemma-generic-point-in-constructible}.
Define a new directed set $I'$ by using the set $I$ with
ordering defined by the rule
$i_1$ is bigger or equal to $i_2$ if and only if $i_1 \geq i'(i_2)$.
Then the sheaves $\mathcal{G}_i$ form a system over $I'$
with colimit $\mathcal{F}$ and the proof is complete.

\medskip\noindent
The proof in case (2) is exactly the same and we omit it.
\end{proof}








\section{Constructible sheaves on Noetherian schemes}
\label{section-noetherian-constructible}

\noindent
If $X$ is a Noetherian scheme then any locally closed subset is a
constructible locally closed subset
(Topology, Lemma \ref{topology-lemma-constructible-Noetherian-space}).
Hence an abelian sheaf $\mathcal{F}$ on $X_\etale$ is constructible
if and only if there exists a finite partition $X = \coprod X_i$
such that $\mathcal{F}|_{X_i}$ is finite locally constant.
(By convention a partition of a topological space has locally
closed parts, see Topology, Section \ref{topology-section-stratifications}.)
In other words, we can omit the adjective ``constructible'' in
Definition \ref{definition-constructible}. Actually, the category
of constructible sheaves on
Noetherian schemes has some additional properties which we will
catalogue in this section.

\begin{proposition}
\label{proposition-constructible-over-noetherian}
Let $X$ be a Noetherian scheme. Let $\Lambda$ be a Noetherian ring.
\begin{enumerate}
\item Any sub or quotient sheaf of a constructible sheaf of sets
is constructible.
\item The category of constructible abelian sheaves on $X_\etale$ is a
(strong) Serre subcategory of $\textit{Ab}(X_\etale)$. In particular,
every sub and quotient sheaf of a constructible abelian sheaf
on $X_\etale$ is constructible.
\item The category of constructible sheaves of $\Lambda$-modules
on $X_\etale$ is a (strong) Serre subcategory of
$\textit{Mod}(X_\etale, \Lambda)$. In particular, every submodule
and quotient module of a constructible sheaf of $\Lambda$-modules
on $X_\etale$ is constructible.
\end{enumerate}
\end{proposition}

\begin{proof}
Proof of (1). Let $\mathcal{G} \subset \mathcal{F}$ with $\mathcal{F}$
a constructible sheaf of sets on $X_\etale$. Let $\eta \in X$ be a generic
point of an irreducible component of $X$. By Noetherian induction
it suffices to find an open neighbourhood $U$ of $\eta$ such that
$\mathcal{G}|_U$ is locally constant. To do this we may replace $X$
by an \'etale neighbourhood of $\eta$.
Hence we may assume $\mathcal{F}$ is constant and $X$ is irreducible.

\medskip\noindent
Say $\mathcal{F} = \underline{S}$ for some finite set $S$.
Then $S' = \mathcal{G}_{\overline{\eta}} \subset S$
say $S' = \{s_1, \ldots, s_t\}$.
Pick an \'etale neighbourhood $(U, \overline{u})$ of $\overline{\eta}$
and sections $\sigma_1, \ldots, \sigma_t \in \mathcal{G}(U)$ which map to
$s_i$ in $\mathcal{G}_{\overline{\eta}} \subset S$.
Since $\sigma_i$ maps to an element
$s_i \in S' \subset S = \Gamma(X, \mathcal{F})$
we see that the two pullbacks of $\sigma_i$ to $U \times_X U$
are the same as sections of $\mathcal{G}$. By the sheaf condition
for $\mathcal{G}$ we find that $\sigma_i$ comes from a section
of $\mathcal{G}$ over the open $\Im(U \to X)$ of $X$.
Shrinking $X$ we may assume
$\underline{S'} \subset \mathcal{G} \subset \underline{S}$.
Then we see that $\underline{S'} = \mathcal{G}$ by
Lemma \ref{lemma-irreducible-subsheaf-constant-zero}.

\medskip\noindent
Let $\mathcal{F} \to \mathcal{Q}$ be a surjection with $\mathcal{F}$
a constructible sheaf of sets on $X_\etale$. Then set
$\mathcal{G} = \mathcal{F} \times_\mathcal{Q} \mathcal{F}$.
By the first part of the proof we see that $\mathcal{G}$ is
constructible as a subsheaf of $\mathcal{F} \times \mathcal{F}$.
This in turn implies that $\mathcal{Q}$ is constructible, see
Lemma \ref{lemma-constructible-abelian}.

\medskip\noindent
Proof of (3). we already know that constructible sheaves of modules
form a weak Serre subcategory, see Lemma \ref{lemma-constructible-abelian}.
Thus it suffices to show the statement on submodules.

\medskip\noindent
Let $\mathcal{G} \subset \mathcal{F}$ be a submodule of a
constructible sheaf of $\Lambda$-modules on $X_\etale$. Let $\eta \in X$
be a generic point of an irreducible component of $X$. By Noetherian induction
it suffices to find an open neighbourhood $U$ of $\eta$ such that
$\mathcal{G}|_U$ is locally constant. To do this we may replace $X$
by an \'etale neighbourhood of $\eta$. Hence we may assume $\mathcal{F}$
is constant and $X$ is irreducible.

\medskip\noindent
Say $\mathcal{F} = \underline{M}$ for some finite $\Lambda$-module $M$.
Then $M' = \mathcal{G}_{\overline{\eta}} \subset M$. Pick finitely
many elements $s_1, \ldots, s_t$ generating $M'$ as a $\Lambda$-module.
(This is possible as $\Lambda$ is Noetherian and $M$ is finite.)
Pick an \'etale neighbourhood $(U, \overline{u})$ of $\overline{\eta}$
and sections $\sigma_1, \ldots, \sigma_t \in \mathcal{G}(U)$ which map to
$s_i$ in $\mathcal{G}_{\overline{\eta}} \subset M$.
Since $\sigma_i$ maps to an element
$s_i \in M' \subset M = \Gamma(X, \mathcal{F})$
we see that the two pullbacks of $\sigma_i$ to $U \times_X U$
are the same as sections of $\mathcal{G}$. By the sheaf condition
for $\mathcal{G}$ we find that $\sigma_i$ comes from a section
of $\mathcal{G}$ over the open $\Im(U \to X)$ of $X$.
Shrinking $X$ we may assume
$\underline{M'} \subset \mathcal{G} \subset \underline{M}$.
Then we see that $\underline{M'} = \mathcal{G}$ by
Lemma \ref{lemma-irreducible-subsheaf-constant-zero}.

\medskip\noindent
Proof of (2). This follows in the usual manner from (3). Details
omitted.
\end{proof}

\noindent
The following lemma tells us that every object of the abelian category of
constructible sheaves on $X$ is ``Noetherian'', i.e., satisfies
a.c.c.\ for subobjects.

\begin{lemma}
\label{lemma-constructible-over-noetherian-noetherian}
Let $X$ be a Noetherian scheme. Let $\Lambda$ be a Noetherian ring.
Consider inclusions
$$
\mathcal{F}_1 \subset \mathcal{F}_2 \subset \mathcal{F}_3 \subset \ldots
\subset \mathcal{F}
$$
in the category of sheaves of sets, abelian groups, or $\Lambda$-modules.
If $\mathcal{F}$ is constructible, then for some $n$
we have $\mathcal{F}_n = \mathcal{F}_{n + 1} = \mathcal{F}_{n + 2} = \ldots$.
\end{lemma}

\begin{proof}
By Proposition \ref{proposition-constructible-over-noetherian}
we see that $\mathcal{F}_i$ and $\colim \mathcal{F}_i$ are constructible.
Then the lemma follows from
Lemma \ref{lemma-colimit-constructible}.
\end{proof}

\begin{lemma}
\label{lemma-constructible-maps-into-constant}
Let $X$ be a Noetherian scheme.
\begin{enumerate}
\item Let $\mathcal{F}$ be a constructible sheaf of sets on $X_\etale$.
There exist an injective map of sheaves
$$
\mathcal{F} \longrightarrow
\prod\nolimits_{i = 1, \ldots, n} f_{i, *}\underline{E_i}
$$
where $f_i : Y_i \to X$ is a finite morphism and $E_i$ is a finite set.
\item Let $\mathcal{F}$ be a constructible abelian sheaf on $X_\etale$.
There exist an injective map of abelian sheaves
$$
\mathcal{F} \longrightarrow
\bigoplus\nolimits_{i = 1, \ldots, n} f_{i, *}\underline{M_i}
$$
where $f_i : Y_i \to X$ is a finite morphism and
$M_i$ is a finite abelian group.
\item Let $\Lambda$ be a Noetherian ring.
Let $\mathcal{F}$ be a constructible sheaf of $\Lambda$-modules on $X_\etale$.
There exist an injective map of sheaves of modules
$$
\mathcal{F} \longrightarrow
\bigoplus\nolimits_{i = 1, \ldots, n} f_{i, *}\underline{M_i}
$$
where $f_i : Y_i \to X$ is a finite morphism and
$M_i$ is a finite $\Lambda$-module.
\end{enumerate}
Moreover, we may assume each $Y_i$ is irreducible, reduced, maps onto
an irreducible and reduced closed subscheme $Z_i \subset X$ such that
$Y_i \to Z_i$ is finite \'etale over a nonempty open of $Z_i$.
\end{lemma}

\begin{proof}
Proof of (1). Because we have the ascending chain condition for
subsheaves of $\mathcal{F}$
(Lemma \ref{lemma-constructible-over-noetherian-noetherian}), it
suffices to show that for every point $x \in X$ we
can find a map $\varphi : \mathcal{F} \to f_*\underline{E}$ where
$f : Y \to X$ is finite and $E$ is a finite set such that
$\varphi_{\overline{x}} : \mathcal{F}_{\overline{x}} \to
(f_*S)_{\overline{x}}$ is injective.
(This argument can be avoided by picking a partition of $X$ as in
Lemma \ref{lemma-constructible-quasi-compact-quasi-separated}
and constructing a $Y_i \to X$ for each irreducible component
of each part.)
Let $Z \subset X$ be the induced reduced scheme structure
(Schemes, Definition \ref{schemes-definition-reduced-induced-scheme})
on $\overline{\{x\}}$.
Since $\mathcal{F}$ is constructible, there is a finite separable
extension $K/\kappa(x)$ such that
$\mathcal{F}|_{\Spec(K)}$ is the constant sheaf with value $E$
for some finite set $E$. Let $Y \to Z$ be the normalization
of $Z$ in $\Spec(K)$.
By Morphisms, Lemma \ref{morphisms-lemma-normal-normalization}
we see that $Y$ is a normal integral scheme.
As $K/\kappa(x)$ is a finite extension, it is clear that $K$ is the function
field of $Y$. Denote $g : \Spec(K) \to Y$ the inclusion.
The map $\mathcal{F}|_{\Spec(K)} \to \underline{E}$ is adjoint
to a map $\mathcal{F}|_Y \to g_*\underline{E} = \underline{E}$
(Lemma \ref{lemma-push-constant-sheaf-from-open}).
This in turn is adjoint to a map
$\varphi : \mathcal{F} \to f_*\underline{E}$.
Observe that the stalk of $\varphi$ at a geometric point
$\overline{x}$ is injective: we may take a lift $\overline{y} \in Y$
of $\overline{x}$ and the commutative diagram
$$
\xymatrix{
\mathcal{F}_{\overline{x}} \ar@{=}[r] \ar[d] &
(\mathcal{F}|_Y)_{\overline{y}} \ar@{=}[d] \\
(f_*\underline{E})_{\overline{x}} \ar[r] &
\underline{E}_{\overline{y}}
}
$$
proves the injectivity. We are not yet done, however, as the
morphism $f : Y \to Z$ is integral but in general not
finite\footnote{If $X$ is a Nagata scheme, for example of finite
type over a field, then $Y \to Z$ is finite.}.

\medskip\noindent
To fix the problem stated in the last sentence of the previous paragraph,
we write $Y = \lim_{i \in I} Y_i$ with $Y_i$ irreducible, integral, and
finite over $Z$. Namely, apply Properties, Lemma
\ref{properties-lemma-integral-algebra-directed-colimit-finite}
to $f_*\mathcal{O}_Y$ viewed as a sheaf of $\mathcal{O}_Z$-algebras
and apply the functor $\underline{\Spec}_Z$.
Then $f_*\underline{E} = \colim f_{i, *}\underline{E}$
by Lemma \ref{lemma-relative-colimit}.
By Lemma \ref{lemma-constructible-is-compact} the map
$\mathcal{F} \to f_*\underline{E}$
factors through $f_{i, *}\underline{E}$ for some $i$.
Since $Y_i \to Z$ is a finite morphism of integral schemes
and since the function field extension
induced by this morphism is finite separable, we see that the
morphism is finite \'etale over a nonempty open of $Z$ (use
Algebra, Lemma \ref{algebra-lemma-smooth-at-generic-point}; details omitted).
This finishes the proof of (1).

\medskip\noindent
The proofs of (2) and (3) are identical to the proof of (1).
\end{proof}

\noindent
In the following lemma we use a standard trick to reduce a very general
statement to the Noetherian case.

\begin{lemma}
\label{lemma-constructible-maps-into-constant-general}
\begin{reference}
\cite[Exposee IX, Proposition 2.14]{SGA4}
\end{reference}
Let $X$ be a quasi-compact and quasi-separated scheme.
\begin{enumerate}
\item Let $\mathcal{F}$ be a constructible sheaf of sets on $X_\etale$.
There exist an injective map of sheaves
$$
\mathcal{F} \longrightarrow
\prod\nolimits_{i = 1, \ldots, n} f_{i, *}\underline{E_i}
$$
where $f_i : Y_i \to X$ is a finite and finitely presented morphism and
$E_i$ is a finite set.
\item Let $\mathcal{F}$ be a constructible abelian sheaf on $X_\etale$.
There exist an injective map of abelian sheaves
$$
\mathcal{F} \longrightarrow
\bigoplus\nolimits_{i = 1, \ldots, n} f_{i, *}\underline{M_i}
$$
where $f_i : Y_i \to X$ is a finite and finitely presented morphism and
$M_i$ is a finite abelian group.
\item Let $\Lambda$ be a Noetherian ring.
Let $\mathcal{F}$ be a constructible sheaf of $\Lambda$-modules on $X_\etale$.
There exist an injective map of sheaves of modules
$$
\mathcal{F} \longrightarrow
\bigoplus\nolimits_{i = 1, \ldots, n} f_{i, *}\underline{M_i}
$$
where $f_i : Y_i \to X$ is a finite and finitely presented morphism and
$M_i$ is a finite $\Lambda$-module.
\end{enumerate}
\end{lemma}

\begin{proof}
We will reduce this lemma to the Noetherian case by absolute Noetherian
approximation. Namely, by
Limits, Proposition \ref{limits-proposition-approximate}
we can write $X = \lim_{t \in T} X_t$ with each $X_t$ of finite type over
$\Spec(\mathbf{Z})$ and with affine transition morphisms. By
Lemma \ref{lemma-category-is-colimit}
the category of constructible sheaves (of sets, abelian groups, or
$\Lambda$-modules) on $X_\etale$ is the colimit of the corresponding
categories for $X_t$. Thus our constructible sheaf $\mathcal{F}$
is the pullback of a similar constructible sheaf $\mathcal{F}_t$
over $X_t$ for some $t$. Then we apply the Noetherian case
(Lemma \ref{lemma-constructible-maps-into-constant})
to find an injection
$$
\mathcal{F}_t \longrightarrow
\prod\nolimits_{i = 1, \ldots, n} f_{i, *}\underline{E_i}
\quad\text{or}\quad
\mathcal{F}_t \longrightarrow
\bigoplus\nolimits_{i = 1, \ldots, n} f_{i, *}\underline{M_i}
$$
over $X_t$ for some finite morphisms $f_i : Y_i \to X_t$.
Since $X_t$ is Noetherian the morphisms $f_i$ are of finite presentation.
Since pullback is exact and since formation of $f_{i, *}$ commutes
with base change
(Lemma \ref{lemma-finite-pushforward-commutes-with-base-change}), we conclude.
\end{proof}

\begin{lemma}
\label{lemma-support-in-subset}
Let $X$ be a Noetherian scheme. Let $E \subset X$ be a subset closed
under specialization.
\begin{enumerate}
\item Let $\mathcal{F}$ be a torsion abelian sheaf on $X_\etale$
whose support is contained in $E$. Then $\mathcal{F} = \colim \mathcal{F}_i$
is a filtered colimit of constructible abelian sheaves $\mathcal{F}_i$
such that for each $i$ the support of $\mathcal{F}_i$ is contained in
a closed subset contained in $E$.
\item Let $\Lambda$ be a Noetherian ring and $\mathcal{F}$ a sheaf
of $\Lambda$-modules on $X_\etale$ whose support is contained in $E$.
Then $\mathcal{F} = \colim \mathcal{F}_i$
is a filtered colimit of constructible sheaves of
$\Lambda$-modules $\mathcal{F}_i$ such that for each $i$
the support of $\mathcal{F}_i$ is contained in a closed subset
contained in $E$.
\end{enumerate}
\end{lemma}

\begin{proof}
Proof of (1). We can write $\mathcal{F} = \colim_{i \in I} \mathcal{F}_i$
with $\mathcal{F}_i$ constructible abelian by
Lemma \ref{lemma-torsion-colimit-constructible}.
By Proposition \ref{proposition-constructible-over-noetherian}
the image $\mathcal{F}'_i \subset \mathcal{F}$
of the map $\mathcal{F}_i \to \mathcal{F}$ is constructible.
Thus $\mathcal{F} = \colim \mathcal{F}'_i$ and the support
of $\mathcal{F}'_i$ is contained in $E$.
Since the support of $\mathcal{F}'_i$ is constructible
(by our definition of constructible sheaves), we see
that its closure is also contained in $E$, see for example
Topology, Lemma
\ref{topology-lemma-constructible-stable-specialization-closed}.

\medskip\noindent
The proof in case (2) is exactly the same and we omit it.
\end{proof}








\section{Specializations and \'etale sheaves}
\label{section-specialization}

\noindent
Topological picture: Let $X$ be a topological space and let $x' \leadsto x$
be a specialization of points in $X$. Then every open neighbourhood of $x$
contains $x'$. Hence for any sheaf $\mathcal{F}$ on $X$ there is a
{\it specialization map}
$$
sp : \mathcal{F}_x \longrightarrow \mathcal{F}_{x'}
$$
of stalks sending the equivalence class of the pair $(U, s)$ in
$\mathcal{F}_x$ to the equivalence class of the pair $(U, s)$ in
$\mathcal{F}_{x'}$; see Sheaves, Section \ref{sheaves-section-stalks}
for the description of stalks in terms of equivalence classes of pairs.
Of course this map is functorial in $\mathcal{F}$, i.e., $sp$
is a transformation of functors.

\medskip\noindent
For sheaves in the \'etale topology we can mimic this construction, see
\cite[Exposee VIII, 7.7, page 397]{SGA4}. To do this suppose we have
a scheme $S$, a geometric point $\overline{s}$ of $S$, and a geometric
point $\overline{t}$ of $\Spec(\mathcal{O}^{sh}_{S, \overline{s}})$.
For any sheaf $\mathcal{F}$ on $S_\etale$ we will construct the
{\it specialization map}
$$
sp : \mathcal{F}_{\overline{s}} \longrightarrow \mathcal{F}_{\overline{t}}
$$
Here we have abused language: instead of writing
$\mathcal{F}_{\overline{t}}$ we should write $\mathcal{F}_{p(\overline{t})}$
where $p : \Spec(\mathcal{O}^{sh}_{S, \overline{s}}) \to S$ is the canonical
morphism. Recall that
$$
\mathcal{F}_{\overline{s}} = \colim_{(U, \overline{u})} \mathcal{F}(U)
$$
where the colimit is over all \'etale neighbourhoods $(U, \overline{u})$
of $(S, \overline{s})$, see Section \ref{section-stalks}. Since
$\mathcal{O}^{sh}_{S, \overline{s}}$ is the stalk of the
structure sheaf, we find for every \'etale neighbourhood
$(U, \overline{u})$ of $(S, \overline{s})$ a canonical map
$\mathcal{O}_{U, u} \to \mathcal{O}^{sh}_{S, \overline{s}}$.
Hence we get a unique factorization
$$
\Spec(\mathcal{O}^{sh}_{S, \overline{s}}) \to U \to S
$$
If $\overline{v}$ denotes the image of $\overline{t}$ in $U$, then we see that
$(U, \overline{v})$ is an \'etale neighbourhood of $(S, \overline{t})$.
This construction defines a functor from the category of \'etale neighbourhoods
of $(S, \overline{s})$ to the category of \'etale neighbourhoods of
$(S, \overline{t})$. Thus we may define the map
$sp : \mathcal{F}_{\overline{s}} \to \mathcal{F}_{\overline{t}}$
by sending the equivalence class of
$(U, \overline{u}, \sigma)$ where $\sigma \in \mathcal{F}(U)$
to the equivalence class of $(U, \overline{v}, \sigma)$.

\medskip\noindent
Let $K \in D(S_\etale)$. With $\overline{s}$ and $\overline{t}$
as above we have the {\it specialization map}
$$
sp : K_{\overline{s}} \longrightarrow K_{\overline{t}}
\quad\text{in}\quad D(\textit{Ab})
$$
Namely, if $K$ is represented by the complex $\mathcal{F}^\bullet$
of abelian sheaves, then we simply that the map
$$
K_{\overline{s}} = \mathcal{F}^\bullet_{\overline{s}}
\longrightarrow
\mathcal{F}^\bullet_{\overline{t}} = K_{\overline{t}}
$$
which is termwise given by the specialization maps for sheaves
constructed above. This is independent of the choice of complex
representing $K$ by the exactness of the stalk functors (i.e.,
taking stalks of complexes is well defined on the derived category).

\medskip\noindent
Clearly the construction is functorial in the sheaf $\mathcal{F}$ on
$S_\etale$. If we think of the stalk functors as morphisms of topoi
$\overline{s}, \overline{t} : \textit{Sets} \to \Sh(S_\etale)$, then
we may think of $sp$ as a $2$-morphism
$$
\xymatrix{
\textit{Sets}
\rrtwocell^{\overline{t}}_{\overline{s}}{\ sp}
&
&
\Sh(S_\etale)
}
$$
of topoi.

\begin{remark}[Alternative description of sp]
\label{remark-alternative-sp}
Let $S$, $\overline{s}$, and $\overline{t}$ be as above.
Another way to describe the specialization map is to use that
$$
\mathcal{F}_{\overline{s}} =
\Gamma(\Spec(\mathcal{O}^{sh}_{S, \overline{s}}), p^{-1}\mathcal{F})
\quad\text{and}\quad
\mathcal{F}_{\overline{t}} = \Gamma(\overline{t},
\overline{t}^{-1}p^{-1}\mathcal{F})
$$
The first equality follows from Theorem \ref{theorem-higher-direct-images}
applied to $\text{id}_S : S \to S$ and the second equality follows
from Lemma \ref{lemma-stalk-pullback}. Then we can think of $sp$
as the map
$$
sp :
\mathcal{F}_{\overline{s}} =
\Gamma(\Spec(\mathcal{O}^{sh}_{S, \overline{s}}), p^{-1}\mathcal{F})
\xrightarrow{\text{pullback by }\overline{t}}
\Gamma(\overline{t}, \overline{t}^{-1}p^{-1}\mathcal{F}) =
\mathcal{F}_{\overline{t}}
$$
\end{remark}

\begin{remark}[Yet another description of sp]
\label{remark-another-sp}
Let $S$, $\overline{s}$, and $\overline{t}$ be as above.
Another alternative is to use the unique morphism
$$
c :
\Spec(\mathcal{O}^{sh}_{S, \overline{t}})
\longrightarrow
\Spec(\mathcal{O}^{sh}_{S, \overline{s}})
$$
over $S$ which is compatible with the given morphism
$\overline{t} \to \Spec(\mathcal{O}^{sh}_{S, \overline{s}})$
and the morphism
$\overline{t} \to \Spec(\mathcal{O}^{sh}_{t, \overline{t}})$.
The uniqueness and existence of the displayed arrow
follows from Algebra, Lemma \ref{algebra-lemma-map-into-henselian-colimit}
applied to $\mathcal{O}_{S, s}$, $\mathcal{O}^{sh}_{S, \overline{t}}$, and
$\mathcal{O}^{sh}_{S, \overline{s}} \to \kappa(\overline{t})$.
We obtain
$$
sp :
\mathcal{F}_{\overline{s}} =
\Gamma(\Spec(\mathcal{O}^{sh}_{S, \overline{s}}), \mathcal{F})
\xrightarrow{\text{pullback by }c}
\Gamma(\Spec(\mathcal{O}^{sh}_{S, \overline{t}}), \mathcal{F}) =
\mathcal{F}_{\overline{t}}
$$
(with obvious notational conventions).
In fact this procedure also works for objects $K$ in $D(S_\etale)$:
the specialization map for $K$ is equal to the map
$$
sp :
K_{\overline{s}} =
R\Gamma(\Spec(\mathcal{O}^{sh}_{S, \overline{s}}), K)
\xrightarrow{\text{pullback by }c}
R\Gamma(\Spec(\mathcal{O}^{sh}_{S, \overline{t}}), K) =
K_{\overline{t}}
$$
The equality signs are valid as taking global sections over
the strictly henselian schemes
$\Spec(\mathcal{O}^{sh}_{S, \overline{s}})$ and
$\Spec(\mathcal{O}^{sh}_{S, \overline{t}})$
is exact (and the same as taking stalks at $\overline{s}$ and $\overline{t}$)
and hence no subtleties related to the fact that $K$ may be unbounded
arise.
\end{remark}

\begin{remark}[Lifting specializations]
\label{remark-can-lift}
Let $S$ be a scheme and let $t \leadsto s$ be a specialization of
point on $S$. Choose geometric points $\overline{t}$ and $\overline{s}$
lying over $t$ and $s$. Since $t$ corresponds to a point of
$\Spec(\mathcal{O}_{S, s})$ by
Schemes, Lemma \ref{schemes-lemma-specialize-points}
and since $\mathcal{O}_{S, s} \to \mathcal{O}^{sh}_{S, \overline{s}}$
is faithfully flat, we can find a point
$t' \in \Spec(\mathcal{O}^{sh}_{S, \overline{s}})$ mapping to $t$.
As $\Spec(\mathcal{O}^{sh}_{S, \overline{s}})$ is a limit of
schemes \'etale over $S$ we see that $\kappa(t')/\kappa(t)$
is a separable algebraic extension (usually not finite of course).
Since $\kappa(\overline{t})$ is algebraically closed, we can
choose an embedding $\kappa(t') \to \kappa(\overline{t})$
as extensions of $\kappa(t)$. This choice gives us a
commutative diagram
$$
\xymatrix{
\overline{t} \ar[d] \ar[r] &
\Spec(\mathcal{O}^{sh}_{S, \overline{s}}) \ar[d] &
\overline{s} \ar[l] \ar[d] \\
t \ar[r] & S & s \ar[l]
}
$$
of points and geometric points. Thus if $t \leadsto s$
we can always ``lift'' $\overline{t}$ to a geometric point
of the strict henselization of $S$ at $\overline{s}$
and get specialization maps as above.
\end{remark}

\begin{lemma}
\label{lemma-specialization-map-pullback}
Let $g : S' \to S$ be a morphism of schemes. Let $\mathcal{F}$ be a sheaf
on $S_\etale$. Let $\overline{s}'$ be a geometric point of $S'$, and let
$\overline{t}'$ be a geometric point of
$\Spec(\mathcal{O}^{sh}_{S', \overline{s}'})$. Denote
$\overline{s} = g(\overline{s}')$ and $\overline{t} = h(\overline{t}')$
where $h : \Spec(\mathcal{O}^{sh}_{S', \overline{s}'}) \to
\Spec(\mathcal{O}^{sh}_{S, \overline{s}})$ is the canonical morphism.
For any sheaf $\mathcal{F}$ on $S_\etale$ the specialization map
$$
sp :
(g^{-1}\mathcal{F})_{\overline{s}'}
\longrightarrow
(g^{-1}\mathcal{F})_{\overline{t}'}
$$
is equal to the specialization map
$sp : \mathcal{F}_{\overline{s}} \to \mathcal{F}_{\overline{t}}$
via the identifications
$(g^{-1}\mathcal{F})_{\overline{s}'} = \mathcal{F}_{\overline{s}}$ and
$(g^{-1}\mathcal{F})_{\overline{t}'} = \mathcal{F}_{\overline{t}}$
of Lemma \ref{lemma-stalk-pullback}.
\end{lemma}

\begin{proof}
Omitted.
\end{proof}

\begin{lemma}
\label{lemma-characterize-locally-constant}
Let $S$ be a scheme such that every quasi-compact open of $S$ has
finite number of irreducible components (for example if $S$ has a
Noetherian underlying topological space, or if $S$ is locally Noetherian).
Let $\mathcal{F}$ be a sheaf of sets on $S_\etale$.
The following are equivalent
\begin{enumerate}
\item $\mathcal{F}$ is finite locally constant, and
\item all stalks of $\mathcal{F}$ are finite sets and all specialization maps
$sp : \mathcal{F}_{\overline{s}} \to \mathcal{F}_{\overline{t}}$
are bijective.
\end{enumerate}
\end{lemma}

\begin{proof}
Assume (2). Let $\overline{s}$ be a geometric point of $S$ lying over
$s \in S$. In order to prove (1) we have to find an \'etale neighbourhood
$(U, \overline{u})$ of $(S, \overline{s})$ such that $\mathcal{F}|_U$
is constant. We may and do assume $S$ is affine.

\medskip\noindent
Since $\mathcal{F}_{\overline{s}}$ is finite, we can choose
$(U, \overline{u})$, $n \geq 0$, and pairwise distinct elements
$\sigma_1, \ldots, \sigma_n \in \mathcal{F}(U)$
such that $\{\sigma_1, \ldots, \sigma_n\} \subset \mathcal{F}(U)$ maps
bijectively to $\mathcal{F}_{\overline{s}}$ via the map
$\mathcal{F}(U) \to \mathcal{F}_{\overline{s}}$.
Consider the map
$$
\varphi : \underline{\{1, \ldots, n\}} \longrightarrow \mathcal{F}|_U
$$
on $U_\etale$ defined by $\sigma_1, \ldots, \sigma_n$.
This map is a bijection on stalks at $\overline{u}$
by construction. Let us consider the subset
$$
E = \{u' \in U \mid \varphi_{\overline{u}'}\text{ is bijective}\} \subset U
$$
Here $\overline{u}'$ is any geometric point of $U$ lying over $u'$
(the condition is independent of the choice by Remark \ref{remark-map-stalks}).
The image $u \in U$ of $\overline{u}$ is in $E$.
By our assumption on the specialization maps for $\mathcal{F}$,
by Remark \ref{remark-can-lift}, and by
Lemma \ref{lemma-specialization-map-pullback}
we see that $E$ is closed under specializations
and generalizations in the topological space $U$.

\medskip\noindent
After shrinking $U$ we may assume $U$ is affine too. By
Descent, Lemma \ref{descent-lemma-locally-finite-nr-irred-local-fppf}
we see that $U$ has a finite number of irreducible components.
After removing the irreducible components which do not pass
through $u$, we may assume every irreducible component of $U$
passes through $u$. Since $U$ is a sober topological space
it follows that $E = U$ and we conclude that $\varphi$ is an isomorphism by
Theorem \ref{theorem-exactness-stalks}. Thus (1) follows.

\medskip\noindent
We omit the proof that (1) implies (2).
\end{proof}

\begin{lemma}
\label{lemma-characterize-locally-constant-module}
Let $S$ be a scheme such that every quasi-compact open of $S$ has
finite number of irreducible components (for example if $S$ has a
Noetherian underlying topological space, or if $S$ is locally Noetherian).
Let $\Lambda$ be a Noetherian ring.
Let $\mathcal{F}$ be a sheaf of $\Lambda$-modules on $S_\etale$.
The following are equivalent
\begin{enumerate}
\item $\mathcal{F}$ is a finite type, locally constant sheaf
of $\Lambda$-modules, and
\item all stalks of $\mathcal{F}$ are finite $\Lambda$-modules and
all specialization maps
$sp : \mathcal{F}_{\overline{s}} \to \mathcal{F}_{\overline{t}}$
are bijective.
\end{enumerate}
\end{lemma}

\begin{proof}
The proof of this lemma is the same as the proof of
Lemma \ref{lemma-characterize-locally-constant}.
Assume (2). Let $\overline{s}$ be a geometric point of $S$ lying over
$s \in S$. In order to prove (1) we have to find an \'etale neighbourhood
$(U, \overline{u})$ of $(S, \overline{s})$ such that $\mathcal{F}|_U$
is constant. We may and do assume $S$ is affine.

\medskip\noindent
Since $M = \mathcal{F}_{\overline{s}}$ is a finite $\Lambda$-module
and $\Lambda$ is Noetherian, we can choose a presentation
$$
\Lambda^{\oplus m} \xrightarrow{A} \Lambda^{\oplus n} \to M \to 0
$$
for some matrix $A = (a_{ji})$ with coefficients in $\Lambda$.
We can choose $(U, \overline{u})$ and elements
$\sigma_1, \ldots, \sigma_n \in \mathcal{F}(U)$
such that $\sum a_{ji}\sigma_i = 0$ in $\mathcal{F}(U)$ and
such that the images of $\sigma_i$ in $\mathcal{F}_{\overline{s}} = M$
are the images of the standard basis element of
$\Lambda^n$ in the presentation of $M$ given above.
Consider the map
$$
\varphi : \underline{M} \longrightarrow \mathcal{F}|_U
$$
on $U_\etale$ defined by $\sigma_1, \ldots, \sigma_n$.
This map is a bijection on stalks at $\overline{u}$
by construction. Let us consider the subset
$$
E = \{u' \in U \mid \varphi_{\overline{u}'}\text{ is bijective}\} \subset U
$$
Here $\overline{u}'$ is any geometric point of $U$ lying over $u'$
(the condition is independent of the choice by Remark \ref{remark-map-stalks}).
The image $u \in U$ of $\overline{u}$ is in $E$.
By our assumption on the specialization maps for $\mathcal{F}$,
by Remark \ref{remark-can-lift}, and by
Lemma \ref{lemma-specialization-map-pullback}
we see that $E$ is closed under specializations
and generalizations in the topological space $U$.

\medskip\noindent
After shrinking $U$ we may assume $U$ is affine too. By
Descent, Lemma \ref{descent-lemma-locally-finite-nr-irred-local-fppf}
we see that $U$ has a finite number of irreducible components.
After removing the irreducible components which do not pass
through $u$, we may assume every irreducible component of $U$
passes through $u$. Since $U$ is a sober topological space
it follows that $E = U$ and we conclude that $\varphi$ is an isomorphism by
Theorem \ref{theorem-exactness-stalks}. Thus (1) follows.

\medskip\noindent
We omit the proof that (1) implies (2).
\end{proof}

\begin{lemma}
\label{lemma-specialization-map-pushforward}
Let $f : X \to S$ be a quasi-compact and quasi-separated morphism of schemes.
Let $K \in D^+(X_\etale)$. Let $\overline{s}$ be a geometric point of $S$
and let $\overline{t}$ be a geometric point of
$\Spec(\mathcal{O}^{sh}_{S, \overline{s}})$. We have a commutative
diagram
$$
\xymatrix{
(Rf_*K)_{\overline{s}} \ar[r]_{sp} \ar@{=}[d] &
(Rf_*K)_{\overline{t}} \ar@{=}[d] \\
R\Gamma(X \times_S \Spec(\mathcal{O}^{sh}_{S, \overline{s}}), K)
\ar[r] &
R\Gamma(X \times_S \Spec(\mathcal{O}^{sh}_{S, \overline{t}}), K)
}
$$
where the bottom horizontal arrow arises as pullback by the morphism
$\text{id}_X \times c$ where
$c : \Spec(\mathcal{O}^{sh}_{S, \overline{t}}) \to
\Spec(\mathcal{O}^{sh}_{S, \overline{s}})$
is the morphism introduced in Remark \ref{remark-another-sp}.
The vertical arrows are given by
Theorem \ref{theorem-higher-direct-images}.
\end{lemma}

\begin{proof}
This follows immediately from the description
of $sp$ in Remark \ref{remark-another-sp}.
\end{proof}

\begin{remark}
\label{remark-specialization-map-and-fibres}
Let $f : X \to S$ be a morphism of schemes. Let $K \in D(X_\etale)$.
Let $\overline{s}$ be a geometric point of $S$ and let $\overline{t}$
be a geometric point of $\Spec(\mathcal{O}^{sh}_{S, \overline{s}})$.
Let $c$ be as in Remark \ref{remark-another-sp}.
We can always make a commutative diagram
$$
\xymatrix{
(Rf_*K)_{\overline{s}} \ar[r] \ar[d]_{sp} &
R\Gamma(X \times_S \Spec(\mathcal{O}_{S, \overline{s}}^{sh}), K)
\ar[r] \ar[d]_{(\text{id}_X \times c)^{-1}} &
R\Gamma(X_{\overline{s}}, K) \\
(Rf_*K)_{\overline{t}} \ar[r] &
R\Gamma(X \times_S \Spec(\mathcal{O}_{S, \overline{t}}^{sh}), K) \ar[r] &
R\Gamma(X_{\overline{t}}, K)
}
$$
where the horizontal arrows are those of Remark \ref{remark-stalk-fibre}.
In general there won't be a vertical map on the right
between the cohomologies of $K$ on the fibres fitting into this
diagram, even in the case of Lemma \ref{lemma-specialization-map-pushforward}.
\end{remark}















\section{Complexes with constructible cohomology}
\label{section-Dc}

\noindent
Let $\Lambda$ be a ring. Denote $D(X_\etale, \Lambda)$ the derived category
of sheaves of $\Lambda$-modules on $X_\etale$.
We denote by $D^b(X_\etale, \Lambda)$ (respectively $D^+$, $D^-$)
the full subcategory of bounded (resp. above, below) complexes in
$D(X_\etale, \Lambda)$.

\begin{definition}
\label{definition-c}
Let $X$ be a scheme. Let $\Lambda$ be a Noetherian ring.
We denote {\it $D_c(X_\etale, \Lambda)$} the full subcategory
of $D(X_\etale, \Lambda)$ of complexes whose cohomology sheaves
are constructible sheaves of $\Lambda$-modules.
\end{definition}

\noindent
This definition makes sense by Lemma \ref{lemma-constructible-abelian} and
Derived Categories, Section \ref{derived-section-triangulated-sub}.
Thus we see that $D_c(X_\etale, \Lambda)$ is a strictly full, saturated
triangulated subcategory of $D(X_\etale, \Lambda)$.

\begin{lemma}
\label{lemma-restrict-and-shriek-from-etale-c}
Let $\Lambda$ be a Noetherian ring.
If $j : U \to X$ is an \'etale morphism of schemes, then
\begin{enumerate}
\item $K|_U \in D_c(U_\etale, \Lambda)$ if $K \in D_c(X_\etale, \Lambda)$, and
\item $j_!M \in D_c(X_\etale, \Lambda)$ if $M \in D_c(U_\etale, \Lambda)$ and
the morphism $j$ is quasi-compact and quasi-separated.
\end{enumerate}
\end{lemma}

\begin{proof}
The first assertion is clear. The second follows from the fact
that $j_!$ is exact and Lemma \ref{lemma-jshriek-constructible}.
\end{proof}

\begin{lemma}
\label{lemma-pullback-c}
Let $\Lambda$ be a Noetherian ring.
Let $f : X \to Y$ be a morphism of schemes. If $K \in D_c(Y_\etale, \Lambda)$
then $Lf^*K \in D_c(X_\etale, \Lambda)$.
\end{lemma}

\begin{proof}
This follows as $f^{-1} = f^*$ is exact and
Lemma \ref{lemma-pullback-constructible}.
\end{proof}

\begin{lemma}
\label{lemma-one-constructible}
Let $X$ be a quasi-compact and quasi-separated scheme.
Let $\Lambda$ be a Noetherian ring. Let $K \in D(X_\etale, \Lambda)$
and $b \in \mathbf{Z}$ such that $H^b(K)$ is constructible.
Then there exist a sheaf $\mathcal{F}$ which is a finite direct sum
of $j_{U!}\underline{\Lambda}$ with $U \in \Ob(X_\etale)$ affine and
a map $\mathcal{F}[-b] \to K$ in $D(X_\etale, \Lambda)$
inducing a surjection $\mathcal{F} \to H^b(K)$.
\end{lemma}

\begin{proof}
Represent $K$ by a complex $\mathcal{K}^\bullet$ of sheaves of
$\Lambda$-modules. Consider the surjection
$$
\Ker(\mathcal{K}^b \to \mathcal{K}^{b + 1})
\longrightarrow
H^b(K)
$$
By Modules on Sites, Lemma
\ref{sites-modules-lemma-module-quotient-direct-sum}
we may choose a surjection
$\bigoplus_{i \in I} j_{U_i!} \underline{\Lambda} \to
\Ker(\mathcal{K}^b \to \mathcal{K}^{b + 1})$
with $U_i$ affine. For $I' \subset I$ finite, denote
$\mathcal{H}_{I'} \subset H^b(K)$ the image of
$\bigoplus_{i \in I'} j_{U_i!} \underline{\Lambda}$. By
Lemma \ref{lemma-colimit-constructible}
we see that $\mathcal{H}_{I'} = H^b(K)$ for some $I' \subset I$ finite.
The lemma follows taking
$\mathcal{F} = \bigoplus_{i \in I'} j_{U_i!} \underline{\Lambda}$.
\end{proof}

\begin{lemma}
\label{lemma-bounded-above-c}
Let $X$ be a quasi-compact and quasi-separated scheme.
Let $\Lambda$ be a Noetherian ring. Let $K \in D^-(X_\etale, \Lambda)$. Then
the following are equivalent
\begin{enumerate}
\item $K$ is in $D_c(X_\etale, \Lambda)$,
\item $K$ can be represented by a bounded above complex
whose terms are finite direct sums of $j_{U!}\underline{\Lambda}$
with $U \in \Ob(X_\etale)$ affine,
\item $K$ can be represented by a bounded above complex
of flat constructible sheaves of $\Lambda$-modules.
\end{enumerate}
\end{lemma}

\begin{proof}
It is clear that (2) implies (3) and that (3) implies (1).
Assume $K$ is in $D_c^-(X_\etale, \Lambda)$.
Say $H^i(K) = 0$ for $i > b$. By induction on $a$
we will construct a complex $\mathcal{F}^a \to \ldots \to \mathcal{F}^b$
such that each $\mathcal{F}^i$ is a finite direct sum
of $j_{U!}\underline{\Lambda}$ with $U \in \Ob(X_\etale)$ affine
and a map $\mathcal{F}^\bullet \to K$ which induces an isomorphism
$H^i(\mathcal{F}^\bullet) \to H^i(K)$ for $i > a$ and a surjection
$H^a(\mathcal{F}^\bullet) \to H^a(K)$.
For $a = b$ this can be done by Lemma \ref{lemma-one-constructible}.
Given such a datum choose a distinguished triangle
$$
\mathcal{F}^\bullet \to K \to L \to \mathcal{F}^\bullet[1]
$$
Then we see that $H^i(L) = 0$ for $i \geq a$. Choose
$\mathcal{F}^{a - 1}[-a +1] \to L$ as in
Lemma \ref{lemma-one-constructible}. The composition
$\mathcal{F}^{a - 1}[-a +1] \to L \to \mathcal{F}^\bullet$
corresponds to a map $\mathcal{F}^{a - 1} \to \mathcal{F}^a$
such that the composition with $\mathcal{F}^a \to \mathcal{F}^{a + 1}$
is zero. By TR4 we obtain a map
$$
(\mathcal{F}^{a - 1} \to \ldots \to \mathcal{F}^b) \to K
$$
in $D(X_\etale, \Lambda)$. This finishes the induction step and the
proof of the lemma.
\end{proof}

\begin{lemma}
\label{lemma-tensor-c}
Let $X$ be a scheme. Let $\Lambda$ be a Noetherian ring.
Let $K, L \in D_c^-(X_\etale, \Lambda)$. Then
$K \otimes_\Lambda^\mathbf{L} L$ is in $D_c^-(X_\etale, \Lambda)$.
\end{lemma}

\begin{proof}
This follows from Lemmas \ref{lemma-bounded-above-c} and
\ref{lemma-tensor-product-constructible}.
\end{proof}




\section{Tor finite with constructible cohomology}
\label{section-ctf}

\noindent
Let $X$ be a scheme and let $\Lambda$ be a Noetherian ring. An often used
subcategory of the derived category $D_c(X_\etale, \Lambda)$ defined
in Section \ref{section-Dc} is the full subcategory consisting of objects
having (locally) finite tor dimension. Here is the formal definition.

\begin{definition}
\label{definition-ctf}
Let $X$ be a scheme. Let $\Lambda$ be a Noetherian ring. We denote
{\it $D_{ctf}(X_\etale, \Lambda)$} the full subcategory of
$D_c(X_\etale, \Lambda)$
consisting of objects having locally finite tor dimension.
\end{definition}

\noindent
This is a strictly full, saturated triangulated subcategory of
$D_c(X_\etale, \Lambda)$ and $D(X_\etale, \Lambda)$. By our conventions, see
Cohomology on Sites, Definition \ref{sites-cohomology-definition-tor-amplitude},
we see that
$$
D_{ctf}(X_\etale, \Lambda) \subset D^b_c(X_\etale, \Lambda)
\subset D^b(X_\etale, \Lambda)
$$
if $X$ is quasi-compact. A good way to think about objects of
$D_{ctf}(X_\etale, \Lambda)$ is given in Lemma \ref{lemma-when-ctf}.

\begin{remark}
\label{remark-different}
Objects in the derived category $D_{ctf}(X_\etale, \Lambda)$ in some sense have
better global properties than the perfect objects in  $D(\mathcal{O}_X)$.
Namely, it can happen that a complex of $\mathcal{O}_X$-modules
is locally quasi-isomorphic to a finite complex of finite locally free
$\mathcal{O}_X$-modules, without being
globally quasi-isomorphic to a bounded complex of locally free
$\mathcal{O}_X$-modules. The following lemma shows this does not
happen for $D_{ctf}$ on a Noetherian scheme.
\end{remark}

\begin{lemma}
\label{lemma-when-ctf}
Let $\Lambda$ be a Noetherian ring. Let $X$ be a quasi-compact
and quasi-separated scheme. Let $K \in D(X_\etale, \Lambda)$. The following
are equivalent
\begin{enumerate}
\item $K \in D_{ctf}(X_\etale, \Lambda)$, and
\item $K$ can be represented by a finite complex of constructible
flat sheaves of $\Lambda$-modules.
\end{enumerate}
In fact, if $K$ has tor amplitude in $[a, b]$ then we can represent
$K$ by a complex $\mathcal{F}^a \to \ldots \to \mathcal{F}^b$ with
$\mathcal{F}^p$ a constructible flat sheaf of $\Lambda$-modules.
\end{lemma}

\begin{proof}
It is clear that a finite complex of constructible
flat sheaves of $\Lambda$-modules has finite tor dimension.
It is also clear that it is an object of $D_c(X_\etale, \Lambda)$.
Thus we see that (2) implies (1).

\medskip\noindent
Assume (1). Choose $a, b \in \mathbf{Z}$ such that
$H^i(K \otimes_\Lambda^\mathbf{L} \mathcal{G}) = 0$ if
$i \not \in [a, b]$ for all sheaves of $\Lambda$-modules $\mathcal{G}$.
We will prove the final assertion holds by induction on $b - a$. If
$a = b$, then $K = H^a(K)[-a]$ is a flat constructible sheaf
and the result holds. Next, assume $b > a$. Represent $K$
by a complex $\mathcal{K}^\bullet$ of sheaves of $\Lambda$-modules.
Consider the surjection
$$
\Ker(\mathcal{K}^b \to \mathcal{K}^{b + 1})
\longrightarrow
H^b(K)
$$
By Lemma \ref{lemma-category-constructible-modules}
we can find finitely many affine schemes $U_i$ \'etale over $X$ and a
surjection $\bigoplus j_{U_i!}\underline{\Lambda}_{U_i} \to H^b(K)$.
After replacing $U_i$ by standard \'etale coverings $\{U_{ij} \to U_i\}$
we may assume this surjection lifts to a map
$\mathcal{F} = \bigoplus j_{U_i!}\underline{\Lambda}_{U_i} \to
\Ker(\mathcal{K}^b \to \mathcal{K}^{b + 1})$.
This map determines a distinguished triangle
$$
\mathcal{F}[-b] \to K \to L \to \mathcal{F}[-b + 1]
$$
in $D(X_\etale, \Lambda)$. Since $D_{ctf}(X_\etale, \Lambda)$ is a triangulated
subcategory we see that $L$ is in it too. In fact $L$ has
tor amplitude in $[a, b - 1]$ as $\mathcal{F}$ surjects onto
$H^b(K)$ (details omitted). By induction hypothesis we can find
a finite complex $\mathcal{F}^a \to \ldots \to \mathcal{F}^{b - 1}$
of flat constructible sheaves of $\Lambda$-modules representing $L$.
The map $L \to \mathcal{F}[-b + 1]$ corresponds to a map
$\mathcal{F}^b \to \mathcal{F}$ annihilating the image
of $\mathcal{F}^{b - 1} \to \mathcal{F}^b$. Then it follows
from axiom TR3 that $K$ is represented by the complex
$$
\mathcal{F}^a \to \ldots \to \mathcal{F}^{b - 1} \to \mathcal{F}^b
$$
which finishes the proof.
\end{proof}

\begin{remark}
\label{remark-projective-each-degree}
Let $\Lambda$ be a Noetherian ring. Let $X$ be a scheme.
For a bounded complex $K^\bullet$ of constructible flat $\Lambda$-modules
on $X_\etale$
each stalk $K^p_{\overline{x}}$ is a finite projective $\Lambda$-module.
Hence the stalks of the complex are perfect complexes of $\Lambda$-modules.
\end{remark}

\begin{lemma}
\label{lemma-restrict-and-shriek-from-etale-ctf}
Let $\Lambda$ be a Noetherian ring.
If $j : U \to X$ is an \'etale morphism of schemes, then
\begin{enumerate}
\item $K|_U \in D_{ctf}(U_\etale, \Lambda)$ if
$K \in D_{ctf}(X_\etale, \Lambda)$, and
\item $j_!M \in D_{ctf}(X_\etale, \Lambda)$ if
$M \in D_{ctf}(U_\etale, \Lambda)$ and
the morphism $j$ is quasi-compact and quasi-separated.
\end{enumerate}
\end{lemma}

\begin{proof}
Perhaps the easiest way to prove this lemma is to reduce to the
case where $X$ is affine and then apply Lemma \ref{lemma-when-ctf}
to translate it into a statement about finite complexes
of flat constructible sheaves of $\Lambda$-modules
where the result follows from Lemma \ref{lemma-jshriek-constructible}.
\end{proof}

\begin{lemma}
\label{lemma-pullback-ctf}
Let $\Lambda$ be a Noetherian ring. Let $f : X \to Y$ be a morphism of schemes.
If $K \in D_{ctf}(Y_\etale, \Lambda)$ then
$Lf^*K \in D_{ctf}(X_\etale, \Lambda)$.
\end{lemma}

\begin{proof}
Apply Lemma \ref{lemma-when-ctf} to reduce this to a question
about finite complexes of flat constructible sheaves of $\Lambda$-modules.
Then the statement follows as $f^{-1} = f^*$ is exact and
Lemma \ref{lemma-pullback-constructible}.
\end{proof}

\begin{lemma}
\label{lemma-connected-ctf-locally-constant}
Let $X$ be a connected scheme. Let $\Lambda$ be a Noetherian ring.
Let $K \in D_{ctf}(X_\etale, \Lambda)$ have locally constant cohomology sheaves.
Then there exists a finite complex of finite projective $\Lambda$-modules
$M^\bullet$ and an \'etale covering $\{U_i \to X\}$ such that
$K|_{U_i} \cong \underline{M^\bullet}|_{U_i}$ in $D(U_{i, \etale}, \Lambda)$.
\end{lemma}

\begin{proof}
Choose an \'etale covering $\{U_i \to X\}$ such that $K|_{U_i}$
is constant, say $K|_{U_i} \cong \underline{M_i^\bullet}_{U_i}$
for some finite complex of finite $\Lambda$-modules $M_i^\bullet$.
See Cohomology on Sites, Lemma
\ref{sites-cohomology-lemma-locally-constant}.
Observe that $U_i \times_X U_j$ is empty if $M_i^\bullet$
is not isomorphic to $M_j^\bullet$ in $D(\Lambda)$.
For each complex of $\Lambda$-modules $M^\bullet$ let
$I_{M^\bullet} =
\{i \in I \mid M_i^\bullet \cong M^\bullet\text{ in }D(\Lambda)\}$.
As \'etale morphisms are open we see that
$U_{M^\bullet} = \bigcup_{i \in I_{M^\bullet}} \Im(U_i \to X)$
is an open subset of $X$. Then $X = \coprod U_{M^\bullet}$ is a disjoint
open covering of $X$. As $X$ is connected only one $U_{M^\bullet}$
is nonempty. As $K$ is in $D_{ctf}(X_\etale, \Lambda)$ we see that $M^\bullet$
is a perfect complex of $\Lambda$-modules, see
More on Algebra, Lemma \ref{more-algebra-lemma-perfect}.
Hence we may assume $M^\bullet$ is a finite complex of finite projective
$\Lambda$-modules.
\end{proof}








\section{Torsion sheaves}
\label{section-torsion}

\noindent
A brief section on torsion abelian sheaves and their \'etale cohomology.
Let $\mathcal{C}$ be a site. We have shown in
Cohomology on Sites, Lemma \ref{sites-cohomology-lemma-torsion}
that any object in $D(\mathcal{C})$ whose cohomology sheaves are
torsion sheaves, can be represented by a complex all of whose terms
are torsion.

\begin{lemma}
\label{lemma-torsion-cohomology}
Let $X$ be a quasi-compact and quasi-separated scheme.
\begin{enumerate}
\item If $\mathcal{F}$ is a torsion abelian sheaf on $X_\etale$, then
$H^n_\etale(X, \mathcal{F})$ is a torsion abelian group for all $n$.
\item If $K$ in $D^+(X_\etale)$ has torsion cohomology sheaves, then
$H^n_\etale(X, K)$ is a torsion abelian group for all $n$.
\end{enumerate}
\end{lemma}

\begin{proof}
To prove (1) we write $\mathcal{F} = \bigcup \mathcal{F}[n]$ where
$\mathcal{F}[d]$ is the $d$-torsion subsheaf. By
Lemma \ref{lemma-colimit} we have
$H^n_\etale(X, \mathcal{F}) = \colim H^n_\etale(X, \mathcal{F}[d])$.
This proves (1) as $H^n_\etale(X, \mathcal{F}[d])$ is annihilated by $d$.

\medskip\noindent
To prove (2) we can use the spectral sequence
$E_2^{p, q} = H^p_\etale(X, H^q(K))$ converging to $H^n_\etale(X, K)$
(Derived Categories, Lemma \ref{derived-lemma-two-ss-complex-functor})
and the result for sheaves.
\end{proof}

\begin{lemma}
\label{lemma-torsion-direct-image}
Let $f : X \to Y$ be a quasi-compact and quasi-separated
morphism of schemes.
\begin{enumerate}
\item If $\mathcal{F}$ is a torsion abelian sheaf on $X_\etale$, then
$R^nf_*\mathcal{F}$ is a torsion abelian sheaf on $Y_\etale$ for all $n$.
\item If $K$ in $D^+(X_\etale)$ has torsion cohomology sheaves, then
$Rf_*K$ is an object of $D^+(Y_\etale)$ whose cohomology sheaves are
torsion abelian sheaves.
\end{enumerate}
\end{lemma}

\begin{proof}
Proof of (1). Recall that $R^nf_*\mathcal{F}$ is the sheaf associated
to the presheaf $V \mapsto H^n_\etale(X \times_Y V, \mathcal{F})$
on $Y_\etale$. See Cohomology on Sites,
Lemma \ref{sites-cohomology-lemma-higher-direct-images}.
If we choose $V$ affine, then $X \times_Y V$ is
quasi-compact and quasi-separated because $f$ is, hence
we can apply Lemma \ref{lemma-torsion-cohomology} to see that
$H^n_\etale(X \times_Y V, \mathcal{F})$ is torsion.

\medskip\noindent
Proof of (2). Recall that $R^nf_*K$ is the sheaf associated
to the presheaf $V \mapsto H^n_\etale(X \times_Y V, K)$
on $Y_\etale$. See Cohomology on Sites,
Lemma \ref{sites-cohomology-lemma-unbounded-describe-higher-direct-images}.
If we choose $V$ affine, then $X \times_Y V$ is
quasi-compact and quasi-separated because $f$ is, hence
we can apply Lemma \ref{lemma-torsion-cohomology} to see that
$H^n_\etale(X \times_Y V, K)$ is torsion.
\end{proof}









\section{Cohomology with support in a closed subscheme}
\label{section-cohomology-support}

\noindent
Let $X$ be a scheme and let $Z \subset X$ be a closed subscheme.
Let $\mathcal{F}$ be an abelian sheaf on $X_\etale$. We let
$$
\Gamma_Z(X, \mathcal{F}) =
\{s \in \mathcal{F}(X) \mid \text{Supp}(s) \subset Z\}
$$
be the sections with support in $Z$ (Definition \ref{definition-support}).
This is a left exact functor which is not exact in general.
Hence we obtain a derived functor
$$
R\Gamma_Z(X, -) : D(X_\etale) \longrightarrow D(\textit{Ab})
$$
and cohomology groups with support in $Z$ defined by
$H^q_Z(X, \mathcal{F}) = R^q\Gamma_Z(X, \mathcal{F})$.

\medskip\noindent
Let $\mathcal{I}$ be an injective abelian sheaf on $X_\etale$.
Let $U = X \setminus Z$.
Then the restriction map $\mathcal{I}(X) \to \mathcal{I}(U)$ is surjective
(Cohomology on Sites, Lemma
\ref{sites-cohomology-lemma-restriction-along-monomorphism-surjective})
with kernel $\Gamma_Z(X, \mathcal{I})$. It immediately follows that
for $K \in D(X_\etale)$ there is a distinguished triangle
$$
R\Gamma_Z(X, K) \to R\Gamma(X, K) \to R\Gamma(U, K) \to R\Gamma_Z(X, K)[1]
$$
in $D(\textit{Ab})$. As a consequence we obtain a long exact cohomology
sequence
$$
\ldots \to H^i_Z(X, K) \to H^i(X, K) \to H^i(U, K) \to
H^{i + 1}_Z(X, K) \to \ldots
$$
for any $K$ in $D(X_\etale)$.

\medskip\noindent
For an abelian sheaf $\mathcal{F}$ on $X_\etale$ we can consider
the {\it subsheaf of sections with support in $Z$}, denoted
$\mathcal{H}_Z(\mathcal{F})$, defined by the rule
$$
\mathcal{H}_Z(\mathcal{F})(U) =
\{s \in \mathcal{F}(U) \mid \text{Supp}(s) \subset U \times_X Z\}
$$
Here we use the support of a section from Definition \ref{definition-support}.
Using the equivalence of
Proposition \ref{proposition-closed-immersion-pushforward}
we may view $\mathcal{H}_Z(\mathcal{F})$ as an abelian sheaf on
$Z_\etale$. Thus we obtain a functor
$$
\textit{Ab}(X_\etale) \longrightarrow \textit{Ab}(Z_\etale),\quad
\mathcal{F} \longmapsto \mathcal{H}_Z(\mathcal{F})
$$
which is left exact, but in general not exact.

\begin{lemma}
\label{lemma-sections-with-support-acyclic}
Let $i : Z \to X$ be a closed immersion of schemes.
Let $\mathcal{I}$ be an injective abelian sheaf on $X_\etale$.
Then $\mathcal{H}_Z(\mathcal{I})$ is an injective abelian sheaf
on $Z_\etale$.
\end{lemma}

\begin{proof}
Observe that for any abelian sheaf $\mathcal{G}$ on $Z_\etale$
we have
$$
\Hom_Z(\mathcal{G}, \mathcal{H}_Z(\mathcal{F})) =
\Hom_X(i_*\mathcal{G}, \mathcal{F})
$$
because after all any section of $i_*\mathcal{G}$ has support in $Z$.
Since $i_*$ is exact (Section \ref{section-closed-immersions}) and as
$\mathcal{I}$ is injective on $X_\etale$ we conclude that
$\mathcal{H}_Z(\mathcal{I})$ is injective on $Z_\etale$.
\end{proof}

\noindent
Denote
$$
R\mathcal{H}_Z : D(X_\etale) \longrightarrow D(Z_\etale)
$$
the derived functor. We set
$\mathcal{H}^q_Z(\mathcal{F}) = R^q\mathcal{H}_Z(\mathcal{F})$ so that
$\mathcal{H}^0_Z(\mathcal{F}) = \mathcal{H}_Z(\mathcal{F})$.
By the lemma above we have a Grothendieck spectral sequence
$$
E_2^{p, q} = H^p(Z, \mathcal{H}^q_Z(\mathcal{F}))
\Rightarrow H^{p + q}_Z(X, \mathcal{F})
$$

\begin{lemma}
\label{lemma-cohomology-with-support-sheaf-on-support}
Let $i : Z \to X$ be a closed immersion of schemes.
Let $\mathcal{G}$ be an injective abelian sheaf on $Z_\etale$.
Then $\mathcal{H}^p_Z(i_*\mathcal{G}) = 0$ for $p > 0$.
\end{lemma}

\begin{proof}
This is true because the functor $i_*$ is exact and transforms
injective abelian sheaves into injective abelian sheaves
(Cohomology on Sites, Lemma
\ref{sites-cohomology-lemma-pushforward-injective-flat}).
\end{proof}

\begin{lemma}
\label{lemma-cohomology-with-support-triangle}
Let $i : Z \to X$ be a closed immersion of schemes.
Let $j : U \to X$ be the inclusion of the complement of $Z$.
Let $\mathcal{F}$ be an abelian sheaf on $X_\etale$.
There is a distinguished triangle
$$
i_*R\mathcal{H}_Z(\mathcal{F}) \to \mathcal{F} \to Rj_*(\mathcal{F}|_U) \to
i_*R\mathcal{H}_Z(\mathcal{F})[1]
$$
in $D(X_\etale)$. This produces an exact sequence
$$
0 \to i_*\mathcal{H}_Z(\mathcal{F}) \to \mathcal{F} \to
j_*(\mathcal{F}|_U) \to i_*\mathcal{H}^1_Z(\mathcal{F}) \to 0
$$
and isomorphisms
$R^pj_*(\mathcal{F}|_U) \cong i_*\mathcal{H}^{p + 1}_Z(\mathcal{F})$
for $p \geq 1$.
\end{lemma}

\begin{proof}
To get the distinguished triangle, choose an injective resolution
$\mathcal{F} \to \mathcal{I}^\bullet$. Then we obtain a short exact
sequence of complexes
$$
0 \to
i_*\mathcal{H}_Z(\mathcal{I}^\bullet) \to \mathcal{I}^\bullet
\to j_*(\mathcal{I}^\bullet|_U) \to 0
$$
by the discussion above. Thus the distinguished triangle by
Derived Categories, Section \ref{derived-section-canonical-delta-functor}.
\end{proof}

\noindent
Let $X$ be a scheme and let $Z \subset X$ be a closed subscheme.
We denote $D_Z(X_\etale)$ the strictly full saturated triangulated
subcategory of $D(X_\etale)$ consisting of complexes whose cohomology
sheaves are supported on $Z$. Note that $D_Z(X_\etale)$ only
depends on the underlying closed subset of $X$.

\begin{lemma}
\label{lemma-complexes-with-support-on-closed}
Let $i : Z \to X$ be a closed immersion of schemes.
The map $Ri_{small, *} = i_{small, *} : D(Z_\etale) \to D(X_\etale)$
induces an equivalence $D(Z_\etale) \to D_Z(X_\etale)$ with quasi-inverse
$$
i_{small}^{-1}|_{D_Z(X_\etale)} = R\mathcal{H}_Z|_{D_Z(X_\etale)}
$$
\end{lemma}

\begin{proof}
Recall that $i_{small}^{-1}$ and $i_{small, *}$ is an adjoint pair of
exact functors such that $i_{small}^{-1}i_{small, *}$ is isomorphic to
the identify functor on abelian sheaves. See
Proposition \ref{proposition-closed-immersion-pushforward} and
Lemma \ref{lemma-stalk-pullback}. Thus
$i_{small, *} : D(Z_\etale) \to D_Z(X_\etale)$ is fully faithful and
$i_{small}^{-1}$ determines
a left inverse. On the other hand, suppose that $K$ is an object of
$D_Z(X_\etale)$ and consider the adjunction map
$K \to i_{small, *}i_{small}^{-1}K$.
Using exactness of $i_{small, *}$ and $i_{small}^{-1}$
this induces the adjunction maps
$H^n(K) \to i_{small, *}i_{small}^{-1}H^n(K)$ on cohomology sheaves.
Since these cohomology
sheaves are supported on $Z$ we see these adjunction maps are isomorphisms
and we conclude that $D(Z_\etale) \to D_Z(X_\etale)$ is an equivalence.

\medskip\noindent
To finish the proof we have to show that $R\mathcal{H}_Z(K) = i_{small}^{-1}K$
if $K$ is an object of $D_Z(X_\etale)$. To do this we can use that
$K = i_{small, *}i_{small}^{-1}K$
as we've just proved this is the case. Then we
can choose a K-injective representative $\mathcal{I}^\bullet$ for
$i_{small}^{-1}K$.
Since $i_{small, *}$ is the right adjoint to the exact functor
$i_{small}^{-1}$, the
complex $i_{small, *}\mathcal{I}^\bullet$ is K-injective
(Derived Categories, Lemma \ref{derived-lemma-adjoint-preserve-K-injectives}).
We see that $R\mathcal{H}_Z(K)$ is computed by
$\mathcal{H}_Z(i_{small, *}\mathcal{I}^\bullet) = \mathcal{I}^\bullet$
as desired.
\end{proof}

\begin{lemma}
\label{lemma-cohomology-with-support-quasi-coherent}
Let $X$ be a scheme. Let $Z \subset X$ be a closed subscheme.
Let $\mathcal{F}$ be a quasi-coherent $\mathcal{O}_X$-module
and denote $\mathcal{F}^a$ the associated quasi-coherent sheaf
on the small \'etale site of $X$
(Proposition \ref{proposition-quasi-coherent-sheaf-fpqc}). Then
\begin{enumerate}
\item $H^q_Z(X, \mathcal{F})$ agrees with $H^q_Z(X_\etale, \mathcal{F}^a)$,
\item if the complement of $Z$ is retrocompact in $X$, then
$i_*\mathcal{H}^q_Z(\mathcal{F}^a)$ is a quasi-coherent sheaf of
$\mathcal{O}_X$-modules equal to $(i_*\mathcal{H}^q_Z(\mathcal{F}))^a$.
\end{enumerate}
\end{lemma}

\begin{proof}
Let $j : U \to X$ be the inclusion of the complement of $Z$.
The statement (1) on cohomology groups follows from the long
exact sequences for cohomology with supports and the agreements
$H^q(X_\etale, \mathcal{F}^a) = H^q(X, \mathcal{F})$ and
$H^q(U_\etale, \mathcal{F}^a) = H^q(U, \mathcal{F})$, see
Theorem \ref{theorem-zariski-fpqc-quasi-coherent}.
If $j : U \to X$ is a quasi-compact morphism, i.e., if $U \subset X$
is retrocompact, then $R^qj_*$ transforms quasi-coherent sheaves
into quasi-coherent sheaves
(Cohomology of Schemes, Lemma
\ref{coherent-lemma-quasi-coherence-higher-direct-images})
and commutes with taking associated
sheaf on \'etale sites
(Descent, Lemma \ref{descent-lemma-higher-direct-images-small-etale}).
We conclude by applying
Lemma \ref{lemma-cohomology-with-support-triangle}.
\end{proof}







\section{Schemes with strictly henselian local rings}
\label{section-strictly-henselian-local-rings}

\noindent
In this section we collect some results about the \'etale cohomology
of schemes whose local rings are strictly henselian. For example,
here is a fun generalization of
Lemma \ref{lemma-vanishing-etale-cohomology-strictly-henselian}.

\begin{lemma}
\label{lemma-local-rings-strictly-henselian}
Let $S$ be a scheme all of whose local rings are strictly henselian.
Then for any abelian sheaf $\mathcal{F}$ on $S_\etale$ we have
$H^i(S_\etale, \mathcal{F}) = H^i(S_{Zar}, \mathcal{F})$.
\end{lemma}

\begin{proof}
Let $\epsilon : S_\etale \to S_{Zar}$ be the morphism of sites given
by the inclusion functor. The Zariski sheaf $R^p\epsilon_*\mathcal{F}$
is the sheaf associated to the presheaf $U \mapsto H^p_\etale(U, \mathcal{F})$.
Thus the stalk at $x \in X$ is
$\colim H^p_\etale(U, \mathcal{F}) =
H^p_\etale(\Spec(\mathcal{O}_{X, x}), \mathcal{G}_x)$
where $\mathcal{G}_x$ denotes the pullback of $\mathcal{F}$
to $\Spec(\mathcal{O}_{X, x})$, see
Lemma \ref{lemma-directed-colimit-cohomology}.
Thus the higher direct images of $R^p\epsilon_*\mathcal{F}$ are
zero by
Lemma \ref{lemma-vanishing-etale-cohomology-strictly-henselian}
and we conclude by the Leray spectral sequence.
\end{proof}

\begin{lemma}
\label{lemma-gabber-criterion}
Let $R$ be a ring all of whose local rings are strictly henselian.
Let $\mathcal{F}$ be a sheaf on $\Spec(R)_\etale$.
Assume that for all $f, g \in R$ the kernel of
$$
H^1_\etale(D(f + g), \mathcal{F})
\longrightarrow
H^1_\etale(D(f(f + g)), \mathcal{F}) \oplus
H^1_\etale(D(g(f + g)), \mathcal{F})
$$
is zero. Then $H^q_\etale(\Spec(R), \mathcal{F}) = 0$ for $q > 0$.
\end{lemma}

\begin{proof}
By Lemma \ref{lemma-local-rings-strictly-henselian} we see that
\'etale cohomology of $\mathcal{F}$ agrees with Zariski cohomology
on any open of $\Spec(R)$. We will prove by induction on $i$ the
statement: for $h \in R$ we have $H^q_\etale(D(h), \mathcal{F}) = 0$
for $1 \leq q \leq i$. The base case $i = 0$ is trivial. Assume $i \geq 1$.

\medskip\noindent
Let $\xi \in H^q_\etale(D(h), \mathcal{F})$ for some $1 \leq q \leq i$
and $h \in R$. If $q < i$ then we are done by induction, so we assume $q = i$.
After replacing $R$ by $R_h$ we may assume
$\xi \in H^i_\etale(\Spec(R), \mathcal{F})$; some details omitted.
Let $I \subset R$ be the set of elements $f \in R$ such that
$\xi|_{D(f)} = 0$. Since $\xi$ is Zariski locally trivial, it follows
that for every prime $\mathfrak p$ of $R$ there exists an $f \in I$
with $f \not \in \mathfrak p$. Thus if we can show that $I$ is an
ideal, then $1 \in I$ and we're done. It is clear that
$f \in I$, $r \in R$ implies $rf \in I$. Thus we assume that $f, g \in I$
and we show that $f + g \in I$. If $q = i = 1$, then this is exactly the
assumption of the lemma! Whence the result for $i = 1$. For
$q = i > 1$, note that
$$
D(f + g) = D(f(f + g)) \cup D(g(f + g))
$$
By Mayer-Vietoris (Cohomology, Lemma \ref{cohomology-lemma-mayer-vietoris}
which applies as \'etale cohomology on open subschemes of $\Spec(R)$ equals
Zariski cohomology) we have an exact sequence
$$
\xymatrix{
H^{i - 1}_\etale(D(fg(f + g)), \mathcal{F}) \ar[d] \\
H^i_\etale(D(f + g), \mathcal{F}) \ar[d] \\
H^i_\etale(D(f(f + g)), \mathcal{F}) \oplus
H^i_\etale(D(g(f + g)), \mathcal{F})
}
$$
and the result follows as the first group is zero by induction.
\end{proof}

\begin{lemma}
\label{lemma-affine-only-closed-points}
Let $S$ be an affine scheme such that
(1) all points are closed, and (2) all residue fields are separably
algebraically closed. Then
for any abelian sheaf $\mathcal{F}$ on $S_\etale$ we have
$H^i(S_\etale, \mathcal{F}) = 0$ for $i > 0$.
\end{lemma}

\begin{proof}
Condition (1) implies that the underlying topological space
of $S$ is profinite, see
Algebra, Lemma \ref{algebra-lemma-ring-with-only-minimal-primes}.
Thus the higher cohomology groups of an abelian sheaf on the topological
space $S$ (i.e., Zariski cohomology) is trivial, see
Cohomology, Lemma \ref{cohomology-lemma-vanishing-for-profinite}.
The local rings are strictly henselian by
Algebra, Lemma \ref{algebra-lemma-local-dimension-zero-henselian}.
Thus \'etale cohomology of $S$ is computed by Zariski cohomology
by Lemma \ref{lemma-local-rings-strictly-henselian}
and the proof is done.
\end{proof}

\noindent
The spectrum of an absolutely integrally closed ring
is an example of a scheme all of whose local rings are
strictly henselian, see More on Algebra, Lemma
\ref{more-algebra-lemma-absolutely-integrally-closed-strictly-henselian}.
It turns out that normal domains with separably closed fraction
fields have an even stronger property as explained in the
following lemma.

\begin{lemma}
\label{lemma-normal-scheme-with-alg-closed-function-field}
Let $X$ be an integral normal scheme with separably closed
function field.
\begin{enumerate}
\item A separated \'etale morphism $U \to X$ is a
disjoint union of open immersions.
\item All local rings of $X$ are strictly henselian.
\end{enumerate}
\end{lemma}

\begin{proof}
Let $R$ be a normal domain whose fraction field is separably algebraically
closed. Let $R \to A$ be an \'etale ring map. Then
$A \otimes_R K$ is as a $K$-algebra a finite product
$\prod_{i = 1, \ldots, n} K$ of copies of $K$. Let $e_i$, $i = 1, \ldots, n$
be the corresponding idempotents of $A \otimes_R K$. Since $A$ is normal
(Algebra, Lemma \ref{algebra-lemma-normal-goes-up})
the idempotents $e_i$ are in $A$
(Algebra, Lemma \ref{algebra-lemma-normal-ring-integrally-closed}).
Hence $A = \prod Ae_i$ and we may assume $A \otimes_R K = K$.
Since $A \subset A \otimes_R K = K$ (by flatness of $R \to A$ and
since $R \subset K$) we conclude that $A$ is a domain.
By the same argument we conclude that
$A \otimes_R A \subset (A \otimes_R A) \otimes_R K = K$.
It follows that the map $A \otimes_R A \to A$ is
injective as well as surjective. Thus $R \to A$ defines an
open immersion by
Morphisms, Lemma \ref{morphisms-lemma-universally-injective}
and
\'Etale Morphisms, Theorem \ref{etale-theorem-etale-radicial-open}.

\medskip\noindent
Let $f : U \to X$ be a separated \'etale morphism. Let $\eta \in X$
be the generic point and let $f^{-1}(\{\eta\}) = \{\xi_i\}_{i \in I}$.
The result of the previous paragraph shows the following:
For any affine open $U' \subset U$ whose image in $X$ is contained in
an affine we have $U' = \coprod_{i \in I} U'_i$ where $U'_i$
is the set of point of $U'$ which are specializations of $\xi_i$.
Moreover, the morphism $U'_i \to X$ is an open immersion.
It follows that $U_i = \overline{\{\xi_i\}}$ is an open and closed
subscheme of $U$ and that $U_i \to X$ is locally on the source
an isomorphism. By Morphisms,
Lemma \ref{morphisms-lemma-distinct-local-rings}
the fact that $U_i \to X$ is separated, implies that
$U_i \to X$ is injective and we conclude that $U_i \to X$
is an open immersion, i.e., (1) holds.

\medskip\noindent
Part (2) follows from part (1) and the description of the strict
henselization of $\mathcal{O}_{X, x}$ as the local ring at $\overline{x}$
on the \'etale site of $X$ (Lemma \ref{lemma-describe-etale-local-ring}).
It can also be proved directly, see
Fundamental Groups, Lemma
\ref{pione-lemma-normal-local-domain-separablly-closed-fraction-field}.
\end{proof}

\begin{lemma}
\label{lemma-Rf-star-zero-normal-with-alg-closed-function-field}
Let $f : X \to Y$ be a morphism of schemes where $X$ is an integral
normal scheme with separably closed function field. Then
$R^qf_*\underline{M} = 0$ for $q > 0$ and any abelian group $M$.
\end{lemma}

\begin{proof}
Recall that $R^qf_*\underline{M}$ is the sheaf associated
to the presheaf $V \mapsto H^q_\etale(V \times_Y X, M)$ on $Y_\etale$, see
Lemma \ref{lemma-higher-direct-images}.
If $V$ is affine, then $V \times_Y X \to X$ is separated and \'etale.
Hence $V \times_Y X = \coprod U_i$ is a disjoint union of open
subschemes $U_i$ of $X$, see
Lemma \ref{lemma-normal-scheme-with-alg-closed-function-field}.
By Lemma \ref{lemma-local-rings-strictly-henselian}
we see that $H^q_\etale(U_i, M)$ is equal to
$H^q_{Zar}(U_i, M)$. This vanishes by
Cohomology, Lemma \ref{cohomology-lemma-irreducible-constant-cohomology-zero}.
\end{proof}

\begin{lemma}
\label{lemma-closed-of-affine-normal-scheme-with-alg-closed-function-field}
Let $X$ be an affine integral normal scheme with separably closed
function field. Let $Z \subset X$ be a closed subscheme. Let
$V \to Z$ be an \'etale morphism with $V$ affine. Then $V$ is a finite
disjoint union of open subschemes of $Z$. If $V \to Z$ is
surjective and finite \'etale, then $V \to Z$ has a section.
\end{lemma}

\begin{proof}
By Algebra, Lemma \ref{algebra-lemma-lift-etale}
we can lift $V$ to an affine scheme $U$ \'etale over $X$.
Apply Lemma \ref{lemma-normal-scheme-with-alg-closed-function-field}
to $U \to X$ to get the first statement.

\medskip\noindent
The final statement is a consequence of the first.
Let $V = \coprod_{i = 1, \ldots, n} V_i$ be a finite
decomposition into open and
closed subschemes with $V_i \to Z$ an open immersion.
As $V \to Z$ is finite we see that $V_i \to Z$ is also closed.
Let $U_i \subset Z$ be the image. Then we have a decomposition
into open and closed subschemes
$$
Z =
\coprod\nolimits_{(A, B)}
\bigcap\nolimits_{i \in A} U_i \cap
\bigcap\nolimits_{i \in B} U_i^c
$$
where the disjoint union is over $\{1, \ldots, n\} = A \amalg B$
where $A$ has at least one element.
Each of the strata is contained in a single $U_i$ and
we find our section.
\end{proof}

\begin{lemma}
\label{lemma-gabber-for-h1-absolutely-algebraically-closed}
Let $X$ be a normal integral affine scheme with separably closed
function field. Let $Z \subset X$ be a closed subscheme.
For any finite abelian group $M$ we have $H^1_\etale(Z, \underline{M}) = 0$.
\end{lemma}

\begin{proof}
By Cohomology on Sites, Lemma \ref{sites-cohomology-lemma-torsors-h1}
an element of $H^1_\etale(Z, \underline{M})$ corresponds to a
$\underline{M}$-torsor $\mathcal{F}$ on $Z_\etale$.
Such a torsor is clearly a finite locally constant sheaf.
Hence $\mathcal{F}$ is representable by a scheme $V$ finite
\'etale over $Z$, Lemma \ref{lemma-characterize-finite-locally-constant}.
Of course $V \to Z$ is surjective as a torsor is locally trivial.
Since $V \to Z$ has a section by
Lemma \ref{lemma-closed-of-affine-normal-scheme-with-alg-closed-function-field}
we are done.
\end{proof}

\begin{lemma}
\label{lemma-gabber-for-absolutely-algebraically-closed}
Let $X$ be a normal integral affine scheme with separably closed
function field. Let $Z \subset X$ be a closed subscheme.
For any finite abelian group $M$ we have
$H^q_\etale(Z, \underline{M}) = 0$ for $q \geq 1$.
\end{lemma}

\begin{proof}
Write $X = \Spec(R)$ and $Z = \Spec(R')$ so that we have a surjection
of rings $R \to R'$. All local rings of $R'$ are strictly henselian by
Lemma \ref{lemma-normal-scheme-with-alg-closed-function-field} and
Algebra, Lemma \ref{algebra-lemma-quotient-strict-henselization}.
Furthermore, we see that for any $f' \in R'$ there is a surjection
$R_f \to R'_{f'}$ where $f \in R$ is a lift of $f'$.
Since $R_f$ is a normal domain with separably closed fraction field
we see that $H^1_\etale(D(f'), \underline{M}) = 0$ by
Lemma \ref{lemma-gabber-for-h1-absolutely-algebraically-closed}.
Thus we may apply Lemma \ref{lemma-gabber-criterion} to $Z = \Spec(R')$
to conclude.
\end{proof}

\begin{lemma}
\label{lemma-integral-cover-trivial-cohomology}
Let $X$ be an affine scheme.
\begin{enumerate}
\item There exists an integral surjective morphism $X' \to X$ such that for
every closed subscheme $Z' \subset X'$, every finite abelian group $M$, and
every $q \geq 1$ we have $H^q_\etale(Z', \underline{M}) = 0$.
\item For any closed subscheme $Z \subset X$, finite abelian group $M$,
$q \geq 1$, and $\xi \in H^q_\etale(Z, \underline{M})$ there exists a
finite surjective morphism $X' \to X$ of finite presentation such that
$\xi$ pulls back to zero in $H^q_\etale(X' \times_X Z, \underline{M})$.
\end{enumerate}
\end{lemma}

\begin{proof}
Write $X = \Spec(A)$. Write $A = \mathbf{Z}[x_i]/J$ for some ideal $J$.
Let $R$ be the integral closure of $\mathbf{Z}[x_i]$ in an algebraic
closure of the fraction field of $\mathbf{Z}[x_i]$. Let
$A' = R/JR$ and set $X' = \Spec(A')$. This gives an example as in (1) by
Lemma \ref{lemma-gabber-for-absolutely-algebraically-closed}.

\medskip\noindent
Proof of (2). Let $X' \to X$ be the integral surjective morphism we found
above. Certainly, $\xi$ maps to zero in
$H^q_\etale(X' \times_X Z, \underline{M})$. We may write $X'$ as a
limit $X' = \lim X'_i$ of schemes finite and of finite presentation
over $X$; this is easy to do in our current affine case, but it
is a special case of the more general Limits, Lemma
\ref{limits-lemma-integral-limit-finite-and-finite-presentation}.
By Lemma \ref{lemma-directed-colimit-cohomology}
we see that $\xi$ maps to zero in $H^q_\etale(X'_i \times_X Z, \underline{M})$
for some $i$ large enough.
\end{proof}








\section{Absolutely integrally closed vanishing}
\label{section-aic-vanishing}

\noindent
Recall that we say a ring $R$ is absolutely integrally closed
if every monic polynomial over $R$ has a root in $R$
(More on Algebra, Definition
\ref{more-algebra-definition-absolutely-integrally-closed}).
In this section we prove that the \'etale cohomology of
$\Spec(R)$ with coefficients in a finite torsion group
vanishes in positive degrees (Proposition \ref{proposition-aic-vanishing})
thereby slightly improving the earlier
Lemma \ref{lemma-gabber-for-absolutely-algebraically-closed}.
We suggest the reader skip this section.

\begin{lemma}
\label{lemma-find-extension}
Let $A$ be a ring. Let $a, b \in A$ such that
$aA + bA = A$ and $a \bmod bA$ is a root of unity.
Then there exists a monogenic extension $A \subset B$
and an element $y \in B$ such that $u = a - by$ is a unit.
\end{lemma}

\begin{proof}
Say $a^n \equiv 1 \bmod bA$. In particular $a^i$ is a unit modulo
$b^mA$ for all $i, m \geq 1$. We claim there exist
$a_1, \ldots, a_n \in A$ such that
$$
1 = a^n + a_1 a^{n - 1}b + a_2 a^{n - 2}b^2 + \ldots + a_n b^n
$$
Namely, since $1 - a^n \in bA$ we can find an element $a_1 \in A$
such that $1 - a^n - a_1 a^{n - 1} b \in b^2A$ using the unit property
of $a^{n - 1}$ modulo $bA$. Next, we can find an element $a_2 \in A$
such that $1 - a^n - a_1 a^{n - 1} b - a_2 a^{n - 2} b^2 \in b^3A$.
And so on. Eventually we find $a_1, \ldots, a_{n - 1} \in A$
such that $1 - (a^n + a_1 a^{n - 1}b + a_2 a^{n - 2}b^2 +
\ldots + a_{n - 1} ab^{n - 1}) \in b^nA$. This allows us to find
$a_n \in A$ such that the displayed equality holds.

\medskip\noindent
With $a_1, \ldots, a_n$ as above we claim that setting
$$
B = A[y]/(y^n + a_1 y^{n - 1} + a_2 y^{n - 2} + \ldots + a_n)
$$
works. Namely, suppose that $\mathfrak q \subset B$ is a prime ideal
lying over $\mathfrak p \subset A$. To get a contradiction
assume $u = a - by$ is in $\mathfrak q$. If $b \in \mathfrak p$
then $a \not \in \mathfrak p$ as $aA + bA = A$ and hence $u$ is not
in $\mathfrak q$. Thus we may assume $b \not \in \mathfrak p$, i.e.,
$b \not \in \mathfrak q$. This implies that $y \bmod \mathfrak q$
is equal to $a/b \bmod \mathfrak q$. However, then we obtain
$$
0 = y^n + a_1 y^{n - 1} + a_2 y^{n - 2} + \ldots + a_n =
b^{-n}(a^n + a_1 a^{n - 1}b + a_2 a^{n - 2}b^2 + \ldots + a_nb^n) =
b^{-n}
$$
a contradiction. This finishes the proof.
\end{proof}

\noindent
In order to explain the proof we need to introduce some group schemes.
Fix a prime number $\ell$. Let
$$
A = \mathbf{Z}[\zeta] =
\mathbf{Z}[x]/(x^{\ell - 1} + x^{\ell - 2} + \ldots + 1)
$$
In other words $A$ is the monogenic extension of $\mathbf{Z}$ generated by
a primitive $\ell$th root of unity $\zeta$. We set
$$
\pi = \zeta - 1
$$
A calculation (omitted) shows that $\ell$ is divisible by $\pi^{\ell - 1}$
in $A$. Our first group scheme over $A$ is
$$
G = \Spec(A[s, \frac{1}{\pi s + 1}])
$$
with group law given by the comultiplication
$$
\mu :
A[s, \frac{1}{\pi s + 1}]
\longrightarrow
A[s, \frac{1}{\pi s + 1}] \otimes_A A[s, \frac{1}{\pi s + 1}],\quad
s \longmapsto \pi s \otimes s + s \otimes 1 + 1 \otimes s
$$
With this choice we have
$$
\mu(\pi s + 1) = (\pi s + 1) \otimes (\pi s + 1)
$$
and hence we indeed have an $A$-algebra map as indicated. We omit the
verification that this indeed defines a group law. Our second group
scheme over $A$ is
$$
H = \Spec(A[t, \frac{1}{\pi^\ell t + 1}])
$$
with group law given by the comultiplication
$$
\mu : A[t, \frac{1}{\pi^\ell t + 1}]
\longrightarrow
A[t, \frac{1}{\pi^\ell t + 1}] \otimes_A A[t, \frac{1}{\pi^\ell t + 1}],\quad
t \longmapsto \pi^\ell t \otimes t + t \otimes 1 + 1 \otimes t
$$
The same verification as before shows that this defines a group law.
Next, we observe that the polynomial
$$
\Phi(s) = \frac{(\pi s + 1)^\ell - 1}{\pi^\ell}
$$
is in $A[s]$ and of degree $\ell$ and monic in $s$. Namely, the coefficicient
of $s^i$ for $0 < i < \ell$ is equal to ${\ell \choose i}\pi^{i - \ell}$
and since $\pi^{\ell - 1}$ divides $\ell$ in $A$ this is an element of $A$.
We obtain a ring map
$$
A[t, \frac{1}{\pi^\ell t + 1}]
\longrightarrow
A[s, \frac{1}{\pi s + 1}],\quad
t \longmapsto \Phi(s)
$$
which the reader easily verifies is compatible with the comultiplications.
Thus we get a morphism of group schemes
$$
f : G \to H
$$
The following lemma in particular shows that this morphism is faithfully
flat (in fact we will see that it is finite \'etale surjective).

\begin{lemma}
\label{lemma-monogenic-one}
We have
$$
A[s, \frac{1}{\pi s + 1}] =
\left(A[t, \frac{1}{\pi^\ell t + 1}]\right)[s]/(\Phi(s) - t)
$$
In particular, the Hopf algebra of $G$ is a monogenic extension
of the Hopf algebra of $H$.
\end{lemma}

\begin{proof}
Follows from the discussion above and the shape of $\Phi(s)$.
In particular, note that using $\Phi(s) = t$ the element
$\frac{1}{\pi^\ell t + 1}$ becomes the element
$\frac{1}{(\pi s + 1)^\ell}$.
\end{proof}

\noindent
Next, let us compute the kernel of $f$. Since the origin of $H$ is
given by $t = 0$ in $H$ we see that the kernel of $f$ is given
by $\Phi(s) = 0$. Now observe that the $A$-valued points
$\sigma_0, \ldots, \sigma_{\ell - 1}$
of $G$ given by
$$
\sigma_i :
s = \frac{\zeta^i - 1}{\pi} = \frac{\zeta^i - 1}{\zeta - 1} =
\zeta^{i - 1} + \zeta^{i - 2} + \ldots + 1,\quad
i = 0, 1, \ldots, \ell - 1
$$
are certainly contained in $\Ker(f)$. Moreover, these are all pairwise
distinct in {\bf all} fibres of $G \to \Spec(A)$. Also, the reader
computes that $\sigma_i +_G \sigma_j = \sigma_{i + j \bmod \ell}$.
Hence we find a closed immersion of group schemes
$$
\underline{\mathbf{Z}/\ell \mathbf{Z}}_A \longrightarrow \Ker(f)
$$
sending $i$ to $\sigma_i$. However, by construction $\Ker(f)$
is finite flat over $\Spec(A)$ of degree $\ell$. Hence we conclude
that this map is an isomorphism. All in all we conclude that
we have a short exact sequence
\begin{equation}
\label{equation-ses}
0 \to
\underline{\mathbf{Z}/\ell \mathbf{Z}}_A
\to G
\to H
\to 0
\end{equation}
of group schemes over $A$.

\begin{lemma}
\label{lemma-lift-points-H-to-G}
Let $R$ be an $A$-algebra which is absolutely integrally closed. Then
$G(R) \to H(R)$ is surjective.
\end{lemma}

\begin{proof}
Let $h \in H(R)$ correspond to the $A$-algebra map
$A[t, \frac{1}{\pi^\ell t + 1}] \to R$ sending $t$ to $a \in A$.
Since $\Phi(s)$ is monic we can find $b \in A$
with $\Phi(b) = a$. By Lemma \ref{lemma-monogenic-one}
sending $s$ to $b$ we obtain a unique $A$-algebra map
$A[s, \frac{1}{\pi s + 1}] \to R$ compatible
with the map $A[t, \frac{1}{\pi^\ell t + 1}] \to R$ above.
This in turn corresponds to an element $g \in G(R)$
mapping to $h \in H(R)$.
\end{proof}

\begin{lemma}
\label{lemma-interpolate}
Let $R$ be an $A$-algebra which is absolutely integrally closed.
Let $I, J \subset R$ be ideals with $I + J = R$.
There exists a $g \in G(R)$ such
that $g \bmod I = \sigma_0$ and $g \bmod J = \sigma_1$.
\end{lemma}

\begin{proof}
Choose $x \in I$ such that $x \equiv 1 \bmod J$.
We may and do replace $I$ by $xR$ and $J$ by $(x - 1)R$.
Then we are looking for an $s \in R$ such that
\begin{enumerate}
\item $1 + \pi s$ is a unit,
\item $s \equiv 0 \bmod xR$, and
\item $s \equiv 1 \bmod (x - 1)R$.
\end{enumerate}
The last two conditions say that $s = x + x(x - 1)y$ for some $y \in R$.
The first condition says that $1 + \pi s = 1 + \pi x + \pi x (x - 1) y$
needs to be a unit of $R$.
However, note that $1 + \pi x$ and $\pi x (x - 1)$ generate the
unit ideal of $R$ and that $1 + \pi x$ is an $\ell$th root of $1$
modulo $\pi x (x - 1)$\footnote{Because $1 + \pi x$ is congruent
to $1$ modulo $\pi$, congruent to $1$ modulo $x$, and congruent to
$1 + \pi = \zeta$ modulo $x - 1$ and because we have
$(\pi) \cap (x) \cap (x - 1) = (\pi x (x - 1))$ in $A[x]$.}.
Thus we win by Lemma \ref{lemma-find-extension} and the
fact that $R$ is absolutely integrally closed.
\end{proof}

\begin{proposition}
\label{proposition-aic-vanishing}
Let $R$ be an absolutely integrally closed ring.
Let $M$ be a finite abelian group.
Then $H^i_\etale(\Spec(R), \underline{M}) = 0$ for $i > 0$.
\end{proposition}

\begin{proof}
Since any finite abelian group has a finite filtration whose
subquotients are cyclic of prime order, we may assume
$M = \mathbf{Z}/\ell\mathbf{Z}$ where $\ell$ is a prime number.

\medskip\noindent
Observe that all local rings of $R$ are strictly henselian, see
More on Algebra, Lemma
\ref{more-algebra-lemma-absolutely-integrally-closed-strictly-henselian}.
Furthermore, any localization of $R$ is also absolutely integrally closed
by More on Algebra, Lemma
\ref{more-algebra-lemma-absolutely-integrally-closed-quotient-localization}.
Thus Lemma \ref{lemma-gabber-criterion} tells us it suffices to
show that the kernel of
$$
H^1_\etale(D(f + g), \mathbf{Z}/\ell\mathbf{Z})
\longrightarrow
H^1_\etale(D(f(f + g)), \mathbf{Z}/\ell\mathbf{Z}) \oplus
H^1_\etale(D(g(f + g)), \mathbf{Z}/\ell\mathbf{Z})
$$
is zero for any $f, g \in R$. After replacing $R$ by $R_{f + g}$ we
reduce to the following claim: given
$\xi \in H^1_\etale(\Spec(R), \mathbf{Z}/\ell\mathbf{Z})$
and an affine open covering $\Spec(R) = U \cup V$
such that $\xi|_U$ and $\xi|_V$ are trivial, then $\xi = 0$.

\medskip\noindent
Let $A = \mathbf{Z}[\zeta]$ as above.
Since $\mathbf{Z} \subset A$ is monogenic, we can find a ring
map $A \to R$. From now on we think of $R$ as an $A$-algebra
and we think of $\Spec(R)$ as a scheme over $\Spec(A)$.
If we base change the short exact sequence (\ref{equation-ses})
to $\Spec(R)$ and take \'etale cohomology we obtain
$$
G(R) \to H(R) \to
H^1_\etale(\Spec(R), \mathbf{Z}/\ell\mathbf{Z}) \to
H^1_\etale(\Spec(R), G)
$$
Please keep this in mind during the rest of the proof.

\medskip\noindent
Let $\tau \in \Gamma(U \cap V, \mathbf{Z}/\ell\mathbf{Z})$
be a section whose boundary in the Mayer-Vietoris sequence
(Lemma \ref{lemma-mayer-vietoris}) gives $\xi$.
For $i = 0, 1, \ldots, \ell - 1$ let $A_i \subset U \cap V$
be the open and closed subset where $\tau$ has the value $i \bmod \ell$.
Thus we have a finite disjoint union decomposition
$$
U \cap V = A_0 \amalg \ldots \amalg A_{\ell - 1}
$$
such that $\tau$ is constant on each $A_i$. For
$i = 0, 1, \ldots, \ell - 1$
denote $\tau_i \in H^0(U \cap V, \mathbf{Z}/\ell\mathbf{Z})$
the element which is equal to $1$ on $A_i$ and
equal to $0$ on $A_j$ for $j \not = i$.
Then $\tau$ is a sum of multiples of the $\tau_i$\footnote{Modulo
calculation errors we have $\tau = \sum i \tau_i$.}.
Hence it suffices to show that the cohomology class corresponding
to $\tau_i$ is trivial. This reduces us to the case where
$\tau$ takes only two distinct values, namely $1$ and $0$.

\medskip\noindent
Assume $\tau$ takes only the values $1$ and $0$. Write
$$
U \cap V = A \amalg B
$$
where $A$ is the locus where $\tau = 0$ and $B$ is the locus
where $\tau = 1$. Then $A$ and $B$ are disjoint closed subsets.
Denote $\overline{A}$ and $\overline{B}$ the closures of $A$ and $B$ in
$\Spec(R)$. Then we have a ``banana'': namely
we have
$$
\overline{A} \cap \overline{B} = Z_1 \amalg Z_2
$$
with $Z_1 \subset U$ and $Z_2 \subset V$ disjoint closed subsets.
Set $T_1 = \Spec(R) \setminus V$ and $T_2 = \Spec(R) \setminus U$.
Observe that $Z_1 \subset T_1 \subset U$, $Z_2 \subset T_2 \subset V$, and
$T_1 \cap T_2 = \emptyset$. Topologically we can write
$$
\Spec(R) = \overline{A} \cup \overline{B} \cup T_1 \cup T_2
$$
We suggest drawing a picture to visualize this.
In order to prove that $\xi$ is zero, we may and do replace $R$ by
its reduction (Proposition \ref{proposition-topological-invariance}).
Below, we think of $A$, $\overline{A}$, $B$, $\overline{B}$, $T_1$, $T_2$
as reduced closed subschemes of $\Spec(R)$. Next, as scheme structures on $Z_1$
and $Z_2$ we use
$$
Z_1 = \overline{A} \cap (\overline{B} \cup T_1)
\quad\text{and}\quad
Z_2 = \overline{A} \cap (\overline{B} \cup T_2)
$$
(scheme theoretic unions and intersections as in Morphisms, Definition
\ref{morphisms-definition-scheme-theoretic-intersection-union}).

\medskip\noindent
Denote $X$ the $G$-torsor over $\Spec(R)$ corresponding to the image
of $\xi$ in $H^1(\Spec(R), G)$. If $X$ is trivial, then $\xi$
comes from an element $h \in H(R)$ (see exact sequence of cohomology
above). However, then by Lemma \ref{lemma-lift-points-H-to-G}
the element $h$ lifts to an element of $G(R)$ and we conclude
$\xi = 0$ as desired. Thus our goal is to prove that $X$ is trivial.

\medskip\noindent
Recall that the embedding $\mathbf{Z}/\ell \mathbf{Z} \to G(R)$
sends $i \bmod \ell$ to $\sigma_i \in G(R)$. Observe that $\overline{A}$
is the spectrum of an absolutely integrally closed ring (namely
a qotient of $R$). By
Lemma \ref{lemma-interpolate} we can find $g \in G(\overline{A})$ with
$g|_{\overline{A} \cap Z_1} = \sigma_0$ and
$g|_{\overline{A} \cap Z_2} = \sigma_1$ (scheme theoretically).
Then we can define
\begin{enumerate}
\item $g_1 \in G(U)$ which is $g$ on $\overline{A} \cap U$, which
is $\sigma_0$ on $\overline{B} \cap U$, and $\sigma_0$ on $T_1$, and
\item $g_2 \in G(V)$ which is $g$ on $\overline{A} \cap V$, which
is $\sigma_1$ on $\overline{B} \cap V$, and $\sigma_1$ on $T_2$.
\end{enumerate}
Namely, to find $g_1$ as in (1) we glue the section
$\sigma_0$ on $\Omega = (\overline{B} \cup T_1) \cap U$
to the restriction of the section $g$ on $\Omega' = \overline{A} \cap U$.
Note that $U = \Omega \cup \Omega'$ (scheme theoretically)
because $U$ is reduced and $\Omega \cap \Omega' = Z_1$ (scheme theoretically)
by our choice of $Z_1$. Hence by
Morphisms, Lemma \ref{morphisms-lemma-scheme-theoretic-union}
we have that $U$ is the pushout of $\Omega$ and $\Omega'$
along $Z_1$. Thus we can find $g_1$.
Similarly for the existence of $g_2$ in (2).
Then we have
$$
\tau = g_2|_{A \cup B} - g_1|_{A \cup B} \quad(\text{addition in group law})
$$
and we see that $X$ is trivial thereby finishing the proof.
\end{proof}




















\section{Affine analog of proper base change}
\label{section-gabber-affine-proper}

\noindent
In this section we discuss a result by Ofer Gabber, see
\cite{gabber-affine-proper}. This was also proved by Roland Huber, see
\cite{Huber-henselian}. We have already done some of the work needed
for Gabber's proof in Section \ref{section-strictly-henselian-local-rings}.

\begin{lemma}
\label{lemma-efface-cohomology-on-closed-by-finite-cover}
Let $X$ be an affine scheme. Let $\mathcal{F}$ be a torsion abelian sheaf
on  $X_\etale$. Let $Z \subset X$ be a closed subscheme. Let
$\xi \in H^q_\etale(Z, \mathcal{F}|_Z)$ for some $q > 0$.
Then there exists an injective map $\mathcal{F} \to \mathcal{F}'$
of torsion abelian sheaves on $X_\etale$ such that
the image of $\xi$ in $H^q_\etale(Z, \mathcal{F}'|_Z)$ is zero.
\end{lemma}

\begin{proof}
By Lemmas \ref{lemma-torsion-colimit-constructible} and \ref{lemma-colimit}
we can find a map $\mathcal{G} \to \mathcal{F}$ with $\mathcal{G}$
a constructible abelian sheaf and $\xi$ coming from an element $\zeta$ of
$H^q_\etale(Z, \mathcal{G}|_Z)$. Suppose we can find an injective map
$\mathcal{G} \to \mathcal{G}'$ of torsion abelian sheaves on $X_\etale$
such that the image of $\zeta$ in $H^q_\etale(Z, \mathcal{G}'|_Z)$ is zero.
Then we can take $\mathcal{F}'$ to be the pushout
$$
\mathcal{F}' = \mathcal{G}' \amalg_{\mathcal{G}} \mathcal{F}
$$
and we conclude the result of the lemma holds. (Observe that restriction
to $Z$ is exact, so commutes with finite limits and colimits and moreover
it commutes with arbitrary colimits as a left adjoint to pushforward.)
Thus we may assume $\mathcal{F}$ is constructible.

\medskip\noindent
Assume $\mathcal{F}$ is constructible. By
Lemma \ref{lemma-constructible-maps-into-constant-general}
it suffices to prove the result when $\mathcal{F}$
is of the form $f_*\underline{M}$ where $M$ is a finite abelian group
and $f : Y \to X$ is a finite morphism of finite presentation
(such sheaves are still constructible by
Lemma \ref{lemma-finite-pushforward-constructible}
but we won't need this).
Since formation of $f_*$ commutes with any base change
(Lemma \ref{lemma-finite-pushforward-commutes-with-base-change})
we see that the restriction of $f_*\underline{M}$ to $Z$ is
equal to the pushforward of $\underline{M}$ via
$Y \times_X Z \to Z$. By the Leray spectral sequence
(Proposition \ref{proposition-leray})
and vanishing of higher direct images
(Proposition \ref{proposition-finite-higher-direct-image-zero}),
we find
$$
H^q_\etale(Z, f_*\underline{M}|_Z) = H^q_\etale(Y \times_X Z, \underline{M}).
$$
By Lemma \ref{lemma-integral-cover-trivial-cohomology}
we can find a finite surjective morphism $Y' \to Y$ of finite presentation
such that $\xi$ maps to zero in $H^q(Y' \times_X Z, \underline{M})$.
Denoting $f' : Y' \to X$ the composition $Y' \to Y \to X$ we claim
the map
$$
f_*\underline{M} \longrightarrow f'_*\underline{M}
$$
is injective which finishes the proof by what was said above.
To see the desired injectivity we can look at stalks. Namely,
if $\overline{x} : \Spec(k) \to X$ is a geometric point, then
$$
(f_*\underline{M})_{\overline{x}} =
\bigoplus\nolimits_{f(\overline{y}) = \overline{x}} M
$$
by Proposition \ref{proposition-finite-higher-direct-image-zero}
and similarly for the other sheaf.
Since $Y' \to Y$ is surjective and finite we see that
the induced map on geometric points lifting $\overline{x}$ is
surjective too and we conclude.
\end{proof}

\noindent
The lemma above will take care of higher cohomology groups in
Gabber's result. The following lemma will be used to deal with
global sections.

\begin{lemma}
\label{lemma-gabber-h0}
Let $X$ be a quasi-compact and quasi-separated scheme.
Let $i : Z \to X$ be a closed immersion. Assume that
\begin{enumerate}
\item for any sheaf $\mathcal{F}$ on $X_{Zar}$ the map
$\Gamma(X, \mathcal{F}) \to \Gamma(Z, i^{-1}\mathcal{F})$
is bijective, and
\item for any finite morphism $X' \to X$ assumption (1) holds
for $Z \times_X X' \to X'$.
\end{enumerate}
Then for any sheaf $\mathcal{F}$ on $X_\etale$ we have
$\Gamma(X, \mathcal{F}) = \Gamma(Z, i^{-1}_{small}\mathcal{F})$.
\end{lemma}

\begin{proof}
Let $\mathcal{F}$ be a sheaf on $X_\etale$. There is a canonical
(base change) map
$$
i^{-1}(\mathcal{F}|_{X_{Zar}})
\longrightarrow
(i_{small}^{-1}\mathcal{F})|_{Z_{Zar}}
$$
of sheaves on $Z_{Zar}$. We will show this map is injective by looking
at stalks. The stalk on the left hand side at $z \in Z$
is the stalk of $\mathcal{F}|_{X_{Zar}}$ at $z$. The stalk on the right
hand side is the colimit over all elementary \'etale neighbourhoods
$(U, u) \to (X, z)$ such that $U \times_X Z \to Z$ has a section over
a neighbourhood of $z$. As \'etale morphisms are open, the image of
$U \to X$ is an open neighbourhood $U_0$ of $z$ in $X$. The map
$\mathcal{F}(U_0) \to \mathcal{F}(U)$ is injective by the sheaf
condition for $\mathcal{F}$ with respect to the \'etale covering $U \to U_0$.
Taking the colimit over all $U$ and $U_0$ we obtain injectivity on stalks.

\medskip\noindent
It follows from this and assumption (1) that the map
$\Gamma(X, \mathcal{F}) \to \Gamma(Z, i^{-1}_{small}\mathcal{F})$
is injective. By (2) the same thing is true on all $X'$ finite over $X$.

\medskip\noindent
Before we prove the surjectivity, let us introduce some notation. For
any object $U$ of $X_\etale$ and $s \in \mathcal{F}(U)$ we
denote $i^{-1}s \in (i_{small}^{-1}\mathcal{F})(U \times_X Z)$
the pullback (more precisely this is the image of $s$ under the map
of Sites, Lemma \ref{sites-lemma-recover} combined with the map from
the presheaf pullback to the sheaf pullback). This construction is
compatible with pullback by morphisms in $X_\etale$ and by (finite)
morphisms $X' \to X$ in a manner which we leave to the reader to clarify.

\medskip\noindent
Let $t \in \Gamma(Z, i^{-1}_{small}\mathcal{F})$. By construction of
$i^{-1}_{small}\mathcal{F}$ there exists an \'etale covering
$\{V_j \to Z\}_{j \in J}$, \'etale morphisms $U_j \to X$, sections
$s_j \in \mathcal{F}(U_j)$ and morphisms $V_j \to U_j$ over $X$
such that $t|_{V_j}$ is the pullback of $i^{-1}s_j$ by
$V_j \to U_j \times_X Z$. Since $V_j \to U_j \times_X Z$ is \'etale,
the image is an open subscheme. Since $U_j \times_X Z \subset U_j$ is
closed, we may, after replacing $U_j$ by an open subscheme, assume
that $V_j \to U_j \times_X Z$ is surjective. Then $\{V_j \to U_j \times_X Z\}$
is an \'etale covering and we conclude that $t|_{U_j \times_X Z} = i^{-1}s_j$.
Observe that every nonempty closed subscheme $T \subset X$ meets $Z$
by assumption (1) applied to the sheaf $(T \to X)_*\underline{\mathbf{Z}}$
for example. Thus we see that $\coprod U_j \to X$ is surjective.
By More on Morphisms, Lemma
\ref{more-morphisms-lemma-there-is-a-scheme-integral-over}
we can find a finite surjective morphism $X' \to X$
such that $X' \to X$ Zariski locally factors through $\coprod U_j \to X$.
Say $X' = \bigcup_{\alpha \in A} W_\alpha$ is an open covering
and $j : A \to J$ is a map, such that there exist morphisms
$g_\alpha : W_\alpha \to U_{j(\alpha)}$ over $X$. Denote
$\mathcal{F}'$ the pullback of $\mathcal{F}$ to $X'_\etale$.
Denote $s'_\alpha \in \mathcal{F}'(W_\alpha)$ the pullback
of $s_{j(\alpha)}$ by $g_\alpha$.
Set $Z' = X' \times_X Z$. Denote $i' : Z' \to X'$
the base change of $i$. Then $(i'_{small})^{-1}\mathcal{F}' =
(Z' \to Z)_{small}^{-1}i_{small}^{-1}\mathcal{F}$. Via this identification
we may denote $t' \in
\Gamma(Z', (i')^{-1}_{small}\mathcal{F}')$
the pullback of $t$ to $Z'$.
Then $t'|_{W_\alpha \times_X Z}$ is equal to $(i')^{-1}s'_\alpha$ by
construction. We conclude that $t'$ is a section of
the subsheaf
$$
(i')^{-1}(\mathcal{F}'|_{X'_{Zar}})
\subset
(i'_{small})^{-1}\mathcal{F}')|_{Z'_{Zar}}
$$
By assumption (2) the section $t'$ comes from a section
$s' \in \Gamma(X', \mathcal{F}'|_{X'_{Zar}}) = \Gamma(X', \mathcal{F}')$.
By the injectivity proved in the second paragraph,
we conclude that the two pullbacks of $s'$ to
$X' \times_X X'$ are the same (after all this is true for
the pullbacks of $t'$ to $Z' \times_Z Z'$). Hence we conclude
$s'$ comes from a section of $\mathcal{F}$ over $X$ by
Remark \ref{remark-cohomological-descent-finite}.
\end{proof}

\begin{lemma}
\label{lemma-connected-topological}
Let $Z \subset X$ be a closed subset of a topological space $X$.
Assume
\begin{enumerate}
\item $X$ is a spectral space
(Topology, Definition \ref{topology-definition-spectral-space}), and
\item for $x \in X$ the intersection $Z \cap \overline{\{x\}}$
is connected (in particular nonempty).
\end{enumerate}
If $Z = Z_1 \amalg Z_2$ with $Z_i$ closed in $Z$,
then there exists a decomposition $X = X_1 \amalg X_2$ with
$X_i$ closed in $X$ and $Z_i = Z \cap X_i$.
\end{lemma}

\begin{proof}
Observe that $Z_i$ is quasi-compact. Hence the set of points
$W_i$ specializing to $Z_i$ is closed in the constructible topology
by Topology, Lemma \ref{topology-lemma-make-spectral-space}.
Assumption (2) implies that $X = W_1 \amalg W_2$.
Let $x \in \overline{W_1}$. By Topology, Lemma
\ref{topology-lemma-constructible-stable-specialization-closed} part (1)
there exists a specialization $x_1 \leadsto x$ with $x_1 \in W_1$.
Thus $\overline{\{x\}} \subset \overline{\{x_1\}}$ and we see
that $x \in W_1$. In other words, setting $X_i = W_i$ does the job.
\end{proof}

\begin{lemma}
\label{lemma-h0-topological}
Let $Z \subset X$ be a closed subset of a topological space $X$.
Assume
\begin{enumerate}
\item $X$ is a spectral space
(Topology, Definition \ref{topology-definition-spectral-space}), and
\item for $x \in X$ the intersection $Z \cap \overline{\{x\}}$
is connected (in particular nonempty).
\end{enumerate}
Then for any sheaf $\mathcal{F}$ on $X$ we have
$\Gamma(X, \mathcal{F}) = \Gamma(Z, \mathcal{F}|_Z)$.
\end{lemma}

\begin{proof}
If $x \leadsto x'$ is a specialization of points, then there is a
canonical map $\mathcal{F}_{x'} \to \mathcal{F}_x$ compatible with
sections over opens and functorial in $\mathcal{F}$. Since every point
of $X$ specializes to a point of $Z$ it follows that
$\Gamma(X, \mathcal{F}) \to \Gamma(Z, \mathcal{F}|_Z)$ is injective.
The difficult part is to show that it is surjective.

\medskip\noindent
Denote $\mathcal{B}$ be the set of all quasi-compact opens of $X$.
Write $\mathcal{F}$ as a filtered colimit $\mathcal{F} = \colim \mathcal{F}_i$
where each $\mathcal{F}_i$ is as in
Modules, Equation (\ref{modules-equation-towards-constructible-sets}).
See Modules, Lemma \ref{modules-lemma-filtered-colimit-constructibles}.
Then $\mathcal{F}|_Z = \colim \mathcal{F}_i|_Z$ as restriction to $Z$
is a left adjoint (Categories, Lemma \ref{categories-lemma-adjoint-exact} and
Sheaves, Lemma \ref{sheaves-lemma-f-map}).
By Sheaves, Lemma \ref{sheaves-lemma-directed-colimits-sections}
the functors $\Gamma(X, -)$ and $\Gamma(Z, -)$ commute with filtered colimits.
Hence we may assume our sheaf $\mathcal{F}$ is as in
Modules, Equation (\ref{modules-equation-towards-constructible-sets}).

\medskip\noindent
Suppose that we have an embedding $\mathcal{F} \subset \mathcal{G}$.
Then we have
$$
\Gamma(X, \mathcal{F}) =
\Gamma(Z, \mathcal{F}|_Z) \cap \Gamma(X, \mathcal{G})
$$
where the intersection takes place in $\Gamma(Z, \mathcal{G}|_Z)$.
This follows from the first remark of the proof because we can check
whether a global section of $\mathcal{G}$ is in $\mathcal{F}$ by looking
at the stalks and because every point of $X$ specializes to a point of $Z$.

\medskip\noindent
By Modules, Lemma \ref{modules-lemma-constructible-in-constant}
there is an injection $\mathcal{F} \to \prod (Z_i \to X)_*\underline{S_i}$
where the product is finite, $Z_i \subset X$ is closed, and $S_i$ is finite.
Thus it suffices to prove surjectivity for the sheaves
$(Z_i \to X)_*\underline{S_i}$. Observe that
$$
\Gamma(X, (Z_i \to X)_*\underline{S_i}) = \Gamma(Z_i, \underline{S_i})
\quad\text{and}\quad
\Gamma(X, (Z_i \to X)_*\underline{S_i}|_Z) =
\Gamma(Z \cap Z_i, \underline{S_i})
$$
Moreover, conditions (1) and (2) are inherited by $Z_i$; this is clear
for (2) and follows from
Topology, Lemma \ref{topology-lemma-spectral-sub} for (1). Thus it
suffices to prove the lemma in the case of a (finite) constant sheaf.
This case is a restatement of Lemma \ref{lemma-connected-topological}
which finishes the proof.
\end{proof}

\begin{example}
\label{example-quinard}
Lemma \ref{lemma-h0-topological} is false if $X$ is not spectral.
Here is an example: Let $Y$ be a $T_1$ topological space, and
$y \in Y$ a non-open point. Let $X = Y \amalg \{ x \}$, endowed with
the topology whose closed sets are $\emptyset$, $\{y\}$, and all
$F \amalg \{ x \}$, where $F$ is a closed subset of $Y$. Then
$Z = \{x, y\}$ is a closed subset of $X$, which satisfies assumption (2)
of Lemma \ref{lemma-h0-topological}. But $X$ is connected, while $Z$ is not.
The conclusion of the lemma thus fails for the constant sheaf
with value $\{0, 1\}$ on $X$.
\end{example}

\begin{lemma}
\label{lemma-h0-henselian-pair}
Let $(A, I)$ be a henselian pair. Set $X = \Spec(A)$ and
$Z = \Spec(A/I)$. For any sheaf $\mathcal{F}$ on $X_\etale$
we have $\Gamma(X, \mathcal{F}) = \Gamma(Z, \mathcal{F}|_Z)$.
\end{lemma}

\begin{proof}
Recall that the spectrum of any ring is a spectral space, see
Algebra, Lemma \ref{algebra-lemma-spec-spectral}. By
More on Algebra, Lemma
\ref{more-algebra-lemma-irreducible-henselian-pair-connected}
we see that $\overline{\{x\}} \cap Z$ is connected for every $x \in X$.
By Lemma \ref{lemma-h0-topological} we see that the statement
is true for sheaves on $X_{Zar}$. For any finite morphism $X' \to X$
we have $X' = \Spec(A')$ and $Z \times_X X' = \Spec(A'/IA')$
with $(A', IA')$ a henselian pair, see More on Algebra, Lemma
\ref{more-algebra-lemma-integral-over-henselian-pair}
and we get the same statement for sheaves on $(X')_{Zar}$.
Thus we can apply Lemma \ref{lemma-gabber-h0} to conclude.
\end{proof}

\noindent
Finally, we can state and prove Gabber's theorem.

\begin{theorem}[Gabber]
\label{theorem-gabber}
Let $(A, I)$ be a henselian pair. Set $X = \Spec(A)$ and
$Z = \Spec(A/I)$. For any torsion abelian sheaf $\mathcal{F}$ on $X_\etale$
we have $H^q_\etale(X, \mathcal{F}) = H^q_\etale(Z, \mathcal{F}|_Z)$.
\end{theorem}

\begin{proof}
The result holds for $q = 0$ by Lemma \ref{lemma-h0-henselian-pair}.
Let $q \geq 1$. Suppose the result has been shown in all degrees $< q$.
Let $\mathcal{F}$ be a torsion abelian sheaf. Let
$\mathcal{F} \to \mathcal{F}'$
be an injective map of torsion abelian sheaves (to be chosen later)
with cokernel $\mathcal{Q}$ so that we have the short exact sequence
$$
0 \to \mathcal{F} \to \mathcal{F}' \to \mathcal{Q} \to 0
$$
of torsion abelian sheaves on $X_\etale$. This gives a map of long exact
cohomology sequences over $X$ and $Z$ part of which looks like
$$
\xymatrix{
H^{q - 1}_\etale(X, \mathcal{F}') \ar[d] \ar[r] &
H^{q - 1}_\etale(X, \mathcal{Q}) \ar[d] \ar[r] &
H^q_\etale(X, \mathcal{F}) \ar[d] \ar[r] &
H^q_\etale(X, \mathcal{F}') \ar[d] \\
H^{q - 1}_\etale(Z, \mathcal{F}'|_Z) \ar[r] &
H^{q - 1}_\etale(Z, \mathcal{Q}|_Z) \ar[r] &
H^q_\etale(Z, \mathcal{F}|_Z) \ar[r] &
H^q_\etale(Z, \mathcal{F}'|_Z)
}
$$
Using this commutative diagram of abelian groups with exact rows
we will finish the proof.

\medskip\noindent
Injectivity for $\mathcal{F}$. Let $\xi$ be a nonzero element of
$H^q_\etale(X, \mathcal{F})$. By
Lemma \ref{lemma-efface-cohomology-on-closed-by-finite-cover} applied with
$Z = X$ (!) we can find $\mathcal{F} \subset \mathcal{F}'$ such that
$\xi$ maps to zero to the right. Then $\xi$ is the image of
an element of $H^{q - 1}_\etale(X, \mathcal{Q})$ and bijectivity
for $q - 1$ implies $\xi$ does not map to zero in
$H^q_\etale(Z, \mathcal{F}|_Z)$.

\medskip\noindent
Surjectivity for $\mathcal{F}$. Let $\xi$ be an element of
$H^q_\etale(Z, \mathcal{F}|_Z)$. By
Lemma \ref{lemma-efface-cohomology-on-closed-by-finite-cover} applied with
$Z = Z$ we can find $\mathcal{F} \subset \mathcal{F}'$ such that
$\xi$ maps to zero to the right. Then $\xi$ is the image of
an element of $H^{q - 1}_\etale(Z, \mathcal{Q}|_Z)$ and bijectivity
for $q - 1$ implies $\xi$ is in the image of the vertical map.
\end{proof}

\begin{lemma}
\label{lemma-vanishing-restriction-injective}
Let $X$ be a scheme with affine diagonal which can be covered by
$n + 1$ affine opens. Let $Z \subset X$ be a closed subscheme.
Let $\mathcal{A}$ be a torsion sheaf of rings on $X_\etale$
and let $\mathcal{I}$ be an injective sheaf of $\mathcal{A}$-modules
on $X_\etale$.
Then $H^q_\etale(Z, \mathcal{I}|_Z) = 0$ for $q > n$.
\end{lemma}

\begin{proof}
We will prove this by induction on $n$. If $n = 0$, then $X$ is affine.
Say $X = \Spec(A)$ and $Z = \Spec(A/I)$. Let $A^h$ be the filtered colimit
of \'etale $A$-algebras $B$ such that $A/I \to B/IB$ is an isomorphism.
Then $(A^h, IA^h)$ is a henselian pair and $A/I = A^h/IA^h$, see
More on Algebra, Lemma \ref{more-algebra-lemma-henselization}
and its proof. Set $X^h = \Spec(A^h)$.
By Theorem \ref{theorem-gabber}
we see that
$$
H^q_\etale(Z, \mathcal{I}|_Z) = H^q_\etale(X^h, \mathcal{I}|_{X^h})
$$
By Theorem \ref{theorem-colimit} we have
$$
H^q_\etale(X^h, \mathcal{I}|_{X^h}) =
\colim_{A \to B} H^q_\etale(\Spec(B), \mathcal{I}|_{\Spec(B)})
$$
where the colimit is over the $A$-algebras $B$ as above.
Since the morphisms $\Spec(B) \to \Spec(A)$ are \'etale,
the restriction $\mathcal{I}|_{\Spec(B)}$ is an injective
sheaf of $\mathcal{A}|_{\Spec(B)}$-modules
(Cohomology on Sites, Lemma \ref{sites-cohomology-lemma-cohomology-of-open}).
Thus the cohomology groups on the right are zero and we get the
result in this case.

\medskip\noindent
Induction step. We can use Mayer-Vietoris to do the induction step.
Namely, suppose that $X = U \cup V$ where $U$ is a union of $n$ affine
opens and $V$ is affine. Then, using that the diagonal of $X$ is affine,
we see that $U \cap V$ is the union of $n$ affine opens. Mayer-Vietoris
gives an exact sequence
$$
H^{q - 1}_\etale(U \cap V \cap Z, \mathcal{I}|_Z) \to
H^q_\etale(Z, \mathcal{I}|_Z) \to
H^q_\etale(U \cap Z, \mathcal{I}|_Z) \oplus
H^q_\etale(V \cap Z, \mathcal{I}|_Z)
$$
and by our induction hypothesis we obtain vanishing for $q > n$ as desired.
\end{proof}





\section{Cohomology of torsion sheaves on curves}
\label{section-vanishing-torsion}

\noindent
The goal of this section is to prove the basic finiteness and vanishing
results for cohomology of torsion sheaves on curves, see
Theorem \ref{theorem-vanishing-affine-curves}.
In Section \ref{section-vanishing-torsion-coefficients}
we will discuss constructible sheaves of torsion modules
over a Noetherian ring.

\begin{situation}
\label{situation-what-to-prove}
Here $k$ is an algebraically closed field, $X$ is a separated, finite type
scheme of dimension $\leq 1$ over $k$, and $\mathcal{F}$ is a torsion
abelian sheaf on $X_\etale$.
\end{situation}

\noindent
In Situation \ref{situation-what-to-prove}
we want to prove the following statements
\begin{enumerate}
\item
\label{item-vanishing}
$H^q_\etale(X, \mathcal{F}) = 0$ for $q > 2$,
\item
\label{item-vanishing-affine}
$H^q_\etale(X, \mathcal{F}) = 0$ for $q > 1$ if $X$ is affine,
\item
\label{item-vanishing-p-p}
$H^q_\etale(X, \mathcal{F}) = 0$ for $q > 1$ if $p = \text{char}(k) > 0$
and $\mathcal{F}$ is $p$-power torsion,
\item
\label{item-finite-prime-to-p}
$H^q_\etale(X, \mathcal{F})$ is finite if $\mathcal{F}$ is
constructible and torsion prime to $\text{char}(k)$,
\item
\label{item-finite-proper}
$H^q_\etale(X, \mathcal{F})$ is finite if $X$ is proper and
$\mathcal{F}$ constructible,
\item
\label{item-base-change-prime-to-p}
$H^q_\etale(X, \mathcal{F}) \to
H^q_\etale(X_{k'}, \mathcal{F}|_{X_{k'}})$ is an isomorphism
for any extension $k'/k$ of algebraically closed fields
if $\mathcal{F}$ is torsion prime to $\text{char}(k)$,
\item
\label{item-base-change-proper}
$H^q_\etale(X, \mathcal{F}) \to
H^q_\etale(X_{k'}, \mathcal{F}|_{X_{k'}})$ is an isomorphism
for any extension $k'/k$ of algebraically closed fields
if $X$ is proper,
\item
\label{item-surjective}
$H^2_\etale(X, \mathcal{F}) \to H^2_\etale(U, \mathcal{F})$
is surjective for all $U \subset X$ open.
\end{enumerate}
Given any Situation \ref{situation-what-to-prove}
we will say that
``statements (\ref{item-vanishing}) -- (\ref{item-surjective}) hold''
if those statements that apply to the given situation are true.
We start the proof with the following consequence of our computation
of cohomology with constant coefficients.

\begin{lemma}
\label{lemma-constant-smooth-statements}
In Situation \ref{situation-what-to-prove}
assume $X$ is smooth and $\mathcal{F} = \underline{\mathbf{Z}/\ell\mathbf{Z}}$
for some prime number $\ell$. Then statements
(\ref{item-vanishing}) -- (\ref{item-surjective}) hold
for $\mathcal{F}$.
\end{lemma}

\begin{proof}
Since $X$ is smooth, we see that $X$ is a finite disjoint union of
smooth curves and points. The higher \'etale cohomology of any sheaf
on a point (the spectrum of $k$) is trivial, for example by
Lemma \ref{lemma-compare-cohomology-point}
or \ref{lemma-affine-only-closed-points}.
Hence we may assume $X$ is a smooth curve.

\medskip\noindent
Case I: $\ell$ different from the characteristic of $k$.
This case follows from
Lemma \ref{lemma-cohomology-smooth-projective-curve}
(projective case) and
Lemma \ref{lemma-vanishing-cohomology-mu-smooth-curve}
(affine case). Statement (\ref{item-base-change-prime-to-p})
on cohomology and extension of algebraically closed ground
field follows from the fact that the genus $g$ and the number
of ``punctures'' $r$ do not change when passing from $k$ to $k'$.
Statement (\ref{item-surjective}) follows as $H^2_\etale(U, \mathcal{F})$
is zero as soon as $U \not = X$, because then $U$ is affine
(Varieties, Lemmas \ref{varieties-lemma-proper-minus-point} and
\ref{varieties-lemma-curve-affine-projective}).

\medskip\noindent
Case II: $\ell$ is equal to the characteristic of $k$.
Vanishing by Lemma \ref{lemma-vanishing-variety-char-p-p}.
Statements (\ref{item-finite-proper}) and (\ref{item-base-change-proper})
follow from
Lemma \ref{lemma-finiteness-proper-variety-char-p-p}.
\end{proof}

\begin{remark}[Invariance under extension of algebraically closed ground field]
\label{remark-invariance}
Let $k$ be an algebraically closed field of characteristic $p > 0$.
In Section \ref{section-artin-schreier} we have seen that there is
an exact sequence
$$
k[x] \to k[x] \to
H^1_\etale(\mathbf{A}^1_k, \mathbf{Z}/p\mathbf{Z}) \to 0
$$
where the first arrow maps $f(x)$ to $f^p - f$. A set of representatives
for the cokernel is formed by the polynomials
$$
\sum\nolimits_{p \not | n} \lambda_n x^n
$$
with $\lambda_n \in k$. (If $k$ is not algebraically closed
you have to add some constants to this as well.) In particular
when $k'/k$ is an algebraically closed extension, then
the map
$$
H^1_\etale(\mathbf{A}^1_k, \mathbf{Z}/p\mathbf{Z})
\to
H^1_\etale(\mathbf{A}^1_{k'}, \mathbf{Z}/p\mathbf{Z})
$$
is not an isomorphism in general. In particular, the map
$\pi_1(\mathbf{A}^1_{k'}) \to \pi_1(\mathbf{A}^1_k)$
between \'etale fundamental groups (insert future reference here)
is not an isomorphism either. Thus the \'etale homotopy type
of the affine line depends on the algebraically closed ground field.
From Lemma \ref{lemma-constant-smooth-statements} above we see that
this is a phenomenon which only happens in characteristic $p$
with $p$-power torsion coefficients.
\end{remark}

\begin{lemma}
\label{lemma-ses-statements}
Let $k$ be an algebraically closed field. Let $X$ be a separated finite
type scheme over $k$ of dimension $\leq 1$. Let
$0 \to \mathcal{F}_1 \to \mathcal{F} \to \mathcal{F}_2 \to 0$
be a short exact sequence of torsion abelian sheaves on $X$.
If statements (\ref{item-vanishing}) -- (\ref{item-surjective}) hold
for $\mathcal{F}_1$ and $\mathcal{F}_2$, then they hold
for $\mathcal{F}$.
\end{lemma}

\begin{proof}
This is mostly immediate from the definitions and the long exact sequence
of cohomology. Also observe that $\mathcal{F}$ is constructible
(resp.\ of torsion prime to the characteristic of $k$) if and only if
both $\mathcal{F}_1$ and $\mathcal{F}_2$ are constructible
(resp.\ of torsion prime to the characteristic of $k$). See
Proposition \ref{proposition-constructible-over-noetherian}.
Some details omitted.
\end{proof}

\begin{lemma}
\label{lemma-finite-pushforward-statements}
Let $k$ be an algebraically closed field. Let $f : X \to Y$ be a
finite morphism of separated finite type schemes over $k$ of
dimension $\leq 1$. Let $\mathcal{F}$ be a torsion abelian sheaf on $X$.
If statements (\ref{item-vanishing}) -- (\ref{item-surjective}) hold
for $\mathcal{F}$, then they hold for $f_*\mathcal{F}$.
\end{lemma}

\begin{proof}
Namely, we have $H^q_\etale(X, \mathcal{F}) = H^q_\etale(Y, f_*\mathcal{F})$
by the vanishing of $R^qf_*$ for $q > 0$
(Proposition \ref{proposition-finite-higher-direct-image-zero}) and
the Leray spectral sequence
(Cohomology on Sites, Lemma \ref{sites-cohomology-lemma-apply-Leray}).
For (\ref{item-surjective}) use that formation of $f_*$ commutes with
arbitrary base change
(Lemma \ref{lemma-finite-pushforward-commutes-with-base-change}).
\end{proof}

\begin{lemma}
\label{lemma-restrict-to-open}
In Situation \ref{situation-what-to-prove} assume $\mathcal{F}$
constructible. Let $j : X' \to X$ be the inclusion of a dense open subscheme.
Then statements
(\ref{item-vanishing}) -- (\ref{item-surjective}) hold for $\mathcal{F}$
if and only if they hold for $j_!j^{-1}\mathcal{F}$.
\end{lemma}

\begin{proof}
Since $X'$ is dense, we see that $Z = X \setminus X'$ has dimension $0$
and hence is a finite set $Z = \{x_1, \ldots, x_n\}$ of $k$-rational points.
Consider the short exact sequence
$$
0 \to j_!j^{-1}\mathcal{F} \to \mathcal{F} \to i_*i^{-1}\mathcal{F} \to 0
$$
of Lemma \ref{lemma-ses-associated-to-open}. Observe that
$H^q_\etale(X, i_*i^{-1}\mathcal{F}) = H^q_\etale(Z, i^*\mathcal{F})$.
Namely, $i : Z \to X$ is a closed immersion, hence finite,
hence we have the vanishing of $R^qi_*$ for $q > 0$ by
Proposition \ref{proposition-finite-higher-direct-image-zero}, and
hence the equality follows from the Leray spectral sequence
(Cohomology on Sites, Lemma \ref{sites-cohomology-lemma-apply-Leray}).
Since $Z$ is a disjoint union of spectra of algebraically closed
fields, we conclude that $H^q_\etale(Z, i^*\mathcal{F}) = 0$
for $q > 0$ and
$$
H^0_\etale(Z, i^{-1}\mathcal{F}) =
\bigoplus\nolimits_{i = 1, \ldots, n} \mathcal{F}_{x_i}
$$
which is finite as $\mathcal{F}_{x_i}$ is finite due to the
assumption that $\mathcal{F}$ is constructible.
The long exact cohomology sequence gives an exact sequence
$$
0 \to
H^0_\etale(X, j_!j^{-1}\mathcal{F}) \to
H^0_\etale(X, \mathcal{F}) \to
H^0_\etale(Z, i^{-1}\mathcal{F}) \to
H^1_\etale(X, j_!j^{-1}\mathcal{F}) \to
H^1_\etale(X, \mathcal{F}) \to 0
$$
and isomorphisms $H^q_\etale(X, j_!j^{-1}\mathcal{F}) \to
H^q_\etale(X, \mathcal{F})$ for $q > 1$.

\medskip\noindent
At this point it is easy to deduce each of 
(\ref{item-vanishing}) -- (\ref{item-surjective})
holds for $\mathcal{F}$ if and only if it holds for
$j_!j^{-1}\mathcal{F}$. We make a few small remarks to help the reader:
(a) if $\mathcal{F}$ is torsion prime to the characteristic
of $k$, then so is $j_!j^{-1}\mathcal{F}$,
(b) the sheaf $j_!j^{-1}\mathcal{F}$ is constructible,
(c) we have $H^0_\etale(Z, i^{-1}\mathcal{F}) =
H^0_\etale(Z_{k'}, i^{-1}\mathcal{F}|_{Z_{k'}})$, and (d)
if $U \subset X$ is an open, then $U' = U \cap X'$ is dense in $U$.
\end{proof}

\begin{lemma}
\label{lemma-even-easier}
In Situation \ref{situation-what-to-prove} assume $X$ is smooth.
Let $j : U \to X$ an open immersion. Let $\ell$ be a prime number.
Let $\mathcal{F} = j_!\underline{\mathbf{Z}/\ell\mathbf{Z}}$.
Then statements (\ref{item-vanishing}) -- (\ref{item-surjective}) hold
for $\mathcal{F}$.
\end{lemma}

\begin{proof}
Since $X$ is smooth, it is a disjoint union of smooth curves
and hence we may assume $X$ is a curve (i.e., irreducible).
Then either $U = \emptyset$ and there is nothing to prove
or $U \subset X$ is dense. In this case the lemma follows from
Lemmas \ref{lemma-constant-smooth-statements} and
\ref{lemma-restrict-to-open}.
\end{proof}

\begin{lemma}
\label{lemma-somewhat-easier}
In Situation \ref{situation-what-to-prove} assume $X$ reduced.
Let $j : U \to X$ an open immersion. Let $\ell$ be a prime number
and $\mathcal{F} = j_!\underline{\mathbf{Z}/\ell\mathbf{Z}}$.
Then statements (\ref{item-vanishing}) -- (\ref{item-surjective}) hold
for $\mathcal{F}$.
\end{lemma}

\begin{proof}
The difference with Lemma \ref{lemma-even-easier} is that here we do not
assume $X$ is smooth. Let $\nu : X^\nu \to X$ be the normalization
morphism. Then $\nu$ is finite
(Varieties, Lemma \ref{varieties-lemma-normalization-locally-algebraic})
and $X^\nu$ is smooth
(Varieties, Lemma \ref{varieties-lemma-regular-point-on-curve}).
Let $j^\nu : U^\nu \to X^\nu$ be the inverse image of $U$.
By Lemma \ref{lemma-even-easier} the result holds for
$j^\nu_!\underline{\mathbf{Z}/\ell\mathbf{Z}}$.
By Lemma \ref{lemma-finite-pushforward-statements}
the result holds for $\nu_*j^\nu_!\underline{\mathbf{Z}/\ell\mathbf{Z}}$.
In general it won't be true that
$\nu_*j^\nu_!\underline{\mathbf{Z}/\ell\mathbf{Z}}$ is equal to
$j_!\underline{\mathbf{Z}/\ell\mathbf{Z}}$
but we can work around this as follows.
As $X$ is reduced the morphism $\nu : X^\nu \to X$
is an isomorphism over a dense open $j' : X' \to X$
(Varieties, Lemma \ref{varieties-lemma-normalization-locally-algebraic}).
Over this open we have agreement
$$
(j')^{-1}(\nu_*j^\nu_!\underline{\mathbf{Z}/\ell\mathbf{Z}})
=
(j')^{-1}(j_!\underline{\mathbf{Z}/\ell\mathbf{Z}})
$$
Using Lemma \ref{lemma-restrict-to-open}
twice for $j' : X' \to X$ and the sheaves above
we conclude.
\end{proof}

\begin{lemma}
\label{lemma-vanishing-easier}
In Situation \ref{situation-what-to-prove} assume $X$ reduced.
Let $j : U \to X$ an open immersion with $U$ connected. Let
$\ell$ be a prime number. Let $\mathcal{G}$ a finite locally
constant sheaf of $\mathbf{F}_\ell$-vector spaces on $U$. Let
$\mathcal{F} = j_!\mathcal{G}$. Then statements
(\ref{item-vanishing}) -- (\ref{item-surjective}) hold for $\mathcal{F}$.
\end{lemma}

\begin{proof}
Let $f : V \to U$ be a finite \'etale morphism of degree prime to $\ell$
as in Lemma \ref{lemma-pullback-filtered}. The discussion in
Section \ref{section-trace-method} gives maps
$$
\mathcal{G} \to f_*f^{-1}\mathcal{G} \to \mathcal{G}
$$
whose composition is an isomorphism. Hence it suffices to prove the
lemma with $\mathcal{F} = j_!f_*f^{-1}\mathcal{G}$.
By Zariski's Main theorem
(More on Morphisms, Lemma
\ref{more-morphisms-lemma-quasi-finite-separated-pass-through-finite})
we can choose a diagram
$$
\xymatrix{
V \ar[r]_{j'} \ar[d]_f & Y \ar[d]^{\overline{f}} \\
U \ar[r]^j & X
}
$$
with $\overline{f} : Y \to X$ finite and $j'$ an open immersion
with dense image. We may replace $Y$ by its reduction (this does
not change $V$ as $V$ is reduced being \'etale over $U$).
Since $f$ is finite and $V$ dense in $Y$ we have $V = U \times_X Y$. By
Lemma \ref{lemma-compatible-shriek-push-finite} we have
$$
j_!f_*f^{-1}\mathcal{G} = \overline{f}_*j'_!f^{-1}\mathcal{G}
$$
By Lemma \ref{lemma-finite-pushforward-statements} it suffices to
consider $j'_!f^{-1}\mathcal{G}$.
The existence of the filtration given by
Lemma \ref{lemma-pullback-filtered},
the fact that $j'_!$ is exact, and
Lemma \ref{lemma-ses-statements}
reduces us to the case
$\mathcal{F} = j'_!\underline{\mathbf{Z}/\ell\mathbf{Z}}$
which is Lemma \ref{lemma-somewhat-easier}.
\end{proof}


%10.20.09

\begin{theorem}
\label{theorem-vanishing-affine-curves}
If $k$ is an algebraically closed field, $X$ is a separated, finite type
scheme of dimension $\leq 1$ over $k$, and $\mathcal{F}$ is a torsion
abelian sheaf on $X_\etale$, then
\begin{enumerate}
\item
$H^q_\etale(X, \mathcal{F}) = 0$ for $q > 2$,
\item
$H^q_\etale(X, \mathcal{F}) = 0$ for $q > 1$ if $X$ is affine,
\item
$H^q_\etale(X, \mathcal{F}) = 0$ for $q > 1$ if $p = \text{char}(k) > 0$
and $\mathcal{F}$ is $p$-power torsion,
\item
$H^q_\etale(X, \mathcal{F})$ is finite if $\mathcal{F}$ is
constructible and torsion prime to $\text{char}(k)$,
\item
$H^q_\etale(X, \mathcal{F})$ is finite if $X$ is proper and
$\mathcal{F}$ constructible,
\item
$H^q_\etale(X, \mathcal{F}) \to
H^q_\etale(X_{k'}, \mathcal{F}|_{X_{k'}})$ is an isomorphism
for any extension $k'/k$ of algebraically closed fields
if $\mathcal{F}$ is torsion prime to $\text{char}(k)$,
\item
$H^q_\etale(X, \mathcal{F}) \to
H^q_\etale(X_{k'}, \mathcal{F}|_{X_{k'}})$ is an isomorphism
for any extension $k'/k$ of algebraically closed fields
if $X$ is proper,
\item
$H^2_\etale(X, \mathcal{F}) \to H^2_\etale(U, \mathcal{F})$
is surjective for all $U \subset X$ open.
\end{enumerate}
\end{theorem}

\begin{proof}
The theorem says that in Situation \ref{situation-what-to-prove}
statements (\ref{item-vanishing}) -- (\ref{item-surjective}) hold.
Our first step is to replace $X$ by its reduction, which is permissible
by Proposition \ref{proposition-topological-invariance}.
By Lemma \ref{lemma-torsion-colimit-constructible} we can write
$\mathcal{F}$ as a filtered colimit of constructible abelian sheaves.
Taking cohomology commutes with colimits, see Lemma \ref{lemma-colimit}.
Moreover, pullback via $X_{k'} \to X$ commutes with colimits as a left
adjoint. Thus it suffices to prove the statements for a constructible sheaf.

\medskip\noindent
In this paragraph we use Lemma \ref{lemma-ses-statements} without further
mention. Writing
$\mathcal{F} = \mathcal{F}_1 \oplus \ldots \oplus \mathcal{F}_r$
where $\mathcal{F}_i$ is $\ell_i$-primary for some prime $\ell_i$, we may
assume that $\ell^n$ kills $\mathcal{F}$ for some prime $\ell$. Now consider
the exact sequence
$$
0 \to \mathcal{F}[\ell] \to \mathcal{F} \to \mathcal{F}/\mathcal{F}[\ell] \to 0.
$$
Thus we see that it suffices to assume that $\mathcal{F}$ is $\ell$-torsion.
This means that $\mathcal{F}$ is a constructible sheaf of
$\mathbf{F}_\ell$-vector spaces for some prime number $\ell$.

\medskip\noindent
By definition this means there is a dense open $U \subset X$
such that $\mathcal{F}|_U$ is finite locally constant sheaf of
$\mathbf{F}_\ell$-vector spaces. Since $\dim(X) \leq 1$ we may
assume, after shrinking $U$, that $U = U_1 \amalg \ldots \amalg U_n$
is a disjoint union of irreducible schemes (just remove the closed
points which lie in the intersections of $\geq 2$ components of $U$).
By Lemma \ref{lemma-restrict-to-open} we reduce to the case
$\mathcal{F} = j_!\mathcal{G}$ where $\mathcal{G}$ is a finite
locally constant sheaf of $\mathbf{F}_\ell$-vector spaces on $U$.

\medskip\noindent
Since we chose $U = U_1 \amalg \ldots \amalg U_n$ with $U_i$ irreducible
we have
$$
j_!\mathcal{G} =
j_{1!}(\mathcal{G}|_{U_1}) \oplus \ldots \oplus
j_{n!}(\mathcal{G}|_{U_n})
$$
where $j_i : U_i \to X$ is the inclusion morphism.
The case of $j_{i!}(\mathcal{G}|_{U_i})$ is handled in
Lemma \ref{lemma-vanishing-easier}.
\end{proof}

\begin{theorem}
\label{theorem-vanishing-curves}
Let $X$ be a finite type, dimension $1$ scheme over an
algebraically closed field $k$. Let $\mathcal{F}$ be a torsion sheaf
on $X_\etale$. Then
$$
H_\etale^q(X, \mathcal{F}) = 0, \quad \forall q \geq 3.
$$
If $X$ affine then also $H_\etale^2(X, \mathcal{F}) = 0$.
\end{theorem}

\begin{proof}
If $X$ is separated, this follows immediately from the more precise
Theorem \ref{theorem-vanishing-affine-curves}.
If $X$ is nonseparated, choose an affine open covering
$X = X_1 \cup \ldots \cup X_n$. By induction on $n$ we may assume
the vanishing holds over $U = X_1 \cup \ldots \cup X_{n - 1}$.
Then Mayer-Vietoris (Lemma \ref{lemma-mayer-vietoris}) gives
$$
H^2_\etale(U, \mathcal{F}) \oplus H^2_\etale(X_n, \mathcal{F}) \to
H^2_\etale(U \cap X_n, \mathcal{F}) \to
H^3_\etale(X, \mathcal{F}) \to 0
$$
However, since $U \cap X_n$ is an open of an affine scheme
and hence affine by our dimension assumption, the group
$H^2_\etale(U \cap X_n, \mathcal{F})$ vanishes
by Theorem \ref{theorem-vanishing-affine-curves}.
\end{proof}

\begin{lemma}
\label{lemma-base-change-dim-1-separably-closed}
Let $k'/k$ be an extension of separably closed fields.
Let $X$ be a proper scheme over $k$ of dimension $\leq 1$.
Let $\mathcal{F}$ be a torsion abelian sheaf on $X$.
Then the map $H^q_\etale(X, \mathcal{F}) \to
H^q_\etale(X_{k'}, \mathcal{F}|_{X_{k'}})$ is an isomorphism
for $q \geq 0$.
\end{lemma}

\begin{proof}
We have seen this for algebraically closed fields in
Theorem \ref{theorem-vanishing-affine-curves}.
Given $k \subset k'$ as in the statement of the lemma we can
choose a diagram
$$
\xymatrix{
k' \ar[r] & \overline{k}' \\
k \ar[u] \ar[r] & \overline{k} \ar[u]
}
$$
where $k \subset \overline{k}$ and $k' \subset \overline{k}'$ are
the algebraic closures. Since $k$ and $k'$ are separably closed
the field extensions
$\overline{k}/k$ and $\overline{k}'/k'$
are algebraic and purely inseparable. In this case the morphisms
$X_{\overline{k}} \to X$ and $X_{\overline{k}'} \to X_{k'}$
are universal homeomorphisms. Thus the cohomology of $\mathcal{F}$
may be computed on $X_{\overline{k}}$ and the cohomology
of $\mathcal{F}|_{X_{k'}}$ may be computed on $X_{\overline{k}'}$,
see Proposition \ref{proposition-topological-invariance}.
Hence we deduce the general case from the case of algebraically
closed fields.
\end{proof}






\section{Cohomology of torsion modules on curves}
\label{section-vanishing-torsion-coefficients}

\noindent
In this section we repeat the arguments of
Section \ref{section-vanishing-torsion}
for constructible sheaves of modules over a Noetherian ring
which are torsion. We start with the most interesting step.

\begin{lemma}
\label{lemma-constant-statements-coefficients}
Let $\Lambda$ be a Noetherian ring, let $M$ be a finite $\Lambda$-module which
is annihilated by an integer $n > 0$, let $k$ be an algebraically closed field,
and let $X$ be a separated, finite type scheme of dimension $\leq 1$ over $k$.
Then
\begin{enumerate}
\item $H^q_\etale(X, \underline{M})$ is a finite $\Lambda$-module
if $n$ is prime to $\text{char}(k)$,
\item $H^q_\etale(X, \underline{M})$ is a finite $\Lambda$-module
if $X$ is proper.
\end{enumerate}
\end{lemma}

\begin{proof}
If $n = \ell n'$ for some prime number $\ell$, then
we get a short exact sequence $0 \to M[\ell] \to M \to M' \to 0$
of finite $\Lambda$-modules and $M'$ is annihilated by $n'$.
This produces a corresponding short exact sequence of constant
sheaves, which in turn gives rise to an exact sequence
of cohomology modules
$$
H^q_\etale(X, \underline{M[n]}) \to
H^q_\etale(X, \underline{M}) \to
H^q_\etale(X, \underline{M'})
$$
Thus, if we can show the result in case $M$ is annihilated by
a prime number, then by induction on $n$ we win.

\medskip\noindent
Let $\ell$ be a prime number such that $\ell$ annihilates $M$.
Then we can replace $\Lambda$ by the $\mathbf{F}_\ell$-algebra
$\Lambda/\ell \Lambda$. Namely, the cohomology of $\mathcal{F}$
as a sheaf of $\Lambda$-modules is the same as the cohomology
of $\mathcal{F}$ as a sheaf of $\Lambda/\ell \Lambda$-modules,
for example by
Cohomology on Sites, Lemma
\ref{sites-cohomology-lemma-cohomology-modules-abelian-agree}.

\medskip\noindent
Assume $\ell$ be a prime number such that $\ell$ annihilates $M$
and $\Lambda$.
Let us reduce to the case where $M$ is a finite free $\Lambda$-module.
Namely, choose a short exact sequence
$$
0 \to N \to \Lambda^{\oplus m} \to M \to 0
$$
This determines an exact sequence
$$
H^q_\etale(X, \underline{\Lambda^{\oplus m}}) \to
H^q_\etale(X, \underline{M}) \to
H^{q + 1}_\etale(X, \underline{N})
$$
By descending induction on $q$ we get the result for $M$
if we know the result for $\Lambda^{\oplus m}$.
Here we use that we know that our cohomology groups
vanish in degrees $> 2$ by
Theorem \ref{theorem-vanishing-affine-curves}.

\medskip\noindent
Let $\ell$ be a prime number and assume that $\ell$ annihilates
$\Lambda$. It remains to show that the cohomology groups
$H^q_\etale(X, \underline{\Lambda})$ are finite $\Lambda$-modules.
We will use a trick to show this; the ``correct'' argument uses
a coefficient theorem which we will show later.
Choose a basis $\Lambda = \bigoplus_{i \in I} \mathbf{F}_\ell e_i$
such that $e_0 = 1$ for some $0 \in I$.
The choice of this basis determines an isomorphism
$$
\underline{\Lambda} = \bigoplus \underline{\mathbf{F}_\ell} e_i
$$
of sheaves on $X_\etale$. Thus we see that
$$
H^q_\etale(X, \underline{\Lambda}) =
H^q_\etale(X, \bigoplus \underline{\mathbf{F}_\ell} e_i) =
\bigoplus H^q_\etale(X, \underline{\mathbf{F}_\ell})e_i
$$
since taking cohomology over $X$ commutes with direct sums
by Theorem \ref{theorem-colimit}
(or Lemma \ref{lemma-colimit} or
Lemma \ref{lemma-direct-sum-bounded-below-cohomology}).
Since we already know that $H^q_\etale(X, \underline{\mathbf{F}_\ell})$
is a finite dimensional $\mathbf{F}_\ell$-vector space
(by Theorem \ref{theorem-vanishing-affine-curves}),
we see that $H^q_\etale(X, \underline{\Lambda})$
is free over $\Lambda$ of the same rank. Namely, given
a basis $\xi_1, \ldots, \xi_m$ of $H^q_\etale(X, \underline{\mathbf{F}_\ell})$
we see that $\xi_1 e_0, \ldots, \xi_m e_0$ form a $\Lambda$-basis for
$H^q_\etale(X, \underline{\Lambda})$.
\end{proof}

\begin{lemma}
\label{lemma-finite-pushforward-coefficients}
Let $\Lambda$ be a Noetherian ring, let $k$ be an algebraically closed field,
let $f : X \to Y$ be a finite morphism of separated finite type schemes
over $k$ of dimension $\leq 1$, and let $\mathcal{F}$ be a sheaf
of $\Lambda$-modules on $X_\etale$. If
$H^q_\etale(X, \mathcal{F})$ is a finite $\Lambda$-module, then
so is $H^q_\etale(Y, f_*\mathcal{F})$.
\end{lemma}

\begin{proof}
Namely, we have $H^q_\etale(X, \mathcal{F}) = H^q_\etale(Y, f_*\mathcal{F})$
by the vanishing of $R^qf_*$ for $q > 0$
(Proposition \ref{proposition-finite-higher-direct-image-zero}) and
the Leray spectral sequence
(Cohomology on Sites, Lemma \ref{sites-cohomology-lemma-apply-Leray}).
\end{proof}

\begin{lemma}
\label{lemma-restrict-to-open-coefficients}
Let $\Lambda$ be a Noetherian ring, let $k$ be an algebraically closed field,
let $X$ be a separated finite type scheme over $k$ of dimension $\leq 1$,
let $\mathcal{F}$ be a constructible sheaf of $\Lambda$-modules on $X_\etale$,
and let $j : X' \to X$ be the inclusion of a dense open subscheme. Then
$H^q_\etale(X, \mathcal{F})$ is a finite $\Lambda$-module if and only if
$H^q_\etale(X, j_!j^{-1}\mathcal{F})$ is a finite $\Lambda$-module.
\end{lemma}

\begin{proof}
Since $X'$ is dense, we see that $Z = X \setminus X'$ has dimension $0$
and hence is a finite set $Z = \{x_1, \ldots, x_n\}$ of $k$-rational points.
Consider the short exact sequence
$$
0 \to j_!j^{-1}\mathcal{F} \to \mathcal{F} \to i_*i^{-1}\mathcal{F} \to 0
$$
of Lemma \ref{lemma-ses-associated-to-open}. Observe that
$H^q_\etale(X, i_*i^{-1}\mathcal{F}) = H^q_\etale(Z, i^*\mathcal{F})$.
Namely, $i : Z \to X$ is a closed immersion, hence finite,
hence we have the vanishing of $R^qi_*$ for $q > 0$ by
Proposition \ref{proposition-finite-higher-direct-image-zero}, and
hence the equality follows from the Leray spectral sequence
(Cohomology on Sites, Lemma \ref{sites-cohomology-lemma-apply-Leray}).
Since $Z$ is a disjoint union of spectra of algebraically closed
fields, we conclude that $H^q_\etale(Z, i^*\mathcal{F}) = 0$
for $q > 0$ and
$$
H^0_\etale(Z, i^{-1}\mathcal{F}) =
\bigoplus\nolimits_{i = 1, \ldots, n} \mathcal{F}_{x_i}
$$
which is a finite $\Lambda$-module $\mathcal{F}_{x_i}$ is finite due to the
assumption that $\mathcal{F}$ is a constructible sheaf of $\Lambda$-modules.
The long exact cohomology sequence gives an exact sequence
$$
0 \to
H^0_\etale(X, j_!j^{-1}\mathcal{F}) \to
H^0_\etale(X, \mathcal{F}) \to
H^0_\etale(Z, i^{-1}\mathcal{F}) \to
H^1_\etale(X, j_!j^{-1}\mathcal{F}) \to
H^1_\etale(X, \mathcal{F}) \to 0
$$
and isomorphisms $H^0_\etale(X, j_!j^{-1}\mathcal{F}) \to
H^0_\etale(X, \mathcal{F})$ for $q > 1$. The lemma follows easily from this.
\end{proof}

\begin{lemma}
\label{lemma-somewhat-easier-coefficients}
Let $\Lambda$ be a Noetherian ring, let $M$ be a finite $\Lambda$-module which
is annihilated by an integer $n > 0$, let $k$ be an algebraically closed field,
let $X$ be a separated, finite type scheme of dimension $\leq 1$ over $k$, and
let $j : U \to X$ be an open immersion. Then
\begin{enumerate}
\item $H^q_\etale(X, j_!\underline{M})$ is a finite $\Lambda$-module
if $n$ is prime to $\text{char}(k)$,
\item $H^q_\etale(X, j_!\underline{M})$ is a finite $\Lambda$-module
if $X$ is proper.
\end{enumerate}
\end{lemma}

\begin{proof}
Since $\dim(X) \leq 1$ there is an open $V \subset X$ which is disjoint from
$U$ such that $X' = U \cup V$ is dense open in $X$ (details omitted).
If $j' : X' \to X$ denotes the inclusion morphism, then we see that
$j_!\underline{M}$ is a direct summand of $j'_!\underline{M}$.
Hence it suffices to prove the lemma in case $U$ is open and dense in $X$.
This case follows from Lemmas \ref{lemma-restrict-to-open-coefficients}
and \ref{lemma-constant-statements-coefficients}.
\end{proof}

\begin{lemma}
\label{lemma-ses-statements-coefficients}
Let $\Lambda$ be a Noetherian ring, let $k$ be an algebraically closed field,
let $X$ be a separated finite type scheme over $k$ of dimension $\leq 1$, and
let $0 \to \mathcal{F}_1 \to \mathcal{F} \to \mathcal{F}_2 \to 0$
be a short exact sequence of sheaves of $\Lambda$-modules on $X_\etale$. If
$H^q_\etale(X, \mathcal{F}_i)$, $i = 1, 2$ are finite $\Lambda$-modules
then $H^q_\etale(X, \mathcal{F})$ is a finite $\Lambda$-module.
\end{lemma}

\begin{proof}
Immediate from the long exact sequence of cohomology.
\end{proof}

\begin{lemma}
\label{lemma-vanishing-easier-coefficients}
Let $\Lambda$ be a Noetherian ring, let $k$ be an algebraically closed field,
let $X$ be a separated, finite type scheme of dimension $\leq 1$ over $k$,
let $j : U \to X$ be an open immersion with $U$ connected, let
$\ell$ be a prime number, let $n > 0$, and let $\mathcal{G}$ be a finite type,
locally constant sheaf of $\Lambda$-modules on $U_\etale$ annihilated by
$\ell^n$. Then
\begin{enumerate}
\item $H^q_\etale(X, j_!\mathcal{G})$ is a finite $\Lambda$-module
if $\ell$ is prime to $\text{char}(k)$,
\item $H^q_\etale(X, j_!\mathcal{G})$ is a finite $\Lambda$-module
if $X$ is proper.
\end{enumerate}
\end{lemma}

\begin{proof}
Let $f : V \to U$ be a finite \'etale morphism of degree prime to $\ell$
as in Lemma \ref{lemma-pullback-filtered-modules}. The discussion in
Section \ref{section-trace-method} gives maps
$$
\mathcal{G} \to f_*f^{-1}\mathcal{G} \to \mathcal{G}
$$
whose composition is an isomorphism. Hence it suffices to prove the
finiteness of $H^q_\etale(X, j_!f_*f^{-1}\mathcal{G})$.
By Zariski's Main theorem
(More on Morphisms, Lemma
\ref{more-morphisms-lemma-quasi-finite-separated-pass-through-finite})
we can choose a diagram
$$
\xymatrix{
V \ar[r]_{j'} \ar[d]_f & Y \ar[d]^{\overline{f}} \\
U \ar[r]^j & X
}
$$
with $\overline{f} : Y \to X$ finite and $j'$ an open immersion with
dense image. Since $f$ is finite and $V$ dense in $Y$ we have
$V = U \times_X Y$. By Lemma \ref{lemma-compatible-shriek-push-finite} we have
$$
j_!f_*f^{-1}\mathcal{G} = \overline{f}_*j'_!f^{-1}\mathcal{G}
$$
By Lemma \ref{lemma-finite-pushforward-coefficients} it suffices to
consider $j'_!f^{-1}\mathcal{G}$. The existence of the filtration given by
Lemma \ref{lemma-pullback-filtered-modules},
the fact that $j'_!$ is exact, and
Lemma \ref{lemma-ses-statements-coefficients}
reduces us to the case
$\mathcal{F} = j'_!\underline{M}$ for a finite $\Lambda$-module $M$
which is Lemma \ref{lemma-somewhat-easier-coefficients}.
\end{proof}

\begin{theorem}
\label{theorem-vanishing-affine-curves-coefficients}
Let $\Lambda$ be a Noetherian ring, let $k$ be an algebraically closed field,
let $X$ be a separated, finite type scheme of dimension $\leq 1$ over $k$,
and let $\mathcal{F}$ be a constructible sheaf of $\Lambda$-modules
on $X_\etale$ which is torsion. Then
\begin{enumerate}
\item
\label{item-finite-prime-to-p-coefficients}
$H^q_\etale(X, \mathcal{F})$ is a finite $\Lambda$-module
if $\mathcal{F}$ is torsion prime to $\text{char}(k)$,
\item
\label{item-finite-proper-coefficients}
$H^q_\etale(X, \mathcal{F})$ is a finite $\Lambda$-module if $X$ is proper.
\end{enumerate}
\end{theorem}

\begin{proof}
without further mention.
Write $\mathcal{F} = \mathcal{F}_1 \oplus \ldots \oplus \mathcal{F}_r$
where $\mathcal{F}_i$ is annihilated by $\ell_i^{n_i}$
for some prime $\ell_i$ and integer $n_i > 0$.
By Lemma \ref{lemma-ses-statements-coefficients} it suffices
to prove the theorem for $\mathcal{F}_i$. Thus we may and do
assume that $\ell^n$ kills $\mathcal{F}$ for some prime $\ell$
and integer $n > 0$.

\medskip\noindent
Since $\mathcal{F}$ is constructible as a sheaf of $\Lambda$-modules,
there is a dense open $U \subset X$ such that $\mathcal{F}|_U$ is a
finite type, locally constant sheaf of $\Lambda$-modules.
Since $\dim(X) \leq 1$ we may assume, after shrinking $U$, that
$U = U_1 \amalg \ldots \amalg U_n$ is a disjoint union of irreducible
schemes (just remove the closed points which lie in the intersections
of $\geq 2$ components of $U$). By
Lemma \ref{lemma-restrict-to-open-coefficients}
we reduce to the case
$\mathcal{F} = j_!\mathcal{G}$ where $\mathcal{G}$ is a finite type,
locally constant sheaf of $\Lambda$-modules on $U$ (and annihilated
by $\ell^n$).

\medskip\noindent
Since we chose $U = U_1 \amalg \ldots \amalg U_n$ with $U_i$ irreducible
we have
$$
j_!\mathcal{G} =
j_{1!}(\mathcal{G}|_{U_1}) \oplus \ldots \oplus
j_{n!}(\mathcal{G}|_{U_n})
$$
where $j_i : U_i \to X$ is the inclusion morphism.
The case of $j_{i!}(\mathcal{G}|_{U_i})$ is handled in
Lemma \ref{lemma-vanishing-easier-coefficients}.
\end{proof}








\section{First cohomology of proper schemes}
\label{section-finite-etale-over-proper}

\noindent
In Fundamental Groups, Section \ref{pione-section-finite-etale-over-proper}
we have seen, in some sense, that taking
$R^1f_*\underline{G}$ commutes with base change if $f : X \to Y$
is a proper morphism and $G$ is a finite group (not necessarily
commutative). In this section
we deduce a useful consequence of these results.

\begin{lemma}
\label{lemma-proper-over-henselian-and-h1}
Let $A$ be a henselian local ring. Let $X$ be a proper scheme over $A$
with closed fibre $X_0$. Let $M$ be a finite abelian group.
Then $H^1_\etale(X, \underline{M}) = H^1_\etale(X_0, \underline{M})$.
\end{lemma}

\begin{proof}
By Cohomology on Sites, Lemma \ref{sites-cohomology-lemma-torsors-h1}
an element of $H^1_\etale(X, \underline{M})$ corresponds to a
$\underline{M}$-torsor $\mathcal{F}$ on $X_\etale$.
Such a torsor is clearly a finite locally constant sheaf.
Hence $\mathcal{F}$ is representable by a scheme $V$ finite
\'etale over $X$, Lemma \ref{lemma-characterize-finite-locally-constant}.
Conversely, a scheme $V$ finite \'etale over $X$ with an $M$-action
which turns it into an $M$-torsor over $X$ gives rise to a cohomology
class. The same translation between cohomology classes over $X_0$ and
torsors finite \'etale over $X_0$ holds. Thus the lemma
is a consequence of the equivalence of categories of
Fundamental Groups, Lemma
\ref{pione-lemma-finite-etale-on-proper-over-henselian}.
\end{proof}

\noindent
The following technical lemma is a key ingredient in the proof of
the proper base change theorem. The argument works word for word
for any proper scheme over $A$ whose special fibre has dimension
$\leq 1$, but in fact the conclusion will be a consequence of the
proper base change theorem and we only need this particular version
in its proof.

\begin{lemma}
\label{lemma-efface-cohomology-on-fibre-by-finite-cover}
Let $A$ be a henselian local ring. Let $X = \mathbf{P}^1_A$.
Let $X_0 \subset X$ be the closed fibre. Let $\ell$ be a prime
number. Let $\mathcal{I}$ be an injective sheaf of
$\mathbf{Z}/\ell\mathbf{Z}$-modules on $X_\etale$. Then
$H^q_\etale(X_0, \mathcal{I}|_{X_0}) = 0$ for $q > 0$.
\end{lemma}

\begin{proof}
Observe that $X$ is a separated scheme which can be covered by $2$
affine opens. Hence for $q > 1$ this follows from Gabber's affine
variant of the proper base change theorem, see
Lemma \ref{lemma-vanishing-restriction-injective}.
Thus we may assume $q = 1$. Let
$\xi \in H^1_\etale(X_0, \mathcal{I}|_{X_0})$.
Goal: show that $\xi$ is $0$.
By Lemmas \ref{lemma-torsion-colimit-constructible} and
\ref{lemma-colimit} we can find a map $\mathcal{F} \to \mathcal{I}$
with $\mathcal{F}$ a constructible sheaf of
$\mathbf{Z}/\ell\mathbf{Z}$-modules
and $\xi$ coming from an element $\zeta$ of
$H^1_\etale(X_0, \mathcal{F}|_{X_0})$. Suppose we have an injective map
$\mathcal{F} \to \mathcal{F}'$ of sheaves of
$\mathbf{Z}/\ell\mathbf{Z}$-modules on $X_\etale$.
Since $\mathcal{I}$ is injective we can extend the given map
$\mathcal{F} \to \mathcal{I}$ to a map $\mathcal{F}' \to \mathcal{I}$.
In this situation we may replace $\mathcal{F}$ by $\mathcal{F}'$
and $\zeta$ by the image of $\zeta$ in $H^1_\etale(X_0, \mathcal{F}'|_{X_0})$.
Also, if $\mathcal{F} = \mathcal{F}_1 \oplus \mathcal{F}_2$ is a direct sum,
then we may replace $\mathcal{F}$ by $\mathcal{F}_i$
and $\zeta$ by the image of $\zeta$ in $H^1_\etale(X_0, \mathcal{F}_i|_{X_0})$.

\medskip\noindent
By Lemma \ref{lemma-constructible-maps-into-constant-general}
and the remarks above we may assume $\mathcal{F}$
is of the form $f_*\underline{M}$ where $M$ is a finite
$\mathbf{Z}/\ell\mathbf{Z}$-module
and $f : Y \to X$ is a finite morphism of finite presentation
(such sheaves are still constructible by
Lemma \ref{lemma-finite-pushforward-constructible}
but we won't need this).
Since formation of $f_*$ commutes with any base change
(Lemma \ref{lemma-finite-pushforward-commutes-with-base-change})
we see that the restriction of $f_*\underline{M}$ to $X_0$ is
equal to the pushforward of $\underline{M}$ via the induced morphism
$Y_0 \to X_0$ of special fibres. By the Leray spectral sequence
(Proposition \ref{proposition-leray})
and vanishing of higher direct images
(Proposition \ref{proposition-finite-higher-direct-image-zero}),
we find
$$
H^1_\etale(X_0, f_*\underline{M}|_{X_0}) = H^1_\etale(Y_0, \underline{M}).
$$
Since $Y \to \Spec(A)$ is proper we can use
Lemma \ref{lemma-proper-over-henselian-and-h1} to see that
the $H^1_\etale(Y_0, \underline{M})$ is equal to
$H^1_\etale(Y, \underline{M})$. Thus we see that our cohomology
class $\zeta$ lifts to a cohomology class
$$
\tilde\zeta \in H^1_\etale(Y, \underline{M}) = H^1_\etale(X, f_*\underline{M})
$$
However, $\tilde \zeta$ maps to zero in
$H^1_\etale(X, \mathcal{I})$ as $\mathcal{I}$ is injective
and by commutativity of
$$
\xymatrix{
H^1_\etale(X, f_*\underline{M}) \ar[r] \ar[d] &
H^1_\etale(X, \mathcal{I}) \ar[d] \\
H^1_\etale(X_0, (f_*\underline{M})|_{X_0}) \ar[r] &
H^1_\etale(X_0, \mathcal{I}|_{X_0})
}
$$
we conclude that the image $\xi$ of $\zeta$ is zero as well.
\end{proof}











\section{Preliminaries on base change}
\label{section-base-change-preliminaries}

\noindent
If you are interested in either the smooth base change theorem
or the proper base change theorem, you should skip directly to
the corresponding sections.
In this section and the next few sections we consider commutative diagrams
$$
\xymatrix{
X \ar[d]_f & Y \ar[l]^h \ar[d]^e \\
S & T \ar[l]_g
}
$$
of schemes; we usually assume this diagram is cartesian,
i.e., $Y = X \times_S T$. A commutative diagram as above
gives rise to a commutative diagram
$$
\xymatrix{
X_\etale \ar[d]_{f_{small}} & Y_\etale \ar[d]^{e_{small}} \ar[l]^{h_{small}} \\
S_\etale & T_\etale \ar[l]_{g_{small}}
}
$$
of small \'etale sites. Let us use the notation
$$
f^{-1} = f_{small}^{-1}, \quad
g_* = g_{small, *}, \quad
e^{-1} = e_{small}^{-1}, \text{ and}\quad
h_* = h_{small, *}.
$$
By Sites, Section \ref{sites-section-pullback}
we get a base change or pullback map
$$
f^{-1}g_*\mathcal{F}
\longrightarrow
h_*e^{-1}\mathcal{F}
$$
for a sheaf $\mathcal{F}$ on $T_\etale$. If $\mathcal{F}$ is an abelian
sheaf on $T_\etale$, then we get a derived base change map
$$
f^{-1}Rg_*\mathcal{F}
\longrightarrow
Rh_*e^{-1}\mathcal{F}
$$
see Cohomology on Sites, Lemma
\ref{sites-cohomology-lemma-base-change-map-flat-case}.
Finally, if $K$ is an arbitrary object of $D(T_\etale)$
there is a base change map
$$
f^{-1}Rg_*K
\longrightarrow
Rh_*e^{-1}K
$$
see
Cohomology on Sites, Remark \ref{sites-cohomology-remark-base-change}.

\begin{lemma}
\label{lemma-base-change-local}
Consider a cartesian diagram of schemes
$$
\xymatrix{
X \ar[d]_f & Y \ar[l]^h \ar[d]^e \\
S & T \ar[l]_g
}
$$
Let $\{U_i \to X\}$ be an \'etale covering such that $U_i \to S$
factors as $U_i \to V_i \to S$ with $V_i \to S$ \'etale
and consider the cartesian diagrams
$$
\xymatrix{
U_i \ar[d]_{f_i} & U_i \times_X Y \ar[l]^{h_i} \ar[d]^{e_i} \\
V_i & V_i \times_S T \ar[l]_{g_i}
}
$$
Let $\mathcal{F}$ be a sheaf on $T_\etale$. Let $K$ in $D(T_\etale)$.
Set $K_i = K|_{V_i \times_S T}$ and
$\mathcal{F}_i = \mathcal{F}|_{V_i \times_S T}$.
\begin{enumerate}
\item If $f_i^{-1}g_{i, *}\mathcal{F}_i = h_{i, *}e_i^{-1}\mathcal{F}_i$
for all $i$, then $f^{-1}g_*\mathcal{F} = h_*e^{-1}\mathcal{F}$.
\item If $f_i^{-1}Rg_{i, *}K_i = Rh_{i, *}e_i^{-1}K_i$
for all $i$, then $f^{-1}Rg_*K = Rh_*e^{-1}K$.
\item If $\mathcal{F}$ is an abelian sheaf and
$f_i^{-1}R^qg_{i, *}\mathcal{F}_i = R^qh_{i, *}e_i^{-1}\mathcal{F}_i$
for all $i$, then
$f^{-1}R^qg_*\mathcal{F} = R^qh_*e^{-1}\mathcal{F}$.
\end{enumerate}
\end{lemma}

\begin{proof}
Proof of (1). First we observe that
$$
(f^{-1}g_*\mathcal{F})|_{U_i} =
f_i^{-1}(g_*\mathcal{F}|_{V_i}) =
f_i^{-1}g_{i, *}\mathcal{F}_i
$$
The first equality because $U_i \to X \to S$ is equal to
$U_i \to V_i \to S$ and the second equality because
$g_*\mathcal{F}|_{V_i} = g_{i, *}\mathcal{F}_i$ by
Sites, Lemma \ref{sites-lemma-localize-morphism-strong}.
Similarly we have
$$
(h_*e^{-1}\mathcal{F})|_{U_i} =
h_{i, *}(e^{-1}\mathcal{F}|_{U_i \times_X Y}) =
h_{i, *}e_i^{-1}\mathcal{F}_i
$$
Thus if the base change maps
$f_i^{-1}g_{i, *}\mathcal{F}_i \to h_{i, *}e_i^{-1}\mathcal{F}_i$
are isomorphisms for all $i$, then the base change map
$f^{-1}g_*\mathcal{F} \to h_*e^{-1}\mathcal{F}$
restricts to an isomorphism over $U_i$ for all $i$
and we conclude it is an isomorphism as $\{U_i \to X\}$
is an \'etale covering.

\medskip\noindent
For the other two statements we replace the appeal to
Sites, Lemma \ref{sites-lemma-localize-morphism-strong}
by an appeal to
Cohomology on Sites, Lemma
\ref{sites-cohomology-lemma-restrict-direct-image-open}.
\end{proof}

\begin{lemma}
\label{lemma-base-change-compose}
Consider a tower of cartesian diagrams of schemes
$$
\xymatrix{
W \ar[d]_i & Z \ar[l]^j \ar[d]^k \\
X \ar[d]_f & Y \ar[l]^h \ar[d]^e \\
S & T \ar[l]_g
}
$$
Let $K$ in $D(T_\etale)$. If
$$
f^{-1}Rg_*K \to Rh_*e^{-1}K
\quad\text{and}\quad
i^{-1}Rh_*e^{-1}K \to Rj_*k^{-1}e^{-1}K
$$
are isomorphisms, then
$(f \circ i)^{-1}Rg_*K \to Rj_*(e \circ k)^{-1}K$
is an isomorphism.
Similarly, if $\mathcal{F}$ is an abelian sheaf on $T_\etale$ and if
$$
f^{-1}R^qg_*\mathcal{F} \to R^qh_*e^{-1}\mathcal{F}
\quad\text{and}\quad
i^{-1}R^qh_*e^{-1}\mathcal{F} \to R^qj_*k^{-1}e^{-1}\mathcal{F}
$$
are isomorphisms, then
$(f \circ i)^{-1}R^qg_*\mathcal{F} \to R^qj_*(e \circ k)^{-1}\mathcal{F}$
is an isomorphism.
\end{lemma}

\begin{proof}
This is formal, provided one checks that the composition of these
base change maps is the base change maps for the outer rectangle, see
Cohomology on Sites, Remark
\ref{sites-cohomology-remark-compose-base-change-horizontal}.
\end{proof}

\begin{lemma}
\label{lemma-base-change-Rf-star-colim}
Let $I$ be a directed set. Consider an inverse system of
cartesian diagrams of schemes
$$
\xymatrix{
X_i \ar[d]_{f_i} & Y_i \ar[l]^{h_i} \ar[d]^{e_i} \\
S_i & T_i \ar[l]_{g_i}
}
$$
with affine transition morphisms and with $g_i$ quasi-compact and
quasi-separated. Set $X = \lim X_i$,
$S = \lim S_i$, $T = \lim T_i$ and $Y = \lim Y_i$ to
obtain the cartesian diagram
$$
\xymatrix{
X \ar[d]_f & Y \ar[l]^h \ar[d]^e \\
S & T \ar[l]_g
}
$$
Let $(\mathcal{F}_i, \varphi_{i'i})$ be a system of sheaves on
$(T_i)$ as in Definition \ref{definition-inverse-system-sheaves}. Set
$\mathcal{F} = \colim p_i^{-1}\mathcal{F}_i$ on $T$
where $p_i : T \to T_i$ is the projection.
Then we have the following
\begin{enumerate}
\item If $f_i^{-1}g_{i, *}\mathcal{F}_i = h_{i, *}e_i^{-1}\mathcal{F}_i$
for all $i$, then
$f^{-1}g_*\mathcal{F} = h_*e^{-1}\mathcal{F}$.
\item If $\mathcal{F}_i$ is an abelian sheaf for all $i$ and
$f_i^{-1}R^qg_{i, *}\mathcal{F}_i = R^qh_{i, *}e_i^{-1}\mathcal{F}_i$
for all $i$, then
$f^{-1}R^qg_*\mathcal{F} = R^qh_*e^{-1}\mathcal{F}$.
\end{enumerate}
\end{lemma}

\begin{proof}
We prove (2) and we omit the proof of (1). We will use without further
mention that pullback of sheaves commutes with colimits as it is a
left adjoint. Observe that $h_i$ is quasi-compact and quasi-separated as a
base change of $g_i$.
Denoting $q_i : Y \to Y_i$ the projections, observe that
$e^{-1}\mathcal{F} = \colim e^{-1}p_i^{-1}\mathcal{F}_i =
\colim q_i^{-1}e_i^{-1}\mathcal{F}_i$.
By Lemma \ref{lemma-relative-colimit-general}
this gives
$$
R^qh_*e^{-1}\mathcal{F} = \colim r_i^{-1}R^qh_{i, *}e_i^{-1}\mathcal{F}_i
$$
where $r_i : X \to X_i$ is the projection.
Similarly, we have
$$
f^{-1}Rg_*\mathcal{F} =
f^{-1}\colim s_i^{-1}R^qg_{i, *}\mathcal{F}_i =
\colim r_i^{-1}f_i^{-1}R^qg_{i, *}\mathcal{F}_i
$$
where $s_i : S \to S_i$ is the projection. The lemma follows.
\end{proof}

\begin{lemma}
\label{lemma-base-change-Rf-star-colim-complexes}
Let $I$, $X_i$, $Y_i$, $S_i$, $T_i$, $f_i$, $h_i$, $e_i$, $g_i$,
$X$, $Y$, $S$, $T$, $f$, $h$, $e$, $g$ be as in the statement
of Lemma \ref{lemma-base-change-Rf-star-colim}.
Let $0 \in I$ and let $K_0 \in D^+(T_{0, \etale})$.
For $i \in I$, $i \geq 0$ denote $K_i$ the pullback of
$K_0$ to $T_i$. Denote $K$ the pullback of $K_0$ to $T$.
If $f_i^{-1}Rg_{i, *}K_i = Rh_{i, *}e_i^{-1}K_i$
for all $i \geq 0$, then $f^{-1}Rg_*K = Rh_*e^{-1}K$.
\end{lemma}

\begin{proof}
It suffices to show that the base change map
$f^{-1}Rg_*K \to Rh_*e^{-1}K$ induces an isomorphism
on cohomology sheaves. In other words, we have to show
that $f^{-1}R^pg_*K \to R^ph_*e^{-1}K$ is an isomorphism
for all $p \in \mathbf{Z}$ if we are given that
$f_i^{-1}R^pg_{i, *}K_i \to R^ph_{i, *}e_i^{-1}K_i$
is an isomorphism for all $i \geq 0$ and $p \in \mathbf{Z}$.
At this point we can argue exactly as in the proof of
Lemma \ref{lemma-base-change-Rf-star-colim}
replacing reference to
Lemma \ref{lemma-relative-colimit-general}
by a reference to
Lemma \ref{lemma-relative-colimit-general-complexes}.
\end{proof}

\begin{lemma}
\label{lemma-base-change-f-star-general-stalks}
Consider a cartesian diagram of schemes
$$
\xymatrix{
X \ar[d]_f & Y \ar[l]^h \ar[d]^e \\
S & T \ar[l]_g
}
$$
where $g : T \to S$ is quasi-compact and quasi-separated.
Let $\mathcal{F}$ be an
abelian sheaf on $T_\etale$. Let $q \geq 0$. The following are equivalent
\begin{enumerate}
\item For every geometric point $\overline{x}$ of $X$ with image
$\overline{s} = f(\overline{x})$ we have
$$
H^q(\Spec(\mathcal{O}^{sh}_{X, \overline{x}}) \times_S T, \mathcal{F})
=
H^q(\Spec(\mathcal{O}^{sh}_{S, \overline{s}}) \times_S T, \mathcal{F})
$$
\item $f^{-1}R^qg_*\mathcal{F} \to R^qh_*e^{-1}\mathcal{F}$
is an isomorphism.
\end{enumerate}
\end{lemma}

\begin{proof}
Since $Y = X \times_S T$ we have
$\Spec(\mathcal{O}^{sh}_{X, \overline{x}}) \times_X Y =
\Spec(\mathcal{O}^{sh}_{X, \overline{x}}) \times_S T$. Thus
the map in (1) is the map of stalks at $\overline{x}$ for the map
in (2) by Theorem \ref{theorem-higher-direct-images} (and
Lemma \ref{lemma-stalk-pullback}).
Thus the result by Theorem \ref{theorem-exactness-stalks}.
\end{proof}

\begin{lemma}
\label{lemma-check-stalks-better}
Let $f : X \to S$ be a morphism of schemes.
Let $\overline{x}$ be a geometric point of $X$ with image $\overline{s}$ in $S$.
Let $\Spec(K) \to \Spec(\mathcal{O}^{sh}_{S, \overline{s}})$
be a morphism with $K$ a separably closed field. Let $\mathcal{F}$ be an
abelian sheaf on $\Spec(K)_\etale$. Let $q \geq 0$. The following are
equivalent
\begin{enumerate}
\item
$H^q(\Spec(\mathcal{O}^{sh}_{X, \overline{x}}) \times_S \Spec(K), \mathcal{F}) =
H^q(\Spec(\mathcal{O}^{sh}_{S, \overline{s}}) \times_S \Spec(K), \mathcal{F})$
\item
$H^q(\Spec(\mathcal{O}^{sh}_{X, \overline{x}})
\times_{\Spec(\mathcal{O}^{sh}_{S, \overline{s}})} \Spec(K), \mathcal{F}) =
H^q(\Spec(K), \mathcal{F})$
\end{enumerate}
\end{lemma}

\begin{proof}
Observe that $\Spec(K) \times_S \Spec(\mathcal{O}^{sh}_{S, \overline{s}})$
is the spectrum of a filtered colimit of \'etale algebras over $K$.
Since $K$ is separably closed, each \'etale $K$-algebra
is a finite product of copies of $K$. Thus we can write
$$
\Spec(K) \times_S \Spec(\mathcal{O}^{sh}_{S, \overline{s}}) =
\lim_{i \in I} \coprod\nolimits_{a \in A_i} \Spec(K)
$$
as a cofiltered limit where each term is a disjoint union
of copies of $\Spec(K)$ over a finite set $A_i$.
Note that $A_i$ is nonempty as we are given
$\Spec(K) \to \Spec(\mathcal{O}^{sh}_{S, \overline{s}})$.
It follows that
\begin{align*}
\Spec(\mathcal{O}^{sh}_{X, \overline{x}}) \times_S \Spec(K)
& =
\Spec(\mathcal{O}^{sh}_{X, \overline{x}})
\times_{\Spec(\mathcal{O}^{sh}_{S, \overline{s}})}
\left(
\Spec(\mathcal{O}^{sh}_{S, \overline{s}})
\times_S \Spec(K)\right) \\
& =
\lim_{i \in I} \coprod\nolimits_{a \in A_i}
\Spec(\mathcal{O}^{sh}_{X, \overline{x}})
\times_{\Spec(\mathcal{O}^{sh}_{S, \overline{s}})} \Spec(K)
\end{align*}
Since taking cohomology in our setting commutes with limits
of schemes (Theorem \ref{theorem-colimit}) we conclude.
\end{proof}













\section{Base change for pushforward}
\label{section-base-change-f-star}

\noindent
This section is preliminary and should be skipped on a first reading.
In this section we discuss for what morphisms $f : X \to S$ we have
$f^{-1}g_* = h_*e^{-1}$ on all sheaves (of sets) for every cartesian
diagram
$$
\xymatrix{
X \ar[d]_f & Y \ar[l]^h \ar[d]^e \\
S & T \ar[l]_g
}
$$
with $g$ quasi-compact and quasi-separated.

\begin{lemma}
\label{lemma-base-change-f-star-general}
Consider the cartesian diagram of schemes
$$
\xymatrix{
X \ar[d]_f & Y \ar[l]^h \ar[d]^e \\
S & T \ar[l]_g
}
$$
Assume that $f$ is flat and every object $U$ of $X_\etale$ has
a covering $\{U_i \to U\}$ such that $U_i \to S$
factors as $U_i \to V_i \to S$ with $V_i \to S$
\'etale and $U_i \to V_i$ quasi-compact with
geometrically connected fibres.
Then for any sheaf $\mathcal{F}$ of sets on $T_\etale$ we have
$f^{-1}g_*\mathcal{F} = h_*e^{-1}\mathcal{F}$.
\end{lemma}

\begin{proof}
Let $U \to X$ be an \'etale morphism such that $U \to S$ factors as
$U \to V \to S$ with $V \to S$ \'etale and $U \to V$ quasi-compact
with geometrically connected fibres. Observe that $U \to V$ is flat
(More on Flatness, Lemma \ref{flat-lemma-etale-flat-up-down}).
We claim that
\begin{align*}
f^{-1}g_*\mathcal{F}(U)
& = g_*\mathcal{F}(V) \\
& = \mathcal{F}(V \times_S T) \\
& = e^{-1}\mathcal{F}(U \times_X Y) \\
& = h_*e^{-1}\mathcal{F}(U)
\end{align*}
Namely, thinking of $U$ as an object of $X_\etale$ and
$V$ as an object of $S_\etale$ we see that the first equality
follows from Lemma \ref{lemma-sections-upstairs}\footnote{Strictly
speaking, we are also using that the restriction of $f^{-1}g_*\mathcal{F}$
to $U_\etale$ is the pullback via $U \to V$ of the restriction of
$g_*\mathcal{F}$ to $V_\etale$. See
Sites, Lemma \ref{sites-lemma-localize-morphism-strong}.}.
Thinking of $V \times_S T$ as an object of $T_\etale$
the second equality follows from the definition of $g_*$.
Observe that $U \times_X Y = U \times_S T$ (because $Y = X \times_S T$)
and hence $U \times_X Y \to V \times_S T$
has geometrically connected fibres as a base change of $U \to V$.
Thinking of $U \times_X Y$ as an object of $Y_\etale$, we see that
the third equality  follows from Lemma \ref{lemma-sections-upstairs}
as before. Finally, the fourth equality follows from the definition
of $h_*$.

\medskip\noindent
Since by assumption every object of $X_\etale$ has an \'etale
covering to which the argument of the previous paragraph applies
we see that the lemma is true.
\end{proof}

\begin{lemma}
\label{lemma-fppf-reduced-fibres-base-change-f-star}
Consider a cartesian diagram of schemes
$$
\xymatrix{
X \ar[d]_f & Y \ar[l]^h \ar[d]^e \\
S & T \ar[l]_g
}
$$
where $f$ is flat and locally of finite presentation
with geometrically reduced fibres.
Then $f^{-1}g_*\mathcal{F} = h_*e^{-1}\mathcal{F}$
for any sheaf $\mathcal{F}$ on $T_\etale$.
\end{lemma}

\begin{proof}
Combine Lemma \ref{lemma-base-change-f-star-general} with
More on Morphisms, Lemma
\ref{more-morphisms-lemma-cover-by-geometrically-connected}.
\end{proof}

\begin{lemma}
\label{lemma-base-change-f-star-field}
Consider the cartesian diagrams of schemes
$$
\xymatrix{
X \ar[d]_f & Y \ar[l]^h \ar[d]^e \\
S & T \ar[l]_g
}
$$
Assume that $S$ is the spectrum of a separably closed field.
Then $f^{-1}g_*\mathcal{F} = h_*e^{-1}\mathcal{F}$
for any sheaf $\mathcal{F}$ on $T_\etale$.
\end{lemma}

\begin{proof}
We may work locally on $X$. Hence we may assume $X$ is affine.
Then we can write $X$ as a cofiltered limit of affine schemes of
finite type over $S$. By Lemma \ref{lemma-base-change-Rf-star-colim}
we may assume that $X$ is of finite type over $S$.
Then Lemma \ref{lemma-base-change-f-star-general}
applies because any scheme of finite
type over a separably closed field is a finite disjoint
union of connected and geometrically connected schemes
(see Varieties, Lemma
\ref{varieties-lemma-separably-closed-field-connected-components}).
\end{proof}

\begin{lemma}
\label{lemma-base-change-f-star-valuation}
Consider a cartesian diagram of schemes
$$
\xymatrix{
X \ar[d]_f & Y \ar[l]^h \ar[d]^e \\
S & T \ar[l]_g
}
$$
Assume that
\begin{enumerate}
\item $f$ is flat and open,
\item the residue fields of $S$ are separably algebraically closed,
\item given an \'etale morphism $U \to X$ with $U$ affine
we can write $U$ as a finite disjoint union of open subschemes
of $X$ (for example if $X$ is a normal integral scheme
with separably closed function field),
\item any nonempty open of a fibre $X_s$ of $f$ is connected
(for example if $X_s$ is irreducible or empty).
\end{enumerate}
Then for any sheaf $\mathcal{F}$ of sets on $T_\etale$ we have
$f^{-1}g_*\mathcal{F} = h_*e^{-1}\mathcal{F}$.
\end{lemma}

\begin{proof}
Omitted. Hint: the assumptions almost trivially imply
the condition of Lemma \ref{lemma-base-change-f-star-general}.
The for example in part (3) follows from
Lemma \ref{lemma-normal-scheme-with-alg-closed-function-field}.
\end{proof}

\noindent
The following lemma doesn't really belong here but there does not
seem to be a good place for it anywhere.

\begin{lemma}
\label{lemma-fppf-reduced-fibres-pullback-products}
Let $f : X \to S$ be a morphism of schemes which is flat and
locally of finite presentation with geometrically reduced fibres.
Then $f^{-1} : \Sh(S_\etale) \to \Sh(X_\etale)$ commutes
with products.
\end{lemma}

\begin{proof}
Let $I$ be a set and let $\mathcal{G}_i$ be a sheaf on $S_\etale$
for $i \in I$.
Let $U \to X$ be an \'etale morphism such that $U \to S$ factors as
$U \to V \to S$ with $V \to S$ \'etale and $U \to V$ flat of finite
presentation with geometrically connected fibres. Then we have
\begin{align*}
f^{-1}(\prod \mathcal{G}_i)(U)
& =
(\prod \mathcal{G}_i)(V) \\
& =
\prod \mathcal{G}_i(V) \\
& =
\prod f^{-1}\mathcal{G}_i(U) \\
& =
(\prod f^{-1}\mathcal{G}_i)(U)
\end{align*}
where we have used Lemma \ref{lemma-sections-upstairs}
in the first and third equality
(we are also using that the restriction of $f^{-1}\mathcal{G}$
to $U_\etale$ is the pullback via $U \to V$ of the restriction of
$\mathcal{G}$ to $V_\etale$, see
Sites, Lemma \ref{sites-lemma-localize-morphism-strong}).
By More on Morphisms, Lemma
\ref{more-morphisms-lemma-cover-by-geometrically-connected}
every object $U$ of $X_\etale$ has an \'etale covering
$\{U_i \to U\}$ such that the discussion in the previous
paragraph applies to $U_i$. The lemma follows.
\end{proof}

\begin{lemma}
\label{lemma-base-change-f-star}
Let $f : X \to S$ be a flat morphism of schemes such
that for every geometric point $\overline{x}$ of $X$ the map
$$
\mathcal{O}_{S, f(\overline{x})}^{sh}
\longrightarrow
\mathcal{O}_{X, \overline{x}}^{sh}
$$
has geometrically connected fibres. Then for every
cartesian diagram of schemes
$$
\xymatrix{
X \ar[d]_f & Y \ar[l]^h \ar[d]^e \\
S & T \ar[l]_g
}
$$
with $g$ quasi-compact and quasi-separated we have
$f^{-1}g_*\mathcal{F} = h_*e^{-1}\mathcal{F}$
for any sheaf $\mathcal{F}$ of sets on $T_\etale$.
\end{lemma}

\begin{proof}
It suffices to check equality on stalks, see
Theorem \ref{theorem-exactness-stalks}.
By Theorem \ref{theorem-higher-direct-images} we have
$$
(h_*e^{-1}\mathcal{F})_{\overline{x}} =
\Gamma(\Spec(\mathcal{O}_{X, \overline{x}}^{sh}) \times_X Y, e^{-1}\mathcal{F})
$$
and we have similarly
$$
(f^{-1}g_*^{-1}\mathcal{F})_{\overline{x}} =
(g_*^{-1}\mathcal{F})_{f(\overline{x})} =
\Gamma(\Spec(\mathcal{O}_{S, f(\overline{x})}^{sh}) \times_S T, \mathcal{F})
$$
These sets are equal by an application of Lemma \ref{lemma-sections-upstairs}
to the morphism
$$
\Spec(\mathcal{O}_{X, \overline{x}}^{sh}) \times_X Y
\longrightarrow
\Spec(\mathcal{O}_{S, f(\overline{x})}^{sh}) \times_S T
$$
which is a base change of
$\Spec(\mathcal{O}_{X, \overline{x}}^{sh}) \to
\Spec(\mathcal{O}_{S, f(\overline{x})}^{sh})$
because $Y = X \times_S T$.
\end{proof}







\section{Base change for higher direct images}
\label{section-base-change}

\noindent
This section is the analogue of Section \ref{section-base-change-f-star}
for higher direct images.
This section is preliminary and should be skipped on a first reading.

\begin{remark}
\label{remark-base-change-holds}
Let $f : X \to S$ be a morphism of schemes. Let $n$ be an integer.
We will say $BC(f, n, q_0)$ is true if for every commutative diagram
$$
\xymatrix{
X \ar[d]_f & X' \ar[l] \ar[d]_{f'} & Y \ar[l]^h \ar[d]^e \\
S & S' \ar[l] & T \ar[l]_g
}
$$
with $X' = X \times_S S'$ and $Y = X' \times_{S'} T$ and
$g$ quasi-compact and quasi-separated, and every abelian sheaf
$\mathcal{F}$ on $T_\etale$ annihilated by $n$ the base change map
$$
(f')^{-1}R^qg_*\mathcal{F}
\longrightarrow
R^qh_*e^{-1}\mathcal{F}
$$
is an isomorphism for $q \leq q_0$.
\end{remark}

\begin{lemma}
\label{lemma-base-change-q-injective}
With $f : X \to S$ and $n$ as in Remark \ref{remark-base-change-holds}
assume for some $q \geq 1$ we have $BC(f, n, q - 1)$. Then
for every commutative diagram
$$
\xymatrix{
X \ar[d]_f & X' \ar[l] \ar[d]_{f'} & Y \ar[l]^h \ar[d]^e \\
S & S' \ar[l] & T \ar[l]_g
}
$$
with $X' = X \times_S S'$ and $Y = X' \times_{S'} T$ and
$g$ quasi-compact and quasi-separated, and every abelian sheaf
$\mathcal{F}$ on $T_\etale$ annihilated by $n$
\begin{enumerate}
\item the base change map
$(f')^{-1}R^qg_*\mathcal{F}\to R^qh_*e^{-1}\mathcal{F}$
is injective,
\item if $\mathcal{F} \subset \mathcal{G}$ where $\mathcal{G}$
on $T_\etale$ is annihilated by $n$, then
$$
\Coker\left(
(f')^{-1}R^qg_*\mathcal{F}\to R^qh_*e^{-1}\mathcal{F}
\right)
\subset
\Coker\left(
(f')^{-1}R^qg_*\mathcal{G}\to R^qh_*e^{-1}\mathcal{G}
\right)
$$
\item if in (2) the sheaf $\mathcal{G}$ is an injective sheaf
of $\mathbf{Z}/n\mathbf{Z}$-modules, then
$$
\Coker\left((f')^{-1}R^qg_*\mathcal{F}\to R^qh_*e^{-1}\mathcal{F} \right)
\subset R^qh_*e^{-1}\mathcal{G}
$$
\end{enumerate}
\end{lemma}

\begin{proof}
Choose a short exact sequence
$0 \to \mathcal{F} \to \mathcal{I} \to \mathcal{Q} \to 0$
where $\mathcal{I}$ is an injective sheaf of $\mathbf{Z}/n\mathbf{Z}$-modules.
Consider the induced diagram
$$
\xymatrix{
(f')^{-1}R^{q - 1}g_*\mathcal{I} \ar[d]_{\cong} \ar[r] &
(f')^{-1}R^{q - 1}g_*\mathcal{Q} \ar[d]_{\cong} \ar[r] &
(f')^{-1}R^qg_*\mathcal{F} \ar[d] \ar[r] &
0 \ar[d] \\
R^{q - 1}h_*e^{-1}\mathcal{I} \ar[r] &
R^{q - 1}h_*e^{-1}\mathcal{Q} \ar[r] &
R^qh_*e^{-1}\mathcal{F} \ar[r] &
R^qh_*e^{-1}\mathcal{I}
}
$$
with exact rows. We have the zero in the right upper corner
as $\mathcal{I}$ is injective. The left two vertical arrows are
isomorphisms by $BC(f, n, q - 1)$. We conclude that part (1) holds.
The above also shows that
$$
\Coker\left(
(f')^{-1}R^qg_*\mathcal{F}\to R^qh_*e^{-1}\mathcal{F}
\right)
\subset
R^qh_*e^{-1}\mathcal{I}
$$
hence part (3) holds. To prove (2) choose
$\mathcal{F} \subset \mathcal{G} \subset \mathcal{I}$.
\end{proof}

\begin{lemma}
\label{lemma-base-change-q-integral-top}
With $f : X \to S$ and $n$ as in Remark \ref{remark-base-change-holds}
assume for some $q \geq 1$ we have $BC(f, n, q - 1)$. Consider
commutative diagrams
$$
\vcenter{
\xymatrix{
X \ar[d]_f &
X' \ar[d]_{f'} \ar[l] &
Y \ar[l]^h \ar[d]^e &
Y' \ar[l]^{\pi'} \ar[d]^{e'} \\
S &
S' \ar[l] &
T \ar[l]_g &
T' \ar[l]_\pi
}
}
\quad\text{and}\quad
\vcenter{
\xymatrix{
X' \ar[d]_{f'} & & Y' \ar[ll]^{h' = h \circ \pi'} \ar[d]^{e'} \\
S' & & T' \ar[ll]_{g' = g \circ \pi}
}
}
$$
where all squares are cartesian, $g$ quasi-compact and quasi-separated, and
$\pi$ is integral surjective. Let $\mathcal{F}$ be an abelian sheaf
on $T_\etale$ annihilated by $n$ and set $\mathcal{F}' = \pi^{-1}\mathcal{F}$.
If the base change map
$$
(f')^{-1}R^qg'_*\mathcal{F}' \longrightarrow R^qh'_*(e')^{-1}\mathcal{F}'
$$
is an isomorphism, then the base change map
$(f')^{-1}R^qg_*\mathcal{F} \to R^qh_*e^{-1}\mathcal{F}$
is an isomorphism.
\end{lemma}

\begin{proof}
Observe that $\mathcal{F} \to \pi_*\pi^{-1}\mathcal{F}'$ is injective
as $\pi$ is surjective (check on stalks). Thus by
Lemma \ref{lemma-base-change-q-injective}
we see that it suffices to show that the base change map
$$
(f')^{-1}R^qg_*\pi_*\mathcal{F}'
\longrightarrow
R^qh_*e^{-1}\pi_*\mathcal{F}'
$$
is an isomorphism. This follows from the assumption because
we have $R^qg_*\pi_*\mathcal{F}' = R^qg'_*\mathcal{F}'$,
we have $e^{-1}\pi_*\mathcal{F}' =\pi'_*(e')^{-1}\mathcal{F}'$, and
we have $R^qh_*\pi'_*(e')^{-1}\mathcal{F}' = R^qh'_*(e')^{-1}\mathcal{F}'$.
This follows from Lemmas
\ref{lemma-integral-pushforward-commutes-with-base-change} and
\ref{lemma-what-integral} and the relative leray spectral sequence
(Cohomology on Sites, Lemma \ref{sites-cohomology-lemma-relative-Leray}).
\end{proof}

\begin{lemma}
\label{lemma-base-change-q-integral-bottom}
With $f : X \to S$ and $n$ as in Remark \ref{remark-base-change-holds}
assume for some $q \geq 1$ we have $BC(f, n, q - 1)$. Consider
commutative diagrams
$$
\vcenter{
\xymatrix{
X \ar[d]_f &
X' \ar[d]_{f'} \ar[l] &
X'' \ar[l]^{\pi'} \ar[d]_{f''} &
Y \ar[l]^{h'} \ar[d]^e \\
S &
S' \ar[l] &
S'' \ar[l]_\pi &
T \ar[l]_{g'}
}
}
\quad\text{and}\quad
\vcenter{
\xymatrix{
X' \ar[d]_{f'} & & Y \ar[ll]^{h = h' \circ \pi'} \ar[d]^e \\
S' & & T \ar[ll]_{g = g' \circ \pi}
}
}
$$
where all squares are cartesian, $g'$ quasi-compact and quasi-separated, and
$\pi$ is integral. Let $\mathcal{F}$ be an abelian sheaf
on $T_\etale$ annihilated by $n$. If the base change map
$$
(f')^{-1}R^qg_*\mathcal{F} \longrightarrow R^qh_*e^{-1}\mathcal{F}
$$
is an isomorphism, then the base change map
$(f'')^{-1}R^qg'_*\mathcal{F} \to R^qh'_*e^{-1}\mathcal{F}$
is an isomorphism.
\end{lemma}

\begin{proof}
Since $\pi$ and $\pi'$ are integral we have $R\pi_* = \pi_*$ and
$R\pi'_* = \pi'_*$, see Lemma \ref{lemma-what-integral}.
We also have $(f')^{-1}\pi_* = \pi'_*(f'')^{-1}$. Thus we see that
$\pi'_*(f'')^{-1}R^qg'_*\mathcal{F} = (f')^{-1}R^qg_*\mathcal{F}$
and
$\pi'_*R^qh'_*e^{-1}\mathcal{F} = R^qh_*e^{-1}\mathcal{F}$.
Thus the assumption means that our map becomes an
isomorphism after applying the functor $\pi'_*$.
Hence we see that it is an isomorphism by Lemma \ref{lemma-what-integral}.
\end{proof}

\begin{lemma}
\label{lemma-formal-argument}
Let $T$ be a quasi-compact and quasi-separated scheme.
Let $P$ be a property for quasi-compact and quasi-separated
schemes over $T$. Assume
\begin{enumerate}
\item If $T'' \to T'$ is a thickening of quasi-compact and
quasi-separated schemes over $T$, then $P(T'')$ if and only if $P(T')$.
\item If $T' = \lim T_i$ is a limit of an inverse system of
quasi-compact and quasi-separated schemes over $T$ with affine
transition morphisms and $P(T_i)$ holds for all $i$, then
$P(T')$ holds.
\item If $Z \subset T'$ is a closed subscheme with
quasi-compact complement $V \subset T'$ and $P(T')$ holds,
then either $P(V)$ or $P(Z)$ holds.
\end{enumerate}
Then $P(T)$ implies $P(\Spec(K))$ for some morphism $\Spec(K) \to T$
where $K$ is a field.
\end{lemma}

\begin{proof}
Consider the set $\mathfrak T$ of closed subschemes $T' \subset T$
such that $P(T')$. By assumption (2) this set has a minimal element,
say $T'$. By assumption (1) we see that $T'$ is reduced.
Let $\eta \in T'$ be the generic point of an irreducible
component of $T'$. Then $\eta = \Spec(K)$ for some field
$K$ and $\eta = \lim V$ where the limit is over the affine
open subschemes $V \subset T'$ containing $\eta$.
By assumption (3) and the minimality of $T'$ we see
that $P(V)$ holds for all these $V$. Hence $P(\eta)$
by (2) and the proof is complete.
\end{proof}

\begin{lemma}
\label{lemma-base-change-does-not-hold-pre}
With $f : X \to S$ and $n$ as in Remark \ref{remark-base-change-holds}
assume for some $q \geq 1$ we have that $BC(f, n, q - 1)$ is true, but
$BC(f, n, q)$ is not. Then there exist a commutative diagram
$$
\xymatrix{
X \ar[d]_f & X' \ar[d]_{f'} \ar[l] & Y \ar[l]^h \ar[d]^e \\
S & S' \ar[l] & \Spec(K) \ar[l]_g
}
$$
where $X' = X \times_S S'$, $Y = X' \times_{S'} \Spec(K)$,
$K$ is a field, and $\mathcal{F}$ is an abelian sheaf
on $\Spec(K)$ annihilated by $n$ such that
$(f')^{-1}R^qg_*\mathcal{F} \to R^qh_*e^{-1}\mathcal{F}$
is not an isomorphism.
\end{lemma}

\begin{proof}
Choose a commutative diagram
$$
\xymatrix{
X \ar[d]_f & X' \ar[l] \ar[d]_{f'} & Y \ar[l]^h \ar[d]^e \\
S & S' \ar[l] & T \ar[l]_g
}
$$
with $X' = X \times_S S'$ and $Y = X' \times_{S'} T$ and
$g$ quasi-compact and quasi-separated, and an abelian sheaf
$\mathcal{F}$ on $T_\etale$ annihilated by $n$ such that
the base change map
$(f')^{-1}R^qg_*\mathcal{F} \to R^qh_*e^{-1}\mathcal{F}$
is not an isomorphism. Of course we may and do replace $S'$
by an affine open of $S'$; this implies that $T$ is quasi-compact
and quasi-separated. By Lemma \ref{lemma-base-change-q-injective} we see
$(f')^{-1}R^qg_*\mathcal{F} \to R^qh_*e^{-1}\mathcal{F}$
is injective. Pick a geometric point $\overline{x}$ of $X'$
and an element $\xi$ of $(R^qh_*q^{-1}\mathcal{F})_{\overline{x}}$
which is not in the image of the map
$((f')^{-1}R^qg_*\mathcal{F})_{\overline{x}} \to
(R^qh_*e^{-1}\mathcal{F})_{\overline{x}}$.

\medskip\noindent
Consider a morphism $\pi : T' \to T$ with $T'$ quasi-compact
and quasi-separated and denote $\mathcal{F}' = \pi^{-1}\mathcal{F}$.
Denote $\pi' : Y' = Y \times_T T' \to Y$ the base change of $\pi$
and $e' : Y' \to T'$ the base change of $e$. Picture
$$
\vcenter{
\xymatrix{
X' \ar[d]_{f'} & Y \ar[l]^h \ar[d]^e & Y' \ar[l]^{\pi'} \ar[d]^{e'} \\
S' & T \ar[l]_g & T' \ar[l]_\pi
}
}
\quad\text{and}\quad
\vcenter{
\xymatrix{
X' \ar[d]_{f'} & & Y' \ar[ll]^{h' = h \circ \pi'} \ar[d]^{e'} \\
S' & & T' \ar[ll]_{g' = g \circ \pi}
}
}
$$
Using pullback maps we obtain a canonical commutative diagram
$$
\xymatrix{
(f')^{-1}R^qg_*\mathcal{F} \ar[r] \ar[d] &
(f')^{-1}R^qg'_*\mathcal{F}' \ar[d] \\
R^qh_*e^{-1}\mathcal{F} \ar[r] &
R^qh'_*(e')^{-1}\mathcal{F}'
}
$$
of abelian sheaves on $X'$. Let $P(T')$ be the property
\begin{itemize}
\item The image $\xi'$ of $\xi$ in
$(Rh'_*(e')^{-1}\mathcal{F}')_{\overline{x}}$ is
not in the image of the map
$(f^{-1}R^qg'_*\mathcal{F}')_{\overline{x}} \to
(R^qh'_*(e')^{-1}\mathcal{F}')_{\overline{x}}$.
\end{itemize}
We claim that hypotheses (1), (2), and (3) of
Lemma \ref{lemma-formal-argument} hold for $P$
which proves our lemma.

\medskip\noindent
Condition (1) of Lemma \ref{lemma-formal-argument}
holds for $P$ because the \'etale topology of a scheme
and a thickening of the scheme is the same. See
Proposition \ref{proposition-topological-invariance}.

\medskip\noindent
Suppose that $I$ is a directed set and that $T_i$
is an inverse system over $I$ of quasi-compact and quasi-separated
schemes over $T$ with affine transition morphisms.
Set $T' = \lim T_i$. Denote $\mathcal{F}'$ and $\mathcal{F}_i$
the pullback of $\mathcal{F}$ to $T'$, resp.\ $T_i$. Consider
the diagrams
$$
\vcenter{
\xymatrix{
X \ar[d]_{f'} & Y \ar[l]^h \ar[d]^e & Y_i \ar[l]^{\pi_i'} \ar[d]^{e_i} \\
S & T \ar[l]_g & T_i \ar[l]_{\pi_i}
}
}
\quad\text{and}\quad
\vcenter{
\xymatrix{
X \ar[d]_{f'} & & Y_i \ar[ll]^{h_i = h \circ \pi_i'} \ar[d]^{e_i} \\
S & & T_i \ar[ll]_{g_i = g \circ \pi_i}
}
}
$$
as in the previous paragraph. It is clear that $\mathcal{F}'$ on
$T'$ is the colimit of the pullbacks of $\mathcal{F}_i$ to $T'$
and that $(e')^{-1}\mathcal{F}'$ is the colimit of the pullbacks
of $e_i^{-1}\mathcal{F}_i$ to $Y'$.
By Lemma \ref{lemma-relative-colimit-general}
we have
$$
R^qh'_*(e')^{-1}\mathcal{F}' = \colim R^qh_{i, *}e_i^{-1}\mathcal{F}_i
\quad\text{and}\quad
(f')^{-1}R^qg'_*\mathcal{F}' = \colim (f')^{-1}R^qg_{i, *}\mathcal{F}_i
$$
It follows that if $P(T_i)$ is true for all $i$, then
$P(T')$ holds. Thus condition (2) of Lemma \ref{lemma-formal-argument}
holds for $P$.

\medskip\noindent
The most interesting is condition (3) of Lemma \ref{lemma-formal-argument}.
Assume $T'$ is a quasi-compact and quasi-separated scheme over $T$
such that $P(T')$ is true.
Let $Z \subset T'$ be a closed subscheme with complement $V \subset T'$
quasi-compact. Consider the diagram
$$
\xymatrix{
Y' \times_{T'} Z \ar[d]_{e_Z} \ar[r]_{i'} &
Y' \ar[d]_{e'} &
Y' \times_{T'} V \ar[l]^{j'} \ar[d]^{e_V} \\
Z \ar[r]^i &
T' &
V \ar[l]_j
}
$$
Choose an injective map $j^{-1}\mathcal{F}' \to \mathcal{J}$
where $\mathcal{J}$ is an injective sheaf of $\mathbf{Z}/n\mathbf{Z}$-modules
on $V$. Looking at stalks we see that the map
$$
\mathcal{F}' \to \mathcal{G} = j_*\mathcal{J} \oplus i_*i^{-1}\mathcal{F}'
$$
is injective. Thus $\xi'$ maps to a nonzero element of
\begin{align*}
& \Coker\left(
((f')^{-1}R^qg'_*\mathcal{G})_{\overline{x}}
\to
(R^qh'_*(e')^{-1}\mathcal{G})_{\overline{x}}
\right) = \\
&
\Coker\left(
((f')^{-1}R^qg'_*j_*\mathcal{J})_{\overline{x}}
\to
(R^qh'_*(e')^{-1}j_*\mathcal{J})_{\overline{x}}
\right) \oplus \\
& \Coker\left(
((f')^{-1}R^qg'_*i_*i^{-1}\mathcal{F}')_{\overline{x}}
\to
(R^qh'_*(e')^{-1}i_*i^{-1}\mathcal{F}')_{\overline{x}}
\right)
\end{align*}
by part (2) of Lemma \ref{lemma-base-change-q-injective}.
If $\xi'$ does not map to zero in the second summand, then
we use
$$
(f')^{-1}R^qg'_*i_*i^{-1}\mathcal{F}' =
(f')^{-1}R^q(g' \circ i)_*i^{-1}\mathcal{F}'
$$
(because $Ri_* = i_*$ by
Proposition \ref{proposition-finite-higher-direct-image-zero}) and
$$
R^qh'_*(e')^{-1}i_*i^{-1}\mathcal{F} =
R^qh'_*i'_*e_Z^{-1}i^{-1}\mathcal{F} =
R^q(h' \circ i')_*e_Z^{-1}i^{-1}\mathcal{F}'
$$
(first equality by
Lemma \ref{lemma-finite-pushforward-commutes-with-base-change}
and the second because
$Ri'_* = i'_*$ by
Proposition \ref{proposition-finite-higher-direct-image-zero})
to we see that we have $P(Z)$.
Finally, suppose $\xi'$ does not map to zero in the first summand.
We have
$$
(e')^{-1}j_*\mathcal{J} = j'_*e_V^{-1}\mathcal{J}
\quad\text{and}\quad
R^aj'_*e_V^{-1}\mathcal{J} = 0, \quad a = 1, \ldots, q - 1
$$
by $BC(f, n, q - 1)$ applied to the diagram
$$
\xymatrix{
X \ar[d]_f & Y' \ar[l] \ar[d]_{e'} & Y \ar[l]^{j'} \ar[d]^{e_V} \\
S & T' \ar[l] & V \ar[l]_j
}
$$
and the fact that $\mathcal{J}$ is injective.
By the relative Leray spectral sequence for $h' \circ j'$
(Cohomology on Sites, Lemma \ref{sites-cohomology-lemma-relative-Leray})
we deduce that
$$
R^qh'_*(e')^{-1}j_*\mathcal{J} =
R^qh'_*j'_*e_V^{-1}\mathcal{J}
\longrightarrow
R^q(h' \circ j')_* e_V^{-1}\mathcal{J}
$$
is injective. Thus $\xi$ maps to a nonzero element of
$(R^q(h' \circ j')_* e_V^{-1}\mathcal{J})_{\overline{x}}$.
Applying part (3) of Lemma \ref{lemma-base-change-q-injective}
to the injection $j^{-1}\mathcal{F}' \to \mathcal{J}$
we conclude that $P(V)$ holds.
\end{proof}

\begin{lemma}
\label{lemma-base-change-does-not-hold}
With $f : X \to S$ and $n$ as in Remark \ref{remark-base-change-holds}
assume for some $q \geq 1$ we have that
$BC(f, n, q - 1)$ is true, but $BC(f, n, q)$ is not.
Then there exist a commutative diagram
$$
\xymatrix{
X \ar[d]_f & X' \ar[d] \ar[l] & Y \ar[l]^h \ar[d] \\
S & S' \ar[l] & \Spec(K) \ar[l]
}
$$
with both squares cartesian, where
\begin{enumerate}
\item $S'$ is affine, integral, and normal with algebraically
closed function field,
\item $K$ is algebraically closed and $\Spec(K) \to S'$
is dominant (in other words $K$ is an extension of
the function field of $S'$)
\end{enumerate}
and there exists an integer $d | n$
such that $R^qh_*(\mathbf{Z}/d\mathbf{Z})$ is nonzero.
\end{lemma}

\noindent
Conversely, nonvanishing of $R^qh_*(\mathbf{Z}/d\mathbf{Z})$
in the lemma implies $BC(f, n, q)$ isn't true as
Lemma \ref{lemma-Rf-star-zero-normal-with-alg-closed-function-field}
shows that $R^q(\Spec(K) \to S')_*\mathbf{Z}/d\mathbf{Z} = 0$.

\begin{proof}
First choose a diagram and $\mathcal{F}$ as in
Lemma \ref{lemma-base-change-does-not-hold-pre}.
We may and do assume $S'$ is affine (this is obvious, but
see proof of the lemma in case of doubt).
By Lemma \ref{lemma-base-change-q-integral-top}
we may assume $K$ is algebraically closed.
Then $\mathcal{F}$ corresponds to a $\mathbf{Z}/n\mathbf{Z}$-module.
Such a modules is a direct sum of copies of $\mathbf{Z}/d\mathbf{Z}$
for varying $d | n$ hence we may assume $\mathcal{F}$ is
constant with value $\mathbf{Z}/d\mathbf{Z}$.
By Lemma \ref{lemma-base-change-q-integral-bottom}
we may replace $S'$ by the normalization
of $S'$ in $\Spec(K)$ which finishes the proof.
\end{proof}










\section{Smooth base change}
\label{section-smooth-base-change}

\noindent
In this section we prove the smooth base change theorem.

\begin{lemma}
\label{lemma-smooth-base-change-fields}
Let $K/k$ be an extension of fields. Let $X$ be a smooth affine curve
over $k$ with a rational point $x \in X(k)$. Let $\mathcal{F}$ be an abelian
sheaf on $\Spec(K)$ annihilated by an integer $n$ invertible in $k$.
Let $q > 0$ and
$$
\xi \in H^q(X_K, (X_K \to \Spec(K))^{-1}\mathcal{F})
$$
There exist
\begin{enumerate}
\item finite extensions $K'/K$ and $k'/k$ with $k' \subset K'$,
\item a finite \'etale Galois cover $Z \to X_{k'}$ with group $G$
\end{enumerate}
such that the order of $G$ divides a power of $n$, such that
$Z \to X_{k'}$ is split over $x_{k'}$, and
such that $\xi$ dies in $H^q(Z_{K'}, (Z_{K'} \to \Spec(K))^{-1}\mathcal{F})$.
\end{lemma}

\begin{proof}
For $q > 1$ we know that $\xi$ dies in
$H^q(X_{\overline{K}}, (X_{\overline{K}} \to \Spec(K))^{-1}\mathcal{F})$
(Theorem \ref{theorem-vanishing-affine-curves}).
By Lemma \ref{lemma-directed-colimit-cohomology} we see that
this means there is a finite extension $K'/K$ such that
$\xi$ dies in $H^q(X_{K'}, (X_{K'} \to \Spec(K))^{-1}\mathcal{F})$.
Thus we can take $k' = k$ and $Z = X$ in this case.

\medskip\noindent
Assume $q = 1$. Recall that $\mathcal{F}$ corresponds to a
discrete module $M$ with continuous $\text{Gal}_K$-action, see
Lemma \ref{lemma-equivalence-abelian-sheaves-point}.
Since $M$ is $n$-torsion, it is the union of finite
$\text{Gal}_K$-stable subgroups. Thus
we reduce to the case where $M$ is a finite abelian group annihilated by $n$,
see Lemma \ref{lemma-colimit}. After replacing $K$ by a finite extension
we may assume that the action of $\text{Gal}_K$ on $M$ is trivial.
Thus we may assume $\mathcal{F} = \underline{M}$ is the constant
sheaf with value a finite abelian group $M$ annihilated by $n$.

\medskip\noindent
We can write $M$ as a direct sum of cyclic groups.
Any two finite \'etale Galois coverings whose Galois groups
have order invertible in $k$, can be dominated by a third one whose
Galois group has order invertible in $k$
(Fundamental Groups, Section \ref{pione-section-finite-etale-under-galois}).
Thus it suffices to prove the lemma when
$M = \mathbf{Z}/d\mathbf{Z}$ where $d | n$.

\medskip\noindent
Assume $M = \mathbf{Z}/d\mathbf{Z}$ where $d | n$.
In this case $\overline{\xi} = \xi|_{X_{\overline{K}}}$ is an element of
$$
H^1(X_{\overline{k}}, \mathbf{Z}/d\mathbf{Z}) =
H^1(X_{\overline{K}}, \mathbf{Z}/d\mathbf{Z})
$$
See Theorem \ref{theorem-vanishing-affine-curves}.
This group classifies $\mathbf{Z}/d\mathbf{Z}$-torsors, see
Cohomology on Sites, Lemma \ref{sites-cohomology-lemma-torsors-h1}.
The torsor corresponding to $\overline{\xi}$ (viewed as a sheaf on
$X_{\overline{k}, \etale}$) in turn gives rise to a finite \'etale morphism
$T \to X_{\overline{k}}$ endowed an action of $\mathbf{Z}/d\mathbf{Z}$
transitive on the fibre of $T$ over $x_{\overline{k}}$, see
Lemma \ref{lemma-characterize-finite-locally-constant}.
Choose a connected component $T' \subset T$ (if $\overline{\xi}$ has
order $d$, then $T$ is already connected).
Then $T' \to X_{\overline{k}}$ is a finite \'etale Galois
cover whose Galois group is a subgroup $G \subset \mathbf{Z}/d\mathbf{Z}$
(small detail omitted). Moreover the element $\overline{\xi}$ maps to zero
under the map $H^1(X_{\overline{k}}, \mathbf{Z}/d\mathbf{Z}) \to
H^1(T', \mathbf{Z}/d\mathbf{Z})$ as this is one
of the defining properties of $T$.

\medskip\noindent
Next, we use a limit argument to choose a finite extension $k'/k$
contained in $\overline{k}$ such that $T' \to X_{\overline{k}}$
descends to a finite \'etale Galois cover $Z \to X_{k'}$ with group $G$.
See Limits, Lemmas \ref{limits-lemma-descend-finite-presentation},
\ref{limits-lemma-descend-finite-finite-presentation}, and
\ref{limits-lemma-descend-etale}.
After increasing $k'$ we may assume that $Z$ splits over $x_{k'}$.
The image of $\xi$ in
$H^1(Z_{\overline{K}}, \mathbf{Z}/d\mathbf{Z})$ is zero by construction.
Thus by Lemma \ref{lemma-directed-colimit-cohomology}
we can find a finite subextension $\overline{K}/K'/K$
containing $k'$ such that $\xi$ dies in $H^1(Z_{K'}, \mathbf{Z}/d\mathbf{Z})$
and this finishes the proof.
\end{proof}

\begin{theorem}[Smooth base change]
\label{theorem-smooth-base-change}
Consider a cartesian diagram of schemes
$$
\xymatrix{
X \ar[d]_f & Y \ar[l]^h \ar[d]^e \\
S & T \ar[l]_g
}
$$
where $f$ is smooth and $g$ quasi-compact and quasi-separated. Then
$$
f^{-1}R^qg_*\mathcal{F} = R^qh_*e^{-1}\mathcal{F}
$$
for any $q$ and any abelian sheaf $\mathcal{F}$
on $T_\etale$ all of whose stalks at geometric points are torsion of
orders invertible on $S$.
\end{theorem}

\begin{proof}[First proof of smooth base change]
This proof is very long but more direct (using less general theory)
than the second proof given below.

\medskip\noindent
The theorem is local on $X_\etale$. More precisely, suppose we have
$U \to X$ \'etale such that $U \to S$ factors as $U \to V \to S$
with $V \to S$ \'etale. Then we can consider the cartesian square
$$
\xymatrix{
U \ar[d]_{f'} & U \times_X Y \ar[l]^{h'} \ar[d]^{e'} \\
V & V \times_S T \ar[l]_{g'}
}
$$
and setting $\mathcal{F}' = \mathcal{F}|_{V \times_S T}$
we have $f^{-1}R^qg_*\mathcal{F}|_U = (f')^{-1}R^qg'_*\mathcal{F}'$
and $R^qh_*e^{-1}\mathcal{F}|_U = R^qh'_*(e')^{-1}\mathcal{F}'$
(as follows from the compatibility of localization with morphisms of sites, see
Sites, Lemma \ref{sites-lemma-localize-morphism-strong} and
and
Cohomology on Sites, Lemma
\ref{sites-cohomology-lemma-restrict-direct-image-open}).
Thus it suffices to produce an \'etale covering of $X$ by
$U \to X$ and factorizations $U \to V \to S$
as above such that the theorem holds for the diagram with
$f'$, $h'$, $g'$, $e'$.

\medskip\noindent
By the local structure of smooth morphisms, see
Morphisms, Lemma \ref{morphisms-lemma-smooth-etale-over-affine-space},
we may assume $X$ and $S$ are affine and $X \to S$
factors through an \'etale morphism $X \to \mathbf{A}^d_S$.
If we have a tower of cartesian diagrams
$$
\xymatrix{
W \ar[d]_i & Z \ar[l]^j \ar[d]^k \\
X \ar[d]_f & Y \ar[l]^h \ar[d]^e \\
S & T \ar[l]_g
}
$$
and the theorem holds for the bottom and top squares, then
the theorem holds for the outer rectangle; this is formal.
Writing $X \to S$ as the composition
$$
X \to \mathbf{A}^{d - 1}_S \to \mathbf{A}^{d - 2}_S \to \ldots \to
\mathbf{A}^1_S \to S
$$
we conclude that it suffices to prove the theorem when $X$ and $S$
are affine and $X \to S$ has relative dimension $1$.

\medskip\noindent
For every $n \geq 1$ invertible on $S$, let $\mathcal{F}[n]$
be the subsheaf of sections of $\mathcal{F}$ annihilated by $n$. Then
$\mathcal{F} = \colim \mathcal{F}[n]$ by our assumption on
the stalks of $\mathcal{F}$. The functors $e^{-1}$ and $f^{-1}$
commute with colimits as they are left adjoints. The functors
$R^qh_*$ and $R^qg_*$ commute with filtered colimits by
Lemma \ref{lemma-relative-colimit}.
Thus it suffices to prove the theorem for $\mathcal{F}[n]$.
From now on we fix an integer $n$, we work with
sheaves of $\mathbf{Z}/n\mathbf{Z}$-modules and
we assume $S$ is a scheme over $\Spec(\mathbf{Z}[1/n])$.

\medskip\noindent
Next, we reduce to the case where $T$ is affine. Since $g$ is quasi-compact
and quasi-separate and $S$ is affine, the scheme $T$ is quasi-compact and
quasi-separated. Thus we can use the induction principle of
Cohomology of Schemes, Lemma \ref{coherent-lemma-induction-principle}.
Hence it suffices to show that if $T = W \cup W'$
is an open covering and the theorem holds for the squares
$$
\xymatrix{
X \ar[d] & e^{-1}(W) \ar[l]^i \ar[d] \\
S & W \ar[l]_a
}
\quad
\xymatrix{
X \ar[d] & e^{-1}(W') \ar[l]^j \ar[d] \\
S & W' \ar[l]_b
}
\quad
\xymatrix{
X \ar[d] & e^{-1}(W \cap W') \ar[l]^-k \ar[d] \\
S & W \cap W' \ar[l]_c
}
$$
then the theorem holds for the original diagram. To see this we consider the
diagram
$$
\xymatrix{
f^{-1}R^{q - 1}c_*\mathcal{F}|_{W \cap W'} \ar[d]^{\cong} \ar[r] &
f^{-1}R^qg_*\mathcal{F} \ar[d] \ar[r] &
f^{-1}R^qa_*\mathcal{F}|_W \oplus
f^{-1}R^qb_*\mathcal{F}|_{W'} \ar[d]_{\cong} \\
R^qk_*e^{-1}\mathcal{F}|_{e^{-1}(W \cap W')} \ar[r] &
R^qh_*e^{-1}\mathcal{F} \ar[r] &
R^qi_*e^{-1}\mathcal{F}|_{e^{-1}(W)} \oplus
R^qj_*e^{-1}\mathcal{F}|_{e^{-1}(W')}
}
$$
whose rows are the long exact sequences of
Lemma \ref{lemma-relative-mayer-vietoris}.
Thus the $5$-lemma gives the desired conclusion.

\medskip\noindent
Summarizing, we may assume $S$, $X$, $T$, and $Y$ affine,
$\mathcal{F}$ is $n$ torsion, $X \to S$ is smooth of relative dimension $1$,
and $S$ is a scheme over $\mathbf{Z}[1/n]$.
We will prove the theorem by induction on $q$. The base case $q = 0$
is handled by Lemma \ref{lemma-fppf-reduced-fibres-base-change-f-star}.
Assume $q > 0$ and the theorem holds for all smaller degrees.
Choose a short exact sequence
$0 \to \mathcal{F} \to \mathcal{I} \to \mathcal{Q} \to 0$
where $\mathcal{I}$ is an injective sheaf of $\mathbf{Z}/n\mathbf{Z}$-modules.
Consider the induced diagram
$$
\xymatrix{
f^{-1}R^{q - 1}g_*\mathcal{I} \ar[d]_{\cong} \ar[r] &
f^{-1}R^{q - 1}g_*\mathcal{Q} \ar[d]_{\cong} \ar[r] &
f^{-1}R^qg_*\mathcal{F} \ar[d] \ar[r] &
0 \ar[d] \\
R^{q - 1}h_*e^{-1}\mathcal{I} \ar[r] &
R^{q - 1}h_*e^{-1}\mathcal{Q} \ar[r] &
R^qh_*e^{-1}\mathcal{F} \ar[r] &
R^qh_*e^{-1}\mathcal{I}
}
$$
with exact rows. We have the zero in the right upper corner
as $\mathcal{I}$ is injective. The left two vertical arrows are
isomorphisms by induction hypothesis. Thus it suffices to prove
that $R^qh_*e^{-1}\mathcal{I} = 0$.

\medskip\noindent
Write $S = \Spec(A)$ and $T = \Spec(B)$ and say the morphism
$T \to S$ is given by the ring map $A \to B$. We can write
$A \to B = \colim_{i \in I} (A_i \to B_i)$ as a filtered
colimit of maps of rings of finite type over $\mathbf{Z}[1/n]$
(see Algebra, Lemma \ref{algebra-lemma-limit-no-condition}).
For $i \in I$ we set $S_i = \Spec(A_i)$ and $T_i = \Spec(B_i)$.
For $i$ large enough we can find a smooth morphism $X_i \to S_i$
of relative dimension $1$ such that $X = X_i \times_{S_i} S$, see
Limits, Lemmas \ref{limits-lemma-descend-finite-presentation},
\ref{limits-lemma-descend-smooth}, and
\ref{limits-lemma-descend-dimension-d}. Set $Y_i = X_i \times_{S_i} T_i$
to get squares
$$
\xymatrix{
X_i \ar[d]_{f_i} & Y_i \ar[l]^{h_i} \ar[d]^{e_i} \\
S_i & T_i \ar[l]_{g_i}
}
$$
Observe that $\mathcal{I}_i = (T \to T_i)_*\mathcal{I}$
is an injective sheaf of $\mathbf{Z}/n\mathbf{Z}$-modules on
$T_i$, see Cohomology on Sites, Lemma
\ref{sites-cohomology-lemma-pushforward-injective-flat}.
We have $\mathcal{I} = \colim (T \to T_i)^{-1}\mathcal{I}_i$
by Lemma \ref{lemma-linus-hamann}. Pulling back by $e$ we get
$e^{-1}\mathcal{I} = \colim (Y \to Y_i)^{-1}e_i^{-1}\mathcal{I}_i$.
By Lemma \ref{lemma-relative-colimit-general} applied to the system
of morphisms $Y_i \to X_i$ with limit $Y \to X$ we have
$$
R^qh_*e^{-1}\mathcal{I} =
\colim (X \to X_i)^{-1} R^qh_{i, *} e_i^{-1}\mathcal{I}_i
$$
This reduces us to the case where $T$ and $S$ are affine of finite
type over $\mathbf{Z}[1/n]$.

\medskip\noindent
Summarizing, we have an integer $q \geq 1$ such that the theorem holds
in degrees $< q$, the schemes $S$ and $T$ affine of finite type
type over $\mathbf{Z}[1/n]$, we have
$X \to S$ smooth of relative dimension $1$ with $X$ affine, and
$\mathcal{I}$ is an injective sheaf of $\mathbf{Z}/n\mathbf{Z}$-modules
and we have to show that $R^qh_*e^{-1}\mathcal{I} = 0$.
We will do this by induction on $\dim(T)$.

\medskip\noindent
The base case is $T = \emptyset$, i.e., $\dim(T) < 0$.
If you don't like this, you can take as your base case
the case $\dim(T) = 0$. In this case $T \to S$ is finite
(in fact even $T \to \Spec(\mathbf{Z}[1/n])$ is finite as the
target is Jacobson; details omitted), so
$h$ is finite too and hence has
vanishing higher direct images (see references below).

\medskip\noindent
Assume $\dim(T) = d \geq 0$ and we know the result for all situations where $T$
has lower dimension. Pick $U$ affine and \'etale over $X$ and a section $\xi$
of $R^qh_*q^{-1}\mathcal{I}$ over $U$. We have to show
that $\xi$ is zero. Of course, we may replace $X$ by $U$
(and correspondingly $Y$ by $U \times_X Y$)
and assume $\xi \in H^0(X, R^qh_*e^{-1}\mathcal{I})$.
Moreover, since $R^qh_*e^{-1}\mathcal{I}$ is a sheaf,
it suffices to prove that $\xi$ is zero locally on $X$.
Hence we may replace $X$ by the members of an \'etale covering.
In particular, using Lemma \ref{lemma-higher-direct-images}
we may assume that $\xi$ is the image of an element
$\tilde \xi \in H^q(Y, e^{-1}\mathcal{I})$.
In terms of $\tilde \xi$ our task is to show that
$\tilde \xi$ dies in $H^q(U_i \times_X Y, e^{-1}\mathcal{I})$
for some \'etale covering $\{U_i \to X\}$.

\medskip\noindent
By More on Morphisms, Lemma \ref{more-morphisms-lemma-cover-smooth-by-special}
we may assume that $X \to S$ factors as $X \to V \to S$
where $V \to S$ is \'etale and $X \to V$ is a smooth morphism
of affine schemes of relative dimension $1$, has a section, and
has geometrically connected fibres. Observe that
$\dim(V \times_S T) \leq \dim(T) = d$
for example by More on Algebra, Lemma
\ref{more-algebra-lemma-dimension-etale-extension}.
Hence we may then replace $S$ by $V$ and $T$ by $V \times_S T$
(exactly as in the discussion in the first paragraph of the proof).
Thus we may assume $X \to S$ is smooth of relative dimension $1$,
geometrically connected fibres, and has a section $\sigma : S \to X$.

\medskip\noindent
Let $\pi : T' \to T$ be a finite surjective morphism.
We will use below that $\dim(T') \leq \dim(T) = d$, see
Algebra, Lemma \ref{algebra-lemma-integral-dim-up}.
Choose an injective map $\pi^{-1}\mathcal{I} \to \mathcal{I}'$
into an injective sheaf of $\mathbf{Z}/n\mathbf{Z}$-modules.
Then $\mathcal{I} \to \pi_*\mathcal{I}'$ is injective
and hence has a splitting (as $\mathcal{I}$ is an injective
sheaf of $\mathbf{Z}/n\mathbf{Z}$-modules).
Denote $\pi' : Y' = Y \times_T T' \to Y$ the base change of $\pi$
and $e' : Y' \to T'$ the base change of $e$. Picture
$$
\xymatrix{
X \ar[d]_f & Y \ar[l]^h \ar[d]^e & Y' \ar[l]^{\pi'} \ar[d]^{e'} \\
S & T \ar[l]_g & T' \ar[l]_\pi
}
$$
By Proposition \ref{proposition-finite-higher-direct-image-zero} and
Lemma \ref{lemma-finite-pushforward-commutes-with-base-change} we have
$R\pi'_*(e')^{-1}\mathcal{I}' = e^{-1}\pi_*\mathcal{I}'$.
Thus by the Leray spectral sequence
(Cohomology on Sites, Lemma \ref{sites-cohomology-lemma-Leray})
we have
$$
H^q(Y', (e')^{-1}\mathcal{I}') =
H^q(Y, e^{-1}\pi_*\mathcal{I}') \supset H^q(Y, e^{-1}\mathcal{I})
$$
and this remains true after base change by any $U \to X$ \'etale.
Thus we may replace $T$ by $T'$, $\mathcal{I}$ by $\mathcal{I}'$
and $\tilde \xi$ by its image in $H^q(Y', (e')^{-1}\mathcal{I}')$.

\medskip\noindent
Suppose we have a factorization $T \to S' \to S$ where
$\pi : S' \to S$ is finite.
Setting $X' = S' \times_S X$ we can consider the induced diagram
$$
\xymatrix{
X \ar[d]_f & X' \ar[l]^{\pi'} \ar[d]^{f'} & Y \ar[l]^{h'} \ar[d]^e \\
S & S' \ar[l]_\pi & T \ar[l]_g
}
$$
Since $\pi'$ has vanishing higher direct images we see that
$R^qh_*e^{-1}\mathcal{I} = \pi'_*R^qh'_*e^{-1}\mathcal{I}$
by the Leray spectral sequence. Hence
$H^0(X, R^qh_*e^{-1}\mathcal{I}) = H^0(X', R^qh'_*e^{-1}\mathcal{I})$.
Thus $\xi$ is zero if and only if the corresponding section of
$R^qh'_*e^{-1}\mathcal{I}$ is zero\footnote{This
step can also be seen another way. Namely, we have to show that
there is an \'etale covering $\{U_i \to X\}$ such that
$\tilde \xi$ dies in $H^q(U_i \times_X Y, e^{-1}\mathcal{I})$.
However, if we prove there is an \'etale covering
$\{U'_j \to X'\}$ such that $\tilde \xi$ dies in
$H^q(U'_i \times_{X'} Y, e^{-1}\mathcal{I})$, then by property (B)
for $X' \to X$ (Lemma \ref{lemma-finite-B}) there exists an \'etale covering
$\{U_i \to X\}$ such that $U_i \times_X X'$ is a disjoint union
of schemes over $X'$ each of which factors through $U'_j$ for some $j$.
Thus we see that $\tilde \xi$ dies in $H^q(U_i \times_X Y, e^{-1}\mathcal{I})$
as desired.}.
Thus we may replace $S$ by $S'$ and $X$ by $X'$.
Observe that $\sigma : S \to X$ base changes to $\sigma' : S' \to X'$
and hence after this replacement it is still true that $X \to S$
has a section $\sigma$ and geometrically connected fibres.

\medskip\noindent
We will use that $S$ and $T$ are Nagata schemes, see
Algebra, Proposition \ref{algebra-proposition-ubiquity-nagata}
which will guarantee
that various normalizations are finite, see
Morphisms, Lemmas \ref{morphisms-lemma-nagata-normalization-finite} and
\ref{morphisms-lemma-nagata-normalization}.
In particular, we may first replace $T$ by its normalization
and then replace $S$ by the normalization of $S$ in $T$.
Then $T \to S$ is a disjoint union of dominant
morphisms of integral normal schemes, see
Morphisms, Lemma \ref{morphisms-lemma-normal-normalization}.
Clearly we may argue one connected component
at a time, hence we may assume $T \to S$ is a
dominant morphism of integral normal schemes.

\medskip\noindent
Let $s \in S$ and $t \in T$ be the generic points. By
Lemma \ref{lemma-smooth-base-change-fields} there exist finite field
extensions $K/\kappa(t)$ and $k/\kappa(s)$ such that $k$ is contained in $K$
and a finite \'etale Galois covering $Z \to X_k$ with Galois group $G$
of order dividing a power of $n$ split over $\sigma(\Spec(k))$
such that $\tilde \xi$ maps to zero in $H^q(Z_K, e^{-1}\mathcal{I}|_{Z_K})$.
Let $T' \to T$ be the normalization of $T$ in $\Spec(K)$
and let $S' \to S$ be the normalization of $S$ in $\Spec(k)$.
Then we obtain a commutative diagram
$$
\xymatrix{
S' \ar[d] & T' \ar[l] \ar[d] \\
S & T \ar[l]
}
$$
whose vertical arrows are finite. By the arguments given above we
may and do replace $S$ and $T$ by $S'$ and $T'$ (and correspondingly
$X$ by $X \times_S S'$ and $Y$ by $Y \times_T T'$). After this replacement
we conclude we have a finite \'etale Galois covering
$Z \to X_s$ of the generic fibre of $X \to S$
with Galois group $G$ of order dividing a power of $n$
split over $\sigma(s)$ such that $\tilde \xi$ maps to zero in
$H^q(Z_t, (Z_t \to Y)^{-1}e^{-1}\mathcal{I})$.
Here $Z_t = Z \times_S t = Z \times_s t = Z \times_{X_s} Y_t$.
Since $n$ is invertible on $S$,
by Fundamental Groups, Lemma \ref{pione-lemma-extend-covering}
we can find a finite \'etale morphism $U \to X$ whose restriction to $X_s$
is $Z$.

\medskip\noindent
At this point we replace $X$ by $U$ and $Y$ by $U \times_X Y$.
After this replacement it may
no longer be the case that the fibres of $X \to S$ are geometrically
connected (there still is a section but we won't use this), but what
we gain is that after this replacement $\tilde \xi$ maps to zero
in $H^q(Y_t, e^{-1}\mathcal{I})$, i.e., $\tilde \xi$ restricts to
zero on the generic fibre of $Y \to T$.

\medskip\noindent
Recall that $t$ is the spectrum of the function field of $T$, i.e.,
as a scheme $t$ is the limit of the nonempty affine open subschemes of $T$.
By Lemma \ref{lemma-directed-colimit-cohomology} we conclude there exists
a nonempty open subscheme $V \subset T$ such that $\tilde \xi$ maps to zero in
$H^q(Y \times_T V, e^{-1}\mathcal{I}|_{Y \times_T V})$.

\medskip\noindent
Denote $Z = T \setminus V$. Consider the diagram
$$
\xymatrix{
Y \times_T Z \ar[d]_{e_Z} \ar[r]_{i'} &
Y \ar[d]_e &
Y \times_T V \ar[l]^{j'} \ar[d]^{e_V} \\
Z \ar[r]^i &
T &
V \ar[l]_j
}
$$
Choose an injection $i^{-1}\mathcal{I} \to \mathcal{I}'$
into an injective sheaf of $\mathbf{Z}/n\mathbf{Z}$-modules on $Z$.
Looking at stalks we see that the map
$$
\mathcal{I} \to j_*\mathcal{I}|_V \oplus i_*\mathcal{I}'
$$
is injective and hence splits as $\mathcal{I}$ is an injective sheaf
of $\mathbf{Z}/n\mathbf{Z}$-modules. Thus it suffices to show that
$\tilde \xi$ maps to zero in
$$
H^q(Y, e^{-1}j_*\mathcal{I}|_V) \oplus
H^q(Y, e^{-1}i_*\mathcal{I}')
$$
at least after replacing $X$ by the members of an \'etale covering.
Observe that
$$
e^{-1}j_*\mathcal{I}|_V = j'_*e_V^{-1}\mathcal{I}|_V,\quad
e^{-1}i_*\mathcal{I}' = i'_*e_Z^{-1}\mathcal{I}'
$$
By induction hypothesis on $q$ we see that
$$
R^aj'_*e_V^{-1}\mathcal{I}|_V = 0, \quad a = 1, \ldots, q - 1
$$
By the Leray spectral sequence for $j'$ and the vanishing
above it follows that
$$
H^q(Y, j'_*(e_V^{-1}\mathcal{I}|_V))
\longrightarrow
H^q(Y \times_T V, e_V^{-1}\mathcal{I}_V) =
H^q(Y \times_T V, e^{-1}\mathcal{I}|_{Y \times_T V})
$$
is injective. Thus the vanishing of the image of $\tilde \xi$
in the first summand above because we know $\tilde \xi$ vanishes
in $H^q(Y \times_T V, e^{-1}\mathcal{I}|_{Y \times_T V})$.
Since $\dim(Z) < \dim(T) = d$ by induction the image of $\tilde \xi$
in the second summand
$$
H^q(Y, e^{-1}i_*\mathcal{I}') =
H^q(Y, i'_*e_Z^{-1}\mathcal{I}') =
H^q(Y \times_T Z, e_Z^{-1}\mathcal{I}')
$$
dies after replacing $X$ by the members of a suitable \'etale covering.
This finishes the proof of the smooth base change theorem.
\end{proof}

\begin{proof}[Second proof of smooth base change]
This proof is the same as the longer first proof;
it is shorter only in that we have split out the
arguments used in a number of lemmas.

\medskip\noindent
The case of $q = 0$ is
Lemma \ref{lemma-fppf-reduced-fibres-base-change-f-star}.
Thus we may assume $q > 0$ and the result is true
for all smaller degrees.

\medskip\noindent
For every $n \geq 1$ invertible on $S$, let $\mathcal{F}[n]$
be the subsheaf of sections of $\mathcal{F}$ annihilated by $n$. Then
$\mathcal{F} = \colim \mathcal{F}[n]$ by our assumption on
the stalks of $\mathcal{F}$. The functors $e^{-1}$ and $f^{-1}$
commute with colimits as they are left adjoints. The functors
$R^qh_*$ and $R^qg_*$ commute with filtered colimits by
Lemma \ref{lemma-relative-colimit}.
Thus it suffices to prove the theorem for $\mathcal{F}[n]$.
From now on we fix an integer $n$ invertible on $S$ and we work with
sheaves of $\mathbf{Z}/n\mathbf{Z}$-modules.

\medskip\noindent
By Lemma \ref{lemma-base-change-local} the question is \'etale local
on $X$ and $S$. By the local structure of smooth morphisms, see
Morphisms, Lemma \ref{morphisms-lemma-smooth-etale-over-affine-space},
we may assume $X$ and $S$ are affine and $X \to S$
factors through an \'etale morphism $X \to \mathbf{A}^d_S$.
Writing $X \to S$ as the composition
$$
X \to \mathbf{A}^{d - 1}_S \to \mathbf{A}^{d - 2}_S \to \ldots \to
\mathbf{A}^1_S \to S
$$
we conclude from Lemma \ref{lemma-base-change-compose}
that it suffices to prove the theorem when $X$ and $S$
are affine and $X \to S$ has relative dimension $1$.

\medskip\noindent
By Lemma \ref{lemma-base-change-does-not-hold} it suffices to show that
$R^qh_*\mathbf{Z}/d\mathbf{Z} = 0$ for $d | n$ whenever we have a
cartesian diagram
$$
\xymatrix{
X \ar[d] & Y \ar[d] \ar[l]^h \\
S & \Spec(K) \ar[l]
}
$$
where $X \to S$ is affine and smooth of relative dimension $1$,
$S$ is the spectrum of a normal
domain $A$ with algebraically closed fraction field $L$,
and $K/L$ is an extension of algebraically closed fields.

\medskip\noindent
Recall that $R^qh_*\mathbf{Z}/d\mathbf{Z}$ is the sheaf associated
to the presheaf
$$
U \longmapsto
H^q(U \times_X Y, \mathbf{Z}/d\mathbf{Z}) =
H^q(U \times_S \Spec(K), \mathbf{Z}/d\mathbf{Z})
$$
on $X_\etale$ (Lemma \ref{lemma-higher-direct-images}).
Thus it suffices to show: given $U$ and
$\xi \in H^q(U \times_S \Spec(K), \mathbf{Z}/d\mathbf{Z})$
there exists an \'etale covering $\{U_i \to U\}$
such that $\xi$ dies in $H^q(U_i \times_S \Spec(K), \mathbf{Z}/d\mathbf{Z})$.

\medskip\noindent
Of course we may take $U$ affine. Then $U \times_S \Spec(K)$
is a (smooth) affine curve over $K$ and hence we have vanishing
for $q > 1$ by Theorem \ref{theorem-vanishing-affine-curves}.

\medskip\noindent
Final case: $q = 1$. We may replace $U$ by the members of an
\'etale covering as in More on Morphisms, Lemma
\ref{more-morphisms-lemma-cover-smooth-by-special}.
Then $U \to S$ factors as $U \to V \to S$ where
$U \to V$ has geometrically connected fibres,
$U$, $V$ are affine, $V \to S$ is \'etale, and there is
a section $\sigma : V \to U$. By
Lemma \ref{lemma-normal-scheme-with-alg-closed-function-field}
we see that $V$ is isomorphic to a (finite) disjoint union of
(affine) open subschemes of $S$.
Clearly we may replace $S$ by one of these
and $X$ by the corresponding component of $U$.
Thus we may assume $X \to S$ has geometrically connected
fibres, has a section $\sigma$, and
$\xi \in H^1(Y, \mathbf{Z}/d\mathbf{Z})$.
Since $K$ and $L$ are algebraically closed we have
$$
H^1(X_L, \mathbf{Z}/d\mathbf{Z}) =
H^1(Y, \mathbf{Z}/d\mathbf{Z})
$$
See Lemma \ref{lemma-base-change-dim-1-separably-closed}.
Thus there is a finite \'etale Galois covering
$Z \to X_L$ with Galois group $G \subset \mathbf{Z}/d\mathbf{Z}$
which annihilates $\xi$.
You can either see this by looking at the statement or
proof of Lemma \ref{lemma-smooth-base-change-fields}
or by using directly that $\xi$ corresponds to a
$\mathbf{Z}/d\mathbf{Z}$-torsor over $X_L$.
Finally, by
Fundamental Groups, Lemma \ref{pione-lemma-extend-covering-general}
we find a (necessarily surjective)
finite \'etale morphism $X' \to X$
whose restriction to $X_L$ is $Z \to X_L$.
Since $\xi$ dies in $X'_K$ this finishes the proof.
\end{proof}

\noindent
The following immediate consequence of the smooth base change
theorem is what is often used in practice.

\begin{lemma}
\label{lemma-smooth-base-change-general}
Let $S$ be a scheme. Let $S' = \lim S_i$ be a directed inverse
limit of schemes $S_i$ smooth over $S$ with affine transition
morphisms. Let $f : X \to S$ be quasi-compact and quasi-separated
and form the fibre square
$$
\xymatrix{
X' \ar[d]_{f'} \ar[r]_{g'} & X \ar[d]^f \\
S' \ar[r]^g & S
}
$$
Then
$$
g^{-1}Rf_*E = R(f')_*(g')^{-1}E
$$
for any $E \in D^+(X_\etale)$ whose cohomology sheaves $H^q(E)$
have stalks which are torsion of orders invertible on $S$.
\end{lemma}

\begin{proof}
Consider the spectral sequences
$$
E_2^{p, q} = R^pf_*H^q(E)
\quad\text{and}\quad
{E'}_2^{p, q} = R^pf'_*H^q((g')^{-1}E) = R^pf'_*(g')^{-1}H^q(E)
$$
converging to $R^nf_*E$ and $R^nf'_*(g')^{-1}E$.
These spectral sequences are constructed in
Derived Categories, Lemma \ref{derived-lemma-two-ss-complex-functor}.
Combining the smooth base change theorem
(Theorem \ref{theorem-smooth-base-change})
with Lemma \ref{lemma-base-change-Rf-star-colim} we see that
$$
g^{-1}R^pf_*H^q(E) = R^p(f')_*(g')^{-1}H^q(E)
$$
Combining all of the above we get the lemma.
\end{proof}












\section{Applications of smooth base change}
\label{section-applications-smooth-base-change}

\noindent
In this section we discuss some more or less immediate
consequences of the smooth base change theorem.

\begin{lemma}
\label{lemma-base-change-field-extension}
Let $L/K$ be an extension of fields. Let $g : T \to S$ be a quasi-compact
and quasi-separated morphism of schemes over $K$. Denote
$g_L : T_L \to S_L$ the base change of $g$ to $\Spec(L)$.
Let $E \in D^+(T_\etale)$ have cohomology sheaves whose stalks
are torsion of orders invertible in $K$. Let $E_L$ be the
pullback of $E$ to $(T_L)_\etale$. Then
$Rg_{L, *}E_L$ is the pullback of $Rg_*E$ to $S_L$.
\end{lemma}

\begin{proof}
If $L/K$ is separable, then $L$ is a filtered colimit of smooth
$K$-algebras, see
Algebra, Lemma \ref{algebra-lemma-colimit-syntomic}.
Thus the lemma in this case follows immediately from
Lemma \ref{lemma-smooth-base-change-general}.
In the general case, let $K'$ and $L'$ be the perfect closures
(Algebra, Definition \ref{algebra-definition-perfection})
of $K$ and $L$. Then $\Spec(K') \to \Spec(K)$ and
$\Spec(L') \to \Spec(L)$ are universal homeomorphisms
as $K'/K$ and $L'/L$ are purely inseparable
(see Algebra, Lemma \ref{algebra-lemma-p-ring-map}).
Thus we have $(T_{K'})_\etale = T_\etale$,
$(S_{K'})_\etale = S_\etale$,
$(T_{L'})_\etale = (T_L)\etale$, and
$(S_{L'})_\etale = (S_L)_\etale$ by
the topological invariance of \'etale cohomology, see
Proposition \ref{proposition-topological-invariance}.
This reduces the lemma to the case of the field
extension $L'/K'$ which is separable (by definition of
perfect fields, see Algebra, Definition \ref{algebra-definition-perfect}).
\end{proof}

\begin{lemma}
\label{lemma-smooth-base-change-separably-closed}
Let $K/k$ be an extension of separably closed fields. Let $X$
be a quasi-compact and quasi-separated scheme over $k$.
Let $E \in D^+(X_\etale)$ have cohomology sheaves whose stalks
are torsion of orders invertible in $k$. Then
\begin{enumerate}
\item the maps $H^q_\etale(X, E) \to H^q_\etale(X_K, E|_{X_K})$
are isomorphisms, and
\item $E \to R(X_K \to X)_*E|_{X_K}$ is an isomorphism.
\end{enumerate}
\end{lemma}

\begin{proof}
Proof of (1). First let $\overline{k}$ and $\overline{K}$
be the algebraic closures
of $k$ and $K$. The morphisms $\Spec(\overline{k}) \to \Spec(k)$ and
$\Spec(\overline{K}) \to \Spec(K)$ are universal homeomorphisms
as $\overline{k}/k$ and $\overline{K}/K$ are purely inseparable
(see Algebra, Lemma \ref{algebra-lemma-p-ring-map}).
Thus $H^q_\etale(X, \mathcal{F}) =
H^q_\etale(X_{\overline{k}}, \mathcal{F}_{X_{\overline{k}}})$ by
the topological invariance of \'etale cohomology, see
Proposition \ref{proposition-topological-invariance}.
Similarly for $X_K$ and $X_{\overline{K}}$.
Thus we may assume $k$ and $K$ are algebraically closed.
In this case $K$ is a limit of smooth $k$-algebras, see
Algebra, Lemma \ref{algebra-lemma-colimit-syntomic}.
We conclude our lemma is a special case of
Theorem \ref{theorem-smooth-base-change} as reformulated in
Lemma \ref{lemma-smooth-base-change-general}.

\medskip\noindent
Proof of (2). For any quasi-compact and quasi-separated $U$ in $X_\etale$
the above shows that the restriction of the map
$E \to R(X_K \to X)_*E|_{X_K}$ determines an isomorphism on cohomology.
Since every object of $X_\etale$ has an \'etale covering by such $U$
this proves the desired statement.
\end{proof}

\begin{lemma}
\label{lemma-base-change-does-not-hold-post}
With $f : X \to S$ and $n$ as in Remark \ref{remark-base-change-holds}
assume $n$ is invertible on $S$ and that for some $q \geq 1$
we have that $BC(f, n, q - 1)$ is true, but $BC(f, n, q)$ is not.
Then there exist a commutative diagram
$$
\xymatrix{
X \ar[d]_f & X' \ar[d] \ar[l] & Y \ar[l]^h \ar[d] \\
S & S' \ar[l] & \Spec(K) \ar[l]
}
$$
with both squares cartesian, where $S'$ is affine, integral, and normal
with algebraically closed function field $K$ and there exists an integer
$d | n$ such that $R^qh_*(\mathbf{Z}/d\mathbf{Z})$ is nonzero.
\end{lemma}

\begin{proof}
First choose a diagram and $\mathcal{F}$ as in
Lemma \ref{lemma-base-change-does-not-hold}.
We may and do assume $S'$ is affine (this is obvious, but
see proof of the lemma in case of doubt).
Let $K'$ be the function field of $S'$
and let $Y' = X' \times_{S'} \Spec(K')$ to get the diagram
$$
\xymatrix{
X \ar[d]_f &
X' \ar[d] \ar[l] &
Y' \ar[l]^{h'} \ar[d] &
Y \ar[l] \ar[d] \\
S &
S' \ar[l] &
\Spec(K') \ar[l] &
\Spec(K) \ar[l]
}
$$
By Lemma \ref{lemma-smooth-base-change-separably-closed}
the total direct image $R(Y \to Y')_*\mathbf{Z}/d\mathbf{Z}$
is isomorphic to $\mathbf{Z}/d\mathbf{Z}$ in $D(Y'_\etale)$; here
we use that $n$ is invertible on $S$.
Thus $Rh'_*\mathbf{Z}/d\mathbf{Z} = Rh_*\mathbf{Z}/d\mathbf{Z}$
by the relative Leray spectral sequence. This finishes the proof.
\end{proof}










\section{The proper base change theorem}
\label{section-proper-base-change}

\noindent
The proper base change theorem is stated and proved in this section.
Our approach follows roughly the proof in \cite[XII, Theorem 5.1]{SGA4}
using Gabber's ideas (from the affine case) to slightly simplify
the arguments.

\begin{lemma}
\label{lemma-zariski-h0-proper-over-henselian-pair}
Let $(A, I)$ be a henselian pair. Let $f : X \to \Spec(A)$ be a proper morphism
of schemes. Let $Z = X \times_{\Spec(A)} \Spec(A/I)$. For any
sheaf $\mathcal{F}$ on the topological space associated to $X$ we
have $\Gamma(X, \mathcal{F}) = \Gamma(Z, \mathcal{F}|_Z)$.
\end{lemma}

\begin{proof}
We will use Lemma \ref{lemma-h0-topological} to prove this. First observe
that the underlying topological space of $X$ is spectral by Properties, Lemma
\ref{properties-lemma-quasi-compact-quasi-separated-spectral}.
Let $Y \subset X$ be an irreducible closed subscheme. To finish the proof
we show that $Y \cap Z = Y \times_{\Spec(A)} \Spec(A/I)$ is connected.
Replacing $X$ by $Y$
we may assume that $X$ is irreducible and we have to show that $Z$
is connected. Let $X \to \Spec(B) \to \Spec(A)$ be the Stein factorization
of $f$ (More on Morphisms, Theorem
\ref{more-morphisms-theorem-stein-factorization-general}).
Then $A \to B$ is integral and $(B, IB)$ is a henselian pair
(More on Algebra, Lemma \ref{more-algebra-lemma-integral-over-henselian-pair}).
Thus we may assume the fibres of $X \to \Spec(A)$ are geometrically
connected. On the other hand, the image $T \subset \Spec(A)$ of $f$
is irreducible and closed as $X$ is proper over $A$. Hence $T \cap V(I)$
is connected by More on Algebra, Lemma
\ref{more-algebra-lemma-irreducible-henselian-pair-connected}.
Now $Y \times_{\Spec(A)} \Spec(A/I) \to T \cap V(I)$
is a surjective closed map with connected fibres.
The result now follows from Topology, Lemma
\ref{topology-lemma-connected-fibres-quotient-topology-connected-components}.
\end{proof}

\begin{lemma}
\label{lemma-h0-proper-over-henselian-pair}
Let $(A, I)$ be a henselian pair. Let $f : X \to \Spec(A)$ be a proper morphism
of schemes. Let $i : Z \to X$ be the closed immersion of
$X \times_{\Spec(A)} \Spec(A/I)$ into $X$. For any
sheaf $\mathcal{F}$ on $X_\etale$ we
have $\Gamma(X, \mathcal{F}) = \Gamma(Z, i_{small}^{-1}\mathcal{F})$.
\end{lemma}

\begin{proof}
This follows from Lemma \ref{lemma-gabber-h0} and
\ref{lemma-zariski-h0-proper-over-henselian-pair}
and the fact that any scheme finite over $X$ is proper over $\Spec(A)$.
\end{proof}

\begin{lemma}
\label{lemma-h0-proper-over-henselian-local}
Let $A$ be a henselian local ring. Let $f : X \to \Spec(A)$
be a proper morphism of schemes. Let $X_0 \subset X$ be the fibre of
$f$ over the closed point. For any sheaf $\mathcal{F}$ on $X_\etale$ we
have $\Gamma(X, \mathcal{F}) = \Gamma(X_0, \mathcal{F}|_{X_0})$.
\end{lemma}

\begin{proof}
This is a special case of Lemma \ref{lemma-h0-proper-over-henselian-pair}.
\end{proof}

\noindent
Let $f : X \to S$ be a morphism of schemes. Let
$\overline{s} : \Spec(k) \to S$ be a geometric point. The fibre
of $f$ at $\overline{s}$ is the scheme
$X_{\overline{s}} = \Spec(k) \times_{\overline{s}, S} X$ viewed
as a scheme over $\Spec(k)$. If $\mathcal{F}$ is a sheaf on
$X_\etale$, then denote
$\mathcal{F}_{\overline{s}} = p_{small}^{-1}\mathcal{F}$
the pullback of $\mathcal{F}$ to $(X_{\overline{s}})_\etale$.
In the following we will consider the set
$$
\Gamma(X_{\overline{s}}, \mathcal{F}_{\overline{s}})
$$
Let $s \in S$ be the image point of $\overline{s}$. Let $\kappa(s)^{sep}$
be the separable algebraic closure of $\kappa(s)$ in $k$ as in
Definition \ref{definition-algebraic-geometric-point}.
By Lemma \ref{lemma-sections-base-field-extension}
pullback defines a bijection
$$
\Gamma(X_{\kappa(s)^{sep}},
p_{sep}^{-1} \mathcal{F})
\longrightarrow
\Gamma(X_{\overline{s}}, \mathcal{F}_{\overline{s}})
$$
where $p_{sep} : X_{\kappa(s)^{sep}} = \Spec(\kappa(s)^{sep}) \times_S X \to X$
is the projection.

\begin{lemma}
\label{lemma-proper-pushforward-stalk}
Let $f : X \to S$ be a proper morphism of schemes. Let
$\overline{s} \to S$ be a geometric point.
For any sheaf $\mathcal{F}$ on $X_\etale$
the canonical map
$$
(f_*\mathcal{F})_{\overline{s}} \longrightarrow
\Gamma(X_{\overline{s}}, \mathcal{F}_{\overline{s}})
$$
is bijective.
\end{lemma}

\begin{proof}
By Theorem \ref{theorem-higher-direct-images} (for sheaves of sets)
we have
$$
(f_*\mathcal{F})_{\overline{s}} =
\Gamma(X \times_S \Spec(\mathcal{O}_{S, \overline{s}}^{sh}),
p_{small}^{-1}\mathcal{F})
$$
where $p : X \times_S \Spec(\mathcal{O}_{S, \overline{s}}^{sh}) \to X$
is the projection. Since the residue field of the strictly henselian
local ring $\mathcal{O}_{S, \overline{s}}^{sh}$ is $\kappa(s)^{sep}$
we conclude from the discussion above the lemma and
Lemma \ref{lemma-h0-proper-over-henselian-local}.
\end{proof}

\begin{lemma}
\label{lemma-proper-base-change-f-star}
Let $f : X \to Y$ be a proper morphism of schemes. Let $g : Y' \to Y$
be a morphism of schemes. Set $X' = Y' \times_Y X$ with projections
$f' : X' \to Y'$ and $g' : X' \to X$. Let $\mathcal{F}$ be any sheaf on
$X_\etale$. Then $g^{-1}f_*\mathcal{F} = f'_*(g')^{-1}\mathcal{F}$.
\end{lemma}

\begin{proof}
There is a canonical map $g^{-1}f_*\mathcal{F} \to f'_*(g')^{-1}\mathcal{F}$.
Namely, it is adjoint to the map
$$
f_*\mathcal{F} \longrightarrow
g_*f'_*(g')^{-1}\mathcal{F} = f_*g'_*(g')^{-1}\mathcal{F}
$$
which is $f_*$ applied to the canonical map
$\mathcal{F} \to g'_*(g')^{-1}\mathcal{F}$. To check this map is an
isomorphism we can compute what happens on stalks.
Let $y' : \Spec(k) \to Y'$ be a geometric point with image $y$ in $Y$.
By Lemma \ref{lemma-proper-pushforward-stalk} the stalks are
$\Gamma(X'_{y'}, \mathcal{F}_{y'})$ and $\Gamma(X_y, \mathcal{F}_y)$
respectively. Here the sheaves $\mathcal{F}_y$ and $\mathcal{F}_{y'}$
are the pullbacks of $\mathcal{F}$ by the projections $X_y \to X$
and $X'_{y'} \to X$. Since $X_y = X'_{y'}$ as schemes
we see that the stalks agree. We omit the
verification that this isomorphism is compatible with our map.
\end{proof}

\noindent
At this point we start discussing the proper base change theorem.
To do so we introduce some notation. consider a commutative diagram
\begin{equation}
\label{equation-base-change-diagram}
\vcenter{
\xymatrix{
X' \ar[r]_{g'} \ar[d]_{f'} & X \ar[d]^f \\
Y' \ar[r]^g & Y
}
}
\end{equation}
of morphisms of schemes. Then we obtain a commutative diagram of sites
$$
\xymatrix{
X'_\etale \ar[r]_{g'_{small}} \ar[d]_{f'_{small}} &
X_\etale \ar[d]^{f_{small}} \\
Y'_\etale \ar[r]^{g_{small}} &
Y_\etale
}
$$
For any object $E$ of $D(X_\etale)$ we obtain a canonical base change map
\begin{equation}
\label{equation-base-change}
g_{small}^{-1}Rf_{small, *}E \longrightarrow Rf'_{small, *}(g'_{small})^{-1}E
\end{equation}
in $D(Y'_\etale)$. See Cohomology on Sites, Remark
\ref{sites-cohomology-remark-base-change} where we use the constant
sheaf $\mathbf{Z}$ as our sheaf of rings.
We will usually omit the subscripts ${}_{small}$ in this formula.
For example, if $E = \mathcal{F}[0]$ where $\mathcal{F}$ is an abelian
sheaf on $X_\etale$, the base change map is a map
\begin{equation}
\label{equation-base-change-sheaf}
g^{-1}Rf_*\mathcal{F} \longrightarrow Rf'_*(g')^{-1}\mathcal{F}
\end{equation}
in $D(Y'_\etale)$.

\medskip\noindent
The map (\ref{equation-base-change}) has no chance of being an isomorphism
in the generality given above. The goal is to show
it is an isomorphism if the diagram (\ref{equation-base-change-diagram})
is cartesian, $f : X \to Y$ proper, the cohomology sheaves of $E$
are torsion, and $E$ is bounded below.
To study this question we introduce the following terminology.
Let us say that {\it cohomology commutes with base change
for $f : X \to Y$} if (\ref{equation-base-change-sheaf})
is an isomorphism for every diagram (\ref{equation-base-change-diagram})
where $X' = Y' \times_Y X$ and every torsion abelian sheaf $\mathcal{F}$.

\begin{lemma}
\label{lemma-proper-base-change-in-terms-of-injectives}
Let $f : X \to Y$ be a proper morphism of schemes.
The following are equivalent
\begin{enumerate}
\item cohomology commutes with base change for $f$ (see above),
\item for every prime number $\ell$ and every injective
sheaf of $\mathbf{Z}/\ell\mathbf{Z}$-modules $\mathcal{I}$
on $X_\etale$ and every diagram (\ref{equation-base-change-diagram})
where $X' = Y' \times_Y X$ the sheaves
$R^qf'_*(g')^{-1}\mathcal{I}$ are zero for $q > 0$.
\end{enumerate}
\end{lemma}

\begin{proof}
It is clear that (1) implies (2). Conversely, assume (2) and let
$\mathcal{F}$ be a torsion abelian sheaf on $X_\etale$. Let $Y' \to Y$
be a morphism of schemes and let $X' = Y' \times_Y X$
with projections $g' : X' \to X$ and $f' : X' \to Y'$ as in
diagram (\ref{equation-base-change-diagram}).
We want to show the maps of sheaves
$$
g^{-1}R^qf_*\mathcal{F} \longrightarrow R^qf'_*(g')^{-1}\mathcal{F}
$$
are isomorphisms for all $q \geq 0$.

\medskip\noindent
For every $n \geq 1$, let $\mathcal{F}[n]$ be the subsheaf of sections
of $\mathcal{F}$ annihilated by $n$. Then
$\mathcal{F} = \colim \mathcal{F}[n]$.
The functors $g^{-1}$ and $(g')^{-1}$ commute with arbitrary colimits
(as left adjoints). Taking higher direct images along $f$ or $f'$
commutes with filtered colimits by Lemma \ref{lemma-relative-colimit}.
Hence we see that
$$
g^{-1}R^qf_*\mathcal{F} = \colim g^{-1}R^qf_*\mathcal{F}[n]
\quad\text{and}\quad
R^qf'_*(g')^{-1}\mathcal{F} =
\colim R^qf'_*(g')^{-1}\mathcal{F}[n]
$$
Thus it suffices to prove the result in case $\mathcal{F}$ is
annihilated by a positive integer $n$.

\medskip\noindent
If $n = \ell n'$ for some prime number $\ell$, then we obtain a short
exact sequence
$$
0 \to \mathcal{F}[\ell] \to \mathcal{F} \to
\mathcal{F}/\mathcal{F}[\ell] \to 0
$$
Observe that $\mathcal{F}/\mathcal{F}[\ell]$ is annihilated by $n'$.
Moreover, if the result holds for both $\mathcal{F}[\ell]$ and
$\mathcal{F}/\mathcal{F}[\ell]$, then the result holds by
the long exact sequence of higher direct images (and the $5$ lemma).
In this way we reduce to the case that $\mathcal{F}$ is annihilated
by a prime number $\ell$.

\medskip\noindent
Assume $\mathcal{F}$ is annihilated by a prime number $\ell$.
Choose an injective resolution $\mathcal{F} \to \mathcal{I}^\bullet$
in $D(X_\etale, \mathbf{Z}/\ell\mathbf{Z})$. Applying assumption
(2) and Leray's acyclicity lemma
(Derived Categories, Lemma \ref{derived-lemma-leray-acyclicity})
we see that
$$
f'_*(g')^{-1}\mathcal{I}^\bullet
$$
computes $Rf'_*(g')^{-1}\mathcal{F}$. We conclude by applying
Lemma \ref{lemma-proper-base-change-f-star}.
\end{proof}

\begin{lemma}
\label{lemma-sandwich}
Let $f : X \to Y$ and $g : Y \to Z$ be proper morphisms of schemes. Assume
\begin{enumerate}
\item cohomology commutes with base change for $f$,
\item cohomology commutes with base change for $g \circ f$, and
\item $f$ is surjective.
\end{enumerate}
Then cohomology commutes with base change for $g$.
\end{lemma}

\begin{proof}
We will use the equivalence of
Lemma \ref{lemma-proper-base-change-in-terms-of-injectives}
without further mention. Let $\ell$ be a prime number.
Let $\mathcal{I}$ be an injective sheaf of
$\mathbf{Z}/\ell\mathbf{Z}$-modules on $Y_\etale$.
Choose an injective map of sheaves $f^{-1}\mathcal{I} \to \mathcal{J}$
where $\mathcal{J}$ is an injective sheaf of
$\mathbf{Z}/\ell\mathbf{Z}$-modules on $X_\etale$.
Since $f$ is surjective the map $\mathcal{I} \to f_*\mathcal{J}$
is injective (look at stalks in geometric points).
Since $\mathcal{I}$ is injective we see that $\mathcal{I}$
is a direct summand of $f_*\mathcal{J}$. Thus it suffices
to prove the desired vanishing for $f_*\mathcal{J}$.

\medskip\noindent
Let $Z' \to Z$ be a morphism of schemes and set
$Y' = Z' \times_Z Y$ and $X' = Z' \times_Z X = Y' \times_ Y X$.
Denote $a : X' \to X$, $b : Y' \to Y$, and $c : Z' \to Z$ the
projections. Similarly for $f' : X' \to Y'$ and $g' : Y' \to Z'$.
By Lemma \ref{lemma-proper-base-change-f-star} we have
$b^{-1}f_*\mathcal{J} = f'_*a^{-1}\mathcal{J}$.
On the other hand, we know that $R^qf'_*a^{-1}\mathcal{J}$ and
$R^q(g' \circ f')_*a^{-1}\mathcal{J}$ are zero for $q > 0$.
Using the spectral sequence
(Cohomology on Sites, Lemma \ref{sites-cohomology-lemma-relative-Leray})
$$
R^pg'_* R^qf'_* a^{-1}\mathcal{J} \Rightarrow
R^{p + q}(g' \circ f')_* a^{-1}\mathcal{J}
$$
we conclude that
$ R^pg'_*(b^{-1}f_*\mathcal{J}) = R^pg'_*(f'_*a^{-1}\mathcal{J}) = 0$
for $p > 0$ as desired.
\end{proof}

\begin{lemma}
\label{lemma-composition}
Let $f : X \to Y$ and $g : Y \to Z$ be proper morphisms of schemes. Assume
\begin{enumerate}
\item cohomology commutes with base change for $f$, and
\item cohomology commutes with base change for $g$.
\end{enumerate}
Then cohomology commutes with base change for $g \circ f$.
\end{lemma}

\begin{proof}
We will use the equivalence of
Lemma \ref{lemma-proper-base-change-in-terms-of-injectives}
without further mention. Let $\ell$ be a prime number.
Let $\mathcal{I}$ be an injective sheaf of
$\mathbf{Z}/\ell\mathbf{Z}$-modules on $X_\etale$.
Then $f_*\mathcal{I}$ is an injective sheaf of
$\mathbf{Z}/\ell\mathbf{Z}$-modules on $Y_\etale$
(Cohomology on Sites, Lemma
\ref{sites-cohomology-lemma-pushforward-injective-flat}).
The result follows formally from this, but we will also
spell it out.

\medskip\noindent
Let $Z' \to Z$ be a morphism of schemes and set
$Y' = Z' \times_Z Y$ and $X' = Z' \times_Z X = Y' \times_ Y X$.
Denote $a : X' \to X$, $b : Y' \to Y$, and $c : Z' \to Z$ the
projections. Similarly for $f' : X' \to Y'$ and $g' : Y' \to Z'$.
By Lemma \ref{lemma-proper-base-change-f-star} we have
$b^{-1}f_*\mathcal{I} = f'_*a^{-1}\mathcal{I}$.
On the other hand, we know that $R^qf'_*a^{-1}\mathcal{I}$ and
$R^q(g')_*b^{-1}f_*\mathcal{I}$ are zero for $q > 0$.
Using the spectral sequence
(Cohomology on Sites, Lemma \ref{sites-cohomology-lemma-relative-Leray})
$$
R^pg'_* R^qf'_* a^{-1}\mathcal{I} \Rightarrow
R^{p + q}(g' \circ f')_* a^{-1}\mathcal{I}
$$
we conclude that $R^p(g' \circ f')_*a^{-1}\mathcal{I} = 0$ for
$p > 0$ as desired.
\end{proof}

\begin{lemma}
\label{lemma-finite}
\begin{slogan}
Proper base change for \'etale cohomology holds for finite morphisms.
\end{slogan}
Let $f : X \to Y$ be a finite morphism of schemes.
Then cohomology commutes with base change for $f$.
\end{lemma}

\begin{proof}
Observe that a finite morphism is proper, see
Morphisms, Lemma \ref{morphisms-lemma-finite-proper}.
Moreover, the base change of a finite morphism is finite, see
Morphisms, Lemma \ref{morphisms-lemma-base-change-finite}.
Thus the result follows from
Lemma \ref{lemma-proper-base-change-in-terms-of-injectives}
combined with
Proposition \ref{proposition-finite-higher-direct-image-zero}.
\end{proof}

\begin{lemma}
\label{lemma-reduce-to-P1}
To prove that cohomology commutes with base change for
every proper morphism of schemes it suffices to prove it
holds for the morphism $\mathbf{P}^1_S \to S$ for every scheme $S$.
\end{lemma}

\begin{proof}
Let $f : X \to Y$ be a proper morphism of schemes.
Let $Y = \bigcup Y_i$ be an affine open covering
and set $X_i = f^{-1}(Y_i)$. If we can prove
cohomology commutes with base change for $X_i \to Y_i$,
then cohomology commutes with base change for $f$.
Namely, the formation of the higher direct images
commutes with Zariski (and even \'etale) localization
on the base, see
Lemma \ref{lemma-higher-direct-images}.
Thus we may assume $Y$ is affine.

\medskip\noindent
Let $Y$ be an affine scheme and let $X \to Y$ be a proper morphism.
By Chow's lemma there exists a commutative diagram
$$
\xymatrix{
X \ar[rd] & X' \ar[d] \ar[l]^\pi \ar[r] & \mathbf{P}^n_Y \ar[dl] \\
& Y &
}
$$
where $X' \to \mathbf{P}^n_Y$ is an immersion, and
$\pi : X' \to X$ is proper and surjective, see
Limits, Lemma \ref{limits-lemma-chow-finite-type}.
Since $X \to Y$ is proper, we find that $X' \to Y$ is proper
(Morphisms, Lemma \ref{morphisms-lemma-composition-proper}).
Hence $X' \to \mathbf{P}^n_Y$ is a closed immersion
(Morphisms, Lemma \ref{morphisms-lemma-image-proper-scheme-closed}).
It follows that $X' \to X \times_Y \mathbf{P}^n_Y = \mathbf{P}^n_X$
is a closed immersion (as an immersion with closed image).

\medskip\noindent
By Lemma \ref{lemma-sandwich}
it suffices to prove cohomology commutes with base change for
$\pi$ and $X' \to Y$. These morphisms both factor as a closed
immersion followed by a projection $\mathbf{P}^n_S \to S$ (for some $S$).
By Lemma \ref{lemma-finite} the result holds for closed
immersions (as closed immersions are finite).
By Lemma \ref{lemma-composition} it suffices to prove the
result for projections $\mathbf{P}^n_S \to S$.

\medskip\noindent
For every $n \geq 1$ there is a finite surjective morphism
$$
\mathbf{P}^1_S \times_S \ldots \times_S \mathbf{P}^1_S
\longrightarrow
\mathbf{P}^n_S
$$
given on coordinates by
$$
((x_1 : y_1), (x_2 : y_2), \ldots, (x_n : y_n))
\longmapsto
(F_0 : \ldots : F_n)
$$
where $F_0, \ldots, F_n$ in $x_1, \ldots, y_n$
are the polynomials with integer coefficients such that
$$
\prod (x_i t + y_i) = F_0 t^n + F_1 t^{n - 1} + \ldots + F_n
$$
Applying
Lemmas \ref{lemma-sandwich}, \ref{lemma-finite}, and \ref{lemma-composition}
one more time we conclude that the lemma is true.
\end{proof}

\begin{theorem}[Proper Base Change]
\label{theorem-proper-base-change}
Let $f : X \to Y$ be a proper morphism of schemes. Let $g : Y' \to Y$ be
a morphism of schemes. Set $X' = Y' \times_Y X$
and consider the cartesian diagram
$$
\xymatrix{
X' \ar[r]_{g'} \ar[d]_{f'} & X \ar[d]^f \\
Y' \ar[r]^g & Y
}
$$
Let $\mathcal{F}$ be an abelian torsion sheaf on $X_\etale$.
Then the base change map
$$
g^{-1}Rf_*\mathcal{F} \longrightarrow Rf'_*(g')^{-1}\mathcal{F}
$$
is an isomorphism.
\end{theorem}

\begin{proof}
In the terminology introduced above, this means that cohomology commutes
with base change for every proper morphism of schemes. By
Lemma \ref{lemma-reduce-to-P1}
it suffices to prove that cohomology commutes with base change
for the morphism $\mathbf{P}^1_S \to S$ for every scheme $S$.

\medskip\noindent
Let $S$ be the spectrum of a strictly henselian local ring with closed
point $s$. Set $X = \mathbf{P}^1_S$ and $X_0 = X_s = \mathbf{P}^1_s$.
Let $\mathcal{F}$ be a sheaf of $\mathbf{Z}/\ell\mathbf{Z}$-modules
on $X_\etale$. The key to our proof is that
$$
H^q_\etale(X, \mathcal{F}) = H^q_\etale(X_0, \mathcal{F}|_{X_0}).
$$
Namely, choose a resolution $\mathcal{F} \to \mathcal{I}^\bullet$
by injective sheaves of $\mathbf{Z}/\ell\mathbf{Z}$-modules.
Then $\mathcal{I}^\bullet|_{X_0}$ is a resolution of $\mathcal{F}|_{X_0}$
by right $H^0_\etale(X_0, -)$-acyclic objects, see
Lemma \ref{lemma-efface-cohomology-on-fibre-by-finite-cover}.
Leray's acyclicity lemma tells us the right hand side is computed by
the complex $H^0_\etale(X_0, \mathcal{I}^\bullet|_{X_0})$
which is equal to $H^0_\etale(X, \mathcal{I}^\bullet)$ by
Lemma \ref{lemma-h0-proper-over-henselian-local}. This complex
computes the left hand side.

\medskip\noindent
Assume $S$ is general and $\mathcal{F}$ is a sheaf of
$\mathbf{Z}/\ell\mathbf{Z}$-modules on $X_\etale$.
Let $\overline{s} : \Spec(k) \to S$ be a geometric point
of $S$ lying over $s \in S$. We have
$$
(R^qf_*\mathcal{F})_{\overline{s}} =
H^q_\etale(\mathbf{P}^1_{\mathcal{O}_{S, \overline{s}}^{sh}},
\mathcal{F}|_{\mathbf{P}^1_{\mathcal{O}_{S, \overline{s}}^{sh}}}) =
H^q_\etale(\mathbf{P}^1_{\kappa(s)^{sep}},
\mathcal{F}|_{\mathbf{P}^1_{\kappa(s)^{sep}}})
$$
where $\kappa(s)^{sep}$ is the residue field of
$\mathcal{O}_{S, \overline{s}}^{sh}$, i.e., the separable algebraic
closure of $\kappa(s)$ in $k$.
The first equality by Theorem \ref{theorem-higher-direct-images}
and the second equality by the displayed formula in the
previous paragraph.

\medskip\noindent
Finally, consider any morphism of schemes $g : T \to S$ where
$S$ and $\mathcal{F}$ are as above.
Set $f' : \mathbf{P}^1_T \to T$ the projection and let
$g' : \mathbf{P}^1_T \to \mathbf{P}^1_S$ the morphism induced
by $g$. Consider the base change map
$$
g^{-1}R^qf_*\mathcal{F}
\longrightarrow
R^qf'_*(g')^{-1}\mathcal{F}
$$
Let $\overline{t}$ be a geometric point of $T$ with image
$\overline{s} = g(\overline{t})$. By our discussion
above the map on stalks at $\overline{t}$ is the map
$$
H^q_\etale(\mathbf{P}^1_{\kappa(s)^{sep}},
\mathcal{F}|_{\mathbf{P}^1_{\kappa(s)^{sep}}})
\longrightarrow
H^q_\etale(\mathbf{P}^1_{\kappa(t)^{sep}},
\mathcal{F}|_{\mathbf{P}^1_{\kappa(t)^{sep}}})
$$
Since $\kappa(s)^{sep} \subset \kappa(t)^{sep}$ this map is an
isomorphism by Lemma \ref{lemma-base-change-dim-1-separably-closed}.

\medskip\noindent
This proves cohomology commutes with base change for
$\mathbf{P}^1_S \to S$ and sheaves of $\mathbf{Z}/\ell\mathbf{Z}$-modules.
In particular, for an injective sheaf of $\mathbf{Z}/\ell\mathbf{Z}$-modules
the higher direct images of any base change are zero.
In other words, condition (2) of
Lemma \ref{lemma-proper-base-change-in-terms-of-injectives}
holds and the proof is complete.
\end{proof}

\begin{lemma}
\label{lemma-proper-base-change}
Let $f : X \to Y$ be a proper morphism of schemes. Let $g : Y' \to Y$ be
a morphism of schemes. Set $X' = Y' \times_Y X$ and denote
$f' : X' \to Y'$ and $g' : X' \to X$ the projections.
Let $E \in D^+(X_\etale)$ have torsion cohomology sheaves.
Then the base change map (\ref{equation-base-change})
$g^{-1}Rf_*E \to Rf'_*(g')^{-1}E$
is an isomorphism.
\end{lemma}

\begin{proof}
This is a simple consequence of the proper base change theorem
(Theorem \ref{theorem-proper-base-change}) using the spectral
sequences
$$
E_2^{p, q} = R^pf_*H^q(E)
\quad\text{and}\quad
{E'}_2^{p, q} = R^pf'_*(g')^{-1}H^q(E)
$$
converging to $R^nf_*E$ and $R^nf'_*(g')^{-1}E$.
The spectral sequences are constructed in
Derived Categories, Lemma \ref{derived-lemma-two-ss-complex-functor}.
Some details omitted.
\end{proof}

\begin{lemma}
\label{lemma-proper-base-change-stalk}
Let $f : X \to Y$ be a proper morphism of schemes. Let $\overline{y} \to Y$
be a geometric point.
\begin{enumerate}
\item For a torsion abelian sheaf $\mathcal{F}$ on $X_\etale$ we have
$(R^nf_*\mathcal{F})_{\overline{y}} =
H^n_\etale(X_{\overline{y}}, \mathcal{F}_{\overline{y}})$.
\item For $E \in D^+(X_\etale)$ with torsion cohomology sheaves we have
$(R^nf_*E)_{\overline{y}} =
H^n_\etale(X_{\overline{y}}, E|_{X_{\overline{y}}})$.
\end{enumerate}
\end{lemma}

\begin{proof}
In the statement, $\mathcal{F}_{\overline{y}}$ denotes the pullback
of $\mathcal{F}$ to the scheme theoretic fibre
$X_{\overline{y}} = \overline{y} \times_Y X$.
Since pulling back by $\overline{y} \to Y$ produces the
stalk of $\mathcal{F}$, the first statement of the lemma
is a special case of Theorem \ref{theorem-proper-base-change}.
The second one is a special case of Lemma \ref{lemma-proper-base-change}.
\end{proof}






\section{Applications of proper base change}
\label{section-applications-proper-base-change}

\noindent
In this section we discuss some more or less immediate
consequences of the proper base change theorem.

\begin{lemma}
\label{lemma-base-change-separably-closed}
Let $K/k$ be an extension of separably closed fields.
Let $X$ be a proper scheme over $k$.
Let $\mathcal{F}$ be a torsion abelian sheaf on $X_\etale$.
Then the map $H^q_\etale(X, \mathcal{F}) \to
H^q_\etale(X_K, \mathcal{F}|_{X_K})$ is an isomorphism
for $q \geq 0$.
\end{lemma}

\begin{proof}
Looking at stalks we see that
this is a special case of
Theorem \ref{theorem-proper-base-change}.
\end{proof}

\begin{lemma}
\label{lemma-cohomological-dimension-proper}
Let $f : X \to Y$ be a proper morphism of schemes
all of whose fibres have dimension $\leq n$.
Then for any abelian torsion sheaf $\mathcal{F}$ on $X_\etale$
we have $R^qf_*\mathcal{F} = 0$ for $q > 2n$.
\end{lemma}

\begin{proof}
We will prove this by induction on $n$ for all proper morphisms.

\medskip\noindent
If $n = 0$, then $f$ is a finite morphism
(More on Morphisms, Lemma \ref{more-morphisms-lemma-characterize-finite})
and the
result is true by Proposition \ref{proposition-finite-higher-direct-image-zero}.

\medskip\noindent
If $n > 0$, then using Lemma \ref{lemma-proper-base-change-stalk}
we see that it suffices to prove $H^i_\etale(X, \mathcal{F}) = 0$
for $i > 2n$ and $X$ a proper scheme, $\dim(X) \leq n$
over an algebraically closed field $k$
and $\mathcal{F}$ is a torsion abelian sheaf on $X$.

\medskip\noindent
If $n = 1$ this follows from Theorem \ref{theorem-vanishing-curves}.
Assume $n > 1$. By Proposition \ref{proposition-topological-invariance}
we may replace $X$ by its reduction.
Let $\nu : X^\nu \to X$ be the normalization.
This is a surjective birational finite morphism
(see Varieties, Lemma \ref{varieties-lemma-normalization-locally-algebraic})
and hence an isomorphism over a dense open $U \subset X$
(Morphisms, Lemma \ref{morphisms-lemma-birational-birational}).
Then we see that
$c : \mathcal{F} \to \nu_*\nu^{-1}\mathcal{F}$ is injective
(as $\nu$ is surjective) and an isomorphism over $U$.
Denote $i : Z \to X$ the inclusion of the complement of $U$.
Since $U$ is dense in $X$ we have $\dim(Z) < \dim(X) = n$.
By Proposition \ref{proposition-closed-immersion-pushforward} have
$\Coker(c) = i_*\mathcal{G}$ for some
abelian torsion sheaf $\mathcal{G}$ on $Z_\etale$.
Then $H^q_\etale(X, \Coker(c)) = H^q_\etale(Z, \mathcal{G})$
(by Proposition \ref{proposition-finite-higher-direct-image-zero}
and the Leray spectral sequence) and by induction hypothesis we conclude that
the cokernel of $c$ has cohomology in degrees $\leq 2(n - 1)$.
Thus it suffices to prove the result for $\nu_*\nu^{-1}\mathcal{F}$.
As $\nu$ is finite this reduces us to showing
that $H^i_\etale(X^\nu, \nu^{-1}\mathcal{F})$ is
zero for $i > 2n$. This case is treated in the next paragraph.

\medskip\noindent
Assume $X$ is integral normal proper scheme over $k$ of dimension $n$.
Choose a nonconstant rational function $f$ on $X$. The graph
$X' \subset X \times \mathbf{P}^1_k$ of $f$ sits into a diagram
$$
X \xleftarrow{b} X' \xrightarrow{f} \mathbf{P}^1_k
$$
Observe that $b$ is an isomorphism over an open subscheme
$U \subset X$ whose complement is a closed subscheme
$Z \subset X$ of codimension $\geq 2$. Namely, $U$ is the
domain of definition of $f$ which contains all codimension $1$
points of $X$, see
Morphisms, Lemmas \ref{morphisms-lemma-rational-map-from-reduced-to-separated}
and \ref{morphisms-lemma-extend-across}
(combined with Serre's criterion for normality, see
Properties, Lemma \ref{properties-lemma-criterion-normal}).
Moreover the fibres of $b$ have dimension $\leq 1$ (as closed subschemes
of $\mathbf{P}^1$). Hence $R^ib_*b^{-1}\mathcal{F}$ is nonzero only
if $i \in \{0, 1, 2\}$ by induction. Choose a distinguished triangle
$$
\mathcal{F} \to Rb_*b^{-1}\mathcal{F} \to Q \to \mathcal{F}[1]
$$
Using that $\mathcal{F} \to b_*b^{-1}\mathcal{F}$ is injective
as before and using what we just said, we see that $Q$ has nonzero
cohomology sheaves only in degrees $0, 1, 2$ sitting on $Z$.
Moreover, these cohomology sheaves are torsion by
Lemma \ref{lemma-torsion-direct-image}.
By induction and the spectral sequence of
Derived Categories, Lemma \ref{derived-lemma-two-ss-complex-functor}
we see that $H^i(X, Q)$ is zero for
$i > 2 + 2\dim(Z) \leq 2 + 2(n - 2) = 2n - 2$. Thus it suffices
to prove that $H^i(X', b^{-1}\mathcal{F}) = 0$ for
$i > 2n$. At this point we use the morphism
$$
f : X' \to \mathbf{P}^1_k
$$
whose fibres have dimension $< n$. Hence by induction we see that
$R^if_*b^{-1}\mathcal{F} = 0$ for $i > 2(n - 1)$.
We conclude by the Leray spectral sequence
$$
H^i(\mathbf{P}^1_k, R^jf_*b^{-1}\mathcal{F})
\Rightarrow
H^{i + j}(X', b^{-1}\mathcal{F})
$$
and the fact that $\dim(\mathbf{P}^1_k) = 1$.
\end{proof}

\noindent
When working with mod $n$ coefficients we can do proper
base change for unbounded complexes.

\begin{lemma}
\label{lemma-proper-base-change-mod-n}
Let $f : X \to Y$ be a proper morphism of schemes. Let $g : Y' \to Y$ be
a morphism of schemes. Set $X' = Y' \times_Y X$ and denote
$f' : X' \to Y'$ and $g' : X' \to X$ the projections.
Let $n \geq 1$ be an integer.
Let $E \in D(X_\etale, \mathbf{Z}/n\mathbf{Z})$.
Then the base change map (\ref{equation-base-change})
$g^{-1}Rf_*E \to Rf'_*(g')^{-1}E$
is an isomorphism.
\end{lemma}

\begin{proof}
It is enough to prove this when $Y$ and $Y'$ are quasi-compact.
By Morphisms, Lemma
\ref{morphisms-lemma-morphism-finite-type-bounded-dimension}
we see that the dimension of the fibres of
$f : X \to Y$ and $f' : X' \to Y'$ are bounded. Thus
Lemma \ref{lemma-cohomological-dimension-proper} implies that
$$
f_* :
\textit{Mod}(X_\etale, \mathbf{Z}/n\mathbf{Z})
\longrightarrow
\textit{Mod}(Y_\etale, \mathbf{Z}/n\mathbf{Z})
$$
and
$$
f'_* :
\textit{Mod}(X'_\etale, \mathbf{Z}/n\mathbf{Z})
\longrightarrow
\textit{Mod}(Y'_\etale, \mathbf{Z}/n\mathbf{Z})
$$
have finite cohomological dimension in the sense of
Derived Categories, Lemma \ref{derived-lemma-unbounded-right-derived}.
Choose a K-injective complex $\mathcal{I}^\bullet$
of $\mathbf{Z}/n\mathbf{Z}$-modules each of whose
terms $\mathcal{I}^n$ is an injective sheaf of
$\mathbf{Z}/n\mathbf{Z}$-modules representing $E$.
See Injectives, Theorem
\ref{injectives-theorem-K-injective-embedding-grothendieck}.
By the usual proper base change theorem we find
that $R^qf'_*(g')^{-1}\mathcal{I}^n = 0$ for $q > 0$, see
Theorem \ref{theorem-proper-base-change}.
Hence we conclude by 
Derived Categories, Lemma \ref{derived-lemma-unbounded-right-derived}
that we may compute $Rf'_*(g')^{-1}E$ by the complex
$f'_*(g')^{-1}\mathcal{I}^\bullet$. Another application
of the usual proper base change theorem shows that
this is equal to $g^{-1}f_*\mathcal{I}^\bullet$ as desired.
\end{proof}

\begin{lemma}
\label{lemma-pull-out-constant}
Let $X$ be a quasi-compact and quasi-separated scheme.
Let $E \in D^+(X_\etale)$ and $K \in D^+(\mathbf{Z})$.
Then
$$
R\Gamma(X, E \otimes_\mathbf{Z}^\mathbf{L} \underline{K}) =
R\Gamma(X, E) \otimes_\mathbf{Z}^\mathbf{L} K
$$
\end{lemma}

\begin{proof}
Say $H^i(E) = 0$ for $i \geq a$ and $H^j(K) = 0$ for $j \geq b$.
We may represent $K$ by a bounded below complex
$K^\bullet$ of torsion free $\mathbf{Z}$-modules.
(Choose a K-flat complex $L^\bullet$ representing $K$ and then
take $K^\bullet = \tau_{\geq b - 1}L^\bullet$. This works because
$\mathbf{Z}$ has global dimension $1$. See
More on Algebra, Lemma \ref{more-algebra-lemma-last-one-flat}.)
We may represent $E$ by a bounded below complex
$\mathcal{E}^\bullet$.
Then $E \otimes_\mathbf{Z}^\mathbf{L} \underline{K}$ is
represented by
$$
\text{Tot}(\mathcal{E}^\bullet \otimes_\mathbf{Z} \underline{K}^\bullet)
$$
Using distinguished triangles
$$
\sigma_{\geq -b + n + 1}K^\bullet
\to K^\bullet \to
\sigma_{\leq -b + n}K^\bullet
$$
and the trivial vanishing
$$
H^n(X,
\text{Tot}(\mathcal{E}^\bullet \otimes_\mathbf{Z}
\sigma_{\geq -a + n + 1}\underline{K}^\bullet) = 0
$$
and
$$
H^n(R\Gamma(X, E)
\otimes_\mathbf{Z}^\mathbf{L} \sigma_{\geq -a + n + 1}K^\bullet) = 0
$$
we reduce to the case where $K^\bullet$ is a bounded complex of
flat $\mathbf{Z}$-modules. Repeating the argument we reduce to the
case where $K^\bullet$ is equal to a single flat $\mathbf{Z}$-module
sitting in some degree. Next, using the stupid
trunctions for $\mathcal{E}^\bullet$
we reduce in exactly the same manner to the case where
$\mathcal{E}^\bullet$ is a single abelian sheaf sitting
in some degree. Thus it suffices to show that
$$
H^n(X, \mathcal{E} \otimes_\mathbf{Z} \underline{M}) =
H^n(X, \mathcal{E}) \otimes_\mathbf{Z} M
$$
when $M$ is a flat $\mathbf{Z}$-module and $\mathcal{E}$
is an abelian sheaf on $X$. In this case we
write $M$ is a filtered colimit of finite free $\mathbf{Z}$-modules
(Lazard's theorem, see Algebra, Theorem \ref{algebra-theorem-lazard}).
By Theorem \ref{theorem-colimit} this reduces us to the case of finite free
$\mathbf{Z}$-module $M$ in which case the result is trivially true.
\end{proof}

\begin{lemma}
\label{lemma-projection-formula-proper}
Let $f : X \to Y$ be a proper morphism of schemes.
Let $E \in D^+(X_\etale)$ have torsion cohomology sheaves.
Let $K \in D^+(Y_\etale)$. Then
$$
Rf_*E \otimes_\mathbf{Z}^\mathbf{L} K =
Rf_*(E \otimes_\mathbf{Z}^\mathbf{L} f^{-1}K)
$$
in $D^+(Y_\etale)$.
\end{lemma}

\begin{proof}
There is a canonical map from left to right by
Cohomology on Sites, Section \ref{sites-cohomology-section-projection-formula}.
We will check the equality on stalks.
Recall that computing derived tensor products commutes with pullbacks.
See Cohomology on Sites, Lemma
\ref{sites-cohomology-lemma-pullback-tensor-product}.
Thus we have
$$
(E \otimes_\mathbf{Z}^\mathbf{L} f^{-1}K)_{\overline{x}}
=
E_{\overline{x}} \otimes_\mathbf{Z}^\mathbf{L} K_{\overline{y}}
$$
where $\overline{y}$ is the image of $\overline{x}$ in $Y$.
Since $\mathbf{Z}$ has global dimension $1$ we see that
this complex has vanishing cohomology in degree $< - 1 + a + b$
if $H^i(E) = 0$ for $i \geq a$ and $H^j(K) = 0$ for $j \geq b$.
Moreover, since $H^i(E)$ is a torsion abelian
sheaf for each $i$, the same is true for the cohomology
sheaves of the complex $E \otimes_\mathbf{Z}^\mathbf{L} K$.
Namely, we have
$$
(E \otimes_\mathbf{Z}^\mathbf{L} f^{-1}K)
\otimes_{\mathbf{Z}}^\mathbf{L} \mathbf{Q} =
(E \otimes_\mathbf{Z}^\mathbf{L} \mathbf{Q})
\otimes_{\mathbf{Q}}^\mathbf{L}
(f^{-1}K \otimes_{\mathbf{Z}}^\mathbf{L} \mathbf{Q})
$$
which is zero in the derived category.
In this way we see that Lemma \ref{lemma-proper-base-change-stalk}
applies to both sides to see that it suffices to show
$$
R\Gamma(X_{\overline{y}},
E|_{X_{\overline{y}}} \otimes_\mathbf{Z}^\mathbf{L}
(X_{\overline{y}} \to \overline{y})^{-1}K_{\overline{y}}) =
R\Gamma(X_{\overline{y}},
E|_{X_{\overline{y}}}) \otimes_\mathbf{Z}^\mathbf{L} K_{\overline{y}}
$$
This is shown in Lemma \ref{lemma-pull-out-constant}.
\end{proof}








\section{Local acyclicity}
\label{section-local-acyclicity}

\noindent
In this section we deduce local acyclicity of smooth morphisms
from the smooth base change theorem. In SGA 4 or SGA 4.5 the authors
first prove a version of local acyclicity for smooth morphisms
and then deduce the smooth base change theorem.

\medskip\noindent
We will use the formulation of local acyclicity given by
Deligne \cite[Definition 2.12, page 242]{SGA4.5}.
Let $f : X \to S$ be a morphism of schemes.
Let $\overline{x}$ be a geometric point of $X$ with image
$\overline{s} = f(\overline{x})$ in $S$.
Let $\overline{t}$ be a geometric point of
$\Spec(\mathcal{O}^{sh}_{S, \overline{s}})$.
We obtain a commutative diagram
$$
\xymatrix{
F_{\overline{x}, \overline{t}} =
\overline{t} \times_{\Spec(\mathcal{O}^{sh}_{S, \overline{s}})}
\Spec(\mathcal{O}^{sh}_{X, \overline{x}}) \ar[r] \ar[d] &
\Spec(\mathcal{O}^{sh}_{X, \overline{x}}) \ar[r] \ar[d] &
X \ar[d] \\
\overline{t} \ar[r] &
\Spec(\mathcal{O}^{sh}_{S, \overline{s}}) \ar[r] &
S
}
$$
The scheme $F_{\overline{x}, \overline{t}}$ is called
a {\it variety of vanishing cycles of $f$ at $\overline{x}$}.
Let $K$ be an object of $D(X_\etale)$. For any morphism of schemes
$g : Y\to X$ we write $R\Gamma(Y, K)$ instead of
$R\Gamma(Y_\etale, g_{small}^{-1}K)$. Since
$\mathcal{O}^{sh}_{X, \overline{x}}$ is strictly henselian
we have
$K_{\overline{x}} = R\Gamma(\Spec(\mathcal{O}^{sh}_{X, \overline{x}}), K)$.
Thus we obtain a canonical map
\begin{equation}
\label{equation-alpha-K}
\alpha_{K, \overline{x}, \overline{t}} :
K_{\overline{x}}
\longrightarrow
R\Gamma(F_{\overline{x}, \overline{t}}, K)
\end{equation}
by pulling back cohomology along
$F_{\overline{x}, \overline{t}} \to \Spec(\mathcal{O}^{sh}_{X, \overline{x}})$.

\begin{definition}
\label{definition-locally-acyclic}
\begin{reference}
\cite[Definition 2.12, page 242]{SGA4.5} and
\cite[Definition (1.3), page 54]{SGA4.5}
\end{reference}
Let $f : X \to S$ be a morphism of schemes.
Let $K$ be an object of $D(X_\etale)$.
\begin{enumerate}
\item
Let $\overline{x}$ be a geometric point of $X$ with image
$\overline{s} = f(\overline{x})$.
We say $f$ is {\it locally acyclic at $\overline{x}$ relative to $K$}
if for every geometric point $\overline{t}$ of
$\Spec(\mathcal{O}^{sh}_{S, \overline{s}})$ the
map (\ref{equation-alpha-K}) is an isomorphism\footnote{We do not
assume $\overline{t}$ is an algebraic geometric point of
$\Spec(\mathcal{O}^{sh}_{S, \overline{s}})$. Often using
Lemma \ref{lemma-smooth-base-change-separably-closed}
one may reduce to this case.}.
\item We say $f$ is {\it locally acyclic relative to $K$}
if $f$ is locally acyclic at $\overline{x}$ relative to $K$
for every geometric point $\overline{x}$ of $X$.
\item We say $f$ is {\it universally locally acyclic relative to $K$}
if for any morphism $S' \to S$ of schemes the base change $f' : X' \to S'$
is locally acyclic relative to the pullback of $K$ to $X'$.
\item We say $f$ is {\it locally acyclic} if for all geometric
points $\overline{x}$ of $X$ and any integer $n$ prime to the characteristic
of $\kappa(\overline{x})$, the morphism $f$ is locally acyclic
at $\overline{x}$ relative to the constant sheaf with value
$\mathbf{Z}/n\mathbf{Z}$.
\item We say $f$ is {\it universally locally acyclic} if
for any morphism $S' \to S$ of schemes the base change $f' : X' \to S'$
is locally acyclic.
\end{enumerate}
\end{definition}

\noindent
Let $M$ be an abelian group. Then local acyclicity of $f : X \to S$
with respect to the constant sheaf $\underline{M}$ boils down
to the requirement that
$$
H^q(F_{\overline{x}, \overline{t}}, \underline{M}) =
\left\{
\begin{matrix}
M & \text{if} & q = 0 \\
0 & \text{if} & q \not = 0
\end{matrix}
\right.
$$
for any geometric point $\overline{x}$ of $X$ and any geometric
point $\overline{t}$ of $\Spec(\mathcal{O}^{sh}_{S, f(\overline{x})})$.
In this way we see that being locally acyclic corresponds to the
vanishing of the higher cohomology groups of the geometric fibres
$F_{\overline{x}, \overline{t}}$ of the maps between the
strict henselizations at $\overline{x}$ and $\overline{s}$.

\begin{proposition}
\label{proposition-smooth-locally-acyclic}
Let $f : X \to S$ be a smooth morphism of schemes.
Then $f$ is universally locally acyclic.
\end{proposition}

\begin{proof}
Since the base change of a smooth morphism is smooth, it suffices
to show that smooth morphisms are locally acyclic.
Let $\overline{x}$ be a geometric point of $X$ with image
$\overline{s} = f(\overline{x})$. Let
$\overline{t}$ be a geometric point of
$\Spec(\mathcal{O}^{sh}_{S, f(\overline{x})})$.
Since we are trying to prove a property of the ring map
$\mathcal{O}^{sh}_{S, \overline{s}} \to \mathcal{O}^{sh}_{X, \overline{x}}$
(see discussion following Definition \ref{definition-locally-acyclic})
we may and do replace $f : X \to S$ by the base change
$X \times_S \Spec(\mathcal{O}^{sh}_{S, \overline{s}}) \to
\Spec(\mathcal{O}^{sh}_{S, \overline{s}})$.
Thus we may and do assume that $S$ is the spectrum of a strictly
henselian local ring and that $\overline{s}$ lies over
the closed point of $S$.

\medskip\noindent
We will apply Lemma \ref{lemma-base-change-f-star-general-stalks}
to the diagram
$$
\xymatrix{
X \ar[d]_f & X_{\overline{t}} \ar[l]^h \ar[d]^e \\
S & \overline{t} \ar[l]_g
}
$$
and the sheaf $\mathcal{F} = \underline{M}$ where $M = \mathbf{Z}/n\mathbf{Z}$
for some integer $n$ prime to the characteristic of the residue field of
$\overline{x}$. We know that the map
$f^{-1}R^qg_*\mathcal{F} \to R^qh_*e^{-1}\mathcal{F}$
is an isomorphism by smooth base change, see
Theorem \ref{theorem-smooth-base-change} (the assumption
on torsion holds by our choice of $n$).
Thus Lemma \ref{lemma-base-change-f-star-general-stalks}
gives us the middle equality in
$$
H^q(F_{\overline{x}, \overline{t}}, \underline{M}) =
H^q(\Spec(\mathcal{O}^{sh}_{X, \overline{x}}) \times_S \overline{t},
\underline{M})
=
H^q(\Spec(\mathcal{O}^{sh}_{S, \overline{s}}) \times_S \overline{t},
\underline{M}) =
H^q(\overline{t}, \underline{M})
$$
For the outer two equalities we use that
$S = \Spec(\mathcal{O}^{sh}_{S, \overline{s}})$.
Since $\overline{t}$ is the spectrum of a separably
closed field we conclude that
$$
H^q(F_{\overline{x}, \overline{t}}, \underline{M}) =
\left\{
\begin{matrix}
M & \text{if} & q = 0 \\
0 & \text{if} & q \not = 0
\end{matrix}
\right.
$$
which is what we had to show (see discussion following
Definition \ref{definition-locally-acyclic}).
\end{proof}

\begin{lemma}
\label{lemma-locally-acyclic-locally-constant}
Let $f : X \to S$ be a morphism of schemes. Let $\mathcal{F}$ be a
locally constant abelian sheaf on $X_\etale$ such that for every geometric
point $\overline{x}$ of $X$ the abelian group
$\mathcal{F}_{\overline{x}}$ is a torsion group all of whose elements have
order prime to the characteristic of the
residue field of $\overline{x}$. If $f$ is locally acyclic, then $f$ is
locally acyclic relative to $\mathcal{F}$.
\end{lemma}

\begin{proof}
Namely, let $\overline{x}$ be a geometric point of $X$.
Since $\mathcal{F}$ is locally constant we see that
the restriction of $\mathcal{F}$ to $\Spec(\mathcal{O}^{sh}_{X, \overline{x}})$
is isomorphic to the constant sheaf $\underline{M}$ with
$M = \mathcal{F}_{\overline{x}}$. By assumption we can write
$M = \colim M_i$ as a filtered colimit of finite abelian groups
$M_i$ of order prime to the characteristic of the residue field of
$\overline{x}$. Consider a geometric point $\overline{t}$ of
$\Spec(\mathcal{O}^{sh}_{S, f(\overline{x})})$.
Since $F_{\overline{x}, \overline{t}}$ is affine, we have
$$
H^q(F_{\overline{x}, \overline{t}}, \underline{M}) =
\colim H^q(F_{\overline{x}, \overline{t}}, \underline{M_i})
$$
by Lemma \ref{lemma-colimit}.
For each $i$ we can write $M_i = \bigoplus \mathbf{Z}/n_{i, j}\mathbf{Z}$
as a finite direct sum for some integers $n_{i, j}$ prime to the
characteristic of the residue field of $\overline{x}$.
Since $f$ is locally acyclic we see that
$$
H^q(F_{\overline{x}, \overline{t}},
\underline{\mathbf{Z}/n_{i, j}\mathbf{Z}}) =
\left\{
\begin{matrix}
\mathbf{Z}/n_{i, j}\mathbf{Z} & \text{if} & q = 0 \\
0 & \text{if} & q \not = 0
\end{matrix}
\right.
$$
See discussion following Definition \ref{definition-locally-acyclic}.
Taking the direct sums and the colimit we conclude that
$$
H^q(F_{\overline{x}, \overline{t}}, \underline{M}) =
\left\{
\begin{matrix}
M & \text{if} & q = 0 \\
0 & \text{if} & q \not = 0
\end{matrix}
\right.
$$
and we win.
\end{proof}

\begin{lemma}
\label{lemma-locally-acyclic-quasi-finite-base-change}
Let
$$
\xymatrix{
X' \ar[r]_{g'} \ar[d]_{f'} & X \ar[d]^f \\
S' \ar[r]^g & S
}
$$
be a cartesian diagram of schemes. Let $K$ be an object of $D(X_\etale)$.
Let $\overline{x}'$ be a geometric point of $X'$ with image $\overline{x}$
in $X$. If
\begin{enumerate}
\item $f$ is locally acyclic at $\overline{x}$ relative to $K$ and
\item $g$ is locally quasi-finite, or $S' = \lim S_i$
is a directed inverse limit of schemes locally quasi-finite over $S$
with affine transition morphisms, or $g : S' \to S$ is integral,
\end{enumerate}
then $f'$ locally acyclic at $\overline{x}'$ relative to $(g')^{-1}K$.
\end{lemma}

\begin{proof}
Denote $\overline{s}'$ and $\overline{s}$ the images of
$\overline{x}'$ and $\overline{x}$ in $S'$ and $S$.
Let $\overline{t}'$ be a geometric point of the spectrum of
$\Spec(\mathcal{O}^{sh}_{S', \overline{s}'})$ and denote
$\overline{t}$ the image in $\Spec(\mathcal{O}^{sh}_{S, \overline{s}})$.
By Algebra, Lemma
\ref{algebra-lemma-base-change-strict-henselization-quasi-finite}
and our assumptions on $g$ we have
$$
\mathcal{O}^{sh}_{X, \overline{x}}
\otimes_{\mathcal{O}^{sh}_{S, \overline{s}}}
\mathcal{O}^{sh}_{S', \overline{s}'}
\longrightarrow
\mathcal{O}^{sh}_{X', \overline{x}'}
$$
is an isomorphism. Since by our conventions
$\kappa(\overline{t}) = \kappa(\overline{t}')$
we conclude that
$$
F_{\overline{x}', \overline{t}'} =
\Spec\left(
\mathcal{O}^{sh}_{X', \overline{x}'}
\otimes_{\mathcal{O}^{sh}_{S', \overline{s}'}}
\kappa(\overline{t}')\right) =
\Spec\left(
\mathcal{O}^{sh}_{X, \overline{x}}
\otimes_{\mathcal{O}^{sh}_{S, \overline{s}}}
\kappa(\overline{t})\right) =
F_{\overline{x}, \overline{t}}
$$
In other words, the varieties of vanishing cycles
of $f'$ at $\overline{x}'$ are examples of varieties of vanishing
cycles of $f$ at $\overline{x}$. The lemma follows
immediately from this and the definitions.
\end{proof}










\section{The cospecialization map}
\label{section-cospecialization}

\noindent
Let $f : X \to S$ be a morphism of schemes. Let $\overline{x}$ be a
geometric point of $X$ with image $\overline{s} = f(\overline{x})$ in $S$.
Let $\overline{t}$ be a geometric point of
$\Spec(\mathcal{O}^{sh}_{S, \overline{s}})$.
Let $K \in D(X_\etale)$. For any morphism $g : Y \to X$ of schemes
we write $K|_Y$ instead of $g_{small}^{-1}K$ and $R\Gamma(Y, K)$
instead of $R\Gamma(Y_\etale, g_{small}^{-1}K)$.
We claim that if
\begin{enumerate}
\item $K$ is bounded below, i.e., $K \in D^+(X_\etale)$,
\item $f$ is locally acyclic relative to $K$
\end{enumerate}
then there is a {\it cospecialization map}
$$
cosp :
R\Gamma(X_{\overline{t}}, K)
\longrightarrow
R\Gamma(X_{\overline{s}}, K)
$$
which will be closely related to the specialization map
considered in Section \ref{section-specialization} and especially
Remark \ref{remark-specialization-map-and-fibres}.

\medskip\noindent
To construct the map we consider the morphisms
$$
X_{\overline{t}}
\xrightarrow{h}
X \times_S \Spec(\mathcal{O}^{sh}_{S, \overline{s}})
\xleftarrow{i}
X_{\overline{s}}
$$
The unit of the adjunction between $h^{-1}$ and $Rh_*$ gives a map
$$
\beta_{K, \overline{s}, \overline{t}} :
K|_{X \times_S \Spec(\mathcal{O}^{sh}_{S, \overline{s}})}
\longrightarrow
Rh_*(K|_{X_{\overline{t}}})
$$
in $D((X \times_S \Spec(\mathcal{O}^{sh}_{S, \overline{s}}))_\etale)$.
Lemma \ref{lemma-beta-is-isomorphism} below shows that the pullback
$i^{-1}\beta_{K, \overline{s}, \overline{t}}$
is an isomorphism under the assumptions above.
Thus we can define the cospecialization map as the composition
\begin{align*}
R\Gamma(X_{\overline{t}}, K)
& =
R\Gamma(X \times_S \Spec(\mathcal{O}^{sh}_{S, \overline{s}}),
Rh_*(K|_{X_{\overline{t}}})) \\
&
\xrightarrow{i^{-1}}
R\Gamma(X_{\overline{s}}, i^{-1}Rh_*(K|_{X_{\overline{t}}})) \\
&
\xrightarrow{(i^{-1}\beta_{K, \overline{s}, \overline{t}})^{-1}}
R\Gamma(X_{\overline{s}},
i^{-1}(K|_{X \times_S \Spec(\mathcal{O}^{sh}_{S, \overline{s}})})) \\
& =
R\Gamma(X_{\overline{s}}, K)
\end{align*}

\begin{lemma}
\label{lemma-beta-is-isomorphism}
The map $i^{-1}\beta_{K, \overline{s}, \overline{t}}$ is an isomorphism.
\end{lemma}

\begin{proof}
The construction of the maps $h$, $i$, $\beta_{K, \overline{s}, \overline{t}}$
only depends on the
base change of $X$ and $K$ to $\Spec(\mathcal{O}^{sh}_{S, \overline{s}})$.
Thus we may and do assume that $S$ is a strictly henselian scheme
with closed point $\overline{s}$. Observe that the local acyclicity
of $f$ relative to $K$ is preserved by this base change (for example
by Lemma \ref{lemma-locally-acyclic-quasi-finite-base-change} or just
directly by comparing strictly henselian rings in this very special case).

\medskip\noindent
Let $\overline{x}$ be a geometric point of $X_{\overline{s}}$.
Or equivalently, let $\overline{x}$ be a geometric point whose image
by $f$ is $\overline{s}$. Let us compute the stalk of
$i^{-1}\beta_{K, \overline{s}, \overline{t}}$
at $\overline{x}$. First, we have
$$
(i^{-1}\beta_{K, \overline{s}, \overline{t}})_{\overline{x}} =
(\beta_{K, \overline{s}, \overline{t}})_{\overline{x}}
$$
since pullback preserves stalks, see Lemma \ref{lemma-stalk-pullback}.
Since we are in the situation $S = \Spec(\mathcal{O}^{sh}_{S, \overline{s}})$
we see that $h : X_{\overline{t}} \to X$ has the property that
$X_{\overline{t}} \times_X \Spec(\mathcal{O}^{sh}_{X, \overline{x}})
= F_{\overline{x}, \overline{t}}$. Thus we see that
$$
(\beta_{K, \overline{s}, \overline{t}})_{\overline{x}} :
K_{\overline{x}}
\longrightarrow
Rh_*(K|_{X_{\overline{t}}})_{\overline{x}} =
R\Gamma(F_{\overline{x}, \overline{t}}, K)
$$
where the equal sign is Theorem \ref{theorem-higher-direct-images}.
It follows that the map $(\beta_{K, \overline{s}, \overline{t}})_{\overline{x}}$
is none other than the map $\alpha_{K, \overline{x}, \overline{t}}$ used
in Definition \ref{definition-locally-acyclic}. The result follows as we may
check whether a map is an isomorphism in stalks by
Theorem \ref{theorem-exactness-stalks}.
\end{proof}

\noindent
The cospecialization map when it exists is trying to be the inverse
of the specialization map.

\begin{lemma}
\label{lemma-specialization-cospecialization}
In the situation above, if in addition
$f$ is quasi-compact and quasi-separated, then the diagram
$$
\xymatrix{
(Rf_*K)_{\overline{s}} \ar[r] \ar[d]_{sp} &
R\Gamma(X_{\overline{s}}, K) \\
(Rf_*K)_{\overline{t}} \ar[r] &
R\Gamma(X_{\overline{t}}, K) \ar[u]_{cosp}
}
$$
is commutative.
\end{lemma}

\begin{proof}
As in the proof of Lemma \ref{lemma-beta-is-isomorphism}
we may replace $S$ by $\Spec(\mathcal{O}^{sh}_{S, \overline{s}})$.
Then our maps simplify to $h : X_{\overline{t}} \to X$,
$i : X_{\overline{s}} \to X$, and
$\beta_{K, \overline{s}, \overline{t}} : K \to Rh_*(K|_{X_{\overline{t}}})$.
Using that $(Rf_*K)_{\overline{s}} = R\Gamma(X, K)$ by
Theorem \ref{theorem-higher-direct-images}
the composition of $sp$ with the base change map
$(Rf_*K)_{\overline{t}} \to R\Gamma(X_{\overline{t}}, K)$
is just pullback of cohomology along $h$.
This is the same as the map
$$
R\Gamma(X, K) \xrightarrow{\beta_{K, \overline{s}, \overline{t}}}
R\Gamma(X, Rh_*(K|_{X_{\overline{t}}})) =
R\Gamma(X_{\overline{t}}, K)
$$
Now the map $cosp$ first inverts the $=$ sign in this displayed
formula, then pulls back along $i$, and finally applies
the inverse of $i^{-1}\beta_{K, \overline{s}, \overline{t}}$.
Hence we get the desired commutativity.
\end{proof}

\begin{lemma}
\label{lemma-sp-isom-proper-torsion-loc-ac}
Let $f : X \to S$ be a morphism of schemes. Let $K \in D(X_\etale)$.
Assume
\begin{enumerate}
\item $K$ is bounded below, i.e., $K \in D^+(X_\etale)$,
\item $f$ is locally acyclic relative to $K$,
\item $f$ is proper, and
\item $K$ has torsion cohomology sheaves.
\end{enumerate}
Then for every geometric point $\overline{s}$ of $S$ and every geometric
point $\overline{t}$ of $\Spec(\mathcal{O}^{sh}_{S, \overline{s}})$
both the specialization map
$sp : (Rf_*K)_{\overline{s}} \to (Rf_*K)_{\overline{t}}$
and the cospecialization map
$cosp : R\Gamma(X_{\overline{t}}, K) \to R\Gamma(X_{\overline{s}}, K)$
are isomorphisms.
\end{lemma}

\begin{proof}
By the proper base change theorem (in the form of
Lemma \ref{lemma-proper-base-change-stalk}) we have
$(Rf_*K)_{\overline{s}} = R\Gamma(X_{\overline{s}}, K)$
and similarly for $\overline{t}$. The ``correct'' proof would be
to show that the argument in
Lemma \ref{lemma-specialization-cospecialization}
shows that $sp$ and $cosp$ are inverse isomorphisms in this case.
Instead we will show directly that $cosp$ is an isomorphism.
From the discussion above we see that $cosp$ is an isomorphism
if and only if pullback by $i$
$$
R\Gamma(X \times_S \Spec(\mathcal{O}^{sh}_{S, \overline{s}}),
Rh_*(K|_{X_{\overline{t}}}))
\longrightarrow
R\Gamma(X_{\overline{s}}, i^{-1}Rh_*(K|_{X_{\overline{t}}}))
$$
is an isomorphism in $D^+(\textit{Ab})$. This is true by the proper
base change theorem for the proper morphism
$f' : X \times_S \Spec(\mathcal{O}^{sh}_{S, \overline{s}}) \to
\Spec(\mathcal{O}^{sh}_{S, \overline{s}})$
by the morphism $\overline{s} \to \Spec(\mathcal{O}^{sh}_{S, \overline{s}})$
and the complex $K' = Rh_*(K|_{X_{\overline{t}}})$. The complex
$K'$ is bounded below and has torsion cohomology sheaves
by Lemma \ref{lemma-torsion-direct-image}.
Since $\Spec(\mathcal{O}^{sh}_{S, \overline{s}})$ is strictly henselian
with $\overline{s}$ lying over the closed point,
we see that the source of the displayed arrow equals
$(Rf'_*K')_{\overline{s}}$ and the target equals
$R\Gamma(X_{\overline{s}}, K')$ and the displayed map is an isomorphism
by the already used
Lemma \ref{lemma-proper-base-change-stalk}.
Thus we see that three out of the four arrows in the diagram
of Lemma \ref{lemma-specialization-cospecialization} are isomorphisms
and we conclude.
\end{proof}

\begin{lemma}
\label{lemma-sp-isom-proper-loc-cst-torsion}
Let $f : X \to S$ be a morphism of schemes. Let $\mathcal{F}$
be an abelian sheaf on $X_\etale$. Assume
\begin{enumerate}
\item $f$ is smooth and proper
\item $\mathcal{F}$ is locally constant, and
\item $\mathcal{F}_{\overline{x}}$ is a torsion group all of
whose elements have order prime to the residue characteristic of
$\overline{x}$ for every geometric point $\overline{x}$ of $X$.
\end{enumerate}
Then for every geometric point $\overline{s}$ of $S$ and every geometric
point $\overline{t}$ of $\Spec(\mathcal{O}^{sh}_{S, \overline{s}})$
the specialization map
$sp : (Rf_*\mathcal{F})_{\overline{s}} \to (Rf_*\mathcal{F})_{\overline{t}}$
is an isomorphism.
\end{lemma}

\begin{proof}
This follows from Lemmas \ref{lemma-sp-isom-proper-torsion-loc-ac}
and \ref{lemma-locally-acyclic-locally-constant} and
Proposition \ref{proposition-smooth-locally-acyclic}.
\end{proof}

















\section{Cohomological dimension}
\label{section-cd}

\noindent
We can deduce some bounds on the cohomological dimension of
schemes and on the cohomological dimension of fields using
the results in Section \ref{section-vanishing-torsion}
and one, seemingly innocuous, application of the proper base
change theorem (in the proof of Proposition \ref{proposition-cd-affine}).

\begin{definition}
\label{definition-cd}
Let $X$ be a quasi-compact and quasi-separated scheme.
The {\it cohomological dimension of $X$} is the smallest
element
$$
\text{cd}(X) \in \{0, 1, 2, \ldots\} \cup \{\infty\}
$$
such that for any abelian torsion sheaf $\mathcal{F}$
on $X_\etale$ we have $H^i_\etale(X, \mathcal{F}) = 0$
for $i > \text{cd}(X)$. If $X = \Spec(A)$ we sometimes
call this the cohomological dimension of $A$.
\end{definition}

\noindent
If the scheme is in characteristic $p$, then we often can
obtain sharper bounds for the vanishing of cohomology of
$p$-power torsion sheaves. We will address this elsewhere
(insert future reference here).

\begin{lemma}
\label{lemma-cd-limit}
Let $X = \lim X_i$ be a directed limit of a system of
quasi-compact and quasi-separated schemes with affine
transition morphisms. Then $\text{cd}(X) \leq \max \text{cd}(X_i)$.
\end{lemma}

\begin{proof}
Denote $f_i : X \to X_i$ the projections.
Let $\mathcal{F}$ be an abelian torsion sheaf on $X_\etale$.
Then we have $\mathcal{F} = \lim f_i^{-1}f_{i, *}\mathcal{F}$
by Lemma \ref{lemma-linus-hamann}.
Thus $H^q_\etale(X, \mathcal{F}) =
\colim H^q_\etale(X_i, f_{i, *}\mathcal{F})$
by Theorem \ref{theorem-colimit}.
The lemma follows.
\end{proof}

\begin{lemma}
\label{lemma-cd-curve-over-field}
Let $K$ be a field. Let $X$ be a $1$-dimensional
affine scheme of finite type over $K$. Then
$\text{cd}(X) \leq 1 + \text{cd}(K)$.
\end{lemma}

\begin{proof}
Let $\mathcal{F}$ be an abelian torsion sheaf on $X_\etale$.
Consider the Leray spectral sequence for the morphism
$f : X \to \Spec(K)$. We obtain
$$
E_2^{p, q} = H^p(\Spec(K), R^qf_*\mathcal{F})
$$
converging to $H^{p + q}_\etale(X, \mathcal{F})$.
The stalk of $R^qf_*\mathcal{F}$ at a geometric point
$\Spec(\overline{K}) \to \Spec(K)$ is the cohomology of the
pullback of $\mathcal{F}$ to $X_{\overline{K}}$.
Hence it vanishes in degrees $\geq 2$ by
Theorem \ref{theorem-vanishing-affine-curves}.
\end{proof}

\begin{lemma}
\label{lemma-cd-field-extension}
Let $L/K$ be a field extension. Then we have
$\text{cd}(L) \leq \text{cd}(K) + \text{trdeg}_K(L)$.
\end{lemma}

\begin{proof}
If $\text{trdeg}_K(L) = \infty$, then this is clear.
If not then we can find a sequence of extensions
$L= L_r/L_{r - 1}/ \ldots /L_1/L_0 = K$ such that
$\text{trdeg}_{L_i}(L_{i + 1}) = 1$ and $r = \text{trdeg}_K(L)$.
Hence it suffices to prove the lemma in the case that $r = 1$.
In this case we can write $L = \colim A_i$
as a filtered colimit of its finite type $K$-subalgebras.
By Lemma \ref{lemma-cd-limit} it suffices to prove that
$\text{cd}(A_i) \leq 1 + \text{cd}(K)$. This follows
from Lemma \ref{lemma-cd-curve-over-field}.
\end{proof}

\begin{lemma}
\label{lemma-strictly-henselian}
Let $K$ be a field. Let $X$ be a scheme of finite type over $K$.
Let $x \in X$. Set $a = \text{trdeg}_K(\kappa(x))$
and $d = \dim_x(X)$. Then there is a map
$$
K(t_1, \ldots, t_a)^{sep} \longrightarrow \mathcal{O}_{X, x}^{sh}
$$
such that
\begin{enumerate}
\item the residue field of $\mathcal{O}_{X, x}^{sh}$ is a purely inseparable
extension of $K(t_1, \ldots, t_a)^{sep}$,
\item $\mathcal{O}_{X, x}^{sh}$ is a filtered colimit of finite
type $K(t_1, \ldots, t_a)^{sep}$-algebras of dimension $\leq d - a$.
\end{enumerate}
\end{lemma}

\begin{proof}
We may assume $X$ is affine. By Noether normalization, after possibly
shrinking $X$ again, we can choose a finite morphism
$\pi : X \to \mathbf{A}^d_K$, see
Algebra, Lemma \ref{algebra-lemma-Noether-normalization-at-point}.
Since $\kappa(x)$ is a finite extension of the residue field of $\pi(x)$,
this residue field has transcendence degree $a$ over $K$ as well.
Thus we can find a finite morphism $\pi' : \mathbf{A}^d_K \to \mathbf{A}^d_K$
such that $\pi'(\pi(x))$ corresponds to the generic point of the linear
subspace $\mathbf{A}^a_K \subset \mathbf{A}^d_K$ given by setting
the last $d - a$ coordinates equal to zero. Hence the composition
$$
X \xrightarrow{\pi' \circ \pi} \mathbf{A}^d_K \xrightarrow{p} \mathbf{A}^a_K
$$
of $\pi' \circ \pi$ and the projection $p$ onto the first $a$ coordinates
maps $x$ to the generic point $\eta \in \mathbf{A}^a_K$. The induced map
$$
K(t_1, \ldots, t_a)^{sep} =
\mathcal{O}_{\mathbf{A}^a_k, \eta}^{sh}
\longrightarrow \mathcal{O}_{X, x}^{sh}
$$
on \'etale local rings satisfies (1) since it is clear that the residue field
of $\mathcal{O}_{X, x}^{sh}$ is an algebraic extension of the separably closed
field $K(t_1, \ldots, t_a)^{sep}$. On the other hand, if $X = \Spec(B)$,
then $\mathcal{O}_{X, x}^{sh} = \colim B_j$ is a filtered colimit
of \'etale $B$-algebras $B_j$. Observe that $B_j$ is quasi-finite over
$K[t_1, \ldots, t_d]$ as $B$ is finite over $K[t_1, \ldots, t_d]$.
We may similarly write
$K(t_1, \ldots, t_a)^{sep} = \colim A_i$ as a filtered colimit
of \'etale $K[t_1, \ldots, t_a]$-algebras. For every $i$ we can
find an $j$ such that
$A_i \to K(t_1, \ldots, t_a)^{sep} \to \mathcal{O}_{X, x}^{sh}$
factors through a map $\psi_{i, j} : A_i \to B_j$. Then $B_j$ is quasi-finite
over $A_i[t_{a + 1}, \ldots, t_d]$. Hence
$$
B_{i, j} = B_j \otimes_{\psi_{i, j}, A_i} K(t_1, \ldots, t_a)^{sep}
$$
has dimension $\leq d - a$ as it is quasi-finite over
$K(t_1, \ldots, t_a)^{sep}[t_{a + 1}, \ldots, t_d]$.
The proof of (2) is now finished as $\mathcal{O}_{X, x}^{sh}$ is a
filtered colimit\footnote{Let $R$ be a ring.
Let $A = \colim_{i \in I} A_i$
be a filtered colimit of finitely presented $R$-algebras.
Let $B = \colim_{j \in J} B_j$ be a filtered colimit of $R$-algebras.
Let $A \to B$ be an $R$-algebra map.
Assume that for all $i \in I$ there is a $j \in J$ and an $R$-algebra map
$\psi_{i, j} : A_i \to B_j$.
Say $(i', j', \psi_{i', j'}) \geq (i, j, \psi_{i, j})$ if
$i' \geq i$, $j' \geq j$, and $\psi_{i, j}$ and $\psi_{i', j'}$
are compatible. Then the collection of triples forms a
directed set and $B = \colim B_j \otimes_{\psi_{i, j} A_i} A$.}
of the algebras $B_{i, j}$. Some details omitted.
\end{proof}

\begin{proposition}
\label{proposition-cd-affine}
Let $K$ be a field. Let $X$ be an affine scheme of finite type over $K$.
Then we have $\text{cd}(X) \leq \dim(X) + \text{cd}(K)$.
\end{proposition}

\begin{proof}
We will prove this by induction on $\dim(X)$.
Let $\mathcal{F}$ be an abelian torsion sheaf on $X_\etale$.

\medskip\noindent
The case $\dim(X) = 0$. In this case the structure morphism
$f : X \to \Spec(K)$ is finite. Hence we see that $R^if_*\mathcal{F} = 0$
for $i > 0$, see Proposition \ref{proposition-finite-higher-direct-image-zero}.
Thus $H^i_\etale(X, \mathcal{F}) = H^i_\etale(\Spec(K), f_*\mathcal{F})$
by the Leray spectral sequence for $f$
(Cohomology on Sites, Lemma \ref{sites-cohomology-lemma-Leray})
and the result is clear.

\medskip\noindent
The case $\dim(X) = 1$. This is Lemma \ref{lemma-cd-curve-over-field}.

\medskip\noindent
Assume $d = \dim(X) > 1$ and the proposition holds for finite type
affine schemes of dimension $< d$ over fields.
By Noether normalization, see for example
Varieties, Lemma \ref{varieties-lemma-noether-normalization-affine},
there exists a finite morphism $f : X \to \mathbf{A}^d_K$.
Recall that $R^if_*\mathcal{F} = 0$ for $i > 0$ by
Proposition \ref{proposition-finite-higher-direct-image-zero}.
By the Leray spectral sequence for $f$
(Cohomology on Sites, Lemma \ref{sites-cohomology-lemma-Leray})
we conclude that it suffices to prove the result for $\pi_*\mathcal{F}$
on $\mathbf{A}^d_K$.

\medskip\noindent
Interlude I. Let $j : X \to Y$ be an open immersion of smooth
$d$-dimensional varieties over $K$ (not necessarily affine)
whose complement is the support of an effective Cartier divisor $D$.
The sheaves $R^qj_*\mathcal{F}$ for $q > 0$ are supported on $D$.
We claim that $(R^qj_*\mathcal{F})_{\overline{y}} = 0$
for $a = \text{trdeg}_K(\kappa(y)) > d - q$. Namely, by
Theorem \ref{theorem-higher-direct-images} we have
$$
(R^qj_*\mathcal{F})_{\overline{y}} =
H^q(\Spec(\mathcal{O}_{Y, y}^{sh}) \times_Y X, \mathcal{F})
$$
Choose a local equation $f \in \mathfrak m_y \subset \mathcal{O}_{Y, y}$
for $D$. Then we have
$$
\Spec(\mathcal{O}_{Y, y}^{sh}) \times_Y X =
\Spec(\mathcal{O}_{Y, y}^{sh}[1/f])
$$
Using Lemma \ref{lemma-strictly-henselian} we get an embedding
$$
K(t_1, \ldots, t_a)^{sep}(x) =
K(t_1, \ldots, t_a)^{sep}[x]_{(x)}[1/x]
\longrightarrow
\mathcal{O}_{Y, y}^{sh}[1/f]
$$
Since the transcendence degree over $K$ of the fraction field of
$\mathcal{O}_{Y, y}^{sh}$ is $d$, we see that $\mathcal{O}_{Y, y}^{sh}[1/f]$
is a filtered colimit of $(d - a - 1)$-dimensional finite type algebras over
the field $K(t_1, \ldots, t_a)^{sep}(x)$ which itself has cohomological
dimension $1$ by Lemma \ref{lemma-cd-field-extension}. Thus by induction
hypothesis and Lemma \ref{lemma-cd-limit}
we obtain the desired vanishing.

\medskip\noindent
Interlude II. Let $Z$ be a smooth variety over $K$ of dimension $d - 1$.
Let $E_a \subset Z$ be the set of points $z \in Z$ with
$\text{trdeg}_K(\kappa(z)) \leq a$. Observe that $E_a$ is closed
under specialization, see
Varieties, Lemma \ref{varieties-lemma-dimension-locally-algebraic}.
Suppose that $\mathcal{G}$ is a torsion abelian sheaf on $Z$
whose support is contained in $E_a$. Then we claim that
$H^b_\etale(Z, \mathcal{G}) = 0$ for $b > a + \text{cd}(K)$.
Namely, we can write $\mathcal{G} = \colim \mathcal{G}_i$
with $\mathcal{G}_i$ a torsion abelian sheaf
supported on a closed subscheme $Z_i$ contained in $E_a$, see
Lemma \ref{lemma-support-in-subset}.
Then the induction hypothesis kicks in to imply
the desired vanishing for $\mathcal{G}_i$\footnote{Here
we first use Proposition \ref{proposition-closed-immersion-pushforward}
to write $\mathcal{G}_i$ as the pushforward of a sheaf on $Z_i$,
the induction hypothesis gives the vanishing for this sheaf on $Z_i$, and
the Leray spectral sequence for $Z_i \to Z$ gives the vanishing
for $\mathcal{G}_i$.}. Finally, we
conclude by Theorem \ref{theorem-colimit}.

\medskip\noindent
Consider the commutative diagram
$$
\xymatrix{
\mathbf{A}^d_K \ar[rd]_f \ar[rr]_-j & &
\mathbf{P}^1_K \times_K \mathbf{A}^{d - 1}_K \ar[ld]^g \\
& \mathbf{A}^{d - 1}_K
}
$$
Observe that $j$ is an open immersion of smooth $d$-dimensional
varieties whose complement is an effective Cartier divisor $D$. Thus
we may use the results obtained in interlude I.
We are going to study the relative Leray spectral sequence
$$
E_2^{p, q} = R^pg_*R^qj_*\mathcal{F} \Rightarrow R^{p + q}f_*\mathcal{F}
$$
Since $R^qj_*\mathcal{F}$ for $q > 0$ is supported on $D$ and since
$g|_D : D \to \mathbf{A}^{d - 1}_K$ is an isomorphism, we find
$R^pg_*R^qj_*\mathcal{F} = 0$ for $p > 0$ and $q > 0$. Moreover, we have
$R^qj_*\mathcal{F} = 0$ for $q > d$. On the other hand, $g$ is a
proper morphism of relative dimension $1$. Hence by
Lemma \ref{lemma-cohomological-dimension-proper}
we see that $R^pg_*j_*\mathcal{F} = 0$ for $p > 2$. Thus the $E_2$-page
of the spectral sequence looks like this
$$
\begin{matrix}
g_*R^dj_*\mathcal{F} & 0 & 0 \\
\ldots & \ldots & \ldots \\
g_*R^2j_*\mathcal{F} & 0 & 0 \\
g_*R^1j_*\mathcal{F} & 0 & 0 \\
g_*j_*\mathcal{F} & R^1g_*j_*\mathcal{F} & R^2g_*j_*\mathcal{F}
\end{matrix}
$$
We conclude that $R^qf_*\mathcal{F} = g_*R^qj_*\mathcal{F}$ for $q > 2$.
By interlude I we see that the support of $R^qf_*\mathcal{F}$ for $q > 2$
is contained in the set of points of $\mathbf{A}^{d - 1}_K$
whose residue field has transcendence degree $\leq d - q$.
By interlude II
$$
H^p(\mathbf{A}^{d - 1}_K, R^qf_*\mathcal{F}) = 0
\text{ for }p > d - q + \text{cd}(K)\text{ and }q > 2
$$
On the other hand, by Theorem \ref{theorem-higher-direct-images}
we have $R^2f_*\mathcal{F}_{\overline{\eta}} =
H^2(\mathbf{A}^1_{\overline{\eta}}, \mathcal{F}) = 0$
(vanishing by the case of dimension $1$) where $\eta$ is the
generic point of $\mathbf{A}^{d - 1}_K$.
Hence by interlude II again we see
$$
H^p(\mathbf{A}^{d - 1}_K, R^2f_*\mathcal{F}) = 0
\text{ for }p > d - 2 + \text{cd}(K)
$$
Finally, we have
$$
H^p(\mathbf{A}^{d - 1}_K, R^qf_*\mathcal{F}) = 0
\text{ for }p > d - 1 + \text{cd}(K)\text{ and }q = 0, 1
$$
by induction hypothesis. Combining everything we just said
with the Leray spectral sequence
$H^p(\mathbf{A}^{d - 1}_K, R^qf_*\mathcal{F}) \Rightarrow
H^{p + q}(\mathbf{A}^d_K, \mathcal{F})$ we conclude.
\end{proof}

\begin{lemma}
\label{lemma-interlude-II}
Let $K$ be a field. Let $X$ be an affine scheme of finite type over $K$.
Let $E_a \subset X$ be the set of points
$x \in X$ with $\text{trdeg}_K(\kappa(x)) \leq a$.
Let $\mathcal{F}$ be an abelian torsion sheaf on $X_\etale$
whose support is contained in $E_a$. Then
$H^b_\etale(X, \mathcal{F}) = 0$ for $b > a + \text{cd}(K)$.
\end{lemma}

\begin{proof}
We can write $\mathcal{F} = \colim \mathcal{F}_i$
with $\mathcal{F}_i$ a torsion abelian sheaf
supported on a closed subscheme $Z_i$ contained in $E_a$, see
Lemma \ref{lemma-support-in-subset}.
Then Proposition \ref{proposition-cd-affine} gives
the desired vanishing for $\mathcal{F}_i$. Details omitted;
hints: first use Proposition \ref{proposition-closed-immersion-pushforward}
to write $\mathcal{F}_i$ as the pushforward of a sheaf on $Z_i$,
use the vanishing for this sheaf on $Z_i$, and use
the Leray spectral sequence for $Z_i \to Z$ to get the vanishing
for $\mathcal{F}_i$. Finally, we
conclude by Theorem \ref{theorem-colimit}.
\end{proof}

\begin{lemma}
\label{lemma-interlude-I}
Let $f : X \to Y$ be an affine morphism of schemes of finite
type over a field $K$. Let $E_a(X)$ be the set of points $x \in X$
with $\text{trdeg}_K(\kappa(x)) \leq a$.
Let $\mathcal{F}$ be an abelian torsion sheaf on $X_\etale$
whose support is contained in $E_a$. Then
$R^qf_*\mathcal{F}$ has support contained in
$E_{a - q}(Y)$.
\end{lemma}

\begin{proof}
The question is local on $Y$ hence we can assume $Y$ is affine.
Then $X$ is affine too and we can choose a diagram
$$
\xymatrix{
X \ar[d]_f \ar[r]_i & \mathbf{A}^{n + m}_K \ar[d]^{\text{pr}} \\
Y \ar[r]^j & \mathbf{A}^n_K
}
$$
where the horizontal arrows are closed immersions and the vertical
arrow on the right is the projection (details omitted).
Then $j_*R^qf_*\mathcal{F} = R^q\text{pr}_*i_*\mathcal{F}$
by the vanishing of the higher direct images of $i$ and $j$, see
Proposition \ref{proposition-finite-higher-direct-image-zero}.
Moreover, the description of the stalks of $j_*$ in the proposition
shows that it suffices to prove the vanishing for $j_*R^qf_*\mathcal{F}$.
Thus we may assume $f$ is the projection
morphism $\text{pr} : \mathbf{A}^{n + m}_K \to \mathbf{A}^n_K$
and an abelian torsion sheaf $\mathcal{F}$ on $\mathbf{A}^{n + m}_K$
satisfying the assumption in the statement of the lemma.

\medskip\noindent
Let $y$ be a point in $\mathbf{A}^n_K$.
By Theorem \ref{theorem-higher-direct-images} we have
$$
(R^q\text{pr}_*\mathcal{F})_{\overline{y}} =
H^q(\mathbf{A}^{n + m}_K \times_{A^n_K} \Spec(\mathcal{O}_{Y, y}^{sh}),
\mathcal{F}) =
H^q(\mathbf{A}^m_{\mathcal{O}_{Y, y}^{sh}}, \mathcal{F})
$$
Say $b = \text{trdeg}_K(\kappa(y))$. From Lemma \ref{lemma-strictly-henselian}
we get an embedding
$$
L = K(t_1, \ldots, t_b)^{sep} \longrightarrow \mathcal{O}_{Y, y}^{sh}
$$
Write $\mathcal{O}_{Y, y}^{sh} = \colim B_i$ as the filtered
colimit of finite type $L$-subalgebras $B_i \subset \mathcal{O}_{Y, y}^{sh}$
containing the ring $K[T_1, \ldots, T_n]$ of regular functions
on $\mathbf{A}^n_K$. Then we get
$$
\mathbf{A}^m_{\mathcal{O}_{Y, y}^{sh}} =
\lim \mathbf{A}^m_{B_i}
$$
If $z \in \mathbf{A}^m_{B_i}$ is a point in the support of
$\mathcal{F}$, then the image $x$ of $z$ in $\mathbf{A}^{m + n}_K$
satisfies $\text{trdeg}_K(\kappa(x)) \leq a$ by our assumption on
$\mathcal{F}$ in the lemma.
Since $\mathcal{O}_{Y, y}^{sh}$ is a filtered colimit of
\'etale algebras over $K[T_1, \ldots, T_n]$ and since
$B_i \subset \mathcal{O}_{Y, y}^{sh}$
we see that $\kappa(z)/\kappa(x)$ is algebraic
(some details omitted). Then $\text{trdeg}_K(\kappa(z)) \leq a$
and hence $\text{trdeg}_L(\kappa(z)) \leq a - b$.
By Lemma \ref{lemma-interlude-II} we see that
$$
H^q(\mathbf{A}^m_{B_i}, \mathcal{F}) = 0\text{ for }q > a - b
$$
Thus by Theorem \ref{theorem-colimit} we get
$(Rf_*\mathcal{F})_{\overline{y}} = 0$ for $q > a - b$
as desired.
\end{proof}








\section{Finite cohomological dimension}
\label{section-finite-cd}

\noindent
We continue the discussion started in Section \ref{section-cd}.

\begin{definition}
\label{definition-cd-f}
Let $f : X \to Y$ be a quasi-compact and quasi-separated
morphism of schemes.
The {\it cohomological dimension of $f$} is the smallest
element
$$
\text{cd}(f) \in \{0, 1, 2, \ldots\} \cup \{\infty\}
$$
such that for any abelian torsion sheaf $\mathcal{F}$
on $X_\etale$ we have $R^if_*\mathcal{F} = 0$
for $i > \text{cd}(f)$.
\end{definition}

\begin{lemma}
\label{lemma-finite-cd}
Let $K$ be a field.
\begin{enumerate}
\item If $f : X \to Y$ is a morphism of finite type schemes over $K$,
then $\text{cd}(f) < \infty$.
\item If $\text{cd}(K) < \infty$, then $\text{cd}(X) < \infty$
for any finite type scheme $X$ over $K$.
\end{enumerate}
\end{lemma}

\begin{proof}
Proof of (1). We may assume $Y$ is affine. We will use the
induction principle of 
Cohomology of Schemes, Lemma \ref{coherent-lemma-induction-principle}
to prove this. If $X$ is affine too, then the result holds by
Lemma \ref{lemma-interlude-I}. Thus it suffices to show that if
$X = U \cup V$ and the result is true for $U \to Y$, $V \to Y$, and
$U \cap V \to Y$, then it is true for $f$. This follows from the
relative Mayer-Vietoris sequence, see
Lemma \ref{lemma-relative-mayer-vietoris}.

\medskip\noindent
Proof of (2). We will use the induction principle of
Cohomology of Schemes, Lemma \ref{coherent-lemma-induction-principle}
to prove this. If $X$ is affine, then the result holds by
Proposition \ref{proposition-cd-affine}. Thus it suffices to show that if
$X = U \cup V$ and the result is true for $U$, $V$, and
$U \cap V $, then it is true for $X$. This follows from the
Mayer-Vietoris sequence, see
Lemma \ref{lemma-mayer-vietoris}.
\end{proof}

\begin{lemma}
\label{lemma-finite-cd-mod-n-direct-sums}
Cohomology and direct sums. Let $n \geq 1$ be an integer.
\begin{enumerate}
\item Let $f : X \to Y$ be a quasi-compact and quasi-separated morphism
of schemes with $\text{cd}(f) < \infty$. Then the functor
$$
Rf_* :
D(X_\etale, \mathbf{Z}/n\mathbf{Z})
\longrightarrow
D(Y_\etale, \mathbf{Z}/n\mathbf{Z})
$$
commutes with direct sums.
\item Let $X$ be a quasi-compact and quasi-separated scheme with
$\text{cd}(X) < \infty$. Then the functor
$$
R\Gamma(X, -) :
D(X_\etale, \mathbf{Z}/n\mathbf{Z})
\longrightarrow
D(\mathbf{Z}/n\mathbf{Z})
$$
commutes with direct sums.
\end{enumerate}
\end{lemma}

\begin{proof}
Proof of (1). Since $\text{cd}(f) < \infty$ we see that
$$
f_* :
\textit{Mod}(X_\etale, \mathbf{Z}/n\mathbf{Z})
\longrightarrow
\textit{Mod}(Y_\etale, \mathbf{Z}/n\mathbf{Z})
$$
has finite cohomological dimension in the sense of
Derived Categories, Lemma \ref{derived-lemma-unbounded-right-derived}.
Let $I$ be a set and for $i \in I$ let $E_i$ be
an object of $D(X_\etale, \mathbf{Z}/n\mathbf{Z})$.
Choose a K-injective complex $\mathcal{I}_i^\bullet$
of $\mathbf{Z}/n\mathbf{Z}$-modules each of whose
terms $\mathcal{I}_i^n$ is an injective sheaf of
$\mathbf{Z}/n\mathbf{Z}$-modules representing $E_i$.
See Injectives, Theorem
\ref{injectives-theorem-K-injective-embedding-grothendieck}.
Then $\bigoplus E_i$ is represented by the complex
$\bigoplus \mathcal{I}_i^\bullet$ (termwise direct sum), see
Injectives, Lemma \ref{injectives-lemma-derived-products}.
By Lemma \ref{lemma-relative-colimit} we have
$$
R^qf_*(\bigoplus \mathcal{I}_i^n) =
\bigoplus R^qf_*(\mathcal{I}_i^n) = 0
$$
for $q > 0$ and any $n$. Hence we conclude by 
Derived Categories, Lemma \ref{derived-lemma-unbounded-right-derived}
that we may compute $Rf_*(\bigoplus E_i)$ by the complex
$$
f_*(\bigoplus \mathcal{I}_i^\bullet) =
\bigoplus f_*(\mathcal{I}_i^\bullet)
$$
(equality again by Lemma \ref{lemma-relative-colimit}) which
represents $\bigoplus Rf_*E_i$ by the already used
Injectives, Lemma \ref{injectives-lemma-derived-products}.

\medskip\noindent
Proof of (2). This is identical to the proof of (1)
and we omit it.
\end{proof}

\begin{lemma}
\label{lemma-proper-mod-n-direct-sums}
Let $f : X \to Y$ be a proper morphism of schemes. Let $n \geq 1$
be an integer. Then the functor
$$
Rf_* :
D(X_\etale, \mathbf{Z}/n\mathbf{Z})
\longrightarrow
D(Y_\etale, \mathbf{Z}/n\mathbf{Z})
$$
commutes with direct sums.
\end{lemma}

\begin{proof}
It is enough to prove this when $Y$ is quasi-compact. By
Morphisms, Lemma \ref{morphisms-lemma-morphism-finite-type-bounded-dimension}
we see that the dimension of the fibres of
$f : X \to Y$ is bounded.
Thus Lemma \ref{lemma-cohomological-dimension-proper}
implies that $\text{cd}(f) < \infty$. Hence the
result by Lemma \ref{lemma-finite-cd-mod-n-direct-sums}.
\end{proof}

\begin{lemma}
\label{lemma-pull-out-constant-mod-n}
Let $X$ be a quasi-compact and quasi-separated scheme
such that $\text{cd}(X) < \infty$. Let $\Lambda$ be a torsion ring.
Let $E \in D(X_\etale, \Lambda)$ and $K \in D(\Lambda)$. Then
$$
R\Gamma(X, E \otimes_\Lambda^\mathbf{L} \underline{K}) =
R\Gamma(X, E) \otimes_\Lambda^\mathbf{L} K
$$
\end{lemma}

\begin{proof}
There is a canonical map from right to left by
Cohomology on Sites, Section \ref{sites-cohomology-section-projection-formula}.
Let $T(K)$ be the property that the statement of the
lemma holds for $K \in D(\Lambda)$.
We will check conditions (1), (2), and (3) of
More on Algebra, Remark \ref{more-algebra-remark-P-resolution}
hold for $T$ to conclude.
Property (1) holds because both sides of the equality
commute with direct sums, see
Lemma \ref{lemma-finite-cd-mod-n-direct-sums}.
Property (2) holds because we are comparing exact functors
between triangulated categories and we can use
Derived Categories, Lemma \ref{derived-lemma-third-isomorphism-triangle}.
Property (3) says the lemma holds when $K = \Lambda[k]$
for any shift $k \in \mathbf{Z}$ and this is obvious.
\end{proof}

\begin{lemma}
\label{lemma-projection-formula-proper-mod-n}
Let $f : X \to Y$ be a proper morphism of schemes. Let $\Lambda$
be a torsion ring. Let $E \in D(X_\etale, \Lambda)$ and
$K \in D(Y_\etale, \Lambda)$. Then
$$
Rf_*E \otimes_\Lambda^\mathbf{L} K =
Rf_*(E \otimes_\Lambda^\mathbf{L} f^{-1}K)
$$
in $D(Y_\etale, \Lambda)$.
\end{lemma}

\begin{proof}
There is a canonical map from left to right by
Cohomology on Sites, Section \ref{sites-cohomology-section-projection-formula}.
We will check the equality on stalks at $\overline{y}$.
By the proper base change (in the form of
Lemma \ref{lemma-proper-base-change-mod-n} where $Y' = \overline{y}$)
this reduces to the case where $Y$ is the spectrum of
an algebraically closed field.
This is shown in Lemma \ref{lemma-pull-out-constant-mod-n}
where we use that $\text{cd}(X) < \infty$ by
Lemma \ref{lemma-cohomological-dimension-proper}.
\end{proof}






\section{K\"unneth in \'etale cohomology}
\label{section-kunneth}

\noindent
We first prove a K\"unneth formula in case one of the factors
is proper. Then we use this formula to prove a base change property
for open immersions. This then gives a
``base change by morphisms towards spectra of fields''
(akin to smooth base change).
Finally we use this to get a more general K\"unneth formula.

\begin{remark}
\label{remark-define-kunneth-map}
Consider a cartesian diagram in the category of schemes:
$$
\xymatrix{
X \times_S Y \ar[d]_p \ar[r]_q \ar[rd]_c & Y \ar[d]^g \\
X \ar[r]^f & S
}
$$
Let $\Lambda$ be a ring and let $E \in D(X_\etale, \Lambda)$
and $K \in D(Y_\etale, \Lambda)$. Then there is a canonical map
$$
Rf_*E \otimes_\Lambda^\mathbf{L} Rg_*K
\longrightarrow
Rc_*(p^{-1}E \otimes_\Lambda^\mathbf{L} q^{-1}K)
$$
For example we can define this using the canonical maps
$Rf_*E \to Rc_*p^{-1}E$ and $Rg_*K \to Rc_*q^{-1}K$ and
the relative cup product defined in Cohomology on Sites,
Remark \ref{sites-cohomology-remark-cup-product}.
Or you can use the adjoint to the map
$$
c^{-1}(Rf_*E \otimes_\Lambda^\mathbf{L} Rg_*K)
=
p^{-1}f^{-1}Rf_*E \otimes_\Lambda^\mathbf{L} q^{-1} g^{-1}Rg_*K
\to
p^{-1}E \otimes_\Lambda^\mathbf{L} q^{-1}K
$$
which uses the adjunction maps $f^{-1}Rf_*E \to E$ and
$g^{-1}Rg_*K \to K$.
\end{remark}


\begin{lemma}
\label{lemma-kunneth-one-proper}
Let $k$ be a separably closed field. Let $X$ be a proper scheme over $k$.
Let $Y$ be a quasi-compact and quasi-separated scheme over $k$.
\begin{enumerate}
\item If $E \in D^+(X_\etale)$ has torsion cohomology sheaves and
$K \in D^+(Y_\etale)$, then
$$
R\Gamma(X \times_{\Spec(k)} Y,
\text{pr}_1^{-1}E
\otimes_\mathbf{Z}^\mathbf{L}
\text{pr}_2^{-1}K
)
=
R\Gamma(X, E)
\otimes_\mathbf{Z}^\mathbf{L}
R\Gamma(Y, K)
$$
\item If $n \geq 1$ is an integer, $Y$ is of finite type over $k$,
$E \in D(X_\etale, \mathbf{Z}/n\mathbf{Z})$, and
$K \in D(Y_\etale, \mathbf{Z}/n\mathbf{Z})$, then
$$
R\Gamma(X \times_{\Spec(k)} Y,
\text{pr}_1^{-1}E
\otimes_{\mathbf{Z}/n\mathbf{Z}}^\mathbf{L}
\text{pr}_2^{-1}K
)
=
R\Gamma(X, E)
\otimes_{\mathbf{Z}/n\mathbf{Z}}^\mathbf{L}
R\Gamma(Y, K)
$$
\end{enumerate}
\end{lemma}

\begin{proof}
Proof of (1). By Lemma \ref{lemma-projection-formula-proper} we have
$$
R\text{pr}_{2, *}(
\text{pr}_1^{-1}E
\otimes_\mathbf{Z}^\mathbf{L}
\text{pr}_2^{-1}K) =
R\text{pr}_{2, *}(\text{pr}_1^{-1}E)
\otimes_\mathbf{Z}^\mathbf{L}
K
$$
By proper base change (in the form of Lemma \ref{lemma-proper-base-change})
this is equal to the object
$$
\underline{R\Gamma(X, E)}
\otimes_\mathbf{Z}^\mathbf{L}
K
$$
of $D(Y_\etale)$. Taking $R\Gamma(Y, -)$ on this object reproduces the
left hand side of the equality in (1) by the Leray spectral sequence
for $\text{pr}_2$. Thus we conclude by
Lemma \ref{lemma-pull-out-constant}.

\medskip\noindent
Proof of (2). This is exactly the same as the proof of (1)
except that we use Lemmas \ref{lemma-projection-formula-proper-mod-n},
\ref{lemma-proper-base-change-mod-n}, and
\ref{lemma-pull-out-constant-mod-n} as well as
$\text{cd}(Y) < \infty$ by Lemma \ref{lemma-finite-cd}.
\end{proof}

\begin{lemma}
\label{lemma-supported-in-closed-points}
Let $K$ be a separably closed field. Let $X$ be a scheme of finite
type over $K$. Let $\mathcal{F}$ be an abelian sheaf on $X_\etale$
whose support is contained in the set of closed points of $X$.
Then $H^q(X, \mathcal{F}) = 0$ for $q > 0$ and $\mathcal{F}$
is globally generated.
\end{lemma}

\begin{proof}
(If $\mathcal{F}$ is torsion, then the vanishing follows immediately
from Lemma \ref{lemma-interlude-II}.)
By Lemma \ref{lemma-support-in-subset} we can write $\mathcal{F}$
as a filtered colimit of constructible sheaves $\mathcal{F}_i$
of $\mathbf{Z}$-modules whose supports $Z_i \subset X$ are finite
sets of closed points. By
Proposition \ref{proposition-closed-immersion-pushforward}
such a sheaf is of the form $(Z_i \to X)_*\mathcal{G}_i$
where $\mathcal{G}_i$ is a sheaf on $Z_i$. As $K$ is separably closed,
the scheme $Z_i$ is a finite disjoint union of spectra of separably closed
fields. Recall that $H^q(Z_i, \mathcal{G}_i) = H^q(X, \mathcal{F}_i)$
by the Leray spectral sequence for $Z_i \to X$ and vanising of higher
direct images for this morphism
(Proposition \ref{proposition-finite-higher-direct-image-zero}).
By Lemmas \ref{lemma-equivalence-abelian-sheaves-point} and
\ref{lemma-compare-cohomology-point}
we see that $H^q(Z_i, \mathcal{G}_i)$ is zero for $q > 0$
and that $H^0(Z_i, \mathcal{G}_i)$ generates $\mathcal{G}_i$.
We conclude the vanishing of $H^q(X, \mathcal{F}_i)$ for $q > 0$
and that $\mathcal{F}_i$ is generated by global sections.
By Theorem \ref{theorem-colimit} we see that
$H^q(X, \mathcal{F}) = 0$ for $q > 0$.
The proof is now done because a
filtered colimit of globally generated sheaves of abelian groups
is globally generated (details omitted).
\end{proof}

\begin{lemma}
\label{lemma-vanishing-closed-points}
Let $K$ be a separably closed field. Let $X$ be a scheme of finite
type over $K$. Let $Q \in D(X_\etale)$. Assume that $Q_{\overline{x}}$
is nonzero only if $x$ is a closed point of $X$. Then
$$
Q = 0 \Leftrightarrow H^i(X, Q) = 0 \text{ for all }i
$$
\end{lemma}

\begin{proof}
The implication from left to right is trivial. Thus we need
to prove the reverse implication.

\medskip\noindent
Assume $Q$ is bounded below; this cases suffices for almost all
applications. If $Q$ is not zero, then we can look at the smallest $i$
such that the cohomology sheaf $H^i(Q)$ is nonzero.
By Lemma \ref{lemma-supported-in-closed-points} we have
$H^i(X, Q) = H^0(X, H^i(Q)) \not = 0$ and we conclude.

\medskip\noindent
General case. Let $\mathcal{B} \subset \Ob(X_\etale)$ be the
quasi-compact objects. By Lemma \ref{lemma-supported-in-closed-points}
the assumptions of Cohomology on Sites, Lemma
\ref{sites-cohomology-lemma-cohomology-over-U-trivial}
are satisfied. We conclude that $H^q(U, Q) = H^0(U, H^q(Q))$
for all $U \in \mathcal{B}$. In particular, this holds for $U = X$.
Thus the conclusion by Lemma \ref{lemma-supported-in-closed-points}
as $Q$ is zero in $D(X_\etale)$ if and only if $H^q(Q)$ is zero
for all $q$.
\end{proof}

\begin{lemma}
\label{lemma-kunneth-localize-on-X}
Let $K$ be a field. Let $j : U \to X$ be an open immersion of
schemes of finite type over $K$. Let $Y$ be a scheme of finite type
over $K$. Consider the diagram
$$
\xymatrix{
Y \times_{\Spec(K)} X \ar[d]_q  &
Y \times_{\Spec(K)} U \ar[l]^h \ar[d]^p \\
X & U \ar[l]_j
}
$$
Then the base change map $q^{-1}Rj_*\mathcal{F} \to Rh_*p^{-1}\mathcal{F}$
is an isomorphism for $\mathcal{F}$ an abelian sheaf on $U_\etale$
whose stalks are torsion of orders invertible in $K$.
\end{lemma}

\begin{proof}
Write $\mathcal{F} = \colim \mathcal{F}[n]$ where the colimit
is over the multiplicative system of integers invertible in $K$.
Since cohomology commutes with filtered colimits in our situation
(for a precise reference see Lemma \ref{lemma-base-change-Rf-star-colim}),
it suffices to prove the lemma for $\mathcal{F}[n]$. Thus we may
assume $\mathcal{F}$ is a sheaf of $\mathbf{Z}/n\mathbf{Z}$-modules
for some $n$ invertible in $K$ (we will use this at the very end of
the proof). In the proof we use the short hand $X \times_K Y$ for the fibre
product over $\Spec(K)$.
We will prove the lemma by induction on $\dim(X) + \dim(Y)$.
The lemma is trivial if $\dim(X) \leq 0$, since in this case
$U$ is an open and closed subscheme of $X$.
Choose a point $z \in X \times_K Y$. We will show
the stalk at $\overline{z}$ is an isomorphism.

\medskip\noindent
Suppose that $z \mapsto x \in X$ and assume $\text{trdeg}_K(\kappa(x)) > 0$.
Set $X' = \Spec(\mathcal{O}_{X, x}^{sh})$
and denote $U' \subset X'$ the inverse image of $U$.
Consider the base change
$$
\xymatrix{
Y \times_K X' \ar[d]_{q'}  &
Y \times_K U' \ar[l]^{h'} \ar[d]^{p'} \\
X' & U' \ar[l]_{j'}
}
$$
of our diagram by $X' \to X$. Observe that $X' \to X$ is a filtered
colimit of \'etale morphisms. By smooth base change in the form of
Lemma \ref{lemma-smooth-base-change-general} the pullback of
$q^{-1}Rj_*\mathcal{F} \to Rh_*p^{-1}\mathcal{F}$ to $X'$ to
$Y \times_K X'$ is the map
$(q')^{-1}Rj'_*\mathcal{F}' \to Rj'_*(p')^{-1}\mathcal{F}'$ where
$\mathcal{F}'$ is the pullback of $\mathcal{F}$ to $U'$.
(In this step it would suffice to use \'etale base change which is
an essentially trivial result.) So it suffices to show
that $(q')^{-1}Rj'_*\mathcal{F}' \to Rj'_*(p')^{-1}\mathcal{F}'$
is an isomorphism in order to prove that our original
map is an isomorphism on stalks at $\overline{z}$.
By Lemma \ref{lemma-strictly-henselian}
there is a separably closed field $L/K$ such that
$X' = \lim X_i$ with $X_i$ affine of finite type over $L$
and $\dim(X_i) < \dim(X)$. For $i$ large enough there
exists an open $U_i \subset X_i$ restricting to $U'$ in $X'$.
We may apply the induction hypothesis to the diagram
$$
\vcenter{
\xymatrix{
Y \times_K X_i \ar[d]_{q_i}  &
Y \times_K U_i \ar[l]^{h_i} \ar[d]^{p_i} \\
X_i & U_i \ar[l]_{j_i}
}
}
\quad\text{equal to}\quad
\vcenter{
\xymatrix{
Y_L \times_L X_i \ar[d]_{q_i}  &
Y_L \times_L U_i \ar[l]^{h_i} \ar[d]^{p_i} \\
X_i & U_i \ar[l]_{j_i}
}
}
$$
over the field $L$ and the pullback of $\mathcal{F}$ to these diagrams.
By Lemma \ref{lemma-base-change-Rf-star-colim} we conclude that the map
$(q')^{-1}Rj'_*\mathcal{F}' \to Rj'_*(p')^{-1}\mathcal{F}$ is an isomorphism.

\medskip\noindent
Suppose that $z \mapsto y \in Y$ and assume $\text{trdeg}_K(\kappa(y)) > 0$.
Let $Y' = \Spec(\mathcal{O}_{X, x}^{sh})$.
By Lemma \ref{lemma-strictly-henselian} there is a separably closed field
$L/K$ such that $Y' = \lim Y_i$ with $Y_i$ affine of finite type over $L$
and $\dim(Y_i) < \dim(Y)$. In particular $Y'$ is a scheme over $L$.
Denote with a subscript $L$ the base change from schemes over $K$
to schemes over $L$. Consider the commutative diagrams
$$
\vcenter{
\xymatrix{
Y' \times_K X \ar[d]_f  &
Y' \times_K U \ar[l]^{h'} \ar[d]^{f'} \\
Y \times_K X \ar[d]_q  &
Y \times_K U \ar[l]^h \ar[d]^p \\
X & U \ar[l]_j
}
}
\quad\text{and}\quad
\vcenter{
\xymatrix{
Y' \times_L X_L \ar[d]_{q'}  &
Y' \times_L U_L \ar[l]^{h'} \ar[d]^{p'} \\
X_L \ar[d] &
U_L \ar[l]^{j_L} \ar[d] \\
X & U \ar[l]_j
}
}
$$
and observe the top and bottom rows are the same on the left and the right.
By smooth base change we see that
$f^{-1}Rh_*p^{-1}\mathcal{F} = Rh'_*(f')^{-1}p^{-1}\mathcal{F}$
(similarly to the previous paragraph).
By smooth base change for $\Spec(L) \to \Spec(K)$
(Lemma \ref{lemma-base-change-field-extension})
we see that $Rj_{L, *}\mathcal{F}_L$ is the pullback of
$Rj_*\mathcal{F}$ to $X_L$. Combining these two observations,
we conclude that it suffices to prove the base change map
for the upper square in the diagram on the right
is an isomorphism in order to prove that our original
map is an isomorphism on stalks at $\overline{z}$\footnote{Here we
use that a ``vertical composition'' of base change maps is a base change
map as explained in Cohomology on Sites, Remark
\ref{sites-cohomology-remark-compose-base-change}.}.
Then using that $Y' = \lim Y_i$ and argueing exactly
as in the previous paragraph we see that the induction
hypothesis forces our map over $Y' \times_K X$ to be
an isomorphism.

\medskip\noindent
Thus any counter example with $\dim(X) + \dim(Y)$ minimal
would only have nonisomorphisms
$q^{-1}Rj_*\mathcal{F} \to Rh_*p^{-1}\mathcal{F}$
on stalks at closed points of $X \times_K Y$ (because a point
$z$ of $X \times_K Y$ is a closed point if and only if
both the image of $z$ in $X$ and in $Y$ are closed).
Since it is enough to prove the isomorphism locally,
we may assume $X$ and $Y$ are affine. However, then
we can choose an open dense immersion $Y \to Y'$ with $Y'$
projective. (Choose a closed immersion $Y \to \mathbf{A}^n_K$
and let $Y'$ be the scheme theoretic closure of $Y$ in
$\mathbf{P}^n_K$.) Then $\dim(Y') = \dim(Y)$ and hence
we get a ``minimal'' counter example with $Y$ projective over $K$.
In the next paragraph we show that this can't happen.

\medskip\noindent
Consider a diagram as in the statement of the lemma
such that $q^{-1}Rj_*\mathcal{F} \to Rh_*p^{-1}\mathcal{F}$
is an isomorphism at all non-closed points of $X \times_K Y$
and such that $Y$ is projective.
The restriction of the map to $(X \times_K Y)_{K^{sep}}$
is the corresponding map for the diagram of the lemma
base changed to $K^{sep}$. Thus we may and do assume $K$
is separably algebraically closed. Choose a distinguished triangle
$$
q^{-1}Rj_*\mathcal{F} \to Rh_*p^{-1}\mathcal{F} \to Q \to
(q^{-1}Rj_*\mathcal{F})[1]
$$
in $D((X \times_K Y)_\etale)$. Since $Q$ is supported in
closed points we see that it suffices to prove
$H^i(X \times_K Y, Q) = 0$ for all $i$, see
Lemma \ref{lemma-vanishing-closed-points}.
Thus it suffices to prove that
$q^{-1}Rj_*\mathcal{F} \to Rh_*p^{-1}\mathcal{F}$
induces an isomorphism on cohomology. Recall that $\mathcal{F}$
is annihilated by $n$ invertible in $K$.
By the K\"unneth formula of Lemma \ref{lemma-kunneth-one-proper}
we have
\begin{align*}
R\Gamma(X \times_K Y, q^{-1}Rj_*\mathcal{F})
&=
R\Gamma(X, Rj_*\mathcal{F})
\otimes_{\mathbf{Z}/n\mathbf{Z}}^\mathbf{L}
R\Gamma(Y, \mathbf{Z}/n\mathbf{Z}) \\
& =
R\Gamma(U, \mathcal{F})
\otimes_{\mathbf{Z}/n\mathbf{Z}}^\mathbf{L}
R\Gamma(Y, \mathbf{Z}/n\mathbf{Z})
\end{align*}
and
$$
R\Gamma(X \times_K Y, Rh_*p^{-1}\mathcal{F}) =
R\Gamma(U \times_K Y, p^{-1}\mathcal{F}) =
R\Gamma(U, \mathcal{F})
\otimes_{\mathbf{Z}/n\mathbf{Z}}^\mathbf{L}
R\Gamma(Y, \mathbf{Z}/n\mathbf{Z})
$$
This finishes the proof.
\end{proof}

\begin{lemma}
\label{lemma-punctual-base-change}
Let $K$ be a field. For any commutative diagram
$$
\xymatrix{
X \ar[d] & X' \ar[l] \ar[d]_{f'} & Y \ar[l]^h \ar[d]^e \\
\Spec(K) & S' \ar[l] & T \ar[l]_g
}
$$
of schemes over $K$ with
$X' = X \times_{\Spec(K)} S'$ and $Y = X' \times_{S'} T$ and
$g$ quasi-compact and quasi-separated, and every abelian sheaf
$\mathcal{F}$ on $T_\etale$ whose stalks are torsion of orders
invertible in $K$ the base change map
$$
(f')^{-1}Rg_*\mathcal{F}
\longrightarrow
Rh_*e^{-1}\mathcal{F}
$$
is an isomorphism.
\end{lemma}

\begin{proof}
The question is local on $X$, hence we may assume $X$ is affine.
By Limits, Lemma \ref{limits-lemma-relative-approximation}
we can write $X = \lim X_i$ as a cofiltered limit with affine
transition morphisms of schemes $X_i$ of finite type over $K$.
Denote $X'_i =  X_i \times_{\Spec(K)} S'$ and $Y_i = X'_i \times_{S'} T$.
By Lemma \ref{lemma-base-change-Rf-star-colim}
it suffices to prove the statement for the squares
with corners $X_i, Y_i, S_i, T_i$.
Thus we may assume $X$ is of finite type over $K$.
Similarly, we may write
$\mathcal{F} = \colim \mathcal{F}[n]$ where the colimit
is over the multiplicative system of integers invertible in $K$.
The same lemma used above reduces us to the case where
$\mathcal{F}$ is a sheaf of $\mathbf{Z}/n\mathbf{Z}$-modules
for some $n$ invertible in $K$.

\medskip\noindent
We may replace $K$ by its algebraic closure $\overline{K}$.
Namely, formation of direct image commutes with base change
to $\overline{K}$ according to
Lemma \ref{lemma-base-change-field-extension} (works for both
$g$ and $h$). And it suffices to prove the agreement after
restriction to $X'_{\overline{K}}$. Next, we may replace
$X$ by its reduction as we have the topological invariance of \'etale
cohomology, see
Proposition \ref{proposition-topological-invariance}.
After this replacement the morphism $X \to \Spec(K)$
is flat, finite presentation, with geometrically reduced fibres
and the same is true for any base change, in particular for $X' \to S'$.
Hence $(f')^{-1}g_*\mathcal{F} \to Rh_*e^{-1}\mathcal{F}$
is an isomorphism by Lemma \ref{lemma-fppf-reduced-fibres-base-change-f-star}.

\medskip\noindent
At this point we may apply
Lemma \ref{lemma-base-change-does-not-hold-post}
to see that it suffices to prove: given a commutative diagram
$$
\xymatrix{
X \ar[d]_f & X' \ar[d] \ar[l] & Y \ar[l]^h \ar[d] \\
\Spec(K) & S' \ar[l] & \Spec(L) \ar[l]
}
$$
with both squares cartesian, where $S'$ is affine, integral, and normal
with algebraically closed function field $K$, then
$R^qh_*(\mathbf{Z}/d\mathbf{Z})$ is zero for $q > 0$ and $d | n$.
Observe that this vanishing is equivalent to the statement that
$$
(f')^{-1}R^q(\Spec(L) \to S')_*\mathbf{Z}/d\mathbf{Z}
\longrightarrow
R^qh_*\mathbf{Z}/d\mathbf{Z}
$$
is an isomorphism, because the left hand side is zero for example by
Lemma \ref{lemma-Rf-star-zero-normal-with-alg-closed-function-field}.

\medskip\noindent
Write $S' = \Spec(B)$ so that $L$ is the fraction field of $B$.
Write $B = \bigcup_{i \in I} B_i$ as the union of its finite type
$K$-subalgebras $B_i$. Let $J$ be the set of pairs $(i, g)$ where
$i \in I$ and $g \in B_i$ nonzero with ordering
$(i', g') \geq (i, g)$ if and only if $i' \geq i$ and
$g$ maps to an invertible element of $(B_{i'})_{g'}$.
Then $L = \colim_{(i, g) \in J} (B_i)_g$.
For $j = (i, g) \in J$ set $S_j = \Spec(B_i)$
and $U_j = \Spec((B_i)_g)$.
Then
$$
\vcenter{
\xymatrix{
X' \ar[d] & Y \ar[l]^h \ar[d] \\
S' & \Spec(L) \ar[l]
}
}
\quad\text{is the colimit of}\quad
\vcenter{
\xymatrix{
X \times_K S_j \ar[d] & X \times_K U_j \ar[l]^{h_j} \ar[d] \\
S_j & U_j \ar[l]
}
}
$$
Thus we may apply Lemma \ref{lemma-base-change-Rf-star-colim}
to see that it suffices to prove
base change holds in the diagrams on the right which is what we
proved in Lemma \ref{lemma-kunneth-localize-on-X}.
\end{proof}

\begin{lemma}
\label{lemma-punctual-base-change-upgrade}
Let $K$ be a field. Let $n \geq 1$ be invertible in $K$.
Consider a commutative diagram
$$
\xymatrix{
X \ar[d] & X' \ar[l]^p \ar[d]_{f'} & Y \ar[l]^h \ar[d]^e \\
\Spec(K) & S' \ar[l] & T \ar[l]_g
}
$$
of schemes with
$X' = X \times_{\Spec(K)} S'$ and $Y = X' \times_{S'} T$ and
$g$ quasi-compact and quasi-separated. The canonical map
$$
p^{-1}E \otimes_{\mathbf{Z}/n\mathbf{Z}}^\mathbf{L} (f')^{-1}Rg_*F
\longrightarrow
Rh_*(h^{-1}p^{-1}E \otimes_{\mathbf{Z}/n\mathbf{Z}}^\mathbf{L} e^{-1}F)
$$
is an isomorphism if $E$ in $D^+(X_\etale, \mathbf{Z}/n\mathbf{Z})$
has tor amplitude in $[a, \infty]$ for some $a \in \mathbf{Z}$ and
$F$ in $D^+(T_\etale, \mathbf{Z}/n\mathbf{Z})$.
\end{lemma}

\begin{proof}
This lemma is a generalization of Lemma \ref{lemma-punctual-base-change}
to objects of the derived category; the assertion of our lemma is true because
in Lemma \ref{lemma-punctual-base-change} the scheme $X$ over $K$
is arbitrary. We strongly urge the reader to skip the laborious proof
(alternative: read only the last paragraph).

\medskip\noindent
We may represent $E$ by a bounded below K-flat complex
$\mathcal{E}^\bullet$ consisting of flat $\mathbf{Z}/n\mathbf{Z}$-modules.
See Cohomology on Sites, Lemma
\ref{sites-cohomology-lemma-bounded-below-tor-amplitude}.
Choose an integer $b$ such that $H^i(F) = 0$ for $i < b$.
Choose a large integer $N$ and consider the short exact sequence
$$
0 \to \sigma_{\geq N + 1}\mathcal{E}^\bullet \to
\mathcal{E}^\bullet \to
\sigma_{\leq N}\mathcal{E}^\bullet \to 0
$$
of stupid truncations. This produces a distinguished triangle
$E'' \to E \to E' \to E''[1]$ in $D(X_\etale, \mathbf{Z}/n\mathbf{Z})$.
For fixed $F$ both sides of the arrow
in the statement of the lemma are exact functors in $E$. Observe that
$$
p^{-1}E'' \otimes_{\mathbf{Z}/n\mathbf{Z}}^\mathbf{L} (f')^{-1}Rg_*F
\quad\text{and}\quad
Rh_*(h^{-1}p^{-1}E'' \otimes_{\mathbf{Z}/n\mathbf{Z}}^\mathbf{L} e^{-1}F)
$$
are sitting in degrees $\geq N + b$. Hence, if we can prove the lemma
for the object $E'$, then we see that the lemma holds in degrees
$\leq N + b$ and we will conclude. Some details omitted.
Thus we may assume $E$ is represented
by a bounded complex of flat $\mathbf{Z}/n\mathbf{Z}$-modules.
Doing another argument of the same nature, we may assume
$E$ is given by a single flat $\mathbf{Z}/n\mathbf{Z}$-module
$\mathcal{E}$.

\medskip\noindent
Next, we use the same arguments for the variable $F$
to reduce to the case where $F$ is given by a single
sheaf of $\mathbf{Z}/n\mathbf{Z}$-modules $\mathcal{F}$.
Say $\mathcal{F}$ is annihilated by an integer $m | n$.
If $\ell$ is a prime number dividing $m$ and $m > \ell$,
then we can look at the short exact sequence
$0 \to \mathcal{F}[\ell] \to \mathcal{F} \to
\mathcal{F}/\mathcal{F}[\ell] \to 0$
and reduce to smaller $m$. This finally reduces us to
the case where $\mathcal{F}$ is annihilated by a prime
number $\ell$ dividing $n$.
In this case observe that
$$
p^{-1}\mathcal{E}
\otimes_{\mathbf{Z}/n\mathbf{Z}}^\mathbf{L}
(f')^{-1}Rg_*\mathcal{F}
=
p^{-1}(\mathcal{E}/\ell \mathcal{E})
\otimes_{\mathbf{F}_\ell}^\mathbf{L}
(f')^{-1}Rg_*\mathcal{F}
$$
by the flatness of $\mathcal{E}$. Similarly for the other term.
This reduces us to the case where we are working with sheaves
of $\mathbf{F}_\ell$-vector spaces which is discussed

\medskip\noindent
Assume $\ell$ is a prime number invertible in $K$.
Assume $\mathcal{E}$, $\mathcal{F}$ are sheaves of
$\mathbf{F}_\ell$-vector spaces on $X_\etale$ and $T_\etale$.
We want to show that
$$
p^{-1}\mathcal{E} \otimes_{\mathbf{F}_\ell} (f')^{-1}R^qg_*\mathcal{F}
\longrightarrow
R^qh_*(h^{-1}p^{-1}\mathcal{E} \otimes_{\mathbf{F}_\ell} e^{-1}\mathcal{F})
$$
is an isomorphism for every $q \geq 0$. This question is local on $X$
hence we may assume $X$ is affine. We can write $\mathcal{E}$
as a filtered colimit of constructible sheaves of
$\mathbf{F}_\ell$-vector spaces on $X_\etale$, see
Lemma \ref{lemma-torsion-colimit-constructible}.
Since tensor products commute
with filtered colimits and since
higher direct images do too (Lemma \ref{lemma-relative-colimit})
we may assume $\mathcal{E}$ is a constructible sheaf of
$\mathbf{F}_\ell$-vector spaces on $X_\etale$.
Then we can choose an integer $m$ and finite and finitely presented morphisms
$\pi_i : X_i \to X$, $i = 1, \ldots, m$
such that there is an injective map
$$
\mathcal{E} \to
\bigoplus\nolimits_{i = 1, \ldots, m}
\pi_{i, *}\mathbf{F}_\ell
$$
See Lemma \ref{lemma-constructible-maps-into-constant-general}.
Observe that the direct sum is a constructible sheaf as well
(Lemma \ref{lemma-finite-pushforward-constructible}).
Thus the cokernel is constructible too
(Lemma \ref{lemma-constructible-abelian}).
By dimension shifting, i.e., induction on $q$,
on the category of constructible sheaves of
$\mathbf{F}_\ell$-vector spaces on $X_\etale$, it suffices to prove the
result for the sheaves $\pi_{i, *}\mathbf{F}_\ell$
(details omitted; hint: start with proving injectivity for $q = 0$
for all constructible $\mathcal{E}$).
To prove this case we extend the diagram of the lemma to
$$
\xymatrix{
X_i \ar[d]^{\pi_i} &
X'_i \ar[l]^{p_i} \ar[d]^{\pi'_i} &
Y_i \ar[l]^{h_i} \ar[d]^{\rho_i} \\
X \ar[d] & X' \ar[l]^p \ar[d]_{f'} & Y \ar[l]^h \ar[d]^e \\
\Spec(K) & S' \ar[l] & T \ar[l]_g
}
$$
with all squares cartesian. In the equations below we are
going to use that $R\pi_{i, *} = \pi_{i, *}$ and similarly
for $\pi'_i$, $\rho_i$, we are going to use the Leray spectral sequence,
we are going to use
Lemma \ref{lemma-finite-pushforward-commutes-with-base-change}, and
we are going to use
Lemma \ref{lemma-projection-formula-proper-mod-n}
(although this lemma is almost trivial for finite morphisms) for
$\pi_i$, $\pi'_i$, $\rho_i$.
Doing so we see that
\begin{align*}
p^{-1}\pi_{i, *}\mathbf{F}_\ell
\otimes_{\mathbf{F}_\ell} (f')^{-1}R^qg_*\mathcal{F}
& =
\pi'_{i, *}\mathbf{F}_\ell
\otimes_{\mathbf{F}_\ell} (f')^{-1}R^qg_*\mathcal{F} \\
& =
\pi'_{i, *}((\pi'_i)^{-1} (f')^{-1}R^qg_*\mathcal{F})
\end{align*}
Similarly, we have
\begin{align*}
R^qh_*(h^{-1}p^{-1} \pi_{i, *}\mathbf{F}_\ell
\otimes_{\mathbf{F}_\ell} e^{-1}\mathcal{F})
& =
R^qh_*(\rho_{i, *}\mathbf{F}_\ell
\otimes_{\mathbf{F}_\ell} e^{-1}\mathcal{F}) \\
& =
R^qh_*(\rho_i^{-1}e^{-1}\mathcal{F}) \\
& =
\pi'_{i, *}R^qh_{i, *} \rho_i^{-1}e^{-1}\mathcal{F})
\end{align*}
Simce
$R^qh_{i, *} \rho_i^{-1}e^{-1}\mathcal{F} = 
(\pi'_i)^{-1} (f')^{-1}R^qg_*\mathcal{F}$ by
Lemma \ref{lemma-punctual-base-change}
we conclude.
\end{proof}

\begin{lemma}
\label{lemma-punctual-base-change-upgrade-unbounded}
Let $K$ be a field. Let $n \geq 1$ be invertible in $K$.
Consider a commutative diagram
$$
\xymatrix{
X \ar[d] & X' \ar[l]^p \ar[d]_{f'} & Y \ar[l]^h \ar[d]^e \\
\Spec(K) & S' \ar[l] & T \ar[l]_g
}
$$
of schemes of finite type over $K$ with
$X' = X \times_{\Spec(K)} S'$ and $Y = X' \times_{S'} T$.
The canonical map
$$
p^{-1}E \otimes_{\mathbf{Z}/n\mathbf{Z}}^\mathbf{L} (f')^{-1}Rg_*F
\longrightarrow
Rh_*(h^{-1}p^{-1}E \otimes_{\mathbf{Z}/n\mathbf{Z}}^\mathbf{L} e^{-1}F)
$$
is an isomorphism for $E$ in $D(X_\etale, \mathbf{Z}/n\mathbf{Z})$
and $F$ in $D(T_\etale, \mathbf{Z}/n\mathbf{Z})$.
\end{lemma}

\begin{proof}
We will reduce this to
Lemma \ref{lemma-punctual-base-change-upgrade}
using that our functors commute with direct sums.
We suggest the reader skip the proof.
Recall that derived tensor product commutes with
direct sums. Recall that (derived) pullback commutes with direct sums.
Recall that $Rh_*$ and $Rg_*$ commute with direct sums, see
Lemmas \ref{lemma-finite-cd} and
\ref{lemma-finite-cd-mod-n-direct-sums}
(this is where we use our schemes are of finite type
over $K$).

\medskip\noindent
To finish the proof we can argue as follows.
First we write $E = \text{hocolim} \tau_{\leq N} E$.
Since our functors commute with direct sums, they commute
with homotopy colimits. Hence it suffices to prove
the lemma for $E$ bounded above. Similarly for $F$ we
may assume $F$ is bounded above.
Then we can represent $E$ by a bounded above complex
$\mathcal{E}^\bullet$ of sheaves of $\mathbf{Z}/n\mathbf{Z}$-modules.
Then
$$
\mathcal{E}^\bullet = \colim \sigma_{\geq -N}\mathcal{E}^\bullet
$$
(stupid truncations).
Thus we may assume $\mathcal{E}^\bullet$ is a bounded
complex of sheaves of $\mathbf{Z}/n\mathbf{Z}$-modules.
For $F$ we choose a bounded above complex of
flat(!) sheaves of $\mathbf{Z}/n\mathbf{Z}$-modules.
Then we reduce to the case where $F$ is represented
by a bounded complex of flat sheaves of $\mathbf{Z}/n\mathbf{Z}$-modules.
At this point Lemma \ref{lemma-punctual-base-change-upgrade}
kicks in and we conclude.
\end{proof}

\begin{lemma}
\label{lemma-kunneth}
Let $k$ be a separably closed field. Let $X$ and $Y$ be
finite type schemes over $k$. Let $n \geq 1$ be an integer
invertible in $k$. Then for
$E \in D(X_\etale, \mathbf{Z}/n\mathbf{Z})$ and
$K \in D(Y_\etale, \mathbf{Z}/n\mathbf{Z})$
we have
$$
R\Gamma(X \times_{\Spec(k)} Y,
\text{pr}_1^{-1}E
\otimes_{\mathbf{Z}/n\mathbf{Z}}^\mathbf{L}
\text{pr}_2^{-1}K
)
=
R\Gamma(X, E)
\otimes_{\mathbf{Z}/n\mathbf{Z}}^\mathbf{L}
R\Gamma(Y, K)
$$
\end{lemma}

\begin{proof}
By Lemma \ref{lemma-punctual-base-change-upgrade-unbounded} we have
$$
R\text{pr}_{1, *}(
\text{pr}_1^{-1}E
\otimes_{\mathbf{Z}/n\mathbf{Z}}^\mathbf{L}
\text{pr}_2^{-1}K) =
E \otimes_{\mathbf{Z}/n\mathbf{Z}}^\mathbf{L}
\underline{R\Gamma(Y, K)}
$$
We conclude by
Lemma \ref{lemma-pull-out-constant-mod-n}
which we may use because
$\text{cd}(X) < \infty$ by Lemma \ref{lemma-finite-cd}.
\end{proof}














\section{Comparing chaotic and Zariski topologies}
\label{section-compare-chaotic-Zariski}

\noindent
When constructing the structure sheaf of an affine scheme, we first
construct the values on affine opens, and then we extend to all opens.
A similar construction is often useful for constructing complexes
of abelian groups on a scheme $X$. Recall that $X_{affine, Zar}$
denotes the category of affine opens of $X$ with topology given
by standard Zariski coverings, see
Topologies, Definition \ref{topologies-definition-big-small-Zariski}.
We remind the reader that the topos of $X_{affine, Zar}$
is the small Zariski topos of $X$, see Topologies, Lemma
\ref{topologies-lemma-alternative-zariski}.
In this section we denote $X_{affine}$ the same underlying
category with the chaotic topology, i.e., such that sheaves
agree with presheaves. We obtain a morphisms of sites
$$
\epsilon : X_{affine, Zar} \longrightarrow X_{affine}
$$
as in Cohomology on Sites, Section \ref{sites-cohomology-section-compare}.

\begin{lemma}
\label{lemma-check-zar}
In the situation above let $K$ be an object of $D^+(X_{affine})$.
Then $K$ is in the essential image of the (fully faithful) functor
$R\epsilon_* ; D(X_{affine, Zar}) \to D(X_{affine})$ if and only
if the following two conditions hold
\begin{enumerate}
\item $R\Gamma(\emptyset, K)$ is zero in $D(\textit{Ab})$, and
\item if $U = V \cup W$ with $U, V, W \subset X$ affine open and
$V, W \subset U$ standard open
(Algebra, Definition \ref{algebra-definition-Zariski-topology}), then
the map $c^K_{U, V, W, V \cap W}$ of
Cohomology on Sites, Lemma \ref{sites-cohomology-lemma-c-square}
is a quasi-isomorphism.
\end{enumerate}
\end{lemma}

\begin{proof}
(The functor $R\epsilon_*$ is fully faithful by the discussion in
Cohomology on Sites, Section \ref{sites-cohomology-section-compare}.)
Except for a snafu having to do with the empty set,
this follows from the very general Cohomology on Sites, Lemma
\ref{sites-cohomology-lemma-descent-squares} whose hypotheses hold by
Schemes, Lemma \ref{schemes-lemma-sheaf-on-affines} and
Cohomology on Sites, Lemma
\ref{sites-cohomology-lemma-descent-squares-helper}.

\medskip\noindent
To get around the snafu, denote $X_{affine, almost-chaotic}$
the site where the empty object $\emptyset$ has two coverings,
namely, $\{\emptyset \to \emptyset\}$ and the empty covering
(see Sites, Example \ref{sites-example-site-topological} for a
discussion). Then we have morphisms of sites
$$
X_{affine, Zar} \to X_{affine, almost-chaotic} \to X_{affine}
$$
The argument above works for the first arrow. Then we leave it
to the reader to see that an object $K$ of $D^+(X_{affine})$
is in the essential image of the (fully faithful) functor
$D(X_{affine, almost-chaotic}) \to D(X_{affine})$ if and only
if $R\Gamma(\emptyset, K)$ is zero in $D(\textit{Ab})$.
\end{proof}









\section{Comparing big and small topoi}
\label{section-compare}

\noindent
Let $S$ be a scheme. In
Topologies, Lemma \ref{topologies-lemma-at-the-bottom-etale}
we have introduced comparison morphisms
$\pi_S : (\Sch/S)_\etale \to S_\etale$ and
$i_S : \Sh(S_\etale) \to \Sh((\Sch/S)_\etale)$
with $\pi_S \circ i_S = \text{id}$ and $\pi_{S, *} = i_S^{-1}$.
More generally, if $f : T \to S$ is an object of $(\Sch/S)_\etale$,
then there is a morphism $i_f : \Sh(T_\etale) \to \Sh((\Sch/S)_\etale)$
such that $f_{small} = \pi_S \circ i_f$, see
Topologies, Lemmas \ref{topologies-lemma-put-in-T-etale} and
\ref{topologies-lemma-morphism-big-small-etale}. In
Descent, Remark \ref{descent-remark-change-topologies-ringed}
we have extended these to a morphism of ringed sites
$$
\pi_S : ((\Sch/S)_\etale, \mathcal{O}) \to (S_\etale, \mathcal{O}_S)
$$
and morphisms of ringed topoi
$$
i_S : (\Sh(S_\etale), \mathcal{O}_S) \to (\Sh((\Sch/S)_\etale), \mathcal{O})
$$
and
$$
i_f : (\Sh(T_\etale), \mathcal{O}_T) \to (\Sh((\Sch/S)_\etale, \mathcal{O}))
$$
Note that the restriction $i_S^{-1} = \pi_{S, *}$ (see
Topologies, Definition \ref{topologies-definition-restriction-small-etale})
transforms $\mathcal{O}$ into $\mathcal{O}_S$.
Similarly, $i_f^{-1}$ transforms $\mathcal{O}$ into $\mathcal{O}_T$.
See Descent, Remark \ref{descent-remark-change-topologies-ringed}.
Hence $i_S^*\mathcal{F} = i_S^{-1}\mathcal{F}$ and
$i_f^*\mathcal{F} = i_f^{-1}\mathcal{F}$ for any $\mathcal{O}$-module
$\mathcal{F}$ on $(\Sch/S)_\etale$. In particular $i_S^*$ and $i_f^*$
are exact functors. The functor $i_S^*$ is often denoted
$\mathcal{F} \mapsto \mathcal{F}|_{S_\etale}$ (and this does not
conflict with the notation in
Topologies, Definition \ref{topologies-definition-restriction-small-etale}).

\begin{lemma}
\label{lemma-compare-injectives}
Let $S$ be a scheme. Let $T$ be an object of $(\Sch/S)_\etale$.
\begin{enumerate}
\item If $\mathcal{I}$ is injective in $\textit{Ab}((\Sch/S)_\etale)$, then
\begin{enumerate}
\item $i_f^{-1}\mathcal{I}$ is injective in $\textit{Ab}(T_\etale)$,
\item $\mathcal{I}|_{S_\etale}$ is injective in $\textit{Ab}(S_\etale)$,
\end{enumerate}
\item If $\mathcal{I}^\bullet$ is a K-injective complex
in $\textit{Ab}((\Sch/S)_\etale)$, then
\begin{enumerate}
\item $i_f^{-1}\mathcal{I}^\bullet$ is a K-injective complex in
$\textit{Ab}(T_\etale)$,
\item $\mathcal{I}^\bullet|_{S_\etale}$ is a K-injective complex in
$\textit{Ab}(S_\etale)$,
\end{enumerate}
\end{enumerate}
The corresponding statements for modules do not hold.
\end{lemma}

\begin{proof}
Parts (1)(b) and (2)(b)
follow formally from the fact that the restriction functor
$\pi_{S, *} = i_S^{-1}$ is a right adjoint of the exact functor
$\pi_S^{-1}$, see
Homology, Lemma \ref{homology-lemma-adjoint-preserve-injectives} and
Derived Categories, Lemma \ref{derived-lemma-adjoint-preserve-K-injectives}.

\medskip\noindent
Parts (1)(a) and (2)(a) can be seen in two ways. First proof: We can use
that $i_f^{-1}$ is a right adjoint of the exact functor $i_{f, !}$.
This functor is constructed in
Topologies, Lemma \ref{topologies-lemma-put-in-T-etale}
for sheaves of sets and for abelian sheaves in
Modules on Sites, Lemma \ref{sites-modules-lemma-g-shriek-adjoint}.
It is shown in Modules on Sites, Lemma
\ref{sites-modules-lemma-exactness-lower-shriek} that it is exact.
Second proof. We can use that $i_f = i_T \circ f_{big}$ as is shown
in Topologies, Lemma \ref{topologies-lemma-morphism-big-small-etale}.
Since $f_{big}$ is a localization, we see that pullback by it
preserves injectives and K-injectives, see
Cohomology on Sites, Lemmas \ref{sites-cohomology-lemma-cohomology-of-open} and
\ref{sites-cohomology-lemma-restrict-K-injective-to-open}.
Then we apply the already proved parts (1)(b) and (2)(b)
to the functor $i_T^{-1}$ to conclude.

\medskip\noindent
Let $S = \Spec(\mathbf{Z})$ and consider the map
$2 : \mathcal{O}_S \to \mathcal{O}_S$. This is an injective map
of $\mathcal{O}_S$-modules on $S_\etale$. However, the pullback
$\pi_S^*(2) : \mathcal{O} \to \mathcal{O}$ is not injective
as we see by evaluating on $\Spec(\mathbf{F}_2)$. Now choose
an injection $\alpha : \mathcal{O} \to \mathcal{I}$ into an injective
$\mathcal{O}$-module $\mathcal{I}$ on $(\Sch/S)_\etale$.
Then consider the diagram
$$
\xymatrix{
\mathcal{O}_S \ar[d]_2 \ar[rr]_{\alpha|_{S_\etale}} & &
\mathcal{I}|_{S_\etale} \\
\mathcal{O}_S \ar@{..>}[rru]
}
$$
Then the dotted arrow cannot exist in the category of $\mathcal{O}_S$-modules
because it would mean
(by adjunction) that the injective map $\alpha$ factors through
the noninjective map $\pi_S^*(2)$ which cannot be the case.
Thus $\mathcal{I}|_{S_\etale}$ is not an injective $\mathcal{O}_S$-module.
\end{proof}

\noindent
Let $f : T \to S$ be a morphism of schemes. The commutative diagram of
Topologies, Lemma \ref{topologies-lemma-morphism-big-small-etale} (3)
leads to a commutative diagram of ringed sites
$$
\xymatrix{
(T_\etale, \mathcal{O}_T) \ar[d]_{f_{small}} &
((\Sch/T)_\etale, \mathcal{O}) \ar[d]^{f_{big}} \ar[l]^{\pi_T} \\
(S_\etale, \mathcal{O}_S) &
((\Sch/S)_\etale, \mathcal{O}) \ar[l]_{\pi_S}
}
$$
as one easily sees by writing out the definitions of
$f_{small}^\sharp$, $f_{big}^\sharp$, $\pi_S^\sharp$, and $\pi_T^\sharp$.
In particular this means that
\begin{equation}
\label{equation-compare-big-small}
(f_{big, *}\mathcal{F})|_{S_\etale} =
f_{small, *}(\mathcal{F}|_{T_\etale})
\end{equation}
for any sheaf $\mathcal{F}$ on $(\Sch/T)_\etale$ and if $\mathcal{F}$
is a sheaf of $\mathcal{O}$-modules, then (\ref{equation-compare-big-small})
is an isomorphism of $\mathcal{O}_S$-modules on $S_\etale$.

\begin{lemma}
\label{lemma-compare-higher-direct-image}
Let $f : T \to S$ be a morphism of schemes.
\begin{enumerate}
\item For $K$ in $D((\Sch/T)_\etale)$ we have
$
(Rf_{big, *}K)|_{S_\etale} = Rf_{small, *}(K|_{T_\etale})
$
in $D(S_\etale)$.
\item For $K$ in $D((\Sch/T)_\etale, \mathcal{O})$ we have
$
(Rf_{big, *}K)|_{S_\etale} = Rf_{small, *}(K|_{T_\etale})
$
in $D(\textit{Mod}(S_\etale, \mathcal{O}_S))$.
\end{enumerate}
More generally, let $g : S' \to S$ be an object of $(\Sch/S)_\etale$.
Consider the fibre product
$$
\xymatrix{
T' \ar[r]_{g'} \ar[d]_{f'} & T \ar[d]^f \\
S' \ar[r]^g & S
}
$$
Then
\begin{enumerate}
\item[(3)] For $K$ in $D((\Sch/T)_\etale)$ we have
$i_g^{-1}(Rf_{big, *}K) = Rf'_{small, *}(i_{g'}^{-1}K)$
in $D(S'_\etale)$.
\item[(4)] For $K$ in $D((\Sch/T)_\etale, \mathcal{O})$ we have
$i_g^*(Rf_{big, *}K) = Rf'_{small, *}(i_{g'}^*K)$
in $D(\textit{Mod}(S'_\etale, \mathcal{O}_{S'}))$.
\item[(5)] For $K$ in $D((\Sch/T)_\etale)$ we have
$g_{big}^{-1}(Rf_{big, *}K) = Rf'_{big, *}((g'_{big})^{-1}K)$
in $D((\Sch/S')_\etale)$.
\item[(6)] For $K$ in $D((\Sch/T)_\etale, \mathcal{O})$ we have
$g_{big}^*(Rf_{big, *}K) = Rf'_{big, *}((g'_{big})^*K)$
in $D(\textit{Mod}(S'_\etale, \mathcal{O}_{S'}))$.
\end{enumerate}
\end{lemma}

\begin{proof}
Part (1) follows from
Lemma \ref{lemma-compare-injectives}
and (\ref{equation-compare-big-small})
on choosing a K-injective complex of abelian sheaves representing $K$.

\medskip\noindent
Part (3) follows from Lemma \ref{lemma-compare-injectives}
and Topologies, Lemma
\ref{topologies-lemma-morphism-big-small-cartesian-diagram-etale}
on choosing a K-injective complex of abelian sheaves representing $K$.

\medskip\noindent
Part (5) is Cohomology on Sites, Lemma
\ref{sites-cohomology-lemma-localize-cartesian-square}.

\medskip\noindent
Part (6) is Cohomology on Sites, Lemma
\ref{sites-cohomology-lemma-localize-cartesian-square-modules}.

\medskip\noindent
Part (2) can be proved as follows. Above we have seen
that $\pi_S \circ f_{big} = f_{small} \circ \pi_T$ as morphisms
of ringed sites. Hence we obtain
$R\pi_{S, *} \circ Rf_{big, *} = Rf_{small, *} \circ R\pi_{T, *}$
by Cohomology on Sites, Lemma
\ref{sites-cohomology-lemma-derived-pushforward-composition}.
Since the restriction functors $\pi_{S, *}$ and $\pi_{T, *}$
are exact, we conclude.

\medskip\noindent
Part (4) follows from part (6) and part (2) applied to $f' : T' \to S'$.
\end{proof}

\noindent
Let $S$ be a scheme and let $\mathcal{H}$ be an abelian sheaf on
$(\Sch/S)_\etale$. Recall that $H^n_\etale(U, \mathcal{H})$ denotes
the cohomology of $\mathcal{H}$ over an object $U$ of $(\Sch/S)_\etale$.

\begin{lemma}
\label{lemma-compare-cohomology}
Let $f : T \to S$ be a morphism of schemes. Then
\begin{enumerate}
\item For $K$ in $D(S_\etale)$ we have
$H^n_\etale(S, \pi_S^{-1}K) = H^n(S_\etale, K)$.
\item For $K$ in $D(S_\etale, \mathcal{O}_S)$ we have
$H^n_\etale(S, L\pi_S^*K) = H^n(S_\etale, K)$.
\item For $K$ in $D(S_\etale)$ we have
$H^n_\etale(T, \pi_S^{-1}K) = H^n(T_\etale, f_{small}^{-1}K)$.
\item For $K$ in $D(S_\etale, \mathcal{O}_S)$ we have
$H^n_\etale(T, L\pi_S^*K) = H^n(T_\etale, Lf_{small}^*K)$.
\item For $M$ in $D((\Sch/S)_\etale)$ we have
$H^n_\etale(T, M) = H^n(T_\etale, i_f^{-1}M)$.
\item For $M$ in $D((\Sch/S)_\etale, \mathcal{O})$ we have
$H^n_\etale(T, M) = H^n(T_\etale, i_f^*M)$.
\end{enumerate}
\end{lemma}

\begin{proof}
To prove (5) represent $M$ by a K-injective complex of abelian sheaves
and apply Lemma \ref{lemma-compare-injectives}
and work out the definitions. Part (3) follows from
this as $i_f^{-1}\pi_S^{-1} = f_{small}^{-1}$. Part (1) is a special
case of (3).

\medskip\noindent
Part (6) follows from the very general Cohomology on Sites, Lemma
\ref{sites-cohomology-lemma-pullback-same-cohomology}. Then part
(4) follows because $Lf_{small}^* = i_f^* \circ L\pi_S^*$.
Part (2) is a special case of (4).
\end{proof}

\begin{lemma}
\label{lemma-cohomological-descent-etale}
Let $S$ be a scheme. For $K \in D(S_\etale)$ the map
$$
K \longrightarrow R\pi_{S, *}\pi_S^{-1}K
$$
is an isomorphism.
\end{lemma}

\begin{proof}
This is true because both $\pi_S^{-1}$ and $\pi_{S, *} = i_S^{-1}$
are exact functors and the composition $\pi_{S, *} \circ \pi_S^{-1}$
is the identity functor.
\end{proof}

\begin{lemma}
\label{lemma-compare-higher-direct-image-proper}
Let $f : T \to S$ be a proper morphism of schemes. Then we have
\begin{enumerate}
\item $\pi_S^{-1} \circ f_{small, *} = f_{big, *} \circ \pi_T^{-1}$
as functors $\Sh(T_\etale) \to \Sh((\Sch/S)_\etale)$,
\item $\pi_S^{-1}Rf_{small, *}K = Rf_{big, *}\pi_T^{-1}K$
for $K$ in $D^+(T_\etale)$ whose cohomology sheaves are torsion,
\item $\pi_S^{-1}Rf_{small, *}K = Rf_{big, *}\pi_T^{-1}K$
for $K$ in $D(T_\etale, \mathbf{Z}/n\mathbf{Z})$, and
\item $\pi_S^{-1}Rf_{small, *}K = Rf_{big, *}\pi_T^{-1}K$
for all $K$ in $D(T_\etale)$ if $f$ is finite.
\end{enumerate}
\end{lemma}

\begin{proof}
Proof of (1). Let $\mathcal{F}$ be a sheaf on $T_\etale$.
Let $g : S' \to S$ be an object of $(\Sch/S)_\etale$. Consider the
fibre product
$$
\xymatrix{
T' \ar[r]_{f'} \ar[d]_{g'} & S' \ar[d]^g \\
T \ar[r]^f & S
}
$$
Then we have
$$
(f_{big, *}\pi_T^{-1}\mathcal{F})(S') =
(\pi_T^{-1}\mathcal{F})(T') =
((g'_{small})^{-1}\mathcal{F})(T')  =
(f'_{small, *}(g'_{small})^{-1}\mathcal{F})(S')
$$
the second equality by Lemma \ref{lemma-describe-pullback}.
On the other hand
$$
(\pi_S^{-1}f_{small, *}\mathcal{F})(S') =
(g_{small}^{-1}f_{small, *}\mathcal{F})(S')
$$
again by Lemma \ref{lemma-describe-pullback}.
Hence by proper base change for sheaves of sets
(Lemma \ref{lemma-proper-base-change-f-star})
we conclude the two sets are canonically isomorphic.
The isomorphism is compatible with restriction mappings
and defines an isomorphism
$\pi_S^{-1}f_{small, *}\mathcal{F} = f_{big, *}\pi_T^{-1}\mathcal{F}$.
Thus an isomorphism of functors
$\pi_S^{-1} \circ f_{small, *} = f_{big, *} \circ \pi_T^{-1}$.

\medskip\noindent
Proof of (2). There is a canonical base change map
$\pi_S^{-1}Rf_{small, *}K \to Rf_{big, *}\pi_T^{-1}K$
for any $K$ in $D(T_\etale)$, see
Cohomology on Sites, Remark \ref{sites-cohomology-remark-base-change}.
To prove it is an isomorphism, it suffices to prove the pull back of
the base change map by $i_g : \Sh(S'_\etale) \to \Sh((\Sch/S)_\etale)$
is an isomorphism for any object $g : S' \to S$ of $(\Sch/S)_\etale$.
Let $T', g', f'$ be as in the previous paragraph.
The pullback of the base change map is
\begin{align*}
g_{small}^{-1}Rf_{small, *}K
& =
i_g^{-1}\pi_S^{-1}Rf_{small, *}K \\
& \to
i_g^{-1}Rf_{big, *}\pi_T^{-1}K \\
& =
Rf'_{small, *}(i_{g'}^{-1}\pi_T^{-1}K) \\
& =
Rf'_{small, *}((g'_{small})^{-1}K)
\end{align*}
where we have used $\pi_S \circ i_g = g_{small}$,
$\pi_T \circ i_{g'} = g'_{small}$, and
Lemma \ref{lemma-compare-higher-direct-image}.
This map is an isomorphism by the proper base change theorem
(Lemma \ref{lemma-proper-base-change}) provided $K$ is bounded
below and the cohomology sheaves of $K$ are torsion.

\medskip\noindent
The proof of part (3) is the same as the proof of part (2), except
we use Lemma \ref{lemma-proper-base-change-mod-n}
instead of Lemma \ref{lemma-proper-base-change}.

\medskip\noindent
Proof of (4). If $f$ is finite, then the functors
$f_{small, *}$ and $f_{big, *}$ are exact. This follows
from Proposition \ref{proposition-finite-higher-direct-image-zero}
for $f_{small}$. Since any base change $f'$ of $f$ is finite too,
we conclude from Lemma \ref{lemma-compare-higher-direct-image} part (3)
that $f_{big, *}$ is exact too (as the higher derived functors are zero).
Thus this case follows from part (1).
\end{proof}




\section{Comparing fppf and \'etale topologies}
\label{section-fppf-etale}

\noindent
A model for this section is the section on the comparison of the
usual topology and the qc topology on locally compact topological
spaces as discussed in Cohomology on Sites, Section
\ref{sites-cohomology-section-cohomology-LC}.
We first review some material from
Topologies, Sections
\ref{topologies-section-change-topologies} and
\ref{topologies-section-etale}.

\medskip\noindent
Let $S$ be a scheme and let $(\Sch/S)_{fppf}$ be an fppf site.
On the same underlying category we have a second topology,
namely the \'etale topology, and hence a second site
$(\Sch/S)_\etale$. The identity functor
$(\Sch/S)_\etale \to (\Sch/S)_{fppf}$ is continuous and defines
a morphism of sites
$$
\epsilon_S : (\Sch/S)_{fppf} \longrightarrow (\Sch/S)_\etale
$$
See Cohomology on Sites, Section \ref{sites-cohomology-section-compare}.
Please note that $\epsilon_{S, *}$ is the identity functor on underlying
presheaves and that $\epsilon_S^{-1}$ associates to an \'etale sheaf the
fppf sheafification. Let $S_\etale$ be the small \'etale site.
There is a morphism of sites
$$
\pi_S : (\Sch/S)_\etale \longrightarrow S_\etale
$$
given by the continuous functor
$S_\etale \to (\Sch/S)_\etale$, $U \mapsto U$.
Namely, $S_\etale$ has fibre products and a final object and the
functor above commutes with these and
Sites, Proposition \ref{sites-proposition-get-morphism} applies.

\begin{lemma}
\label{lemma-describe-pullback-pi-fppf}
With notation as above.
Let $\mathcal{F}$ be a sheaf on $S_\etale$. The rule
$$
(\Sch/S)_{fppf} \longrightarrow \textit{Sets},\quad
(f : X \to S) \longmapsto \Gamma(X, f_{small}^{-1}\mathcal{F})
$$
is a sheaf and a fortiori a sheaf on $(\Sch/S)_\etale$.
In fact this sheaf is equal to
$\pi_S^{-1}\mathcal{F}$ on $(\Sch/S)_\etale$ and
$\epsilon_S^{-1}\pi_S^{-1}\mathcal{F}$ on $(\Sch/S)_{fppf}$.
\end{lemma}

\begin{proof}
The statement about the \'etale topology is the content
of Lemma \ref{lemma-describe-pullback}. To finish the proof it
suffices to show that $\pi_S^{-1}\mathcal{F}$ is a sheaf for the fppf
topology. This is shown in Lemma \ref{lemma-describe-pullback}
as well.
\end{proof}

\noindent
In the situation of Lemma \ref{lemma-describe-pullback-pi-fppf}
the composition of $\epsilon_S$ and $\pi_S$ and the equality
determine a morphism of sites
$$
a_S : (\Sch/S)_{fppf} \longrightarrow S_\etale
$$

\begin{lemma}
\label{lemma-push-pull-fppf-etale}
With notation as above.
Let $f : X \to Y$ be a morphism of $(\Sch/S)_{fppf}$.
Then there are commutative diagrams of topoi
$$
\xymatrix{
\Sh((\Sch/X)_{fppf}) \ar[rr]_{f_{big, fppf}} \ar[d]_{\epsilon_X} & &
\Sh((\Sch/Y)_{fppf}) \ar[d]^{\epsilon_Y} \\
\Sh((\Sch/X)_\etale) \ar[rr]^{f_{big, \etale}} & &
\Sh((\Sch/Y)_\etale)
}
$$
and
$$
\xymatrix{
\Sh((\Sch/X)_{fppf}) \ar[rr]_{f_{big, fppf}} \ar[d]_{a_X} & &
\Sh((\Sch/Y)_{fppf}) \ar[d]^{a_Y} \\
\Sh(X_\etale) \ar[rr]^{f_{small}} & &
\Sh(Y_\etale)
}
$$
with $a_X = \pi_X \circ \epsilon_X$ and $a_Y = \pi_X \circ \epsilon_X$.
\end{lemma}

\begin{proof}
The commutativity of the diagrams follows from the discussion in
Topologies, Section \ref{topologies-section-change-topologies}.
\end{proof}

\begin{lemma}
\label{lemma-proper-push-pull-fppf-etale}
In Lemma \ref{lemma-push-pull-fppf-etale} if $f$ is proper, then we have
$a_Y^{-1} \circ f_{small, *} = f_{big, fppf, *} \circ a_X^{-1}$.
\end{lemma}

\begin{proof}
You can prove this by repeating the proof of
Lemma \ref{lemma-compare-higher-direct-image-proper} part (1);
we will instead deduce the result from this.
As $\epsilon_{Y, *}$ is the identity functor on underlying presheaves,
it reflects isomorphisms. The description
in Lemma \ref{lemma-describe-pullback-pi-fppf}
shows that $\epsilon_{Y, *} \circ a_Y^{-1} = \pi_Y^{-1}$
and similarly for $X$. To show that the canonical map
$a_Y^{-1}f_{small, *}\mathcal{F} \to f_{big, fppf, *}a_X^{-1}\mathcal{F}$
is an isomorphism, it suffices to show that
\begin{align*}
\pi_Y^{-1}f_{small, *}\mathcal{F}
& =
\epsilon_{Y, *}a_Y^{-1}f_{small, *}\mathcal{F} \\
& \to 
\epsilon_{Y, *}f_{big, fppf, *}a_X^{-1}\mathcal{F} \\
& =
f_{big, \etale, *} \epsilon_{X, *}a_X^{-1}\mathcal{F} \\
& =
f_{big, \etale, *}\pi_X^{-1}\mathcal{F}
\end{align*}
is an isomorphism. This is part
(1) of Lemma \ref{lemma-compare-higher-direct-image-proper}.
\end{proof}

\begin{lemma}
\label{lemma-descent-sheaf-fppf-etale}
In Lemma \ref{lemma-push-pull-fppf-etale} assume
$f$ is flat, locally of finite presentation, and surjective.
Then the functor
$$
\Sh(Y_\etale) \longrightarrow
\left\{
(\mathcal{G}, \mathcal{H}, \alpha)
\middle|
\begin{matrix}
\mathcal{G} \in \Sh(X_\etale),\ \mathcal{H} \in \Sh((\Sch/Y)_{fppf}), \\
\alpha : a_X^{-1}\mathcal{G} \to f_{big, fppf}^{-1}\mathcal{H}
\text{ an isomorphism}
\end{matrix}
\right\}
$$
sending $\mathcal{F}$ to
$(f_{small}^{-1}\mathcal{F}, a_Y^{-1}\mathcal{F}, can)$ is an equivalence.
\end{lemma}

\begin{proof}
The functor $a_X^{-1}$ is fully faithful (as $a_{X, *}a_X^{-1} = \text{id}$ by
Lemma \ref{lemma-describe-pullback-pi-fppf}). Hence the forgetful functor
$(\mathcal{G}, \mathcal{H}, \alpha) \mapsto \mathcal{H}$ identifies the
category of triples with a full subcategory of $\Sh((\Sch/Y)_{fppf})$.
Moreover, the functor $a_Y^{-1}$ is fully faithful, hence the functor
in the lemma is fully faithful as well.

\medskip\noindent
Suppose that we have an \'etale covering $\{Y_i \to Y\}$.
Let $f_i : X_i \to Y_i$ be the base change of $f$.
Denote $f_{ij} = f_i \times f_j : X_i \times_X X_j  \to Y_i \times_Y Y_j$.
Claim: if the lemma is true for $f_i$ and $f_{ij}$ for all $i, j$, then
the lemma is true for $f$. To see this, note that the given \'etale covering
determines an \'etale covering of the final object in each of
the four sites $Y_\etale, X_\etale, (\Sch/Y)_{fppf}, (\Sch/X)_{fppf}$.
Thus the category of sheaves is equivalent to the category of
glueing data for this covering
(Sites, Lemma \ref{sites-lemma-mapping-property-glue})
in each of the four cases. A huge commutative diagram of
categories then finishes the proof of the claim. We omit the details.
The claim shows that we may work \'etale locally on $Y$.

\medskip\noindent
Note that $\{X \to Y\}$ is an fppf covering. Working \'etale locally on $Y$,
we may assume there exists a morphism $s : X' \to X$ such that the composition
$f' = f \circ s : X' \to Y$ is surjective finite locally free, see
More on Morphisms, Lemma
\ref{more-morphisms-lemma-dominate-fppf-etale-locally}.
Claim: if the lemma is true for $f'$, then it is true for $f$.
Namely, given a triple $(\mathcal{G}, \mathcal{H}, \alpha)$
for $f$, we can pullback by $s$ to get a triple
$(s_{small}^{-1}\mathcal{G}, \mathcal{H}, s_{big, fppf}^{-1}\alpha)$
for $f'$. A solution for this triple gives a sheaf
$\mathcal{F}$ on $Y_\etale$ with $a_Y^{-1}\mathcal{F} = \mathcal{H}$.
By the first paragraph of the proof this means the triple is
in the essential image. This reduces us to
the case described in the next paragraph.

\medskip\noindent
Assume $f$ is surjective finite locally free.
Let $(\mathcal{G}, \mathcal{H}, \alpha)$ be a triple.
In this case consider the triple
$$
(\mathcal{G}_1, \mathcal{H}_1, \alpha_1) =
(f_{small}^{-1}f_{small, *}\mathcal{G},
f_{big, fppf, *}f_{big, fppf}^{-1}\mathcal{H}, \alpha_1)
$$
where $\alpha_1$ comes from the identifications
\begin{align*}
a_X^{-1}f_{small}^{-1}f_{small, *}\mathcal{G}
& =
f_{big, fppf}^{-1}a_Y^{-1}f_{small, *}\mathcal{G} \\
& =
f_{big, fppf}^{-1}f_{big, fppf, *}a_X^{-1}\mathcal{G} \\
& \to
f_{big, fppf}^{-1}f_{big, fppf, *}f_{big, fppf}^{-1}\mathcal{H}
\end{align*}
where the third equality is Lemma \ref{lemma-proper-push-pull-fppf-etale}
and the arrow is given by $\alpha$.
This triple is in the image of our functor because
$\mathcal{F}_1 = f_{small, *}\mathcal{F}$ is a solution
(to see this use Lemma \ref{lemma-proper-push-pull-fppf-etale} again;
details omitted). There is a canonical map of triples
$$
(\mathcal{G}, \mathcal{H}, \alpha)
\to
(\mathcal{G}_1, \mathcal{H}_1, \alpha_1)
$$
which uses the unit $\text{id} \to f_{big, fppf, *}f_{big, fppf}^{-1}$
on the second entry (it is enough to prescribe morphisms on the
second entry by the first paragraph of the proof). Since
$\{f : X \to Y\}$ is an fppf covering the map
$\mathcal{H} \to \mathcal{H}_1$ is injective (details omitted).
Set
$$
\mathcal{G}_2 = \mathcal{G}_1 \amalg_\mathcal{G} \mathcal{G}_1\quad
\mathcal{H}_2 = \mathcal{H}_1 \amalg_\mathcal{H} \mathcal{H}_1
$$
and let $\alpha_2$ be the induced isomorphism (pullback functors
are exact, so this makes sense). Then $\mathcal{H}$ is the
equalizer of the two maps $\mathcal{H}_1 \to \mathcal{H}_2$.
Repeating the arguments above for the triple
$(\mathcal{G}_2, \mathcal{H}_2, \alpha_2)$
we find an injective morphism of triples
$$
(\mathcal{G}_2, \mathcal{H}_2, \alpha_2)
\to
(\mathcal{G}_3, \mathcal{H}_3, \alpha_3)
$$
such that this last triple is in the image of our functor.
Say it corresponds to $\mathcal{F}_3$ in $\Sh(Y_\etale)$.
By fully faithfulness we obtain two maps
$\mathcal{F}_1 \to \mathcal{F}_3$ and we can let
$\mathcal{F}$ be the equalizer of these two maps.
By exactness of the pullback functors involved we
find that $a_Y^{-1}\mathcal{F} = \mathcal{H}$ as desired.
\end{proof}

\begin{lemma}
\label{lemma-compare-fppf-etale}
Consider the comparison morphism
$\epsilon : (\Sch/S)_{fppf} \to (\Sch/S)_\etale$.
Let $\mathcal{P}$ denote the class of finite morphisms of schemes.
For $X$ in $(\Sch/S)_\etale$ denote
$\mathcal{A}'_X \subset \textit{Ab}((\Sch/X)_\etale)$
the full subcategory consisting of sheaves of the form
$\pi_X^{-1}\mathcal{F}$ with $\mathcal{F}$ in $\textit{Ab}(X_\etale)$.
Then Cohomology on Sites, Properties
(\ref{sites-cohomology-item-base-change-P}),
(\ref{sites-cohomology-item-restriction-A}),
(\ref{sites-cohomology-item-A-sheaf}),
(\ref{sites-cohomology-item-A-and-P}), and
(\ref{sites-cohomology-item-refine-tau-by-P})
of Cohomology on Sites, Situation
\ref{sites-cohomology-situation-compare} hold.
\end{lemma}

\begin{proof}
We first show that $\mathcal{A}'_X \subset \textit{Ab}((\Sch/X)_\etale)$
is a weak Serre subcategory by checking conditions (1), (2), (3), and (4)
of Homology, Lemma \ref{homology-lemma-characterize-weak-serre-subcategory}.
Parts (1), (2), (3) are immediate as $\pi_X^{-1}$ is exact and
fully faithful for example by Lemma \ref{lemma-cohomological-descent-etale}. If
$0 \to \pi_X^{-1}\mathcal{F} \to \mathcal{G} \to \pi_X^{-1}\mathcal{F}' \to 0$
is a short exact sequence in $\textit{Ab}((\Sch/X)_\etale)$
then $0 \to \mathcal{F} \to \pi_{X, *}\mathcal{G} \to \mathcal{F}' \to 0$
is exact by Lemma \ref{lemma-cohomological-descent-etale}.
Hence $\mathcal{G} = \pi_X^{-1}\pi_{X, *}\mathcal{G}$ is in
$\mathcal{A}'_X$ which checks the final condition.

\medskip\noindent
Cohomology on Sites, Property (\ref{sites-cohomology-item-base-change-P}) holds
by the existence of fibre products of schemes
and the fact that the base change of a finite morphism of
schemes is a finite morphism of schemes, see
Morphisms, Lemma \ref{morphisms-lemma-base-change-finite}.

\medskip\noindent
Cohomology on Sites, Property (\ref{sites-cohomology-item-restriction-A})
follows from the commutative diagram (3) in
Topologies, Lemma \ref{topologies-lemma-morphism-big-small-etale}.

\medskip\noindent
Cohomology on Sites, Property (\ref{sites-cohomology-item-A-sheaf}) is
Lemma \ref{lemma-describe-pullback-pi-fppf}.

\medskip\noindent
Cohomology on Sites, Property (\ref{sites-cohomology-item-A-and-P}) holds by
Lemma \ref{lemma-compare-higher-direct-image-proper} part (4).

\medskip\noindent
Cohomology on Sites, Property (\ref{sites-cohomology-item-refine-tau-by-P})
is implied by
More on Morphisms, Lemma
\ref{more-morphisms-lemma-dominate-fppf-etale-locally}.
\end{proof}

\begin{lemma}
\label{lemma-V-C-all-n-etale-fppf}
With notation as above.
\begin{enumerate}
\item For $X \in \Ob((\Sch/S)_{fppf})$ and an abelian sheaf $\mathcal{F}$
on $X_\etale$ we have
$\epsilon_{X, *}a_X^{-1}\mathcal{F} = \pi_X^{-1}\mathcal{F}$
and $R^i\epsilon_{X, *}(a_X^{-1}\mathcal{F}) = 0$ for $i > 0$.
\item For a finite morphism $f : X \to Y$ in $(\Sch/S)_{fppf}$
and abelian sheaf $\mathcal{F}$ on $X$ we have
$a_Y^{-1}(R^if_{small, *}\mathcal{F}) =
R^if_{big, fppf, *}(a_X^{-1}\mathcal{F})$
for all $i$.
\item For a scheme $X$ and $K$ in $D^+(X_\etale)$ the map
$\pi_X^{-1}K \to R\epsilon_{X, *}(a_X^{-1}K)$ is an isomorphism.
\item For a finite morphism $f : X \to Y$ of schemes
and $K$ in $D^+(X_\etale)$ we have
$a_Y^{-1}(Rf_{small, *}K) = Rf_{big, fppf, *}(a_X^{-1}K)$.
\item For a proper morphism $f : X \to Y$ of schemes
and $K$ in $D^+(X_\etale)$ with torsion cohomology sheaves we have
$a_Y^{-1}(Rf_{small, *}K) = Rf_{big, fppf, *}(a_X^{-1}K)$.
\end{enumerate}
\end{lemma}

\begin{proof}
By Lemma \ref{lemma-compare-fppf-etale} the lemmas in
Cohomology on Sites, Section \ref{sites-cohomology-section-compare-general}
all apply to our current setting. To translate the results
observe that the category $\mathcal{A}_X$ of
Cohomology on Sites, Lemma \ref{sites-cohomology-lemma-A}
is the essential image of
$a_X^{-1} : \textit{Ab}(X_\etale) \to \textit{Ab}((\Sch/X)_{fppf})$.

\medskip\noindent
Part (1) is equivalent to $(V_n)$ for all $n$ which holds by
Cohomology on Sites, Lemma \ref{sites-cohomology-lemma-V-C-all-n-general}.

\medskip\noindent
Part (2) follows by applying $\epsilon_Y^{-1}$ to the conclusion of
Cohomology on Sites, Lemma \ref{sites-cohomology-lemma-V-implies-C-general}.

\medskip\noindent
Part (3) follows from Cohomology on Sites, Lemma
\ref{sites-cohomology-lemma-V-C-all-n-general} part (1)
because $\pi_X^{-1}K$ is in $D^+_{\mathcal{A}'_X}((\Sch/X)_\etale)$
and $a_X^{-1} = \epsilon_X^{-1} \circ a_X^{-1}$.

\medskip\noindent
Part (4) follows from Cohomology on Sites, Lemma
\ref{sites-cohomology-lemma-V-C-all-n-general} part (2)
for the same reason.

\medskip\noindent
Part (5). We use that
\begin{align*}
R\epsilon_{Y, *}Rf_{big, fppf, *}a_X^{-1}K
& =
Rf_{big, \etale, *}R\epsilon_{X, *}a_X^{-1}K \\
& =
Rf_{big, \etale, *}\pi_X^{-1}K \\
& =
\pi_Y^{-1}Rf_{small, *}K \\
& =
R\epsilon_{Y, *} a_Y^{-1}Rf_{small, *}K
\end{align*}
The first equality by the commutative diagram in
Lemma \ref{lemma-push-pull-fppf-etale}
and Cohomology on Sites, Lemma
\ref{sites-cohomology-lemma-derived-pushforward-composition}.
The second equality is (3). The third is
Lemma \ref{lemma-compare-higher-direct-image-proper} part (2).
The fourth is (3) again. Thus the base change map
$a_Y^{-1}(Rf_{small, *}K) \to Rf_{big, fppf, *}(a_X^{-1}K)$
induces an isomorphism
$$
R\epsilon_{Y, *}a_Y^{-1}Rf_{small, *}K \to
R\epsilon_{Y, *}Rf_{big, fppf, *}a_X^{-1}K
$$
The proof is finished by the following remark: a map
$\alpha : a_Y^{-1}L \to M$ with $L$ in $D^+(Y_\etale)$
and $M$ in $D^+((\Sch/Y)_{fppf})$ such that $R\epsilon_{Y, *}\alpha$
is an isomorphism, is an isomorphism. Namely, 
we show by induction on $i$ that $H^i(\alpha)$ is an isomorphism.
This is true for all sufficiently small $i$.
If it holds for $i \leq i_0$, then we see that
$R^j\epsilon_{Y, *}H^i(M) = 0$ for $j > 0$ and $i \leq i_0$
by (1) because $H^i(M) = a_Y^{-1}H^i(L)$ in this range.
Hence $\epsilon_{Y, *}H^{i_0 + 1}(M) = H^{i_0 + 1}(R\epsilon_{Y, *}M)$
by a spectral sequence argument.
Thus $\epsilon_{Y, *}H^{i_0 + 1}(M) = \pi_Y^{-1}H^{i_0 + 1}(L) =
\epsilon_{Y, *}a_Y^{-1}H^{i_0 + 1}(L)$.
This implies $H^{i_0 + 1}(\alpha)$ is an isomorphism
(because $\epsilon_{Y, *}$ reflects isomorphisms as it is the
identity on underlying presheaves) as desired.
\end{proof}

\begin{lemma}
\label{lemma-cohomological-descent-etale-fppf}
Let $X$ be a scheme. For $K \in D^+(X_\etale)$ the map
$$
K \longrightarrow Ra_{X, *}a_X^{-1}K
$$
is an isomorphism with $a_X : \Sh((\Sch/X)_{fppf}) \to \Sh(X_\etale)$
as above.
\end{lemma}

\begin{proof}
We first reduce the statement to the case where
$K$ is given by a single abelian sheaf. Namely, represent $K$
by a bounded below complex $\mathcal{F}^\bullet$. By the case of a
sheaf we see that
$\mathcal{F}^n = a_{X, *} a_X^{-1} \mathcal{F}^n$
and that the sheaves $R^qa_{X, *}a_X^{-1}\mathcal{F}^n$
are zero for $q > 0$. By Leray's acyclicity lemma
(Derived Categories, Lemma \ref{derived-lemma-leray-acyclicity})
applied to $a_X^{-1}\mathcal{F}^\bullet$
and the functor $a_{X, *}$ we conclude. From now on assume $K = \mathcal{F}$.

\medskip\noindent
By Lemma \ref{lemma-describe-pullback-pi-fppf} we have
$a_{X, *}a_X^{-1}\mathcal{F} = \mathcal{F}$. Thus it suffices to show that
$R^qa_{X, *}a_X^{-1}\mathcal{F} = 0$ for $q > 0$.
For this we can use $a_X = \epsilon_X \circ \pi_X$ and
the Leray spectral sequence
(Cohomology on Sites, Lemma \ref{sites-cohomology-lemma-relative-Leray}).
By Lemma \ref{lemma-V-C-all-n-etale-fppf}
we have $R^i\epsilon_{X, *}(a_X^{-1}\mathcal{F}) = 0$ for $i > 0$
and
$\epsilon_{X, *}a_X^{-1}\mathcal{F} = \pi_X^{-1}\mathcal{F}$.
By Lemma \ref{lemma-cohomological-descent-etale} we have
$R^j\pi_{X, *}(\pi_X^{-1}\mathcal{F}) = 0$ for $j > 0$.
This concludes the proof.
\end{proof}

\begin{lemma}
\label{lemma-compare-cohomology-etale-fppf}
For a scheme $X$ and $a_X : \Sh((\Sch/X)_{fppf}) \to \Sh(X_\etale)$
as above:
\begin{enumerate}
\item $H^q(X_\etale, \mathcal{F}) = H^q_{fppf}(X, a_X^{-1}\mathcal{F})$
for an abelian sheaf $\mathcal{F}$ on $X_\etale$,
\item $H^q(X_\etale, K) = H^q_{fppf}(X, a_X^{-1}K)$ for $K \in D^+(X_\etale)$.
\end{enumerate}
Example: if $A$ is an abelian group, then
$H^q_\etale(X, \underline{A}) = H^q_{fppf}(X, \underline{A})$.
\end{lemma}

\begin{proof}
This follows from Lemma \ref{lemma-cohomological-descent-etale-fppf}
by Cohomology on Sites, Remark \ref{sites-cohomology-remark-before-Leray}.
\end{proof}







\section{Comparing fppf and \'etale topologies: modules}
\label{section-fppf-etale-modules}

\noindent
We continue the discussion in Section \ref{section-fppf-etale} but in this
section we briefly discuss what happens for sheaves of modules.

\medskip\noindent
Let $S$ be a scheme. The morphisms of sites $\epsilon_S$, $\pi_S$, and
their composition $a_S$ introduced in Section \ref{section-fppf-etale}
have natural enhancements to morphisms of ringed sites. The first
is written as
$$
\epsilon_S :
((\Sch/S)_{fppf}, \mathcal{O})
\longrightarrow
((\Sch/S)_\etale, \mathcal{O})
$$
Note that we can use the same symbol for the structure sheaf as indeed
the sheaves have the same underlying presheaf. The second is
$$
\pi_S :
((\Sch/S)_\etale, \mathcal{O})
\longrightarrow
(S_\etale, \mathcal{O}_S)
$$
The third is the morphism
$$
a_S :
((\Sch/S)_{fppf}, \mathcal{O})
\longrightarrow
(S_\etale, \mathcal{O}_S)
$$
We already know that the category of quasi-coherent modules on the
scheme $S$ is the same as the category of quasi-coherent modules
on $(S_\etale, \mathcal{O}_S)$, see
Descent, Proposition \ref{descent-proposition-equivalence-quasi-coherent}.
Since we are interested in stating a comparison between
\'etale and fppf cohomology, we will in the rest of this
section think of quasi-coherent sheaves in terms of the
small \'etale site.
Let us review what we already know about quasi-coherent
modules on these sites.

\begin{lemma}
\label{lemma-review-quasi-coherent}
Let $S$ be a scheme. Let $\mathcal{F}$ be a quasi-coherent
$\mathcal{O}_S$-module on $S_\etale$.
\begin{enumerate}
\item The rule
$$
\mathcal{F}^a : (\Sch/S)_\etale \longrightarrow \textit{Ab},\quad
(f : T \to S) \longmapsto \Gamma(T, f_{small}^*\mathcal{F})
$$
satisfies the sheaf condition for fppf and a fortiori \'etale coverings,
\item $\mathcal{F}^a = \pi_S^*\mathcal{F}$ on $(\Sch/S)_\etale$,
\item $\mathcal{F}^a = a_S^*\mathcal{F}$ on $(\Sch/S)_{fppf}$,
\item the rule $\mathcal{F} \mapsto \mathcal{F}^a$ defines
an equivalence between quasi-coherent $\mathcal{O}_S$-modules
and quasi-coherent modules on
$((\Sch/S)_\etale, \mathcal{O})$,
\item the rule $\mathcal{F} \mapsto \mathcal{F}^a$ defines
an equivalence between quasi-coherent $\mathcal{O}_S$-modules
and quasi-coherent modules on
$((\Sch/S)_{fppf}, \mathcal{O})$,
\item we have $\epsilon_{S, *}a_S^*\mathcal{F} = \pi_S^*\mathcal{F}$
and $a_{S, *}a_S^*\mathcal{F} = \mathcal{F}$,
\item we have $R^i\epsilon_{S, *}(a_S^*\mathcal{F}) = 0$
and $R^ia_{S, *}(a_S^*\mathcal{F}) = 0$ for $i > 0$.
\end{enumerate}
\end{lemma}

\begin{proof}
We urge the reader to find their own proof of these results
based on the material in
Descent, Sections \ref{descent-section-quasi-coherent-sheaves},
\ref{descent-section-quasi-coherent-cohomology}, and
\ref{descent-section-quasi-coherent-sheaves-bis}.

\medskip\noindent
We first explain why the notation in this lemma is consistent with our
earlier use of the notation $\mathcal{F}^a$ in
Sections \ref{section-quasi-coherent} and
\ref{section-cohomology-quasi-coherent}
and in
Descent, Section \ref{descent-section-quasi-coherent-sheaves}.
Namely, we know by
Descent, Proposition \ref{descent-proposition-equivalence-quasi-coherent}
that there exists a quasi-coherent module
$\mathcal{F}_0$ on the scheme $S$ (in other words on the small
Zariski site) such that $\mathcal{F}$ is the restriction of the
rule
$$
\mathcal{F}_0^a : (\Sch/S)_\etale \longrightarrow \textit{Ab},\quad
(f : T \to S) \longmapsto \Gamma(T, f^*\mathcal{F})
$$
to the subcategory $S_\etale \subset (\Sch/S)_\etale$
where here $f^*$ denotes usual pullback of sheaves of modules on schemes.
Since $\mathcal{F}_0^a$ is pullback by the morphism of ringed
sites
$$
((\Sch/S)_\etale, \mathcal{O}) \longrightarrow (S_{Zar}, \mathcal{O}_{S_{Zar}})
$$
by Descent, Remark \ref{descent-remark-change-topologies-ringed-sites}
it follows immediately (from composition of pullbacks) that
$\mathcal{F}^a = \mathcal{F}_0^a$. This proves the sheaf property
even for fpqc coverings by
Descent, Lemma \ref{descent-lemma-sheaf-condition-holds} (see also
Proposition \ref{proposition-quasi-coherent-sheaf-fpqc}).
Then (2) and (3) follow
again by Descent, Remark \ref{descent-remark-change-topologies-ringed-sites}
and (4) and (5) follow from
Descent, Proposition \ref{descent-proposition-equivalence-quasi-coherent}
(see also the meta result
Theorem \ref{theorem-quasi-coherent}).

\medskip\noindent
Part (6) is immediate from the description of the sheaf
$\mathcal{F}^a = \pi_S^*\mathcal{F} = a_S^*\mathcal{F}$.

\medskip\noindent
For any abelian $\mathcal{H}$ on $(\Sch/S)_{fppf}$ the
higher direct image $R^p\epsilon_{S, *}\mathcal{H}$ is the sheaf
associated to the presheaf $U \mapsto H^p_{fppf}(U, \mathcal{H})$
on $(\Sch/S)_\etale$. See
Cohomology on Sites, Lemma \ref{sites-cohomology-lemma-higher-direct-images}.
Hence to prove
$R^p\epsilon_{S, *}a_S^*\mathcal{F} = R^p\epsilon_{S, *}\mathcal{F}^a = 0$
for $p > 0$ it suffices to show that any scheme $U$ over $S$
has an \'etale covering $\{U_i \to U\}_{i \in I}$ such that
$H^p_{fppf}(U_i, \mathcal{F}^a) = 0$ for $p > 0$.
If we take an open covering by affines, then the required
vanishing follows from comparison with usual cohomology
(Descent, Proposition \ref{descent-proposition-same-cohomology-quasi-coherent}
or
Theorem \ref{theorem-zariski-fpqc-quasi-coherent})
and the vanishing of cohomology of quasi-coherent sheaves
on affine schemes afforded by Cohomology of Schemes, Lemma
\ref{coherent-lemma-quasi-coherent-affine-cohomology-zero}.

\medskip\noindent
To show that $R^pa_{S, *}a_S^{-1}\mathcal{F} = R^pa_{S, *}\mathcal{F}^a = 0$
for $p > 0$ we argue in exactly the same manner. This finishes the proof.
\end{proof}

\begin{lemma}
\label{lemma-cohomological-descent-etale-fppf-modules}
Let $S$ be a scheme. For $\mathcal{F}$ a quasi-coherent
$\mathcal{O}_S$-module on $S_\etale$ the maps
$$
\pi_S^*\mathcal{F} \longrightarrow R\epsilon_{S, *}(a_S^*\mathcal{F})
\quad\text{and}\quad
\mathcal{F} \longrightarrow Ra_{S, *}(a_S^*\mathcal{F})
$$
are isomorphisms with
$a_S : \Sh((\Sch/S)_{fppf}) \to \Sh(S_\etale)$ as above.
\end{lemma}

\begin{proof}
This is an immediate consequence of
parts (6) and (7) of
Lemma \ref{lemma-review-quasi-coherent}.
\end{proof}

\begin{lemma}
\label{lemma-cohomological-descent-complex-modules}
Let $S = \Spec(A)$ be an affine scheme. Let $M^\bullet$ be a complex
of $A$-modules. Consider the complex $\mathcal{F}^\bullet$ of
presheaves of $\mathcal{O}$-modules on
$(\textit{Aff}/S)_{fppf}$ given by the rule
$$
(U/S) = (\Spec(B)/\Spec(A)) \longmapsto M^\bullet \otimes_A B
$$
Then this is a complex of modules and the canonical map
$$
M^\bullet \longrightarrow
R\Gamma((\textit{Aff}/S)_{fppf}, \mathcal{F}^\bullet)
$$
is a quasi-isomorphism.
\end{lemma}

\begin{proof}
Each $\mathcal{F}^n$ is a sheaf of modules as it agrees with the
restriction of the module $\mathcal{G}^n = (\widetilde{M}^n)^a$ of
Lemma \ref{lemma-review-quasi-coherent} to
$(\textit{Aff}/S)_{fppf} \subset (\Sch/S)_{fppf}$.
Since this inclusion defines an equivalence of ringed topoi
(Topologies, Lemma \ref{topologies-lemma-affine-big-site-fppf}),
we have
$$
R\Gamma((\textit{Aff}/S)_{fppf}, \mathcal{F}^\bullet) =
R\Gamma((\Sch/S)_{fppf}, \mathcal{G}^\bullet)
$$
We observe that $M^\bullet = R\Gamma(S, \widetilde{M}^\bullet)$
for example by Derived Categories of Schemes, Lemma
\ref{perfect-lemma-affine-compare-bounded}. Hence we are trying
to show the comparison map
$$
R\Gamma(S, \widetilde{M}^\bullet)
\longrightarrow
R\Gamma((\Sch/S)_{fppf}, (\widetilde{M}^\bullet)^a)
$$
is an isomorphism. If $M^\bullet$ is bounded below, then this
holds by Descent, Proposition
\ref{descent-proposition-same-cohomology-quasi-coherent}
and the first spectral sequence of Derived Categories, Lemma
\ref{derived-lemma-two-ss-complex-functor}.
For the general case, let us write $M^\bullet = \lim M_n^\bullet$
with $M_n^\bullet = \tau_{\geq -n}M^\bullet$. Whence the system
$M_n^p$ is eventually constant with value $M^p$. We claim that
$$
(\widetilde{M}^\bullet)^a = R\lim (\widetilde{M}_n^\bullet)^a
$$
Namely, it suffices to show that the natural map from left to right
induces an isomorphism on cohomology over any affine object
$U = \Spec(B)$ of $(\Sch/S)_{fppf}$. For $i \in \mathbf{Z}$ and
$n > |i|$ we have
$$
H^i(U, (\widetilde{M}_n^\bullet)^a) =
H^i(\tau_{\geq -n}M^\bullet \otimes_A B) =
H^i(M^\bullet \otimes_A B)
$$
The first equality holds by the bounded below case treated above.
Thus we see from Cohomology on Sites, Lemma
\ref{sites-cohomology-lemma-RGamma-commutes-with-Rlim}
that the claim holds. Then we finally get
\begin{align*}
R\Gamma((\Sch/S)_{fppf}, (\widetilde{M}^\bullet)^a)
& =
R\Gamma((\Sch/S)_{fppf}, R\lim (\widetilde{M}_n^\bullet)^a) \\
& =
R\lim R\Gamma((\Sch/S)_{fppf}, (\widetilde{M}_n^\bullet)^a) \\
& =
R\lim M_n^\bullet \\
& =
M^\bullet
\end{align*}
as desired. The second equality holds because $R\lim$ commutes with
$R\Gamma$, see Cohomology on Sites, Lemma
\ref{sites-cohomology-lemma-RGamma-commutes-with-Rlim}.
\end{proof}










\section{Comparing ph and \'etale topologies}
\label{section-ph-etale}

\noindent
A model for this section is the section on the comparison of the
usual topology and the qc topology on locally compact topological
spaces as discussed in Cohomology on Sites, Section
\ref{sites-cohomology-section-cohomology-LC}.
We first review some material from
Topologies, Sections
\ref{topologies-section-change-topologies} and
\ref{topologies-section-etale}.

\medskip\noindent
Let $S$ be a scheme and let $(\Sch/S)_{ph}$ be a ph site.
On the same underlying category we have a second topology,
namely the \'etale topology, and hence a second site
$(\Sch/S)_\etale$. The identity functor
$(\Sch/S)_\etale \to (\Sch/S)_{ph}$ is continuous
(by More on Morphisms, Lemma \ref{more-morphisms-lemma-fppf-ph}
and Topologies, Lemma
\ref{topologies-lemma-zariski-etale-smooth-syntomic-fppf})
and defines a morphism of sites
$$
\epsilon_S : (\Sch/S)_{ph} \longrightarrow (\Sch/S)_\etale
$$
See Cohomology on Sites, Section \ref{sites-cohomology-section-compare}.
Please note that $\epsilon_{S, *}$ is the identity functor on underlying
presheaves and that $\epsilon_S^{-1}$ associates to an \'etale sheaf the
ph sheafification.
Let $S_\etale$ be the small \'etale site.
There is a morphism of sites
$$
\pi_S : (\Sch/S)_\etale \longrightarrow S_\etale
$$
given by the continuous functor
$S_\etale \to (\Sch/S)_\etale$, $U \mapsto U$.
Namely, $S_\etale$ has fibre products and a final object and the
functor above commutes with these and
Sites, Proposition \ref{sites-proposition-get-morphism} applies.

\begin{lemma}
\label{lemma-describe-pullback-pi-ph}
With notation as above.
Let $\mathcal{F}$ be a sheaf on $S_\etale$. The rule
$$
(\Sch/S)_{ph} \longrightarrow \textit{Sets},\quad
(f : X \to S) \longmapsto \Gamma(X, f_{small}^{-1}\mathcal{F})
$$
is a sheaf and a fortiori a sheaf on $(\Sch/S)_\etale$.
In fact this sheaf is equal to
$\pi_S^{-1}\mathcal{F}$ on $(\Sch/S)_\etale$ and
$\epsilon_S^{-1}\pi_S^{-1}\mathcal{F}$ on $(\Sch/S)_{ph}$.
\end{lemma}

\begin{proof}
The statement about the \'etale topology is the content
of Lemma \ref{lemma-describe-pullback}. To finish the proof it
suffices to show that $\pi_S^{-1}\mathcal{F}$ is a sheaf for the ph
topology. By Topologies, Lemma \ref{topologies-lemma-characterize-sheaf}
it suffices to show that given a proper surjective morphism
$V \to U$ of schemes over $S$ we have an equalizer diagram
$$
\xymatrix{
(\pi_S^{-1}\mathcal{F})(U) \ar[r] &
(\pi_S^{-1}\mathcal{F})(V) \ar@<1ex>[r] \ar@<-1ex>[r] &
(\pi_S^{-1}\mathcal{F})(V \times_U V)
}
$$
Set $\mathcal{G} = \pi_S^{-1}\mathcal{F}|_{U_\etale}$.
Consider the commutative diagram
$$
\xymatrix{
V \times_U V \ar[r] \ar[rd]_g \ar[d] & V \ar[d]^f \\
V \ar[r]^f & U
}
$$
We have
$$
(\pi_S^{-1}\mathcal{F})(V) = \Gamma(V, f^{-1}\mathcal{G}) =
\Gamma(U, f_*f^{-1}\mathcal{G})
$$
where we use $f_*$ and $f^{-1}$ to denote functorialities between
small \'etale sites. Second, we have
$$
(\pi_S^{-1}\mathcal{F})(V \times_U V) =
\Gamma(V \times_U V, g^{-1}\mathcal{G}) =
\Gamma(U, g_*g^{-1}\mathcal{G})
$$
The two maps in the equalizer diagram come from the two maps
$$
f_*f^{-1}\mathcal{G} \longrightarrow g_*g^{-1}\mathcal{G}
$$
Thus it suffices to prove $\mathcal{G}$ is
the equalizer of these two maps of sheaves.
Let $\overline{u}$ be a geometric point of $U$. Set
$\Omega = \mathcal{G}_{\overline{u}}$.
Taking stalks at $\overline{u}$ by
Lemma \ref{lemma-proper-pushforward-stalk}
we obtain the two maps
$$
H^0(V_{\overline{u}}, \underline{\Omega}) \longrightarrow
H^0((V \times_U V)_{\overline{u}}, \underline{\Omega}) =
H^0(V_{\overline{u}} \times_{\overline{u}} V_{\overline{u}},
\underline{\Omega})
$$
where $\underline{\Omega}$ indicates the constant sheaf with value
$\Omega$. Of course these maps are the pullback by the projection maps.
Then it is clear that the sections coming from pullback
by projection onto the first factor are constant on the fibres of
the first projection, and sections coming from pullback
by projection onto the first factor are constant on the fibres of
the first projection. The sections in the intersection of the images
of these pullback maps are constant on all of
$V_{\overline{u}} \times_{\overline{u}} V_{\overline{u}}$, i.e.,
these come from elements of $\Omega$ as desired.
\end{proof}

\noindent
In the situation of Lemma \ref{lemma-describe-pullback-pi-ph}
the composition of $\epsilon_S$ and $\pi_S$ and the equality
determine a morphism of sites
$$
a_S : (\Sch/S)_{ph} \longrightarrow S_\etale
$$

\begin{lemma}
\label{lemma-push-pull-ph-etale}
With notation as above.
Let $f : X \to Y$ be a morphism of $(\Sch/S)_{ph}$.
Then there are commutative diagrams of topoi
$$
\xymatrix{
\Sh((\Sch/X)_{ph}) \ar[rr]_{f_{big, ph}} \ar[d]_{\epsilon_X} & &
\Sh((\Sch/Y)_{ph}) \ar[d]^{\epsilon_Y} \\
\Sh((\Sch/X)_\etale) \ar[rr]^{f_{big, \etale}} & &
\Sh((\Sch/Y)_\etale)
}
$$
and
$$
\xymatrix{
\Sh((\Sch/X)_{ph}) \ar[rr]_{f_{big, ph}} \ar[d]_{a_X} & &
\Sh((\Sch/Y)_{ph}) \ar[d]^{a_Y} \\
\Sh(X_\etale) \ar[rr]^{f_{small}} & &
\Sh(Y_\etale)
}
$$
with $a_X = \pi_X \circ \epsilon_X$ and $a_Y = \pi_X \circ \epsilon_X$.
\end{lemma}

\begin{proof}
The commutativity of the diagrams follows from the discussion in
Topologies, Section \ref{topologies-section-change-topologies}.
\end{proof}

\begin{lemma}
\label{lemma-proper-push-pull-ph-etale}
In Lemma \ref{lemma-push-pull-ph-etale} if $f$ is proper, then we have
$a_Y^{-1} \circ f_{small, *} = f_{big, ph, *} \circ a_X^{-1}$.
\end{lemma}

\begin{proof}
You can prove this by repeating the proof of
Lemma \ref{lemma-compare-higher-direct-image-proper} part (1);
we will instead deduce the result from this.
As $\epsilon_{Y, *}$ is the identity functor on underlying presheaves,
it reflects isomorphisms. The description
in Lemma \ref{lemma-describe-pullback-pi-ph}
shows that $\epsilon_{Y, *} \circ a_Y^{-1} = \pi_Y^{-1}$
and similarly for $X$. To show that the canonical map
$a_Y^{-1}f_{small, *}\mathcal{F} \to f_{big, ph, *}a_X^{-1}\mathcal{F}$
is an isomorphism, it suffices to show that
\begin{align*}
\pi_Y^{-1}f_{small, *}\mathcal{F}
& =
\epsilon_{Y, *}a_Y^{-1}f_{small, *}\mathcal{F} \\
& \to 
\epsilon_{Y, *}f_{big, ph, *}a_X^{-1}\mathcal{F} \\
& =
f_{big, \etale, *} \epsilon_{X, *}a_X^{-1}\mathcal{F} \\
& =
f_{big, \etale, *}\pi_X^{-1}\mathcal{F}
\end{align*}
is an isomorphism. This is part
(1) of Lemma \ref{lemma-compare-higher-direct-image-proper}.
\end{proof}

\begin{lemma}
\label{lemma-compare-ph-etale}
Consider the comparison morphism
$\epsilon : (\Sch/S)_{ph} \to (\Sch/S)_\etale$.
Let $\mathcal{P}$ denote the class of proper morphisms of schemes.
For $X$ in $(\Sch/S)_\etale$ denote
$\mathcal{A}'_X \subset \textit{Ab}((\Sch/X)_\etale)$
the full subcategory consisting of sheaves of the form
$\pi_X^{-1}\mathcal{F}$ where $\mathcal{F}$ is a
torsion abelian sheaf on $X_\etale$
Then Cohomology on Sites, Properties
(\ref{sites-cohomology-item-base-change-P}),
(\ref{sites-cohomology-item-restriction-A}),
(\ref{sites-cohomology-item-A-sheaf}),
(\ref{sites-cohomology-item-A-and-P}), and
(\ref{sites-cohomology-item-refine-tau-by-P})
of Cohomology on Sites, Situation
\ref{sites-cohomology-situation-compare} hold.
\end{lemma}

\begin{proof}
We first show that $\mathcal{A}'_X \subset \textit{Ab}((\Sch/X)_\etale)$
is a weak Serre subcategory by checking conditions (1), (2), (3), and (4)
of Homology, Lemma \ref{homology-lemma-characterize-weak-serre-subcategory}.
Parts (1), (2), (3) are immediate as $\pi_X^{-1}$ is exact and
fully faithful for example by Lemma \ref{lemma-cohomological-descent-etale}. If
$0 \to \pi_X^{-1}\mathcal{F} \to \mathcal{G} \to \pi_X^{-1}\mathcal{F}' \to 0$
is a short exact sequence in $\textit{Ab}((\Sch/X)_\etale)$
then $0 \to \mathcal{F} \to \pi_{X, *}\mathcal{G} \to \mathcal{F}' \to 0$
is exact by Lemma \ref{lemma-cohomological-descent-etale}.
In particular we see that $\pi_{X, *}\mathcal{G}$ is an abelian
torsion sheaf on $X_\etale$.
Hence $\mathcal{G} = \pi_X^{-1}\pi_{X, *}\mathcal{G}$ is in
$\mathcal{A}'_X$ which checks the final condition.

\medskip\noindent
Cohomology on Sites, Property (\ref{sites-cohomology-item-base-change-P}) holds
by the existence of fibre products of schemes
and the fact that the base change of a proper morphism of
schemes is a proper morphism of schemes, see
Morphisms, Lemma \ref{morphisms-lemma-base-change-proper}.

\medskip\noindent
Cohomology on Sites, Property (\ref{sites-cohomology-item-restriction-A})
follows from the commutative diagram (3) in
Topologies, Lemma \ref{topologies-lemma-morphism-big-small-etale}.

\medskip\noindent
Cohomology on Sites, Property (\ref{sites-cohomology-item-A-sheaf}) is
Lemma \ref{lemma-describe-pullback-pi-ph}.

\medskip\noindent
Cohomology on Sites, Property (\ref{sites-cohomology-item-A-and-P}) holds by
Lemma \ref{lemma-compare-higher-direct-image-proper} part (2)
and the fact that $R^if_{small}\mathcal{F}$
is torsion if $\mathcal{F}$ is an abelian torsion
sheaf on $X_\etale$, see Lemma \ref{lemma-torsion-direct-image}.

\medskip\noindent
Cohomology on Sites, Property (\ref{sites-cohomology-item-refine-tau-by-P})
follows from More on Morphisms, Lemma
\ref{more-morphisms-lemma-dominate-fppf-etale-locally}
combined with the fact that a finite morphism is proper
and a surjective proper morphism defines a ph covering, see
Topologies, Lemma \ref{topologies-lemma-surjective-proper-ph}.
\end{proof}

\begin{lemma}
\label{lemma-V-C-all-n-etale-ph}
With notation as above.
\begin{enumerate}
\item For $X \in \Ob((\Sch/S)_{ph})$ and an abelian torsion sheaf $\mathcal{F}$
on $X_\etale$ we have
$\epsilon_{X, *}a_X^{-1}\mathcal{F} = \pi_X^{-1}\mathcal{F}$
and $R^i\epsilon_{X, *}(a_X^{-1}\mathcal{F}) = 0$ for $i > 0$.
\item For a proper morphism $f : X \to Y$ in $(\Sch/S)_{ph}$
and abelian torsion sheaf $\mathcal{F}$ on $X$ we have
$a_Y^{-1}(R^if_{small, *}\mathcal{F}) =
R^if_{big, ph, *}(a_X^{-1}\mathcal{F})$
for all $i$.
\item For a scheme $X$ and $K$ in $D^+(X_\etale)$ with torsion
cohomology sheaves the map
$\pi_X^{-1}K \to R\epsilon_{X, *}(a_X^{-1}K)$ is an isomorphism.
\item For a proper morphism $f : X \to Y$ of schemes
and $K$ in $D^+(X_\etale)$ with torsion cohomology sheaves we have
$a_Y^{-1}(Rf_{small, *}K) = Rf_{big, ph, *}(a_X^{-1}K)$.
\end{enumerate}
\end{lemma}

\begin{proof}
By Lemma \ref{lemma-compare-ph-etale} the lemmas in
Cohomology on Sites, Section \ref{sites-cohomology-section-compare-general}
all apply to our current setting. To translate the results
observe that the category $\mathcal{A}_X$ of
Cohomology on Sites, Lemma \ref{sites-cohomology-lemma-A}
is the full subcategory of $\textit{Ab}((\Sch/X)_{ph})$
consisting of sheaves of the form $a_X^{-1}\mathcal{F}$
where $\mathcal{F}$ is an abelian torsion sheaf on $X_\etale$.

\medskip\noindent
Part (1) is equivalent to $(V_n)$ for all $n$ which holds by
Cohomology on Sites, Lemma \ref{sites-cohomology-lemma-V-C-all-n-general}.

\medskip\noindent
Part (2) follows by applying $\epsilon_Y^{-1}$ to the conclusion of
Cohomology on Sites, Lemma \ref{sites-cohomology-lemma-V-implies-C-general}.

\medskip\noindent
Part (3) follows from Cohomology on Sites, Lemma
\ref{sites-cohomology-lemma-V-C-all-n-general} part (1)
because $\pi_X^{-1}K$ is in $D^+_{\mathcal{A}'_X}((\Sch/X)_\etale)$
and $a_X^{-1} = \epsilon_X^{-1} \circ a_X^{-1}$.

\medskip\noindent
Part (4) follows from Cohomology on Sites, Lemma
\ref{sites-cohomology-lemma-V-C-all-n-general} part (2)
for the same reason.
\end{proof}

\begin{lemma}
\label{lemma-cohomological-descent-etale-ph}
Let $X$ be a scheme. For $K \in D^+(X_\etale)$ with torsion cohomology
sheaves the map
$$
K \longrightarrow Ra_{X, *}a_X^{-1}K
$$
is an isomorphism with $a_X : \Sh((\Sch/X)_{ph}) \to \Sh(X_\etale)$ as above.
\end{lemma}

\begin{proof}
We first reduce the statement to the case where
$K$ is given by a single abelian sheaf. Namely, represent $K$
by a bounded below complex $\mathcal{F}^\bullet$ of torsion
abelian sheaves. This is possible by Cohomology on Sites, Lemma
\ref{sites-cohomology-lemma-torsion}. By the case of a
sheaf we see that
$\mathcal{F}^n = a_{X, *} a_X^{-1} \mathcal{F}^n$
and that the sheaves $R^qa_{X, *}a_X^{-1}\mathcal{F}^n$
are zero for $q > 0$. By Leray's acyclicity lemma
(Derived Categories, Lemma \ref{derived-lemma-leray-acyclicity})
applied to $a_X^{-1}\mathcal{F}^\bullet$
and the functor $a_{X, *}$ we conclude. From now on assume $K = \mathcal{F}$
where $\mathcal{F}$ is a torsion abelian sheaf.

\medskip\noindent
By Lemma \ref{lemma-describe-pullback-pi-ph} we have
$a_{X, *}a_X^{-1}\mathcal{F} = \mathcal{F}$. Thus it suffices to show that
$R^qa_{X, *}a_X^{-1}\mathcal{F} = 0$ for $q > 0$.
For this we can use $a_X = \epsilon_X \circ \pi_X$ and
the Leray spectral sequence
(Cohomology on Sites, Lemma \ref{sites-cohomology-lemma-relative-Leray}).
By Lemma \ref{lemma-V-C-all-n-etale-ph}
we have $R^i\epsilon_{X, *}(a_X^{-1}\mathcal{F}) = 0$ for $i > 0$
and $\epsilon_{X, *}a_X^{-1}\mathcal{F} = \pi_X^{-1}\mathcal{F}$.
By Lemma \ref{lemma-cohomological-descent-etale} we have
$R^j\pi_{X, *}(\pi_X^{-1}\mathcal{F}) = 0$ for $j > 0$.
This concludes the proof.
\end{proof}

\begin{lemma}
\label{lemma-compare-cohomology-etale-ph}
For a scheme $X$ and $a_X : \Sh((\Sch/X)_{ph}) \to \Sh(X_\etale)$
as above:
\begin{enumerate}
\item $H^q(X_\etale, \mathcal{F}) = H^q_{ph}(X, a_X^{-1}\mathcal{F})$
for a torsion abelian sheaf $\mathcal{F}$ on $X_\etale$,
\item $H^q(X_\etale, K) = H^q_{ph}(X, a_X^{-1}K)$
for $K \in D^+(X_\etale)$ with torsion cohomology sheaves.
\end{enumerate}
Example: if $A$ is a torsion abelian group, then
$H^q_\etale(X, \underline{A}) = H^q_{ph}(X, \underline{A})$.
\end{lemma}

\begin{proof}
This follows from Lemma \ref{lemma-cohomological-descent-etale-ph}
by Cohomology on Sites, Remark \ref{sites-cohomology-remark-before-Leray}.
\end{proof}












\section{Comparing h and \'etale topologies}
\label{section-h-etale}

\noindent
A model for this section is the section on the comparison of the
usual topology and the qc topology on locally compact topological
spaces as discussed in
Cohomology on Sites, Section \ref{sites-cohomology-section-cohomology-LC}.
Moreover, this section is almost word for word the same as the
section comparing the ph and \'etale topologies.
We first review some material from
Topologies, Sections
\ref{topologies-section-change-topologies} and
\ref{topologies-section-etale} and
More on Flatness, Section \ref{flat-section-h}.

\medskip\noindent
Let $S$ be a scheme and let $(\Sch/S)_h$ be an h site.
On the same underlying category we have a second topology,
namely the \'etale topology, and hence a second site
$(\Sch/S)_\etale$. The identity functor
$(\Sch/S)_\etale \to (\Sch/S)_h$ is continuous
(by More on Flatness, Lemma \ref{flat-lemma-zariski-h}
and Topologies, Lemma
\ref{topologies-lemma-zariski-etale-smooth-syntomic-fppf})
and defines a morphism of sites
$$
\epsilon_S : (\Sch/S)_h \longrightarrow (\Sch/S)_\etale
$$
See Cohomology on Sites, Section \ref{sites-cohomology-section-compare}.
Please note that $\epsilon_{S, *}$ is the identity functor on underlying
presheaves and that $\epsilon_S^{-1}$ associates to an \'etale sheaf the
h sheafification. Let $S_\etale$ be the small \'etale site.
There is a morphism of sites
$$
\pi_S : (\Sch/S)_\etale \longrightarrow S_\etale
$$
given by the continuous functor
$S_\etale \to (\Sch/S)_\etale$, $U \mapsto U$.
Namely, $S_\etale$ has fibre products and a final object and the
functor above commutes with these and
Sites, Proposition \ref{sites-proposition-get-morphism} applies.

\begin{lemma}
\label{lemma-describe-pullback-pi-h}
With notation as above.
Let $\mathcal{F}$ be a sheaf on $S_\etale$. The rule
$$
(\Sch/S)_h \longrightarrow \textit{Sets},\quad
(f : X \to S) \longmapsto \Gamma(X, f_{small}^{-1}\mathcal{F})
$$
is a sheaf and a fortiori a sheaf on $(\Sch/S)_\etale$.
In fact this sheaf is equal to
$\pi_S^{-1}\mathcal{F}$ on $(\Sch/S)_\etale$ and
$\epsilon_S^{-1}\pi_S^{-1}\mathcal{F}$ on $(\Sch/S)_h$.
\end{lemma}

\begin{proof}
The statement about the \'etale topology is the content
of Lemma \ref{lemma-describe-pullback}. To finish the proof it
suffices to show that $\pi_S^{-1}\mathcal{F}$ is a sheaf for the h
topology. However, in Lemma \ref{lemma-describe-pullback-pi-ph}
we have shown that
$\pi_S^{-1}\mathcal{F}$ is a sheaf even in the stronger ph
topology.
\end{proof}

\noindent
In the situation of Lemma \ref{lemma-describe-pullback-pi-h}
the composition of $\epsilon_S$ and $\pi_S$ and the equality
determine a morphism of sites
$$
a_S : (\Sch/S)_h \longrightarrow S_\etale
$$

\begin{lemma}
\label{lemma-push-pull-h-etale}
With notation as above.
Let $f : X \to Y$ be a morphism of $(\Sch/S)_h$.
Then there are commutative diagrams of topoi
$$
\xymatrix{
\Sh((\Sch/X)_h) \ar[rr]_{f_{big, h}} \ar[d]_{\epsilon_X} & &
\Sh((\Sch/Y)_h) \ar[d]^{\epsilon_Y} \\
\Sh((\Sch/X)_\etale) \ar[rr]^{f_{big, \etale}} & &
\Sh((\Sch/Y)_\etale)
}
$$
and
$$
\xymatrix{
\Sh((\Sch/X)_h) \ar[rr]_{f_{big, h}} \ar[d]_{a_X} & &
\Sh((\Sch/Y)_h) \ar[d]^{a_Y} \\
\Sh(X_\etale) \ar[rr]^{f_{small}} & &
\Sh(Y_\etale)
}
$$
with $a_X = \pi_X \circ \epsilon_X$ and $a_Y = \pi_X \circ \epsilon_X$.
\end{lemma}

\begin{proof}
The commutativity of the diagrams follows similarly to what was said in
Topologies, Section \ref{topologies-section-change-topologies}.
\end{proof}

\begin{lemma}
\label{lemma-proper-push-pull-h-etale}
In Lemma \ref{lemma-push-pull-h-etale} if $f$ is proper, then we have
$a_Y^{-1} \circ f_{small, *} = f_{big, h, *} \circ a_X^{-1}$.
\end{lemma}

\begin{proof}
You can prove this by repeating the proof of
Lemma \ref{lemma-compare-higher-direct-image-proper} part (1);
we will instead deduce the result from this.
As $\epsilon_{Y, *}$ is the identity functor on underlying presheaves,
it reflects isomorphisms. The description
in Lemma \ref{lemma-describe-pullback-pi-h}
shows that $\epsilon_{Y, *} \circ a_Y^{-1} = \pi_Y^{-1}$
and similarly for $X$. To show that the canonical map
$a_Y^{-1}f_{small, *}\mathcal{F} \to f_{big, h, *}a_X^{-1}\mathcal{F}$
is an isomorphism, it suffices to show that
\begin{align*}
\pi_Y^{-1}f_{small, *}\mathcal{F}
& =
\epsilon_{Y, *}a_Y^{-1}f_{small, *}\mathcal{F} \\
& \to 
\epsilon_{Y, *}f_{big, h, *}a_X^{-1}\mathcal{F} \\
& =
f_{big, \etale, *} \epsilon_{X, *}a_X^{-1}\mathcal{F} \\
& =
f_{big, \etale, *}\pi_X^{-1}\mathcal{F}
\end{align*}
is an isomorphism. This is part
(1) of Lemma \ref{lemma-compare-higher-direct-image-proper}.
\end{proof}

\begin{lemma}
\label{lemma-compare-h-etale}
Consider the comparison morphism $\epsilon : (\Sch/S)_h \to (\Sch/S)_\etale$.
Let $\mathcal{P}$ denote the class of proper morphisms.
For $X$ in $(\Sch/S)_\etale$ denote
$\mathcal{A}'_X \subset \textit{Ab}((\Sch/X)_\etale)$
the full subcategory consisting of sheaves of the form
$\pi_X^{-1}\mathcal{F}$ where $\mathcal{F}$ is a
torsion abelian sheaf on $X_\etale$
Then Cohomology on Sites, Properties
(\ref{sites-cohomology-item-base-change-P}),
(\ref{sites-cohomology-item-restriction-A}),
(\ref{sites-cohomology-item-A-sheaf}),
(\ref{sites-cohomology-item-A-and-P}), and
(\ref{sites-cohomology-item-refine-tau-by-P})
of Cohomology on Sites, Situation
\ref{sites-cohomology-situation-compare} hold.
\end{lemma}

\begin{proof}
We first show that $\mathcal{A}'_X \subset \textit{Ab}((\Sch/X)_\etale)$
is a weak Serre subcategory by checking conditions (1), (2), (3), and (4)
of Homology, Lemma \ref{homology-lemma-characterize-weak-serre-subcategory}.
Parts (1), (2), (3) are immediate as $\pi_X^{-1}$ is exact and
fully faithful for example by Lemma \ref{lemma-cohomological-descent-etale}. If
$0 \to \pi_X^{-1}\mathcal{F} \to \mathcal{G} \to \pi_X^{-1}\mathcal{F}' \to 0$
is a short exact sequence in $\textit{Ab}((\Sch/X)_\etale)$
then $0 \to \mathcal{F} \to \pi_{X, *}\mathcal{G} \to \mathcal{F}' \to 0$
is exact by Lemma \ref{lemma-cohomological-descent-etale}.
In particular we see that $\pi_{X, *}\mathcal{G}$ is an abelian
torsion sheaf on $X_\etale$.
Hence $\mathcal{G} = \pi_X^{-1}\pi_{X, *}\mathcal{G}$ is in
$\mathcal{A}'_X$ which checks the final condition.

\medskip\noindent
Cohomology on Sites, Property (\ref{sites-cohomology-item-base-change-P}) holds
by the existence of fibre products of schemes,
the fact that the base change of a proper morphism of
schemes is a proper morphism of schemes, see
Morphisms, Lemma \ref{morphisms-lemma-base-change-proper}, and
the fact that the base change of a morphism of finite presentation
is a morphism of finite presentation, see
Morphisms, Lemma \ref{morphisms-lemma-base-change-finite-presentation}.

\medskip\noindent
Cohomology on Sites, Property (\ref{sites-cohomology-item-restriction-A})
follows from the commutative diagram (3) in
Topologies, Lemma \ref{topologies-lemma-morphism-big-small-etale}.

\medskip\noindent
Cohomology on Sites, Property (\ref{sites-cohomology-item-A-sheaf}) is
Lemma \ref{lemma-describe-pullback-pi-h}.

\medskip\noindent
Cohomology on Sites, Property (\ref{sites-cohomology-item-A-and-P}) holds by
Lemma \ref{lemma-compare-higher-direct-image-proper} part (2)
and the fact that $R^if_{small}\mathcal{F}$
is torsion if $\mathcal{F}$ is an abelian torsion
sheaf on $X_\etale$, see Lemma \ref{lemma-torsion-direct-image}.

\medskip\noindent
Cohomology on Sites, Property (\ref{sites-cohomology-item-refine-tau-by-P})
is implied by More on Morphisms, Lemma
\ref{more-morphisms-lemma-dominate-fppf-etale-locally}
combined with the fact that a surjective finite locally free morphism
is surjective, proper, and of finite presentation and hence
defines a h-covering by More on Flatness, Lemma
\ref{flat-lemma-surjective-proper-finite-presentation-h}.
\end{proof}

\begin{lemma}
\label{lemma-V-C-all-n-etale-h}
With notation as above.
\begin{enumerate}
\item For $X \in \Ob((\Sch/S)_{h})$ and an abelian torsion sheaf $\mathcal{F}$
on $X_\etale$ we have
$\epsilon_{X, *}a_X^{-1}\mathcal{F} = \pi_X^{-1}\mathcal{F}$
and $R^i\epsilon_{X, *}(a_X^{-1}\mathcal{F}) = 0$ for $i > 0$.
\item For a proper morphism $f : X \to Y$ in $(\Sch/S)_h$
and abelian torsion sheaf $\mathcal{F}$ on $X$ we have
$a_Y^{-1}(R^if_{small, *}\mathcal{F}) =
R^if_{big, h, *}(a_X^{-1}\mathcal{F})$
for all $i$.
\item For a scheme $X$ and $K$ in $D^+(X_\etale)$ with torsion
cohomology sheaves the map
$\pi_X^{-1}K \to R\epsilon_{X, *}(a_X^{-1}K)$ is an isomorphism.
\item For a proper morphism $f : X \to Y$ of schemes
and $K$ in $D^+(X_\etale)$ with torsion cohomology sheaves we have
$a_Y^{-1}(Rf_{small, *}K) = Rf_{big, h, *}(a_X^{-1}K)$.
\end{enumerate}
\end{lemma}

\begin{proof}
By Lemma \ref{lemma-compare-h-etale} the lemmas in
Cohomology on Sites, Section \ref{sites-cohomology-section-compare-general}
all apply to our current setting. To translate the results
observe that the category $\mathcal{A}_X$ of
Cohomology on Sites, Lemma \ref{sites-cohomology-lemma-A}
is the full subcategory of $\textit{Ab}((\Sch/X)_h)$
consisting of sheaves of the form $a_X^{-1}\mathcal{F}$
where $\mathcal{F}$ is an abelian torsion sheaf on $X_\etale$.

\medskip\noindent
Part (1) is equivalent to $(V_n)$ for all $n$ which holds by
Cohomology on Sites, Lemma \ref{sites-cohomology-lemma-V-C-all-n-general}.

\medskip\noindent
Part (2) follows by applying $\epsilon_Y^{-1}$ to the conclusion of
Cohomology on Sites, Lemma \ref{sites-cohomology-lemma-V-implies-C-general}.

\medskip\noindent
Part (3) follows from Cohomology on Sites, Lemma
\ref{sites-cohomology-lemma-V-C-all-n-general} part (1)
because $\pi_X^{-1}K$ is in $D^+_{\mathcal{A}'_X}((\Sch/X)_\etale)$
and $a_X^{-1} = \epsilon_X^{-1} \circ a_X^{-1}$.

\medskip\noindent
Part (4) follows from Cohomology on Sites, Lemma
\ref{sites-cohomology-lemma-V-C-all-n-general} part (2)
for the same reason.
\end{proof}

\begin{lemma}
\label{lemma-cohomological-descent-etale-h}
Let $X$ be a scheme. For $K \in D^+(X_\etale)$ with torsion cohomology
sheaves the map
$$
K \longrightarrow Ra_{X, *}a_X^{-1}K
$$
is an isomorphism with $a_X : \Sh((\Sch/X)_h) \to \Sh(X_\etale)$ as above.
\end{lemma}

\begin{proof}
We first reduce the statement to the case where
$K$ is given by a single abelian sheaf. Namely, represent $K$
by a bounded below complex $\mathcal{F}^\bullet$ of torsion
abelian sheaves. This is possible by Cohomology on Sites, Lemma
\ref{sites-cohomology-lemma-torsion}. By the case of a
sheaf we see that
$\mathcal{F}^n = a_{X, *} a_X^{-1} \mathcal{F}^n$
and that the sheaves $R^qa_{X, *}a_X^{-1}\mathcal{F}^n$
are zero for $q > 0$. By Leray's acyclicity lemma
(Derived Categories, Lemma \ref{derived-lemma-leray-acyclicity})
applied to $a_X^{-1}\mathcal{F}^\bullet$
and the functor $a_{X, *}$ we conclude. From now on assume $K = \mathcal{F}$
where $\mathcal{F}$ is a torsion abelian sheaf.

\medskip\noindent
By Lemma \ref{lemma-describe-pullback-pi-h} we have
$a_{X, *}a_X^{-1}\mathcal{F} = \mathcal{F}$. Thus it suffices to show that
$R^qa_{X, *}a_X^{-1}\mathcal{F} = 0$ for $q > 0$.
For this we can use $a_X = \epsilon_X \circ \pi_X$ and
the Leray spectral sequence
(Cohomology on Sites, Lemma \ref{sites-cohomology-lemma-relative-Leray}).
By Lemma \ref{lemma-V-C-all-n-etale-h}
we have $R^i\epsilon_{X, *}(a_X^{-1}\mathcal{F}) = 0$ for $i > 0$
and $\epsilon_{X, *}a_X^{-1}\mathcal{F} = \pi_X^{-1}\mathcal{F}$.
By Lemma \ref{lemma-cohomological-descent-etale} we have
$R^j\pi_{X, *}(\pi_X^{-1}\mathcal{F}) = 0$ for $j > 0$.
This concludes the proof.
\end{proof}

\begin{lemma}
\label{lemma-compare-cohomology-etale-h}
For a scheme $X$ and $a_X : \Sh((\Sch/X)_h) \to \Sh(X_\etale)$
as above:
\begin{enumerate}
\item $H^q(X_\etale, \mathcal{F}) = H^q_h(X, a_X^{-1}\mathcal{F})$
for a torsion abelian sheaf $\mathcal{F}$ on $X_\etale$,
\item $H^q(X_\etale, K) = H^q_h(X, a_X^{-1}K)$
for $K \in D^+(X_\etale)$ with torsion cohomology sheaves.
\end{enumerate}
Example: if $A$ is a torsion abelian group, then
$H^q_\etale(X, \underline{A}) = H^q_h(X, \underline{A})$.
\end{lemma}

\begin{proof}
This follows from Lemma \ref{lemma-cohomological-descent-etale-h}
by Cohomology on Sites, Remark \ref{sites-cohomology-remark-before-Leray}.
\end{proof}











\section{Descending \'etale sheaves}
\label{section-glueing-etale}

\noindent
We prove that \'etale sheaves ``glue'' in the fppf and h topology
and related results. We have already shown the following related
results
\begin{enumerate}
\item Lemma \ref{lemma-describe-pullback} tells us that a sheaf
on the small \'etale site of a scheme $S$ determines a sheaf on the
big \'etale site of $S$ satisfying the sheaf condition for
fpqc coverings (and a fortiori for Zariski, \'etale, smooth, syntomic,
and fppf coverings),
\item Lemma \ref{lemma-describe-pullback-pi-fppf} is a restatement of
the previous point for the fppf topology,
\item Lemma \ref{lemma-describe-pullback-pi-ph} proves the same for the
ph topology,
\item Lemma \ref{lemma-describe-pullback-pi-h} proves the same for the
h topology,
\item Lemma \ref{lemma-descent-sheaf-fppf-etale} is a version of
fppf descent for \'etale sheaves, and
\item Remark \ref{remark-cohomological-descent-finite} tells us that
we have descent of \'etale sheaves for finite surjective morphisms
(we will clarify and strengthen this below).
\end{enumerate}
In the chapter on simplicial spaces we will prove some additional
results on this, see for example
Simplicial Spaces, Sections \ref{spaces-simplicial-section-fppf-hypercovering}
and \ref{spaces-simplicial-section-proper-hypercovering-spaces}.

\medskip\noindent
In order to conveniently express our results we need some notation.
Let $\mathcal{U} = \{f_i : X_i \to X\}$
be a family of morphisms of schemes with fixed target.
A {\it descent datum} for \'etale sheaves with respect to
$\mathcal{U}$ is a family
$((\mathcal{F}_i)_{i \in I}, (\varphi_{ij})_{i, j \in I})$ where
\begin{enumerate}
\item $\mathcal{F}_i$ is in $\Sh(X_{i, \etale})$, and
\item $\varphi_{ij} :
\text{pr}_{0, small}^{-1} \mathcal{F}_i
\longrightarrow
\text{pr}_{1, small}^{-1} \mathcal{F}_j$
is an isomorphism in $\Sh((X_i \times_X X_j)_\etale)$
\end{enumerate}
such that the {\it cocycle condition} holds: the diagrams
$$
\xymatrix{
\text{pr}_{0, small}^{-1}\mathcal{F}_i
\ar[dr]_{\text{pr}_{02, small}^{-1}\varphi_{ik}}
\ar[rr]^{\text{pr}_{01, small}^{-1}\varphi_{ij}} & &
\text{pr}_{1, small}^{-1}\mathcal{F}_j
\ar[dl]^{\text{pr}_{12, small}^{-1}\varphi_{jk}} \\
& \text{pr}_{2, small}^{-1}\mathcal{F}_k
}
$$
commute in $\Sh((X_i \times_X X_j \times_X X_k)_\etale)$.
There is an obvious notion of {\it morphisms of descent data}
and we obtain a category of descent data.
A descent datum
$((\mathcal{F}_i)_{i \in I}, (\varphi_{ij})_{i, j \in I})$
is called {\it effective} if there exist a
$\mathcal{F}$ in $\Sh(X_\etale)$ and isomorphisms
$\varphi_i : f_{i, small}^{-1} \mathcal{F} \to \mathcal{F}_i$
in $\Sh(X_{i, \etale})$ compatible with the $\varphi_{ij}$, i.e.,
such that
$$
\varphi_{ij} =
\text{pr}_{1, small}^{-1} (\varphi_j) \circ
\text{pr}_{0, small}^{-1} (\varphi_i^{-1})
$$
Another way to say this is the following. Given an object $\mathcal{F}$
of $\Sh(X_\etale)$ we obtain the {\it canonical descent datum}
$(f_{i, small}^{-1}\mathcal{F}_i, c_{ij})$ where $c_{ij}$
is the canonical isomorphism
$$
c_{ij} : \text{pr}_{0, small}^{-1} f_{i, small}^{-1}\mathcal{F}
\longrightarrow
\text{pr}_{1, small}^{-1} f_{j, small}^{-1}\mathcal{F}
$$
The descent datum
$((\mathcal{F}_i)_{i \in I}, (\varphi_{ij})_{i, j \in I})$
is effective if and only if it is isomorphic to the canonical
descent datum associated to some $\mathcal{F}$ in $\Sh(X_\etale)$.

\medskip\noindent
If the family consists of a single morphism $\{X \to Y\}$,
then we think of a descent datum as a pair $(\mathcal{F}, \varphi)$
where $\mathcal{F}$ is an object of $\Sh(X_\etale)$ and
$\varphi$ is an isomorphism
$$
\text{pr}_{0, small}^{-1} \mathcal{F}
\longrightarrow
\text{pr}_{1, small}^{-1} \mathcal{F}
$$
in $\Sh((X \times_Y X)_\etale)$ such that the cocycle condition holds:
$$
\xymatrix{
\text{pr}_{0, small}^{-1}\mathcal{F}
\ar[dr]_{\text{pr}_{02, small}^{-1}\varphi}
\ar[rr]^{\text{pr}_{01, small}^{-1}\varphi} & &
\text{pr}_{1, small}^{-1}\mathcal{F}
\ar[dl]^{\text{pr}_{12, small}^{-1}\varphi} \\
& \text{pr}_{2, small}^{-1}\mathcal{F}
}
$$
commutes in $\Sh((X \times_Y X \times_Y X)_\etale)$.
There is a notion of morphisms of descent data and effectivity
exactly as before.

\medskip\noindent
We first prove effective descent for surjective integral morphisms.

\begin{lemma}
\label{lemma-glue-etale-sheaf-section}
Let $f : X \to Y$ be a morphism of schemes which has a section. Then the
functor
$$
\Sh(Y_\etale)
\longrightarrow
\text{descent data for \'etale sheaves wrt }\{X \to Y\}
$$
sending $\mathcal{G}$ in $\Sh(Y_\etale)$ to the canonical descent datum
is an equivalence of categories.
\end{lemma}

\begin{proof}
This is formal and depends only on functoriality of the pullback
functors. We omit the details. Hint: If $s : Y \to X$ is a section,
then a quasi-inverse is the functor sending $(\mathcal{F}, \varphi)$
to $s_{small}^{-1}\mathcal{F}$.
\end{proof}

\begin{lemma}
\label{lemma-glue-etale-sheaf-integral-surjective}
Let $f : X \to Y$ be a surjective integral morphism of schemes.
The functor
$$
\Sh(Y_\etale)
\longrightarrow
\text{descent data for \'etale sheaves wrt }\{X \to Y\}
$$
is an equivalence of categories.
\end{lemma}

\begin{proof}
In this proof we drop the subscript ${}_{small}$ from our pullback
and pushforward functors.
Denote $X_1 = X \times_Y X$ and denote $f_1 : X_1 \to Y$ the
morphism $f \circ \text{pr}_0 = f \circ \text{pr}_1$.
Let $(\mathcal{F}, \varphi)$ be a descent datum for $\{X \to Y\}$.
Let us set $\mathcal{F}_1 = \text{pr}_0^{-1}\mathcal{F}$.
We may think of $\varphi$ as defining an isomorphism
$\mathcal{F}_1 \to \text{pr}_1^{-1}\mathcal{F}$.
We claim that the rule which sends a descent datum
$(\mathcal{F}, \varphi)$
to the sheaf
$$
\mathcal{G} =
\text{Equalizer}\left(
\xymatrix{
f_*\mathcal{F}
\ar@<1ex>[r] \ar@<-1ex>[r] &
f_{1, *}\mathcal{F}_1
}
\right)
$$
is a quasi-inverse to the functor in the statement of the lemma.
The first of the two arrows comes from the map
$$
f_*\mathcal{F} \to
f_*\text{pr}_{0, *}\text{pr}_0^{-1}\mathcal{F} =
f_{1, *}\mathcal{F}_1
$$
and the second arrow comes from the map
$$
f_*\mathcal{F} \to
f_* \text{pr}_{1, *}\text{pr}_1^{-1}\mathcal{F}
\xleftarrow{\varphi}
f_* \text{pr}_{0, *} \text{pr}_0^{-1}\mathcal{F} =
f_{1, *}\mathcal{F}_1
$$
where the arrow pointing left is invertible.
To prove this works we have to show
that the canonical map $f^{-1}\mathcal{G} \to \mathcal{F}$
is an isomorphism; details omitted. In order to prove this it
suffices to check after pulling back by any collection of morphisms
$\Spec(k) \to Y$ where $k$ is an algebraically closed field.
Namely, the corresponding base changes $X_k \to X$ are jointly
surjective and we can check whether a map of sheaves on $X_\etale$
is an isomorphism by looking at stalks on geometric points, see
Theorem \ref{theorem-exactness-stalks}.
By Lemma \ref{lemma-integral-pushforward-commutes-with-base-change}
the construction of $\mathcal{G}$ from the descent datum
$(\mathcal{F}, \varphi)$ commutes with any base change.
Thus we may assume $Y$ is the spectrum of an algebraically
closed field (note that base change preserves the properties
of the morphism $f$, see
Morphisms, Lemma \ref{morphisms-lemma-base-change-surjective}
and \ref{morphisms-lemma-base-change-finite}).
In this case the morphism $X \to Y$ has a section, so we
know that the functor is an equivalence by
Lemma \ref{lemma-glue-etale-sheaf-section}.
However, the reader may show that the functor is an equivalence
if and only if the construction above is a quasi-inverse;
details omitted. This finishes the proof.
\end{proof}

\begin{lemma}
\label{lemma-glue-etale-sheaf-proper-surjective}
Let $f : X \to Y$ be a surjective proper morphism of schemes.
The functor
$$
\Sh(Y_\etale)
\longrightarrow
\text{descent data for \'etale sheaves wrt }\{X \to Y\}
$$
is an equivalence of categories.
\end{lemma}

\begin{proof}
The exact same proof as given in
Lemma \ref{lemma-glue-etale-sheaf-integral-surjective}
works, except the appeal to
Lemma \ref{lemma-integral-pushforward-commutes-with-base-change}
should be replaced by an appeal to
Lemma \ref{lemma-proper-base-change-f-star}.
\end{proof}

\begin{lemma}
\label{lemma-glue-etale-sheaf-check-after-base-change}
Let $f : X \to Y$ be a morphism of schemes. Let $Z \to Y$ be a surjective
integral morphism of schemes or a surjective proper morphism of schemes.
If the functors
$$
\Sh(Z_\etale)
\longrightarrow
\text{descent data for \'etale sheaves wrt }\{X \times_Y Z \to Z\}
$$
and
$$
\Sh((Z \times_Y Z)_\etale)
\longrightarrow
\text{descent data for \'etale sheaves wrt }
\{X \times_Y (Z \times_Y Z) \to Z \times_Y Z\}
$$
are equivalences of categories, then
$$
\Sh(Y_\etale)
\longrightarrow
\text{descent data for \'etale sheaves wrt }\{X \to Y\}
$$
is an equivalence.
\end{lemma}

\begin{proof}
Formal consequence of the definitions and
Lemmas \ref{lemma-glue-etale-sheaf-integral-surjective} and
\ref{lemma-glue-etale-sheaf-proper-surjective}. Details omitted.
\end{proof}

\begin{lemma}
\label{lemma-glue-etale-sheaf-fppf-cover}
Let $f : X \to Y$ be a morphism of schemes which is
surjective, flat, locally of finite presentation.
The functor
$$
\Sh(Y_\etale)
\longrightarrow
\text{descent data for \'etale sheaves wrt }\{X \to Y\}
$$
is an equivalence of categories.
\end{lemma}

\begin{proof}
Exactly as in the proof of
Lemma \ref{lemma-glue-etale-sheaf-integral-surjective}
we claim a quasi-inverse is given by the functor sending
$(\mathcal{F}, \varphi)$ to
$$
\mathcal{G} =
\text{Equalizer}\left(
\xymatrix{
f_*\mathcal{F}
\ar@<1ex>[r] \ar@<-1ex>[r] &
f_{1, *}\mathcal{F}_1
}
\right)
$$
and in order to prove this it suffices to show that
$f^{-1}\mathcal{G} \to \mathcal{F}$ is an isomorphism.
This we may check locally, hence we may and do assume $Y$
is affine. Then we can find a finite surjective morphism
$Z \to Y$ such that there exists an open covering
$Z = \bigcup W_i$ such that $W_i \to Y$ factors through $X$.
See
More on Morphisms, Lemma \ref{more-morphisms-lemma-dominate-fppf-finite}.
Applying Lemma
\ref{lemma-glue-etale-sheaf-check-after-base-change}
we see that it suffices to prove
the lemma after replacing $Y$ by $Z$ and $Z \times_Y Z$ and $f$
by its base change. Thus we may assume $f$ has sections Zariski locally.
Of course, using that the problem is local on $Y$ we reduce
to the case where we have a section which is
Lemma \ref{lemma-glue-etale-sheaf-section}.
\end{proof}

\begin{lemma}
\label{lemma-glue-etale-sheaf-fppf}
Let $\{f_i : X_i \to X\}$ be an fppf covering of schemes.
The functor
$$
\Sh(X_\etale)
\longrightarrow
\text{descent data for \'etale sheaves wrt }\{f_i : X_i \to X\}
$$
is an equivalence of categories.
\end{lemma}

\begin{proof}
We have Lemma \ref{lemma-glue-etale-sheaf-fppf-cover}
for the morphism $f : \coprod X_i \to X$.
Then a formal argument shows that descent data for $f$
are the same thing as descent data for the covering, compare
with Descent, Lemma \ref{descent-lemma-family-is-one}.
Details omitted.
\end{proof}

\begin{lemma}
\label{lemma-glue-etale-sheaf-modification}
Let $f : X' \to X$ be a proper morphism of schemes. Let $i : Z \to X$
be a closed immersion. Set $E = Z \times_X X'$. Picture
$$
\xymatrix{
E \ar[d]_g \ar[r]_j & X' \ar[d]^f \\
Z \ar[r]^i & X
}
$$
If $f$ is an isomorphism over $X \setminus Z$, then the functor
$$
\Sh(X_\etale)
\longrightarrow
\Sh(X'_\etale) \times_{\Sh(E_\etale)} \Sh(Z_\etale)
$$
is an equivalence of categories.
\end{lemma}

\begin{proof}
We will work with the $2$-fibre product category as constructed in
Categories, Example \ref{categories-example-2-fibre-product-categories}.
The functor sends $\mathcal{F}$ to the triple
$(f^{-1}\mathcal{F}, i^{-1}\mathcal{F}, c)$ where
$c : j^{-1}f^{-1}\mathcal{F} \to g^{-1}i^{-1}\mathcal{F}$
is the canonical isomorphism. We will construct a quasi-inverse functor. Let
$(\mathcal{F}', \mathcal{G}, \alpha)$ be an object
of the right hand side of the arrow.
We obtain an isomorphism
$$
i^{-1}f_*\mathcal{F}' = g_*j^{-1}\mathcal{F}'
\xrightarrow{g_*\alpha}
g_*g^{-1}\mathcal{G}
$$
The first equality is Lemma \ref{lemma-proper-base-change-f-star}.
Using this we obtain maps
$i_*\mathcal{G} \to i_*g_*g^{-1}\mathcal{G}$
and $f'_*\mathcal{F}' \to i_*g_*g^{-1}\mathcal{G}$. We set
$$
\mathcal{F} = f_*\mathcal{F}' \times_{i_*g_*g^{-1}\mathcal{G}} i_*\mathcal{G}
$$
and we claim that $\mathcal{F}$ is an object of the left hand side
of the arrow whose image in the right hand side is isomorphic to
the triple we started out with. Let us compute the stalk of $\mathcal{F}$
at a geometric point $\overline{x}$ of $X$.

\medskip\noindent
If $\overline{x}$ is not
in $Z$, then on the one hand $\overline{x}$ comes from a unique
geometric point $\overline{x}'$ of $X'$ and
$\mathcal{F}'_{\overline{x}'} = (f_*\mathcal{F}')_{\overline{x}}$
and on the other hand we have $(i_*\mathcal{G})_{\overline{x}}$
and $(i_*g_*g^{-1}\mathcal{G})_{\overline{x}}$ are singletons.
Hence we see that $\mathcal{F}_{\overline{x}}$ equals
$\mathcal{F}'_{\overline{x}'}$.

\medskip\noindent
If $\overline{x}$ is in $Z$, i.e., $\overline{x}$ is the image of
a geometric point $\overline{z}$ of $Z$, then we obtain
$(i_*\mathcal{G})_{\overline{x}} = \mathcal{G}_{\overline{z}}$
and
$$
(i_*g_*g^{-1}\mathcal{G})_{\overline{x}} =
(g_*g^{-1}\mathcal{G})_{\overline{z}} =
\Gamma(E_{\overline{z}}, g^{-1}\mathcal{G}|_{E_{\overline{z}}})
$$
(by the proper base change for pushforward used above)
and similarly
$$
(f_*\mathcal{F}')_{\overline{x}} =
\Gamma(X'_{\overline{x}}, \mathcal{F}'|_{X'_{\overline{x}}})
$$
Since we have the identification
$E_{\overline{z}} = X'_{\overline{x}}$ and since $\alpha$
defines an isomorphism between the sheaves
$\mathcal{F}'|_{X'_{\overline{x}}}$ and
$g^{-1}\mathcal{G}|_{E_{\overline{z}}}$
we conclude that we get
$$
\mathcal{F}_{\overline{x}} = \mathcal{G}_{\overline{z}}
$$
in this case.

\medskip\noindent
To finish the proof, we observe that there are canonical maps
$i^{-1}\mathcal{F} \to \mathcal{G}$ and $f^{-1}\mathcal{F} \to \mathcal{F}'$
compatible with $\alpha$ which on stalks produce the isomorphisms
we saw above. We omit the careful construction of these maps.
\end{proof}

\begin{lemma}
\label{lemma-h-descent-etale-sheaves}
Let $S$ be a scheme. Then the category fibred in groupoids
$$
p : \mathcal{S} \longrightarrow (\Sch/S)_h
$$
whose fibre category over $U$ is the category $\Sh(U_\etale)$
of sheaves on the small \'etale site of $U$ is a stack in groupoids.
\end{lemma}

\begin{proof}
To prove the lemma we will check conditions (1), (2), and (3) of
More on Flatness, Lemma \ref{flat-lemma-refine-check-h-stack}.

\medskip\noindent
Condition (1) holds because we have glueing for sheaves (and
Zariski coverings are \'etale coverings). See
Sites, Lemma \ref{sites-lemma-glue-sheaves}.

\medskip\noindent
To see condition (2), suppose that $f : X \to Y$ is a surjective,
flat, proper morphism of finite presentation over $S$ with $Y$ affine.
Then we have descent for $\{X \to Y\}$ by either
Lemma \ref{lemma-glue-etale-sheaf-fppf-cover} or
Lemma \ref{lemma-glue-etale-sheaf-proper-surjective}.

\medskip\noindent
Condition (3) follows immediately from the more general
Lemma \ref{lemma-glue-etale-sheaf-modification}.
\end{proof}



















\section{Blow up squares and \'etale cohomology}
\label{section-blow-up-square}

\noindent
Blow up squares are introduced in
More on Flatness, Section \ref{flat-section-blow-up-ph}.
Using the proper base change theorem we can see that
we have a Mayer-Vietoris type result for blow up squares.

\begin{lemma}
\label{lemma-blow-up-square-cohomological-descent}
Let $X$ be a scheme and let $Z \subset X$ be a closed subscheme
cut out by a quasi-coherent ideal of finite type. Consider the
corresponding blow up square
$$
\xymatrix{
E \ar[d]_\pi \ar[r]_j & X' \ar[d]^b \\
Z \ar[r]^i & X
}
$$
For $K \in D^+(X_\etale)$ with torsion cohomology sheaves
we have a distinguished triangle
$$
K \to Ri_*(K|_Z) \oplus Rb_*(K|_{X'}) \to Rc_*(K|_E) \to K[1]
$$
in $D(X_\etale)$ where $c = i \circ \pi = b \circ j$.
\end{lemma}

\begin{proof}
The notation $K|_{X'}$ stands for $b_{small}^{-1}K$.
Choose a bounded below complex $\mathcal{F}^\bullet$
of abelian sheaves representing $K$. Observe that
$i_*(\mathcal{F}^\bullet|_Z)$ represents $Ri_*(K|_Z)$
because $i_*$ is exact
(Proposition \ref{proposition-finite-higher-direct-image-zero}).
Choose a quasi-isomorphism
$b_{small}^{-1}\mathcal{F}^\bullet \to \mathcal{I}^\bullet$
where $\mathcal{I}^\bullet$ is a bounded below complex of injective
abelian sheaves on $X'_\etale$. This map is adjoint to a map
$\mathcal{F}^\bullet \to b_*(\mathcal{I}^\bullet)$ and
$b_*(\mathcal{I}^\bullet)$ represents $Rb_*(K|_{X'})$.
We have $\pi_*(\mathcal{I}^\bullet|_E) = (b_*\mathcal{I}^\bullet)|_Z$
by Lemma \ref{lemma-proper-base-change-f-star} and by
Lemma \ref{lemma-proper-base-change} this complex represents
$R\pi_*(K|_E)$. Hence the map
$$
Ri_*(K|_Z) \oplus Rb_*(K|_{X'}) \to Rc_*(K|_E)
$$
is represented by the surjective map of bounded below complexes
$$
i_*(\mathcal{F}^\bullet|_Z) \oplus
b_*(\mathcal{I}^\bullet)
\to
i_*\left(b_*(\mathcal{I}^\bullet)|_Z\right)
$$
To get our distinguished triangle it suffices to show that
the canonical map
$\mathcal{F}^\bullet \to i_*(\mathcal{F}^\bullet|_Z) \oplus
b_*(\mathcal{I}^\bullet)$
maps quasi-isomorphically onto the kernel of the map
of complexes displayed above (namely a short exact sequence
of complexes determines a distinguished triangle in the derived
category, see
Derived Categories, Section \ref{derived-section-canonical-delta-functor}).
We may check this on stalks at a geometric point $\overline{x}$ of $X$.
If $\overline{x}$ is not in $Z$, then $X' \to X$ is an isomorphism
over an open neighbourhood of $\overline{x}$. Thus, if $\overline{x}'$
denotes the corresponding geometric point of $X'$ in this case, then
we have to show that
$$
\mathcal{F}^\bullet_{\overline{x}} \to \mathcal{I}^\bullet_{\overline{x}'}
$$
is a quasi-isomorphism. This is true by our choice of $\mathcal{I}^\bullet$.
If $\overline{x}$ is in $Z$, then
$b_(\mathcal{I}^\bullet)_{\overline{x}} \to 
i_*\left(b_*(\mathcal{I}^\bullet)|_Z\right)_{\overline{x}}$
is an isomorphism of complexes of abelian groups. Hence the
kernel is equal to
$i_*(\mathcal{F}^\bullet|_Z)_{\overline{x}} =
\mathcal{F}^\bullet_{\overline{x}}$ as desired.
\end{proof}

\begin{lemma}
\label{lemma-blow-up-square-etale-cohomology}
Let $X$ be a scheme and let $K \in D^+(X_\etale)$ have
torsion cohomology sheaves. Let $Z \subset X$ be a closed subscheme
cut out by a quasi-coherent ideal of finite type. Consider the
corresponding blow up square
$$
\xymatrix{
E \ar[d] \ar[r] & X' \ar[d]^b \\
Z \ar[r] & X
}
$$
Then there is a canonical long exact sequence
$$
H^p_\etale(X, K) \to
H^p_\etale(X', K|_{X'}) \oplus
H^p_\etale(Z, K|_Z) \to
H^p_\etale(E, K|_E) \to
H^{p + 1}_\etale(X, K)
$$
\end{lemma}

\begin{proof}[First proof]
This follows immediately from
Lemma \ref{lemma-blow-up-square-cohomological-descent}
and the fact that
$$
R\Gamma(X, Rb_*(K|_{X'})) = R\Gamma(X', K|_{X'})
$$
(see Cohomology on Sites, Section \ref{sites-cohomology-section-leray})
and similarly for the others.
\end{proof}

\begin{proof}[Second proof]
By Lemma \ref{lemma-compare-cohomology-etale-ph}
these cohomology groups are the cohomology of
$X, X', E, Z$ with values in some complex of abelian sheaves
on the site $(\Sch/X)_{ph}$. (Namely, the object
$a_X^{-1}K$ of the derived category, see
Lemma \ref{lemma-describe-pullback-pi-ph} above
and recall that $K|_{X'} = b_{small}^{-1}K$.)
By More on Flatness, Lemma \ref{flat-lemma-blow-up-square-ph}
the ph sheafification of the diagram of representable
presheaves is cocartesian. Thus the lemma follows from the very general
Cohomology on Sites, Lemma \ref{sites-cohomology-lemma-square-triangle}
applied to the site $(\Sch/X)_{ph}$ and the commutative diagram
of the lemma.
\end{proof}

\begin{lemma}
\label{lemma-blow-up-square-equivalence}
Let $X$ be a scheme and let $Z \subset X$ be a closed subscheme
cut out by a quasi-coherent ideal of finite type. Consider the
corresponding blow up square
$$
\xymatrix{
E \ar[d]_\pi \ar[r]_j & X' \ar[d]^b \\
Z \ar[r]^i & X
}
$$
Suppose given
\begin{enumerate}
\item an object $K'$ of $D^+(X'_\etale)$ with torsion cohomology sheaves,
\item an object $L$ of $D^+(Z_\etale)$ with torsion cohomology sheaves, and
\item an isomorphism $\gamma : K'|_E \to L|_E$.
\end{enumerate}
Then there exists an object $K$ of $D^+(X_\etale)$
and isomorphisms $f : K|_{X'} \to K'$, $g : K|_Z \to L$ such
that $\gamma = g|_E \circ f^{-1}|_E$.
Moreover, given
\begin{enumerate}
\item an object $M$ of $D^+(X_\etale)$ with torsion cohomology sheaves,
\item a morphism $\alpha : K' \to M|_{X'}$ of $D(X'_\etale)$,
\item a morphism $\beta : L \to M|_Z$ of $D(Z_\etale)$,
\end{enumerate}
such that
$$
\alpha|_E  = \beta|_E \circ \gamma.
$$
Then there exists a morphism $M \to K$ in $D(X_\etale)$
whose restriction to $X'$ is $a \circ f$
and whose restriction to $Z$ is $b \circ g$.
\end{lemma}

\begin{proof}
If $K$ exists, then Lemma \ref{lemma-blow-up-square-cohomological-descent}
tells us a distinguished triangle that it fits in. Thus we simply choose
a distinguished triangle
$$
K \to Ri_*(L) \oplus Rb_*(K') \to Rc_*(L|_E) \to K[1]
$$
where $c = i \circ \pi = b \circ j$. Here the map $Ri_*(L) \to Rc_*(L|_E)$
is $Ri_*$ applied to the adjunction mapping $E \to R\pi_*(L|_E)$.
The map $Rb_*(K') \to Rc_*(L|_E)$ is the composition of the canonical map
$Rb_*(K') \to Rc_*(K'|_E)) = R$ and $Rc_*(\gamma)$.
The maps $g$ and $f$ of the statement of the lemma are the adjoints
of these maps. If we restrict this distinguished triangle to $Z$
then the map $Rb_*(K) \to Rc_*(L|_E)$ becomes an isomorphism
by the base change theorem (Lemma \ref{lemma-proper-base-change}) and hence
the map $g : K|_Z \to L$ is an isomorphism.
Looking at the distinguished triangle we see that $f : K|_{X'} \to K'$
is an isomorphism over $X' \setminus E = X \setminus Z$.
Moreover, we have $\gamma \circ f|_E = g|_E$ by construction.
Then since $\gamma$ and $g$ are isomorphisms we conclude
that $f$ induces isomorphisms on stalks at geometric points of $E$
as well. Thus $f$ is an isomorphism.

\medskip\noindent
For the final statement, we may replace $K'$ by $K|_{X'}$,
$L$ by $K|_Z$, and $\gamma$ by the canonical identification.
Observe that $\alpha$ and $\beta$ induce a commutative square
$$
\xymatrix{
K \ar[r] \ar@{..>}[d] &
Ri_*(K|_Z) \oplus Rb_*(K|_{X'}) \ar[r] \ar[d]_{\beta \oplus \alpha} &
Rc_*(K|_E) \ar[r] \ar[d]_{\alpha|_E} &
K[1] \ar@{..>}[d] \\
M \ar[r] &
Ri_*(M|_Z) \oplus Rb_*(M|_{X'}) \ar[r] &
Rc_*(M|_E) \ar[r] &
M[1]
}
$$
Thus by the axioms of a derived category we get a dotted
arrow producing a morphism of distinguished triangles.
\end{proof}














\section{Almost blow up squares and the h topology}
\label{section-blow-up-h}

\noindent
In this section we continue the discussion in
More on Flatness, Section \ref{flat-section-blow-up-h}.
For the convenience of the reader we recall that an
almost blow up square is a commutative diagram
\begin{equation}
\label{equation-almost-blow-up-square}
\vcenter{
\xymatrix{
E \ar[d] \ar[r] & X' \ar[d]^b \\
Z \ar[r] & X
}
}
\end{equation}
of schemes satisfying the following conditions:
\begin{enumerate}
\item $Z \to X$ is a closed immersion of finite presentation,
\item $E = b^{-1}(Z)$ is a locally principal closed subscheme of $X'$,
\item $b$ is proper and of finite presentation,
\item the closed subscheme $X'' \subset X'$ cut out by the quasi-coherent
ideal of sections of $\mathcal{O}_{X'}$ supported on $E$
(Properties, Lemma \ref{properties-lemma-sections-supported-on-closed-subset})
is the blow up of $X$ in $Z$.
\end{enumerate}
It follows that the morphism $b$ induces an isomorphism
$X' \setminus E \to X \setminus Z$.

\medskip\noindent
We are going to give a criterion for ``h sheafiness'' for
objects in the derived category of the big fppf site
$(\Sch/S)_{fppf}$ of a scheme $S$. On the same underlying category
we have a second topology, namely the h topology
(More on Flatness, Section \ref{flat-section-h}).
Recall that fppf coverings are h coverings
(More on Flatness, Lemma \ref{flat-lemma-zariski-h}). Hence we may
consider the morphism
$$
\epsilon : (\Sch/S)_h \longrightarrow (\Sch/S)_{fppf}
$$
See Cohomology on Sites, Section \ref{sites-cohomology-section-compare}.
In particular, we have a fully faithful functor
$$
R\epsilon_* : D((\Sch/S)_h) \longrightarrow D((\Sch/S)_{fppf})
$$
and we can ask: what is the essential image of this functor?

\begin{lemma}
\label{lemma-blow-up-square-h}
With notation as above, if $K$ is in the essential image
of $R\epsilon_*$, then the maps $c^K_{X, Z, X', E}$ of
Cohomology on Sites, Lemma \ref{sites-cohomology-lemma-c-square}
are quasi-isomorphisms.
\end{lemma}

\begin{proof}
Denote ${}^\#$ sheafification in the h topology.
We have seen in More on Flatness, Lemma \ref{flat-lemma-blow-up-square-h}
that $h_X^\# = h_Z^\# \amalg_{h_E^\#} h_{X'}^\#$. On the other hand,
the map $h_E^\# \to h_{X'}^\#$ is injective as $E \to X'$ is a
monomorphism. Thus this lemma is a special case of
Cohomology on Sites, Lemma \ref{sites-cohomology-lemma-descent-squares-helper}
(which itself is a formal consequence of
Cohomology on Sites, Lemma \ref{sites-cohomology-lemma-square-triangle}).
\end{proof}

\begin{proposition}
\label{proposition-check-h}
Let $K$ be an object of $D^+((\Sch/S)_{fppf})$.
Then $K$ is in the essential image of
$R\epsilon_* : D((\Sch/S)_h) \to D((\Sch/S)_{fppf})$
if and only if $c^K_{X, X', Z, E}$ is a quasi-isomorphism
for every almost blow up square (\ref{equation-almost-blow-up-square})
in $(\Sch/S)_h$ with $X$ affine.
\end{proposition}

\begin{proof}
We prove this by applying
Cohomology on Sites, Lemma \ref{sites-cohomology-lemma-descent-squares}
whose hypotheses hold by
Lemma \ref{lemma-blow-up-square-h} and
More on Flatness, Proposition \ref{flat-proposition-check-h}.
\end{proof}

\begin{lemma}
\label{lemma-refine-check-h}
Let $K$ be an object of $D^+((\Sch/S)_{fppf})$. Then $K$ is in the
essential image of $R\epsilon_* : D((\Sch/S)_h) \to D((\Sch/S)_{fppf})$
if and only if $c^K_{X, X', Z, E}$ is a quasi-isomorphism
for every almost blow up square as in
More on Flatness, Examples \ref{flat-example-one-generator} and
\ref{flat-example-two-generators}.
\end{lemma}

\begin{proof}
We prove this by applying
Cohomology on Sites, Lemma \ref{sites-cohomology-lemma-descent-squares}
whose hypotheses hold by
Lemma \ref{lemma-blow-up-square-h} and
More on Flatness, Lemma \ref{flat-lemma-refine-check-h}
\end{proof}






\section{Cohomology of the structure sheaf in the h topology}
\label{section-cohomology-O-h}

\noindent
Let $p$ be a prime number. Let $(\mathcal{C}, \mathcal{O})$ be a ringed site
with $p\mathcal{O} = 0$. Then we set $\colim_F \mathcal{O}$
equal to the colimit in the category of sheaves of rings of the system
$$
\mathcal{O} \xrightarrow{F} \mathcal{O} \xrightarrow{F}
\mathcal{O} \xrightarrow{F} \ldots
$$
where $F : \mathcal{O} \to \mathcal{O}$, $f \mapsto f^p$
is the Frobenius endomorphism.

\begin{lemma}
\label{lemma-h-sheaf-colim-F}
Let $p$ be a prime number. Let $S$ be a scheme over $\mathbf{F}_p$.
Consider the sheaf $\mathcal{O}^{perf} = \colim_F \mathcal{O}$
on $(\Sch/S)_{fppf}$. Then $\mathcal{O}^{perf}$ is in the essential
image of $R\epsilon_* : D((\Sch/S)_h) \to D((\Sch/S)_{fppf})$.
\end{lemma}

\begin{proof}
We prove this using the criterion of Lemma \ref{lemma-refine-check-h}.
Before check the conditions, we note that for a
quasi-compact and quasi-separated object $X$ of
$(\Sch/S)_{fppf}$ we have
$$
H^i_{fppf}(X, \mathcal{O}^{perf}) = \colim_F H^i_{fppf}(X, \mathcal{O})
$$
See Cohomology on Sites,
Lemma \ref{sites-cohomology-lemma-colim-works-over-collection}.
We will also use that $H^i_{fppf}(X, \mathcal{O}) = H^i(X, \mathcal{O})$, see
Descent, Proposition \ref{descent-proposition-same-cohomology-quasi-coherent}.

\medskip\noindent
Let $A, f, J$ be as in
More on Flatness, Example \ref{flat-example-one-generator}
and consider the associated almost blow up square.
Since $X$, $X'$, $Z$, $E$ are affine, we have no
higher cohomology of $\mathcal{O}$. Hence we only
have to check that
$$
0 \to
\mathcal{O}^{perf}(X) \to
\mathcal{O}^{perf}(X') \oplus \mathcal{O}^{perf}(Z) \to
\mathcal{O}^{perf}(E) \to 0
$$
is a short exact sequence. This was shown in (the proof of)
More on Flatness, Lemma \ref{flat-lemma-h-sheaf-colim-F}.

\medskip\noindent
Let $X, X', Z, E$ be as in
More on Flatness, Example \ref{flat-example-two-generators}.
Since $X$ and $Z$ are affine we have
$H^p(X, \mathcal{O}_X) = H^p(Z, \mathcal{O}_X) = 0$ for $p > 0$.
By More on Flatness, Lemma \ref{flat-lemma-funny-blow-up}
we have $H^p(X', \mathcal{O}_{X'}) = 0$ for $p > 0$.
Since $E = \mathbf{P}^1_Z$ and $Z$ is affine we also have
$H^p(E, \mathcal{O}_E) = 0$ for $p > 0$. As in the previous
paragraph we reduce to checking that
$$
0 \to
\mathcal{O}^{perf}(X) \to
\mathcal{O}^{perf}(X') \oplus \mathcal{O}^{perf}(Z) \to
\mathcal{O}^{perf}(E) \to 0
$$
is a short exact sequence. This was shown in (the proof of)
More on Flatness, Lemma \ref{flat-lemma-h-sheaf-colim-F}.
\end{proof}

\begin{proposition}
\label{proposition-h-cohomology-structure-sheaf}
Let $p$ be a prime number. Let $S$ be a quasi-compact and quasi-separated
scheme over $\mathbf{F}_p$. Then
$$
H^i((\Sch/S)_h, \mathcal{O}^h) =
\colim_F H^i(S, \mathcal{O})
$$
Here on the left hand side by $\mathcal{O}^h$ we mean
the h sheafification of the structure sheaf.
\end{proposition}

\begin{proof}
This is just a reformulation of Lemma \ref{lemma-h-sheaf-colim-F}.
Recall that
$\mathcal{O}^h = \mathcal{O}^{perf} = \colim_F \mathcal{O}$, see
More on Flatness, Lemma \ref{flat-lemma-char-p}.
By Lemma \ref{lemma-h-sheaf-colim-F} we see that
$\mathcal{O}^{perf}$ viewed as an object of $D((\Sch/S)_{fppf})$
is of the form $R\epsilon_*K$ for some $K \in D((\Sch/S)_h)$.
Then $K = \epsilon^{-1}\mathcal{O}^{perf}$ which is actually
equal to $\mathcal{O}^{perf}$ because $\mathcal{O}^{perf}$ is an h sheaf. See
Cohomology on Sites, Section \ref{sites-cohomology-section-compare}.
Hence $R\epsilon_*\mathcal{O}^{perf} = \mathcal{O}^{perf}$
(with apologies for the confusing notation).
Thus the lemma now follows from Leray
$$
R\Gamma_h(S, \mathcal{O}^{perf}) =
R\Gamma_{fppf}(S, R\epsilon_*\mathcal{O}^{perf}) =
R\Gamma_{fppf}(S, \mathcal{O}^{perf})
$$
and the fact that
$$
H^i_{fppf}(S, \mathcal{O}^{perf}) =
H^i_{fppf}(S, \colim_F \mathcal{O}) =
\colim_F H^i_{fppf}(S, \mathcal{O})
$$
as $S$ is quasi-compact and quasi-separated
(see proof of Lemma \ref{lemma-h-sheaf-colim-F}).
\end{proof}












\begin{multicols}{2}[\section{Autres chapitres}]
\noindent
Préliminaires
\begin{enumerate}
\item \hyperref[introduction-section-phantom]{Introduction}
\item \hyperref[conventions-section-phantom]{Conventions}
\item \hyperref[sets-section-phantom]{Théorie des ensembles}
\item \hyperref[categories-section-phantom]{Catégories}
\item \hyperref[topology-section-phantom]{Topologie}
\item \hyperref[sheaves-section-phantom]{Faisceaux sur les espaces}
\item \hyperref[sites-section-phantom]{Sites et faisceaux}
\item \hyperref[stacks-section-phantom]{Champs}
\item \hyperref[fields-section-phantom]{Corps}
\item \hyperref[algebra-section-phantom]{Algèbre commutative}
\item \hyperref[brauer-section-phantom]{Groupes de Brauer}
\item \hyperref[homology-section-phantom]{Algèbre homologique}
\item \hyperref[derived-section-phantom]{Catégories dérivées}
\item \hyperref[simplicial-section-phantom]{Méthodes simpliciales}
\item \hyperref[more-algebra-section-phantom]{Compléments d'algèbre}
\item \hyperref[smoothing-section-phantom]{Lissage des morphismes d'anneaux}
\item \hyperref[modules-section-phantom]{Faisceaux de modules}
\item \hyperref[sites-modules-section-phantom]{Modules sur les sites}
\item \hyperref[injectives-section-phantom]{Objets injectifs}
\item \hyperref[cohomology-section-phantom]{Cohomologie des faisceaux}
\item \hyperref[sites-cohomology-section-phantom]{Cohomologie sur les sites}
\item \hyperref[dga-section-phantom]{Algèbre différentielle graduée}
\item \hyperref[dpa-section-phantom]{Algèbre à puissances divisées}
\item \hyperref[sdga-section-phantom]{Faisceaux différentiels gradués}
\item \hyperref[hypercovering-section-phantom]{Hyperrecouvrements}
\end{enumerate}
Schémas
\begin{enumerate}
\setcounter{enumi}{25}
\item \hyperref[schemes-section-phantom]{Schémas}
\item \hyperref[constructions-section-phantom]{Constructions de schémas}
\item \hyperref[properties-section-phantom]{Propriétés des schémas}
\item \hyperref[morphisms-section-phantom]{Morphismes de schémas}
\item \hyperref[coherent-section-phantom]{Cohomologie des schémas}
\item \hyperref[divisors-section-phantom]{Diviseurs}
\item \hyperref[limits-section-phantom]{Limites de schémas}
\item \hyperref[varieties-section-phantom]{Variétés}
\item \hyperref[topologies-section-phantom]{Topologies sur les schémas}
\item \hyperref[descent-section-phantom]{Descente}
\item \hyperref[perfect-section-phantom]{Catégories dérivées de schémas}
\item \hyperref[more-morphisms-section-phantom]{Compléments sur les morphismes}
\item \hyperref[flat-section-phantom]{Compléments sur la platitude}
\item \hyperref[groupoids-section-phantom]{Schémas en groupoïdes}
\item \hyperref[more-groupoids-section-phantom]{Compléments sur les schémas en groupoïdes}
\item \hyperref[etale-section-phantom]{Morphismes étales de schémas}
\end{enumerate}
Thèmes de théorie des schémas
\begin{enumerate}
\setcounter{enumi}{41}
\item \hyperref[chow-section-phantom]{Homologie de Chow}
\item \hyperref[intersection-section-phantom]{Théorie de l'intersection}
\item \hyperref[pic-section-phantom]{Schémas de Picard des courbes}
\item \hyperref[weil-section-phantom]{Théories cohomologiques de Weil}
\item \hyperref[adequate-section-phantom]{Modules adéquats}
\item \hyperref[dualizing-section-phantom]{Complexes dualisants}
\item \hyperref[duality-section-phantom]{Dualité pour les schémas}
\item \hyperref[discriminant-section-phantom]{Discriminants et différentes}
\item \hyperref[derham-section-phantom]{Cohomologie de de Rham}
\item \hyperref[local-cohomology-section-phantom]{Cohomologie locale}
\item \hyperref[algebraization-section-phantom]{Géométrie algébrique et formelle}
\item \hyperref[curves-section-phantom]{Courbes algébriques}
\item \hyperref[resolve-section-phantom]{Résolution des surfaces}
\item \hyperref[models-section-phantom]{Réduction semi-stable}
\item \hyperref[functors-section-phantom]{Foncteurs et morphismes}
\item \hyperref[equiv-section-phantom]{Catégories dérivées de variétés}
\item \hyperref[pione-section-phantom]{Groupes fondamentaux de schémas}
\item \hyperref[etale-cohomology-section-phantom]{Cohomologie étale}
\item \hyperref[crystalline-section-phantom]{Cohomologie cristalline}
\item \hyperref[proetale-section-phantom]{Cohomologie pro-étale}
\item \hyperref[relative-cycles-section-phantom]{Cycles relatifs}
\item \hyperref[more-etale-section-phantom]{Compléments de cohomologie étale}
\item \hyperref[trace-section-phantom]{La formule des traces}
\end{enumerate}
Espaces algébriques
\begin{enumerate}
\setcounter{enumi}{64}
\item \hyperref[spaces-section-phantom]{Espaces algébriques}
\item \hyperref[spaces-properties-section-phantom]{Propriétés des espaces algébriques}
\item \hyperref[spaces-morphisms-section-phantom]{Morphismes d'espaces algébriques}
\item \hyperref[decent-spaces-section-phantom]{Espaces algébriques décents}
\item \hyperref[spaces-cohomology-section-phantom]{Cohomologie des espaces algébriques}
\item \hyperref[spaces-limits-section-phantom]{Limites d'espaces algébriques}
\item \hyperref[spaces-divisors-section-phantom]{Diviseurs sur les espaces algébriques}
\item \hyperref[spaces-over-fields-section-phantom]{Espaces algébriques sur des corps}
\item \hyperref[spaces-topologies-section-phantom]{Topologies sur les espaces algébriques}
\item \hyperref[spaces-descent-section-phantom]{Descente et espaces algébriques}
\item \hyperref[spaces-perfect-section-phantom]{Catégories dérivées d'espaces}
\item \hyperref[spaces-more-morphisms-section-phantom]{Compléments sur les morphismes d'espaces}
\item \hyperref[spaces-flat-section-phantom]{Platitude sur les espaces algébriques}
\item \hyperref[spaces-groupoids-section-phantom]{Groupoïdes dans les espaces algébriques}
\item \hyperref[spaces-more-groupoids-section-phantom]{Compléments sur les groupoïdes dans les espaces}
\item \hyperref[bootstrap-section-phantom]{Amorçage}
\item \hyperref[spaces-pushouts-section-phantom]{Sommes amalgamées d'espaces algébriques}
\end{enumerate}
Thèmes de géométrie
\begin{enumerate}
\setcounter{enumi}{81}
\item \hyperref[spaces-chow-section-phantom]{Groupes de Chow des espaces}
\item \hyperref[groupoids-quotients-section-phantom]{Quotients de groupoïdes}
\item \hyperref[spaces-more-cohomology-section-phantom]{Compléments sur la cohomologie des espaces}
\item \hyperref[spaces-simplicial-section-phantom]{Espaces simpliciaux}
\item \hyperref[spaces-duality-section-phantom]{Dualité pour les espaces}
\item \hyperref[formal-spaces-section-phantom]{Espaces algébriques formels}
\item \hyperref[restricted-section-phantom]{Algébrisation des espaces formels}
\item \hyperref[spaces-resolve-section-phantom]{Résolution des surfaces revisitée}
\end{enumerate}
Théorie des déformations
\begin{enumerate}
\setcounter{enumi}{89}
\item \hyperref[formal-defos-section-phantom]{Théorie formelle des déformations}
\item \hyperref[defos-section-phantom]{Théorie des déformations}
\item \hyperref[cotangent-section-phantom]{Le complexe cotangent}
\item \hyperref[examples-defos-section-phantom]{Problèmes de déformation}
\end{enumerate}
Champs algébriques
\begin{enumerate}
\setcounter{enumi}{93}
\item \hyperref[algebraic-section-phantom]{Champs algébriques}
\item \hyperref[examples-stacks-section-phantom]{Exemples de champs}
\item \hyperref[stacks-sheaves-section-phantom]{Faisceaux sur les champs algébriques}
\item \hyperref[criteria-section-phantom]{Critères de représentabilité}
\item \hyperref[artin-section-phantom]{Axiomes d'Artin}
\item \hyperref[quot-section-phantom]{Espaces de Quot et de Hilbert}
\item \hyperref[stacks-properties-section-phantom]{Propriétés des champs algébriques}
\item \hyperref[stacks-morphisms-section-phantom]{Morphismes de champs algébriques}
\item \hyperref[stacks-limits-section-phantom]{Limites de champs algébriques}
\item \hyperref[stacks-cohomology-section-phantom]{Cohomologie des champs algébriques}
\item \hyperref[stacks-perfect-section-phantom]{Catégories dérivées de champs}
\item \hyperref[stacks-introduction-section-phantom]{Introduction aux champs algébriques}
\item \hyperref[stacks-more-morphisms-section-phantom]{Compléments sur les morphismes de champs}
\item \hyperref[stacks-geometry-section-phantom]{La géométrie des champs}
\end{enumerate}
Thèmes de théorie des espaces de modules
\begin{enumerate}
\setcounter{enumi}{107}
\item \hyperref[moduli-section-phantom]{Champs de modules}
\item \hyperref[moduli-curves-section-phantom]{Espaces de modules de courbes}
\end{enumerate}
Divers
\begin{enumerate}
\setcounter{enumi}{109}
\item \hyperref[examples-section-phantom]{Exemples}
\item \hyperref[exercises-section-phantom]{Exercices}
\item \hyperref[guide-section-phantom]{Guide bibliographique}
\item \hyperref[desirables-section-phantom]{Desiderata}
\item \hyperref[coding-section-phantom]{Style de programmation}
\item \hyperref[obsolete-section-phantom]{Obsolète}
\item \hyperref[fdl-section-phantom]{Licence de documentation libre GNU}
\item \hyperref[index-section-phantom]{Index généré automatiquement}
\end{enumerate}
\end{multicols}


\bibliography{my}
\bibliographystyle{amsalpha}

\end{document}
